% ============ ``古典调和分析''中文翻译 LaTeX 重排版 ============
% 原书:藪田公三・佐藤圓治・田中仁・宮地晶彦・中路貴彦
%      ``古典調和解析''(解析学百科 I),朝倉書店,2008
% 模板:AJbook.cls(李文威``代数学方法''书系模板,见 templates 说明)
% 结构:前言 + 第 1-4 章(每章末附参考文献)+ 索引
%
% 编号体系(原书约定,AJbook.cls 默认行为即符合):
%   - 章:第一章,第二章...(中文数字,见 titles-setup.tex)
%   - 节:1.1,1.2...(章.节,book 类默认,每章重置)
%   - 小节:1.2.1,1.2.2...(章.节.小节)
%   - 定理/定义/命题/引理/推论/例:共用一个计数器,按节编号(定理 1.1.4,定义 1.3.1,例 1.1.10)
%   - 公式:按节编号(式 (1.2.7))
%   - 注/注意:不编号(用 remarknonum 环境)
\documentclass[geometry=a4, fontsetup=font-setup-open.tex, titlesetup=titles-setup, coverpage=cover]{AJbook}

% ============ 宏包 ============
\usepackage{mycommand}      % 常用数学命令(\dd \R \Z \Q \CC \e \ii 等)
\usepackage{myarrows}       % 可伸缩箭头
\usepackage{chngcntr}       % \numberwithin
\usepackage{imakeidx}
\usepackage{hyperref}
\usetikzlibrary{calc}          % 封面坐标计算

% ============ 编号体系 ============
\numberwithin{equation}{section}   % 公式按节编号 (1.2.7)

% 编号的``注意''环境:原书``注意''与定理共享计数器并编号(如 注意1.4.2,注意3.1.43)
% AJbook.cls 在 \AtBeginDocument 中定义定理类环境(theorem 计数器),
% 故此处也须用 \AtBeginDocument,且先切回 plain 样式(类 hook 末尾是 nonumberplain)
\AtBeginDocument{%
  \theoremstyle{plain}%
  \newtheorem{note}[theorem]{注意}%
  \theoremstyle{nonumberplain}%
  \newtheorem{remarknonum}{注}%
}

% ============ 超链接 ============
\hypersetup{
    pdfauthor={薮田公三, 佐藤圓治, 田中仁, 宮地晶彦, 中路貴彦},
    pdftitle={古典调和分析(解析学百科 I)},
    pdfsubject={调和分析},
    bookmarksnumbered=true
}

% ============ 自定义命令 ============
\newcommand{\C}{\mathbb{C}}        % 复数域(mycommand 用 \CC)
\newcommand{\Fourier}{\mathcal{F}} % Fourier 变换
\newcommand{\Hilbert}{\mathcal{H}} % Hilbert 变换
\newcommand{\esssup}{\operatorname{ess\,sup}}
\newcommand{\diam}{\operatorname{diam}}

% 参考文献:保留正文手动 \section*{参考文献} 标题,抑制 thebibliography 自动生成的章标题
\makeatletter
\patchcmd{\thebibliography}{\chapter*{\bibname}}{}{}{}
\makeatother

% 索引
\makeindex

% 全局断行弹性:消除边界性 Overfull(与 Gelbart 项目同款)
\emergencystretch=3em

% ============ 文档开始 ============
\begin{document}
\frontmatter        % 封面(cover.tex) + 目录(模板自动处理,罗马页码)

% 前言(原书まえがき)
% ============ 前言(原书まえがき,原书 pp.3-8 = PDF pp.3-8) ============
\chapter*{前言}\addcontentsline{toc}{chapter}{前言}

调和分析一词在与正交级数展开相关的地方被广泛使用,但本书所处理的,是在某种意义下直接与三角级数或 Fourier 变换相关的古典调和分析.不过,所谓古典,是指题材是自古以来就被处理的,这一领域目前仍在被积极地研究,其成果和方法也被活跃地用于分析学的其他领域;此外,它还包含具有重要意义的一些未解决问题,是可能参与 21 世纪整个分析学的本质性进展的领域.

本书以在大学数学专业中掌握了复函数论,Lebesgue 积分和泛函分析的基本内容,并打算在研究生阶段专攻分析学的学生们为对象,其主要目的是引导年轻人进入古典调和分析的研究.同时,也期望已经从事分析学各领域研究的年轻人,能把这一领域的理论作为一种素养来熟悉.本书不是对古典调和分析的包罗万象的详细解说,而是作为对这一领域的邀请,解说若干话题.各章相互独立撰写.下面,对本书的内容作稍详细的说明.

第\ref{ch:1}章是对目前常被称为调和分析的实函数论方法的理论的解说.单变量的 $p$ 次可积函数的 Fourier 级数的部分和是否在 $L^{p}$ 范数意义下收敛于原函数的问题,可归结为被称为共轭函数变换或 Hilbert 变换的某个奇异积分算子的 $L^{p}$ 范数有界性.此外,对函数 $f$,把 Laplace 方程 $\Delta u = f$ 的解 $u$ 的二阶偏导数 $\partial^{2}u/\partial x_{i}\partial x_{j}$ 与之对应的线性算子,也是同类的奇异积分算子.这些奇异积分算子的 $L^{p}$ 理论,是 1950 年前后由 A.\,P.\ Calderón 和 A. Zygmund 开创的.作为这一理论的钥匙而被使用的,是分解函数的方法和把算子的弱型估计进行插值的方法,这些是实函数论的方法.实函数论的方法在 1970 年代由 C. Fefferman 和 E.\,M.\ Stein 等人发展成实函数论式的 $H^{p}$ 空间理论.这些成果和方法目前也被有效地用于偏微分方程论等领域.第\ref{ch:1}章就单变量函数解说这一实函数论方法.

第\ref{ch:2}章解说单变量全纯函数的 Hardy 空间理论.构成 Fourier 级数的函数 $e^{in\theta}$,在 $n \geq 0$ 时是全纯函数 $z^{n} = (re^{i\theta})^{n}$ 在 $r = 1$ 处的边界值.由此,把单位圆周上的函数与单位圆盘内的全纯函数对应起来,利用全纯函数的性质来研究 Fourier 级数的性质,这在 20 世纪前半叶曾被积极地研究,被称为 Fourier 级数的复解析方法.例如,单位圆周上的复值测度 $\mu$ 的 Fourier 系数 $\hat{\mu}(n)$ 对负的 $n$ 全为 0 时,$\mu$ 关于 Lebesgue 测度绝对连续----这个 Riesz 兄弟定理可用复解析方法简明地证明.此外,本书虽不涉及,但关于 Hilbert 变换的加权 $L^2$ 范数有界性的 Helson--Szeg\"o 定理,可利用单位圆周上满足 $\log w$ 可积的正值函数 $w$ 等于单位圆盘内全纯函数的边界值的绝对值这一事实,从复函数的性质推导出来.在这个复解析方法中扮演重要角色的全纯函数 Hardy 空间,作为一个函数空间的结构等,其本身提供了许多有趣的问题,得到了大量研究.关于 Hardy 空间各种形式推广的研究,即使在最近也仍是古典调和分析的重要研究课题.

第\ref{ch:3}章是 Fourier 分析中 Banach 环理论的解说.这一理论的发端是 N. Wiener 定理:具有绝对收敛 Fourier 级数的连续函数,若处处不为零,则其倒函数仍具有绝对收敛的 Fourier 级数.这一结果在 Fourier 分析中提供了两个话题.其一是谱综合问题,例如,对于具有绝对收敛 Fourier 级数的连续函数的空间 $A(\mathbf{T})$,可以表述为:"对单位圆的任意闭集,在其上为 0 的 $A(\mathbf{T})$ 的函数,能否由在该闭集邻域内为 0 的 $A(\mathbf{T})$ 的函数在范数意义下逼近(谱综合集)?"此后,这一问题经 I. Gelfand 的理论被推广到交换 Banach 环.1948 年,L. Schwartz 证明了 3 维以上 Euclid 空间的单位球面不是 $A(\mathbf{R}^{n})$ 的谱综合集.决定性的结果是 1959 年 P. Malliavin 关于"$A(\mathbf{T})$ 中非谱综合集的存在性".在与 Fourier 变换有关的 Banach 环中,他们的研究是结合几何性质以及(Lebesgue 测度为零的)稀疏集的性质发展起来的.其二是被称为作用函数的问题,例如,对 $A(\mathbf{T})$:"$[-1,1]$ 上的怎样的函数作用于 $A(\mathbf{T})$?"1930 年代初期,N. Wiener 与 P. Levy 证明了"解析函数作用于 $A(\mathbf{T})$",而 1958 年 Y. Katznelson 证明了其逆也成立.在 P. Malliavin 与 Y. Katznelson 的两个定理中,证明的本质部分从互为共轭的立场进行范数估计,这一点非常有趣.作用函数问题也被推广到交换 Banach 环并得到发展.谱综合与作用函数的研究,是通过考察交换 Banach 环的一个特征而发展起来的;而最近,本书虽不涉及,但已结合调和分析,在各种函数空间中开展研究.

第\ref{ch:4}章是对被称为 Kakeya 问题的,古典 Fourier 分析中若干重要未解决问题的解说.Kakeya 问题,是关于多变量函数或高维图形的一些互相关联的问题的总称.例如,确定多变量 Fourier 级数的被称为 Bochner--Riesz 平均的求和法在 $L^{p}$ 范数意义下收敛的 $p$ 的范围的问题;又如,对高维 Euclid 空间 $\mathbb{R}^{d}$ 上的函数 $f$,把 $f$ 的 Fourier 变换 $\hat{f}$ 到球面 $S$ 的限制 $\hat{f}|S$ 与之对应的算子 $f \mapsto \hat{f}|S$ 成为从 $L^{p}(\mathbb{R}^{d})$ 到 $L^{q}(S)$ 的有界算子的指数 $(p,q)$ 的确定问题.近年来,已明确这些问题与被称为 Nikodym 极大算子的 $\mathbb{R}^{d}$ 上算子的 $L^{p}$ 有界性问题,被称为 Kakeya 集的 $\mathbb{R}^{d}$ 子集的 Hausdorff 维数估计问题等深刻相关.这些问题在二维情形均已解决,但在三维以上情形,与其表面上的简单相反,已弄清楚是非常困难的问题.近年 Wolff 和 Bourgain 等人取得了重要成果,但尚未达到解决.这些问题的核心常常归结为估计满足某条件的有限集合的元素个数的问题,对此的讨论被称为组合论证.这些问题得到解决,并且其中组合论证作为可理解的东西被纳入分析学,尚有待将来.

本书的内容就是以上 4 章,但作为对古典调和分析的邀请,也许还可以设立一章关于概率论的.关于间隙 Fourier 级数与独立随机变量序列的相似性,自古以来就有大量研究.Burkholder,Gundy 与 Silverstein 首次给出用 Brown 运动定义的最大函数刻画全纯函数 Hardy 空间中函数的实部的定理时,该最大函数成为关键.另外,内山明人用实函数论方法证明 BMO 的 Fefferman--Stein 分解,其最初是在鞅的设定中进行的.奇异积分,实函数论 Hardy 空间,BMO 空间等,与对实变函数并行的理论一样,也可在概率论的鞅形式下建立.关于这些,请参考以下文献等.

[1] D. L. Burkholder, R. F. Gundy and M. L. Silverstein, A maximal function characterization of the class $H^{p}$, Trans. Amer. Math. Soc. 157 (1971), 137--153.

[2] A. M. Garsia, Martingale Inequalities, Seminar Notes on Recent Progress, Mathematics Lecture Note Series, W. A. Benjamin, 1973.

[3] K. E. Petersen, Brownian Motion, Hardy Spaces and Bounded Mean Oscillation, London Mathematical Society Lecture Note Series 28, Cambridge University Press, 1977.

[4] D. L. Burkholder, Martingale theory and harmonic analysis in Euclidean spaces, in Harmonic Analysis in Euclidean Spaces, G. Weiss and S. Wainger eds., part 2, Proc. Symp. Pure Math. Vol. 35, American Math. Soc., 1979, pp. 283--301.

[5] A. Uchiyama, A constructive proof of the Fefferman--Stein decomposition of BMO on simple martingales, Conference on Harmonic Analysis in honor of Antoni Zygmund, Volume II, W. Beckner et al. eds, pp. 495--505, 1983.

各章的执笔:第\ref{ch:1}章为薮田公三,第\ref{ch:2}章为中路貴彦,第\ref{ch:3}章为佐藤圓治,第\ref{ch:4}章为田中仁.全书由宮地晶彦负责统稿.

中井英一先生与富田直人先生通读了第\ref{ch:1}章,井上純治先生与山本隆範先生通读了第\ref{ch:2}章,均从原稿阶段起仔细阅读并提出了各种宝贵意见.另外,朝倉書店编辑部从本书的策划阶段起给予了多方面的帮助.在此谨表深切谢意.

2008 年 2 月

宮地晶彦


\mainmatter         % 正文(阿拉伯页码从 1 开始)

% ============ 第1章 奇异积分入门（薮田公三，原书 pp.1-123 = PDF pp.13-135） ============
\chapter{奇异积分入门（薮田公三）}\label{ch:1}

% ============ 第1章 §1.1 准备(原书 pp.1-9 = PDF pp.13-21) ============
\section{准备}\label{sec:1.1}
\index{Lebesgue 收敛定理}
\index{积分号下的微分定理}
\index{Fubini 定理}
\index{Hölder 不等式}
\index{Cauchy--Schwarz 不等式}
\index{Young 不等式}
\index{Schwartz 函数}
\index{速降函数}


首先, 列举经常使用的记号.以下, $A := B$ 表示将 $B$ 定义为 $A$.

$\mathbb{R} := (-\infty, \infty)$, $\mathbb{R}_{+} := (0, \infty)$, $\mathbb{C}$:全体复数

\begin{equation*}
\mathbb{N} := \{1, 2, 3, \dots \}, \quad
\mathbb{N}_{0} := \{0, 1, 2, 3, \dots \}, \quad
\mathbb{Z} := \{\dots, -2, -1, 0, 1, 2, \dots \}
\end{equation*}

$C(\mathbb{R})$:$\mathbb{R}$ 上连续函数的全体

$C_0(\mathbb{R})$:$\mathbb{R}$ 上满足 $\lim_{|x|\to\infty}f(x) = 0$ 的连续函数 $f$ 的全体

$C_{c}(\mathbb{R})$:$\mathbb{R}$ 上支集为有界闭集的连续函数的全体

$C^{k}(\mathbb{R})$:$\mathbb{R}$ 上 $k$ 次连续可微函数的全体

$L^{\infty}(\mathbb{R})$:$\mathbb{R}$ 上本性有界函数的全体

\begin{equation*}
\| f \|_{L^{\infty}(\mathbb{R})} = \| f \|_{L^{\infty}} = \| f \|_{\infty} := \operatorname{ess\,sup}_{x \in \mathbb{R}} |f(x)|
\end{equation*}

$L^{p}(\mathbb{R})$:$\mathbb{R}$ 上 $p$ 次可积函数的全体

\begin{equation*}
\| f \|_{L^{p}(\mathbb{R})} = \| f \|_{L^{p}} = \| f \|_{p} := \left(\int_{\mathbb{R}} |f(x)|^{p}\dd x\right)^{1/p}
\end{equation*}

$\lim_{j\to\infty}\|f_{j}-f\|_{L^{p}} = 0$ 记为 $f_{j}\to f$ in $L^{p}$($j\to\infty$).

对几乎所有的 $x \in \mathbb{R}$ 成立 $f(x) = g(x)$ 记为
$f(x) = g(x)$(a.e. $x\in\mathbb{R}$).

$L_{\mathrm{loc}}^{p}(\mathbb{R})$:在 $\mathbb{R}$ 的任意有界区间上 $p$ 次可积的函数的全体

对 $\mathbb{R}$ 上的可测集 $E$, 用 $|E|$ 表示 $E$ 的 Lebesgue 测度.

\begin{definition}\label{def:1.1.1}
$f \in \mathcal{S}(\mathbb{R})$, 或者说 $f$ 是 Schwartz 函数, 或者说 $f$ 是速降函数, 是指 $f \in C^{\infty}(\mathbb{R})$, 并且对任意 $k, \ell \in \mathbb{N}_0$, 存在正数 $C_{k,\ell}$ 使得
\begin{equation*}
\left| (1 + |x|)^{k} f^{(\ell)}(x) \right| \leq C_{k,\ell} \quad (x \in \mathbb{R}).
\end{equation*}
\end{definition}

接下来, 叙述 Lebesgue 积分中的基本事实.

\begin{theorem}[Lebesgue 收敛定理]\label{thm:1.1.2}
$f_{j}(x)$($j = 1, 2, \ldots$)是可测函数, $\varphi(x) \in L^{1}(\mathbb{R})$, 且
\begin{align*}
|f_{j}(x)| \leq \varphi(x) & (\text{a.e. } x \in \mathbb{R},\ j = 1, 2, \ldots) \\
\lim_{j \to \infty} f_{j}(x) = f(x) & (\text{a.e. } x \in \mathbb{R})
\end{align*}
时, 
\begin{equation*}
\lim_{j \rightarrow \infty} \int_{-\infty}^{\infty} f_{j}(x)\dd x = \int_{-\infty}^{\infty} f(x)\dd x.
\end{equation*}
\end{theorem}

\begin{corollary}\label{cor:1.1.3}
设 $1 \leq p < \infty$, $f_{j}(x)$($j = 1, 2, \ldots$)是可测函数, $\varphi(x) \in L^{p}(\mathbb{R})$, 且
\begin{align*}
|f_{j}(x)| \leq \varphi(x) & (\text{a.e. } x \in \mathbb{R},\ j = 1, 2, \ldots) \\
\lim_{j \to \infty} f_{j}(x) = f(x) & (\text{a.e. } x \in \mathbb{R})
\end{align*}
时, 
\begin{align*}
\lim_{j \to \infty} \left(\int_{-\infty}^{\infty} |f_{j}(x) - f(x)|^{p}\dd x\right)^{1/p} &= 0, \\
\lim_{j \to \infty} \left(\int_{-\infty}^{\infty} |f_{j}(x)|^{p}\dd x\right)^{1/p} &= \left(\int_{-\infty}^{\infty} |f(x)|^{p}\dd x\right)^{1/p}.
\end{align*}
\end{corollary}

\begin{theorem}[积分号下的微分定理]\label{thm:1.1.4}
设 $f(x,t)$ 关于 $t$ 可偏导, 对每个 $t$, $f(x,t)$ 与 $f_{t}(x,t)$ 关于 $x$ 可积, 并且存在可积函数 $\varphi(x)$ 使得 $|f_{t}(x,t)| \leq \varphi(x)$.于是 $\int_{-\infty}^{\infty} f(x,t)\dd x$ 关于 $t$ 可微, 且
\begin{equation*}
\frac{\dd}{\dd t}\left(\int_{-\infty}^{\infty} f(x,t)\dd x\right) = \int_{-\infty}^{\infty} f_{t}(x,t)\dd x.
\end{equation*}
这里, 若对几乎所有的 $x$, $f_{t}(x,t)$ 关于 $t$ 连续, 则微分得到的函数也连续.
\end{theorem}

\begin{proof}
令 $F(t) = \int_{-\infty}^{\infty} f(x,t)\dd x$, 则
\begin{equation*}
\frac{F(t+h) - F(t)}{h} = \int_{-\infty}^{\infty} \frac{f(x,t+h) - f(x,t)}{h}\dd x.
\end{equation*}
由中值定理, 
\begin{equation*}
\left| \frac{f(x,t+h) - f(x,t)}{h} \right| \leq \varphi(x),
\end{equation*}
又由假设
\begin{equation*}
\lim_{h \rightarrow 0} \frac{f(x,t+h) - f(x,t)}{h} = f_{t}(x,t),
\end{equation*}
故由 Lebesgue 收敛定理
\begin{equation*}
F^{\prime}(t) = \lim_{h \rightarrow 0} \frac{F(t+h) - F(t)}{h} = \int_{-\infty}^{\infty} f_{t}(x,t)\dd x.
\end{equation*}
最后一步也是 Lebesgue 收敛定理的应用.
\end{proof}

\begin{theorem}[Fubini 定理]\label{thm:1.1.5}
设 $f(x,y)$ 关于 $(x,y)$ 可测.若
\begin{align*}
&\iint_{\mathbb{R}^{2}} |f(x,y)|\dd x\dd y, \quad \int_{-\infty}^{\infty}\left(\int_{-\infty}^{\infty} |f(x,y)|\dd x\right)\dd y, \\
&\int_{-\infty}^{\infty}\left(\int_{-\infty}^{\infty} |f(x,y)|\dd y\right)\dd x
\end{align*}
三者之一有限, 则其余两个也有限, 并且成立如下结论:

(i) 对几乎所有的 $y \in \mathbb{R}$, $\int_{-\infty}^{\infty} f(x,y)\dd x$ 存在且作为 $y$ 的函数可测.

(ii) 对几乎所有的 $x \in \mathbb{R}$, $\int_{-\infty}^{\infty} f(x,y)\dd y$ 存在且作为 $x$ 的函数可测.

(iii)
\begin{align*}
\iint_{\mathbb{R}^{2}} f(x,y)\dd x\dd y
& = \int_{-\infty}^{\infty}\left(\int_{-\infty}^{\infty} f(x,y)\dd x\right)\dd y \\
& = \int_{-\infty}^{\infty}\left(\int_{-\infty}^{\infty} f(x,y)\dd y\right)\dd x.
\end{align*}
\end{theorem}

\begin{corollary}\label{cor:1.1.6}
设 $0 < p < \infty$.则
\begin{equation*}
\int_{-\infty}^{\infty} |f(x)|^{p}\dd x = p\int_{0}^{\infty} |\{x;\ |f(x)| > t\}|\,t^{p-1}\dd t.
\end{equation*}
\end{corollary}

\begin{theorem}[Hölder 不等式]\label{thm:1.1.7}
设 $1 \leq p \leq \infty$, $1/p + 1/p' = 1$.若 $f(x) \in L^{p}(\mathbb{R})$, $g(x) \in L^{p'}(\mathbb{R})$, 则 $f(x)g(x)$ 可积, 且
\begin{equation*}
\left| \int_{-\infty}^{\infty} f(x)g(x)\dd x \right| \leq \| f \|_{L^{p}} \| g \|_{L^{p'}}.
\end{equation*}
\end{theorem}

Hölder 不等式在 $p = 2$ 时称为 Cauchy--Schwarz 不等式.此后也常以此名使用.

\begin{theorem}[Young 不等式的一个特殊情形]\label{thm:1.1.8}
设 $1 \leq p \leq \infty$.若 $f(x) \in L^{1}(\mathbb{R})$, $g(x) \in L^{p}(\mathbb{R})$, 则对几乎所有的 $x$, $\int f(x-y)g(y)\dd y$ 存在, 且下式成立:
\begin{equation*}
\left(\int_{-\infty}^{\infty}\left| \int_{-\infty}^{\infty} f(x-y)g(y)\dd y \right|^{p}\dd x\right)^{1/p} \leq \| f \|_{L^{1}} \| g \|_{L^{p}}.
\end{equation*}
\end{theorem}

\begin{proof}
\begin{align*}
& \left| \int_{-\infty}^{\infty} f(x-y)g(y)\dd y \right| \\
& \qquad \leq \int_{-\infty}^{\infty} |f(x-y)g(y)|\dd y
\end{align*}
\begin{align*}
&\leq \left(\int_{-\infty}^{\infty} |f(x-y)|\dd y\right)^{1/p'}\left(\int_{-\infty}^{\infty} |f(x-y)||g(y)|^{p}\dd y\right)^{1/p} \\
&= (\| f \|_{L^{1}})^{1/p'}\left(\int_{-\infty}^{\infty} |f(x-y)||g(y)|^{p}\dd y\right)^{1/p}.
\end{align*}
又由 Fubini 定理
\begin{align*}
\int_{-\infty}^{\infty} & \left(\int_{-\infty}^{\infty} |f(x-y)||g(y)|^{p}\dd y\right)\dd x
= \int_{-\infty}^{\infty}\left(\int_{-\infty}^{\infty} |f(x-y)||g(y)|^{p}\dd x\right)\dd y \\
& = \| f \|_{L^{1}} \int_{-\infty}^{\infty} |g(y)|^{p}\dd y < \infty.
\end{align*}
从而, 由 Fubini 定理, 对几乎所有的 $x$, $\int_{-\infty}^{\infty} f(x-y)g(y)\dd y$ 存在.由上面的两个不等式即得所证不等式.
\end{proof}

\begin{theorem}\label{thm:1.1.9}
设 $K(t,x)$ 满足如下三条:
\begin{enumerate}[(i)]
\item $\int_{-\infty}^{\infty}|K(t,x)|\dd x \leq A \quad (t > 0)$, 
\item $\int_{-\infty}^{\infty} K(t,x)\dd x = 1 \quad (t > 0)$, 
\item 对每个 $\delta > 0$, 
\begin{equation*}
\int_{|x| > \delta} |K(t,x)|\dd x \rightarrow 0 \quad (t \rightarrow 0).
\end{equation*}
\end{enumerate}
则

(a) 若 $f(x) \in L^{\infty}(\mathbb{R})$ 且在点 $x_{0}$ 连续, 则
\begin{equation*}
\int_{-\infty}^{\infty} K(t, x_{0} - y)f(y)\dd y \rightarrow f(x_{0}) \quad (t \rightarrow 0).
\end{equation*}

(b) 若 $f(x) \in L^{\infty}(\mathbb{R}) \cap C(\mathbb{R})$, 则在任意有限区间上一致地
\begin{equation*}
\int_{-\infty}^{\infty} K(t, x - y)f(y)\dd y \rightarrow f(x) \quad (t \rightarrow 0).
\end{equation*}

(c) 若 $1 \leq p < \infty$ 且 $f \in L^{p}(\mathbb{R})$, 则
\begin{equation*}
\left\| \int_{-\infty}^{\infty} K(t, x - y)f(y)\dd y - f(x) \right\|_{L^{p}} \rightarrow 0 \quad (t \rightarrow 0).
\end{equation*}

(d) 若 $1 \leq p \leq \infty$ 且 $f \in L^{p}(\mathbb{R})$, 则
\begin{equation*}
\left\| \int_{-\infty}^{\infty} K(t, x - y)f(y)\dd y \right\|_{L^{p}} \leq A\| f \|_{L^{p}}.
\end{equation*}
\end{theorem}

\begin{proof}
(a) 首先利用 (ii), 
\begin{align*}
\int_{-\infty}^{\infty} K(t, x_{0} - y)f(y)\dd y - f(x_{0})
& = \int_{-\infty}^{\infty} K(t, x_{0} - y)f(y)\dd y - f(x_{0})\int_{-\infty}^{\infty} K(t, x_{0} - y)\dd y \\
& = \int_{-\infty}^{\infty} K(t, x_{0} - y)\bigl(f(y) - f(x_{0})\bigr)\dd y.
\end{align*}
从而, 
\begin{align*}
& \left| \int_{-\infty}^{\infty} K(t, x_{0} - y)f(y)\dd y - f(x_{0}) \right| \\
& \qquad \leq \int_{|y - x_{0}| > \delta} |K(t, x_{0} - y)||f(y) - f(x_{0})|\dd y \\
& \qquad\quad + \int_{|y - x_{0}| \leq \delta} |K(t, x_{0} - y)||f(y) - f(x_{0})|\dd y \\
& \qquad \leq 2\| f \|_{\infty} \int_{|y - x_{0}| > \delta} |K(t, x_{0} - y)|\dd y \\
& \qquad\quad + \sup_{|x_{0} - y| \leq \delta} |f(y) - f(x_{0})| \int_{-\infty}^{\infty} |K(t, x_{0} - y)|\dd y.
\end{align*}
于是, 任取 $\varepsilon > 0$.因 $f(x)$ 在点 $x_0$ 连续, 故存在具有如下性质的 $\delta > 0$:
\begin{equation*}
|y - x_{0}| \leq \delta \quad \text{则} \quad |f(y) - f(x_{0})| < \varepsilon.
\end{equation*}
又由 (iii), 对此 $\delta$, 存在 $t_{0} > 0$ 使得
\begin{equation*}
0 < t < t_{0} \quad \text{则} \quad \int_{|y - x_{0}| > \delta} |K(t, x_{0} - y)|\dd y = \int_{|z| > \delta} |K(t, z)|\dd z < \varepsilon.
\end{equation*}
综合运用 (i) 可得, 对任意 $\varepsilon > 0$ 存在 $t_0 > 0$, 使得当 $0 < t < t_0$ 时
\begin{equation*}
\left| \int_{-\infty}^{\infty} K(t, x_{0} - y)f(y)\dd y - f(x_{0}) \right| \leq (2\| f \|_{\infty} + A)\varepsilon
\end{equation*}
故 (a) 得证.

(b) 令 $I = [-a, a]$, 则由假设 $f(x)$ 在 $I$ 上一致连续.由此与 (a) 的证法可得结论.

(c) 设 $f \in L^{p}(\mathbb{R})$.任取 $\varepsilon > 0$.利用 $C_{c}(\mathbb{R})$ 在 $L^{p}(\mathbb{R})$ 中稠密, 存在连续函数 $f_{0}$ 与正数 $R$ 使得
\begin{equation*}
\| f - f_{0} \|_{L^{p}} < \varepsilon
\end{equation*}
且 $\supp f_{0} \subset \{x;\ |x| \leq R\}$.利用 $L^{p}$ 的三角不等式
\begin{align*}
& \left\| \int_{-\infty}^{\infty} K(t, x - y)f(y)\dd y - f(x) \right\|_{L^{p}} \\
& \quad \leq \left\| \int_{-\infty}^{\infty} K(t, x - y)(f(y) - f_{0}(y))\dd y \right\|_{L^{p}} \\
& \quad\quad + \left\| \int_{-\infty}^{\infty} K(t, x - y)f_{0}(y)\dd y - f_{0}(x) \right\|_{L^{p}} + \| f_{0}(x) - f(x) \|_{L^{p}} \\
& \quad = I_{1} + I_{2} + I_{3}
\end{align*}
由 Young 不等式与 (i)
\begin{equation*}
I_{1} \leq \| K(t, \cdot) \|_{L^{1}} \| f - f_{0} \|_{L^{p}} \leq A\| f - f_{0} \|_{L^{p}}.
\end{equation*}
故
\begin{equation*}
I_{1} + I_{3} < (A + 1)\varepsilon.
\end{equation*}
关于 $I_{2}$, 利用 (b), 存在正数 $t_{0}$, 使得当 $0 < t < t_{0}$ 时, 对有界闭区间 $[-R-1, R+1]$ 上的任意 $x$, 
\begin{equation*}
\left| \int_{-\infty}^{\infty} K(t, x - y)f_{0}(y)\dd y - f_{0}(x) \right| < \frac{\varepsilon}{2(R+1)^{1/p}}.
\end{equation*}
从而
\begin{equation*}
\int_{|x| \leq R + 1}\left| \int_{-\infty}^{\infty} K(t, x - y)f_{0}(y)\dd y - f_{0}(x) \right|^{p}\dd x < \varepsilon^{p}.
\end{equation*}
其次, 利用 $f_{0}(x) = 0$($|x| \geq R$)与 Hölder 不等式,(i), 以及 Fubini 定理, 
\begin{align*}
J_{t} &:= \int_{|x| \geq R + 1}\left| \int_{-\infty}^{\infty} K(t, x - y)f_{0}(y)\dd y - f_{0}(x) \right|^{p}\dd x \\
&= \int_{|x| \geq R + 1}\left| \int_{-\infty}^{\infty} K(t, x - y)f_{0}(y)\dd y \right|^{p}\dd x \\
&\leq \int_{|x| \geq R + 1}\left(\int_{-\infty}^{\infty} |K(t, x - y)|\dd y\right)^{p/p'}\left(\int_{-\infty}^{\infty} |K(t, x - y)||f_{0}(y)|^{p}\dd y\right)\dd x \\
&\leq A^{p/p'}\int_{|x| \geq R + 1}\left(\int_{-\infty}^{\infty} |K(t, x - y)||f_{0}(y)|^{p}\dd y\right)\dd x \\
&= A^{p/p'}\int_{|x| \geq R + 1}\left(\int_{|y| \leq R} |K(t, x - y)||f_{0}(y)|^{p}\dd y\right)\dd x \\
&= A^{p/p'}\int_{|y| \leq R}\left(\int_{|x| \geq R + 1} |K(t, x - y)|\dd x\right)|f_{0}(y)|^{p}\dd y
\end{align*}
这里注意到当 $|y| \leq R$, $|x| \geq R+1$ 时 $|x-y| \geq |x| - |y| \geq (R+1) - R = 1$, 
\begin{align*}
J_{t} &\leq A^{p/p'}\int_{|y| \leq R}\left(\int_{|x - y| \geq 1} |K(t, x - y)|\dd x\right)|f_{0}(y)|^{p}\dd y \\
&= A^{p/p'}\int_{|y| \leq R}\left(\int_{|z| \geq 1} |K(t, z)|\dd z\right)|f_{0}(y)|^{p}\dd y \\
&= A^{p/p'}\int_{|y| \leq R} |f_{0}(y)|^{p}\dd y \int_{|z| \geq 1} |K(t, z)|\dd z \\
&= A^{p/p'}\| f_{0} \|_{L^{p}}^{p}\int_{|x| \geq 1} |K(t, x)|\dd x.
\end{align*}
由假设 (iii), 存在 $t_{1} > 0$ 使得
\begin{equation*}
J_{t} < \varepsilon^{p} \quad (0 < t < t_{1}).
\end{equation*}
故当 $0 < t < \min(t_{0}, t_{1})$ 时
\begin{equation*}
I_{2}^{p} < \varepsilon^{p} + \varepsilon^{p} < (2\varepsilon)^{p}.
\end{equation*}
从而
\begin{equation*}
I_{1} + I_{2} + I_{3} < (A + 3)\varepsilon
\end{equation*}
由此证得
\begin{equation*}
\lim_{t \rightarrow 0}\left\| \int_{-\infty}^{\infty} K(t, x - y)f(y)\dd y - f(x) \right\|_{L^{p}} = 0.
\end{equation*}

(d) 与 (c) 中 $I_{1}$ 的估计一样, 由 Young 不等式与 (i) 即得结论.
\end{proof}

\begin{example}\label{exa:1.1.10}
设 $\varphi(x) \in L^{1}(\mathbb{R})$, $\int_{-\infty}^{\infty} \varphi(x)\dd x = 1$, 并令 $\varphi_{t}(x) = t^{-1}\varphi(t^{-1}x)$.则 $K(t,x) = \varphi_{t}(x)$ 满足定理~\ref{thm:1.1.9} 的假设 (i),(ii),(iii).

事实上, (i) 由变量替换 $y = x/t$, 
\begin{equation*}
\int_{-\infty}^{\infty} |\varphi_{t}(x)|\dd x = \int_{-\infty}^{\infty}\left| \frac{1}{t}\varphi\left(\frac{x}{t}\right) \right|\dd x = \int_{-\infty}^{\infty} |\varphi(y)|\dd y.
\end{equation*}

(ii) 与 (i) 类似, 
\begin{equation*}
\int_{-\infty}^{\infty} \varphi_{t}(x)\dd x = \int_{-\infty}^{\infty} \frac{1}{t}\varphi\left(\frac{x}{t}\right)\dd x = \int_{-\infty}^{\infty} \varphi(y)\dd y = 1.
\end{equation*}

(iii) 因 $\varphi(x) \in L^{1}(\mathbb{R})$, 
\begin{equation*}
\int_{|x| > \delta} |\varphi_{t}(x)|\dd x = \int_{|x| > \delta}\left| \frac{1}{t}\varphi\left(\frac{x}{t}\right)\right|\dd x = \int_{|y| > \delta/t} |\varphi(y)|\dd y \rightarrow 0 \quad (t \rightarrow 0).
\end{equation*}
\end{example}

% ============ 第1章 §1.2 Fourier变换(原书 pp.9-29 = PDF pp.21-41) ============
\section{Fourier 变换}\label{sec:1.2}
\index{Fourier 变换}
\index{Fourier 变换!逆 Fourier 变换}
\index{Fourier 变换!(L2(R) 意义下的)}
\index{逆 Fourier 变换!(L2(R) 意义下的)}
\index{Riemann--Lebesgue 定理}
\index{Hausdorff--Young 不等式}
\index{Parseval 等式}
\index{Plancherel 等式}
\index{卷积}
\index{Riesz--Thorin 定理}
\index{Young 不等式}
\index{插值定理!M. Riesz--Thorin 的}


\subsection{\texorpdfstring{$L^1(\mathbb{R})$ 函数与速降函数的 Fourier 变换}{L1(R) 函数与速降函数的 Fourier 变换}}\label{sec:1.2.1}

\begin{definition}\label{def:1.2.1}
对 $f(x) \in L^{1}(\mathbb{R})$,
\begin{equation*}
\hat{f}(\xi) = \mathcal{F}f(\xi) = \int_{-\infty}^{\infty} f(x)e^{-ix\xi}\dd x
\end{equation*}
称为 $f$ 的 Fourier 变换.此外,
\begin{equation*}
\check{f}(x) = \mathcal{F}^{-1}f(x) = \frac{1}{2\pi}\int_{-\infty}^{\infty} f(\xi)e^{ix\xi}\dd \xi
\end{equation*}
称为 $f$ 的逆 Fourier 变换.由定义
\begin{equation*}
\mathcal{F}^{-1}f(x) = \frac{1}{2\pi}\mathcal{F}\bigl(f(-\,\cdot)\bigr)(x) = \frac{1}{2\pi}\mathcal{F}f(-x).
\end{equation*}
\end{definition}

下面举出以后经常使用的 Fourier 变换的具体例子.

\begin{example}\label{exa:1.2.2}
\begin{equation}\tag{1.2.1}
\mathcal{F}\bigl(e^{-\frac{x^{2}}{2}}\bigr)(\xi) = \sqrt{2\pi}\,e^{-\frac{\xi^{2}}{2}}.
\label{eq:1.2.1}\end{equation}
\begin{equation}\tag{1.2.2}
\mathcal{F}^{-1}\mathcal{F}\bigl(e^{-\frac{x^{2}}{2}}\bigr) = e^{-\frac{x^{2}}{2}}.
\label{eq:1.2.2}\end{equation}
\end{example}

\begin{proof}
\begin{align*}
\mathcal{F}\bigl(e^{-\frac{x^{2}}{2}}\bigr)(\xi) & = \int_{-\infty}^{\infty} e^{-\frac{x^{2}}{2}}e^{-ix\xi}\dd x = \int_{-\infty}^{\infty} e^{-\frac{1}{2}(x^{2} + 2ix\xi)}\dd x \\
& = \int_{-\infty}^{\infty} e^{-\frac{1}{2}(x + i\xi)^{2} - \frac{\xi^{2}}{2}}\dd x = e^{-\frac{\xi^{2}}{2}}\int_{-\infty}^{\infty} e^{-\frac{1}{2}(x + i\xi)^{2}}\dd x \\
& = e^{-\frac{\xi^{2}}{2}}\int_{-\infty}^{\infty} e^{-\frac{x^{2}}{2}}\dd x = \sqrt{2\pi}\,e^{-\frac{\xi^{2}}{2}}.
\end{align*}
第 5 个等式是简单的复函数论习题.

等式~\eqref{eq:1.2.2} 使用等式~\eqref{eq:1.2.1} 与定义~\ref{def:1.2.1} 中逆 Fourier 变换的定义,得到
\begin{equation*}
\begin{aligned}
\mathcal{F}^{-1}\mathcal{F}\bigl(e^{-\frac{x^{2}}{2}}\bigr) &= \mathcal{F}^{-1}\bigl(\sqrt{2\pi}\,e^{-\frac{\xi^{2}}{2}}\bigr)(x) = \sqrt{2\pi}\,\mathcal{F}^{-1}\bigl(e^{-\frac{\xi^{2}}{2}}\bigr)(x) \\
&= \sqrt{2\pi}\cdot\frac{1}{2\pi}\mathcal{F}\bigl(e^{-\frac{(-\xi)^{2}}{2}}\bigr)(x) = e^{-\frac{x^{2}}{2}}.
\end{aligned}
\end{equation*}
\end{proof}

作为 $L^{1}$ 函数 Fourier 变换的基本定理,从下面的定理开始.

\begin{theorem}\label{thm:1.2.3}
若 $f \in L^{1}(\mathbb{R})$,则
\begin{enumerate}[(i)]
\item $|\mathcal{F}f(\xi)| \leq \int_{-\infty}^{\infty}|f(x)|\dd x$($\xi \in \mathbb{R}$),且 $\mathcal{F}f(\xi)$ 在 $\mathbb{R}$ 上一致连续.
\item $\lim_{|\xi| \to \infty}\mathcal{F}f(\xi) = 0$(Riemann--Lebesgue 定理).
\end{enumerate}
\end{theorem}

\begin{proof}
(i) 第一个估计是显然的,我们来证明一致连续性.
\begin{align*}
|\hat{f}(\xi + h) - \hat{f}(\xi)| & = \left|\int_{-\infty}^{\infty}\bigl(e^{-ix(\xi+h)} - e^{-ix\xi}\bigr)f(x)\dd x\right| \\
& \leq \int_{-\infty}^{\infty}\left|e^{-ix(\xi+h)} - e^{-ix\xi}\right||f(x)|\dd x \\
& = \int_{-\infty}^{\infty}\left|e^{-ixh} - 1\right||f(x)|\dd x.
\end{align*}
于是令 $h \to 0$.由于 $|(e^{ixh} - 1)f(x)| \leq 2|f(x)| \in L^{1}(\mathbb{R})$,故可用 Lebesgue 收敛定理,且与 $\xi$ 无关地
\begin{equation*}
\lim_{h \to 0}|\hat{f}(\xi + h) - \hat{f}(\xi)| = 0.
\end{equation*}
因此,$\mathcal{F}f(\xi)$ 在 $\mathbb{R}$ 上一致连续.

(ii) 任意固定正数 $\varepsilon$.由于 $C_{c}^{\infty}(\mathbb{R})$ 在 $L^{1}(\mathbb{R})$ 中稠密,存在满足
\begin{equation*}
\| f - f_{0} \|_{L^{1}} < \varepsilon
\end{equation*}
的 $f_{0} \in C_{c}^{1}(\mathbb{R})$.
\begin{align*}
\mathcal{F}f(\xi) & = \int_{-\infty}^{\infty}e^{-ix\xi}\bigl(f(x) - f_{0}(x)\bigr)\dd x + \int_{-\infty}^{\infty}e^{-ix\xi}f_{0}(x)\dd x \\
& = \int_{-\infty}^{\infty}e^{-ix\xi}\bigl(f(x) - f_{0}(x)\bigr)\dd x + \frac{1}{i\xi}\int_{-\infty}^{\infty}e^{-ix\xi}f_{0}'(x)\dd x.
\end{align*}
因此,若 $|\xi| > \frac{1}{\varepsilon}\int_{-\infty}^{\infty}|f_{0}'(x)|\dd x$,则
\begin{align*}
|\mathcal{F}f(\xi)| &\leq \int_{-\infty}^{\infty}|f(x) - f_{0}(x)|\dd x + \frac{1}{|\xi|}\int_{-\infty}^{\infty}|f_{0}'(x)|\dd x \\
&< \varepsilon + \varepsilon = 2\varepsilon.
\end{align*}
因此
\begin{equation*}
\lim_{|\xi| \to \infty}\mathcal{F}f(\xi) = 0.
\end{equation*}
\end{proof}

\begin{lemma}\label{lem:1.2.4}
若 $f, g \in L^{1}(\mathbb{R})$,则
\begin{equation*}
\int f(x)\mathcal{F}g(x)\dd x = \int \mathcal{F}f(x)g(x)\dd x.
\end{equation*}
\end{lemma}

\begin{proof}
由于 $e^{-ix\xi}f(x)g(\xi) \in L^{1}(\mathbb{R}\times\mathbb{R})$,由 Fubini 定理
\begin{align*}
\int f(x)\mathcal{F}g(x)\dd x & = \int f(x)\left(\int e^{-ix\xi}g(\xi)\dd \xi\right)\dd x \\
& = \int g(\xi)\left(\int e^{-ix\xi}f(x)\dd x\right)\dd \xi = \int \mathcal{F}f(x)g(x)\dd x.
\end{align*}
\end{proof}

\begin{theorem}\label{thm:1.2.5}
若 $f \in L^{1}(\mathbb{R})$ 且 $\mathcal{F}f \in L^{1}(\mathbb{R})$,则
\begin{equation*}
f(x) = \mathcal{F}^{-1}\mathcal{F}f(x) \quad (\text{a.e. } x \in \mathbb{R}),
\end{equation*}
并且,在测度为 $0$ 的集合上修正 $f(x)$ 后,
\begin{equation*}
f(x) \in C_{0}(\mathbb{R}) \cap L^{\infty}(\mathbb{R}).
\end{equation*}
\end{theorem}

\begin{proof}
设 $g \in L^{1}(\mathbb{R})$ 且 $\mathcal{F}g \in L^{1}(\mathbb{R})$.于是,由于 $f \in L^{1}(\mathbb{R})$,$\mathcal{F}g \in L^{1}(\mathbb{R})$,故 $e^{-iy\xi}f(y)e^{ix\xi}\mathcal{F}g(\xi) \in L^{1}(\mathbb{R}_{y}\times\mathbb{R}_{\xi})$.因此,使用 Fubini 定理
\begin{align}
\int e^{ix\xi}\mathcal{F}f(\xi)\mathcal{F}g(\xi)\dd\xi
&= \int e^{ix\xi}\mathcal{F}g(\xi)\left(\int e^{-iy\xi}f(y)\dd y\right)\dd\xi = \int\left(\int e^{i(x-y)\xi}\mathcal{F}g(\xi)\dd\xi\right)f(y)\dd y \tag{1.2.3}\label{eq:1.2.3}\\
&= 2\pi\int\mathcal{F}^{-1}(\mathcal{F}g)(x-y)f(y)\dd y = 2\pi\int\mathcal{F}^{-1}(\mathcal{F}g)(y)f(x-y)\dd y. \tag{1.2.5}\label{eq:1.2.5}
\end{align}

再者,由于 $g \in L^{1}(\mathbb{R})$,$\mathcal{F}f \in L^{1}(\mathbb{R})$,交换 $f$ 与 $g$ 的角色,得到下述等式:
\begin{equation}\tag{1.2.6}
\int\mathcal{F}^{-1}(\mathcal{F}g)(x-y)f(y)\dd y = \int\mathcal{F}^{-1}(\mathcal{F}f)(y)g(x-y)\dd y.
\label{eq:1.2.6}\end{equation}

这里,令 $g(x) = e^{-\frac{1}{2}\left(\frac{x}{t}\right)^{2}}$($t > 0$).由于 $g \in L^{1}(\mathbb{R})$,$\mathcal{F}g \in L^{1}(\mathbb{R})$,等式~\eqref{eq:1.2.6} 成立.进一步,由例~\ref{exa:1.2.2} 与变量替换得
\begin{equation*}
\mathcal{F}g(\xi) = t\mathcal{F}\bigl(e^{-x^{2}/2}\bigr)(t\xi) = \sqrt{2\pi}\,t\,e^{-\frac{(t\xi)^{2}}{2}}.
\end{equation*}
再次使用变量替换与例~\ref{exa:1.2.2},
\begin{align*}
\mathcal{F}^{-1}(\mathcal{F}g)(x) & = \frac{1}{2\pi}\mathcal{F}(\mathcal{F}g)(-x) = \frac{t}{\sqrt{2\pi}}\mathcal{F}\bigl(e^{-\frac{(t\xi)^{2}}{2}}\bigr)(-x) \\
& = \frac{t}{\sqrt{2\pi}}\cdot\frac{1}{t}\mathcal{F}\bigl(e^{-\frac{\xi^{2}}{2}}\bigr)\left(-\frac{x}{t}\right) = \frac{1}{\sqrt{2\pi}}\cdot\sqrt{2\pi}\,e^{-\frac{1}{2}\left(\frac{x}{t}\right)^{2}} = e^{-\frac{1}{2}\left(\frac{x}{t}\right)^{2}}.
\end{align*}
因此,由式~\eqref{eq:1.2.6}
\begin{equation}\tag{1.2.7}
\frac{1}{\sqrt{2\pi}\,t}\int e^{-\frac{1}{2}\left(\frac{x-y}{t}\right)^{2}}f(y)\dd y = \frac{1}{\sqrt{2\pi}\,t}\int e^{-\frac{1}{2}\left(\frac{x-y}{t}\right)^{2}}\mathcal{F}^{-1}(\mathcal{F}f)(y)\dd y.
\label{eq:1.2.7}\end{equation}
这里,令 $\varphi(x) = \frac{1}{\sqrt{2\pi}}e^{-\frac{x^{2}}{2}}$,$\varphi_{t}(x) = \frac{1}{t}\varphi\left(\frac{x}{t}\right)$,则
\begin{equation*}
\int|\varphi(x)|\dd x = \int\varphi(x)\dd x = \frac{1}{\sqrt{2\pi}}\int e^{-\frac{x^{2}}{2}}\dd x = 1
\end{equation*}
所以,如例~\ref{exa:1.1.10} 所述,$\varphi(x)$ 满足定理~\ref{thm:1.1.9} 的假设.于是应用定理~\ref{thm:1.1.9},在式~\eqref{eq:1.2.7} 中令 $t \to 0$,得
\begin{equation*}
f(x) = \mathcal{F}^{-1}(\mathcal{F}f)(x) \quad (\text{a.e. } x \in \mathbb{R}).
\end{equation*}
由于 $\mathcal{F}f \in L^{1}(\mathbb{R})$,由定理~\ref{thm:1.2.3},$\mathcal{F}^{-1}\mathcal{F}f \in C_{0}(\mathbb{R}) \cap L^{\infty}(\mathbb{R})$.
\end{proof}

\begin{lemma}\label{lem:1.2.6}
若 $f \in L^{1}(\mathbb{R})$ 且 $\mathcal{F}f = 0$,则 $f = 0$.
\end{lemma}

\begin{proof}
由于 $f \in L^{1}(\mathbb{R})$ 且 $\mathcal{F}f = 0 \in L^{1}(\mathbb{R})$,由定理~\ref{thm:1.2.5},
\begin{equation*}
f(x) = \mathcal{F}^{-1}(\mathcal{F}f)(x) = 0 \quad (\text{a.e. } x \in \mathbb{R}).
\end{equation*}
\end{proof}

为了研究速降函数 Fourier 变换的基本性质,先从下面这个简单的事实开始.

\begin{lemma}\label{lem:1.2.7}
设 $k \in \mathbb{N}$.若 $f(x) \in \mathcal{S}(\mathbb{R})$,$g(x) \in C^{k}(\mathbb{R})$,且 $|g^{(k-1)}(x)| \leq C_{k-1}(1 + |x|)^{m_{k-1}}$($C_{k-1}$ 为正常数,$m_{k-1} \in \mathbb{N}_{0}$),则对 $0 \leq \ell \leq k$,
\begin{equation}\tag{1.2.8}
\int_{-\infty}^{\infty} f^{(\ell)}(x)g(x)\dd x = (-1)^{\ell}\int_{-\infty}^{\infty} f(x)g^{(\ell)}(x)\dd x.
\label{eq:1.2.8}\end{equation}
\end{lemma}

\begin{proof}
由 $\mathcal{S}(\mathbb{R})$ 的定义与 $|g^{(k-1)}(x)| \leq C_{k-1}(1 + |x|)^{m_{k-1}}$,可得 $|g^{(\ell)}(x)| \leq C_{\ell}(1 + |x|)^{m_{\ell}}$($0 \leq \ell \leq k - 1$),从而注意到 $\lim_{|x| \to \infty}|f^{(m)}(x)g^{(\ell-m)}(x)| = 0$($0 \leq m < \ell$),分部积分即可.
\end{proof}

\begin{theorem}\label{thm:1.2.8}
对 $f \in \mathcal{S}(\mathbb{R})$ 成立下列性质:
\begin{enumerate}[(i)]
\item $|\mathcal{F}f(\xi)| \leq \| f \|_{L^{1}}$ $(\xi \in \mathbb{R})$,
\item $\mathcal{F}(f^{(k)})(\xi) = (i\xi)^{k}\mathcal{F}f(\xi)$ $(k \in \mathbb{N}_{0}, \xi \in \mathbb{R})$,
\item $(\mathcal{F}f)^{(k)}(\xi) = \mathcal{F}\bigl((-ix)^{k}f(x)\bigr)(\xi)$ $(k \in \mathbb{N}_{0}, \xi \in \mathbb{R})$,
\item $\mathcal{F}\bigl(f(x)e^{ixh}\bigr)(\xi) = \mathcal{F}f(\xi - h)$ $(h \in \mathbb{R}, \xi \in \mathbb{R})$,
\item $\mathcal{F}\bigl(f(x - h)\bigr)(\xi) = e^{-ih\xi}\mathcal{F}f(\xi)$ $(h \in \mathbb{R}, \xi \in \mathbb{R})$,
\item $\mathcal{F}(f(tx))(\xi) = \frac{1}{|t|}\mathcal{F}f\left(\frac{\xi}{t}\right)$ $(t \in \mathbb{R}\setminus\{0\}, \xi \in \mathbb{R})$,
\item $\mathcal{F}f \in \mathcal{S}(\mathbb{R})$.
\end{enumerate}
\end{theorem}

\begin{proof}
(i)
\begin{equation*}
|\mathcal{F}f(\xi)| = \left|\int_{-\infty}^{\infty} f(x)e^{-ix\xi}\dd x\right| \leq \int_{-\infty}^{\infty}|f(x)|\dd x = \| f \|_{L^{1}}.
\end{equation*}

(ii) 使用引理~\ref{lem:1.2.7},
\begin{align*}
\mathcal{F}(f')(\xi) & = \int_{-\infty}^{\infty} e^{-ix\xi}f'(x)\dd x = -\int_{-\infty}^{\infty}\left(\frac{\dd}{\dd x}e^{-ix\xi}\right)f(x)\dd x \\
& = i\xi\int_{-\infty}^{\infty} e^{-ix\xi}f(x)\dd x = i\xi\mathcal{F}f(\xi).
\end{align*}
以下对一般的 $k$ 用归纳法即可证明.

(iii)
\begin{equation*}
\left|e^{-ix\xi}f(x)\right| = |f(x)| \in L^{1}(\mathbb{R}),
\end{equation*}
\begin{equation*}
\left|\frac{\partial}{\partial \xi}e^{-ix\xi}f(x)\right| = |-ix\,e^{-ix\xi}f(x)| = |xf(x)| \leq \frac{C_{3,0}|x|}{(1 + |x|)^{3}} \in L^{1}(\mathbb{R})
\end{equation*}
所以可用积分号下的微分定理~\ref{thm:1.1.4},
\begin{align*}
(\mathcal{F}f)'(\xi) & = \frac{\dd}{\dd \xi}\left(\int_{-\infty}^{\infty} e^{-ix\xi}f(x)\dd x\right) = \int_{-\infty}^{\infty}\frac{\partial}{\partial \xi}\bigl(e^{-ix\xi}f(x)\bigr)\dd x \\
& = \int_{-\infty}^{\infty} -ix\,e^{-ix\xi}f(x)\dd x = \mathcal{F}\bigl(-ixf(x)\bigr)(\xi).
\end{align*}
以下对一般的 $k$ 用归纳法即可证明.

(iv)
\begin{align*}
\mathcal{F}\bigl(f(x)e^{ixh}\bigr)(\xi) & = \int_{-\infty}^{\infty} e^{-ix\xi}f(x)e^{ixh}\dd x \\
& = \int_{-\infty}^{\infty} e^{-ix(\xi - h)}f(x)\dd x = \mathcal{F}f(\xi - h).
\end{align*}

(v) 使用变量替换 $y = x - h$,
\begin{align*}
\mathcal{F}\bigl(f(x - h)\bigr)(\xi) & = \int_{-\infty}^{\infty} e^{-ix\xi}f(x - h)\dd x = \int_{-\infty}^{\infty} e^{-i(y+h)\xi}f(y)\dd y \\
& = e^{-ih\xi}\int_{-\infty}^{\infty} e^{-ix\xi}f(x)\dd x = e^{-ih\xi}\mathcal{F}f(\xi).
\end{align*}

(vi) 与 (v) 相同,使用变量替换 $y = tx$,
\begin{align*}
\mathcal{F}(f(tx))(\xi) & = \int_{-\infty}^{\infty} e^{-ix\xi}f(tx)\dd x = \int_{-\infty}^{\infty} e^{-i\frac{y}{t}\xi}f(y)\frac{\dd y}{|t|} \\
& = \frac{1}{|t|}\int_{-\infty}^{\infty} e^{-iy\frac{\xi}{t}}f(y)\dd y = \frac{1}{|t|}\mathcal{F}f\left(\frac{\xi}{t}\right).
\end{align*}

(vii) 由 (iii) 可知 $\mathcal{F}f \in C^{\infty}(\mathbb{R})$.又,对任意 $k, \ell \in \mathbb{N}_{0}$,由 (ii),(iii),
\begin{equation*}
(i\xi)^{k}(\mathcal{F}f)^{(\ell)}(\xi) = (i\xi)^{k}\mathcal{F}\bigl((-ix)^{\ell}f(x)\bigr)(\xi) = \mathcal{F}\bigl(\bigl[(-ix)^{\ell}f(x)\bigr]^{(k)}\bigr)(\xi).
\end{equation*}
因此,使用 (i) 与 Leibniz 公式,并注意 $\mathcal{S}(\mathbb{R})$ 的定义,则
\begin{align*}
|(i\xi)^{k}(\mathcal{F}f)^{(\ell)}(\xi)|&\leq \left\| \bigl[(-ix)^{\ell}f(x)\bigr]^{(k)} \right\|_{L^{1}} = \left\| \sum_{j=0}^{k}\binom{k}{j}\bigl[(-ix)^{\ell}\bigr]^{(j)}f^{(k-j)} \right\|_{L^{1}}  \\
&\leq C\left\| (1 + |x|)^{\ell}\sum_{j=0}^{k}|f^{(j)}(x)| \right\|_{L^{1}} \leq C\left\| (1 + |x|)^{\ell}\sum_{j=0}^{k}\frac{C_{\ell+2,j}}{(1 + |x|)^{\ell+2}} \right\|_{L^{1}} \\
&\leq C\sum_{j=0}^{k}C_{\ell+2,j}\left\| \frac{1}{(1 + |x|)^{2}} \right\|_{L^{1}} < \infty
\end{align*}
由此证明了 $\mathcal{F}f \in \mathcal{S}(\mathbb{R})$.
\end{proof}

\begin{theorem}\label{thm:1.2.9}
(i) $\mathcal{F}^{-1}\mathcal{F}f = f$,$\mathcal{F}\mathcal{F}^{-1}f = f$ $(f \in \mathcal{S}(\mathbb{R}))$.

(ii) $\int_{-\infty}^{\infty} f(x)\overline{g(x)}\dd x = \frac{1}{2\pi}\int_{-\infty}^{\infty}\mathcal{F}f(\xi)\overline{\mathcal{F}g(\xi)}\dd \xi$ $(f, g \in \mathcal{S}(\mathbb{R}))$(Parseval 等式).

(iii) $\int_{-\infty}^{\infty}|f(x)|^{2}\dd x = \frac{1}{2\pi}\int_{-\infty}^{\infty}|\mathcal{F}f(\xi)|^{2}\dd \xi$ $(f \in \mathcal{S}(\mathbb{R}))$(Plancherel 等式).
\end{theorem}

\begin{proof}
(i) 由于 $f \in \mathcal{S}(\mathbb{R})$,由定理~\ref{thm:1.2.8}(vii) 有 $\mathcal{F}f \in \mathcal{S}(\mathbb{R})$;又因 $\mathcal{S}(\mathbb{R}) \subset L^{1}(\mathbb{R})$,由定理~\ref{thm:1.2.5},
\begin{equation*}
f(x) = \mathcal{F}^{-1}(\mathcal{F}f)(x) \quad (\text{a.e. } x \in \mathbb{R}).
\end{equation*}
但是,$f$ 连续,且 $\mathcal{F}f \in \mathcal{S}(\mathbb{R}) \subset L^{1}(\mathbb{R})$,故由定理~\ref{thm:1.2.3},$\mathcal{F}^{-1}(\mathcal{F}f) \in C(\mathbb{R})$,因此
\begin{equation*}
f(x) = \mathcal{F}^{-1}(\mathcal{F}f)(x) \quad (x \in \mathbb{R}).
\end{equation*}
此外,
\begin{align*}
\mathcal{F}(\mathcal{F}^{-1}f)(x) & = \mathcal{F}\left(\frac{1}{2\pi}\mathcal{F}f(-\xi)\right)(x) = \frac{1}{2\pi}\mathcal{F}\bigl(\mathcal{F}f(\xi)\bigr)(-x) \\
& = \mathcal{F}^{-1}(\mathcal{F}f)(x) = f(x).
\end{align*}

(ii) 由于 $f, g \in \mathcal{S}(\mathbb{R})$,它们满足定理~\ref{thm:1.2.5} 证明中关于 $f$,$g$ 的条件,故在式~\eqref{eq:1.2.5} 中令 $x = 0$,并利用 $\mathcal{F}^{-1}(\mathcal{F}g) = g$,得
\begin{equation*}
\int\mathcal{F}f(\xi)\mathcal{F}g(\xi)\dd \xi = 2\pi\int f(y)\mathcal{F}^{-1}(\mathcal{F}g)(-y)\dd y = 2\pi\int f(y)g(-y)\dd y.
\end{equation*}
因此
\begin{equation*}
\int f(y)\overline{g(y)}\dd y = \int f(y)\overline{g(-(-y))}\dd y = \frac{1}{2\pi}\int\mathcal{F}f(\xi)\mathcal{F}\bigl(\overline{g(-x)}\bigr)(\xi)\dd \xi.
\end{equation*}
这里,由 Fourier 变换的定义与变量替换 $y = -x$,
\begin{align*}
\mathcal{F}\bigl(\overline{g(-x)}\bigr)(\xi) &= \int e^{-ix\xi}\overline{g(-x)}\dd x = \int e^{iy\xi}\overline{g(y)}\dd y = \overline{\int e^{-iy\xi}g(y)\dd y} \\
&= \overline{\mathcal{F}g(\xi)}.
\end{align*}
因此
\begin{equation*}
\int f(y)\overline{g(y)}\dd y = \frac{1}{2\pi}\int\mathcal{F}f(\xi)\overline{\mathcal{F}g(\xi)}\dd \xi
\end{equation*}
成立.

(iii) 在 (ii) 中令 $g = f$ 即可.
\end{proof}

\begin{definition}\label{def:1.2.10}
对 $f, g \in L^{1}(\mathbb{R})$,
\begin{equation*}
f * g(x) := \int f(y)g(x - y)\dd y = \int f(x - y)g(y)\dd y
\end{equation*}
称为 $f$ 与 $g$ 的卷积(convolution).
\end{definition}

\begin{lemma}\label{lem:1.2.11}
若 $f, g \in L^{1}(\mathbb{R})$,则
\begin{equation*}
\mathcal{F}(f * g)(\xi) = \mathcal{F}f(\xi)\mathcal{F}g(\xi).
\end{equation*}
\end{lemma}

\begin{proof}
设 $f, g \in L^{1}(\mathbb{R})$.由于
\begin{equation*}
\iint |e^{-ix\xi}f(x - y)g(y)|\dd x\dd y = \int\left(\int |f(x - y)|\dd x\right)|g(y)|\dd y = \int |f(x)|\dd x\int |g(y)|\dd y < \infty,
\end{equation*}
使用 Fubini 定理后,作变量替换 $z = x - y$,得
\begin{align*}
\mathcal{F}(f * g)(\xi) & = \int e^{-ix\xi}\left(\int f(x - y)g(y)\dd y\right)\dd x \\
& = \int\left(\int f(x - y)e^{-ix\xi}\dd x\right)g(y)\dd y \\
& = \int\left(\int f(x - y)e^{-i(x-y)\xi}\dd x\right)e^{-iy\xi}g(y)\dd y
\end{align*}
\begin{align*}
&= \int\left(\int f(z)e^{-iz\xi}\dd z\right)e^{-iy\xi}g(y)\dd y \\
&= \int\mathcal{F}f(\xi)e^{-iy\xi}g(y)\dd y = \mathcal{F}f(\xi)\mathcal{F}g(\xi).
\end{align*}
\end{proof}

\begin{theorem}\label{thm:1.2.12}
对 $f, g \in \mathcal{S}(\mathbb{R})$,
\begin{enumerate}[(i)]
\item $\mathcal{F}(f * g)(\xi) = (\mathcal{F}f)(\xi)(\mathcal{F}g)(\xi)$,
\begin{equation*}
\mathcal{F}^{-1}(f * g)(x) = 2\pi(\mathcal{F}^{-1}f)(x)(\mathcal{F}^{-1}g)(x);
\end{equation*}
\item $\mathcal{F}(fg)(\xi) = \frac{1}{2\pi}\mathcal{F}f * \mathcal{F}g(\xi)$,$\mathcal{F}^{-1}(fg)(x) = \mathcal{F}^{-1}f * \mathcal{F}^{-1}g(x)$.
\end{enumerate}
\end{theorem}

\begin{proof}
设 $f, g \in \mathcal{S}(\mathbb{R})$.

(i) 由于 $f, g \in L^{1}(\mathbb{R})$,由引理~\ref{lem:1.2.11} 得第一个等式.

再者,因 $\mathcal{F}^{-1}f(x) = \frac{1}{2\pi}\mathcal{F}f(-x)$,故
\begin{align*}
\mathcal{F}^{-1}(f * g)(x) & = \frac{1}{2\pi}\mathcal{F}(f * g)(-x) \\
& = \frac{1}{2\pi}\mathcal{F}f(-x)\mathcal{F}g(-x) = 2\pi(\mathcal{F}^{-1}f)(x)(\mathcal{F}^{-1}g)(x).
\end{align*}

(ii) 由于 $\mathcal{F}f, \mathcal{F}g \in \mathcal{S}(\mathbb{R})$,由 (i),
\begin{equation*}
\mathcal{F}^{-1}(\mathcal{F}f * \mathcal{F}g)(x) = 2\pi\mathcal{F}^{-1}\mathcal{F}f(x)\mathcal{F}^{-1}\mathcal{F}g(x) = 2\pi f(x)g(x).
\end{equation*}
又因 $f, g \in \mathcal{S}(\mathbb{R})$,故 $\mathcal{F}f, \mathcal{F}g \in L^{1}(\mathbb{R})$,$fg \in \mathcal{S}(\mathbb{R}) \subset L^{1}(\mathbb{R})$;进而 $\mathcal{F}f * \mathcal{F}g \in C(\mathbb{R})$ 也容易看出.因此,如 (i) 中所用,若 $h$,$\mathcal{F}h \in L^{1}(\mathbb{R})$ 则 $\mathcal{F}\mathcal{F}^{-1}h = h$(a.e.),于是
\begin{equation*}
\mathcal{F}(fg)(\xi) = \frac{1}{2\pi}\mathcal{F}\bigl(\mathcal{F}^{-1}(\mathcal{F}f * \mathcal{F}g)\bigr)(\xi) = \frac{1}{2\pi}(\mathcal{F}f * \mathcal{F}g)(\xi)
\end{equation*}
同样地,
\begin{align*}
\mathcal{F}^{-1}(fg)(x) & = \mathcal{F}^{-1}\bigl((\mathcal{F}\mathcal{F}^{-1}f)(\mathcal{F}\mathcal{F}^{-1}g)\bigr)(x) = \mathcal{F}^{-1}\bigl(\mathcal{F}(\mathcal{F}^{-1}f * \mathcal{F}^{-1}g)\bigr)(x) \\
& = \bigl(\mathcal{F}^{-1}\mathcal{F}(\mathcal{F}^{-1}f * \mathcal{F}^{-1}g)\bigr)(x) = \mathcal{F}^{-1}f * \mathcal{F}^{-1}g(x).
\end{align*}
\end{proof}

\subsection{\texorpdfstring{$L^2$ 函数的 Fourier 变换}{L2 函数的 Fourier 变换}}\label{sec:1.2.2}

\begin{lemma}\label{lem:1.2.13}
若 $f \in L^{2}(\mathbb{R})$,则对一切满足 $\lim_{j \to \infty}\|f_{j} - f\|_{L^{2}} = 0$ 的 $f_{j} \in \mathcal{S}(\mathbb{R})$($j = 1, 2, \ldots$),存在唯一的 $g \in L^{2}(\mathbb{R})$ 使得 $\lim_{j \to \infty}\|\mathcal{F}f_{j} - g\|_{L^{2}} = 0$.
\end{lemma}

\begin{proof}
设 $f_{j} \in \mathcal{S}(\mathbb{R})$ 且 $\lim_{j \to \infty}\|f_{j} - f\|_{L^{2}} = 0$.由对 $\mathcal{S}(\mathbb{R})$ 的 Plancherel 等式(定理~\ref{thm:1.2.9}(iii)),
\begin{align*}
\lim_{j, \ell \to \infty}\|\mathcal{F}f_{j} - \mathcal{F}f_{\ell}\|_{L^{2}} & = \lim_{j, \ell \to \infty}\sqrt{2\pi}\,\|f_{j} - f_{\ell}\|_{L^{2}} \\
& \leq \lim_{j \to \infty}\sqrt{2\pi}\,\|f_{j} - f\|_{L^{2}} + \lim_{\ell \to \infty}\sqrt{2\pi}\,\|f - f_{\ell}\|_{L^{2}} \\
& = 0.
\end{align*}
因此,函数列 $\{\mathcal{F}f_{j}\}_{j=1}^{\infty}$ 是 $L^{2}(\mathbb{R})$ 中的 Cauchy 列.由于 $L^{2}(\mathbb{R})$ 完备,存在 $g \in L^{2}(\mathbb{R})$ 使得
\begin{equation*}
\lim_{j \to \infty}\|\mathcal{F}f_{j} - g\|_{L^{2}} = 0.
\end{equation*}

其次,设 $h_{j} \in \mathcal{S}(\mathbb{R})$ 且 $\lim_{j \to \infty}\|h_{j} - f\|_{L^{2}} = 0$.与上面相同,存在 $k \in L^{2}(\mathbb{R})$ 使得 $\lim_{j \to \infty}\|\mathcal{F}h_{j} - k\|_{L^{2}} = 0$.于是,由对 $\mathcal{S}(\mathbb{R})$ 的 Plancherel 等式,
\begin{align*}
\| g - k \|_{L^{2}} & = \lim_{j \to \infty}\|\mathcal{F}f_{j} - \mathcal{F}h_{j}\|_{L^{2}} = \lim_{j \to \infty}\sqrt{2\pi}\,\|f_{j} - h_{j}\|_{L^{2}} \\
& \leq \lim_{j \to \infty}\sqrt{2\pi}\,\|f_{j} - f\|_{L^{2}} + \lim_{j \to \infty}\sqrt{2\pi}\,\|f - h_{j}\|_{L^{2}} = 0.
\end{align*}
因此,
\begin{equation*}
g(x) = k(x) \quad (\text{a.e. } x \in \mathbb{R}).
\end{equation*}
\end{proof}

由于 $\mathcal{S}(\mathbb{R})$ 在 $L^{2}(\mathbb{R})$ 中稠密,对 $f \in L^{2}(\mathbb{R})$ 可如下定义 $f$ 的 Fourier 变换.

\begin{definition}\label{def:1.2.14}
对 $f \in L^{2}(\mathbb{R})$,取满足 $\lim_{j \to \infty}\|f_{j} - f\|_{L^{2}} = 0$ 的 $f_{j} \in \mathcal{S}(\mathbb{R})$,把满足 $\lim_{j \to \infty}\|\mathcal{F}f_{j} - g\|_{L^{2}} = 0$ 的 $g \in L^{2}(\mathbb{R})$ 称为 $f$ 的(在 $L^{2}(\mathbb{R})$ 意义下的)Fourier 变换,记作 $\mathcal{F}_{L^{2}}f$.又,把满足 $\lim_{j \to \infty}\|\mathcal{F}^{-1}f_{j} - g\|_{L^{2}} = 0$ 的 $g \in L^{2}(\mathbb{R})$ 称为 $f$ 的(在 $L^{2}(\mathbb{R})$ 意义下的)逆 Fourier 变换,记作 $\mathcal{F}_{L^{2}}^{-1}f$.
\end{definition}

\begin{theorem}\label{thm:1.2.15}
成立下列性质.
\begin{enumerate}[(i)]
\item $\mathcal{F}_{L^{2}}: L^{2}(\mathbb{R}) \to L^{2}(\mathbb{R})$ 是从 $L^{2}(\mathbb{R})$ 映到 $L^{2}(\mathbb{R})$ 上的一一且连续的线性映射.
\item $\mathcal{F}_{L^{2}}^{-1}: L^{2}(\mathbb{R}) \to L^{2}(\mathbb{R})$ 也是从 $L^{2}(\mathbb{R})$ 映到 $L^{2}(\mathbb{R})$ 上的一一且连续的线性映射.
\item $\mathcal{F}_{L^{2}}^{-1}\mathcal{F}_{L^{2}}f = f$,$\mathcal{F}_{L^{2}}\mathcal{F}_{L^{2}}^{-1}f = f$ $(f \in L^{2}(\mathbb{R}))$.
\item $\int_{-\infty}^{\infty} f(x)\overline{g(x)}\dd x = \frac{1}{2\pi}\int_{-\infty}^{\infty}\mathcal{F}_{L^{2}}f(\xi)\overline{\mathcal{F}_{L^{2}}g(\xi)}\dd \xi$ $(f, g \in L^{2}(\mathbb{R}))$(Parseval 等式).
\item $\int_{-\infty}^{\infty}|f(x)|^{2}\dd x = \frac{1}{2\pi}\int_{-\infty}^{\infty}|\mathcal{F}_{L^{2}}f(\xi)|^{2}\dd \xi$ $(f \in L^{2}(\mathbb{R}))$(Plancherel 等式).
\end{enumerate}
\end{theorem}

\begin{proof}
(iii) 首先确认 $\mathcal{F}_{L^{2}}^{-1}\mathcal{F}_{L^{2}}f = f$.设 $f \in L^{2}(\mathbb{R})$,取满足 $\lim_{j \to \infty}\|f_{j} - f\|_{L^{2}} = 0$ 的 $f_{j} \in \mathcal{S}(\mathbb{R})$.由 $L^{2}$ 函数 Fourier 变换的定义,
\begin{equation*}
\lim_{j \to \infty}\|\mathcal{F}f_{j} - \mathcal{F}_{L^{2}}f\|_{L^{2}} = 0
\end{equation*}
因此,再由 $L^{2}$ 函数逆 Fourier 变换的定义,
\begin{equation*}
\lim_{j \to \infty}\|\mathcal{F}^{-1}\mathcal{F}f_{j} - \mathcal{F}_{L^{2}}^{-1}\mathcal{F}_{L^{2}}f\|_{L^{2}} = 0.
\end{equation*}
但是,由于 $f_{j} \in \mathcal{S}(\mathbb{R})$,由定理~\ref{thm:1.2.9}(i),$\mathcal{F}^{-1}\mathcal{F}f_{j} = f_{j}$.因此
\begin{equation*}
f(x) = \mathcal{F}_{L^{2}}^{-1}\mathcal{F}_{L^{2}}f(x) \quad (\text{a.e. } x \in \mathbb{R}),
\end{equation*}
即 $\mathcal{F}_{L^{2}}^{-1}\mathcal{F}_{L^{2}}f = f$.同样地,可证 $\mathcal{F}_{L^{2}}\mathcal{F}_{L^{2}}^{-1}f = f$.

(iv) 设 $f, g \in L^{2}(\mathbb{R})$,取满足 $\lim_{j \to \infty}\|f_{j} - f\|_{L^{2}} = 0$,$\lim_{j \to \infty}\|g_{j} - g\|_{L^{2}} = 0$ 的 $f_{j}, g_{j} \in \mathcal{S}(\mathbb{R})$.由于 $\|f\|_{L^{2}} - \|f_{j} - f\|_{L^{2}} \leq \|f_{j}\|_{L^{2}} \leq \|f_{j} - f\|_{L^{2}} + \|f\|_{L^{2}}$,故 $\lim_{j \to \infty}\|f_{j}\|_{L^{2}} = \|f\|_{L^{2}}$.因此,使用 Cauchy--Schwarz 不等式,
\begin{align*}
& \left|\int_{\mathbb{R}} f(x)\overline{g(x)}\dd x - \int_{\mathbb{R}} f_{j}(x)\overline{g_{j}(x)}\dd x\right| \\
& \quad \leq \left|\int_{\mathbb{R}} (f(x) - f_{j}(x))\overline{g(x)}\dd x\right| + \left|\int_{\mathbb{R}} f_{j}(x)\overline{(g(x) - g_{j}(x))}\dd x\right| \\
& \quad \leq \| f - f_{j} \|_{L^{2}}\| g \|_{L^{2}} + \| f_{j} \|_{L^{2}}\| g - g_{j} \|_{L^{2}}
\end{align*}
令 $j \to \infty$,得
\begin{equation*}
\lim_{j \to \infty}\int_{\mathbb{R}} f_{j}(x)\overline{g_{j}(x)}\dd x = \int_{\mathbb{R}} f(x)\overline{g(x)}\dd x.
\end{equation*}
同样地,由引理~\ref{lem:1.2.13},$\lim_{j \to \infty}\|\mathcal{F}f_{j} - \mathcal{F}_{L^{2}}f\|_{L^{2}} = 0$,$\lim_{j \to \infty}\|\mathcal{F}g_{j} - \mathcal{F}_{L^{2}}g\|_{L^{2}} = 0$,所以
\begin{equation*}
\int_{\mathbb{R}}\mathcal{F}_{L^{2}}f(\xi)\overline{\mathcal{F}_{L^{2}}g(\xi)}\dd \xi = \lim_{j \to \infty}\int_{\mathbb{R}}\mathcal{F}f_{j}(\xi)\overline{\mathcal{F}g_{j}(\xi)}\dd \xi.
\end{equation*}
由于 $f_{j}, g_{j} \in \mathcal{S}(\mathbb{R})$,由对 $\mathcal{S}(\mathbb{R})$ 的 Plancherel 等式(定理~\ref{thm:1.2.9}(iii)),
\begin{align*}
\int_{\mathbb{R}} f(x)\overline{g(x)}\dd x & = \lim_{j \to \infty}\int_{\mathbb{R}} f_{j}(x)\overline{g_{j}(x)}\dd x \\
& = \lim_{j \to \infty}\frac{1}{2\pi}\int_{\mathbb{R}}\mathcal{F}f_{j}(\xi)\overline{\mathcal{F}g_{j}(\xi)}\dd \xi \\
& = \frac{1}{2\pi}\int_{\mathbb{R}}\mathcal{F}_{L^{2}}f(\xi)\overline{\mathcal{F}_{L^{2}}g(\xi)}\dd \xi.
\end{align*}

(v) 在 (iv) 中令 $g(x) = f(x)$ 即可.

(i) $\mathcal{F}_{L^{2}}$ 的线性是显然的.连续性由 (v) 得出.由 (iii) 可知它是从 $L^{2}(\mathbb{R})$ 映到 $L^{2}(\mathbb{R})$ 上的一一对应.
\end{proof}

\begin{theorem}\label{thm:1.2.16}
(i) 若 $f \in L^{1}(\mathbb{R}) \cap L^{2}(\mathbb{R})$,则 $\mathcal{F}f(\xi) = \mathcal{F}_{L^{2}}f(\xi)$(a.e. $\xi \in \mathbb{R}$).因此,自本定理的证明以后,$\mathcal{F}_{L^{2}}$ 也记作 $\mathcal{F}$.

(ii) 若 $f \in L^{2}(\mathbb{R})$,则
\begin{equation*}
\lim_{A \to \infty}\left\| \int_{-A}^{A} e^{-ix\xi}f(x)\dd x - \mathcal{F}_{L^{2}}f(\xi) \right\|_{L^{2}} = 0.
\end{equation*}
\end{theorem}

\begin{proof}
(i) 设 $f \in L^{1}(\mathbb{R}) \cap L^{2}(\mathbb{R})$.
\begin{equation*}
f_{j}(x) = \left\{ \begin{array}{l l}
f(x) & (|x| \leq j) \\
0 & (|x| > j)
\end{array} \right.
\end{equation*}
则
\begin{equation*}
\lim_{j \to \infty}\|f_{j} - f\|_{L^{1}} = 0, \quad \lim_{j \to \infty}\|f_{j} - f\|_{L^{2}} = 0.
\end{equation*}
这里,取一个满足 $\int_{\mathbb{R}}\varphi(x)\dd x = 1$ 的 $\varphi \in C_{c}^{\infty}(\mathbb{R})$.由例~\ref{exa:1.1.10},对每个 $j \in \mathbb{N}$,存在 $t_{j} > 0$ 使得
\begin{equation*}
\| f_{j} - \varphi_{t_{j}} * f_{j} \|_{L^{1}} \leq \frac{1}{j}, \quad \| f_{j} - \varphi_{t_{j}} * f_{j} \|_{L^{2}} \leq \frac{1}{j}.
\end{equation*}
令 $g_{j} = \varphi_{t_{j}} * f_{j}$,则 $g_{j} \in C_{c}^{\infty}(\mathbb{R}) \subset \mathcal{S}(\mathbb{R})$,且
\begin{equation*}
\| g_{j} - f \|_{L^{1}} \leq \| g_{j} - f_{j} \|_{L^{1}} + \| f_{j} - f \|_{L^{1}} \leq \frac{1}{j} + \| f_{j} - f \|_{L^{1}}
\end{equation*}
\begin{equation*}
\| g_{j} - f \|_{L^{2}} \leq \| g_{j} - f_{j} \|_{L^{2}} + \| f_{j} - f \|_{L^{2}} \leq \frac{1}{j} + \| f_{j} - f \|_{L^{2}}
\end{equation*}
故
\begin{equation*}
\lim_{j \to \infty}\| g_{j} - f \|_{L^{1}} = 0, \quad \lim_{j \to \infty}\| g_{j} - f \|_{L^{2}} = 0.
\end{equation*}
由 $L^{1}$ 收敛性,根据定理~\ref{thm:1.2.3},对一切 $\xi \in \mathbb{R}$,
\begin{equation*}
\mathcal{F}g_{j}(\xi) \to \mathcal{F}f(\xi) = \int_{-\infty}^{\infty} f(x)e^{ix\xi}\dd x \quad (j \to \infty)
\end{equation*}
又由 $L^{2}$ 收敛性,根据对 $L^{2}(\mathbb{R})$ 的 Plancherel 等式(定理~\ref{thm:1.2.15}(v)),
\begin{equation*}
\|\mathcal{F}g_{j} - \mathcal{F}_{L^{2}}f\|_{L^{2}} \to 0 \quad (j \to \infty)
\end{equation*}
所以
\begin{equation*}
\mathcal{F}_{L^{2}}f(\xi) = \mathcal{F}f(\xi) = \int_{-\infty}^{\infty} e^{-ix\xi}f(x)\dd x \quad (\text{a.e. } \xi \in \mathbb{R}).
\end{equation*}

(ii) 设 $f \in L^{2}(\mathbb{R})$.
\begin{equation*}
f_{A}(x) = \left\{ \begin{array}{l l}
f(x) & (|x| \leq A) \\
0 & (|x| > A)
\end{array} \right.
\end{equation*}
则
\begin{equation*}
f_{A} \in L^{1}(\mathbb{R}) \cap L^{2}(\mathbb{R}) \quad \text{且} \quad \lim_{A \to \infty}\|f_{A} - f\|_{L^{2}} = 0.
\end{equation*}
因此,由 (i),
\begin{equation*}
\mathcal{F}_{L^{2}}f_{A}(\xi) = \mathcal{F}f_{A}(\xi) = \int_{-\infty}^{\infty} f_{A}(x)e^{-ix\xi}\dd x = \int_{-A}^{A} e^{-ix\xi}f(x)\dd x.
\end{equation*}
另一方面,由于 $\lim_{A \to \infty}\|f_{A} - f\|_{L^{2}} = 0$,根据对 $L^{2}(\mathbb{R})$ 的 Plancherel 等式(定理~\ref{thm:1.2.15}(v)),
\begin{equation*}
\lim_{A \to \infty}\|\mathcal{F}_{L^{2}}f_{A} - \mathcal{F}_{L^{2}}f\|_{L^{2}} = 0.
\end{equation*}
因此
\begin{equation*}
\lim_{A \to \infty}\left\| \int_{-A}^{A} e^{-ix\xi}f(x)\dd x - \mathcal{F}_{L^{2}}f(\xi) \right\|_{L^{2}} = 0.
\end{equation*}
\end{proof}

在本节的最后,请注意以下事实.由定理~\ref{thm:1.2.3}(i),Fourier 变换是从 $L^{1}(\mathbb{R})$ 到 $L^{\infty}(\mathbb{R})$ 的有界线性算子;由定理~\ref{thm:1.2.15},它又是从 $L^{2}(\mathbb{R})$ 到 $L^{2}(\mathbb{R})$ 的有界线性算子.由这两个事实可知,当 $1 \leq p \leq 2$ 时,Fourier 变换是从 $L^{p}(\mathbb{R})$ 到 $L^{p'}(\mathbb{R})$ 的有界线性算子(其中 $1/p + 1/p' = 1$).这一点由下面的 Riesz--Thorin 插值定理来保证.

\begin{theorem}[M. Riesz--Thorin 插值定理]\label{thm:1.2.17}
设 $(X, \mu)$,$(Y, \nu)$ 是测度空间,$S(X)$,$S(Y)$ 分别是 $X$ 上($Y$ 上)的 $\mu$-($\nu$-)可积简单函数的全体,$L(Y)$ 是 $\nu$-可测函数的全体.

设 $T$ 是从 $S(X)$ 到 $L(Y)$ 的线性算子,即满足
\begin{equation}\tag{1.2.9}
T(af + bg) = aT(f) + bT(g) \quad (a, b \in \mathbb{C}, \ f, g \in S(X)).
\label{eq:1.2.9}\end{equation}
进一步,设 $1 \leq p_{1}, p_{2} \leq \infty$,$1 \leq q_{1}, q_{2} \leq \infty$,且
\begin{equation}\tag{1.2.10}
\| Tf \|_{L^{q_{1}}(Y, \nu)} \leq M_{1}\| f \|_{L^{p_{1}}(X, \mu)} \quad (f \in S(X)),
\label{eq:1.2.10}\end{equation}
\begin{equation}\tag{1.2.11}
\| Tf \|_{L^{q_{2}}(Y, \nu)} \leq M_{2}\| f \|_{L^{p_{2}}(X, \mu)} \quad (f \in S(X)).
\label{eq:1.2.11}\end{equation}
此时,对 $0 \leq t \leq 1$,由 $1/p = (1 - t)/p_{1} + t/p_{2}$,$1/q = (1 - t)/q_{1} + t/q_{2}$ 定义 $p$,$q$,则
\begin{equation}\tag{1.2.12}
\| Tf \|_{L^{q}(Y, \nu)} \leq M_{1}^{1-t}M_{2}^{t}\| f \|_{L^{p}(X, \mu)} \quad (f \in S(X)).
\label{eq:1.2.12}\end{equation}

特别地,当 $1 \leq p < \infty$ 时,$T$ 可唯一扩张为从 $L^{p}(X, \mu)$ 到 $L^{q}(Y, \nu)$ 的有界线性算子.
\end{theorem}

\begin{proof}
当 $t = 0$ 或 $t = 1$ 时无须证明,故设 $0 < t < 1$,并对 $z \in \mathbb{C}$ 令
\begin{equation*}
\alpha(z) = \frac{1 - z}{p_{1}} + \frac{z}{p_{2}}, \quad \beta(z) = \frac{1 - z}{q_{1}} + \frac{z}{q_{2}}.
\end{equation*}
则 $\alpha(0) = 1/p_{1}$,$\alpha(1) = 1/p_{2}$,$\alpha(t) = 1/p$;$\beta(0) = 1/q_{1}$,$\beta(1) = 1/q_{2}$,$\beta(t) = 1/q$.令 $B = \{z \in \mathbb{C};\ 0 < \operatorname{Re} z < 1\}$,$\overline{B} = \{z \in \mathbb{C};\ 0 \leq \operatorname{Re} z \leq 1\}$.对 $f \in S(X)$,
\begin{equation*}
\| Tf \|_{L^{q}(Y, \nu)} = \sup_{g \in S(Y),\, \| g \|_{L^{q'}(Y, \nu)} = 1}\left|\int_{Y} Tf(y)g(y)\dd \nu(y)\right|
\end{equation*}
所以,只需对满足 $\|f\|_{L^{p}(X, \mu)} = 1$ 的 $f \in S(X)$ 与满足 $\|g\|_{L^{q'}(Y, \nu)} = 1$ 的 $g \in S(Y)$ 证明
\begin{equation*}
\left|\int_{Y} Tf(y)g(y)\dd \nu(y)\right| \leq M_{1}^{1-t}M_{2}^{t}.
\end{equation*}

(イ)首先处理 $1 \leq p < \infty$,$1 < q \leq \infty$ 的情形.可设 $f$,$g$ 都不是零函数,故 $f$,$g$ 可设为如下形状:
\begin{equation*}
f(x) = \sum_{j=1}^{\ell} a_{j}\chi_{E_{j}}(x), \quad g(y) = \sum_{k=1}^{m} b_{k}\chi_{F_{k}}(y).
\end{equation*}
其中,$a_{j}, b_{k} \in \mathbb{C}\setminus\{0\}$;$E_{j}$ 是满足 $0 < \mu(E_{j}) < \infty$ 的 $\mu$ 可测集,$F_{k}$ 是满足 $0 < \nu(F_{k}) < \infty$ 的 $\nu$ 可测集;且 $E_{j}$ 两两互不相交,$F_{k}$ 也两两互不相交($j = 1, \ldots, \ell;\ k = 1, \ldots, m$).于是令
\begin{align}
F_{z}(x) & = \sum_{j=1}^{\ell}|a_{j}|^{p\alpha(z)}e^{i\arg a_{j}}\chi_{E_{j}}(x), \notag \\
G_{z}(y) & = \sum_{k=1}^{m}|b_{k}|^{q'(1 - \beta(z))}e^{i\arg b_{k}}\chi_{F_{k}}(y) \tag{1.2.13}\label{eq:1.2.13}
\end{align}
并考虑
\begin{equation*}
\Phi(z) = \int_{Y} T(F_{z})(y)G_{z}(y)\dd \nu(y).
\end{equation*}
由于 $F_{t}(x) = f(x)$,$G_{t}(y) = g(y)$,
\begin{equation}\tag{1.2.14}
\Phi(t) = \int_{Y} T(F_{t})(y)G_{t}(y)\dd \nu(y) = \int_{Y} T(f)(y)g(y)\dd \nu(y).
\label{eq:1.2.14}\end{equation}
进一步,由 $T$ 与积分的线性,
\begin{equation*}
\Phi(z) = \sum_{j=1}^{\ell}\sum_{k=1}^{m}|a_{j}|^{p\alpha(z)}e^{i\arg a_{j}}|b_{k}|^{q'(1 - \beta(z))}e^{i\arg b_{k}}\int_{Y} T(\chi_{E_{j}})(y)\chi_{F_{k}}(y)\dd \nu(y)
\end{equation*}
且
\begin{align*}
\left|\int_{Y} T(\chi_{E_{j}})(y)\chi_{F_{k}}(y)\dd \nu(y)\right| & \leq \| T(\chi_{E_{j}}) \|_{L^{q_{1}}(Y, \nu)}\| \chi_{F_{k}} \|_{L^{q_{1}'}(Y, \nu)} \\
& \leq M_{1}\| \chi_{E_{j}} \|_{L^{p_{1}}(X, \mu)}\| \chi_{F_{k}} \|_{L^{q_{1}'}(Y, \nu)} \\
& \leq M_{1}(\mu(E_{j}))^{1/p_{1}}(\nu(F_{k}))^{1/q_{1}'} < \infty
\end{align*}
所以 $\Phi(z)$ 在 $\overline{B}$ 上有界连续,在 $B$ 内正则.因此,由复函数论中的 Hadamard 三线定理,
\begin{equation}\tag{1.2.15}
|\Phi(t)| \leq \left(\sup_{\operatorname{Re} z = 0}|\Phi(z)|\right)^{1-t}\left(\sup_{\operatorname{Re} z = 1}|\Phi(z)|\right)^{t}.
\label{eq:1.2.15}\end{equation}
现在,当 $\operatorname{Re} z = 0$ 时,
\begin{equation*}
|\Phi(z)| \leq \| T(F_{z}) \|_{L^{q_{1}}(Y, \nu)}\| G_{z} \|_{L^{q_{1}'}(Y, \nu)} \leq M_{1}\| F_{z} \|_{L^{p_{1}}(X, \mu)}\| G_{z} \|_{L^{q_{1}'}(Y, \nu)}
\end{equation*}
且
\begin{align*}
\| F_{z} \|_{L^{p_{1}}(X, \mu)} &= \left(\int_{X}\sum_{j=1}^{\ell}\left(|a_{j}|^{p/p_{1}}\right)^{p_{1}}\chi_{E_{j}}(x)\dd \mu(x)\right)^{1/p_{1}} \\
&= \left(\int_{X}|f(x)|^{p}\dd \mu(x)\right)^{1/p_{1}} = \left(\| f \|_{L^{p}(X, \mu)}\right)^{p/p_{1}} = 1^{p/p_{1}} = 1,
\end{align*}
\begin{align*}
\| G_{z} \|_{L^{q_{1}'}(Y, \nu)} & = \left(\int_{Y}\sum_{k=1}^{m}\left(|b_{k}|^{q'/q_{1}'}\right)^{q_{1}'}\chi_{F_{k}}(y)\dd \nu(y)\right)^{1/q_{1}'} \\
& = \left(\int_{Y}|g(y)|^{q'}\dd \nu(y)\right)^{1/q_{1}'} = \left(\| g \|_{L^{q'}(Y, \nu)}\right)^{q'/q_{1}'} = 1^{q'/q_{1}'} = 1
\end{align*}
所以
\begin{equation*}
|\Phi(z)| \leq M_{1} \quad (\operatorname{Re} z = 0).
\end{equation*}
同样地,也能证明
\begin{equation*}
|\Phi(z)| \leq M_{2} \quad (\operatorname{Re} z = 1),
\end{equation*}
故由 \eqref{eq:1.2.14},\eqref{eq:1.2.15},
\begin{equation*}
\left|\int_{Y} T(f)(y)g(y)\dd \nu(y)\right| = |\Phi(t)| \leq M_{1}^{1-t}M_{2}^{t},
\end{equation*}
从而 $1 \leq p < \infty$,$1 < q \leq \infty$ 的情形得证.

(ロ)$q = 1$ 的情形.此时 $q_{1} = q_{2} = q = 1$.若 $p_{1} = p_{2} = \infty$,则 $p = \infty$,定理以 $M_{1} = M_{2}$ 成立.所以可设 $p_{1}$,$p_{2}$ 中至少有一个不是 $\infty$.因此设 $1 \leq p < \infty$.于是,$F_{z}$ 保持(イ)中式~\eqref{eq:1.2.13} 不变,令 $G_{z} = g$,其余按(イ)的相同步骤进行即可.

(ハ)$p = \infty$ 的情形.此时 $p_{1} = p_{2} = p = \infty$.若 $q_{1} = q_{2} = 1$,则 $q = 1$,定理以 $M_{1} = M_{2}$ 成立.所以可设 $q_{1}$,$q_{2}$ 中至少有一个不是 $1$.因此设 $1 < q \leq \infty$.于是,令 $F_{z} = f$,$G_{z}$ 保持(イ)中式~\eqref{eq:1.2.13} 不变,其余按(イ)的相同步骤进行即可.

(ニ)当 $1 \leq p < \infty$ 时,$S(X)$ 在 $L^{p}(X, \mu)$ 中稠密,故如泛函分析中所熟知的那样,$T$ 可唯一扩张为从 $L^{p}(X, \mu)$ 到 $L^{q}(Y, \nu)$ 的有界线性算子,且
\begin{equation*}
\| Tf \|_{L^{q}(Y, \nu)} \leq M_{1}^{1-t}M_{2}^{t}\| f \|_{L^{p}(X, \mu)} \quad (f \in L^{p}(X, \mu)).
\end{equation*}
\end{proof}

由此导出:当 $1 \leq p \leq 2$ 时,Fourier 变换是从 $L^{p}(\mathbb{R})$ 到 $L^{p'}(\mathbb{R})$ 的有界线性算子,即 Hausdorff--Young 不等式.

\begin{theorem}[Hausdorff--Young 不等式]\label{thm:1.2.18}
设 $1 \leq p \leq 2$.则 Fourier 变换是从 $L^{p}(\mathbb{R})$ 到 $L^{p'}(\mathbb{R})$ 的有界线性算子,且
\begin{equation*}
\| \mathcal{F}f \|_{L^{p'}} \leq (2\pi)^{1/p'}\| f \|_{L^{p}} \quad (f \in L^{p}(\mathbb{R})).
\end{equation*}
其中,$1/p + 1/p' = 1$.
\end{theorem}

\begin{proof}
$p = 1$ 与 $p = 2$ 的情形已知,故设 $1 < p < 2$.Fourier 变换 $\mathcal{F}$ 的线性是显然的;由定理~\ref{thm:1.2.3},它是从可积简单函数的全体到 $L^{\infty}(\mathbb{R})$ 的线性算子,因而也是到可测函数全体的线性算子,且以 $p_{1} = 1$,$q_{1} = \infty$,$M_{1} = 1$ 满足 Riesz--Thorin 插值定理的条件 \eqref{eq:1.2.10}.又,可积简单函数也是平方可积的,故由定理~\ref{thm:1.2.15},它以 $p_{2} = q_{2} = 2$,$M_{2} = \sqrt{2\pi}$ 满足 Riesz--Thorin 插值定理的条件 \eqref{eq:1.2.11}.当 $1 < p < 2$ 时,由于 $1/p = (1 - 2(1 - 1/p)) + 2(1 - 1/p)/2$,故对 $1/q = (1 - 2(1 - 1/p))/\infty + 2(1 - 1/p)/2 = 1 - 1/p = 1/p'$,可用 Riesz--Thorin 插值定理,得
\begin{equation*}
\| \mathcal{F}f \|_{L^{p'}} \leq (2\pi)^{1/p'}\| f \|_{L^{p}}.
\end{equation*}
\end{proof}

\begin{note}\label{note:1.2.19}
当 $1 < p < 2$ 时,对 $f \in L^{p}(\mathbb{R})$,作分解 $f = f\chi_{\{|f| \leq 1\}} + f\chi_{\{|f| > 1\}}$.由于 $f\chi_{\{|f| \leq 1\}} \in L^{2}(\mathbb{R})$,$f\chi_{\{|f| > 1\}} \in L^{1}(\mathbb{R})$,由 $\mathcal{F}$ 作为从可积简单函数的全体到 $L^{p}(\mathbb{R})$ 的有界线性算子的扩张的唯一性,
\begin{equation*}
\mathcal{F}f = \mathcal{F}_{L^{2}}(f\chi_{\{|f| \leq 1\}}) + \mathcal{F}_{L^{1}}(f\chi_{\{|f| > 1\}}).
\end{equation*}
\end{note}

作为 Riesz--Thorin 插值定理的另一个应用,我们给出关于卷积的 Young 不等式.它是定理~\ref{thm:1.1.8} 的一般形式.

\begin{theorem}[Young 不等式]\label{thm:1.2.20}
设 $1 \leq p, q \leq \infty$,$0 \leq 1/r = 1/p + 1/q - 1 \leq 1$,则
\begin{equation*}
\| f * g \|_{L^{r}} \leq \| f \|_{L^{p}}\| g \|_{L^{q}} \quad (f \in L^{p}(\mathbb{R}),\ g \in L^{q}(\mathbb{R})).
\end{equation*}
\end{theorem}

\begin{proof}
设 $1 \leq q \leq \infty$,$g \in L^{q}(\mathbb{R})$.对这个 $g$ 令
\begin{equation*}
T(f) = f * g.
\end{equation*}
则,由定理~\ref{thm:1.1.8},$T$ 是从 $L^{1}(\mathbb{R})$ 到 $L^{q}(\mathbb{R})$ 的有界线性算子,且
\begin{equation*}
\| T(f) \|_{L^{q}} \leq \| g \|_{L^{q}}\| f \|_{L^{1}}.
\end{equation*}
又,应用 Hölder 不等式可知,$T$ 是从 $L^{q'}(\mathbb{R})$ 到 $L^{\infty}(\mathbb{R})$ 的有界线性算子,且
\begin{equation*}
\| T(f) \|_{L^{\infty}} \leq \| g \|_{L^{q}}\| f \|_{L^{q'}}.
\end{equation*}
当 $1 < p < q'$ 时,若 $f \in L^{p}(\mathbb{R})$,则可分解为 $f = f_{1} + f_{2}$,$f_{1} \in L^{1}(\mathbb{R})$,$f_{2} \in L^{q'}(\mathbb{R})$,故可定义 $T(f) = T(f_{1}) + T(f_{2})$,并且可看出这与 $f$ 的分解方式无关.进一步,当 $f$ 是可积简单函数时,$T(f) = f * g$.由于 $1/q > 1/p + 1/q - 1 = 1/p - 1/q' > 0$,令 $1/r = 1/p + 1/q - 1$,则 $q < r < \infty$;又令 $1/p = (1 - t) + t/q'$,则 $t = q(1 - 1/p)$,故 $(1 - t)/q = 1/q - (1 - 1/p) = 1/r$,因此由 Riesz--Thorin 插值定理,$T$ 是从 $L^{p}(\mathbb{R})$ 到 $L^{r}(\mathbb{R})$ 的有界线性算子,且
\begin{equation*}
\| f * g \|_{L^{r}} = \| T(f) \|_{L^{r}} \leq \| g \|_{L^{q}}\| f \|_{L^{p}}.
\end{equation*}
\end{proof}

% ============ 第1章 §1.3 Hilbert变换的L²有界性(原书 pp.29-33 = PDF pp.41-45) ============
\section{\texorpdfstring{Hilbert 变换的 $L^2$ 有界性}{Hilbert 变换的 L2 有界性}}\label{sec:1.3}
\index{Hilbert 变换}
\index{Hilbert 变换!L2 有界性}


\begin{definition}\label{def:1.3.1}
对几乎所有的 $x \in \mathbb{R}$, 
\begin{equation*}
Hf(x) = \mathrm{p.v.}\,\frac{1}{\pi}\int_{-\infty}^{\infty}\frac{f(y)}{x - y}\dd y := \lim_{\varepsilon \to 0}\frac{1}{\pi}\int_{|x - y| > \varepsilon}\frac{f(y)}{x - y}\dd y
\end{equation*}
存在时, 称 $Hf$ 为 $f$ 的 Hilbert 变换.p.v. 是 Cauchy 意义下主值的缩写.
\end{definition}

\begin{lemma}\label{lem:1.3.2}
若 $f \in C^{1}(\mathbb{R})$, $|f'(x)| \leq \frac{A_{1}}{1 + |x|}$ 且 $|f(x)| \leq \frac{A_{2}}{(1 + |x|)^{2}}$($x \in \mathbb{R}$), 则 p.v. $\frac{1}{\pi}\int_{-\infty}^{\infty}\frac{f(y)}{x - y}\dd y$ 对一切 $x \in \mathbb{R}$ 存在, 且存在满足下述不等式的正常数 $C_{1}$:
\begin{equation*}
|Hf(x)| \leq \frac{C_{1}(A_{1} + A_{2})}{1 + |x|}.
\end{equation*}
\end{lemma}

\begin{proof}
令
\begin{align*}
H_{\varepsilon}f(x) & := \frac{1}{\pi}\int_{|y - x| > \varepsilon}\frac{f(y)}{x - y}\dd y \\
& = \frac{1}{\pi}\int_{\varepsilon < |y - x| < 1}\frac{f(y)}{x - y}\dd y + \frac{1}{\pi}\int_{|y - x| \geq 1}\frac{f(y)}{x - y}\dd y \\
& =: I_{1,\varepsilon}(x) + I_{2}(x).
\end{align*}
由于 $\int_{\varepsilon < |y - x| < 1}\frac{1}{x - y}\dd y = 0$, 
\begin{equation*}
I_{1,\varepsilon}(x) = \frac{1}{\pi}\int_{\varepsilon < |y - x| < 1}\frac{f(y) - f(x)}{x - y}\dd y.
\end{equation*}
又, 由中值定理与对 $f'$ 的估计, 对 $0 < |x - y| < 1$, 
\begin{equation*}
\left|\frac{f(y) - f(x)}{x - y}\right| \leq \max_{|z - x| \leq 1}|f'(z)| \leq \max_{|z - x| \leq 1}\frac{A_{1}}{1 + |z|} \leq \frac{2A_{1}}{1 + |x|}.
\end{equation*}
因此, $\frac{f(y) - f(x)}{x - y}$ 在区间 $(x - 1, x)$,$(x, x + 1)$ 上可积, 故对一切 $x \in \mathbb{R}$, $\lim_{\varepsilon \to 0}\int_{\varepsilon < |y - x| < 1}\frac{f(y) - f(x)}{x - y}\dd y$ 存在, 且
\begin{equation*}
\left|\lim_{\varepsilon \to 0}\int_{\varepsilon < |y - x| < 1}\frac{f(y) - f(x)}{x - y}\dd y\right| \leq \frac{4A_{1}}{1 + |x|}.
\end{equation*}

关于 $I_{2}$, 分情形讨论.

(i) $|x| \leq 2$ 的情形.此时注意 $1 + |x| \leq 3$.由 Cauchy--Schwarz 不等式, 
\begin{align*}
\pi I_{2}(x) &\leq \int_{|z| \geq 1}\frac{A_{2}}{|z|(1 + |x - z|)^{2}}\dd z \\
\quad &\leq A_{2}\left(\int_{|z| \geq 1}\frac{\dd z}{|z|^{2}}\right)^{1/2}\left(\int_{|z| \geq 1}\frac{\dd z}{(1 + |x - z|)^{4}}\right)^{1/2} \\
\quad &\leq A_{2}\left(2\int_{1}^{\infty}\frac{\dd t}{t^{2}}\right)^{1/2}\left(2\int_{0}^{\infty}\frac{\dd t}{(1 + t)^{4}}\right)^{1/2} = \frac{2A_{2}}{\sqrt{3}} < \frac{2\sqrt{3}\,A_{2}}{1 + |x|}.
\end{align*}

(ii) $|x| \geq 2$ 的情形.此时 $3|x| \geq 2 + 2|x|$, 注意 $|x| > \frac{2}{3}(1 + |x|)$.首先, 当 $1 \leq |z| \leq \frac{|x|}{2}$ 时, $|x - z| \geq |x| - |z| \geq \frac{|x|}{2} \geq |z|$, 故
\begin{align*}
& \int_{1 \leq |z| \leq |x|/2}\frac{1}{|z|(1 + |x - z|)^{2}}\dd z \leq \frac{1}{1 + \frac{|x|}{2}}\int_{1 \leq |z| \leq |x|/2}\frac{1}{|z|(1 + |z|)}\dd z \\
& \qquad < \frac{4}{1 + |x|}\int_{1}^{\infty}\frac{\dd t}{t^{2}} = \frac{4}{1 + |x|}.
\end{align*}
此外, 
\begin{align*}
\int_{|z| > |x|/2}\frac{1}{|z|(1 + |x - z|)^{2}}\dd z & < \frac{2}{|x|}\int_{-\infty}^{\infty}\frac{1}{(1 + |x - z|)^{2}}\dd z \\
& = \frac{2}{|x|}\int_{-\infty}^{\infty}\frac{1}{(1 + |t|)^{2}}\dd t = \frac{4}{|x|} \leq \frac{6}{1 + |x|}.
\end{align*}

由 (i),(ii), 
\begin{equation*}
\pi I_{2}(x) \leq A_{2}\int_{|z| \geq 1}\frac{1}{|z|(1 + |x - z|)^{2}}\dd z \leq \frac{10A_{2}}{1 + |x|}.
\end{equation*}
与对 $\lim_{\varepsilon \to 0}I_{1,\varepsilon}$ 的估计合并, 得
\begin{equation*}
|Hf(x)| \leq \frac{(4A_{1} + 10A_{2})}{\pi(1 + |x|)}.
\end{equation*}
\end{proof}

利用上述引理, 可证明 $\mathcal{S}(\mathbb{R})$ 中函数的 Hilbert 变换的 $L^{2}$ 有界性.

\begin{theorem}\label{thm:1.3.3}
\begin{equation*}
\| Hf \|_{L^{2}} = \| f \|_{L^{2}} \quad (f \in \mathcal{S}(\mathbb{R})).
\end{equation*}
\end{theorem}

\begin{proof}
设 $f \in \mathcal{S}(\mathbb{R})$.由上述引理, $Hf \in L^{2}(\mathbb{R})$.现在, 对 $\varphi \in \mathcal{S}(\mathbb{R})$, 
\begin{equation*}
I := \int Hf(x)\overline{\varphi(x)}\dd x = \frac{1}{\pi}\int\left(\lim_{\varepsilon \to 0,\, A \to \infty}\int_{\varepsilon < |x - y| < A}\frac{f(y)}{x - y}\dd y\right)\overline{\varphi(x)}\dd x.
\end{equation*}
由上述引理的证明, 存在与 $\varepsilon$,$A$ 无关的常数 $C$, 使得
\begin{equation*}
\left|\int_{\varepsilon < |x - y| < A}\frac{f(y)}{x - y}\dd y\right| \leq C\frac{1}{1 + |x|}
\end{equation*}
所以
\begin{equation*}
\left|\left(\int_{\varepsilon < |x - y| < A}\frac{f(y)}{x - y}\dd y\right)\overline{\varphi(x)}\right| \leq C\frac{|\varphi(x)|}{1 + |x|} \in L^{1}(\mathbb{R}).
\end{equation*}
因此可用 Lebesgue 收敛定理, 
\begin{align*}
& \pi I \\
& = \lim_{\varepsilon \to 0,\, A \to \infty}\int\left(\int_{\varepsilon < |x - y| < A}\frac{f(y)}{x - y}\dd y\right)\overline{\varphi(x)}\dd x \\
& = \lim_{\varepsilon \to 0,\, A \to \infty}\int\left(\int_{\varepsilon < |z| < A}\frac{f(x - z)}{z}\dd z\right)\overline{\varphi(x)}\dd x \quad (\text{变量替换 } z = x - y) \\
& = \lim_{\varepsilon \to 0,\, A \to \infty}\int_{\varepsilon < |z| < A}\left(\int f(x - z)\overline{\varphi(x)}\dd x\right)\frac{\dd z}{z} \quad (\text{Fubini})
\end{align*}
\begin{align*}
& = \lim_{\varepsilon \to 0,\, A \to \infty}\int_{\varepsilon < |z| < A}\left(\frac{1}{2\pi}\int \mathcal{F}(f(x - z))(\xi)\overline{\mathcal{F}\varphi(\xi)}\dd \xi\right)\frac{\dd z}{z} \quad (\text{Parseval}) \\
& = \lim_{\varepsilon \to 0,\, A \to \infty}\int_{\varepsilon < |z| < A}\left(\frac{1}{2\pi}\int \mathcal{F}f(\xi)e^{-iz\xi}\overline{\mathcal{F}\varphi(\xi)}\dd \xi\right)\frac{\dd z}{z} \quad (\text{定理~\ref{thm:1.2.8}(v)}) \\
& = \lim_{\varepsilon \to 0,\, A \to \infty}\frac{1}{2\pi}\int\left(\int_{\varepsilon < |z| < A}\frac{e^{-iz\xi}}{z}\dd z\right)\mathcal{F}f(\xi)\overline{\mathcal{F}\varphi(\xi)}\dd \xi \quad (\text{Fubini}).
\end{align*}
这里
\begin{equation*}
I(\varepsilon, A; \xi) := \int_{\varepsilon < |z| < A}\frac{e^{-iz\xi}}{z}\dd z = -2i\int_{\varepsilon}^{A}\frac{\sin z\xi}{z}\dd z = -2i\int_{\varepsilon\xi}^{A\xi}\frac{\sin z}{z}\dd z
\end{equation*}
且, 如所周知
\begin{equation*}
\lim_{\varepsilon \to 0,\, A \to \infty} I(\varepsilon, A; \xi) = \left\{ \begin{array}{l l}
-2i\int_{0}^{\infty}\frac{\sin z}{z}\dd z = -\pi i & (\xi > 0),\\
2i\int_{0}^{\infty}\frac{\sin z}{z}\dd z = \pi i & (\xi < 0).
\end{array} \right.
\end{equation*}
又, 容易看出
\begin{equation*}
|I(\varepsilon, A; \xi)| \leq 2\int_{0}^{\pi}\frac{\sin t}{t}\dd t < \pi + \frac{4}{\pi}.
\end{equation*}
因此, 用 Lebesgue 收敛定理, 
\begin{equation*}
\pi I = -\frac{i}{2}\int\frac{\xi}{|\xi|}\mathcal{F}f(\xi)\overline{\mathcal{F}\varphi(\xi)}\dd \xi.
\end{equation*}
另一方面, 注意到 $Hf \in L^{2}(\mathbb{R})$, 由 Parseval 等式, 
\begin{equation*}
\int Hf(x)\overline{\varphi(x)}\dd x = \frac{1}{2\pi}\int\mathcal{F}(Hf)(\xi)\overline{\mathcal{F}\varphi(\xi)}\dd \xi
\end{equation*}
所以
\begin{equation*}
\int\mathcal{F}(Hf)(\xi)\overline{\mathcal{F}\varphi(\xi)}\dd \xi = \int -i\frac{\xi}{|\xi|}\mathcal{F}f(\xi)\overline{\mathcal{F}\varphi(\xi)}\dd \xi.
\end{equation*}
由 $\mathcal{F}(Hf)$,$\frac{\xi}{|\xi|}\mathcal{F}f \in L^{2}(\mathbb{R})$ 以及 $\mathcal{S}(\mathbb{R})$ 在 $L^{2}(\mathbb{R})$ 中稠密, 得
\begin{equation*}
\mathcal{F}(Hf)(\xi) = -i\frac{\xi}{|\xi|}\mathcal{F}f(\xi).
\end{equation*}
因此, 
\begin{equation*}
\| Hf \|_{L^{2}} = \frac{1}{\sqrt{2\pi}}\|\mathcal{F}(Hf)\|_{L^{2}} = \frac{1}{\sqrt{2\pi}}\left\|\frac{\xi}{|\xi|}\mathcal{F}f\right\|_{L^{2}} = \frac{1}{\sqrt{2\pi}}\|\mathcal{F}f\|_{L^{2}} = \| f \|_{L^{2}}
\end{equation*}
证毕.
\end{proof}

\begin{definition}\label{def:1.3.4}
由于 $\mathcal{S}(\mathbb{R})$ 在 $L^{2}(\mathbb{R})$ 中稠密, Hilbert 变换 $H$ 可唯一地扩张到整个 $L^{2}(\mathbb{R})$.扩张到整个 $L^{2}(\mathbb{R})$ 后的变换仍记作 $H$.
\end{definition}

于是, 

\begin{theorem}\label{thm:1.3.5}
\begin{equation*}
\| Hf \|_{L^{2}} = \| f \|_{L^{2}} \quad (f \in L^{2}(\mathbb{R})).
\end{equation*}
\end{theorem}

\begin{note}\label{note:1.3.6}
如后文所述, 当 $f \in L^{2}(\mathbb{R})$ 时, 对几乎所有的 $x \in \mathbb{R}$, $\lim_{\varepsilon \to 0}H_{\varepsilon}f(x)$ 存在, 且
\begin{equation*}
\lim_{\varepsilon \to 0}H_{\varepsilon}f(x) = Hf(x) \quad (\text{a.e. } x \in \mathbb{R}).
\end{equation*}
\end{note}

% ============ 第1章 §1.4 Hilbert变换及其锐极大函数估计(原书 pp.33-37 = PDF pp.45-49) ============
\section{Hilbert 变换及其锐极大函数估计}\label{sec:1.4}
\index{极大函数!锐极大函数}


本节验证:$f$ 的 Hilbert 变换 $Hf$ 的某种平均化可用 $f$ 的同类平均化来估计.为此定义如下极大函数.

\begin{definition}\label{def:1.4.1}
设 $1 \leq p < \infty$.对 $f \in L_{\mathrm{loc}}^{p}(\mathbb{R})$,
\begin{equation*}
M_{p}f(x) := \sup_{x \in Q}\left(\frac{1}{|Q|}\int_{Q}|f(y)|^{p}\dd y\right)^{1/p}
\end{equation*}
(其中,上确界取遍包含 $x$ 的区间 $Q$)称为 $p$ 次 Hardy--Littlewood 极大函数.又,
\begin{equation*}
\tilde{M}_{p}f(x) := \sup_{r > 0}\left(\frac{1}{2r}\int_{x-r}^{x+r}|f(y)|^{p}\dd y\right)^{1/p}
\end{equation*}
称为 $p$ 次 Hardy--Littlewood 中心极大函数.

当 $p = 1$ 时,常常省略 $1$,记为 $Mf$,$\tilde{M}f$ 等.
\end{definition}

\begin{note}\label{note:1.4.2}
由定义并使用 Hölder 不等式,可知
\begin{equation*}
M_{p_{1}}f(x) \leq M_{p_{2}}f(x) \quad (1 \leq p_{1} < p_{2}).
\end{equation*}
\end{note}

\begin{definition}\label{def:1.4.3}
对 $f \in L_{\mathrm{loc}}^{1}(\mathbb{R})$,
\begin{equation*}
f^{\#}(x) := \sup_{x \in Q}\frac{1}{|Q|}\int_{Q}|f(y) - f_{Q}|\dd y
\end{equation*}
(其中,上确界取遍包含 $x$ 的区间 $Q$,$f_{Q} = |Q|^{-1}\int_{Q}f(y)\dd y$)称为 $f$ 的 $\#$ 极大函数(锐极大函数).
\end{definition}

首先,$M_{p}f(x)$ 与 $\tilde{M}_{p}f(x)$ 之间有如下关系.

\begin{lemma}\label{lem:1.4.4}
\begin{equation*}
2^{-1/p}M_{p}f(x) \leq \tilde{M}_{p}f(x) \leq M_{p}f(x).
\end{equation*}
\end{lemma}

\begin{proof}
设 $x \in Q = (a - r, a + r)$.则
\begin{equation*}
Q = (a - r, a + r) \subset (x - 2r, x + 2r) =: Q'.
\end{equation*}
因此
\begin{align*}
\left(\frac{1}{|Q|}\int_{Q}|f(y)|^{p}\dd y\right)^{1/p} &\leq \left(\frac{|Q'|}{|Q|}\cdot\frac{1}{|Q'|}\int_{Q'}|f(y)|^{p}\dd y\right)^{1/p} \\
&= 2^{1/p}\left(\frac{1}{4r}\int_{x-2r}^{x+2r}|f(y)|^{p}\dd y\right)^{1/p} \leq 2^{1/p}\tilde{M}_{p}f(x).
\end{align*}
因此
\begin{equation*}
2^{-1/p}M_{p}f(x) \leq \tilde{M}_{p}f(x).
\end{equation*}
不等式 $\tilde{M}_{p}f(x) \leq M_{p}f(x)$ 是显然的.
\end{proof}

现在,本节的主题是如下的命题.

\begin{proposition}\label{prop:1.4.5}
对 $f \in \mathcal{S}(\mathbb{R})$,
\begin{equation*}
(Hf)^{\#}(x) \leq 2\left(\sqrt{2} + \frac{8}{\pi}\right)M_{2}f(x) \quad (x \in \mathbb{R}).
\end{equation*}
\end{proposition}

\begin{proof}
设 $f \in \mathcal{S}(\mathbb{R})$,任意固定 $x_{0} \in \mathbb{R}$.又设
\begin{equation*}
x_{0} \in Q = (a - r, a + r), \quad \tilde{Q} = (a - 2r, a + 2r),
\end{equation*}
并令
\begin{equation*}
f_{1}(x) = f(x)\chi_{\tilde{Q}}(x), \quad f_{2}(x) = f(x) - f_{1}(x) = f(x)\chi_{\tilde{Q}^{c}}.
\end{equation*}
则对 $x \in Q$,
\begin{equation*}
Hf_{2}(x) = \frac{1}{\pi}\int_{|y - a| > 2r}\frac{1}{x - y}f_{2}(y)\dd y,
\end{equation*}
且
\begin{align*}
& \frac{1}{|Q|}\int_{Q}|Hf(x) - Hf_{2}(a)|\dd x \\
& \quad \leq \frac{1}{|Q|}\int_{Q}|Hf_{1}(x)|\dd x + \frac{1}{|Q|}\int_{Q}|Hf_{2}(x) - Hf_{2}(a)|\dd x =: I_{1} + I_{2}.
\end{align*}
于是
\begin{align*}
I_{1} & \leq \left(\frac{1}{|Q|}\int_{Q}|Hf_{1}(x)|^{2}\dd x\right)^{1/2} \quad \text{(Cauchy--Schwarz)} \\
& \leq \left(\frac{1}{|Q|}\int_{-\infty}^{\infty}|Hf_{1}(x)|^{2}\dd x\right)^{1/2} \\
& \leq \left(\frac{1}{|Q|}\int_{-\infty}^{\infty}|f_{1}(x)|^{2}\dd x\right)^{1/2} \quad \text{(Hilbert 变换的 $L^{2}$ 有界性(定理~\ref{thm:1.3.5}))} \\
& = \left(\frac{|\tilde{Q}|}{|Q|}\cdot\frac{1}{|\tilde{Q}|}\int_{\tilde{Q}}|f(x)|^{2}\dd x\right)^{1/2} \\
& \leq \sqrt{2}\,M_{2}f(x_{0}).
\end{align*}

再者,对 $x \in Q$,$y \in \tilde{Q}^{c}$,有 $|x - a| < r < |y - a|/2$,故 $|x - y| \geq |y - a| - |x - a| > |y - a|/2 > r$.因此
\begin{equation*}
\left|\frac{1}{x - y} - \frac{1}{a - y}\right| = \frac{|x - a|}{|x - y||a - y|} \leq 2\frac{|x - a|}{|y - a|^{2}}
\end{equation*}
注意到这一点,则
\begin{align*}
I_{2} &= \frac{1}{|Q|}\int_{Q}\left|\frac{1}{\pi}\int_{|y - a| > 2r}\left(\frac{1}{x - y} - \frac{1}{a - y}\right)f_{2}(y)\dd y\right|\dd x \\
&\leq \frac{1}{\pi|Q|}\int_{Q}\left(\int_{|y - a| > 2r}\left|\frac{1}{x - y} - \frac{1}{a - y}\right||f(y)|\dd y\right)\dd x \\
&\leq \frac{2}{\pi|Q|}\int_{Q}\left(\int_{|y - a| > 2r}\frac{|x - a|}{|y - a|^{2}}|f(y)|\dd y\right)\dd x \\
&= \frac{2}{\pi|Q|}\int_{Q}\left(\sum_{j=1}^{\infty}\int_{2^{j}r < |y - a| \leq 2^{j+1}r}\frac{|x - a|}{|y - a|^{2}}|f(y)|\dd y\right)\dd x \\
&\leq \frac{2}{\pi|Q|}\int_{Q}\left(\sum_{j=1}^{\infty}\int_{2^{j}r < |y - a| \leq 2^{j+1}r}\frac{1}{2^{2j}r}|f(y)|\dd y\right)\dd x \\
&\leq \frac{2^{3}}{\pi}\cdot\frac{1}{|Q|}\int_{Q}\left(\sum_{j=1}^{\infty}\frac{1}{2^{j}}\cdot\frac{1}{2^{j+2}r}\int_{|y - a| \leq 2^{j+1}r}|f(y)|\dd y\right)\dd x \\
&\leq \frac{8}{\pi}\cdot\frac{1}{|Q|}\int_{Q}\left(\sum_{j=1}^{\infty}\frac{1}{2^{j}}M_{1}f(x_{0})\right)\dd x \quad (\because x_{0} \in Q \subset (a - 2^{j+1}r, a + 2^{j+1}r)) \\
&\leq \frac{8}{\pi}M_{2}f(x_{0}) \quad (\because M_{p_{1}}f(x) \leq M_{p_{2}}f(x) \ (1 \leq p_{1} < p_{2})).
\end{align*}

因此,令 $Hf_{2}(a) = c$,合并 $I_{1}$ 与 $I_{2}$ 的估计,得
\begin{equation*}
\frac{1}{|Q|}\int_{Q}|Hf(x) - c|\dd x \leq \left(\sqrt{2} + \frac{8}{\pi}\right)M_{2}f(x_{0}).
\end{equation*}
由此,
\begin{align*}
\frac{1}{|Q|}\int_{Q}|Hf(x) - (Hf)_{Q}|\dd x&\leq \frac{1}{|Q|}\int_{Q}|Hf(x) - c|\dd x + \frac{1}{|Q|}\int_{Q}|(Hf)_{Q} - c|\dd x  \\
&= \frac{1}{|Q|}\int_{Q}|Hf(x) - c|\dd x + \left|\frac{1}{|Q|}\int_{Q}(Hf(y) - c)\dd y\right|
\end{align*}
\begin{equation*}
\leq 2\cdot\frac{1}{|Q|}\int_{Q}|Hf(x) - c|\dd x \leq 2\left(\sqrt{2} + \frac{8}{\pi}\right)M_{2}f(x_{0}).
\end{equation*}
对满足 $x_{0} \in Q$ 的区间 $Q$ 取上确界,得
\begin{equation*}
(Hf)^{\#}(x_{0}) \leq 2\left(\sqrt{2} + \frac{8}{\pi}\right)M_{2}f(x_{0}).
\end{equation*}
由于 $x_{0} \in \mathbb{R}$ 是任意的,所求估计得证.
\end{proof}

\begin{note}\label{note:1.4.6}
如上述证明的最后一段所示,成立下面的结论.这一事实将在 1.13 节再次提及.
\begin{equation*}
\sup_{x \in Q}\inf_{c \in \mathbb{C}}\frac{1}{|Q|}\int_{Q}|f(y) - c|\dd y \leq f^{\#}(x) \leq 2\sup_{x \in Q}\inf_{c \in \mathbb{C}}\frac{1}{|Q|}\int_{Q}|f(y) - c|\dd y.
\end{equation*}
\end{note}

% ============ 第1章 §1.5 覆盖定理与Hardy-Littlewood极大函数(原书 pp.37-44 = PDF pp.49-56) ============
\section{覆盖定理与 Hardy--Littlewood 极大函数}\label{sec:1.5}
\index{Hardy--Littlewood 极大函数}
\index{Hardy--Littlewood 中心极大函数}
\index{覆盖定理!Besicovitch 的}
\index{弱 (1,1) 型!Hardy--Littlewood 极大函数的}
\index{Chebyshev 不等式}


为了研究 Hardy--Littlewood 极大函数的基本性质,先来研究覆盖定理.以下,用 $\ell(Q)$ 表示区间 $Q$ 的长度.又,$Q(x, r)$ 表示中心为 $x$,长度为 $r$ 的区间 $(x - r/2, x + r/2)$.

\begin{theorem}[Besicovitch 覆盖定理]\label{thm:1.5.1}
设 $A$ 是 $\mathbb{R}$ 内的有界集,对每个 $x \in A$ 对应一个以 $x$ 为中心的区间 $Q(x)$,且 $\sup_{x \in A}\ell(Q(x)) < \infty$.则可以从区间族 $\{Q(x)\}_{x \in A}$ 中选出至多可数个区间列 $\{Q_{k}\}$,使其满足以下 3 个条件:

\begin{enumerate}[(i)]
\item $A \subset \bigcup_{k}Q_{k}$,
\item $\sum_{k}\chi_{Q_{k}}(x) \leq 4$,即包含 $x$($\in A$)的 $Q_{k}$ 至多 4 个,
\item 区间列 $\{Q_{k}\}$ 可分解为至多 9 组互不相交的子区间列 $\{Q_{k_{1}}\}$,$\{Q_{k_{2}}\}$,$\ldots$,$\{Q_{k_{\ell}}\}$($\ell \leq 9$),即各 $\{Q_{k_{j}}\}$ 互不相交,且 $\{Q_{k_{1}}\} \cup \{Q_{k_{2}}\} \cup \cdots \cup \{Q_{k_{\ell}}\} = \{Q_{k}\}$.
\end{enumerate}
\end{theorem}

\begin{proof}
设 $\alpha_{0} := \sup_{x \in A}\ell(Q(x)) < \infty$.首先选取 $x_{1} \in A$ 使 $\ell(Q(x_{1})) > \frac{1}{2}\sup_{x \in A}\ell(Q(x))$.其次,若 $A \setminus Q(x_{1}) \neq \emptyset$,选取 $x_{2} \in A \setminus Q(x_{1})$ 使 $\ell(Q(x_{2})) > \frac{1}{2}\sup_{x \in A \setminus Q(x_{1})}\ell(Q(x))$.以下归纳地进行:若 $A \setminus \bigcup_{j=1}^{k}Q(x_{j}) \neq \emptyset$,选取 $x_{k+1} \in A \setminus \bigcup_{j=1}^{k}Q(x_{j})$ 使 $\ell(Q(x_{k+1})) > \frac{1}{2}\sup_{x \in A \setminus \bigcup_{j=1}^{k}Q(x_{j})}\ell(Q(x))$.若 $A \setminus \bigcup_{j=1}^{k}Q(x_{j}) = \emptyset$,则在该阶段终止此操作.对如此得到的区间列 $\{Q(x_{k})\}$,令 $Q_{k} = Q(x_{k})$.下面证明这个区间列 $\{Q_{k}\}$ 满足 (i),(ii),(iii).

首先,成立以下 3 条.

(a) 若 $j < k$,则 $x_{k} \notin Q_{j}$,

(b) 若 $j < k$,则 $\ell(Q_{j}) > \frac{1}{2}\ell(Q_{k})$,

(c) $Q\bigl(x_{j}, \ell(Q_{j})/3\bigr) \cap Q\bigl(x_{k}, \ell(Q_{k})/3\bigr) = \emptyset$($j \neq k$).

(a),(b) 由 $Q_{k}$ 的构造方式是显然的.为验证 (c),可设 $j < k$.于是由 (a),$|x_{k} - x_{j}| \geq \ell(Q_{j})/2$;使用 (b),得 $\frac{1}{6}\ell(Q_{j}) + \frac{1}{6}\ell(Q_{k}) < \frac{1}{6}\ell(Q_{j}) + \frac{1}{3}\ell(Q_{j}) = \frac{1}{2}\ell(Q_{j}) \leq |x_{k} - x_{j}|$,即区间 $Q\bigl(x_{j}, \ell(Q_{j})/3\bigr)$ 与 $Q\bigl(x_{k}, \ell(Q_{k})/3\bigr)$ 的中心间的距离大于各自长度之和的一半,故 $Q\bigl(x_{j}, \ell(Q_{j})/3\bigr)$ 与 $Q\bigl(x_{k}, \ell(Q_{k})/3\bigr)$ 没有公共部分.

现在回到 (i) 的证明.若构造区间列 $\{Q(x_{k})\}$ 的操作有限次结束,则 (i) 已证,故设它无限继续下去.此时,$\lim_{j \to \infty}\ell(Q_{j}) = 0$.事实上,若设 $\limsup_{j \to \infty}\ell(Q_{j}) = a > 0$,则存在子列 $\{Q_{j_{k}}\}$ 使得 $\ell(Q_{j_{k}}) > a/2$.令 $Q'_{j_{k}} = Q\bigl(x_{j_{k}}, \frac{1}{3}\ell(Q_{j_{k}})\bigr)$,则
\begin{equation}\tag{1.5.1}
\sum_{k=1}^{\infty}\ell(Q'_{j_{k}}) > \sum_{k=1}^{\infty}\frac{a}{6} = +\infty.
\label{eq:1.5.1}\end{equation}
另一方面,$Q'_{j_{k}} \subset Q_{j_{k}} \subset Q\bigl(x_{1}, \mathrm{diam}(A) + \alpha_{0}\bigr)$,且由 (c),$Q'_{j_{k}}$ 两两不相交,所以
\begin{equation}\tag{1.5.2}
\sum_{k=1}^{\infty}\ell(Q'_{j_{k}}) < \mathrm{diam}(A) + \alpha_{0}.
\label{eq:1.5.2}\end{equation}
式~\eqref{eq:1.5.1} 与式~\eqref{eq:1.5.2} 矛盾.因此 $\lim_{j \to \infty}\ell(Q_{j}) = 0$.

于是,假设存在 $x_{0} \in A \setminus \bigcup_{j=1}^{\infty}Q_{j}$.则由 $\lim_{j \to \infty}\ell(Q_{j}) = 0$,存在 $Q_{j_{0}}$ 使得 $\ell(Q_{j_{0}}) < \frac{1}{2}\ell(Q_{x_{0}})$.由 $x_{0} \in A \setminus \bigcup_{j=1}^{\infty}Q_{j}$,有 $x_{0} \in A \setminus \bigcup_{j=1}^{j_{0}-1}Q_{j}$,故由 $Q_{j_{0}}$ 的定义,
\begin{equation*}
\ell(Q_{j_{0}}) > \frac{1}{2}\sup_{x \in A \setminus \bigcup_{j=1}^{j_{0}-1}Q(x_{j})}\ell(Q(x)) \geq \frac{1}{2}\ell(Q(x_{0})).
\end{equation*}
这与 $\ell(Q_{j_{0}}) < \frac{1}{2}\ell(Q_{x_{0}})$ 矛盾.因此
\begin{equation*}
A \setminus \bigcup_{j=1}^{\infty}Q_{j} = \emptyset \quad \text{即} \quad A \subset \bigcup_{j=1}^{\infty}Q_{j}.
\end{equation*}

(ii) 的证明.由 (c),$x_{i} \neq x_{j}$($i \neq j$).对给定的 $x$,设满足 $x \in Q_{j}$ 且 $x \leq x_{j}$ 的编号 $j$ 至少有 2 个,并设其中最小的编号为 $j_{1}$,第二小的编号为 $j_{2}$.今设 $m > j_{2}$ 且 $x < x_{m}$.由 (a),$x_{j_{2}}, x_{m} \notin Q_{j_{1}}$,$x_{m} \notin Q_{j_{2}}$.又由 (b),
\begin{equation}\tag{1.5.3}
\frac{1}{2}\ell(Q_{m}) < \ell(Q_{j_{1}}).
\label{eq:1.5.3}\end{equation}
进一步,由 $x \in Q_{j_{2}}$,$x \leq x_{j_{1}} < x_{j_{2}}$,$x_{j_{2}} \notin Q_{j_{1}}$,得 $\ell(Q_{j_{2}})/2 > \ell(Q_{j_{1}})/2$.因此,再利用 $x_{j_{2}} \notin Q_{j_{1}}$,$x_{m} \notin Q_{j_{2}}$ 以及 $x \leq x_{j_{1}} < x_{j_{2}} < x_{m}$,得
\begin{equation}\tag{1.5.4}
|x_{m} - x| \geq |x_{m} - x_{j_{2}}| + |x_{j_{2}} - x_{j_{1}}| > \frac{\ell(Q_{j_{2}})}{2} + \frac{\ell(Q_{j_{1}})}{2} > \ell(Q_{j_{1}}).
\label{eq:1.5.4}\end{equation}
由 \eqref{eq:1.5.3},\eqref{eq:1.5.4} 可知 $x \notin Q_{m}$.因此,中心在 $x$ 右侧且包含 $x$ 的 $Q_{j}$ 至多 2 个.同样地,中心在 $x$ 左侧且包含 $x$ 的 $Q_{j}$ 也至多 2 个.所以,包含 $x$ 的 $Q_{j}$ 至多 4 个.

(iii) 的证明.设 $j \in \mathbb{N}$.若 $m < j$ 且 $Q_{m} \cap Q_{j} \neq \emptyset$,则 $Q_{m}$ 要么包含 $Q_{j}$ 的端点,要么不包含端点.若不包含端点,则 $Q_{m} \subset Q_{j}$,且由 (b),$\ell(Q_{m}) > \frac{1}{2}\ell(Q_{j})$,故 $x_{j} \in Q_{m}$.但这与 (a) 矛盾.因此 $Q_{m}$ 包含 $Q_{j}$ 的两个端点之一.于是,满足 $m < j$ 且 $Q_{m} \cap Q_{j} \neq \emptyset$ 的编号 $m$ 由 (ii) 至多为 $4 \times 2 = 8$ 个.于是,把 (i) 的证明中构造的区间列 $\{Q(x_{k})\}$ 按如下方式分成 9 个类 $I_{1}, I_{2}, \ldots, I_{9}$.首先,把 $Q_{1}$ 放入类 $I_{1}$,$Q_{2}$ 放入类 $I_{2}$,$\ldots$,$Q_{9}$ 放入类 $I_{9}$.与 $Q_{10}$ 相交的 $Q_{j}$($j < 10$)至多 8 个,故至少有一个与 $Q_{10}$ 不相交的 $Q_{\ell}$($\ell \leq 9$).把 $Q_{10}$ 放入该编号的类 $I_{\ell}$.其次,与 $Q_{11}$ 相交的 $Q_{j}$($j < 11$)至多 8 个,故存在编号 $\ell$($\ell \leq 9$)使得
\begin{equation*}
Q_{11} \cap \bigcup_{Q_{j} \in I_{\ell},\, j < 11}Q_{j} = \emptyset.
\end{equation*}
把 $Q_{11}$ 放入该编号的类 $I_{\ell}$.继续这一操作,则每个 $Q_{i}$ 都落入 $I_{1}, I_{2}, \ldots, I_{9}$ 中的某一类,且各同类 $I_{\ell}$ 中的 $Q_{j}$ 两两不相交.即,这种分类正是 (iii) 所求的分类.
\end{proof}

利用这个覆盖定理的 (i),(ii),可证明称为 Hardy--Littlewood 极大函数弱 $(1,1)$ 型的定理.

\begin{theorem}\label{thm:1.5.2}
存在常数 $C$,使得对任意 $f \in L^{1}(\mathbb{R})$,
\begin{equation*}
|\{x \in \mathbb{R};\ (M_{1}f)(x) > \lambda\}| \leq C\frac{\| f \|_{L^{1}}}{\lambda} \quad (\lambda > 0).
\end{equation*}
\end{theorem}

\begin{proof}
首先,对中心极大函数 $\tilde{M}_{1}$ 证明
\begin{equation*}
|\{x \in \mathbb{R};\ (\tilde{M}_{1}f)(x) > \lambda\}| \leq 4\frac{\| f \|_{L^{1}}}{\lambda} \quad (\lambda > 0).
\end{equation*}
则由引理~\ref{lem:1.4.4},$M_{1}f(x) \leq 2\tilde{M}_{1}f(x)$,定理即得证.

现在,设 $\lambda > 0$,任意固定 $B > 0$.定义集合 $A_{B}$ 为
\begin{equation*}
A_{B} := (-B, B) \cap \{x \in \mathbb{R};\ (\tilde{M}_{1}f)(x) > \lambda\}.
\end{equation*}
对每个 $x \in A_{B}$,由 $(\tilde{M}_{1}f)(x)$ 的定义,存在 $r_{x} > 0$ 使得
\begin{equation*}
\frac{1}{|Q(x, r_{x})|}\int_{Q(x, r_{x})}|f(y)|\dd y > \lambda,
\end{equation*}
故对应一个区间 $Q(x, r_{x})$.于是,由 Besicovitch 覆盖定理,可从区间族 $\{Q(x, r_{x})\}_{x \in A_{B}}$ 中选出满足
(i) $A_{B} \subset \bigcup_{k}Q_{k}$,
(ii) $\sum_{k}\chi_{Q_{k}}(x) \leq 4$
的区间列 $\{Q_{k}\}$.于是,由 (i),
\begin{align*}
|A_{B}| &\leq \sum_{k}|Q_{k}| \leq \sum_{k}\frac{1}{\lambda}\int_{Q_{k}}|f(y)|\dd y = \frac{1}{\lambda}\int_{\mathbb{R}}\sum_{k}\chi_{Q_{k}}(y)|f(y)|\dd y \\
&\leq \frac{4}{\lambda}\int_{\mathbb{R}}|f(y)|\dd y.
\end{align*}
因此,令 $B \to \infty$,得
\begin{equation*}
|\{x \in \mathbb{R};\ (\tilde{M}_{1}f)(x) > \lambda\}| \leq 4\frac{\| f \|_{L^{1}}}{\lambda},
\end{equation*}
定理证毕.
\end{proof}

\begin{note}\label{note:1.5.3}
令 $f(x) = \chi_{[-1,1]}$,则
\begin{equation*}
M_{1}f(x) = \left\{ \begin{array}{l l}
1 & (|x| \leq 1), \\
\frac{2}{|x| + 1} & (|x| > 1)
\end{array} \right.
\end{equation*}
所以对 $0 < \lambda < 1$,
\begin{equation*}
|\{x \in \mathbb{R};\ (M_{1}f)(x) > \lambda\}| = \frac{4}{\lambda} - 2.
\end{equation*}
\end{note}

\begin{note}\label{note:1.5.4}
定理~\ref{thm:1.5.2} 也可以看作对 Chebyshev 不等式
\begin{equation*}
|\{x \in \mathbb{R};\ |f(x)| > \lambda\}| \leq \frac{\| f \|_{L^{1}}}{\lambda} \quad (\lambda > 0)
\end{equation*}
的精细化.
\end{note}

由定理~\ref{thm:1.5.2} 可导出 Hardy--Littlewood 极大函数的 $L^{p}$ 有界性.

\begin{theorem}\label{thm:1.5.5}
设 $1 < p \leq \infty$.则存在只与 $p$ 有关的正常数 $C_{p}$,使得
\begin{equation*}
\| M_{1}f \|_{L^{p}} \leq C_{p}\| f \|_{L^{p}} \quad (f \in L^{p}(\mathbb{R})).
\end{equation*}
\end{theorem}

\begin{proof}
(i) $p = \infty$ 的情形.若 $f \in L^{\infty}(\mathbb{R})$,则
\begin{equation*}
\frac{1}{|Q|}\int_{Q}|f(y)|\dd y \leq \| f \|_{L^{\infty}}\frac{1}{|Q|}\int_{Q}\dd y = \| f \|_{L^{\infty}}
\end{equation*}
所以对任意 $x \in \mathbb{R}$,
\begin{equation*}
M_{1}f(x) \leq \| f \|_{L^{\infty}}.
\end{equation*}
因此
\begin{equation*}
\| M_{1}f \|_{L^{\infty}} \leq \| f \|_{L^{\infty}} \quad (f \in L^{\infty}(\mathbb{R})),
\end{equation*}
故 $p = \infty$ 时定理成立.

(ii) $1 < p < \infty$ 的情形.设 $f \in L^{p}(\mathbb{R})$.对 $\lambda > 0$,令 $E_{\lambda} = \{x \in \mathbb{R};\ M_{1}f(x) > \lambda\}$,$\sigma(\lambda) = |E_{\lambda}|$,则
\begin{equation*}
\int_{\mathbb{R}}|M_{1}f(x)|^{p}\dd x = -\int_{0}^{\infty}\lambda^{p}\dd \sigma(\lambda) = p\int_{0}^{\infty}\lambda^{p-1}\sigma(\lambda)\dd \lambda.
\end{equation*}
进一步,令
\begin{equation*}
f_{1}(x) = \left\{ \begin{array}{l l}
f(x) & (|f(x)| > \lambda/2) \\
0 & (|f(x)| \leq \lambda/2),
\end{array} \right. \quad f_{2}(x) = f(x) - f_{1}(x)
\end{equation*}
则 $|f(x)| \leq |f_{1}(x)| + \frac{\lambda}{2}$,故 $M_{1}f(x) \leq M_{1}f_{1}(x) + \frac{\lambda}{2}$.因此
\begin{equation*}
E_{\lambda} \subset \left\{x \in \mathbb{R};\ M_{1}f_{1}(x) > \frac{\lambda}{2}\right\}.
\end{equation*}
于是,$\sigma(\lambda) \leq |\{x \in \mathbb{R};\ M_{1}f_{1}(x) > \lambda/2\}|$,故使用定理~\ref{thm:1.5.2},
\begin{equation*}
\sigma(\lambda) \leq \frac{16}{\lambda}\int_{\mathbb{R}}|f_{1}(y)|\dd y = \frac{16}{\lambda}\int_{\{y;\ |f(y)| > \lambda/2\}}|f(y)|\dd y.
\end{equation*}
因此
\begin{align*}
\int_{\mathbb{R}}|M_{1}f(x)|^{p}\dd x & = p\int_{0}^{\infty}\lambda^{p-1}\sigma(\lambda)\dd \lambda \\
& \leq 16p\int_{0}^{\infty}\lambda^{p-1}\left(\frac{1}{\lambda}\int_{\{y;\ |f(y)| > \lambda/2\}}|f(y)|\dd y\right)\dd \lambda \\
& = 16p\int_{0}^{\infty}\lambda^{p-2}\left(\int_{\{y;\ |f(y)| > \lambda/2\}}|f(y)|\dd y\right)\dd \lambda \\
& = 16p\int_{\mathbb{R}}|f(y)|\left(\int_{0}^{2|f(y)|}\lambda^{p-2}\dd \lambda\right)\dd y \quad \text{(Fubini)} \\
& = \frac{16p}{p - 1}\int_{\mathbb{R}}2^{p-1}|f(y)|^{p}\dd y.
\end{align*}
由此,
\begin{equation*}
\| M_{1}f \|_{L^{p}} \leq 2\left(\frac{8p}{p - 1}\right)^{1/p}\| f \|_{L^{p}}.
\end{equation*}
\end{proof}

由这个定理可得出 $p$ 次 Hardy--Littlewood 极大函数的 $L^{r}$ 有界性.

\begin{corollary}\label{cor:1.5.6}
设 $1 \leq p < r < \infty$.则存在只与 $p$,$r$ 有关的正常数 $C_{p,r}$,使得
\begin{equation*}
\| M_{p}f \|_{L^{r}} \leq C_{p,r}\| f \|_{L^{r}} \quad (f \in L^{r}(\mathbb{R})).
\end{equation*}
\end{corollary}

\begin{proof}
由 $M_{p}$ 的定义与定理~\ref{thm:1.5.5},
\begin{align*}
\left(\int_{\mathbb{R}}M_{p}f(x)^{r}\dd x\right)^{1/r} &= \left[\left(\int_{\mathbb{R}}M_{1}(|f|^{p})(x)^{r/p}\dd x\right)^{p/r}\right]^{1/p} \\
&\leq \left[C_{r/p}\left(\int_{\mathbb{R}}\left(|f(x)|^{p}\right)^{r/p}\dd x\right)^{p/r}\right]^{1/p} \\
&= (C_{r/p})^{1/p}\| f \|_{L^{r}}.
\end{align*}
\end{proof}

% ============ 第1章 §1.6 锐极大函数与Hardy-Littlewood极大函数的关系(原书 pp.44-51 = PDF pp.56-63) ============
\section{锐极大函数与 Hardy--Littlewood 极大函数的关系}\label{sec:1.6}
\index{good $\lambda$-不等式}
\index{Whitney 型覆盖定理}
\index{二进区间}


\begin{definition}\label{def:1.6.1}
形如 $[\ell 2^{k}, (\ell + 1)2^{k})$(其中 $k, \ell \in \mathbb{Z}$)的区间称为二进区间.令 $\mathbb{D}_{k} = \{[\ell 2^{k}, (\ell + 1)2^{k})\}_{\ell \in \mathbb{Z}}$,并用 $\mathbb{D}$ 表示全体二进区间.
\end{definition}

对二进区间,成立如下基本性质.

\begin{lemma}\label{lem:1.6.2}
(i) 对每个 $x \in \mathbb{R}$ 与每个 $k \in \mathbb{Z}$,包含 $x$ 的二进区间 $Q \in \mathbb{D}_{k}$ 恰有一个.

(ii) 若 $Q, Q' \in \mathbb{D}$,则 $Q \subset Q'$,$Q' \subset Q$ 或 $Q \cap Q' = \emptyset$ 三者恰有一个发生.

(iii) 对每个 $Q \in \mathbb{D}_{k}$,包含 $Q$ 的 $Q' \in \mathbb{D}_{k+1}$ 恰有一个.
\end{lemma}

由这个几乎自明的事实可得出下面的事实.

\begin{lemma}\label{lem:1.6.3}
设 $Q_{j} \in \mathbb{D}$($j = 1, 2, \ldots$),且在 $\{Q_{j}\}_{j \in \mathbb{N}}$ 中不存在无限扩张的子区间列.则可以从 $\{Q_{j}\}_{j \in \mathbb{N}}$ 中选出子区间列 $\{Q_{j_{\ell}}\}_{\ell}$(有限个或无限个),使其满足
\begin{enumerate}[(i)]
\item $\bigcup_{j \in \mathbb{N}}Q_{j} = \bigcup_{\ell}Q_{j_{\ell}}$,
\item $Q_{j_{\ell}} \cap Q_{j_{k}} = \emptyset$($\ell \neq k$),
\item $Q_{j_{\ell}}$ 在 $\{Q_{j}\}_{j \in \mathbb{N}}$ 中是极大的,即 $\{Q_{j}\}_{j \in \mathbb{N}}$ 中不存在真包含 $Q_{j_{\ell}}$ 的区间.
\end{enumerate}
\end{lemma}

\begin{proof}
由引理~\ref{lem:1.6.2}(ii),(iii),对每个 $Q_{j}$,包含 $Q_{j}$ 的区间 $Q_{k}$ 的全体按包含关系排成一列.由假设,这一列是有限列,故最宽的区间唯一确定,它在包含关系下是极大的.把在包含关系下极大的这些 $Q_{k}$ 的全体记为 $\{Q_{j_{\ell}}\}$,则
\begin{equation*}
\bigcup_{j \in \mathbb{N}}Q_{j} = \bigcup_{\ell}Q_{j_{\ell}},
\end{equation*}
且由极大性,
\begin{equation*}
Q_{j_{\ell}} \cap Q_{j_{k}} = \emptyset \quad (\ell \neq k),
\end{equation*}
引理的结论全部得到满足.
\end{proof}

利用二进区间的这些基本性质,可以证明如下的 Whitney 型覆盖定理.

\begin{theorem}[Whitney 型覆盖定理]\label{thm:1.6.4}
设 $\Omega$ 是 $\mathbb{R}$ 的非空真开子集.则存在闭区间列 $\{Q_{j}\}$(有限个或无限个)满足:
\begin{enumerate}[(i)]
\item $\Omega = \bigcup_{j}Q_{j}$ 且 $Q_{j}^{\circ} \cap Q_{k}^{\circ} = \emptyset$($j \neq k$);
\item $\ell(Q_{j}) \leq \mathrm{dist}(Q_{j}, \Omega^{c}) \leq 4\ell(Q_{j})$;
\item 若 $Q_{j}$ 与 $Q_{k}$ 相接,则
\begin{equation*}
\frac{1}{4} \leq \frac{\ell(Q_{j})}{\ell(Q_{k})} \leq 4.
\end{equation*}
\end{enumerate}
\end{theorem}

\begin{proof}
对 $k \in \mathbb{Z}$,令
\begin{equation*}
\Omega_{k} = \{x \in \Omega;\ 2\cdot 2^{-k} < \mathrm{dist}(x, \Omega^{c}) \leq 4\cdot 2^{-k}\},
\end{equation*}
则
\begin{equation*}
\Omega = \bigcup_{k \in \mathbb{Z}}\Omega_{k}.
\end{equation*}
设 $E_{k}$ 是满足 $Q \cap \Omega_{k} \neq \emptyset$ 且 $Q \in \mathbb{D}_{-k}$ 的二进区间的全体,并令 $E = \bigcup_{k \in \mathbb{Z}}E_{k}$.

若 $Q \in E$,则对某个 $k \in \mathbb{Z}$ 有 $Q \in E_{k}$,故对这个 $k$,存在 $x \in Q \cap \Omega_{k}$,且 $\ell(Q) = 2^{-k}$.由 $x \in \Omega_{k}$,$2\cdot 2^{-k} < \mathrm{dist}(x, \Omega^{c}) \leq 4\cdot 2^{-k}$,所以
\begin{equation*}
\ell(Q) < \mathrm{dist}(x, \Omega^{c}) - \ell(Q) \leq \mathrm{dist}(Q, \Omega^{c}) \leq \mathrm{dist}(x, \Omega^{c}) \leq 4\cdot 2^{-k} = 4\ell(Q),
\end{equation*}
故 $Q$ 满足条件 (ii).又,$\mathrm{dist}(Q, \Omega^{c}) \geq \ell(Q) > 0$,故 $Q \subset \Omega$.因此 $\bigcup_{Q \in E}Q \subset \Omega$.

反之,若 $x \in \Omega$,则存在 $k \in \mathbb{Z}$ 使得 $x \in \Omega_{k}$.进一步,由于 $\bigcup_{Q \in \mathbb{D}_{-k}}Q = \mathbb{R}$,存在 $Q \in \mathbb{D}_{-k}$ 使得 $x \in Q$.由 $E_{k}$ 的定义,这个 $Q$ 属于 $E_{k}$,从而属于 $E$.因此,$\Omega \subset \bigcup_{Q \in E}Q$,得
\begin{equation*}
\Omega = \bigcup_{Q \in E}Q.
\end{equation*}

又,若 $Q, Q' \in E$,则分别有
\begin{equation*}
\ell(Q) < \mathrm{dist}(Q, \Omega^{c}) \leq 4\ell(Q), \quad \ell(Q') < \mathrm{dist}(Q', \Omega^{c}) \leq 4\ell(Q'),
\end{equation*}
进一步设 $Q \subset Q'$,注意到 $\mathrm{dist}(Q, \Omega^{c}) \geq \mathrm{dist}(Q', \Omega^{c})$,则
\begin{equation*}
\ell(Q') \leq 4\ell(Q).
\end{equation*}
因此,$E$ 中包含给定 $Q \in E$ 的二进区间至多 2 个.也就是说,$E$ 中的二进区间满足引理~\ref{lem:1.6.3} 的假设.于是,设由引理~\ref{lem:1.6.3} 从 $E$ 中得到的极大二进区间列为 $\{Q_{j}\}$,则
\begin{equation*}
\Omega = \bigcup_{j}Q_{j} \quad \text{且} \quad Q_{j} \cap Q_{k} = \emptyset \quad (j \neq k).
\end{equation*}

现在来确认 (iii).若 $Q_{j}$ 与 $Q_{k}$ 相接,则
\begin{equation*}
\begin{aligned}
\ell(Q_{j}) &\leq \mathrm{dist}(Q_{j}, \Omega^{c}) \leq \mathrm{diam}(Q_{k}) + \mathrm{dist}(Q_{k}, \Omega^{c}) \\
&\leq \ell(Q_{k}) + 4\ell(Q_{k}) \leq 5\ell(Q_{k}).
\end{aligned}
\end{equation*}
因此 $\ell(Q_{j})/\ell(Q_{k}) \leq 5$.但是,$Q_{j}$,$Q_{k}$ 都是二进区间,故 $\ell(Q_{j})/\ell(Q_{k})$ 具有 $2^{m}$($m \in \mathbb{Z}$)的形状,所以 $\ell(Q_{j})/\ell(Q_{k}) \leq 4$.交换 $Q_{j}$ 与 $Q_{k}$ 的角色,得 $\ell(Q_{k})/\ell(Q_{j}) \leq 4$,从而 $1/4 \leq \ell(Q_{j})/\ell(Q_{k}) \leq 4$.若把 $Q_{j}$ 换成加上其右端点的闭区间(仍记作 $Q_{j}$),则 $\{Q_{j}\}$ 满足全部所需性质.
\end{proof}

利用这个覆盖定理,可证明给出 Hardy--Littlewood 极大函数与锐极大函数关系的如下 ``good $\lambda$-不等式''.

\begin{lemma}\label{lem:1.6.5}
若 $f \in L_{\mathrm{loc}}^{1}(\mathbb{R})$,则对任意 $\lambda > 0$,$\gamma > 0$,
\begin{equation*}
\left|\left\{x \in \mathbb{R};\ M_{1}f(x) > 2\lambda,\ f^{\#}(x) < \gamma\lambda\right\}\right| \leq 96\gamma\left|\left\{x \in \mathbb{R};\ M_{1}f(x) > \lambda\right\}\right|
\end{equation*}
成立.
\end{lemma}

\begin{proof}
以下,将 $M_{1}f$ 简记为 $Mf$.设 $\lambda > 0$,$\gamma > 0$,令 $E_{\lambda} = \{x \in \mathbb{R};\ Mf(x) > \lambda\}$.若 $E_{\lambda} = \mathbb{R}$,则引理的结论正确,故以下设 $E_{\lambda} \neq \mathbb{R}$.$E_{\lambda}$ 是开集.事实上,若 $x_{0} \in E_{\lambda}$,由上确界的定义,存在包含 $x_{0}$ 的区间 $Q_{0}$ 使得
\begin{equation*}
\frac{1}{|Q_{0}|}\int_{Q_{0}}|f(x)|\dd x > \lambda.
\end{equation*}
由积分的性质,可取一个包含 $Q_{0}$ 的开区间 $Q_{1}$ 仍满足上述性质.若 $x \in Q_{1}$,则由 $\frac{1}{|Q_{1}|}\int_{Q_{1}}|f(x)|\dd x > \lambda$ 得 $Mf(x) > \lambda$,故 $Q_{1} \subset E_{\lambda}$.因此 $E_{\lambda}$ 是开集.又因 $E_{\lambda} \neq \mathbb{R}$,$E_{\lambda}$ 满足定理~\ref{thm:1.6.4} 的假设,故可取如下区间列 $\{Q_{j}\}$:
\begin{equation*}
(\mathrm{i})\ E_{\lambda} = \bigcup_{j}\mathring{Q}_{j} \quad \text{且} \quad Q_{j}^{\circ} \cap Q_{k}^{\circ} = \emptyset \quad (j \neq k).
\end{equation*}
\begin{equation*}
\text{(ii)}\ \ell(Q_{j}) \leq \mathrm{dist}(Q_{j}, E_{\lambda}^{c}) \leq 4\ell(Q_{j}).
\end{equation*}
固定其中任意一个 $Q_{j}$,为简单起见记作 $Q$.

设 $\tilde{Q} = 12Q$.由于 $\mathrm{dist}(Q, E_{\lambda}^{c}) \leq 4\ell(Q) = \frac{1}{3}\ell(\tilde{Q})$,故 $\tilde{Q} \cap E_{\lambda}^{c} \neq \emptyset$,由 $E_{\lambda}$ 的定义,
\begin{equation*}
\frac{1}{|\tilde{Q}|}\int_{\tilde{Q}}|f(x)|\dd x \leq \lambda.
\end{equation*}

其次,若 $x \in Q$ 且 $Mf(x) > 2\lambda$,则存在包含 $x$ 的区间 $R$ 使得 $|R|^{-1}\int_{R}|f(y)|\dd y > 2\lambda$.于是,设 $R'$ 是以 $x$ 为中心,长度为 $2\ell(R)$ 的区间,则
\begin{equation*}
\frac{1}{|R'|}\int_{R'}|f(y)|\dd y \geq \frac{|R|}{|R'|}\cdot\frac{1}{|R|}\int_{R}|f(y)|\dd y > \lambda.
\end{equation*}
若 $R' \cap E_{\lambda}^{c} \neq \emptyset$,则 $|R'|^{-1}\int_{R'}|f(y)|\dd y \leq \lambda$,与上面的不等式矛盾,故 $R' \cap E_{\lambda}^{c} = \emptyset$.又由 $\mathrm{dist}(Q, E_{\lambda}^{c}) \leq 4\ell(Q)$,
\begin{equation*}
\mathrm{dist}(x, E_{\lambda}^{c}) \leq \ell(Q) + \mathrm{dist}(Q, E_{\lambda}^{c}) \leq 5\ell(Q).
\end{equation*}
因此 $R' \subset Q(x, 11\ell(Q))$,从而 $R \subset Q(x, 11\ell(Q))$.用 $x_{Q}$ 表示 $Q$ 的中心,则 $|x - x_{Q}| \leq \ell(Q)/2$,故 $R \subset Q(x, 11\ell(Q)) \subset Q(x_{Q}, 12\ell(Q)) = \tilde{Q}$.因此 $Q \subset \tilde{Q}$ 且 $R \subset \tilde{Q}$.于是,
\begin{equation*}
\frac{1}{|R|}\int_{R}|f(y)\chi_{\tilde{Q}}(y)|\dd y = \frac{1}{|R|}\int_{R}|f(y)|\dd y > 2\lambda
\end{equation*}
所以
\begin{align*}
\frac{1}{|R|}\int_{R}|f(y) - f_{\tilde{Q}}|\chi_{\tilde{Q}}(y)\dd y&= \frac{1}{|R|}\int_{R}|f(y) - f_{\tilde{Q}}|\dd y \geq \frac{1}{|R|}\int_{R}|f(y)|\dd y - |f_{\tilde{Q}}|  \\
\qquad &\geq \frac{1}{|R|}\int_{R}|f(y)|\dd y - \frac{1}{|\tilde{Q}|}\int_{\tilde{Q}}|f(y)|\dd y > 2\lambda - \lambda = \lambda.
\end{align*}
这表明
\begin{equation*}
M\bigl((f - f_{\tilde{Q}})\chi_{\tilde{Q}}\bigr)(x) > \lambda.
\end{equation*}

对 $M\bigl((f - f_{\tilde{Q}})\chi_{\tilde{Q}}\bigr)$ 使用定理~\ref{thm:1.5.2},
\begin{align*}
|\{x \in \mathbb{R};\ M\bigl((f - f_{\tilde{Q}})\chi_{\tilde{Q}}\bigr)(x) > \lambda\}|&\leq \frac{8}{\lambda}\int_{\tilde{Q}}|f(y) - f_{\tilde{Q}}|\dd y \leq \frac{8|\tilde{Q}|}{\lambda}\cdot\frac{1}{|\tilde{Q}|}\int_{\tilde{Q}}|f(y) - f_{\tilde{Q}}|\dd y  \\
&\leq \frac{8\cdot 12|Q|}{\lambda}\inf_{x \in \tilde{Q}}f^{\#}(x) \leq \frac{96|Q|}{\lambda}\inf_{x \in Q}f^{\#}(x).
\end{align*}
因此,若 $\inf_{x \in Q}f^{\#}(x) < \gamma\lambda$,则
\begin{align*}
& |\{x \in Q;\ Mf(x) > 2\lambda,\ f^{\#}(x) < \gamma\lambda\}| \\
& \qquad \leq |\{x \in Q;\ M\bigl((f - f_{\tilde{Q}})\chi_{\tilde{Q}}\bigr)(x) > \lambda,\ f^{\#}(x) < \gamma\lambda\}| \\
& \qquad \leq \frac{96|Q|}{\lambda}\inf_{x \in Q}f^{\#}(x) = 96\gamma|Q|.
\end{align*}
又,若 $\inf_{x \in Q}f^{\#}(x) \geq \gamma\lambda$,则上式左端内的集合为空,故上式成立.无论哪种情形,
\begin{equation}\tag{1.6.1}
\left|\{x \in Q;\ Mf(x) > 2\lambda,\ f^{\#}(x) < \gamma\lambda\}\right| \leq 96\gamma|Q|
\label{eq:1.6.1}\end{equation}
得证.

把对 $Q_{j}$ 的估计式~\eqref{eq:1.6.1} 对 $Q_{j}$ 求和,由于 $Q_{j}$ 的内部两两不相交,
\begin{align*}
& |\{x \in \mathbb{R};\ Mf(x) > 2\lambda,\ f^{\#}(x) < \gamma\lambda\}| \\
& \qquad = |\{x \in E_{\lambda};\ Mf(x) > 2\lambda,\ f^{\#}(x) < \gamma\lambda\}| \\
& \qquad \leq 96\gamma|E_{\lambda}| = 96\gamma|\{x \in \mathbb{R};\ Mf(x) > \lambda\}|.
\end{align*}
\end{proof}

下面给出联系 ``good $\lambda$-不等式'' 与 $L^{p}$ 不等式的引理.

\begin{lemma}\label{lem:1.6.6}
设 $0 \leq w(x) \in L_{\mathrm{loc}}^{1}(\mathbb{R})$,$0 < p < \infty$.设 $f$,$g$ 是可测函数,$|f(x)|^{p}w(x) \in L^{1}(\mathbb{R})$,且存在满足以下两个条件的 $\varepsilon > 0$,$\gamma > 0$:
\begin{enumerate}[(i)]
\item $(1 + \varepsilon)^{p}\gamma < 1$,
\item $w(\{x \in \mathbb{R};\ |f(x)| > (1 + \varepsilon)\lambda,\ |g(x)| \leq \lambda\}) \leq \gamma w(\{x \in \mathbb{R};\ |f(x)| > \lambda\})$($\lambda > 0$).
\end{enumerate}
其中,$w(E) = \int_{E}w(x)\dd x$.

则存在常数 $C(p, \varepsilon, \gamma)$ 使得
\begin{equation*}
\| f \|_{L^{p}(w)} \leq C(p, \varepsilon, \gamma)\| g \|_{L^{p}(w)}.
\end{equation*}
其中,$\| f \|_{L^{p}(w)} = \left(\int_{\mathbb{R}}|f(x)|^{p}w(x)\dd x\right)^{1/p}$.
\end{lemma}

\begin{proof}
由假设 (ii),
\begin{align*}
w(\{x \in \mathbb{R};\ |f(x)| > (1 + \varepsilon)\lambda\})&\leq w(\{x \in \mathbb{R};\ |g(x)| > \lambda\}) + \gamma w(\{x \in \mathbb{R};\ |f(x)| > \lambda\})
 \end{align*}
所以
\begin{align*}
\| f \|_{L^{p}(w)}^{p} &= \int_{\mathbb{R}}|f(x)|^{p}w(x)\dd x \\
\qquad &= \int_{0}^{\infty}p\lambda^{p-1}w(\{x \in \mathbb{R};\ |f(x)| > \lambda\})\dd \lambda \\
\qquad &= (1 + \varepsilon)^{p}\int_{0}^{\infty}p\lambda^{p-1}w(\{x \in \mathbb{R};\ |f(x)| > (1 + \varepsilon)\lambda\})\dd \lambda \\
\qquad &\leq (1 + \varepsilon)^{p}\int_{0}^{\infty}p\lambda^{p-1}w(\{x \in \mathbb{R};\ |g(x)| > \lambda\})\dd \lambda \\
&\qquad\quad + \gamma(1 + \varepsilon)^{p}\int_{0}^{\infty}p\lambda^{p-1}w(\{x \in \mathbb{R};\ |f(x)| > \lambda\})\dd \lambda \\
\qquad &= (1 + \varepsilon)^{p}\| g \|_{L^{p}(w)}^{p} + \gamma(1 + \varepsilon)^{p}\| f \|_{L^{p}(w)}^{p}.
\end{align*}
因此,由 $|f(x)|^{p}w(x) \in L^{1}(\mathbb{R})$ 得 $\| f \|_{L^{p}(w)} < \infty$,从两边减去 $\gamma(1 + \varepsilon)^{p}\| f \|_{L^{p}(w)}^{p}$,得
\begin{equation*}
\left(1 - \gamma(1 + \varepsilon)^{p}\right)\| f \|_{L^{p}(w)}^{p} \leq (1 + \varepsilon)^{p}\| g \|_{L^{p}(w)}^{p},
\end{equation*}
由条件 $(1 + \varepsilon)^{p}\gamma < 1$,
\begin{equation*}
\| f \|_{L^{p}(w)} \leq \frac{(1 + \varepsilon)}{(1 - \gamma(1 + \varepsilon)^{p})^{1/p}}\| g \|_{L^{p}(w)}.
\end{equation*}
\end{proof}

\begin{note}\label{note:1.6.7}
条件 $f \in L^{p}(w)$ 也可以用 $\inf(1, |f|) \in L^{p}(w)$ 代替.进一步,更一般地,也可以用存在 $\alpha > 0$ 与 $C > 0$ 使得
\begin{equation*}
w(\{x \in \mathbb{R};\ |f(x)| > \lambda\}) \leq \frac{C}{\lambda^{\alpha}} \quad (0 < \lambda < 1)
\end{equation*}
这一条件来代替(宫地,薮田,On good $\lambda$-inequality, Bull.\ Fac.\ Sci.\ Ibaraki Univ.\ 16 (1984), 1--11).
\end{note}

利用关于 Hardy--Littlewood 极大函数与锐极大函数的 ``good $\lambda$-不等式'' 以及上面的引理,可证明给出 Hardy--Littlewood 极大函数与锐极大函数 $L^{p}$ 关系的如下定理.

\begin{theorem}\label{thm:1.6.8}
设 $1 \leq p < \infty$.若 $f \in L_{\mathrm{loc}}^{1}(\mathbb{R})$ 且 $Mf \in L^{p}(\mathbb{R})$,则存在与 $f$ 无关的常数 $C_{p} > 0$ 使得
\begin{equation*}
\| Mf \|_{L^{p}} \leq C_{p}\| f^{\#} \|_{L^{p}}.
\end{equation*}
\end{theorem}

\begin{proof}
由引理~\ref{lem:1.6.5},
\begin{equation*}
|\{x \in \mathbb{R};\ Mf(x) > 2\lambda,\ \frac{f^{\#}(x)}{\gamma} < \lambda\}| \leq 96\gamma|\{x \in \mathbb{R};\ Mf(x) > \lambda\}| \quad (\lambda > 0).
\end{equation*}
于是固定一个 $\gamma$ 使得 $2^{p} \times 96\gamma < 1$.则,由引理~\ref{lem:1.6.6},
\begin{align*}
\| Mf \|_{L^{p}} & \leq \frac{2}{(1 - 2^{p} \times 96\gamma)^{1/p}}\left\| \frac{f^{\#}(x)}{\gamma} \right\|_{L^{p}} \\
& \leq \frac{2}{\gamma(1 - 2^{p} \times 96\gamma)^{1/p}}\| f^{\#} \|_{L^{p}}.
\end{align*}
\end{proof}

\begin{note}\label{note:1.6.9}
条件 $Mf \in L^{p}(\mathbb{R})$ 也可以用 $\inf(1, |f|) \in L^{p}(\mathbb{R})$ 代替.进一步,更一般地,也可以用
\begin{equation*}
|\{x \in \mathbb{R};\ |f(x)| > \lambda\}| < +\infty \quad (0 < \lambda < 1)
\end{equation*}
这一条件来代替(参见注意~\ref{note:1.6.7} 的文献).
\end{note}

% ============ 第1章 §1.7 Hilbert变换的L^p有界性(原书 pp.51-54 = PDF pp.63-66) ============
\section{\texorpdfstring{Hilbert 变换的 $L^p$ 有界性($1<p<\infty$)}{Hilbert 变换的 Lp 有界性(1<p<∞)}}\label{sec:1.7}
\index{Hilbert 变换!Lp 有界性}


关于 Hilbert 变换的 $L^{p}$ 有界性,$L^{2}$ 有界性已经证明,下面先处理 $2 < p < \infty$ 的情形.

\begin{proposition}\label{prop:1.7.1}
设 $2 \leq p < \infty$.存在常数 $C_{p}$ 使得
\begin{equation*}
\| Hf \|_{L^{p}} \leq C_{p}\| f \|_{L^{p}} \quad (f \in \mathcal{S}(\mathbb{R})).
\end{equation*}
\end{proposition}

\begin{proof}
由于 $f \in \mathcal{S}(\mathbb{R})$,由引理~\ref{lem:1.3.2},$|Hf(x)| \leq C(1 + |x|)^{-1}$.因此,容易看出
\begin{equation*}
M(Hf)(x) \leq \frac{C}{1 + |x|}\log(1 + |x|) \in L^{r}(\mathbb{R}) \quad (1 < r < \infty).
\end{equation*}
由此,可用定理~\ref{thm:1.6.8},
\begin{equation*}
\| M(Hf) \|_{L^{p}} \leq C_{p}\| (Hf)^{\#} \|_{L^{p}} \quad (1 < p < \infty).
\end{equation*}
结合命题~\ref{prop:1.4.5},
\begin{equation*}
\| M(Hf) \|_{L^{p}} \leq C_{p}\| M_{2}f \|_{L^{p}}.
\end{equation*}
进一步使用推论~\ref{cor:1.5.6},当 $p > 2$ 时,
\begin{equation*}
\| Hf \|_{L^{p}} \leq C_{p}'\| f \|_{L^{p}}.
\end{equation*}
\end{proof}

由 $2 \leq p < \infty$ 时的 $L^{p}$ 有界性与对偶性,可证明 $1 < p < \infty$ 时的 $L^{p}$ 有界性.

\begin{theorem}\label{thm:1.7.2}
设 $1 < p < \infty$.存在常数 $C_{p}$ 使得
\begin{equation*}
\| Hf \|_{L^{p}} \leq C_{p}\| f \|_{L^{p}} \quad (f \in \mathcal{S}(\mathbb{R})).
\end{equation*}
\end{theorem}

\begin{proof}
$2 \leq p < \infty$ 的情形已经证明,故设 $1 < p < 2$.对 $f, g \in \mathcal{S}(\mathbb{R})$,
\begin{align*}
(Hf, g) & = \int_{-\infty}^{\infty} Hf(x)\overline{g(x)}\dd x \\
& = \frac{1}{2\pi}\int_{-\infty}^{\infty}\widehat{Hf}(\xi)\overline{\hat{g}(\xi)}\dd \xi \qquad \text{(Parseval)} \\
& = \frac{1}{2\pi}\int_{-\infty}^{\infty} -i\frac{\xi}{|\xi|}\hat{f}(\xi)\overline{\hat{g}(\xi)}\dd \xi \\
& = \frac{1}{2\pi}\int_{-\infty}^{\infty} \hat{f}(\xi)\overline{i\frac{\xi}{|\xi|}\hat{g}(\xi)}\dd \xi \\
& = \int_{-\infty}^{\infty} f(x)\overline{(-1)Hg(x)}\dd x = (f, -Hg) \qquad \text{(Parseval)}
\end{align*}
设 $p'$ 是满足 $1/p + 1/p' = 1$ 的数,则 $2 < p' < \infty$,故使用命题~\ref{prop:1.7.1},
\begin{equation*}
|(Hf, g)| = |(f, -Hg)| \leq \| f \|_{L^{p}}\| Hg \|_{L^{p'}} \leq C_{p'}\| f \|_{L^{p}}\| g \|_{L^{p'}}.
\end{equation*}
因此,由引理~\ref{lem:1.3.2} 有 $Hf \in L^{p}(\mathbb{R})$,且由于 $\mathcal{S}(\mathbb{R})$ 在 $L^{p'}(\mathbb{R})$ 中稠密,
\begin{equation*}
\left|\int_{-\infty}^{\infty} Hf(x)\overline{g(x)}\dd x\right| \leq C_{p'}\| f \|_{L^{p}}\| g \|_{L^{p'}} \quad (g \in L^{p'}(\mathbb{R})).
\end{equation*}
因此,由 $L^{p}$ 与 $L^{p'}$ 的对偶性,
\begin{equation*}
\| Hf \|_{L^{p}} \leq C_{p'}\| f \|_{L^{p}}.
\end{equation*}

(不使用对偶性,注意到由引理~\ref{lem:1.3.2} 有 $|Hf(x)| \leq C(1 + |x|)^{-1}$,在最后一个不等式中令 $g(x) = |Hf(x)|^{p-1}\frac{Hf(x)}{|Hf(x)|}\chi_{\operatorname{supp}(Hf)}(x)$,也可以直接验证 $\| Hf \|_{L^{p}} \leq C_{p'}\| f \|_{L^{p}}$.)
\end{proof}

\begin{note}\label{note:1.7.3}
由上面的定理,由于 $\mathcal{S}(\mathbb{R})$ 在 $L^{p}(\mathbb{R})$ 中稠密,Hilbert 变换可唯一地扩张到整个 $L^{p}(\mathbb{R})$;把扩张后的变换仍记作 $Hf$,则 Hilbert 变换是从 $L^{p}(\mathbb{R})$ 到 $L^{p}(\mathbb{R})$ 的有界线性算子,且
\begin{equation*}
\| Hf \|_{L^{p}} \leq C_{p}\| f \|_{L^{p}} \quad (f \in L^{p}(\mathbb{R})).
\end{equation*}
\end{note}

\begin{corollary}\label{cor:1.7.4}
设 $1 < p < \infty$.则存在常数 $C_{p}$ 使得
\begin{equation*}
(Hf)^{\#}(x) \leq C_{p}M_{p}f(x) \quad (x \in \mathbb{R},\ f \in \mathcal{S}(\mathbb{R})).
\end{equation*}
\end{corollary}

\begin{proof}
对命题~\ref{prop:1.4.5} 的证明稍加修改即可证明.事实上,在命题~\ref{prop:1.4.5} 的证明中,
\begin{equation*}
I_{2} \leq CM_{1}f(x).
\end{equation*}
关于 $I_{1}$,
\begin{align*}
& I_{1} \leq \left(\frac{1}{|Q|}\int_{Q}|Hf_{1}(x)|^{p}\dd x\right)^{1/p} \quad \text{(Hölder)} \\
& \quad \leq \left(\frac{1}{|Q|}\int_{-\infty}^{\infty}|Hf_{1}(x)|^{p}\dd x\right)^{1/p} \\
& \quad \leq C_{p}\left(\frac{1}{|Q|}\int_{-\infty}^{\infty}|f_{1}(x)|^{p}\dd x\right)^{1/p} \quad \text{(Hilbert 变换的 $L^{p}$ 有界性)} \\
& \quad = C_{p}\left(\frac{|Q'|}{|Q|}\cdot\frac{1}{|Q'|}\int_{Q'}|f(x)|^{p}\dd x\right)^{1/p} \\
& \quad \leq 2^{1/p}C_{p}M_{p}f(x_{0}).
\end{align*}
其余与命题~\ref{prop:1.4.5} 的证明相同.
\end{proof}

% ============ 第1章 §1.8 Calderón-Zygmund分解与Hilbert变换的弱(1,1)性(原书 pp.54-70) ============
\section{\texorpdfstring{Calder\'on--Zygmund 分解与 Hilbert 变换的弱 $(1,1)$ 性}{Calderon--Zygmund 分解与 Hilbert 变换的弱 (1,1) 性}}\label{sec:1.8}
\index{Calderón--Zygmund 分解}
\index{弱 (1,1) 型!Hilbert 变换的}
\index{Kolmogorov 不等式}
\index{插值定理!Marcinkiewicz 的}


当 $1 < p < \infty$ 时,我们已经看到 Hilbert 变换是 $L^{p}$ 有界线性算子.那么 $p = 1$ 的情形如何呢:$L^{1}$ 有界性不成立,但称为弱 $(1,1)$ 性的结果由 A. P. Calderón 与 A. Zygmund 给出.其关键就是下面的 Calderón--Zygmund 分解.

\begin{lemma}[Calderón--Zygmund 分解]\label{lem:1.8.1}
设 $0 \leq f(x) \in L^{1}(\mathbb{R})$.则对任意 $\lambda > 0$,存在 $\mathbb{R}$ 的如下分解:

(i) $\mathbb{R} = F \cup \Omega$,$F \cap \Omega = \emptyset$;

(ii) $f(x) \leq \lambda$(a.e. $x \in F$);

(iii) $\Omega = \bigcup_{j} Q_{j}$,$Q_{j} \in \mathbb{D}$,$Q_{j} \cap Q_{k} = \emptyset$($j \neq k$),且
\begin{align*}
\text{(a)} & \lambda < \dfrac{1}{|Q_{j}|}\displaystyle\int_{Q_{j}} f(x)\dd x \leq 2\lambda, \\[4pt]
\text{(b)} & |\Omega| = \displaystyle\sum_{j}|Q_{j}| \leq \dfrac{1}{\lambda}\displaystyle\int_{\mathbb{R}} f(x)\dd x.
\end{align*}
\end{lemma}

\begin{proof}
设 $0 \leq f(x) \in L^{1}(\mathbb{R})$.固定一个满足
\begin{equation*}
\int_{\mathbb{R}} f(x)\dd x \leq 2^{k_{0}}\lambda
\end{equation*}
的整数 $k_{0}$.则对 $Q \in \mathbb{D}_{k_{0}}$,
\begin{equation}\tag{1.8.1}
\frac{1}{|Q|}\int_{Q} f(x)\dd x \leq \lambda.
\label{eq:1.8.1}\end{equation}
$\mathbb{D}_{k_{0}-1}$ 由 $\mathbb{D}_{k_{0}}$ 的二进区间二等分而得到;将 $\mathbb{D}_{k_{0}-1}$ 的二进区间 $Q$ 中满足
\begin{equation}\tag{1.8.2}
\frac{1}{|Q|}\int_{Q} f(x)\dd x > \lambda
\label{eq:1.8.2}\end{equation}
的区间称为属于第 1 群.再将不属于第 1 群的 $\mathbb{D}_{k_{0}-1}$ 的二进区间二等分所得 $\mathbb{D}_{k_{0}-2}$ 的二进区间中,满足 \eqref{eq:1.8.2} 的区间称为属于第 1 群.以下归纳地重复这一操作.如此所得属于第 1 群的二进区间的全体,因二进区间的全体是可数集,故当然可数,记为 $\{Q_{j}\}_{j}$.由作法,$Q_{j} \cap Q_{k} = \emptyset$($j \neq k$),且包含 $Q_{j}$ 的最小二进区间 $Q$ 不属于第 1 群,故满足 \eqref{eq:1.8.1}.于是
\begin{equation*}
\lambda < \frac{1}{|Q_{j}|}\int_{Q_{j}} f(x)\dd x \leq \frac{|Q|}{|Q_{j}|}\cdot\frac{1}{|Q|}\int_{Q} f(x)\dd x \leq 2\lambda.
\end{equation*}
又因 $Q_{j}$ 满足 \eqref{eq:1.8.2},
\begin{equation*}
\sum_{j}|Q_{j}| \leq \sum_{j}\frac{1}{\lambda}\int_{Q_{j}} f(x)\dd x \leq \frac{1}{\lambda}\int_{-\infty}^{\infty} f(x)\dd x.
\end{equation*}
故令 $\Omega = \bigcup_{j} Q_{j}$,$F = \Omega^{c}$,则 (i),(iii) 成立.

关于 (ii).设 $x \in F = \Omega^{c}$,则对每个 $k \leq k_{0}$,存在二进区间 $R_{k} \in \mathbb{D}_{k}$ 使 $x \in R_{k}$.于是对几乎所有的 $x \in F$,
\begin{equation*}
f(x) = \lim_{k \to -\infty}\frac{1}{|R_{k}|}\int_{R_{k}} f(y)\dd y \leq \lambda.
\end{equation*}
\end{proof}

利用这个分解,可以证明 Hilbert 变换的弱 $(1,1)$ 性.

\begin{theorem}\label{thm:1.8.2}
对 $f \in L^{1}(\mathbb{R}) \cap L^{2}(\mathbb{R})$,存在常数 $C > 0$ 使得
\begin{equation*}
\left|\{x \in \mathbb{R};\ |Hf(x)| > \lambda\}\right| \leq \frac{C}{\lambda}\| f \|_{L^{1}} \quad (\lambda > 0).
\end{equation*}
\end{theorem}

\begin{proof}
设 $f \in L^{1}(\mathbb{R}) \cap L^{2}(\mathbb{R})$,$\lambda > 0$.对 $|f(x)|$ 应用 Calderón--Zygmund 分解的引理~\ref{lem:1.8.1},定义
\begin{align*}
g(x) &= \begin{cases} f(x) & (x \in F) \\ \dfrac{1}{|Q_{j}|}\displaystyle\int_{Q_{j}} f(x)\dd x & (x \in Q_{j},\ j = 1, 2, \dots) \end{cases} \\[10pt]
f_{j}(x) &= \left(f(x) - \dfrac{1}{|Q_{j}|}\displaystyle\int_{Q_{j}} f(y)\dd y\right)\chi_{Q_{j}} \\[8pt]
b(x) &= \displaystyle\sum_{j} f_{j}(x).
\end{align*}
则 $f(x) = g(x) + b(x)$,且
\begin{equation*}
\begin{aligned}
|g(x)| &\leq 2\lambda \quad (\text{a.e. } x \in \mathbb{R}), \\
\int_{\mathbb{R}}|g(x)|\dd x &\leq \int_{\mathbb{R}}|f(x)|\dd x.
\end{aligned}
\end{equation*}
由此 $g \in L^{1}(\mathbb{R}) \cap L^{\infty}(\mathbb{R})$,故亦有 $g \in L^{1}(\mathbb{R}) \cap L^{2}(\mathbb{R})$.另外
\begin{equation*}
\int_{Q_{j}} f_{j}(x)\dd x = 0.
\end{equation*}
其次,设 $Q_{j}$ 的中心为 $x_{j}$,
\begin{equation*}
Q_{j}^{*} = 2Q_{j} = Q(x_{j}, 2\ell(Q_{j})), \quad \Omega^{*} = \bigcup_{j} Q_{j}^{*},
\end{equation*}
则
\begin{equation}\tag{1.8.3}
|\Omega^{*}| \leq \sum_{j}|Q_{j}^{*}| = \sum_{j} 2|Q_{j}| = 2|\Omega| \leq \frac{2}{\lambda}\| f \|_{L^{1}}.
\label{eq:1.8.3}\end{equation}
现在,由 $f(x) = g(x) + b(x)$ 得 $Hf(x) = Hg(x) + Hb(x)$,若 $|Hg(x)| \leq \lambda/2$ 且 $|Hb(x)| \leq \lambda/2$,则 $|Hf(x)| \leq |Hg(x)| + |Hb(x)| \leq \lambda$,故
\begin{align*}
\{x \in \mathbb{R};\ |Hf(x)| > \lambda,\ x \in (\Omega^{*})^{c}\}&\subset \{x \in (\Omega^{*})^{c};\ |Hg(x)| > \lambda/2\} \cup \{x \in (\Omega^{*})^{c};\ |Hb(x)| > \lambda/2\}.
 \end{align*}
令
\begin{equation*}
E_{1} = \{x \in (\Omega^{*})^{c};\ |Hg(x)| > \lambda/2\}, \quad E_{2} = \{x \in (\Omega^{*})^{c};\ |Hb(x)| > \lambda/2\}.
\end{equation*}
首先,利用 $H$ 的 $L^{2}$ 有界性(定理~\ref{thm:1.3.5}),
\begin{align*}
|E_{1}| &\leq \dfrac{1}{(\lambda/2)^{2}}\displaystyle\int_{E_{1}}|Hg(x)|^{2}\dd x \leq \dfrac{4}{\lambda^{2}}\displaystyle\int_{\mathbb{R}}|Hg(x)|^{2}\dd x \\
\quad &\leq \dfrac{4}{\lambda^{2}}\displaystyle\int_{\mathbb{R}}|g(x)|^{2}\dd x = \dfrac{4}{\lambda^{2}}\left(\displaystyle\int_{\Omega}|g(x)|^{2}\dd x + \displaystyle\int_{\Omega^{c}}|g(x)|^{2}\dd x\right) \\
\quad &= \dfrac{4}{\lambda^{2}}\left(\displaystyle\sum_{j}|Q_{j}||f_{Q_{j}}|^{2} + \lambda\displaystyle\int_{\Omega^{c}}|g(x)|\dd x\right) \quad (|g(x)| \leq \lambda,\ \text{a.e. } x \in \Omega^{c}) \\
\quad &\leq \dfrac{4}{\lambda^{2}}\left(\displaystyle\sum_{j}|Q_{j}|(2\lambda)^{2} + \lambda\displaystyle\int_{\mathbb{R}}|f(x)|\dd x\right) \\
&\quad \quad \quad \quad \quad \quad \quad \quad \quad \quad \quad \quad (|f_{Q_{j}}| \leq 2\lambda,\ g(x) = f(x)\ (x \in \Omega^{c})) \\
&= 16\displaystyle\sum_{j}|Q_{j}| + \dfrac{4}{\lambda}\| f \|_{L^{1}} \leq \dfrac{20}{\lambda}\| f \|_{L^{1}}.
\end{align*}
其次,若 $x \in (Q_{j}^{*})^{c}$,则
\begin{align*}
Hf_{j}(x) & = \displaystyle\lim_{\varepsilon \to 0}\frac{1}{\pi}\int_{|y - x| > \varepsilon}\frac{1}{x - y}f_{j}(y)\dd y \\
& = \dfrac{1}{\pi}\displaystyle\int_{Q_{j}}\frac{1}{x - y}f_{j}(y)\dd y - \dfrac{1}{\pi}\int_{Q_{j}}\frac{1}{x - x_{j}}f_{j}(y)\dd y \\
& \quad \left(\because \displaystyle\int_{Q_{j}} f_{j}(y)\dd y = 0\right) \\
& = \dfrac{1}{\pi}\displaystyle\int_{Q_{j}}\left(\frac{1}{x - y} - \frac{1}{x - x_{j}}\right)f_{j}(y)\dd y
\end{align*}
又若 $y \in Q_{j}$,$x \in (Q_{j}^{*})^{c}$,则 $|x - x_{j}| \geq \ell(Q_{j}) > 2|x_{j} - y|$,故 $|x - y| \geq |x - x_{j}| - |x_{j} - y| \geq |x - x_{j}|/2$,$|x - x_{j}|^{2} + \ell(Q_{j})^{2} \leq 2|x - x_{j}|^{2}$,从而
\begin{equation*}
\left|\frac{1}{x - y} - \frac{1}{x - x_{j}}\right| = \frac{|y - x_{j}|}{|x - y||x - x_{j}|} \leq \frac{2\ell(Q_{j})}{|x - x_{j}|^{2} + \ell(Q_{j})^{2}}.
\end{equation*}
于是若 $x \in (Q_{j}^{*})^{c}$,则
\begin{align*}
|Hb(x)| &\leq \displaystyle\sum_{j}|Hf_{j}(x)| \leq \sum_{j}\frac{1}{\pi}\int_{Q_{j}}\frac{2\ell(Q_{j})}{|x - x_{j}|^{2} + \ell(Q_{j})^{2}}|f_{j}(y)|\dd y \\
&\leq \dfrac{8\lambda}{\pi}\displaystyle\sum_{j}\frac{\ell(Q_{j})^{2}}{|x - x_{j}|^{2} + \ell(Q_{j})^{2}}.
\end{align*}
由此
\begin{align*}
|E_{2}| &\leq \dfrac{1}{\lambda/2}\displaystyle\int_{E_{2}}|Hb(x)|\dd x \leq \dfrac{2}{\lambda}\int_{(\Omega^{*})^{c}}|Hb(x)|\dd x \\
\quad &\leq \dfrac{16}{\pi}\displaystyle\int_{(\Omega^{*})^{c}}\sum_{j}\frac{\ell(Q_{j})^{2}}{|x - x_{j}|^{2} + \ell(Q_{j})^{2}}\dd x \\
\quad &\leq \dfrac{16}{\pi}\displaystyle\sum_{j}\ell(Q_{j})\int_{-\infty}^{\infty}\frac{\ell(Q_{j})}{|x - x_{j}|^{2} + \ell(Q_{j})^{2}}\dd x \\
\quad &= \dfrac{16}{\pi}\displaystyle\sum_{j}|Q_{j}|\int_{-\infty}^{\infty}\frac{1}{t^{2} + 1}\dd t \quad (t = (x - x_{j})/\ell(Q_{j})) \\
\quad &\leq \dfrac{16}{\lambda}\| f \|_{L^{1}}.
\end{align*}
因此
\begin{equation}\tag{1.8.4}
\left|\left\{x \in \mathbb{R};\ |Hf(x)| > \lambda,\ x \in (\Omega^{*})^{c}\right\}\right| \leq |E_{1}| + |E_{2}| \leq \frac{20 + 16}{\lambda}\| f \|_{L^{1}}.
\label{eq:1.8.4}\end{equation}
由 \eqref{eq:1.8.3},\eqref{eq:1.8.4},
\begin{align*}
|\{x \in \mathbb{R};\ |Hf(x)| > \lambda\}| & \leq |\Omega^{*}| + |\{x \in \mathbb{R};\ |Hf(x)| > \lambda,\ x \in (\Omega^{*})^{c}\}| \\
& \leq \dfrac{2 + 20 + 16}{\lambda}\| f \|_{L^{1}} = \dfrac{38}{\lambda}\| f \|_{L^{1}}.
\end{align*}
\end{proof}

\begin{note}\label{note:1.8.3}
把定理~\ref{thm:1.8.2} 证明开头所述 $f(x) = g(x) + b(x)$ 的分解,也常称为 Calderón--Zygmund 分解.另外,定理~\ref{thm:1.8.2} 可唯一地扩张到 $L^{1}(\mathbb{R})$ 上,使它对一切 $f \in L^{1}(\mathbb{R})$ 成立.这一点由下面的引理保证.
\end{note}

\begin{lemma}\label{lem:1.8.4}
\begin{sloppypar}
设 $(X, \mu)$ 为完备测度空间,$E \subset L^{1}(\mu)$ 且 $E$ 在 $L^{1}(\mu)$ 中稠密.若 $T$ 是从 $E$ 到 $\{\mu\text{-可测函数的全体}\}$ 的线性映射,满足
\end{sloppypar}
\begin{equation*}
\mu(\{x \in X;\ |Tf(x)| > \lambda\}) \leq \frac{C}{\lambda}\| f \|_{L^{1}(\mu)} \quad (\lambda > 0,\ f \in E),
\end{equation*}
则 $T$ 可唯一扩张为满足弱 $(1,1)$ 性的线性映射到 $L^{1}(\mu)$ 上.扩张后的映射仍记为 $T$,则

(i)
\begin{equation*}
\mu(\{x \in X;\ |Tf(x)| > \lambda\}) \leq \frac{C}{\lambda}\| f \|_{L^{1}(\mu)} \quad (\lambda > 0,\ f \in L^{1}(\mu)).
\end{equation*}

(ii) 若 $\lim_{j \to \infty}\| f_{j} - f \|_{L^{1}(\mu)} = 0$,则对任意 $\varepsilon > 0$ 有 $\lim_{j \to \infty}\mu(\{x \in X;\ |Tf_{j}(x) - Tf(x)| > \varepsilon\}) = 0$,即当 $j \to \infty$ 时 $Tf_{j}$ 依测度收敛于 $Tf$.
\end{lemma}

\begin{proof}
设 $f \in L^{1}(\mu) \setminus E$.因 $E$ 在 $L^{1}(\mu)$ 中稠密,存在 $f_{j} \in E$($j = 1, 2, \ldots$)满足
\begin{equation*}
\| f_{j} - f_{j+1} \|_{L^{1}(\mu)} \leq 2^{-2j}, \quad \lim_{j \to \infty}\| f_{j} - f \|_{L^{1}(\mu)} = 0.
\end{equation*}
对这个 $\{f_{j}\}$ 令
\begin{equation*}
F_{j} := \{x \in X;\ |Tf_{j}(x) - Tf_{j+1}(x)| > 2^{-j}\}.
\end{equation*}
则
\begin{equation*}
\mu\left(\bigcup_{j = k}^{\infty} F_{j}\right) \leq \sum_{j = k}^{\infty}\mu(F_{j}) \leq \sum_{j = k}^{\infty} C2^{j}\| f_{j} - f_{j+1} \|_{L^{1}(\mu)} \leq C\sum_{j = k}^{\infty}2^{-j} = C2^{-k+1}
\end{equation*}
故
\begin{equation*}
\mu\left(\bigcap_{k = 1}^{\infty}\bigcup_{j = k}^{\infty} F_{j}\right) = 0.
\end{equation*}
于是令
\begin{equation*}
N = \bigcap_{k = 1}^{\infty}\bigcup_{j = k}^{\infty} F_{j}.
\end{equation*}
若 $x \in N^{c}$,则存在编号 $k_{0}$ 使 $x \in \bigcap_{j = k_{0}}^{\infty}F_{j}^{c}$,故
\begin{equation*}
|Tf_{j}(x) - Tf_{j+1}(x)| \leq 2^{-j} \quad (j \geq k_{0}).
\end{equation*}
由此,由复数的完备性,数列 $\{Tf_{j}(x)\}_{j=1}^{\infty}$ 在 $j \to \infty$ 时收敛到某个数.将其值记为 $Tf(x)$.又当 $x \in N$ 时令 $Tf(x) = 0$,则 $Tf(x)$ 对所有 $x \in X$ 有定义,且 $\mu$-可测函数 $Tf_{j}$ 在几乎所有的点上收敛到 $Tf(x)$,故 $Tf$ 也是 $\mu$-可测函数.现在,若 $x \in \bigcap_{j = k}^{\infty}F_{j}^{c} \subset N^{c}$,则
\begin{align*}
|Tf_{k}(x) - Tf(x)| & = \displaystyle\lim_{j \to \infty}|Tf_{k}(x) - Tf_{j}(x)| \\
& \leq \displaystyle\lim_{j \to \infty}\sum_{\ell = k}^{j-1}|Tf_{\ell}(x) - Tf_{\ell+1}(x)| \leq \sum_{\ell = k}^{\infty}2^{-\ell} = 2^{1-k}.
\end{align*}
故
\begin{equation*}
\mu(\{x \in X;\ |Tf_{k}(x) - Tf(x)| > 2 \cdot 2^{-k}\}) \leq \mu\left(\bigcup_{j = k}^{\infty} F_{j}\right) \leq 2C2^{-k}.
\end{equation*}
其次,设 $\lambda > 0$,$g \in E$.取编号 $k$ 满足
\begin{align*}
\| f - g \|_{L^{1}(\mu)} &> \max(\lambda, 1)2^{-k+1}, \quad \lambda - 2^{-k+1} > \dfrac{\lambda}{2}, \\
\| f_{k} - f \|_{L^{1}(\mu)} &\leq \| f - g \|_{L^{1}(\mu)}.
\end{align*}
则当 $|Tf(x) - Tf_{k}(x)| \leq 2^{-k+1}$ 且 $|Tf_{k}(x) - Tg(x)| \leq \lambda - 2^{-k+1}$ 时,$|Tf(x) - Tg(x)| \leq |Tf(x) - Tf_{k}(x)| + |Tf_{k}(x) - Tg(x)| \leq \lambda$,故
\begin{align*}
& \mu(\{x \in X;\ |Tf(x) - Tg(x)| > \lambda\}) \\
& \qquad \leq \mu(\{x \in X;\ |Tf(x) - Tf_{k}(x)| > 2^{-k+1}\}) \\
& \qquad\qquad + \mu(\{x \in X;\ |Tf_{k}(x) - Tg(x)| > \lambda - 2^{-k+1}\}) \\
& \qquad \leq 2C2^{-k} + \dfrac{C}{\lambda - 2^{-k+1}}\| f_{k} - g \|_{L^{1}(\mu)} \\
& \qquad \leq \dfrac{C\| f - g \|_{L^{1}(\mu)}}{\max(\lambda, 1)} + \dfrac{2C}{\lambda}\| f_{k} - g \|_{L^{1}(\mu)} \\
& \qquad \leq \dfrac{C\| f - g \|_{L^{1}(\mu)}}{\lambda} + \dfrac{2C}{\lambda}\left(\| f - g \|_{L^{1}(\mu)} + \| f_{k} - f \|_{L^{1}(\mu)}\right)
\end{align*}
\begin{equation*}
\leq \frac{5C}{\lambda}\| f - g \|_{L^{1}(\mu)}.
\end{equation*}
也就是说,若 $g \in E$ 依 $L^{1}(\mu)$ 收敛于 $f$,则 $Tg$ 依测度收敛于 $Tf$.因此,定义在 $E$ 上的 $T$ 由稠密性可唯一扩张到 $L^{1}(\mu)$ 上.

进而,对 $0 < \varepsilon < 1$,因 $\{x \in X;\ |Tf(x)| > \lambda\} \subset \{x \in X;\ |Tf(x) - Tg(x)| > \varepsilon\lambda\} \cup \{x \in X;\ |Tg(x)| > (1 - \varepsilon)\lambda\}$,故
\begin{align*}
\mu(\{x \in X;\ |Tf(x)| > \lambda\}) &\leq \mu(\{x \in X;\ |Tf(x) - Tg(x)| > \varepsilon\lambda\})  \\
& \quad + \mu(\{x \in X;\ |Tg(x)| > (1 - \varepsilon)\lambda\}) \\
& \leq \dfrac{5C}{\varepsilon\lambda}\| f - g \|_{L^{1}(\mu)} + \dfrac{C}{(1 - \varepsilon)\lambda}\| g \|_{L^{1}(\mu)} \\
& \leq \left(\dfrac{5C}{\varepsilon\lambda} + \dfrac{C}{(1 - \varepsilon)\lambda}\right)\| f - g \|_{L^{1}(\mu)} + \dfrac{C}{(1 - \varepsilon)\lambda}\| f \|_{L^{1}(\mu)}.
\end{align*}
于是取 $g_{j} \in E$ 收敛于 $f$,令 $j \to \infty$ 后再令 $\varepsilon \to 0$,得
\begin{equation*}
\mu(\{x \in X;\ |Tf(x)| > \lambda\}) \leq \frac{C}{\lambda}\| f \|_{L^{1}(\mu)}.
\end{equation*}
$T$ 的线性性也容易确认.
\end{proof}

作为相关的话题,我们来叙述下面的 Kolmogorov 不等式.

\begin{theorem}[Kolmogorov 不等式]\label{thm:1.8.5}
设 $(X, \mu)$ 为测度空间.若存在常数 $A > 0$ 使得
\begin{equation*}
\mu(\{x \in X;\ |g(x)| > \lambda\}) \leq \frac{A}{\lambda} \quad (\lambda > 0),
\end{equation*}
则对任意 $0 < \delta < 1$ 与 $\mu(E) < \infty$ 的 $\mu$-可测集 $E$,存在常数 $C_{\delta}$ 使得
\begin{equation*}
\int_{E}|g(x)|^{\delta}\dd \mu \leq C_{\delta}A^{\delta}\mu(E)^{1-\delta}.
\end{equation*}
\end{theorem}

\begin{proof}
不妨设 $A = 1$.对 $k \in \mathbb{Z}$ 令
\begin{equation*}
E_{k} = \{x \in E;\ |g(x)| > 2^{k}\},
\end{equation*}
并令 $k_{0}$ 为使 $2^{-k_{0}} \geq \mu(E)$ 的最大整数.则 $\mu(E_{k}) \leq 2^{-k}$,且
\begin{align*}
\displaystyle\int_{E}|g(x)|^{\delta}\dd \mu &= \sum_{k = -\infty}^{\infty}\int_{E_{k} \setminus E_{k+1}}|g(x)|^{\delta}\dd \mu \\
&\leq \displaystyle\sum_{k = -\infty}^{\infty}2^{(k+1)\delta}\left(\mu(E_{k}) - \mu(E_{k+1})\right) \\
&= \displaystyle\sum_{k = -\infty}^{\infty}(2^{\delta} - 1)2^{k\delta}\mu(E_{k}) \\
&= \displaystyle\sum_{k < k_{0}}(2^{\delta} - 1)2^{k\delta}\mu(E_{k}) + \sum_{k \geq k_{0}}(2^{\delta} - 1)2^{k\delta}\mu(E_{k}) \\
&\leq \displaystyle\sum_{k < k_{0}}(2^{\delta} - 1)2^{k\delta}\mu(E) + \sum_{k \geq k_{0}}(2^{\delta} - 1)2^{k\delta}2^{-k} \\
&= (2^{\delta} - 1)\dfrac{2^{(k_{0}-1)\delta}}{1 - 2^{-\delta}}\mu(E) + (2^{\delta} - 1)\dfrac{2^{-k_{0}(1-\delta)}}{1 - 2^{\delta-1}} \\
&\leq (2^{\delta} - 1)\left(\dfrac{2^{-\delta}}{1 - 2^{-\delta}} + \dfrac{2^{1-\delta}}{1 - 2^{\delta-1}}\right)\mu(E)^{1-\delta}.
\end{align*}
\end{proof}

在本节最后注意如下.历史上,Hilbert 变换的 $L^{p}$ 有界性的证明顺序是:$L^{2}$ 有界性之后先证明弱 $(1,1)$ 性,对 $1 < p < 2$ 用 Marcinkiewicz 插值法证明 $L^{p}$ 有界性,对 $2 < p < \infty$ 则用对偶性来证明.用这个方法,Calderón 与 Zygmund 在 $n$ 维 Euclid 空间 $\mathbb{R}^{n}$ 上建立了``Calderón--Zygmund''奇异积分理论.Marcinkiewicz 插值法也是在证明 Hardy--Littlewood 极大函数的 $L^{p}$ 有界性时实质使用的方法,这里叙述其一般形式.

\begin{theorem}[Marcinkiewicz 插值定理]\label{thm:1.8.6}
设 $(X, \mu)$,$(Y, \nu)$ 为测度空间,$1 \leq p_{1} \leq q_{1} \leq \infty$,$1 \leq p_{2} \leq q_{2} \leq \infty$,且 $q_{1} \neq q_{2}$.设 $T$ 是从 $L^{p_{1}}(\mu) + L^{p_{2}}(\mu)$ 到 $\nu$-可测函数全体构成的函数空间的映射,满足

(i)
\begin{equation*}
|T(f_{1} + f_{2})(y)| \leq \kappa\left(|Tf_{1}(y)| + |Tf_{2}(y)|\right) \quad (\nu\text{-a.e. } y \in Y);
\end{equation*}

(ii)
\begin{align*}
& \nu(\{y \in Y;\ |Tf(y)| > \lambda\}) \leq \left(\dfrac{M_{1}}{\lambda}\| f \|_{L^{p_{1}}(\mu)}\right)^{q_{1}} \quad (\lambda > 0,\ f \in L^{p_{1}}(\mu)) \\[4pt]
& \left(q_{1} = \infty \text{ 时 } \| Tf \|_{L^{\infty}(\nu)} \leq M_{1}\| f \|_{L^{p_{1}}(\mu)}\right);
\end{align*}

(iii)
\begin{align*}
& \nu(\{y \in Y;\ |Tf(y)| > \lambda\}) \leq \left(\dfrac{M_{2}}{\lambda}\| f \|_{L^{p_{2}}(\mu)}\right)^{q_{2}} \quad (\lambda > 0,\ f \in L^{p_{2}}(\mu)) \\[4pt]
& \left(q_{2} = \infty \text{ 时 } \| Tf \|_{L^{\infty}(\nu)} \leq M_{2}\| f \|_{L^{p_{2}}(\mu)}\right).
\end{align*}

则对 $0 < t < 1$,对满足 $\dfrac{1}{p} = \dfrac{1 - t}{p_{1}} + \dfrac{t}{p_{2}}$,$\dfrac{1}{q} = \dfrac{1 - t}{q_{1}} + \dfrac{t}{q_{2}}$ 的 $p, q$,存在常数 $K$ 使得
\begin{equation*}
\| Tf \|_{L^{q}(\nu)} \leq KM_{1}^{1-t}M_{2}^{t}\| f \|_{L^{p}(\mu)} \quad (f \in L^{p}(\mu)).
\end{equation*}
此常数 $K$ 依赖于 $t, \kappa, p_{1}, p_{2}, q_{1}, q_{2}$;若 $0 < t_{1} < t < t_{2} < 1$,则仅依赖于 $t_{1}, t_{2}, \kappa, p_{1}, p_{2}, q_{1}, q_{2}$.
\end{theorem}

\begin{proof}
不妨设 $p_{1} \leq p_{2}$.设 $f \in L^{p}(\mu)$,且 $\dfrac{1}{p} = \dfrac{1 - t}{p_{1}} + \dfrac{t}{p_{2}}$,$\dfrac{1}{q} = \dfrac{1 - t}{q_{1}} + \dfrac{t}{q_{2}}$.

(a-1) $p_{1} < p_{2} < \infty$,$q_{1} < q_{2} < \infty$ 的情形.对 $a > 0$ 定义
\begin{equation*}
f_{1,a}(x) = \begin{cases} 0 & (|f(x)| \leq a) \\ f(x) & (|f(x)| > a), \end{cases}
\end{equation*}
\begin{equation*}
f_{2,a}(x) = \begin{cases} f(x) & (|f(x)| \leq a) \\ 0 & (|f(x)| > a). \end{cases}
\end{equation*}
因 $p_{1} < p$,$a^{p-p_{1}}|f_{1,a}(x)|^{p_{1}} \leq |f_{1,a}(x)|^{p}$,$|f_{1,a}(x)| \leq |f(x)| \in L^{p}(\mu)$,故
\begin{equation}\tag{1.8.5}
\int_{X}|f_{1,a}(x)|^{p_{1}}\dd \mu \leq a^{p_{1}-p}\int_{X}|f_{1,a}(x)|^{p}\dd \mu \leq a^{p_{1}-p}\| f \|_{L^{p}(\mu)}^{p} < \infty.
\label{eq:1.8.5}\end{equation}
又因 $p < p_{2}$,$|f_{2,a}(x)| \leq a$,$|f_{2,a}(x)| \leq |f(x)| \in L^{p}(\mu)$,故
\begin{equation}\tag{1.8.6}
\int_{X}|f_{2,a}(x)|^{p_{2}}\dd \mu \leq a^{p_{2}-p}\int_{X}|f_{2,a}(x)|^{p}\dd \mu \leq a^{p_{2}-p}\| f \|_{L^{p}(\mu)}^{p} < \infty.
\label{eq:1.8.6}\end{equation}
于是 $f = f_{1,a} + f_{2,a}$,$f_{1,a} \in L^{p_{1}}(\mu)$,$f_{2,a} \in L^{p_{2}}(\mu)$,故 $T(f) = T(f_{1,a} + f_{2,a})$ 有定义,由假设 (i),
\begin{equation}\tag{1.8.7}
|Tf(y)| \leq \kappa\left(|Tf_{1,a}(y)| + |Tf_{2,a}(y)|\right).
\label{eq:1.8.7}\end{equation}
这里令
\begin{equation*}
\begin{aligned}
m(\lambda) &= \mu(\{x \in X;\ |f(x)| > \lambda\}), \\
m_{1,a}(\lambda) &= \mu(\{x \in X;\ |f_{1,a}(x)| > \lambda\}), \\
m_{2,a}(\lambda) &= \mu(\{x \in X;\ |f_{2,a}(x)| > \lambda\}), \\
n(\lambda) &= \nu(\{y \in Y;\ |Tf(y)| > \lambda\}),
\end{aligned}
\end{equation*}
则
\begin{equation}\tag{1.8.8}
\| f \|_{L^{p}(\mu)}^{p} = \int_{X}|f(x)|^{p}\dd \mu = p\int_{0}^{\infty}\lambda^{p-1}m(\lambda)\dd \lambda,
\label{eq:1.8.8}\end{equation}
\begin{equation}\tag{1.8.9}
\begin{aligned}
\| Tf \|_{L^{q}(\nu)}^{q} &= \displaystyle\int_{Y}|Tf(y)|^{q}\dd \nu = q\int_{0}^{\infty}\lambda^{q-1}n(\lambda)\dd \lambda \\
&= (2\kappa)^{q}q\displaystyle\int_{0}^{\infty}\lambda^{q-1}n(2\kappa\lambda)\dd \lambda,
\end{aligned}
\label{eq:1.8.9}\end{equation}
且
\begin{equation}\tag{1.8.10}
m_{1,a}(s) = \begin{cases} m(a) & (0 < s < a), \\ m(s) & (s \geq a), \end{cases} \quad
m_{2,a}(s)\begin{cases} \leq m(s) & (0 < s < a), \\ = 0 & (s \geq a). \end{cases}
\label{eq:1.8.10}\end{equation}
其次,设 $\xi, A$ 为后面确定的常数,对 $\lambda > 0$ 令 $a = (\lambda/A)^{\xi}$.则由式~\eqref{eq:1.8.7},若 $|Tf_{1,a}(y)| \leq \lambda$ 且 $|Tf_{2,a}(y)| \leq \lambda$,则 $|Tf(y)| \leq 2\kappa\lambda$,故
\begin{equation*}
n(2\kappa\lambda) \leq \nu(\{y \in Y;\ |Tf_{1,a}(y)| > \lambda\}) + \nu(\{y \in Y;\ |Tf_{2,a}(y)| > \lambda\}).
\end{equation*}
因此,由假设 (ii),(iii) 与 \eqref{eq:1.8.5},\eqref{eq:1.8.6},
\begin{align*}
n(2\kappa\lambda) & \leq M_{1}^{q_{1}}\lambda^{-q_{1}}\| f_{1,a} \|_{L^{p_{1}}(\mu)}^{q_{1}} + M_{2}^{q_{2}}\lambda^{-q_{2}}\| f_{2,a} \|_{L^{p_{2}}(\mu)}^{q_{2}} \\
& \leq M_{1}^{q_{1}}\lambda^{-q_{1}}\left(a^{p_{1}-p}\| f \|_{L^{p}(\mu)}^{p}\right)^{q_{1}/p_{1}-1}\| f_{1,a} \|_{L^{p_{1}}(\mu)}^{p_{1}} \\
& \quad + M_{2}^{q_{2}}\lambda^{-q_{2}}\left(a^{p_{2}-p}\| f \|_{L^{p}(\mu)}^{p}\right)^{q_{2}/p_{2}-1}\| f_{2,a} \|_{L^{p_{2}}(\mu)}^{p_{2}}.
\end{align*}
由此与 \eqref{eq:1.8.9},
\begin{align*}
\| Tf \|_{L^{q}(\nu)}^{q} & = (2\kappa)^{q}qM_{1}^{q_{1}}\| f \|_{L^{p}(\mu)}^{p(q_{1}/p_{1}-1)} \\
& \quad \times \displaystyle\int_{0}^{\infty}\lambda^{q-q_{1}-1}a^{(p_{1}-p)(q_{1}/p_{1}-1)}\| f_{1,a} \|_{L^{p_{1}}(\mu)}^{p_{1}}\dd \lambda \\
& \quad + (2\kappa)^{q}qM_{2}^{q_{2}}\| f \|_{L^{p}(\mu)}^{p(q_{2}/p_{2}-1)} \\
& \quad \times \displaystyle\int_{0}^{\infty}\lambda^{q-q_{2}-1}a^{(p_{2}-p)(q_{2}/p_{2}-1)}\| f_{2,a} \|_{L^{p_{2}}(\mu)}^{p_{2}}\dd \lambda \\
& =: (2\kappa)^{q}qM_{1}^{q_{1}}\| f \|_{L^{p}(\mu)}^{p(q_{1}/p_{1}-1)}I_{1} + (2\kappa)^{q}qM_{2}^{q_{2}}\| f \|_{L^{p}(\mu)}^{p(q_{2}/p_{2}-1)}I_{2}.
\end{align*}
(最后一行表示把倒数第二行中的积分部分分别记为 $I_{1}, I_{2}$.)

这里令 $\xi = \dfrac{1/q_{1} - 1/q_{2}}{1/p_{1} - 1/p_{2}}\cdot\dfrac{q}{p}$.此时,
\begin{equation}\tag{1.8.11}
\xi = \frac{q - q_{1}}{p - p_{1}}\cdot\frac{p_{1}}{q_{1}} = \frac{q - q_{2}}{p - p_{2}}\cdot\frac{p_{2}}{q_{2}},
\label{eq:1.8.11}\end{equation}
\begin{equation}\tag{1.8.12}
\frac{q - q_{1}}{\xi} + (p_{1} - p)\left(\frac{q_{1}}{p_{1}} - 1\right) + p_{1} = p,
\label{eq:1.8.12}\end{equation}
\begin{equation}\tag{1.8.13}
\frac{q - q_{2}}{\xi} + (p_{2} - p)\left(\frac{q_{2}}{p_{2}} - 1\right) + p_{2} = p
\label{eq:1.8.13}\end{equation}
成立(实际上,正是为了满足上式才这样确定 $\xi$).因令 $a = (\lambda/A)^{\xi}$,故利用 \eqref{eq:1.8.10} 得
\begin{align*}
I_{1}&= A^{-\xi(p_{1}-p)(q_{1}/p_{1}-1)}\displaystyle\int_{0}^{\infty}\lambda^{q-q_{1}-1+\xi(p_{1}-p)(q_{1}/p_{1}-1)}\left(\int_{0}^{\infty}p_{1}s^{p_{1}-1}m_{1,a}(s)\dd s\right)\dd \lambda  \\
&= A^{-\xi(p_{1}-p)(q_{1}/p_{1}-1)}\displaystyle\int_{0}^{\infty}\lambda^{q-q_{1}-1+\xi(p_{1}-p)(q_{1}/p_{1}-1)}\left(\int_{(\lambda/A)^{\xi}}^{\infty}p_{1}s^{p_{1}-1}m(s)\dd s\right)\dd \lambda \\
&\quad + A^{-\xi(p_{1}-p)(q_{1}/p_{1}-1)}\displaystyle\int_{0}^{\infty}\lambda^{q-q_{1}-1+\xi(p_{1}-p)(q_{1}/p_{1}-1)}\left(\frac{\lambda}{A}\right)^{\xi p_{1}}m\left(\left(\frac{\lambda}{A}\right)^{\xi}\right)\dd \lambda \\
&= A^{-\xi(p_{1}-p)(q_{1}/p_{1}-1)}\displaystyle\int_{0}^{\infty}p_{1}s^{p_{1}-1}m(s)\left(\int_{0}^{As^{1/\xi}}\lambda^{q-q_{1}-1+\xi(p_{1}-p)(q_{1}/p_{1}-1)}\dd \lambda\right)\dd s \\
&\quad + A^{-\xi(p_{1}-p)(q_{1}/p_{1}-1)}\displaystyle\int_{0}^{\infty}\lambda^{q-q_{1}-1+\xi(p_{1}-p)(q_{1}/p_{1}-1)}\left(\frac{\lambda}{A}\right)^{\xi p_{1}}m\left(\left(\frac{\lambda}{A}\right)^{\xi}\right)\dd \lambda \\
&= \frac{p_{1}A^{q-q_{1}}}{q - q_{1} + \xi(p_{1}-p)(q_{1}/p_{1}-1)}\displaystyle\int_{0}^{\infty}s^{p_{1}-1+(q-q_{1})/\xi+(p_{1}-p)(q_{1}/p_{1}-1)}m(s)\dd s
\end{align*}
\begin{equation*}
+ \frac{A^{q-q_{1}}}{\xi}\int_{0}^{\infty}s^{p_{1}-1+(q-q_{1})/\xi+(p_{1}-p)(q_{1}/p_{1}-1)}m(s)\dd s.
\end{equation*}
(在由倒数第二个式子的第一个积分项变为最后一个式子的过程中,除了对倒数第二,第三个式子的第一项使用 Fubini 定理外,还利用了 \eqref{eq:1.8.12},\eqref{eq:1.8.11} 推出的 $q - q_{1} - 1 + \xi(p_{1}-p)\left(\frac{q_{1}}{p_{1}}-1\right) = (q-q_{1})p_{1}/q_{1} - 1 > -1$.最后一步中作了变量替换 $\lambda = As^{1/\xi}$.)

对 $I_{2}$ 作同样的考察,
\begin{align*}
I_{2}&= A^{-\xi(p_{2}-p)(q_{2}/p_{2}-1)}\displaystyle\int_{0}^{\infty}\lambda^{q-q_{2}-1+\xi(p_{2}-p)(q_{2}/p_{2}-1)}\left(\int_{0}^{\infty}p_{2}s^{p_{2}-1}m_{2,a}(s)\dd s\right)\dd \lambda  \\
&\leq A^{-\xi(p_{2}-p)(q_{2}/p_{2}-1)}\displaystyle\int_{0}^{\infty}\lambda^{q-q_{2}-1+\xi(p_{2}-p)(q_{2}/p_{2}-1)}\left(\int_{0}^{(\lambda/A)^{\xi}}p_{2}s^{p_{2}-1}m(s)\dd s\right)\dd \lambda \\
&= A^{-\xi(p_{2}-p)(q_{2}/p_{2}-1)}\displaystyle\int_{0}^{\infty}p_{2}s^{p_{2}-1}m(s)\left(\int_{As^{1/\xi}}^{\infty}\lambda^{q-q_{2}-1+\xi(p_{2}-p)(q_{2}/p_{2}-1)}\dd \lambda\right)\dd s \\
&= \frac{-p_{2}A^{q-q_{2}}}{q - q_{2} + \xi(p_{2}-p)(q_{2}/p_{2}-1)}\displaystyle\int_{0}^{\infty}s^{p_{2}-1+(q-q_{2})/\xi+(p_{2}-p)(q_{2}/p_{2}-1)}m(s)\dd s
\end{align*}
因此,利用 $\xi$ 的定义,\eqref{eq:1.8.11},\eqref{eq:1.8.12},\eqref{eq:1.8.13} 整理得
\begin{align*}
I_{1} &= \dfrac{1/p_{1} - 1/p_{2}}{1/q_{1} - 1/q_{2}}\dfrac{pA^{q-q_{1}}}{q(p - p_{1})}\| f \|_{L^{p}(\mu)}^{p}, \\
I_{2} &\leq \dfrac{1/p_{1} - 1/p_{2}}{1/q_{1} - 1/q_{2}}\dfrac{p_{2}A^{q-q_{2}}}{q(p_{2} - p)}\| f \|_{L^{p}(\mu)}^{p}.
\end{align*}
由 $I_{1}, I_{2}$ 的估计式,
\begin{align*}
\| Tf \|_{L^{q}(\nu)}^{q} & \leq (2\kappa)^{q}M_{1}^{q_{1}}A^{q-q_{1}}\dfrac{1/p_{1} - 1/p_{2}}{1/q_{1} - 1/q_{2}}\dfrac{p}{p - p_{1}}\| f \|_{L^{p}(\mu)}^{pq_{1}/p_{1}} \\
& \quad + (2\kappa)^{q}M_{2}^{q_{2}}A^{q-q_{2}}\dfrac{1/p_{1} - 1/p_{2}}{1/q_{1} - 1/q_{2}}\dfrac{p_{2}}{p_{2} - p}\| f \|_{L^{p}(\mu)}^{pq_{2}/p_{2}}.
\end{align*}
于是,选取 $A$ 使
\begin{equation*}
M_{1}^{q_{1}}A^{q-q_{1}}\| f \|_{L^{p}(\mu)}^{pq_{1}/p_{1}} = M_{2}^{q_{2}}A^{q-q_{2}}\| f \|_{L^{p}(\mu)}^{pq_{2}/p_{2}},
\end{equation*}
即令
\begin{equation*}
A = M_{1}^{q_{1}/(q_{1}-q_{2})}M_{2}^{q_{2}/(q_{2}-q_{1})}\| f \|_{L^{p}(\mu)}^{p(q_{2}/p_{2}-q_{1}/p_{1})/(q_{2}-q_{1})},
\end{equation*}
并注意到
\begin{equation*}
\frac{q - q_{1}}{q_{2} - q_{1}}q_{2} = qt, \qquad
\frac{pq_{1}}{p_{1}} + \frac{q - q_{1}}{q_{2} - q_{1}}\left(\frac{q_{2}}{p_{2}} - \frac{q_{1}}{p_{1}}\right)p = q,
\end{equation*}
则
\begin{equation*}
\| Tf \|_{L^{q}(\nu)}^{q} \leq (2\kappa)^{q}\frac{1/p_{1} - 1/p_{2}}{1/q_{1} - 1/q_{2}}\left(\frac{p}{p - p_{1}} + \frac{p_{2}}{p_{2} - p}\right)M_{1}^{q(1-t)}M_{2}^{qt}\| f \|_{L^{p}(\mu)}^{q}
\end{equation*}
因此
\begin{equation*}
\| Tf \|_{L^{q}(\nu)} \leq (2\kappa)\left(\frac{1/p_{1} - 1/p_{2}}{1/q_{1} - 1/q_{2}}\left(\frac{p}{p - p_{1}} + \frac{p_{2}}{p_{2} - p}\right)\right)^{1/q}M_{1}^{1-t}M_{2}^{t}\| f \|_{L^{p}(\mu)}
\end{equation*}

(a-2) $p_{1} < p_{2} < \infty$,$q_{2} < q_{1} < \infty$ 的情形.与 (a-1) 的情形几乎相同,但令 $\xi = \dfrac{1/q_{1} - 1/q_{2}}{1/p_{1} - 1/p_{2}}\cdot\dfrac{q}{p}$ 时,$\xi$ 为负这一差别会出现.从那里开始,
\begin{align*}
I_{1}&= A^{-\xi(p_{1}-p)(q_{1}/p_{1}-1)}\displaystyle\int_{0}^{\infty}p_{1}s^{p_{1}-1}m(s)\left(\int_{As^{1/\xi}}^{\infty}\lambda^{q-q_{1}-1+\xi(p_{1}-p)(q_{1}/p_{1}-1)}\dd \lambda\right)\dd s  \\
&\quad + A^{-\xi(p_{1}-p)(q_{1}/p_{1}-1)}\displaystyle\int_{0}^{\infty}\lambda^{q-q_{1}-1+\xi(p_{1}-p)(q_{1}/p_{1}-1)}\left(\frac{\lambda}{A}\right)^{\xi p_{1}}m\left(\left(\frac{\lambda}{A}\right)^{\xi}\right)\dd \lambda \\
&= \frac{-p_{1}A^{q-q_{1}}}{q - q_{1} + \xi(p_{1}-p)(q_{1}/p_{1}-1)}\displaystyle\int_{0}^{\infty}s^{p_{1}-1+(q-q_{1})/\xi+(p_{1}-p)(q_{1}/p_{1}-1)}m(s)\dd s \\
&\quad - \frac{A^{q-q_{1}}}{\xi}\displaystyle\int_{0}^{\infty}s^{p_{1}-1+(q-q_{1})/\xi+(p_{1}-p)(q_{1}/p_{1}-1)}m(s)\dd s \\
&= \dfrac{1/p_{1} - 1/p_{2}}{1/q_{2} - 1/q_{1}}\dfrac{pA^{q-q_{1}}}{q(p - p_{1})}\| f \|_{L^{p}(\mu)}^{p}. \\
&I_{2} \\
&\leq A^{-\xi(p_{2}-p)(q_{2}/p_{2}-1)}\displaystyle\int_{0}^{\infty}p_{2}s^{p_{2}-1}m(s)\left(\int_{0}^{As^{1/\xi}}\lambda^{q-q_{2}-1+\xi(p_{2}-p)(q_{2}/p_{2}-1)}\dd \lambda\right)\dd s \\
&= \dfrac{p_{2}A^{q-q_{2}}}{q - q_{2} + \xi(p_{2}-p)(q_{2}/p_{2}-1)}\displaystyle\int_{0}^{\infty}s^{p_{2}-1+(q-q_{2})/\xi+(p_{2}-p)(q_{2}/p_{2}-1)}m(s)\dd s
\end{align*}
\begin{equation*}
= \frac{1/p_{1} - 1/p_{2}}{1/q_{2} - 1/q_{1}}\frac{p_{2}A^{q-q_{2}}}{q(p_{2} - p)}\| f \|_{L^{p}(\mu)}^{p}
\end{equation*}
以下与 (a-1) 同样选取 $A$,得
\begin{equation*}
\| Tf \|_{L^{q}(\nu)} \leq (2\kappa)\left(\frac{1/p_{1} - 1/p_{2}}{1/q_{2} - 1/q_{1}}\left(\frac{p}{p - p_{1}} + \frac{p_{2}}{p_{2} - p}\right)\right)^{1/q}M_{1}^{1-t}M_{2}^{t}\| f \|_{L^{p}(\mu)}.
\end{equation*}

(a-3) $p_{1} < p_{2} < \infty$,$q_{1} < q_{2} = \infty$ 的情形.在 (a-1) 中令 $a = (\lambda/A)^{\xi}$,$\xi = p_{2}/(p_{2} - p)$,$A = M_{2}\| f \|_{L^{p}(\mu)}^{p/p_{2}}$,则由假设 (iii),
\begin{equation*}
\| Tf_{2,a}(y) \|_{L^{\infty}(\nu)} \leq M_{2}\| f_{2,a} \|_{L^{p_{2}}(\mu)} \leq M_{2}\left(a^{p_{2}-p}\| f \|_{L^{p}(\mu)}^{p}\right)^{1/p_{2}} \leq \lambda,
\end{equation*}
故 $\nu(\{y \in Y;\ |Tf_{2,a}(y)| > \lambda\}) = 0$.因此 $n(2\kappa\lambda) \leq \nu(\{y \in Y;\ |Tf_{1,a}(y)| > \lambda\})$,从而
\begin{align*}
\| Tf \|_{L^{q}(\nu)}^{q} & \leq (2\kappa)^{q}M_{1}^{q_{1}}A^{q-q_{1}}\dfrac{1/p_{1} - 1/p_{2}}{1/q_{1}}\dfrac{p}{p - p_{1}}\| f \|_{L^{p}(\mu)}^{pq_{1}/p_{1}} \\
& = (2\kappa)^{q}M_{1}^{q_{1}}M_{2}^{q-q_{1}}\dfrac{1/p_{1} - 1/p_{2}}{1/q_{1}}\dfrac{p}{p - p_{1}}\| f \|_{L^{p}(\mu)}^{p(q_{1}/p_{1} + (q-q_{1})/p_{2})}.
\end{align*}
注意 $1/q = (1-t)/q_{1}$,$q_{1}/(qp_{1}) + (1 - q_{1}/q)/p_{2} = (1-t)/p_{1} + t/p_{2} = 1/p$,计算得
\begin{equation*}
\| Tf \|_{L^{q}(\nu)} \leq 2\kappa M_{1}^{1-t}M_{2}^{t}\left(\frac{1/p_{1} - 1/p_{2}}{1/q_{1}}\frac{p}{p - p_{1}}\right)^{1/q}\| f \|_{L^{p}(\mu)}.
\end{equation*}

(a-4) $p_{1} < p_{2} < \infty$,$q_{2} < q_{1} = \infty$ 的情形.在 (a-1) 中令 $a = (\lambda/A)^{\xi}$,$\xi = p_{1}/(p_{1} - p)$,$A = M_{1}\| f \|_{L^{p}(\mu)}^{p/p_{1}}$,则由假设 (ii),
\begin{equation*}
\| Tf_{1,a}(y) \|_{L^{\infty}(\nu)} \leq M_{1}\| f_{1,a} \|_{L^{p_{1}}(\mu)} \leq M_{1}\left(a^{p_{1}-p}\| f \|_{L^{p}(\mu)}^{p}\right)^{1/p_{1}} = \lambda,
\end{equation*}
故 $\nu(\{y \in Y;\ |Tf_{1,a}(y)| > \lambda\}) = 0$.因此 $n(2\kappa\lambda) \leq \nu(\{y \in Y;\ |Tf_{2,a}(y)| > \lambda\})$,从而
\begin{align*}
\| Tf \|_{L^{q}(\nu)}^{q} & \leq (2\kappa)^{q}M_{2}^{q_{2}}A^{q-q_{2}}\dfrac{1/p_{1} - 1/p_{2}}{1/q_{2}}\dfrac{p_{2}}{p_{2} - p}\| f \|_{L^{p}(\mu)}^{pq_{2}/p_{2}} \\
& = (2\kappa)^{q}M_{1}^{q-q_{2}}M_{2}^{q_{2}}\dfrac{1/p_{1} - 1/p_{2}}{1/q_{2}}\dfrac{p}{p_{2} - p}\| f \|_{L^{p}(\mu)}^{p(q-q_{2})/p_{1} + pq_{2}/p_{2}}.
\end{align*}
注意 $1/q = t/q_{2}$,$(1 - q_{2}/q)/p_{1} + q_{2}/(qp_{2}) = (1-t)/p_{1} + t/p_{2} = 1/p$,计算得
\begin{equation*}
\| Tf \|_{L^{q}(\nu)} \leq 2\kappa M_{1}^{1-t}M_{2}^{t}\left(\frac{1/p_{1} - 1/p_{2}}{1/q_{2}}\frac{p}{p_{2} - p}\right)^{1/q}\| f \|_{L^{p}(\mu)}.
\end{equation*}

(b) $p_{1} < p_{2} = \infty$,$q_{1} < q_{2} = \infty$ 的情形.在 (a-1) 中令 $a = (\lambda/A)^{\xi}$,$\xi = 1$,$A = M_{2}$,则由假设 (iii),
\begin{equation*}
\| Tf_{2,a}(y) \|_{L^{\infty}(\nu)} \leq M_{2}\| f_{2,a} \|_{L^{\infty}(\mu)} \leq M_{2}a = \lambda,
\end{equation*}
以下与 (a-3) 同样进行,得
\begin{equation*}
\| Tf \|_{L^{q}(\nu)} \leq 2\kappa M_{1}^{1-t}M_{2}^{t}\left(\frac{q_{1}}{p_{1}}\frac{p}{p - p_{1}}\right)^{1/q}\| f \|_{L^{p}(\mu)}.
\end{equation*}

(c-1) $p_{1} = p_{2} < \infty$,$q_{1} \neq q_{2}$,$q_{1}, q_{2} < \infty$ 的情形.此时不妨设 $p_{1} = p_{2} = p$,$q_{1} < q_{2}$.由假设 (ii),(iii),
\begin{equation*}
n(\lambda) \leq \left(\frac{M_{1}}{\lambda}\| f \|_{L^{p}(\mu)}\right)^{q_{1}}, \quad n(\lambda) \leq \left(\frac{M_{2}}{\lambda}\| f \|_{L^{p}(\mu)}\right)^{q_{2}},
\end{equation*}
故对 $A > 0$,
\begin{align*}
\| Tf \|_{L^{q}(\nu)}^{q} & = \displaystyle\int_{0}^{A}q\lambda^{q-1}n(\lambda)\dd \lambda + \int_{A}^{\infty}q\lambda^{q-1}n(\lambda)\dd \lambda \\
& \leq \displaystyle\int_{0}^{A}q\lambda^{q-q_{1}-1}n(\lambda)\dd \lambda\left(M_{1}\| f \|_{L^{p}(\mu)}\right)^{q_{1}} \\
& \quad + \displaystyle\int_{A}^{\infty}q\lambda^{q-q_{2}-1}n(\lambda)\dd \lambda\left(M_{2}\| f \|_{L^{p}(\mu)}\right)^{q_{2}} \\
& = M_{1}^{q_{1}}\dfrac{q}{q - q_{1}}A^{q-q_{1}}\| f \|_{L^{p}(\mu)}^{q_{1}} + M_{2}^{q_{2}}\dfrac{q}{q_{2} - q}A^{q-q_{2}}\| f \|_{L^{p}(\mu)}^{q_{2}}.
\end{align*}
与 (a-1) 的情形一样,令 $A = M_{1}^{q_{1}/(q-q_{1})}M_{2}^{-q_{2}/(q-q_{1})}\| f \|_{L^{p}(\mu)}^{q_{1}-q_{2}}$,得
\begin{equation*}
\| Tf \|_{L^{q}(\nu)} \leq \left(\frac{1}{q - q_{1}} + \frac{1}{q_{2} - q}\right)^{1/q}M_{1}^{1-t}M_{2}^{t}\| f \|_{L^{p}(\mu)}.
\end{equation*}

(c-2) $p_{1} = p_{2} < \infty$,$q_{1}$ 或 $q_{2} = \infty$ 的情形.此时不妨设 $p_{1} = p_{2} = p$,$q_{1} < q_{2} = \infty$.由假设 (ii),(iii),
\begin{equation*}
n(\lambda) \leq \left(\frac{M_{1}}{\lambda}\| f \|_{L^{p}(\mu)}\right)^{q_{1}}, \quad \| Tf(y) \|_{L^{\infty}(\nu)} \leq M_{2}\| f \|_{L^{p}(\mu)},
\end{equation*}
故当 $\lambda > M_{2}\| f \|_{L^{p}(\mu)}$ 时 $\{y \in Y;\ |Tf(y)| > \lambda\} = \emptyset$.因此
\begin{align*}
\| Tf \|_{L^{q}(\nu)}^{q} &= \displaystyle\int_{0}^{\infty}q\lambda^{q-1}n(\lambda)\dd \lambda \\
&= \displaystyle\int_{0}^{M_{2}\| f \|_{L^{p}(\mu)}}q\lambda^{q-q_{1}-1}n(\lambda)\dd \lambda\left(M_{1}\| f \|_{L^{p}(\mu)}\right)^{q_{1}} \\
&= q\left(M_{1}\| f \|_{L^{p}(\mu)}\right)^{q_{1}}\dfrac{\left(M_{2}\| f \|_{L^{p}(\mu)}\right)^{q-q_{1}}}{q - q_{1}} = \dfrac{qM_{1}^{q_{1}}M_{2}^{q-q_{1}}}{q - q_{1}}\| f \|_{L^{p}(\mu)}^{q}.
\end{align*}
注意 $1/q = (1-t)/q_{1}$,计算得
\begin{equation*}
\| Tf \|_{L^{q}(\nu)} \leq \left(\frac{q}{q - q_{1}}\right)^{1/q}M_{1}^{1-t}M_{2}^{t}\| f \|_{L^{p}(\mu)}.
\end{equation*}

(c-3) $p_{1} = p_{2} = \infty$,$q_{1} < q_{2} = \infty$ 的情形.此时与 (c-2) 中 $p_{1} = p_{2} = p$,$q_{1} < q_{2} = \infty$ 情形的证明相同.
\end{proof}

% ============ 第1章 §1.9 Hilbert变换的极大算子与主值的逐点收敛(原书 pp.70-78) ============
\section{Hilbert 变换的极大算子与主值的逐点收敛}\label{sec:1.9}
\index{Hilbert 变换!极大算子}
\index{Cotlar 不等式}
\index{Cotlar--Knapp--Stein 概正交引理}


当 $1 \leq p < \infty$ 时,若 $f \in L^{p}(\mathbb{R})$,则对一切 $\varepsilon > 0$,
\begin{equation*}
H_{\varepsilon}f(x) := \frac{1}{\pi}\int_{|y - x| > \varepsilon}\frac{1}{x - y}f(y)\dd y
\end{equation*}
存在.这是因为对每个 $x$,$\frac{1}{x-y}\chi_{|x-y|>\varepsilon}(y) \in L^{p'}(\mathbb{R})$,由 Hölder 不等式即可保证.于是,定义 Hilbert 变换的极大算子 $H_{*}$ 为
\begin{equation*}
H_{*}f(x) = \sup_{\varepsilon > 0}|H_{\varepsilon}f(x)|.
\end{equation*}
于是得到称为 Cotlar 不等式的下述不等式.

\begin{theorem}[Cotlar 不等式]\label{thm:1.9.1}
设 $1 \leq p < \infty$,$0 < \delta \leq 1$.则对任意 $f \in L^{p}(\mathbb{R})$,存在常数 $C_{\delta}$ 与 $C_{2}$ 使得
\begin{equation*}
H_{*}f(x) \leq C_{\delta}\left\{\left(M\left(|Hf|^{\delta}\right)(x)\right)^{1/\delta} + C_{2}Mf(x)\right\} \quad (x \in \mathbb{R}).
\end{equation*}
\end{theorem}

\begin{proof}
设 $f \in L^{p}(\mathbb{R})$.首先考虑 $\delta = 1$ 的情形.令 $x = 0$ 不失一般性,只需证明对任意正数 $\varepsilon$,
\begin{equation}\tag{1.9.1}
|H_{\varepsilon}f(0)| \leq C_{1}M(Hf)(0) + C_{2}Mf(0).
\label{eq:1.9.1}\end{equation}
令 $Q = (-\varepsilon/2, \varepsilon/2)$,$\tilde{Q} = (-\varepsilon, \varepsilon)$,$f_{1} = f\chi_{\tilde{Q}}$,$f_{2} = f - f_{1} = f\chi_{\tilde{Q}^{c}}$.则
\begin{equation*}
Hf_{2}(0) = H_{\varepsilon}f(0),
\end{equation*}
且对 $x \in Q$,若 $|y| \geq \varepsilon$,则 $|y - x| \geq |y| - |x| > \max(|y|/2, \varepsilon/2)$,故
\begin{align*}
\pi|Hf_{2}(0) - Hf_{2}(x)|&= \left|\displaystyle\int_{|y| \geq \varepsilon}\frac{f(y)}{-y}\dd y - \int_{|y| \geq \varepsilon}\frac{f(y)}{x - y}\dd y\right|  \\
&\leq \displaystyle\int_{|y| \geq \varepsilon}\left|\frac{1}{-y} - \frac{1}{x - y}\right||f(y)|\dd y = \int_{|y| \geq \varepsilon}\frac{|x|}{|y||x - y|}|f(y)|\dd y \\
&\leq 2\displaystyle\int_{|y| \geq \varepsilon}\frac{|x|}{|y|^{2}}|f(y)|\dd y \leq \varepsilon\sum_{j = 0}^{\infty}\int_{2^{j}\varepsilon \leq |y| < 2^{j+1}\varepsilon}\frac{|f(y)|}{|y|^{2}}\dd y \\
&\leq \displaystyle\sum_{j = 0}^{\infty}\frac{1}{2^{j-2}}\cdot\frac{1}{2^{j+2}\varepsilon}\int_{|y| < 2^{j+1}\varepsilon}|f(y)|\dd y \\
&\leq \displaystyle\sum_{j = 0}^{\infty}\frac{1}{2^{j-2}}Mf(0) = 8Mf(0).
\end{align*}
因此
\begin{equation}\tag{1.9.2}
|H_{\varepsilon}f(0)| \leq |Hf(x)| + |Hf_{1}(x)| + \frac{8}{\pi}Mf(0) \quad (x \in Q).
\label{eq:1.9.2}\end{equation}
若 $H_{\varepsilon}f(0) = 0$,则 \eqref{eq:1.9.1} 成立,故设 $H_{\varepsilon}f(0) \neq 0$.任取满足 $0 < \lambda < |H_{\varepsilon}f(0)|$ 的 $\lambda$.则对每个 $x \in Q$,$|Hf(x)|$,$|Hf_{1}(x)|$,$\frac{8}{\pi}Mf(0)$ 中必有一个大于 $\lambda/3$.因此

(i) $\lambda < \dfrac{24}{\pi}Mf(0)$

或

(ii) $Q \subset \{x \in Q;\ |Hf(x)| > \lambda/3\} \cup \{x \in Q;\ |Hf_{1}(x)| > \lambda/3\}$

二者之一成立.这里由 Chebyshev 不等式,
\begin{equation*}
|\{x \in Q;\ |Hf(x)| > \lambda/3\}| \leq \frac{3}{\lambda}\int_{Q}|Hf(x)|\dd x \leq \frac{3|Q|}{\lambda}M(Hf)(0),
\end{equation*}
由 Hilbert 变换的弱 $(1,1)$ 性,
\begin{align*}
|\{x \in Q;\ |Hf_{1}(x)| > \lambda/3\}| & \leq \dfrac{3C}{\lambda}\| f_{1} \|_{L^{1}} = \dfrac{3C}{\lambda}\int_{|x| < \varepsilon}|f(x)|\dd x \\
& \leq \dfrac{6C|Q|}{\lambda}Mf(0),
\end{align*}
故在 (ii) 的情形,
\begin{equation*}
\lambda < 3M(Hf)(0) + 6CMf(0).
\end{equation*}
(i),(ii) 合并,
\begin{equation*}
\lambda < 3M(Hf)(0) + \left(6C + \frac{24}{\pi}\right)Mf(0).
\end{equation*}
因 $\lambda$ 是满足 $0 < \lambda < |H_{\varepsilon}f(0)|$ 的任意数,令 $\lambda \to |H_{\varepsilon}f(0)|$,
\begin{equation*}
|H_{\varepsilon}f(0)| \leq 3M(Hf)(0) + \left(6C + \frac{24}{\pi}\right)Mf(0).
\end{equation*}

其次转向 $0 < \delta < 1$ 的情形.一般地,$(a + b)^{\delta} < a^{\delta} + b^{\delta}$($a, b > 0$),故由 \eqref{eq:1.9.2},
\begin{equation*}
|H_{\varepsilon}f(0)|^{\delta} \leq |Hf(x)|^{\delta} + |Hf_{1}(x)|^{\delta} + \left(\frac{8}{\pi}Mf(0)\right)^{\delta} \quad (x \in Q).
\end{equation*}
对两边在 $Q$ 上取积分平均后再取 $1/\delta$ 次幂,
\begin{align*}
|H_{\varepsilon}f(0)|&\leq 3^{1/\delta}\left\{\left(\dfrac{1}{|Q|}\displaystyle\int_{Q}|Hf(x)|^{\delta}\dd x\right)^{1/\delta} + \left(\frac{1}{|Q|}\int_{Q}|Hf_{1}(x)|^{\delta}\dd x\right)^{1/\delta}\right. \\
&\quad + \left.\left(\dfrac{1}{|Q|}\displaystyle\int_{Q}\left(\frac{8}{\pi}Mf(0)\right)^{\delta}\dd x\right)^{1/\delta}\right\}  \\
&\leq 3^{1/\delta}\left\{\left(M\left(|Hf|^{\delta}\right)(0)\right)^{1/\delta} + \left(\frac{1}{|Q|}\int_{Q}|Hf_{1}(x)|^{\delta}\dd x\right)^{1/\delta} + \frac{8}{\pi}Mf(0)\right\}.
\end{align*}
最后一行中间的一项,由 Hilbert 变换的弱 $(1,1)$ 性与 Kolmogorov 不等式,
\begin{align*}
\left(\dfrac{1}{|Q|}\displaystyle\int_{Q}|Hf_{1}(x)|^{\delta}\dd x\right)^{1/\delta} & \leq \left(\dfrac{C_{\delta}'}{|Q|}\left(C_{1}'\| f_{1} \|_{L^{1}}\right)^{\delta}|Q|^{1-\delta}\right)^{1/\delta} \\
& \leq (C_{\delta}')^{1/\delta}C_{1}'\dfrac{1}{|Q|}\int_{\tilde{Q}}|f(x)|\dd x \\
& \leq 2(C_{\delta}')^{1/\delta}C_{1}'Mf(0).
\end{align*}
因此
\begin{equation*}
|H_{\varepsilon}f(0)| \leq C_{\delta}\left\{\left(M\left(|Hf|^{\delta}\right)(0)\right)^{1/\delta} + C_{2}Mf(0)\right\},
\end{equation*}
证明完毕.
\end{proof}

\begin{corollary}\label{cor:1.9.2}
(i) 当 $1 < p < \infty$ 时,
\begin{equation*}
\| H_{*}f \|_{L^{p}} \leq C_{p}\| f \|_{L^{p}} \quad (f \in L^{p}(\mathbb{R})).
\end{equation*}

(ii) $p = 1$ 时,
\begin{equation*}
|\{x \in \mathbb{R};\ H_{*}f(x) > \lambda\}| \leq C_{1}\frac{\| f \|_{L^{1}}}{\lambda} \quad (f \in L^{1}(\mathbb{R}),\ \lambda > 0).
\end{equation*}
\end{corollary}

\begin{proof}
(i) $1 < p < \infty$ 的情形.设 $f \in L^{p}(\mathbb{R})$.由 Cotlar 不等式的 $\delta = 1$ 情形,
\begin{equation*}
H_{*}f(x) \leq CM(Hf)(x) + CMf(x),
\end{equation*}
故由 Minkowski 不等式,极大函数的 $L^{p}$ 不等式(定理~\ref{thm:1.5.5})以及 Hilbert 变换的 $L^{p}$ 不等式(定理~\ref{thm:1.7.2}),
\begin{align*}
\| H_{*}f \|_{L^{p}} & \leq C\| M(Hf) \|_{L^{p}} + C\| Mf \|_{L^{p}} \\
& \leq CC_{p}\| Hf \|_{L^{p}} + CC_{p}\| f \|_{L^{p}} \\
& \leq CC_{p}C_{p}'\| f \|_{L^{p}} + CC_{p}\| f \|_{L^{p}} = C''\| f \|_{L^{p}}.
\end{align*}

(ii) $p = 1$ 的情形.设 $f \in L^{1}(\mathbb{R})$.固定 $0 < \delta < \eta < 1$,使用 Cotlar 不等式(为简单起见令 $C_{\delta}C_{2} = C_{3}$),
\begin{equation}\tag{1.9.3}\begin{aligned}
\{x \in \mathbb{R};\ H_{*}f(x) > \lambda\}&\subset \left\{x \in \mathbb{R};\ C_{\delta}\left(M\left[(Hf)^{\delta}\right](x)\right)^{1/\delta} > \frac{\lambda}{2}\right\} \cup \left\{x \in \mathbb{R};\ C_{3}Mf(x) > \frac{\lambda}{2}\right\}.
 \end{aligned}\label{eq:1.9.3}\end{equation}
现在,对非负值 $g \in L_{\mathrm{loc}}^{1}(\mathbb{R})$,由 $M(g^{\delta}\chi_{\{g \leq \lambda/2\}})(x) \leq (\lambda/2)^{\delta}$,
\begin{align*}
& \{x \in \mathbb{R};\ \left(M(g^{\delta})(x)\right)^{1/\delta} > \lambda\} \\
& \qquad \subset \left\{x \in \mathbb{R};\ M(g^{\delta}\chi_{\{g > \lambda/2\}})(x) + M(g^{\delta}\chi_{\{g \leq \lambda/2\}})(x) > \lambda^{\delta}\right\} \\
& \qquad \subset \left\{x \in \mathbb{R};\ M(g^{\delta}\chi_{\{g > \lambda/2\}})(x) > (1 - 1/2^{\delta})\lambda^{\delta}\right\} \\
& \qquad = \left\{x \in \mathbb{R};\ \left(M(g^{\delta}\chi_{\{g > \lambda/2\}})(x)\right)^{\eta/\delta} > (1 - 1/2^{\delta})^{\eta/\delta}\lambda^{\eta}\right\}
\end{align*}
故注意 $\eta/\delta > 1$,依次使用 Chebyshev 不等式与极大函数的 $L^{p}$ 不等式,
\begin{align*}
|\{x \in \mathbb{R};\ \left(M(g^{\delta})(x)\right)^{1/\delta} > \lambda\}|&\leq \dfrac{1}{(1 - 1/2^{\delta})^{\eta/\delta}\lambda^{\eta}}\displaystyle\int_{\mathbb{R}}\left(M(g^{\delta}\chi_{\{g > \lambda/2\}})(x)\right)^{\eta/\delta}\dd x  \\
&\leq \dfrac{C^{\eta/\delta}}{(1 - 1/2^{\delta})^{\eta/\delta}\lambda^{\eta}}\displaystyle\int_{\mathbb{R}}\left(g^{\delta}\chi_{\{g > \lambda/2\}}\right)^{\eta/\delta}(x)\dd x \\
&\leq \dfrac{C}{\lambda^{\eta}}\displaystyle\int_{\{g > \lambda/2\}}g^{\eta}(x)\dd x.
\end{align*}
因此
\begin{equation}\tag{1.9.4}
|\{x \in \mathbb{R};\ C_{\delta}\left(M\left[(Hf)^{\delta}\right](x)\right)^{1/\delta} > \lambda/2\}|
\label{eq:1.9.4}\end{equation}
\begin{equation}\tag{1.9.5}
\leq \frac{C}{\lambda^{\eta}}\int_{\{|Hf| > \lambda/(4C_{\delta})\}}|Hf(x)|^{\eta}\dd x
\label{eq:1.9.5}\end{equation}
\begin{equation}\tag{1.9.6}
\leq \frac{C}{\lambda^{\eta}}\left|\{|Hf| > \lambda/(4C_{\delta})\}\right|^{1-\eta}(C_{H}\| f \|_{L^{1}})^{\eta}
\label{eq:1.9.6}\end{equation}
\begin{equation}\tag{1.9.7}
\leq \frac{C}{\lambda^{\eta}}\left(4C_{\delta}C_{H}\| f \|_{L^{1}}/\lambda\right)^{1-\eta}(C_{H}\| f \|_{L^{1}})^{\eta}
\label{eq:1.9.7}\end{equation}
\begin{equation}\tag{1.9.8}
\leq C_{\delta,\eta,H}\frac{\| f \|_{L^{1}}}{\lambda}.
\label{eq:1.9.8}\end{equation}
从 \eqref{eq:1.9.5} 到 \eqref{eq:1.9.6},注意 $|\{x \in \mathbb{R};\ |Hf(x)| > \lambda\}| \leq C_{H}\| f \|_{L^{1}}/\lambda$,并应用了 Kolmogorov 不等式(定理~\ref{thm:1.8.5}).

另外,由极大函数的弱 $(1,1)$ 性,
\begin{equation}\tag{1.9.9}
\left|\left\{x \in \mathbb{R};\ C_{3}Mf(x) > \frac{\lambda}{2}\right\}\right| \leq \frac{2C_{3}}{\lambda}C'\| f \|_{L^{1}}.
\label{eq:1.9.9}\end{equation}
合并 \eqref{eq:1.9.3},\eqref{eq:1.9.8},\eqref{eq:1.9.9},
\begin{equation*}
|\{x \in \mathbb{R};\ H_{*}f(x) > \lambda\}| \leq C\frac{\| f \|_{L^{1}}}{\lambda}.
\end{equation*}
\end{proof}

由这一事实,$H_{\varepsilon}f(x)$ 到 $Hf(x)$ 的 $L^{p}$ 收敛性与逐点收敛性可分别由下面两个推论证明.

\begin{corollary}\label{cor:1.9.3}
下列命题成立.

(i) 当 $1 < p < \infty$ 时,

(a) $\| H_{\varepsilon}f \|_{L^{p}} \leq C_{p}\| f \|_{L^{p}}$($f \in L^{p}(\mathbb{R}),\ \varepsilon > 0$).

(b) $\| H_{\varepsilon}f - Hf \|_{L^{p}} \to 0$($\varepsilon \to 0$).

(ii) 当 $p = 1$ 时,对任意正数 $\lambda$,
\begin{equation*}
\lim_{\varepsilon \to 0}|\{x \in \mathbb{R};\ |H_{\varepsilon}f(x) - Hf(x)| > \lambda\}| = 0,
\end{equation*}
即当 $\varepsilon \to 0$ 时,$H_{\varepsilon}f$ 依测度收敛于 $Hf$.
\end{corollary}

\begin{proof}
(i) 的 (a) 由 $|H_{\varepsilon}f(x)| \leq H_{*}f(x)$ 与系~\ref{cor:1.9.2} 显然.

(i) 的 (b).设 $f \in L^{p}(\mathbb{R})$.其次任取 $\eta > 0$.则由 $\mathcal{S}(\mathbb{R})$ 在 $L^{p}(\mathbb{R})$ 中稠密,存在 $f_{0} \in \mathcal{S}(\mathbb{R})$ 满足
\begin{equation*}
\| f - f_{0} \|_{L^{p}} < \eta.
\end{equation*}
又因 $f_{0} \in \mathcal{S}(\mathbb{R})$,由引理~\ref{lem:1.3.2},存在正数 $\varepsilon_{0}$ 使当 $0 < \varepsilon < \varepsilon_{0}$ 时
\begin{equation*}
\| H_{\varepsilon}f_{0} - Hf_{0} \|_{L^{p}} < \eta.
\end{equation*}
于是当 $0 < \varepsilon < \varepsilon_{0}$ 时,也利用 Hilbert 变换的 $L^{p}$ 有界性,
\begin{align*}
\| H_{\varepsilon}f - Hf \|_{L^{p}} & \leq \| H_{\varepsilon}f - H_{\varepsilon}f_{0} \|_{L^{p}} + \| H_{\varepsilon}f_{0} - Hf_{0} \|_{L^{p}} + \| Hf_{0} - Hf \|_{L^{p}} \\
& \leq C_{p}\| f - f_{0} \|_{L^{p}} + \eta + C_{p}\| f - f_{0} \|_{L^{p}} \\
& \leq (2C_{p} + 1)\eta.
\end{align*}
这表明
\begin{equation*}
\lim_{\varepsilon \to 0}\| H_{\varepsilon}f - Hf \|_{L^{p}} = 0.
\end{equation*}

(ii) 设 $f \in L^{1}(\mathbb{R})$,固定任意正数 $\lambda$.其次任取正数 $\eta$.然后与 (i) 的 (b) 一样,取 $f_{0} \in \mathcal{S}(\mathbb{R})$ 使
\begin{equation*}
\| f - f_{0} \|_{L^{1}} < \eta.
\end{equation*}
则由引理~\ref{lem:1.3.2},
\begin{equation}\tag{1.9.10}
|H_{\varepsilon}f_{0}(x)|,\ |Hf_{0}(x)| \leq C_{0}/(1 + |x|), \quad \lim_{\varepsilon \to 0}H_{\varepsilon}f_{0}(x) = Hf_{0}(x),
\label{eq:1.9.10}\end{equation}
故由前者,当 $1 + |x| > 6C_{0}/\lambda$ 时 $|H_{\varepsilon}f_{0}(x) - Hf_{0}(x)| \leq 2C_{0}/(1 + |x|) \leq \lambda/3$,故令 $A = 6C_{0}/\lambda - 1$,
\begin{align*}
& \{x \in \mathbb{R};\ |H_{\varepsilon}f_{0}(x) - Hf_{0}(x)| > \lambda/3\} \\
& \qquad \subset \{x \in \mathbb{R};\ |x| \leq A,\ |H_{\varepsilon}f_{0}(x) - Hf_{0}(x)| > \lambda/3\}.
\end{align*}
因此由 Chebyshev 不等式,
\begin{equation*}
\left|\left\{x \in \mathbb{R};\ \left|H_{\varepsilon}f_{0}(x) - Hf_{0}(x)\right| > \lambda/3\right\}\right| \leq \frac{3}{\lambda}\int_{-A}^{A}\left|H_{\varepsilon}f_{0}(x) - Hf_{0}(x)\right|\dd x.
\end{equation*}
对 \eqref{eq:1.9.10} 的右端使用 Lebesgue 收敛定理,可知存在正数 $\varepsilon_{0}$ 使当 $0 < \varepsilon < \varepsilon_{0}$ 时
\begin{equation*}
\left|\left\{x \in \mathbb{R};\ \left|H_{\varepsilon}f_{0}(x) - Hf_{0}(x)\right| > \lambda/3\right\}\right| < \eta.
\end{equation*}
于是再利用 $H$ 与 $H_{*}$ 的弱 $(1,1)$ 性(定理~\ref{thm:1.8.2},系~\ref{cor:1.9.2}),
\begin{align*}
& |\{x \in \mathbb{R};\ |H_{\varepsilon}f(x) - Hf(x)| > \lambda\}| \\
& \quad \leq |\{x \in \mathbb{R};\ |H_{\varepsilon}f(x) - H_{\varepsilon}f_{0}(x)| > \lambda/3\}| \\
& \quad + |\{x \in \mathbb{R};\ |H_{\varepsilon}f_{0}(x) - Hf_{0}(x)| > \lambda/3\}| \\
& \quad + |\{x \in \mathbb{R};\ |Hf_{0}(x) - Hf(x)| > \lambda/3\}| \\
& \quad \leq \frac{3C}{\lambda}\| f - f_{0} \|_{L^{1}} + \eta + \frac{3C}{\lambda}\| f - f_{0} \|_{L^{1}} \\
& \quad \leq (6C/\lambda + 1)\eta.
\end{align*}
因此
\begin{equation*}
\lim_{\varepsilon \to 0}|\{x \in \mathbb{R};\ |H_{\varepsilon}f(x) - Hf(x)| > \lambda\}| = 0.
\end{equation*}
\end{proof}

\begin{corollary}\label{cor:1.9.4}
设 $1 \leq p < \infty$.若 $f \in L^{p}(\mathbb{R})$,则当 $\varepsilon \to 0$ 时
\begin{equation*}
H_{\varepsilon}f(x) \to Hf(x) \quad (\text{a.e. } x \in \mathbb{R}).
\end{equation*}
\end{corollary}

\begin{proof}
设 $f \in L^{p}(\mathbb{R})$.为证明起见,只需对实部与虚部分别证明即可,故从一开始就假定为实值函数.首先,由系~\ref{cor:1.9.2} 得
\begin{equation}\tag{1.9.11}
|\{x \in \mathbb{R};\ H_{*}f(x) > \lambda\}| \leq C\left(\frac{\| f \|_{L^{p}}}{\lambda}\right)^{p} \quad (\lambda > 0).
\label{eq:1.9.11}\end{equation}
为证明 $\lim_{\varepsilon \to 0}H_{\varepsilon}f(x) = Hf(x)$,考虑下面的函数:
\begin{equation*}
\Delta f(x) := \limsup_{\varepsilon \to 0}H_{\varepsilon}f(x) - \liminf_{\varepsilon \to 0}H_{\varepsilon}f(x).
\end{equation*}
其次任取 $\eta > 0$.则由 $\mathcal{S}(\mathbb{R})$ 在 $L^{p}(\mathbb{R})$ 中稠密,存在 $g \in \mathcal{S}(\mathbb{R})$ 满足
\begin{equation*}
\| f - g \|_{L^{p}} < \eta.
\end{equation*}
对 $g \in \mathcal{S}(\mathbb{R})$,已知 $\lim_{\varepsilon \to 0}H_{\varepsilon}g(x) = Hg(x)$($x \in \mathbb{R}$),故
\begin{align*}
\limsup_{\varepsilon \to 0}(H_{\varepsilon}f(x) - H_{\varepsilon}g(x)) &= \limsup_{\varepsilon \to 0}H_{\varepsilon}f(x) \\
\liminf_{\varepsilon \to 0}(H_{\varepsilon}f(x) - H_{\varepsilon}g(x)) &= \liminf_{\varepsilon \to 0}H_{\varepsilon}f(x).
\end{align*}
故 $\Delta(f - g)(x) = \Delta f(x)$.因此由 $|\Delta(f - g)(x)| \leq 2H_{*}(f - g)(x)$,再利用 \eqref{eq:1.9.11},
\begin{align*}
|\{x \in \mathbb{R};\ |\Delta f(x)| > 1/N\}| & = |\{x \in \mathbb{R};\ |\Delta(f - g)(x)| > 1/N\}| \\
& \leq |\{x \in \mathbb{R};\ 2H_{*}(f - g)(x) > 1/N\}| \\
& \leq C(2N)^{p}\| f - g \|_{L^{p}}^{p} = C(2N)^{p}\eta^{p}.
\end{align*}
因 $\eta > 0$ 任意,令 $\eta \to 0$,
\begin{equation*}
|\{x \in \mathbb{R};\ |\Delta f(x)| > 1/N\}| = 0.
\end{equation*}
于是因 $\{x \in \mathbb{R};\ |\Delta f(x)| > 0\} = \bigcup_{N=1}^{\infty}\{x \in \mathbb{R};\ |\Delta f(x)| > 1/N\}$,
\begin{equation*}
|\{x \in \mathbb{R};\ |\Delta f(x)| > 0\}| \leq \sum_{N=1}^{\infty}\left|\{x \in \mathbb{R};\ |\Delta f(x)| > 1/N\}\right| = 0.
\end{equation*}
即对几乎所有的 $x \in \mathbb{R}$,$\lim_{\varepsilon \to 0}H_{\varepsilon}f(x)$ 存在.另一方面,由系~\ref{cor:1.9.3},这个极限函数是 $Hf(x)$,故
\begin{equation*}
\lim_{\varepsilon \to 0}H_{\varepsilon}f(x) = Hf(x) \quad (\text{a.e. } x \in \mathbb{R}).
\end{equation*}
\end{proof}

% ============ 第1章 §1.10 Hilbert变换的L²有界性(再访)(原书 pp.78-88) ============
\section{\texorpdfstring{Hilbert 变换的 $L^2$ 有界性(再访)}{Hilbert 变换的 L2 有界性(再访)}}\label{sec:1.10}
\index{Hilbert 变换!L2 有界性}


Hilbert 变换的 $L^{2}$ 有界性是用 Fourier 变换证明的,但在满足与 Hilbert 变换类似的积分核条件但无齐次性的情形下,有时不能使用 Fourier 变换.这种情形下,作为证明 $L^{2}$ 有界性的有效方法,人们开发了称为 Cotlar--Knapp--Stein 概正交性引理的方法.本节介绍它,并举例说明用它可证明 Hilbert 变换的 $L^{2}$ 有界性.

\begin{lemma}[Cotlar--Knapp--Stein 概正交性引理]\label{lem:1.10.1}
设 $\{T_{j}\}_{j \in \mathbb{Z}}$ 是 Hilbert 空间 $H$ 上的有界线性算子列,设存在数列 $\{a(j)\}_{j \in \mathbb{Z}}$ 满足
\begin{equation}\tag{1.10.1}
\| T_{j}^{*}T_{k} \|_{L(H)} + \| T_{j}T_{k}^{*} \|_{L(H)} \leq a(j - k) \quad (j, k \in \mathbb{Z})
\label{eq:1.10.1}\end{equation}
与
\begin{equation}\tag{1.10.2}
A := \sum_{j \in \mathbb{Z}}\sqrt{a(j)} < \infty.
\label{eq:1.10.2}\end{equation}
则下面结论成立:

(i) 对 $\mathbb{Z}$ 的任意有限子集 $\Lambda$,
\begin{equation*}
\left\| \sum_{j \in \Lambda}T_{j} \right\|_{L(H)} \leq A.
\end{equation*}

(ii) 对所有 $x \in H$,
\begin{equation*}
\sum_{j \in \mathbb{Z}}\| T_{j}(x) \|_{H}^{2} \leq A^{2}\| x \|_{H}^{2}.
\end{equation*}

(iii) 对每个 $x \in H$,当 $N \to \infty$ 时 $\sum_{|j| \leq N}T_{j}(x)$ 依 $H$ 范数收敛到某个 $T(x) \in H$,如此定义的 $T$ 是 $H$ 上的有界线性算子,且
\begin{equation*}
\| T(x) \|_{H} \leq A\| x \|_{H} \quad (x \in H).
\end{equation*}

其中 $T_{j}^{*}$ 是 $T_{j}$ 的伴随算子,$\| T_{j} \|_{L(H)}$ 表示 $T_{j}$ 的算子范数.
\end{lemma}

\begin{proof}
(i) 一般地,对 $H$ 上的有界线性算子 $S$,$\| S \|_{L(H)}^{2} = \| SS^{*} \|_{L(H)}$,故 $\| S \|_{L(H)}^{2} = \| [SS^{*}]^{2^{\ell}} \|_{L(H)}^{2^{-\ell}}$($\ell \in \mathbb{N}$),于是
\begin{equation*}
B^{2} := \left\| \sum_{j \in \Lambda}T_{j} \right\|_{L(H)}^{2} = \left\| \left[\left(\sum_{j \in \Lambda}T_{j}\right)\left(\sum_{j \in \Lambda}T_{j}^{*}\right)\right]^{2^{\ell}} \right\|_{L(H)}^{2^{-\ell}}.
\end{equation*}
因此令 $m = 2^{\ell}$,展开右边范数记号内的部分,
\begin{equation}\tag{1.10.3}\begin{aligned}
B^{2} &= \left\| \displaystyle\sum_{j_{1}, j_{2}, \dots, j_{2m} \in \Lambda}T_{j_{1}}T_{j_{2}}^{*}\cdots T_{j_{2m-1}}T_{j_{2m}}^{*} \right\|_{L(H)}^{1/m} \\
&\leq \left(\displaystyle\sum_{j_{1}, j_{2}, \dots, j_{2m} \in \Lambda}\left\| T_{j_{1}}T_{j_{2}}^{*}\cdots T_{j_{2m-1}}T_{j_{2m}}^{*} \right\|_{L(H)}\right)^{1/m}.
\end{aligned}\label{eq:1.10.3}\end{equation}
由 \eqref{eq:1.10.1},
\begin{align*}
\left\| T_{j_{1}}T_{j_{2}}^{*}\cdots T_{j_{2m-1}}T_{j_{2m}}^{*} \right\|_{L(H)} & \leq \| T_{j_{1}}T_{j_{2}}^{*} \|_{L(H)}\cdots\| T_{j_{2m-1}}T_{j_{2m}}^{*} \|_{L(H)} \\
& \leq a(j_{1} - j_{2})a(j_{3} - j_{4})\cdots a(j_{2m-1} - j_{2m}),
\end{align*}
且
\begin{align*}
& \left\| T_{j_{1}}T_{j_{2}}^{*}\cdots T_{j_{2m-1}}T_{j_{2m}}^{*} \right\|_{L(H)} \\
& \quad \leq \| T_{j_{1}} \|_{L(H)}\| T_{j_{2}}^{*}T_{j_{3}} \|_{L(H)}\cdots\| T_{j_{2m-2}}^{*}T_{j_{2m-1}} \|_{L(H)}\| T_{j_{2m}}^{*} \|_{L(H)} \\
& \quad \leq \sqrt{a(0)}a(j_{2} - j_{3})a(j_{4} - j_{5})\cdots a(j_{2m-2} - j_{2m-1})\sqrt{a(0)},
\end{align*}
故取两者的几何平均,
\begin{align*}
\left\| T_{j_{1}}T_{j_{2}}^{*}\cdots T_{j_{2m-1}}T_{j_{2m}}^{*} \right\|_{L(H)}&\leq \sqrt{a(0)}\sqrt{a(j_{1} - j_{2})}\sqrt{a(j_{2} - j_{3})}\cdots\sqrt{a(j_{2m-1} - j_{2m})}.
 \end{align*}
于是由 \eqref{eq:1.10.2},
\begin{align*}
& \displaystyle\sum_{j_{1}, j_{2}, \ldots, j_{2m} \in \Lambda}\left\| T_{j_{1}}T_{j_{2}}^{*}\cdots T_{j_{2m-1}}T_{j_{2m}}^{*} \right\|_{L(H)} \\
& \leq \sqrt{a(0)}\displaystyle\sum_{j_{2}, \ldots, j_{2m} \in \Lambda}\left(\sum_{j_{1} \in \Lambda}\sqrt{a(j_{1} - j_{2})}\right)\sqrt{a(j_{2} - j_{3})}\cdots\sqrt{a(j_{2m-1} - j_{2m})} \\
& \ldots \\
& \leq \sqrt{a(0)}A^{2m-1}\sum_{j_{2m} \in \Lambda}1 = \sqrt{a(0)}A^{2m-1}(\Lambda \text{ 的元素个数}).
\end{align*}
与 \eqref{eq:1.10.3} 合并,
\begin{equation*}
B^{2} \leq a(0)^{1/(2m)}A^{2-1/m}(\Lambda \text{ 的元素个数})^{1/m}.
\end{equation*}
令 $\ell \to \infty$,则 $m = 2^{\ell} \to \infty$,故
\begin{equation*}
\left\| \sum_{j \in \Lambda}T_{j} \right\|_{L(H)} = B \leq A.
\end{equation*}

(ii) 设 $x \in H$.对实数 $\theta$,算子列 $\{e^{2\pi ij\theta}T_{j}\}_{j \in \mathbb{Z}}$ 也满足 \eqref{eq:1.10.1},\eqref{eq:1.10.2},故由 (i),
\begin{equation}\tag{1.10.4}
\left\| \sum_{|j| \leq N}e^{2\pi ij\theta}T_{j}(x) \right\|_{H} \leq A\| x \|_{H}.
\label{eq:1.10.4}\end{equation}
另一方面,
\begin{align*}
\displaystyle\int_{0}^{1}\left\| \sum_{|j| \leq N}e^{2\pi ij\theta}T_{j}(x) \right\|_{H}^{2}\dd\theta&= \displaystyle\int_{0}^{1}\left(\sum_{|j| \leq N}e^{2\pi ij\theta}T_{j}(x), \sum_{|k| \leq N}e^{2\pi ik\theta}T_{k}(x)\right)\dd\theta  \\
&= \displaystyle\sum_{|j| \leq N}\sum_{|k| \leq N}\int_{0}^{1}e^{2\pi i(j-k)\theta}\dd\theta\left(T_{j}(x), T_{k}(x)\right) \\
&= \displaystyle\sum_{|j| \leq N}\| T_{j}(x) \|_{H}^{2}.
\end{align*}
与 \eqref{eq:1.10.4} 合并,
\begin{equation*}
\sum_{|j| \leq N}\| T_{j}(x) \|_{H}^{2} \leq A^{2}\| x \|_{H}^{2}. \quad \text{令 } N \to \infty \text{ 得 } \sum_{j \in \mathbb{Z}}\| T_{j}(x) \|_{H}^{2} \leq A^{2}\| x \|_{H}^{2}.
\end{equation*}

(iii) 设 $x \in H$.首先证明 $\left\{\sum_{|j| \leq N}T_{j}(x)\right\}_{N=1}^{\infty}$ 是 $H$ 中的 Cauchy 列.若否定这一点,则存在正数 $\varepsilon$ 与自然数列 $N_{1} < N_{2} < \cdots$ 使
\begin{equation}\tag{1.10.5}
\left\| \sum_{N_{k} < |j| \leq N_{k+1}}T_{j}(x) \right\|_{H} \geq \varepsilon \quad (k = 1, 2, \dots).
\label{eq:1.10.5}\end{equation}
另一方面,令
\begin{equation*}
S_{j} = e^{2\pi ik\theta}T_{j} \quad (N_{k} < |j| \leq N_{k+1},\ k = 1, 2, \dots),
\end{equation*}
容易看出 $\{S_{j}\}_{j \in \mathbb{Z}}$ 也满足 \eqref{eq:1.10.1},\eqref{eq:1.10.2},故由 (i),
\begin{equation}\tag{1.10.6}
\left\| \sum_{k=1}^{K}e^{2\pi ik\theta}\sum_{N_{k} < |j| \leq N_{k+1}}T_{j}(x) \right\|_{H} \leq A\| x \|_{H}.
\label{eq:1.10.6}\end{equation}
与 (ii) 中的计算一样,
\begin{equation}\tag{1.10.7}
\int_{0}^{1}\left\| \sum_{k=1}^{K}e^{2\pi ik\theta}\sum_{N_{k} < |j| \leq N_{k+1}}T_{j}(x) \right\|_{H}^{2}\dd\theta = \sum_{k=1}^{K}\left\| \sum_{N_{k} < |j| \leq N_{k+1}}T_{j}(x) \right\|_{H}^{2}.
\label{eq:1.10.7}\end{equation}
由 \eqref{eq:1.10.5},\eqref{eq:1.10.6},\eqref{eq:1.10.7},
\begin{equation*}
K\varepsilon^{2} \leq \sum_{k=1}^{K}\left\| \sum_{N_{k} < |j| \leq N_{k+1}}T_{j}(x) \right\|_{H}^{2} \leq A^{2}\| x \|_{H}^{2},
\end{equation*}
令 $K \to \infty$ 即得矛盾.因此 $\left\{\sum_{|j| \leq N}T_{j}(x)\right\}_{N=1}^{\infty}$ 是 $H$ 中的 Cauchy 列.由 Hilbert 空间的完备性,当 $N \to \infty$ 时 $\sum_{|j| \leq N}T_{j}(x)$ 依 $H$ 范数收敛到某个 $T(x) \in H$.另外,对 $\alpha, \beta \in \mathbb{C}$,$x, y \in H$,
\begin{align*}
\left\| \displaystyle\sum_{|j| \leq N}T_{j}(\alpha x + \beta y) - \alpha T(x) - \beta T(y) \right\|_{H}&= \left\| \alpha\left(\displaystyle\sum_{|j| \leq N}T_{j}(x) - T(x)\right) + \beta\left(\sum_{|j| \leq N}T_{j}(y) - T(y)\right) \right\|_{H}  \\
&\leq |\alpha|\left\| \displaystyle\sum_{|j| \leq N}T_{j}(x) - T(x) \right\|_{H} + |\beta|\left\| \sum_{|j| \leq N}T_{j}(y) - T(y) \right\|_{H},
\end{align*}
故当 $N \to \infty$ 时 $\sum_{|j| \leq N}T_{j}(\alpha x + \beta y)$ 收敛到 $\alpha T(x) + \beta T(y)$.另一方面,$\sum_{|j| \leq N}T_{j}(\alpha x + \beta y)$ 本来收敛到 $T(\alpha x + \beta y)$,故 $T(\alpha x + \beta y) = \alpha T(x) + \beta T(y)$,即 $T$ 是从 $H$ 到 $H$ 的线性算子.进而,
\begin{align*}
\| T(x) \|_{H} & \leq \left\| T(x) - \displaystyle\sum_{|j| \leq N}T_{j}(x) \right\|_{H} + \left\| \sum_{|j| \leq N}T_{j}(x) \right\|_{H} \\
& \leq \left\| T(x) - \displaystyle\sum_{|j| \leq N}T_{j}(x) \right\|_{H} + A\| x \|_{H},
\end{align*}
故令 $N \to \infty$ 得
\begin{equation*}
\| T(x) \|_{H} \leq A\| x \|_{H}.
\end{equation*}
因此 $T$ 是 $H$ 上的有界线性算子.
\end{proof}

利用这个引理,可证明如下奇异积分的 $L^{2}$ 有界性.

\begin{theorem}\label{thm:1.10.2}
对二元函数 $K(x, y)$,设存在满足以下 4 个条件的正数 $C_{1}, C_{2}, C_{3}$ 与 $0 < \delta \leq 1$:

(i)
\begin{equation*}
|K(x, y)| \leq \frac{C_{1}}{|x - y|} \quad (x \neq y),
\end{equation*}

(ii)
\begin{equation*}
|K(x, y) - K(x, z)| \leq \frac{C_{2}}{|x - y|}\left(\frac{|y - z|}{|x - y|}\right)^{\delta} \quad (|y - z| < |x - y|/2),
\end{equation*}
\begin{equation*}
|K(x, y) - K(z, y)| \leq \frac{C_{3}}{|x - y|}\left(\frac{|x - z|}{|x - y|}\right)^{\delta} \quad (|x - z| < |x - y|/2),
\end{equation*}

(iii)
\begin{equation*}
\lim_{\varepsilon \to 0}\int_{\varepsilon < |x - y| < 2^{m}}K(x, y)\dd y = 0 \quad (m \in \mathbb{Z},\ x \in \mathbb{R}),
\end{equation*}

(iv)
\begin{equation*}
\int_{2^{m} \leq |x - y| < 2^{m+1}}K(x, y)\dd x = 0 \quad (m \in \mathbb{Z},\ y \in \mathbb{R}).
\end{equation*}

则当 $f \in C^{1}(\mathbb{R})$ 且 $|f(x)| \leq C_{4}/(1 + |x|)^{2}$,$|f'(x)| \leq C_{5}/(1 + |x|)$ 时,
\begin{equation*}
Tf(x) := \mathrm{p.v.}\int K(x, y)f(y)\dd y = \lim_{\varepsilon \to 0}\int_{|x - y| > \varepsilon}K(x, y)f(y)\dd y
\end{equation*}
对每个 $x \in \mathbb{R}$ 存在,并对某个常数 $C_{6}$ 有
\begin{equation*}
|Tf(x)| \leq \frac{C_{6}}{1 + |x|},
\end{equation*}
且存在不依赖于 $f$ 的常数 $C$ 使
\begin{equation*}
\| Tf \|_{L^{2}} \leq C\| f \|_{L^{2}} \quad (f \in \mathcal{S}(\mathbb{R})).
\end{equation*}
\end{theorem}

\begin{proof}
注意由条件 (iii) 得
\begin{equation}\tag{1.10.8}
\int_{2^{m} \leq |x - y| < 2^{m+1}}K(x, y)\dd y = 0 \quad (m \in \mathbb{Z},\ x \in \mathbb{R}).
\label{eq:1.10.8}\end{equation}
现在,设 $f \in C^{1}(\mathbb{R})$ 且 $|f(x)| \leq C_{4}/(1 + |x|)^{2}$,$|f'(x)| \leq C_{5}/(1 + |x|)$,则
\begin{align*}
T_{\varepsilon}f(x)&:= \displaystyle\int_{|x - y| > \varepsilon}K(x, y)f(y)\dd y  \\
&= f(x)\displaystyle\int_{\varepsilon < |x - y| < 1}K(x, y)\dd y + \int_{\varepsilon < |x - y| < 1}K(x, y)\left(f(y) - f(x)\right)\dd y \\
&\quad + \displaystyle\int_{|x - y| \geq 1}K(x, y)f(y)\dd y.
\end{align*}
当 $0 < |x - y| \leq 1$ 时,
\begin{equation*}
|K(x, y)\left(f(y) - f(x)\right)| \leq C_{1}\max_{|x - z| \leq 1}|f'(z)| \leq \max_{|x - z| \leq 1}\frac{C_{1}C_{5}}{1 + |z|} \leq \frac{2C_{1}C_{5}}{1 + |x|}.
\end{equation*}
当 $|x - y| > 1$ 时,
\begin{equation*}
|K(x, y)f(y)| \leq \frac{C_{1}C_{4}}{|x - y|(1 + |y|)^{2}} \leq \frac{C_{1}C_{4}}{(1 + |y|)^{2}}.
\end{equation*}
故 $K(x, y)\left(f(y) - f(x)\right)$ 在 $[x - 1, x + 1]$ 上可积,$K(x, y)f(y)$ 在 $(-\infty, x - 1)$,$(x + 1, \infty)$ 上可积,故利用条件 (iii),
\begin{align*}
\lim_{\varepsilon \to 0}T_{\varepsilon}f(x) &= \displaystyle\int_{|x - y| < 1}K(x, y)\left(f(y) - f(x)\right)\dd y + \int_{|x - y| \geq 1}K(x, y)f(y)\dd y \\
&=: I_{1}(x) + I_{2}(x).
\end{align*}
以下与 Hilbert 变换的引理~\ref{lem:1.3.2} 一样,可证明存在常数 $C_{6}$ 使
\begin{equation*}
I_{1}(x) + I_{2}(x) \leq \frac{C_{6}}{1 + |x|} \quad (x \in \mathbb{R}).
\end{equation*}

这里,对每个 $j \in \mathbb{Z}$ 令
\begin{equation*}
K_{j}(x, y) := K(x, y)\chi_{\{2^{j} \leq |x - y| < 2^{j+1}\}}(x, y).
\end{equation*}
对每个 $j \in \mathbb{Z}$,$K_{j}(x, y)$ 作为 $y$ 的函数属于 $L^{2}(\mathbb{R})$,故当 $f \in L^{2}(\mathbb{R})$ 时可定义
\begin{equation}\tag{1.10.9}
T_{j}f(x) := \int_{\mathbb{R}}K_{j}(x, y)f(y)\dd y.
\label{eq:1.10.9}\end{equation}
对这个 $T_{j}$,
\begin{equation}\tag{1.10.10}
\int_{\mathbb{R}}|K_{j}(x, y)|\dd y,\ \int_{\mathbb{R}}|K_{j}(x, y)|\dd x \leq \int_{2^{j} \leq |x - y| < 2^{j+1}}\frac{C_{1}}{|x - y|}\dd y \leq 2C_{1}\log 2,
\label{eq:1.10.10}\end{equation}
故使用 Cauchy--Schwarz 不等式与 Fubini 定理,
\begin{align*}
\| T_{j}f \|_{L^{2}}&\leq \left(\displaystyle\int_{\mathbb{R}}\left[\left(\int_{\mathbb{R}}|K_{j}(x, y)|\dd y\right)^{1/2}\left(\int_{\mathbb{R}}|K_{j}(x, y)||f(y)|^{2}\dd y\right)^{1/2}\right]^{2}\dd x\right)^{1/2}  \\
&\leq \sqrt{2C_{1}\log 2}\left(\displaystyle\int_{\mathbb{R}}\int_{\mathbb{R}}|K_{j}(x, y)||f(y)|^{2}\dd y\dd x\right)^{1/2} \\
&\leq 2C_{1}\log 2\| f \|_{L^{2}}.
\end{align*}
因此
\begin{equation}\tag{1.10.11}
\| T_{j}T_{k}^{*} \|_{L(H)} + \| T_{j}^{*}T_{k} \|_{L(H)} \leq 8C_{1}^{2}\log^{2}2 \quad (j, k \in \mathbb{Z}).
\label{eq:1.10.11}\end{equation}
其次,对 $j, k \in \mathbb{Z}$ 令
\begin{equation*}
S_{j,k}(x, y) := \int_{\mathbb{R}}K_{j}(x, z)\overline{K_{k}(y, z)}\dd z,
\end{equation*}
则对 $f \in \mathcal{S}(\mathbb{R})$,
\begin{equation*}
T_{j}T_{k}^{*}f(x) = \int_{\mathbb{R}}K_{j}(x, z)\left(\int_{\mathbb{R}}\overline{K_{k}(y, z)}f(y)\dd y\right)\dd z = \int_{\mathbb{R}}S_{j,k}(x, y)f(y)\dd y.
\end{equation*}
于是来考察 $S_{j,k}(x, y)$.

(i) $k \leq j - 3$ 的情形.

(a) 当 $|x - y| \geq 2^{j+1} + 2^{k+1}$ 时,若 $2^{j} \leq |x - z| < 2^{j+1}$ 且 $2^{k} \leq |y - z| < 2^{k+1}$,则 $|x - y| \leq |x - z| + |y - z| < 2^{j+1} + 2^{k+1}$,故 $K_{j}(x, z)\overline{K_{k}(y, z)} = 0$,
\begin{equation*}
S_{j,k}(x, y) = 0.
\end{equation*}

(b) 当 $2^{j+1} - 2^{k+1} \leq |x - y| < 2^{j+1} + 2^{k+1}$ 时,若 $2^{k} \leq |y - z| < 2^{k+1}$,则 $|x - z| \geq |x - y| - |y - z| > 2^{j+1} - 2^{k+1} - 2^{k+1} > 2^{j}$,故
\begin{equation*}
|S_{j,k}(x, y)| \leq \int_{2^{k} \leq |y - z| < 2^{k+1}}\frac{C_{1}}{|x - z|}\frac{C_{1}}{|y - z|}\dd z \leq 2C_{1}^{2}2^{-j}\log 2.
\end{equation*}

(c) 当 $2^{j} + 2^{k+1} \leq |x - y| < 2^{j+1} - 2^{k+1}$ 时,若 $2^{k} \leq |y - z| < 2^{k+1}$,则 $|y - z| < |x - y|/2$,且 $|x - z| \geq |x - y| - |y - z| \geq 2^{j} + 2^{k+1} - 2^{k+1} = 2^{j}$,$|x - z| \leq |x - y| + |y - z| \leq 2^{j+1} - 2^{k+1} + 2^{k+1} = 2^{j+1}$,故 $K_{j}(x, z) = K(x, z)$.因此由 \eqref{eq:1.10.8},
\begin{equation*}
S_{j,k}(x, y) = \int_{\mathbb{R}}\left(K(x, z) - K(x, y)\right)\overline{K_{k}(y, z)}\dd z.
\end{equation*}
又已知 $|y - z| < |x - y|/2$,故由条件 (ii),
\begin{align*}
|S_{j,k}(x, y)| &\leq \displaystyle\int_{2^{k} \leq |y - z| < 2^{k+1}}\frac{C_{2}}{|x - y|}\left(\frac{|y - z|}{|x - y|}\right)^{\delta}\frac{C_{1}}{|y - z|}\dd z \\
&\leq 2C_{1}C_{2}2^{-j + \delta(k + 1 - j)}\log 2.
\end{align*}

(d) 当 $2^{j} - 2^{k+1} \leq |x - y| < 2^{j} + 2^{k+1}$ 时,若 $2^{k} \leq |y - z| < 2^{k+1}$,则 $|x - z| \geq |x - y| - |y - z| > 2^{j} - 2^{k+1} - 2^{k+1} \geq 2^{j} - 2^{j-1} = 2^{j-1}$.因此与 (b) 的情形一样,
\begin{equation*}
|S_{j,k}(x, y)| \leq 4C_{1}^{2}2^{-j}\log 2.
\end{equation*}

(e) 当 $|x - y| \leq 2^{j} - 2^{k+1}$ 时,若 $2^{j} \leq |x - z| < 2^{j+1}$ 且 $2^{k} \leq |y - z| < 2^{k+1}$,则 $|x - y| \geq |x - z| - |y - z| > 2^{j} - 2^{k+1}$,故 $K_{j}(x, z)\overline{K_{k}(y, z)} = 0$,
\begin{equation*}
S_{j,k}(x, y) = 0.
\end{equation*}

由 (a),(b),(c),(d),(e),
\begin{align*}
& \displaystyle\int_{\mathbb{R}}|S_{j,k}(x, y)|\dd y \\
& \leq 2(2C_{1}^{2}2^{-j}\log 2)2^{k+2} + 2(2C_{1}C_{2}2^{-j+\delta(k+1-j)}\log 2)(2^{j+1} - 2^{j} + 2^{k+2}) \\
& \quad + 2(4C_{1}^{2}2^{-j}\log 2)2^{k+2} \\
& \leq (48C_{1}^{2}\log 2)2^{k-j} + (8C_{1}C_{2}\log 2)2^{\delta}2^{\delta(k-j)} \\
& \leq 8C_{1}(6C_{1} + C_{2})(\log 2)2^{\delta}2^{-\delta|j-k|}.
\end{align*}
同样,
\begin{equation*}
\int_{\mathbb{R}}|S_{j,k}(x, y)|\dd x \leq 8C_{1}(6C_{1} + C_{2})(\log 2)2^{\delta}2^{-\delta|j-k|}.
\end{equation*}
于是与证明 $T_{j}$ 的 $L^{2}$ 有界性时一样,
\begin{equation}\tag{1.10.12}
\| T_{j}T_{k}^{*} \|_{L(H)} \leq 8C_{1}(6C_{1} + C_{2})(\log 2)2^{\delta}2^{-\delta|j-k|}.
\label{eq:1.10.12}\end{equation}

(ii) $k \geq j + 3$ 的情形.在 (i) 中把 $j$ 与 $k$ 的角色互换,可知上述估计 \eqref{eq:1.10.12} 成立.

由 (i),(ii) 的证明可见,对 $T_{j}^{*}T_{k}$ 也有同样的估计.

因此,令
\begin{equation*}
a(j) = \begin{cases} 16C_{1}^{2}\log^{2}2 & (j = -2, -1, 0, 1, 2) \\ 16C_{1}(6C_{1} + C_{2})(\log 2)2^{\delta}2^{-\delta|j|} & (|j| \geq 3), \end{cases}
\end{equation*}
则
\begin{equation*}
\| T_{j}T_{k}^{*} \|_{L(H)} + \| T_{j}^{*}T_{k} \|_{L(H)} \leq a(j - k) \quad (j, k \in \mathbb{Z}),
\end{equation*}
且
\begin{equation*}
A := \sum_{j \in \mathbb{Z}}\sqrt{a(j)} < \infty.
\end{equation*}
另一方面,前段中已证对 $f \in \mathcal{S}(\mathbb{R})$,
\begin{equation*}
\lim_{N \to \infty}\sum_{j = -N}^{N}T_{j}f(x) = Tf(x),
\end{equation*}
故由 Cotlar--Knapp--Stein 引理,
\begin{equation*}
\| Tf \|_{L^{2}} \leq A\| f \|_{L^{2}}.
\end{equation*}
\end{proof}

Hilbert 变换的积分核满足上面定理的全部假设,故重新作为

\begin{corollary}\label{cor:1.10.3}
对 Hilbert 变换 $H$,
\begin{equation*}
\| Hf \|_{L^{2}} \leq C\| f \|_{L^{2}} \quad (f \in \mathcal{S}(\mathbb{R})).
\end{equation*}
\end{corollary}

% ============ 第1章 §1.11 加权范数不等式(原书 pp.88-108) ============
\section{加权范数不等式}\label{sec:1.11}
\index{权}
\index{Ap 类}
\index{加权 Lp 估计!Hardy--Littlewood 极大函数的}
\index{加权 Lp 估计!Hilbert 变换的}
\index{逆 Hölder 不等式}
\index{加权中心极大函数}


在微分方程的边值问题等中,$|x|^{a}$ 形式的权重有时起着重要作用.对这种形式的权重 $w(x)$,Hilbert 变换有时是 $L^{p}(w(x)\dd x)$ 有界的;更一般地,作为使 Hilbert 变换成为 $L^{p}(w(x)\dd x)$ 有界的权重族(类),有 B. Muckenhoupt 发现的 $A_{p}$ 类.本节介绍这个 $A_{p}$ 类.

到前一节为止 Hardy--Littlewood 极大函数也起着重要作用,本节就在与此的关联中展开讨论.本节中也将把函数 $f$ 的 Hardy--Littlewood 极大函数 $M_{1}f$ 简记为 $Mf$.

首先,$\mathbb{R}$ 上的权重是指 $\mathbb{R}$ 上恒不为零,局部可积且非负值的 Lebesgue 可测函数.从下面的定理开始.

\begin{theorem}\label{thm:1.11.1}
设 $1 \leq p < \infty$.设 $\mathbb{R}$ 上的正则正测度 $\mu$ 存在正数 $C_{0}$ 使得
\begin{equation}\tag{1.11.1}
\mu(\{x \in \mathbb{R};\ Mf(x) > \lambda\})^{1/p} \leq \frac{C_{0}}{\lambda}\left(\int_{\mathbb{R}}|f(x)|^{p}\dd\mu(x)\right)^{1/p} \quad (\lambda > 0)
\label{eq:1.11.1}\end{equation}
对一切 $f \in L^{p}(\dd\mu)$ 成立(Hardy--Littlewood 极大函数的弱 $L^{p}(\dd\mu)$ 性).则

(i) $\mu$ 关于 Lebesgue 测度绝对连续(即 $|E| = 0$ 则 $\mu(E) = 0$),存在权重 $w(x)$ 使 $\dd\mu(x) = w(x)\dd x$.且 $w(x) > 0$(a.e. $x \in \mathbb{R}$).

(ii) 对一切开区间 $Q$ 与一切 $f \in L^{p}(\dd\mu)$,
\begin{equation}\tag{1.11.2}
\frac{1}{|Q|}\int_{Q}|f(x)|\dd x \leq C_{0}\left(\frac{1}{\mu(Q)}\int_{Q}|f(x)|^{p}\dd\mu(x)\right)^{1/p}.
\label{eq:1.11.2}\end{equation}

(iii) 对 (i) 中确定的权重 $w$,当 $1 < p < \infty$ 时,
\begin{equation}\tag{1.11.3}
\left(\frac{1}{|Q|}\int_{Q}w(x)\dd x\right)\left(\frac{1}{|Q|}\int_{Q}w(x)^{-1/(p-1)}\dd x\right)^{p-1} \leq C_{0}^{p},
\label{eq:1.11.3}\end{equation}
当 $p = 1$ 时,
\begin{equation}\tag{1.11.4}
\frac{1}{|Q|}\int_{Q}w(x)\dd x \leq C_{0}\operatorname{ess}\inf_{x \in Q}w(x)
\label{eq:1.11.4}\end{equation}
对任意开区间 $Q$ 成立.

(iv) 对每个开区间 $Q$ 与 $Q$ 内的一切 Lebesgue 可测集 $E$,
\begin{equation}\tag{1.11.5}
\frac{|E|}{|Q|} \leq C_{0}\left(\frac{\mu(E)}{\mu(Q)}\right)^{1/p}.
\label{eq:1.11.5}\end{equation}
\end{theorem}

\begin{proof}
(i) 设 $|E| = 0$.为证明 $\mu(E) = 0$,因 $\mu$ 是正则测度,不妨设 $E$ 为紧集.对正数 $\varepsilon$,由 $\mu$ 的正则性,存在开集 $O$ 使 $E \subset O$ 且 $\mu(O \setminus E) = \mu(O) - \mu(E) < \varepsilon$.令 $f(x) = \chi_{O\setminus E}(x)$,则 $f \in L^{p}(\dd\mu)$ 且 $\| f \|_{L^{p}(\dd\mu)} = \left(\int_{\mathbb{R}}\left(\chi_{O\setminus E}(y)\right)^{p}\dd\mu(y)\right)^{1/p} = \mu(O \setminus E)^{1/p} < \varepsilon^{1/p}$.其次,若 $x \in E$,则 $x \in E \subset O$,因 $O$ 是开集,存在开区间 $Q$ 使 $x \in Q \subset O$.因此,因 $|E| = 0$,$|Q|^{-1}\int_{Q}|f(y)|\dd y = |(O \setminus E) \cap Q|/|Q| = |Q \setminus (E \cap Q)|/|Q| = |Q|/|Q| = 1$,故 $Mf(x) = 1$.因此由 \eqref{eq:1.11.1},
\begin{equation*}
\mu(E) \leq \mu(\{x \in \mathbb{R};\ Mf(x) > 1/2\}) \leq 2^{p}C_{0}^{p}\| f \|_{L^{p}(\dd\mu)}^{p} < 2^{p}C_{0}^{p}\varepsilon.
\end{equation*}
因正数 $\varepsilon$ 任意,令 $\varepsilon \to 0$ 得 $\mu(E) = 0$.由此可知 $\mu$ 关于 Lebesgue 测度绝对连续,故存在非负值可测函数 $w(x) \in L_{\mathrm{loc}}^{1}(\mathbb{R})$ 使 $\dd\mu(x) = w(x)\dd x$.为证明 $w(x) > 0$(a.e. $x \in \mathbb{R}$),只需证明对任何具有正测度的可测集 $E$ 有 $\mu(E) > 0$.对使 $E \subset (-b, b)$ 的 $b > 0$,
\begin{equation*}
\frac{1}{2b}\int_{-b}^{b}\chi_{E}(y)\dd y = \frac{|E|}{2b} \leq M\chi_{E}(x) \quad (-b < x < b),
\end{equation*}
故
\begin{equation}\tag{1.11.6}\begin{aligned}
\mu((-b, b)) & \leq \mu(\{x \in \mathbb{R};\ M\chi_{E}(x) > |E|/(4b)\}) \\
& \leq C_{0}^{p}\left(\dfrac{4b}{|E|}\right)^{p}\displaystyle\int_{\mathbb{R}}\left(\chi_{E}(x)\right)^{p}\dd\mu(x) = C_{0}^{p}\left(\frac{4b}{|E|}\right)^{p}\mu(E).
\end{aligned}\label{eq:1.11.6}\end{equation}
因 $\mu$ 是正测度,至少对充分大的 $b$ 有 $\mu((-b, b)) > 0$.因此由 \eqref{eq:1.11.6},$\mu(E) > 0$.

(ii) 设 $Q$ 为开区间,$f \in L^{p}(\dd\mu)$.因 $|Q|^{-1}\int_{Q}|f(y)|\dd y \leq M(f\chi_{Q})(x)$($x \in Q$),
\begin{equation*}
|f|_{Q} := \frac{1}{|Q|}\int_{Q}|f(y)|\dd y \leq \operatorname{essinf}_{x \in Q}M(f\chi_{Q})(x).
\end{equation*}
故由 \eqref{eq:1.11.1},对任意 $0 < a < 1$,
\begin{align*}
\mu(Q) &\leq \mu(\{x \in \mathbb{R};\ M(f\chi_{Q})(x) > a|f|_{Q}\}) \\
\quad &\leq \left(\dfrac{C_{0}}{a|f|_{Q}}\right)^{p}\| f\chi_{Q} \|_{L^{p}(\dd\mu)}^{p} = \left(\dfrac{C_{0}}{a|f|_{Q}}\right)^{p}\displaystyle\int_{Q}|f(y)|^{p}\dd\mu(y).
\end{align*}
由此,因 $0 < a < 1$ 任意,令 $a \to 1$ 得
\begin{equation*}
\frac{1}{|Q|}\int_{Q}|f(x)|\dd x \leq C_{0}\left(\frac{1}{\mu(Q)}\int_{Q}|f(x)|^{p}\dd\mu(x)\right)^{1/p}.
\end{equation*}

(iii) 设 $Q$ 为开区间.首先考虑 $1 < p < \infty$ 的情形.任取正数 $\varepsilon$,则 $(w(x) + \varepsilon)^{-1/(p-1)} \in L^{\infty}(\mathbb{R})$,故
\begin{equation*}
\int_{Q}\left((w(x) + \varepsilon)^{-1/(p-1)}\right)^{p}w(x)\dd x < +\infty.
\end{equation*}
因此在 (ii) 中取 $f = (w(x) + \varepsilon)^{-1/(p-1)}$,
\begin{align*}
\displaystyle\int_{Q}w(x)\dd x\left(\frac{1}{|Q|}\int_{Q}(w(x) + \varepsilon)^{-1/(p-1)}\dd x\right)^{p}&\leq C_{0}^{p}\displaystyle\int_{Q}(w(x) + \varepsilon)^{-p/(p-1)}w(x)\dd x  \\
&= C_{0}^{p}\displaystyle\int_{Q}(w(x) + \varepsilon)^{-1/(p-1)}\frac{w(x)}{w(x) + \varepsilon}\dd x \\
&\leq C_{0}^{p}\displaystyle\int_{Q}(w(x) + \varepsilon)^{-1/(p-1)}\dd x.
\end{align*}
因此
\begin{equation*}
\left(\frac{1}{|Q|}\int_{Q}w(x)\dd x\right)\left(\frac{1}{|Q|}\int_{Q}(w(x) + \varepsilon)^{-1/(p-1)}\dd x\right)^{p-1} \leq C_{0}^{p}.
\end{equation*}
令 $\varepsilon \to 0$,
\begin{equation*}
\left(\frac{1}{|Q|}\int_{Q}w(x)\dd x\right)\left(\frac{1}{|Q|}\int_{Q}w(x)^{-1/(p-1)}\dd x\right)^{p-1} \leq C_{0}^{p}.
\end{equation*}

其次考虑 $p = 1$ 的情形.在 (ii) 中对 $E \subset Q$ 取 $f = \chi_{E}$,
\begin{equation}\tag{1.11.7}
\frac{|E|}{|Q|} \leq C_{0}\frac{\mu(E)}{\mu(Q)} \quad \text{换句话说} \quad \frac{1}{|Q|}\int_{Q}w(x)\dd x \leq C_{0}\frac{1}{|E|}\int_{E}w(x)\dd x.
\label{eq:1.11.7}\end{equation}
这里,若 $a > \operatorname{ess}\inf_{x \in Q}w(x) = \inf\{b > 0;\ |\{x \in Q;\ w(x) < b\}| > 0\}$,则 $|\{x \in Q;\ w(x) < a\}| > 0$,故在 \eqref{eq:1.11.7} 中取 $E = \{x \in Q;\ w(x) < a\}$,
\begin{align*}
\dfrac{1}{|Q|}\displaystyle\int_{Q}w(x)\dd x &\leq C_{0}\frac{1}{|\{y \in Q;\ w(y) < a\}|}\int_{\{y \in Q;\ w(y) < a\}}w(x)\dd x \\
\qquad \qquad \qquad &\leq C_{0}\dfrac{1}{|\{y \in Q;\ w(y) < a\}|}\int_{\{y \in Q;\ w(y) < a\}}a\dd x \leq C_{0}a.
\end{align*}
令 $a \to \operatorname{ess}\inf_{x \in Q}w(x)$,
\begin{equation*}
\frac{1}{|Q|}\int_{Q}w(x)\dd x \leq C_{0}\operatorname{ess}\inf_{x \in Q}w(x).
\end{equation*}

(iv) 对 $E \subset Q$,在 (ii) 中令 $f = \chi_{E}$,
\begin{equation*}
\frac{|E|}{|Q|} \leq C_{0}\left(\frac{\mu(E)}{\mu(Q)}\right)^{1/p}.
\end{equation*}
\end{proof}

后面会证明,由上述 (iii) 可推出 Hardy--Littlewood 极大函数的弱 $L^{p}(w(x)\dd x)$ 性,故引入下面的权重类.

\begin{definition}\label{def:1.11.2}
$\mathbb{R}$ 上的权重 $w(x)$ 满足:对任意区间 $Q$,当 $1 < p < \infty$ 时
\begin{equation}\tag{1.11.8}
\left(\frac{1}{|Q|}\int_{Q}w(x)\dd x\right)\left(\frac{1}{|Q|}\int_{Q}w(x)^{-1/(p-1)}\dd x\right)^{p-1} \leq A^{p},
\label{eq:1.11.8}\end{equation}
当 $p = 1$ 时
\begin{equation}\tag{1.11.9}
\frac{1}{|Q|}\int_{Q}w(x)\dd x \leq A\operatorname{ess}\inf_{x \in Q}w(x)
\label{eq:1.11.9}\end{equation}
若存在满足上述条件的常数 $A$,则称 $w$ 属于(Muckenhoupt 的)$A_{p}$ 类,记为 $w \in A_{p}$.另外,上述常数 $A$ 的下确界称为 $w$ 的 $A_{p}$ 范数,记为 $[w]_{A_{p}}$.

此外,令 $A_{\infty} = \bigcup_{p \geq 1}A_{p}$.
\end{definition}

注意,若权重 $w(x)$ 满足 \eqref{eq:1.11.8} 或 \eqref{eq:1.11.9},则对充分大的区间 $Q$ 有 $\int_{Q}w(x)\dd x > 0$,故由 \eqref{eq:1.11.8} 或 \eqref{eq:1.11.9} 可知 $w(x) > 0$(a.e. $x \in \mathbb{R}$).

另外,把关于 $A_{p}$ 的注意以引理的形式列出.

\begin{lemma}\label{lem:1.11.3}
(i) 若 $1 \leq p < q < \infty$,则 $A_{p} \subset A_{q}$,$[w]_{A_{q}} \leq [w]_{A_{p}}$.

(ii) $[w]_{A_{p}} \geq 1$.

(iii) 若 $1 \leq p < \infty$,$w \in A_{p}$,则对一切开区间 $Q$ 与一切 $f \in L^{p}(\dd\mu)$,
\begin{equation}\tag{1.11.10}
\frac{1}{|Q|}\int_{Q}|f(x)|\dd x \leq [w]_{A_{p}}\left(\frac{1}{w(Q)}\int_{Q}|f(x)|^{p}w(x)\dd x\right)^{1/p}.
\label{eq:1.11.10}\end{equation}

(iv) 设 $1 \leq p < \infty$,$w \in A_{p}$.则对任意开区间 $Q$ 与一切可测集 $E \subset Q$,
\begin{equation*}
\frac{|E|}{|Q|} \leq [w]_{A_{p}}\left(\frac{w(E)}{w(Q)}\right)^{1/p}.
\end{equation*}

(v) 设 $1 \leq p < \infty$,$w \in A_{p}$.又设 $0 < \alpha < 1$.则对任意开区间 $Q$ 与一切满足 $|E| < \alpha|Q|$ 的可测集 $E \subset Q$,存在常数 $0 < \beta < 1$ 使
\begin{equation*}
w(E) \leq \beta w(Q).
\end{equation*}

(vi) 若 $1 \leq p < \infty$,$w \in A_{p}$,则 $w^{-1/(p-1)} \in A_{p'}$($1/p + 1/p' = 1$),且
\begin{equation*}
[w^{-1/(p-1)}]_{A_{p'}} = [w]_{A_{p}}.
\end{equation*}

这里,$w(E) = \int_{E}w(x)\dd x$.
\end{lemma}

\begin{proof}
(i) 设 $w \in A_{p}$,$Q$ 为开区间.当 $p > 1$ 时,因 $(q - 1)/(p - 1) > 1$,由 Hölder 不等式,
\begin{align*}
\left(\dfrac{1}{|Q|}\displaystyle\int_{Q}w(x)^{-1/(q-1)}\dd x\right)^{q-1}&= \left(\dfrac{1}{|Q|}\displaystyle\int_{Q}\left(w(x)^{-1/(p-1)}\right)^{(p-1)/(q-1)}\dd x\right)^{q-1}  \\
&\leq \left(\dfrac{1}{|Q|}\displaystyle\int_{Q}w(x)^{-1/(p-1)}\dd x\right)^{p-1}.
\end{align*}
由此立即得结论.

当 $p = 1$ 时,因 $w(x) \geq [w]_{A_{1}}^{-1}|Q|^{-1}\int_{Q}w(y)\dd y$(a.e. $x \in Q$),故
\begin{align*}
\left(\dfrac{1}{|Q|}\displaystyle\int_{Q}w(x)\dd x\right)\left(\dfrac{1}{|Q|}\int_{Q}w(x)^{-1/(q-1)}\dd x\right)^{q-1}&\leq \left(\dfrac{1}{|Q|}\displaystyle\int_{Q}w(x)\dd x\right)\left(\dfrac{1}{|Q|}\int_{Q}[w]_{A_{1}}^{1/(q-1)}\left(\dfrac{1}{|Q|}\int_{Q}w(x)\dd x\right)^{-1/(q-1)}\dd x\right)^{q-1}  \\
&= [w]_{A_{1}}.
\end{align*}

(ii) 当 $p > 1$ 时,在 \eqref{eq:1.11.8} 中把 $Q$ 缩向 $x$,则对几乎所有的 $x$,$|Q|^{-1}\int_{Q}w(y)\dd y$ 收敛到 $w(x)$,$|Q|^{-1}\int_{Q}w(x)^{-1/(p-1)}\dd x$ 收敛到 $w(x)^{-1/(p-1)}$,故 $1 \leq [w]_{A_{p}}$.当 $p = 1$ 时,注意 $\operatorname{ess}\inf_{x \in Q}w(x) \leq |Q|^{-1}\int_{Q}w(x)\dd x$,则与 $p > 1$ 的情形相同.

(iii) 设 $Q$ 为开区间,$f \in L^{p}(\dd\mu)$.当 $p > 1$ 时,由 Hölder 不等式,
\begin{align*}
\dfrac{1}{|Q|}\displaystyle\int_{Q}|f(x)|\dd x&= \dfrac{1}{|Q|}\displaystyle\int_{Q}|f(x)|w(x)^{1/p}w(x)^{-1/p}\dd x  \\
&\leq \left(\dfrac{1}{|Q|}\displaystyle\int_{Q}|f(x)|^{p}w(x)\dd x\right)^{1/p}\left(\dfrac{1}{|Q|}\int_{Q}w(x)^{-p'/p}\dd x\right)^{1/p'} \\
&= \left(\dfrac{1}{|Q|}\displaystyle\int_{Q}|f(x)|^{p}w(x)\dd x\right)^{1/p}\left(\dfrac{1}{|Q|}\int_{Q}w(x)^{-1/(p-1)}\dd x\right)^{1-1/p} \\
&\leq \left(\dfrac{1}{|Q|}\displaystyle\int_{Q}|f(x)|^{p}w(x)\dd x\right)^{1/p}[w]_{A_{p}}\left(\dfrac{1}{|Q|}\int_{Q}w(x)\dd x\right)^{-1/p} \\
&= [w]_{A_{p}}\left(\dfrac{1}{w(Q)}\displaystyle\int_{Q}|f(x)|^{p}w(x)\dd x\right)^{1/p}.
\end{align*}

当 $p = 1$ 时,利用 $w(x) \geq [w]_{A_{1}}^{-1}|Q|^{-1}\int_{Q}w(y)\dd y$(a.e. $x \in Q$),
\begin{align*}
\dfrac{1}{|Q|}\displaystyle\int_{Q}|f(x)|\dd x & = \dfrac{1}{|Q|}\int_{Q}|f(x)|w(x)w(x)^{-1}\dd x \\
& \leq \dfrac{1}{|Q|}\displaystyle\int_{Q}|f(x)|w(x)\dd x[w]_{A_{1}}\left(\dfrac{1}{|Q|}\int_{Q}w(y)\dd y\right)^{-1} \\
& = [w]_{A_{1}}\dfrac{1}{w(Q)}\displaystyle\int_{Q}|f(x)|w(x)\dd x.
\end{align*}

(iv) 在 (iii) 中取 $f(x) = \chi_{E}$ 即可.

(v) 设 $0 < \alpha < 1$,$Q$ 为开区间,$E$ 为满足 $E \subset Q$ 与 $|E| < \alpha|Q|$ 的可测集.在式~\eqref{eq:1.11.10} 中取 $f = \chi_{Q\setminus E}$,
\begin{align*}
& \left(\dfrac{|Q \setminus E|}{|Q|}\right)^{p} \leq [w]_{A_{p}}^{p}\dfrac{w(Q \setminus E)}{w(Q)} \quad \text{改写为} \\
& \left(1 - \dfrac{|E|}{|Q|}\right)^{p} \leq [w]_{A_{p}}^{p}\left(1 - \dfrac{w(E)}{w(Q)}\right).
\end{align*}
因此由 $|E| < \alpha|Q|$,
\begin{equation*}
(1 - \alpha)^{p} \leq [w]_{A_{p}}^{p}\left(1 - \frac{w(E)}{w(Q)}\right) \quad \text{整理得} \quad \frac{w(E)}{w(Q)} \leq 1 - \frac{(1 - \alpha)^{p}}{[w]_{A_{p}}^{p}}.
\end{equation*}
由 (ii) 知 $[w]_{A_{p}} \geq 1$,又 $0 < \alpha < 1$,故令
\begin{equation*}
\beta = 1 - \frac{(1 - \alpha)^{p}}{[w]_{A_{p}}^{p}},
\end{equation*}
则 $0 < \beta < 1$,满足所求条件.

(vi) 因 $p' = p/(p - 1)$,故 $(p' - 1)(p - 1) = 1$.由它与
\begin{equation}\tag{1.11.11}
\left(\frac{1}{|Q|}\int_{Q}w(x)\dd x\right)\left(\frac{1}{|Q|}\int_{Q}w(x)^{-1/(p-1)}\dd x\right)^{p-1} \leq [w]_{A_{p}}^{p}
\label{eq:1.11.11}\end{equation}
得
\begin{align*}
\left(\dfrac{1}{|Q|}\displaystyle\int_{Q}w(x)^{-1/(p-1)}\dd x\right)\left(\dfrac{1}{|Q|}\int_{Q}\left(w(x)^{-1/(p-1)}\right)^{-1/(p'-1)}\dd x\right)^{p'-1}&\leq [w]_{A_{p}}^{p/(p-1)} = [w]_{A_{p}}^{p'}.
 \end{align*}
因此 $w^{-1/(p-1)} \in A_{p'}$,且 $[w^{-1/(p-1)}]_{A_{p'}} \leq [w]_{A_{p}}$.把 $p$ 与 $p'$ 的角色互换,得 $[w]_{A_{p}} \leq [w^{-1/(p-1)}]_{A_{p'}} \leq [w]_{A_{p}}$,故 $[w^{-1/(p-1)}]_{A_{p'}} = [w]_{A_{p}}$.
\end{proof}

其次,考虑加权中心极大函数,考察其弱 $(1,1)$ 性.

\begin{lemma}\label{lem:1.11.4}
设 $\mu$ 为 $\mathbb{R}$ 上的正测度,满足 $\mu(Q(x, r)) > 0$($x \in \mathbb{R}$,$r > 0$).对 $g \in L_{\mathrm{loc}}^{1}(\mathbb{R}, \dd\mu)$ 定义
\begin{equation*}
\tilde{M}_{\mu}(g)(x) = \sup_{r > 0}\frac{1}{\mu(Q(x, r))}\int_{Q(x, r)}|g(y)|\dd\mu(y).
\end{equation*}
则对一切 $f \in L^{1}(\mathbb{R}, \dd\mu)$,
\begin{equation*}
\mu(\{x \in \mathbb{R};\ \tilde{M}_{\mu}(f)(x) > \lambda\}) \leq \frac{4}{\lambda}\int_{\mathbb{R}}|f(y)|\dd\mu(y) \quad (\lambda > 0).
\end{equation*}
\end{lemma}

\begin{proof}
与定理~\ref{thm:1.5.2} 中 Hardy--Littlewood 中心极大函数 $\tilde{M}f$ 的弱 $(1,1)$ 性的证明完全相同,故省略.
\end{proof}

由上述两个引理,可以证明当 $w \in A_{p}$ 类时 Hardy--Littlewood 极大函数的弱 $L^{p}(w\dd x)$ 性.

\begin{theorem}\label{thm:1.11.5}
设 $1 \leq p < \infty$ 且 $w \in A_{p}$.则对一切 $f \in L^{p}(w(x)\dd x)$,
\begin{equation*}
w(\{x \in \mathbb{R};\ Mf(x) > \lambda\}) \leq \frac{4\cdot 2^{p}[w]_{A_{p}}^{p}}{\lambda^{p}}\int_{\mathbb{R}}|f(x)|^{p}w(x)\dd x \quad (\lambda > 0).
\end{equation*}
\end{theorem}

\begin{proof}
对 $g \in L_{\mathrm{loc}}^{1}(\mathbb{R}, w(x)\dd x)$ 定义
\begin{equation*}
\tilde{M}_{w}g(x) = \sup_{r > 0}\frac{1}{w(Q(x, r))}\int_{Q(x, r)}|g(y)|w(y)\dd y,
\end{equation*}
则由引理~\ref{lem:1.4.4} 与引理~\ref{lem:1.11.3} (iii),
\begin{equation*}
Mf(x) \leq 2[w]_{A_{p}}\left(\tilde{M}_{w}(|f|^{p})(x)\right)^{1/p} \quad (x \in \mathbb{R}).
\end{equation*}
故由引理~\ref{lem:1.11.4},
\begin{align*}
w(\{x \in \mathbb{R};\ Mf(x) > \lambda\}) &\leq w\left(\left\{x \in \mathbb{R};\ \tilde{M}_{w}(|f|^{p})(x) > \lambda^{p}/(2[w]_{A_{p}})^{p}\right\}\right) \\
&\leq \dfrac{4\cdot 2^{p}[w]_{A_{p}}^{p}}{\lambda^{p}}\displaystyle\int_{\mathbb{R}}|f(y)|^{p}w(y)\dd y \quad (\lambda > 0).
\end{align*}
\end{proof}

由这些结果,可以证明对 $A_{p}$ 权重称为逆 Hölder 不等式的下述事实.

\begin{theorem}\label{thm:1.11.6}
设 $1 \leq p < \infty$ 且 $w \in A_{p}$.则存在常数 $C_{0} > 0$,$\gamma > 0$ 使
\begin{equation}\tag{1.11.12}
\left(\frac{1}{|Q|}\int_{Q}w(y)^{1+\gamma}\dd y\right)^{1/(1+\gamma)} \leq \frac{C_{0}}{|Q|}\int_{Q}w(y)\dd y
\label{eq:1.11.12}\end{equation}
对任意区间 $Q$ 成立.这里 $C_{0}, \gamma$ 仅依赖于 $[w]_{A_{p}}$ 与 $p$.
\end{theorem}

\begin{proof}
设 $1 \leq p < \infty$,$w \in A_{p}$,令 $Q = [a, b)$.又设 $0 < \alpha < 1$.进一步令
\begin{equation*}
\alpha_{0} = \frac{1}{|Q|}\int_{Q}w(x)\dd x, \quad \alpha_{k} = \left(\frac{2}{\alpha}\right)^{k}\alpha_{0} \quad (k \in \mathbb{N}_{0}).
\end{equation*}
对 $k \in \mathbb{N}_{0}$,把 $Q$ 作二进分解,对 $w(x)$ 与 $\alpha_{k}$ 作与 Calderón--Zygmund 分解时相同的操作,把所得 $Q$ 的二进分解区间列记为 $\{Q_{k,j}\}$.则

(1) $\alpha_{k} < \dfrac{1}{|Q_{k,j}|}\displaystyle\int_{Q_{k,j}}w(x)\dd x \leq 2\alpha_{k}$.

(2) $w(x) \leq \alpha_{k}$(a.e. $x \in \left(\bigcup_{j}Q_{k,j}\right)^{c}$).

(3) 每个 $Q_{k+1,j}$ 包含在某个 $Q_{k,\ell}$ 中.

(1),(2) 已在 Calderón--Zygmund 分解处证明.关于 (3),因
\begin{equation*}
\frac{1}{|Q_{k+1,j}|}\int_{Q_{k+1,j}}w(x)\dd x > \alpha_{k+1} > \alpha_{k},
\end{equation*}
故在包含 $Q_{k+1,j}$ 的 $Q$ 的二进分解区间 $R$ 中,若不存在满足 $|R|^{-1}\int_{R}w(x)\dd x > \alpha_{k}$ 者,则 $Q_{k+1,j}$ 自身就是某个 $Q_{k,\ell}$;若存在满足 $|R|^{-1}\int_{R}w(x)\dd x > \alpha_{k}$ 者,则其中宽度最大者就是某个 $Q_{k,\ell}$,且 $Q_{k+1,j}$ 包含在其中.

现在,令 $U_{m} = \bigcup_{j}Q_{m,j}$($m \in \mathbb{N}_{0}$),则对 $k \in \mathbb{N}_{0}$,
\begin{align*}
2\alpha_{k} &\geq \dfrac{1}{|Q_{k,\ell}|}\displaystyle\int_{Q_{k,\ell} \cap U_{k+1}}w(y)\dd y \\
&= \dfrac{1}{|Q_{k,\ell}|}\displaystyle\sum_{Q_{k+1,j} \subset Q_{k,\ell}}|Q_{k+1,j}|\frac{1}{|Q_{k+1,j}|}\int_{Q_{k+1,j}}w(y)\dd y \\
&> \dfrac{|Q_{k,\ell} \cap U_{k+1}|}{|Q_{k,\ell}|}\alpha_{k+1} = \dfrac{|Q_{k,\ell} \cap U_{k+1}|}{|Q_{k,\ell}|}\frac{2}{\alpha}\alpha_{k}.
\end{align*}
因此 $|Q_{k,\ell} \cap U_{k+1}| < \alpha|Q_{k,\ell}|$.由引理~\ref{lem:1.11.3} (v),
\begin{equation*}
\frac{w(Q_{k,\ell} \cap U_{k+1})}{w(Q_{k,\ell})} < \beta = 1 - \frac{(1 - \alpha)^{p}}{[w]_{A_{p}}^{p}}.
\end{equation*}
由此,考虑 (3) 得
\begin{align*}
w(U_{k+1}) &= w\left(\left(\bigcup_{\ell}Q_{k,\ell}\right) \cap U_{k+1}\right) \\
&= \displaystyle\sum_{\ell}w(Q_{k,\ell} \cap U_{k+1}) \\
&< \displaystyle\sum_{\ell}\beta w(Q_{k,\ell}) = \beta w(U_{k}).
\end{align*}
故 $w(U_{k}) < \beta^{k}w(U_{0})$.另外,同样可得 $|U_{k}| < \alpha^{k}|U_{0}|$,故 $|U_{k}| \to 0$($k \to \infty$).因此 $|\bigcap_{k}U_{k}| = 0$.由此
\begin{equation}\tag{1.11.13}
\left|Q \setminus \left((Q \setminus U_{0}) \cup \left(\bigcup_{k=0}^{\infty}(U_{k} \setminus U_{k+1})\right)\right)\right| = 0.
\label{eq:1.11.13}\end{equation}
这里,取正数 $\gamma$ 使
\begin{equation*}
(2\alpha^{-1})^{\gamma}\beta < 1.
\end{equation*}
则由 \eqref{eq:1.11.13} 与 (2),
\begin{align*}
\displaystyle\int_{Q}w(y)^{1+\gamma}\dd y &= \int_{Q \setminus U_{0}}w(y)^{\gamma}w(y)\dd y + \sum_{k=0}^{\infty}\int_{U_{k} \setminus U_{k+1}}w(y)^{\gamma}w(y)\dd y \\
\qquad &\leq \alpha_{0}^{\gamma}w(Q \setminus U_{0}) + \displaystyle\sum_{k=0}^{\infty}\alpha_{k+1}^{\gamma}w(U_{k}) \\
\qquad &\leq \alpha_{0}^{\gamma}w(Q \setminus U_{0}) + \displaystyle\sum_{k=0}^{\infty}\left((2\alpha^{-1})^{k+1}\alpha_{0}\right)^{\gamma}\beta^{k}w(U_{0}) \\
\qquad &\leq \alpha_{0}^{\gamma}\left(1 + (2\alpha^{-1})^{\gamma}\sum_{k=0}^{\infty}(2\alpha^{-1})^{\gamma k}\beta^{k}\right)\left(w(Q \setminus U_{0}) + w(U_{0})\right) \\
\qquad &= \left(\dfrac{1}{|Q|}\displaystyle\int_{Q}w(y)\dd y\right)^{\gamma}\left(1 + \dfrac{(2\alpha^{-1})^{\gamma}}{1 - (2\alpha^{-1})^{\gamma}\beta}\right)\int_{Q}w(y)\dd y.
\end{align*}
因此,选取正数 $\gamma, C$ 满足
\begin{align*}
\gamma &< \dfrac{-\log\beta}{\log 2 - \log\alpha} = \dfrac{\log[w]_{A_{p}}^{p} - \log\left([w]_{A_{p}}^{p} - (1 - \alpha)^{p}\right)}{\log 2 - \log\alpha}, \\
C^{1+\gamma} &= 1 + \dfrac{(2\alpha^{-1})^{\gamma}}{1 - (2\alpha^{-1})^{\gamma}\beta} = 1 + \dfrac{(2\alpha^{-1})^{\gamma}}{1 - (2\alpha^{-1})^{\gamma}\left(1 - (1 - \alpha)^{p}/[w]_{A_{p}}^{p}\right)},
\end{align*}
则
\begin{equation*}
\left(\frac{1}{|Q|}\int_{Q}w(y)^{1+\gamma}\dd y\right)^{1/(1+\gamma)} \leq \frac{C}{|Q|}\int_{Q}w(y)\dd y.
\end{equation*}
\end{proof}

作为逆 Hölder 不等式的直接推论,

\begin{lemma}\label{lem:1.11.7}
(i) 设 $1 \leq p < \infty$,$w \in A_{p}$,$\gamma, C_{0}$ 为定理~\ref{thm:1.11.6} 中的常数.则对一切区间 $Q$ 与任意可测集 $E \subset Q$,
\begin{equation*}
\frac{w(E)}{w(Q)} \leq C_{0}\left(\frac{|E|}{|Q|}\right)^{\gamma/(\gamma+1)}.
\end{equation*}

(ii) 设 $1 < p < \infty$,$w \in A_{p}$.则存在仅与 $[w]_{A_{p}}$ 和 $p$ 有关的常数 $p_{0}$($1 < p_{0} < p$)使 $w \in A_{p_{0}}$.
\end{lemma}

\begin{proof}
(i) 设 $Q$ 为区间,$E$ 为其可测子集.由 Hölder 不等式与定理~\ref{thm:1.11.6},
\begin{align*}
\displaystyle\int_{E}w(x)\dd x &\leq \left(\int_{E}\dd x\right)^{\gamma/(\gamma+1)}\left(\int_{E}w(x)^{1+\gamma}\dd x\right)^{1/(1+\gamma)} \\
&\leq |E|^{\gamma/(\gamma+1)}|Q|^{1/(1+\gamma)}\dfrac{C_{0}}{|Q|}\int_{Q}w(x)\dd x.
\end{align*}
整理得
\begin{equation*}
\frac{w(E)}{w(Q)} \leq C_{0}\left(\frac{|E|}{|Q|}\right)^{\gamma/(1+\gamma)}.
\end{equation*}

(ii) 由引理~\ref{lem:1.11.3} (vi),$w^{-1/(p-1)} \in A_{p'}$($1/p + 1/p' = 1$)且 $[w^{-1/(p-1)}]_{A_{p'}} = [w]_{A_{p}}$.对这个 $w^{-1/(p-1)}$ 使用逆 Hölder 不等式(定理~\ref{thm:1.11.6}),把此时的常数记为 $\gamma$,$C_{0}$,用 $(1 + \gamma)/(p - 1) = 1/(p_{0} - 1)$ 确定 $p_{0}$,则 $1 < p_{0} < p$ 且
\begin{equation*}
\begin{aligned}
\left(\dfrac{1}{|Q|}\displaystyle\int_{Q}w(x)^{-1/(p_{0}-1)}\dd x\right)^{p_{0}-1}&= \left(\dfrac{1}{|Q|}\displaystyle\int_{Q}\left(w(x)^{-1/(p-1)}\right)^{1+\gamma}\dd x\right)^{(p-1)/(1+\gamma)}
 \end{aligned}
\end{equation*}
\begin{align*}
&\leq \left(C_{0}\dfrac{1}{|Q|}\displaystyle\int_{Q}w(x)^{-1/(p-1)}\dd x\right)^{p-1} \\
&\leq C_{0}^{p-1}[w]_{A_{p}}^{p}\left(\dfrac{1}{|Q|}\displaystyle\int_{Q}w(x)\dd x\right)^{-1}.
\end{align*}
也就是说,
\begin{equation*}
w \in A_{p_{0}}, \quad [w]_{A_{p_{0}}} \leq C_{0}^{(p-1)/p_{0}}[w]_{A_{p}}^{p/p_{0}}.
\end{equation*}
\end{proof}

作为这个应用,当 $1 < p < \infty$ 时,得到强化定理~\ref{thm:1.11.5} 的 Hardy--Littlewood 极大函数的加权 $L^{p}$ 估计.

\begin{lemma}\label{lem:1.11.8}
设 $1 < p < \infty$,$w \in A_{p}$.则存在常数 $C$ 使
\begin{equation*}
\| Mf \|_{L^{p}(w\dd x)} \leq C\| f \|_{L^{p}(w\dd x)} \quad (f \in L^{p}(w\dd x)).
\end{equation*}
\end{lemma}

\begin{proof}
设 $1 < p < \infty$,$w \in A_{p}$,则由引理~\ref{lem:1.11.7} (ii),存在 $p_{0}$($1 < p_{0} < p$)使 $w \in A_{p_{0}}$.因此由定理~\ref{thm:1.11.5},存在常数 $C_{0}$ 使
\begin{align*}
w(\{x \in \mathbb{R};\ Mf(x) > \lambda\}) & \leq \dfrac{C_{0}}{\lambda^{p_{0}}}\displaystyle\int_{\mathbb{R}}|f(x)|^{p_{0}}w(x)\dd x \\
& \quad (f \in L^{p_{0}}(w\dd x),\ \lambda > 0).
\end{align*}
另外,显然
\begin{equation*}
\| Mf \|_{L^{\infty}(w\dd x)} \leq \| f \|_{L^{\infty}(w\dd x)}.
\end{equation*}
所以在 $p_{0}$ 与 $\infty$ 之间应用 Marcinkiewicz 插值定理~\ref{thm:1.8.6},
\begin{equation*}
\| Mf \|_{L^{p}(w\dd x)} \leq C\| f \|_{L^{p}(w\dd x)} \quad (f \in L^{p}(w\dd x)).
\end{equation*}
\end{proof}

另外,可以证明 Hardy--Littlewood 极大函数与锐极大函数之间的加权``good $\lambda$-不等式'',从而也得到 Hilbert 变换的加权 $L^{p}$ 结论.

\begin{lemma}\label{lem:1.11.9}
设 $1 \leq p < \infty$,$w \in A_{p}$.则当 $f \in L_{\mathrm{loc}}^{1}(\mathbb{R})$ 时,对任意 $\lambda > 0$,$\gamma > 0$,存在常数 $C > 0$ 与 $\delta > 0$ 使
\begin{equation*}
w(\{x \in \mathbb{R};\ Mf(x) > 2\lambda,\ f^{\#}(x) < \gamma\lambda\}) \leq C\gamma^{\delta}w(\{x \in \mathbb{R};\ Mf(x) > \lambda\}).
\end{equation*}
\end{lemma}

\begin{proof}
证明只需对无权重情形的引理~\ref{lem:1.6.5} 的证明作部分修改即可.到式~\eqref{eq:1.6.1},即
\begin{equation*}
|\{x \in Q;\ Mf(x) > 2\lambda,\ f^{\#}(x) < \gamma\lambda\}| \leq 96\gamma|Q|
\end{equation*}
为止完全相同.由这个不等式与引理~\ref{lem:1.11.7} (i)(把那里的常数 $C_{0}$,$\gamma$ 记为 $C_{0}$,$\beta$),得
\begin{equation}\tag{1.11.14}
w(\{x \in Q;\ Mf(x) > 2\lambda,\ f^{\#}(x) < \gamma\lambda\}) \leq C_{0}(96\gamma)^{\beta/(1+\beta)}w(Q).
\label{eq:1.11.14}\end{equation}
因此,把对 $Q_{j}$ 的估计式~\eqref{eq:1.11.14} 对 $Q_{j}$ 求和,因 $Q_{j}$ 的内部互不相交,
\begin{align*}
w(\{x \in \mathbb{R};\ Mf(x) > 2\lambda,\ f^{\#}(x) < \gamma\lambda\})&= w(\{x \in E_{\lambda};\ Mf(x) > 2\lambda,\ f^{\#}(x) < \gamma\lambda\})  \\
&\leq C_{0}(96\gamma)^{\beta/(1+\beta)}w(E_{\lambda}) = C_{0}(96\gamma)^{\beta/(1+\beta)}w(\{x \in \mathbb{R};\ Mf(x) > \lambda\}).
\end{align*}
\end{proof}

把这个加权的``good $\lambda$-不等式''与引理~\ref{lem:1.6.6} 组合,就得到定理~\ref{thm:1.6.8} 的加权版本.

\begin{lemma}\label{lem:1.11.10}
设 $1 \leq p < \infty$,$w \in A_{p}$.则当 $Mf \in L^{p}(w\dd x)$ 时,存在常数 $C$ 使
\begin{equation*}
\| Mf \|_{L^{p}(w\dd x)} \leq C\| f^{\#} \|_{L^{p}(w\dd x)} \quad (f \in L^{p}(w\dd x)).
\end{equation*}
\end{lemma}

\begin{proof}
只需对无权重情形的引理~\ref{lem:1.6.6} 的证明稍作修改即可,故省略证明.
\end{proof}

有了以上准备,就能给出 Hilbert 变换的加权 $L^{p}$ 估计.

\begin{theorem}\label{thm:1.11.11}
(i) 当 $1 < p < \infty$,$w \in A_{p}$ 时,存在常数 $C$ 使
\begin{equation*}
\| Hf \|_{L^{p}(w\dd x)} \leq C\| f \|_{L^{p}(w\dd x)} \quad (f \in L^{p}(w\dd x)).
\end{equation*}

(ii) 当 $w \in A_{1}$ 时,存在常数 $C$ 使
\begin{equation*}
w(\{x \in \mathbb{R};\ |Hf(x)| > \lambda\}) \leq \frac{C}{\lambda}\| f \|_{L^{1}(w\dd x)} \quad (f \in L^{1}(w\dd x)).
\end{equation*}
\end{theorem}

\begin{proof}
(i) 设 $f \in \mathcal{S}(\mathbb{R})$.已知 $|Hf(x)| \leq C_{1}(1 + |x|)^{-1}$,故 $M(Hf)(x) \leq 4C_{1}(1 + |x|)^{-1}\log(e + |x|)$(这个计算留给读者).于是 $M(Hf) \in L^{p}(w\dd x)$.实际上,由引理~\ref{lem:1.11.7} (ii) 存在 $1 < p_{0} < p$ 使 $w \in A_{p_{0}}$,故利用引理~\ref{lem:1.11.3} (iv),对 $j \geq 1$ 可导出 $w(\{|x| < 2^{j+1}\})/w(\{|x| < 2\}) \leq [w]_{A_{p_{0}}}^{p_{0}}(2^{j+1}/2)^{p_{0}} \leq [w]_{A_{p}}^{p_{0}}2^{jp_{0}}$,故当 $j \geq 1$ 时,
\begin{align*}
\displaystyle\int_{2^{j} \leq |x| < 2^{j+1}}M(Hf)(x)^{p}w(x)\dd x&\leq \left(4C_{1}\dfrac{\log(e + 2^{j+1})}{1 + 2^{j}}\right)^{p}\displaystyle\int_{2^{j} \leq |x| < 2^{j+1}}w(x)\dd x  \\
&\leq (4C_{1}\log 2)^{p}(j + 2)^{p}2^{-jp}\displaystyle\int_{|x| < 2^{j+1}}w(x)\dd x \\
&= (4C_{1}\log 2)^{p}(j + 2)^{p}2^{-jp}\dfrac{w(\{|x| < 2^{j+1}\})}{w(\{|x| < 2\})}w(\{|x| < 2\}) \\
&\leq (4C_{1}\log 2)^{p}(j + 2)^{p}2^{-jp}C_{2}2^{jp_{0}}w(\{|x| < 2\}) \\
&\leq C_{3}j^{p}2^{-(p-p_{0})j}w(\{|x| < 2\}).
\end{align*}
因此,
\begin{equation*}
\int_{\mathbb{R}}M(Hf)(x)^{p}w(x)\dd x
\end{equation*}
\begin{align*}
&= \displaystyle\int_{|x| < 2}M(Hf)(x)^{p}w(x)\dd x + \sum_{j=1}^{\infty}\int_{2^{j} \leq |x| < 2^{j+1}}M(Hf)(x)^{p}w(x)\dd x \\
&\leq \left(4C_{1}\log(e + 2)\right)^{p}w(\{|x| < 2\}) + \displaystyle\sum_{j=1}^{\infty}C_{3}j^{p}2^{-(p-p_{0})j}w(\{|x| < 2\}) < +\infty.
\end{align*}
由此利用引理~\ref{lem:1.11.10},
\begin{equation}\tag{1.11.15}
\| Hf \|_{L^{p}(w\dd x)} \leq \| M(Hf) \|_{L^{p}(w\dd x)} \leq C_{4}\| (Hf)^{\#} \|_{L^{p}(w\dd x)}.
\label{eq:1.11.15}\end{equation}
另一方面,由引理~\ref{lem:1.11.7} (ii) 存在 $1 < p_{0} < p$ 使 $w \in A_{p_{0}}$,又由系~\ref{cor:1.7.4},
\begin{equation*}
(Hf)^{\#}(x) \leq C_{5}M_{p/p_{0}}f(x) \quad (x \in \mathbb{R}),
\end{equation*}
故由引理~\ref{lem:1.11.8},
\begin{equation}\tag{1.11.16}\begin{aligned}
\| (Hf)^{\#} \|_{L^{p}(w\dd x)} & \leq C_{5}\| M_{p/p_{0}}f \|_{L^{p}(w\dd x)} \\
& = C_{5}\left(\displaystyle\int_{\mathbb{R}}\left(M(|f|^{p/p_{0}})(x)\right)^{p_{0}}w(x)\dd x\right)^{1/p} \\
& \leq C_{5}\left[C_{6}\left(\displaystyle\int_{\mathbb{R}}\left(|f|^{p/p_{0}}(x)\right)^{p_{0}}w(x)\dd x\right)^{1/p_{0}}\right]^{p_{0}/p} \\
& \leq C_{6}\| f \|_{L^{p}(w\dd x)}.
\end{aligned}\label{eq:1.11.16}\end{equation}
合并 \eqref{eq:1.11.15},\eqref{eq:1.11.16},
\begin{equation*}
\| Hf \|_{L^{p}(w\dd x)} \leq C_{4}C_{6}\| f \|_{L^{p}(w\dd x)}.
\end{equation*}
可以证明 $\mathcal{S}(\mathbb{R})$ 在 $L^{p}(w\dd x)$ 中稠密,故得到所求结论.

(ii) 设 $w \in A_{1}$,$f \in L^{1}(w\dd x) \cap L^{2}(w\dd x)$,$\lambda > 0$.首先确认 $f \in L_{\mathrm{loc}}^{1}(\mathbb{R})$.设 $Q$ 为区间,则
\begin{align}
\frac{1}{|Q|}\int_{Q}|f(x)|\dd x &\leq \frac{1}{|Q|}\int_{Q}|f(x)|w(x)\dd x\cdot\frac{1}{\operatorname{essinf}_{x \in Q}w(x)} \notag \\
&\leq \frac{1}{|Q|}\int_{Q}|f(x)|w(x)\dd x\cdot\frac{[w]_{A_{1}}}{\frac{1}{|Q|}\int_{Q}w(x)\dd x} \notag \\
&= \frac{[w]_{A_{1}}}{w(Q)}\int_{Q}|f(x)|w(x)\dd x \tag{1.11.17}\label{eq:1.11.17} \\
&\leq \frac{[w]_{A_{1}}}{w(Q)}\| f \|_{L^{1}(w\dd x)} < +\infty. \tag{1.11.18}\label{eq:1.11.18}
\end{align}
因此 $f \in L_{\mathrm{loc}}^{1}(\mathbb{R})$.其次,对 Calderón--Zygmund 分解的引理~\ref{lem:1.8.1} 的证明稍作修改,确认现在的情形下也能作同样的分解.引入二进极大函数
\begin{equation*}
M_{d}f(x) = \sup_{x \in Q \in \mathbb{D}}\frac{1}{|Q|}\int_{Q}|f(y)|\dd y,
\end{equation*}
令
\begin{equation*}
\Omega = \{x \in \mathbb{R};\ M_{d}f(x) > \lambda\}.
\end{equation*}
若 $x \in \Omega$,则由上确界的性质,存在包含 $x$ 且满足 $1/|Q|\int_{Q}|f(y)|\dd y > \lambda$ 的二进区间 $Q$.反之,若 $1/|Q|\int_{Q}|f(y)|\dd y > \lambda$,则对所有 $x \in Q$ 有 $M_{d}f(x) > \lambda$.因此,把满足 $1/|Q|\int_{Q}|f(y)|\dd y > \lambda$ 的 $Q \in \mathbb{D}$ 的全体记为 $\{R_{j}\}$,则
\begin{equation*}
\Omega = \bigcup_{j}R_{j}.
\end{equation*}
对 $\{R_{j}\}$ 中不断扩大的子列 $S_{1} \subset S_{2} \subset \cdots$,由 \eqref{eq:1.11.18} 与引理~\ref{lem:1.11.7},
\begin{align*}
\dfrac{1}{|S_{j}|}\displaystyle\int_{S_{j}}|f(x)|\dd x & \leq [w]_{A_{1}}\dfrac{w(S_{1})}{w(S_{j})}\dfrac{\| f \|_{L^{1}(w\dd x)}}{w(S_{1})} \\
& \leq C_{0}[w]_{A_{1}}\left(\dfrac{|S_{1}|}{|S_{j}|}\right)^{\gamma/(1+\gamma)}\dfrac{\| f \|_{L^{1}(w\dd x)}}{w(S_{1})},
\end{align*}
故可知这种子列只有有限个.因此由引理~\ref{lem:1.6.3},在 $\{R_{j}\}$ 中取极大二进区间 $\{Q_{j}\}$,则
\begin{equation*}
\Omega = \bigcup_{j}Q_{j}, \quad Q_{k} \cap Q_{\ell} = \emptyset \quad (k \neq \ell),
\end{equation*}
且由 $Q_{j}$ 的极大性,
\begin{equation*}
\lambda < \frac{1}{|Q_{j}|}\int_{Q_{j}}|f(x)|\dd x = |f|_{Q_{j}} \leq 2\lambda.
\end{equation*}
也利用 \eqref{eq:1.11.17},
\begin{equation*}
\lambda < \frac{[w]_{A_{1}}}{w(Q_{j})}\int_{Q_{j}}|f(x)|w(x)\dd x,
\end{equation*}
故
\begin{equation}\tag{1.11.19}
w(\Omega) = \sum_{j}w(Q_{j}) < \sum_{j}\frac{[w]_{A_{1}}}{\lambda}\int_{Q_{j}}|f(x)|w(x)\dd x = \frac{[w]_{A_{1}}}{\lambda}\| f \|_{L^{1}(w\dd x)}.
\label{eq:1.11.19}\end{equation}
另外,由 $M_{d}f(x) \leq \lambda$($x \in \Omega^{c}$),
\begin{equation*}
|f(x)| \leq \lambda \quad (\text{a.e. } x \in \Omega^{c}).
\end{equation*}
于是,像定理~\ref{thm:1.8.2} 的证明那样定义
\begin{align*}
g(x) &= \begin{cases} f(x) & (x \in \Omega^{c}) \\ \dfrac{1}{|Q_{j}|}\displaystyle\int_{Q_{j}}f(x)\dd x & (x \in Q_{j},\ j = 1, 2, \dots) \end{cases} \\[10pt]
f_{j}(x) &= \left(f(x) - \dfrac{1}{|Q_{j}|}\displaystyle\int_{Q_{j}}f(y)\dd y\right)\chi_{Q_{j}} \\[8pt]
b(x) &= \displaystyle\sum_{j}f_{j}(x).
\end{align*}
则 $f(x) = g(x) + b(x)$,且
\begin{equation*}
|g(x)| \leq 2\lambda \quad (\text{a.e. } x \in \mathbb{R}).
\end{equation*}
进而,由 \eqref{eq:1.11.17} 与 $[w]_{A_{1}} \geq 1$,
\begin{align*}
\displaystyle\int_{\mathbb{R}}|g(x)|w(x)\dd x&\leq \displaystyle\int_{\Omega^{c}}|f(x)|w(x)\dd x + \sum_{j}\int_{Q_{j}}\left(\frac{1}{|Q_{j}|}\int_{Q_{j}}|f(y)|\dd y\right)w(x)\dd x  \\
&= \displaystyle\int_{\Omega^{c}}|f(x)|w(x)\dd x + \sum_{j}w(Q_{j})\frac{1}{|Q_{j}|}\int_{Q_{j}}|f(y)|\dd y \\
&\leq \displaystyle\int_{\Omega^{c}}|f(x)|w(x)\dd x + \sum_{j}[w]_{A_{1}}\int_{Q_{j}}|f(y)|w(y)\dd y \\
&\leq [w]_{A_{1}}\displaystyle\int_{\mathbb{R}}|f(x)|w(x)\dd x.
\end{align*}
由此 $g \in L^{1}(w\dd x) \cap L^{\infty}(w\dd x)$,故亦有 $g \in L^{1}(w\dd x) \cap L^{2}(w\dd x)$.

其次,设 $Q_{j}$ 的中心为 $x_{j}$,
\begin{equation*}
Q_{j}^{*} = 2Q_{j} = Q(x_{j}, 2\ell(Q_{j})), \quad \Omega^{*} = \bigcup_{j}Q_{j}^{*},
\end{equation*}
则由引理~\ref{lem:1.11.3} (iv) 与 \eqref{eq:1.11.19},令 $C_{1} = 2[w]_{A_{1}}$,
\begin{equation}\tag{1.11.20}
w(\Omega^{*}) \leq \sum_{j}w(Q_{j}^{*}) = \sum_{j}C_{1}w(Q_{j}) = C_{1}w(\Omega) \leq \frac{C_{2}}{\lambda}\| f \|_{L^{1}(w\dd x)}.
\label{eq:1.11.20}\end{equation}
现在,由 $f(x) = g(x) + b(x)$ 得 $Hf(x) = Hg(x) + Hb(x)$,若 $|Hg(x)| \leq \lambda/2$ 且 $|Hb(x)| \leq \lambda/2$,则 $|Hf(x)| \leq |Hg(x)| + |Hb(x)| \leq \lambda$,故
\begin{align*}
& \{x \in \mathbb{R};\ |Hf(x)| > \lambda,\ x \in (\Omega^{*})^{c}\} \\
& \qquad \subset \{x \in (\Omega^{*})^{c};\ |Hg(x)| > \lambda/2\} \cup \{x \in (\Omega^{*})^{c};\ |Hb(x)| > \lambda/2\}.
\end{align*}
令 $E_{1} = \{x \in (\Omega^{*})^{c};\ |Hg(x)| > \lambda/2\}$,$E_{2} = \{x \in (\Omega^{*})^{c};\ |Hb(x)| > \lambda/2\}$.首先,利用 $H$ 的 $L^{2}(w\dd x)$ 有界性与 \eqref{eq:1.11.19},
\begin{equation}\tag{1.11.21}\begin{aligned}
w(E_{1}) &\leq \dfrac{1}{(\lambda/2)^{2}}\displaystyle\int_{E_{1}}|Hg(x)|^{2}w(x)\dd x \leq \dfrac{4}{\lambda^{2}}\int_{\mathbb{R}}|Hg(x)|^{2}w(x)\dd x \\
&\leq \dfrac{C_{3}}{\lambda^{2}}\displaystyle\int_{\mathbb{R}}|g(x)|^{2}w(x)\dd x \\
&= \dfrac{C_{3}}{\lambda^{2}}\left(\displaystyle\int_{\Omega}|g(x)|^{2}w(x)\dd x + \int_{\Omega^{c}}|g(x)|^{2}w(x)\dd x\right) \\
&= \dfrac{C_{3}}{\lambda^{2}}\left(\displaystyle\sum_{j}\int_{Q_{j}}|f_{Q_{j}}|^{2}w(x)\dd x + \lambda\int_{\Omega^{c}}|g(x)|w(x)\dd x\right) \\
&\quad \quad \quad \quad \quad \quad \quad \quad \quad \quad \quad \quad \quad \quad \quad (|g(x)| \leq \lambda,\ \text{a.e. } x \in \Omega^{c}) \\
&\leq \dfrac{C_{3}}{\lambda^{2}}\left(\displaystyle\sum_{j}w(Q_{j})(2\lambda)^{2} + \lambda\int_{\mathbb{R}}|f(x)|w(x)\dd x\right) \\
&\quad \quad \quad \quad \quad \quad \quad \quad \quad \quad (|f_{Q_{j}}| \leq 2\lambda,\ g(x) = f(x)\ (x \in \Omega^{c})) \\
&= 4C_{3}\displaystyle\sum_{j}w(Q_{j}) + \dfrac{C_{3}}{\lambda}\| f \|_{L^{1}(w\dd x)} \leq \dfrac{C_{4}}{\lambda}\| f \|_{L^{1}(w\dd x)}.
\end{aligned}\label{eq:1.11.21}\end{equation}
其次,若 $x \in (Q_{j}^{*})^{c}$,则
\begin{equation*}
|Hb(x)| \leq \sum_{j}|Hf_{j}(x)| \leq \frac{8\lambda}{\pi}\sum_{j}\frac{\ell(Q_{j})^{2}}{|x - x_{j}|^{2} + \ell(Q_{j})^{2}},
\end{equation*}
故注意 $\ell(Q_{j}) = |Q_{j}|$,
\begin{equation}\tag{1.11.22}\begin{aligned}
w(E_{2}) &\leq \dfrac{1}{\lambda/2}\displaystyle\int_{E_{2}}|Hb(x)|w(x)\dd x \leq \dfrac{2}{\lambda}\int_{(\Omega^{*})^{c}}|Hb(x)|w(x)\dd x \\
\quad &\leq \dfrac{16}{\pi}\displaystyle\int_{(\Omega^{*})^{c}}\sum_{j}\frac{|Q_{j}|^{2}}{|x - x_{j}|^{2} + |Q_{j}|^{2}}w(x)\dd x \\
\quad &\leq \dfrac{16}{\pi}\displaystyle\sum_{j}|Q_{j}|\int_{-\infty}^{\infty}\frac{|Q_{j}|}{|x - x_{j}|^{2} + |Q_{j}|^{2}}w(x)\dd x.
\end{aligned}\label{eq:1.11.22}\end{equation}
这里,令 $2^{\ell}Q_{j} = \left[x_{j} - 2^{\ell-1}|Q_{j}|,\ x_{j} + 2^{\ell-1}|Q_{j}|\right)$,则由 $w \in A_{1}$,
\begin{align*}
& \displaystyle\int_{-\infty}^{\infty}\frac{|Q_{j}|}{|x - x_{j}|^{2} + |Q_{j}|^{2}}w(x)\dd x \\
& = \displaystyle\int_{|x - x_{j}| < |Q_{j}|/2}\frac{|Q_{j}|}{|x - x_{j}|^{2} + |Q_{j}|^{2}}w(x)\dd x \\
& \qquad + \displaystyle\sum_{m=0}^{\infty}\int_{2^{m-1}|Q_{j}| \leq |x - x_{j}| < 2^{m}|Q_{j}|}\frac{|Q_{j}|}{|x - x_{j}|^{2} + |Q_{j}|^{2}}w(x)\dd x \\
& \leq \dfrac{1}{|Q_{j}|}\displaystyle\int_{|x - x_{j}| < |Q_{j}|/2}w(x)\dd x \\
& \qquad + \displaystyle\sum_{m=0}^{\infty}\frac{1}{2^{m-1}}\int_{|x - x_{j}| < 2^{m}|Q_{j}|}\frac{1}{2^{m-1}|Q_{j}|}w(x)\dd x \\
& \leq \dfrac{w(Q_{j})}{|Q_{j}|} + \displaystyle\sum_{m=0}^{\infty}\frac{1}{2^{m-2}}\frac{w(2^{m}Q_{j})}{|2^{m}Q_{j}|} \\
& \leq \dfrac{w(Q_{j})}{|Q_{j}|} + \displaystyle\sum_{m=0}^{\infty}\frac{1}{2^{m-2}}[w]_{A_{1}}\operatorname{essinf}_{x \in 2^{m}Q_{j}}w(x) \\
& \leq \dfrac{w(Q_{j})}{|Q_{j}|} + \displaystyle\sum_{m=0}^{\infty}\frac{1}{2^{m-2}}[w]_{A_{1}}\operatorname{essinf}_{x \in Q_{j}}w(x) \\
& \leq \dfrac{w(Q_{j})}{|Q_{j}|} + \displaystyle\sum_{m=0}^{\infty}\frac{1}{2^{m-2}}[w]_{A_{1}}\frac{w(Q_{j})}{|Q_{j}|} \\
& = (1 + 8[w]_{A_{1}})\dfrac{w(Q_{j})}{|Q_{j}|} \leq 9[w]_{A_{1}}\dfrac{w(Q_{j})}{|Q_{j}|}.
\end{align*}
因此由 \eqref{eq:1.11.22},\eqref{eq:1.11.19},
\begin{equation*}
w(E_{2}) \leq \frac{16}{\pi}\sum_{j}9[w]_{A_{1}}w(Q_{j}) \leq C_{5}w(\Omega) \leq \frac{C_{6}}{\lambda}\| f \|_{L^{1}(w\dd x)}.
\end{equation*}
因此与 \eqref{eq:1.11.21} 合并,
\begin{equation}\tag{1.11.23}\begin{aligned}
w(\{x \in \mathbb{R};\ |Hf(x)| > \lambda,\ x \in (\Omega^{*})^{c}\}) & \leq w(E_{1}) + w(E_{2}) \\
& \leq \dfrac{C_{4} + C_{6}}{\lambda}\| f \|_{L^{1}(w\dd x)}.
\end{aligned}\label{eq:1.11.23}\end{equation}
由 \eqref{eq:1.11.20},\eqref{eq:1.11.23},
\begin{align*}
w(\{x \in \mathbb{R};\ |Hf(x)| > \lambda\}) & \leq w(\Omega^{*}) + w(\{x \in \mathbb{R};\ |Hf(x)| > \lambda,\ x \in (\Omega^{*})^{c}\}) \\
& \leq \dfrac{C_{2} + C_{4} + C_{6}}{\lambda}\| f \|_{L^{1}(w\dd x)}.
\end{align*}
由引理~\ref{lem:1.8.4},对任意 $f \in L^{1}(w\dd x)$,$Hf$ 被唯一地定义,且此估计式成立.
\end{proof}

% ============ 第1章 §1.12 Hardy空间(原书 pp.108-114) ============
\section{Hardy 空间}\label{sec:1.12}
\index{Hardy 空间!(R 上的)}
\index{H1 范数}


在 $p = 1$ 的情形,Hilbert 变换在 $L^{1}(\mathbb{R})$ 上不是有界的,但本节引入 Hilbert 变换在其上有界的空间----Hardy 空间 $H^{1}(\mathbb{R})$.设 $f \in L^{1}(\mathbb{R})$,且存在 $g \in L^{1}(\mathbb{R})$ 使 $-i\xi|\xi|^{-1}\hat{f}(\xi) = \hat{g}(\xi)$($\xi \in \mathbb{R}$).设 $\varphi \in \mathcal{S}(\mathbb{R})$,$\int_{\mathbb{R}}\varphi(x)\dd x = 1$,$\varphi_{\varepsilon}(x) = \varphi(x/\varepsilon)/\varepsilon$.由例~\ref{exa:1.1.10} 与定理~\ref{thm:1.1.9},当 $\varepsilon \to 0$ 时,$f * \varphi_{\varepsilon}$,$g * \varphi_{\varepsilon}$ 分别依 $L^{1}(\mathbb{R})$ 范数收敛到 $f$,$g$.另外,因 $f, g, \varphi_{\varepsilon} \in L^{1}(\mathbb{R})$,
\begin{align*}
\mathcal{F}(g * \varphi_{\varepsilon}) & = \mathcal{F}(g)\mathcal{F}(\varphi_{\varepsilon}) = -i\dfrac{\xi}{|\xi|}\mathcal{F}(f)\mathcal{F}(\varphi_{\varepsilon}) \\
& = -i\dfrac{\xi}{|\xi|}\mathcal{F}(f * \varphi_{\varepsilon}) = \mathcal{F}\left(H(f * \varphi_{\varepsilon})\right).
\end{align*}
因 $f, g \in L^{1}(\mathbb{R})$,$\varphi_{\varepsilon} \in L^{2}(\mathbb{R})$,故 $f * \varphi_{\varepsilon}, g * \varphi_{\varepsilon} \in L^{2}(\mathbb{R})$,因此 $H(f * \varphi_{\varepsilon}) \in L^{2}(\mathbb{R})$,
\begin{equation*}
g * \varphi_{\varepsilon} = H(f * \varphi_{\varepsilon}).
\end{equation*}
令 $\varepsilon \to 0$,则左边依 $L^{1}(\mathbb{R})$ 范数收敛到 $g$,右边依测度收敛到 $Hf$.因此 $g = Hf$.

考虑这一点,定义下面的函数空间.

\begin{definition}\label{def:1.12.1}
使 $-i\xi|\xi|^{-1}\hat{f}(\xi) = \hat{g}(\xi)$($\xi \in \mathbb{R}$)成立的 $g \in L^{1}(\mathbb{R})$ 存在的 $f \in L^{1}(\mathbb{R})$ 的全体,称为 Hardy 空间 $H^{1}(\mathbb{R})$.对 $f \in H^{1}(\mathbb{R})$,
\begin{equation*}
\| f \|_{H^{1}(\mathbb{R})} = \| f \|_{H^{1}} = \| f \|_{L^{1}(\mathbb{R})} + \| Hf \|_{L^{1}(\mathbb{R})}
\end{equation*}
称为 $f$ 的 $H^{1}$ 范数.
\end{definition}

$H^{1}$ 范数是范数是显然的.设 $\{f_{j}\}_{j=1}^{\infty}$ 为 $H^{1}(\mathbb{R})$ 的 Cauchy 列.则因 $L^{1}(\mathbb{R})$ 完备,存在 $f, g \in L^{1}(\mathbb{R})$ 使 $\| f_{j} - f \|_{L^{1}} \to 0$,$\| Hf_{j} - g \|_{L^{1}} \to 0$($j \to \infty$).由定理~\ref{thm:1.2.3},
\begin{equation*}
\lim_{j \to \infty}\mathcal{F}f_{j}(\xi) = \mathcal{F}f(\xi), \quad \lim_{j \to \infty}\mathcal{F}(Hf_{j})(\xi) = \mathcal{F}g(\xi).
\end{equation*}
另一方面,因 $\mathcal{F}(Hf_{j})(\xi) = -i\xi|\xi|^{-1}\mathcal{F}f_{j}(\xi)$,
\begin{equation*}
-i\frac{\xi}{|\xi|}\mathcal{F}f(\xi) = \mathcal{F}g(\xi).
\end{equation*}
因此 $f \in H^{1}(\mathbb{R})$ 且 $\lim_{j \to \infty}\| f_{j} - f \|_{H^{1}} = 0$,可知 $H^{1}(\mathbb{R})$ 是 Banach 空间.

若 $f \in L^{1}(\mathbb{R}) \cap L^{2}(\mathbb{R})$ 且 $Hf \in L^{1}(\mathbb{R})$,则因 $f \in L^{2}(\mathbb{R})$,$\mathcal{F}(Hf) = -i\xi|\xi|^{-1}\mathcal{F}(f)$,故 $f \in H^{1}(\mathbb{R})$.下面考虑更具体的例子.

\begin{example}\label{exa:1.12.2}
若 $f \in C^{1}(\mathbb{R})$ 且 $|f'(x)| \leq A_{1}(1 + |x|)^{-2}$,$|f(x)| \leq A_{2}(1 + |x|)^{-3}$($x \in \mathbb{R}$),且 $\int_{\mathbb{R}}f(x)\dd x = 0$,则 $f \in H^{1}(\mathbb{R})$.

因此,特别是当 $f \in \mathcal{S}(\mathbb{R})$ 且 $\int_{\mathbb{R}}f(x)\dd x = 0$ 时,$f \in H^{1}(\mathbb{R})$.

实际上,由引理~\ref{lem:1.3.2},$|Hf(x)| \leq C_{1}(A_{1} + A_{2})(1 + |x|)^{-1}$.当 $|x| \geq 1$ 时,由 $\int_{\mathbb{R}}f(x)\dd x = 0$,
\begin{equation*}
Hf(x) = \frac{1}{\pi}\mathrm{p.v.}\int_{\mathbb{R}}\frac{1}{x - y}f(y)\dd y - \frac{1}{\pi}\int_{\mathbb{R}}\frac{1}{x}f(y)\dd y
\end{equation*}
\begin{equation*}
= \frac{1}{\pi x}\mathrm{p.v.}\int_{\mathbb{R}}\frac{1}{x - y}yf(y)\dd y,
\end{equation*}
故注意 $|(xf(x))'| \leq (A_{1} + A_{2})(1 + |x|)^{-1}$,$|xf(x)| \leq A_{2}(1 + |x|)^{-2}$,再使用引理~\ref{lem:1.3.2},得 $|Hf(x)| \leq 2C_{1}(A_{1} + A_{2})(1 + |x|)^{-2} \in L^{1}(\mathbb{R})$.又 $f \in L^{1}(\mathbb{R}) \cap L^{2}(\mathbb{R})$,故 $f \in H^{1}(\mathbb{R})$.这里为以后起见注意 $|H_{\varepsilon}f(x)| \leq 2C_{1}(A_{1} + A_{2})(1 + |x|)^{-2}$($\varepsilon > 0$).
\end{example}

由 $H^{1}(\mathbb{R})$ 的定义,若 $f \in H^{1}(\mathbb{R})$,则 $Hf \in L^{1}(\mathbb{R})$ 且 $-i\xi|\xi|^{-1}\mathcal{F}f(\xi) = \mathcal{F}(Hf)(\xi)$,故 $-\mathcal{F}(f)(\xi) = -i\xi|\xi|^{-1}\mathcal{F}(Hf)(\xi)$.因此,与 $f \in L^{1}(\mathbb{R})$ 合并,$Hf \in H^{1}(\mathbb{R})$ 且 $\| Hf \|_{H^{1}} = \| f \|_{H^{1}}$.总结为

\begin{theorem}\label{thm:1.12.3}
Hilbert 变换在 $H^{1}(\mathbb{R})$ 上有界,且
\begin{equation*}
\| Hf \|_{H^{1}} = \| f \|_{H^{1}} \quad (f \in H^{1}(\mathbb{R})).
\end{equation*}
\end{theorem}

下面讨论在 $H^{1}(\mathbb{R})$ 中稠密的好的空间.

\begin{lemma}\label{lem:1.12.4}
$H_{00}^{1}(\mathbb{R}) = \{f \in L^{1}(\mathbb{R});\ \hat{f} \in C^{\infty}(\mathbb{R}) \text{ 且 } \hat{f} \text{ 的支集是有界闭集且不包含原点}\}$ 在 $H^{1}(\mathbb{R})$ 中稠密.因此,$\{f \in \mathcal{S}(\mathbb{R});\ \int_{\mathbb{R}}f(x)\dd x = 0\}$ 在 $H^{1}(\mathbb{R})$ 中稠密,$H^{1}(\mathbb{R}) \cap L^{2}(\mathbb{R})$ 也在 $H^{1}(\mathbb{R})$ 中稠密.
\end{lemma}

\begin{proof}
设 $f \in H^{1}(\mathbb{R})$,任取 $\varepsilon > 0$.由 $\mathcal{F}(Hf)(\xi) = -i\xi|\xi|^{-1}\hat{f}(\xi)$,且 $Hf \in L^{1}(\mathbb{R})$,故 $-i\xi|\xi|^{-1}\hat{f}(\xi)$ 是连续函数.因此 $\hat{f}(0) = 0$,$\mathcal{F}(Hf)(0) = 0$.于是取 $\varphi \in \mathcal{S}(\mathbb{R})$ 使 $\hat{\varphi}(\xi) = 1$($|\xi| \leq 1$),$\hat{\varphi}(\xi) = 0$($|\xi| \geq 2$).首先证明 $\lim_{A \to \infty}\| \varphi_{A} * f \|_{L^{1}} = 0$.由 $\hat{f}(0) = 0$ 得 $\int_{\mathbb{R}}f(x)\dd x = 0$,故
\begin{align*}
\displaystyle\int_{\mathbb{R}}|\varphi_{A} * f(x)|\dd x & = \int_{\mathbb{R}}\left|\int_{\mathbb{R}}\left(\varphi_{A}(x - y) - \varphi_{A}(x)\right)f(y)\dd y\right|\dd x \\
& \leq \displaystyle\int_{\mathbb{R}}\left(\int_{\mathbb{R}}\frac{1}{A}\left|\varphi\left(\frac{x - y}{A}\right) - \varphi\left(\frac{x}{A}\right)\right|\dd x\right)|f(y)|\dd y \\
& = \displaystyle\int_{\mathbb{R}}\left(\int_{\mathbb{R}}|\varphi(x - y/A) - \varphi(x)|\dd x\right)|f(y)|\dd y.
\end{align*}
这里,$\int_{\mathbb{R}}|\varphi(x - y/A) - \varphi(x)|\dd x \leq 2\int_{\mathbb{R}}|\varphi(x)|\dd x$,且因 $\varphi \in L^{1}(\mathbb{R})$,$\lim_{A \to \infty}\int_{\mathbb{R}}|\varphi(x - y/A) - \varphi(x)|\dd x = 0$.因此由 Lebesgue 收敛定理,
\begin{equation*}
\lim_{A \to \infty}\int_{\mathbb{R}}\left(\int_{\mathbb{R}}|\varphi(x - y/A) - \varphi(x)|\dd x\right)|f(y)|\dd y = 0.
\end{equation*}
由此 $\lim_{A \to \infty}\| \varphi_{A} * f \|_{L^{1}} = 0$.同样,$\lim_{A \to \infty}\| \varphi_{A} * (Hf) \|_{L^{1}} = 0$.于是取 $A > 0$ 使 $\| \varphi_{A} * f \|_{L^{1}} < \varepsilon$,$\| \varphi_{A} * (Hf) \|_{L^{1}} < \varepsilon$,令 $g = f - \varphi_{A} * f$,则 $\| g - f \|_{L^{1}} < \varepsilon$,$\| Hg - Hf \|_{L^{1}} = \| \varphi_{A} * (Hf) \|_{L^{1}} < \varepsilon$.另外,由 $\hat{\varphi}(\xi) = 1$($|\xi| \leq 1$)与 $\hat{g}(\xi) = \hat{f}(\xi) - \hat{\varphi}(A\xi)\hat{f}(\xi)$,$\hat{g}(\xi) = 0$($|\xi| \leq 1/A$).其次,由逼近定理~\ref{thm:1.1.9},$\lim_{t \to 0}\| \varphi_{t} * g - g \|_{L^{1}} = 0$,$\lim_{t \to 0}\| \varphi_{t} * (Hg) - Hg \|_{L^{1}} = 0$,故取 $0 < a < A$ 使 $\| \varphi_{a} * g - g \|_{L^{1}} < \varepsilon$,$\| \varphi_{a} * (Hg) - Hg \|_{L^{1}} < \varepsilon$,令 $h = \varphi_{a} * g$.由 $\hat{h}(\xi) = \hat{\varphi}(a\xi)\hat{g}(\xi)$,$\hat{g}(\xi) = 0$($|\xi| \leq 1/A$),$\hat{\varphi}(\xi) = 0$($|\xi| \geq 2$),故 $\hat{h}(\xi) = 0$($|\xi| \leq 1/A$ 或 $|\xi| \geq 2/a$).注意 $\hat{h} \in L^{1}(\mathbb{R})$.

最后,取 $\psi \in \mathcal{S}(\mathbb{R})$ 使 $\psi(\xi) = 0$($|\xi| \geq 1$),$\mathcal{F}^{-1}\psi(0) = (2\pi)^{-1}\int_{\mathbb{R}}\psi(\xi)\dd\xi = 1$.则 $|h(x)\mathcal{F}^{-1}\psi(tx) - h(x)| \leq (\| \psi \|_{L^{1}} + 1)|h(x)|$,且对每个 $x \in \mathbb{R}$,$\lim_{t \to 0}\left(h(x)\mathcal{F}^{-1}\psi(tx) - h(x)\right) = 0$,故由 Lebesgue 收敛定理,$\lim_{t \to 0}\int_{\mathbb{R}}|h(x)\mathcal{F}^{-1}\psi(tx) - h(x)|\dd x = 0$.另外,因 $\hat{h} \in L^{1}(\mathbb{R})$,$\mathcal{F}^{-1}(\hat{h} * \psi_{t})(x) = 2\pi\mathcal{F}^{-1}\hat{h}(x)\mathcal{F}^{-1}(\psi_{t})(x) = 2\pi h(x)\mathcal{F}^{-1}\psi(tx)$,因此 $\mathcal{F}\left(2\pi h(\cdot)\mathcal{F}^{-1}\psi(t\cdot)\right) = \hat{h} * \psi_{t}$.由此 $\mathcal{F}\left(H\left(h(\cdot)\mathcal{F}^{-1}\psi(t\cdot)\right)\right)(\xi) = -i\xi|\xi|^{-1}\hat{h} * \psi_{t}(\xi)$.

现在,设 $0 < t < 1/(2A)$.若 $|\xi|, |\eta| < 1/(2A)$,则 $|\xi - \eta| < 1/A$,且 $\hat{h}(\xi - \eta) = 0$($|\xi - \eta| < 1/A$),$\psi_{t}(\eta) = 0$($|\eta| > t$),故 $\int_{\mathbb{R}}\hat{h}(\xi - \eta)\psi_{t}(\eta)\dd\eta = 0$($|\xi| < 1/(2A)$).同样,若 $|\xi| > 2/a + 1/(2A)$,$|\eta| < 1/(2A)$,则 $|\xi - \eta| > 2/a$,且 $\hat{h}(\xi - \eta) = 0$($|\xi - \eta| > 2/a$),$\psi_{t}(\eta) = 0$($|\eta| > t$),故 $\int_{\mathbb{R}}\hat{h}(\xi - \eta)\psi_{t}(\eta)\dd\eta = 0$($|\xi| > 2/a + 1/(2A)$).于是取 $k \in \mathcal{S}(\mathbb{R})$ 使 $k(\xi) = -i\xi|\xi|^{-1}$($1/(2A) < |\xi| < 2/a + 1/(2A)$),令 $K(x) = \mathcal{F}^{-1}k(x)$.则 $K \in \mathcal{S}(\mathbb{R})$,且因 $\hat{h}(\xi) = 0$($|\xi| \leq 1/A$ 或 $|\xi| \geq 2/a$),$\hat{h} * \psi_{t}(\xi) = 0$($|\xi| \leq 1/(2A)$ 或 $|\xi| \geq 2/a + 1/(2A)$),故 $k(\xi)\hat{h}(\xi) = -i\xi|\xi|^{-1}\hat{h}(\xi)$,$k(\xi)\hat{h} * \psi_{t}(\xi) = -i\xi|\xi|^{-1}\hat{h} * \psi_{t}(\xi)$,
\begin{equation*}
Hh(x) = K * h(x), \quad H\left(h(\cdot)\mathcal{F}^{-1}\psi(t\cdot)\right)(x) = K * \left(h\mathcal{F}^{-1}(\psi_{t})\right)(x).
\end{equation*}
由此
\begin{align*}
\| Hh - H\left(h(\cdot)\mathcal{F}^{-1}\psi(t\cdot)\right) \|_{L^{1}} & = \| K * \left(h - h\mathcal{F}^{-1}(\psi_{t})\right) \|_{L^{1}} \\
& \leq \| K \|_{L^{1}}\| h - h\mathcal{F}^{-1}(\psi_{t}) \|_{L^{1}}.
\end{align*}
因此,因 $\lim_{t \to 0}\| h - h\mathcal{F}^{-1}(\psi_{t}) \|_{L^{1}} = 0$,取 $0 < b < 1/(2A)$ 使
\begin{equation*}
\| h - h\mathcal{F}^{-1}(\psi_{b}) \|_{L^{1}}, \quad \| Hh - H\left(h(\cdot)\mathcal{F}^{-1}\psi(b\cdot)\right) \|_{L^{1}} < \varepsilon,
\end{equation*}
令 $f_{0} = h\mathcal{F}^{-1}(\psi_{b})$.则
\begin{equation*}
\hat{f}_{0}(\xi) = 0 \quad \left(|\xi| < 1/(2A) \text{ 或 } |\xi| > 2/a + 1/(2A)\right),
\end{equation*}
\begin{equation*}
\| f - f_{0} \|_{L^{1}} \leq \| f - g \|_{L^{1}} + \| g - h \|_{L^{1}} + \| h - f_{0} \|_{L^{1}} < 3\varepsilon,
\end{equation*}
\begin{align*}
\| Hf - Hf_{0} \|_{L^{1}} &\leq \| Hf - Hg \|_{L^{1}} + \| Hg - Hh \|_{L^{1}} + \| Hh - Hf_{0} \|_{L^{1}} \\
&< 3\varepsilon.
\end{align*}
即 $\| f - f_{0} \|_{H^{1}} < 6\varepsilon$.另外,因 $\hat{f}_{0}(\xi) = \hat{h} * \psi_{b}(\xi)$,$\psi_{b} \in \mathcal{S}(\mathbb{R})$,$\hat{h} \in L^{1}(\mathbb{R})$,故 $\hat{f}_{0} \in C^{\infty}(\mathbb{R})$.由这些可知 $H_{00}^{1}(\mathbb{R})$ 在 $H^{1}(\mathbb{R})$ 中稠密.
\end{proof}

其次,考察复平面上半平面中的正则函数与 Hardy 空间的关系.
\begin{equation*}
P(x) = \frac{1}{\pi}\frac{1}{1 + x^{2}}, \quad P(x, y) = \frac{1}{\pi}\frac{y}{y^{2} + x^{2}}, \quad Q(x, y) = \frac{1}{\pi}\frac{x}{y^{2} + x^{2}} \quad (y > 0,\ x \in \mathbb{R})
\end{equation*}
令($y > 0$,$x \in \mathbb{R}$).此时,$P_{t}(x) = P(x/t)/t = P(x, t)$,且由 $\int_{0}^{\infty}e^{-t\xi}e^{ix\xi}\dd\xi = (t - ix)^{-1}$,
\begin{equation*}
\mathcal{F}^{-1}(e^{-t|\xi|})(x) = P(x, t), \quad \mathcal{F}^{-1}\left(-i\frac{\xi}{|\xi|}e^{-t|\xi|}\right)(x) = Q(x, t) \quad (t > 0).
\end{equation*}
对 $g \in L^{1}(\mathbb{R})$,$h \in L^{2}(\mathbb{R})$,$\mathcal{F}(g * h)(\xi) = \mathcal{F}(g)(\xi)\mathcal{F}(h)(\xi)$ 容易证明,故注意 $Q(x, t)$ 关于 $x$ 属于 $L^{2}(\mathbb{R})$,则若 $f \in H^{1}(\mathbb{R})$ 且为实值函数,
\begin{equation*}
\begin{aligned}
\mathcal{F}(P_{t} * (Hf))(\xi) &= \mathcal{F}(P_{t})(\xi)\mathcal{F}(Hf)(\xi) = e^{-t|\xi|}\cdot\left(-i\dfrac{\xi}{|\xi|}\right)\hat{f}(\xi) \\
&= \mathcal{F}(Q(\cdot, t))(\xi)\hat{f}(\xi) = \mathcal{F}(Q(\cdot, t) * f)(\xi),
\end{aligned}
\end{equation*}
故对 $z = x + iy$,
\begin{equation*}
F(z) := \left(P_{y} * (f + iHf)\right)(x) = P(\cdot, y) * f(x) + iQ(\cdot, y) * f(x).
\end{equation*}
$P(x, y)$,$Q(x, y)$ 满足 Cauchy--Riemann 关系式容易验证,且注意 $f \in L^{1}(\mathbb{R})$,可知 $P(\cdot, y) * f(x)$,$Q(\cdot, y) * f(x)$ 满足 Cauchy--Riemann 关系式.因此 $F(z)$ 在上半平面正则.进而,
\begin{equation*}
\int_{\mathbb{R}}|F(x + iy)|\dd x \leq \int_{\mathbb{R}}P_{y}(x)\dd x\int_{\mathbb{R}}|f(x) + iHf(x)|\dd x \leq \| f \|_{H^{1}} \quad (y > 0).
\end{equation*}
实际上,其逆也成立.这里不给出证明,但下式成立.

\begin{theorem}\label{thm:1.12.5}
若 $F(z)$ 在上半平面正则且
\begin{equation*}
\| F \|_{H^{1}(\mathbb{R}_{+}^{2})} := \sup_{y > 0}\int_{\mathbb{R}}|F(x + iy)|\dd x < +\infty,
\end{equation*}
则 $f(x) = \lim_{y \to 0}\operatorname{Re}F(x + iy)$ 对几乎所有的 $x \in \mathbb{R}$ 存在,$f \in H^{1}(\mathbb{R})$,且
\begin{equation*}
F(z) = P_{y} * (f + iHf)(x),
\end{equation*}
$\| f \|_{H^{1}}$ 与 $\| F \|_{H^{1}(\mathbb{R}_{+}^{2})}$ 等价.
\end{theorem}

上述定义中,把 $|F(x + iy)|$ 换成 $|F(x + iy)|^{p}$($p > 0$)者,就是上半平面的一般 Hardy 空间 $H^{p}(\mathbb{R}_{+}^{2})$,它是单位圆盘上的经典 Hardy 空间 $H^{p}(U)$ 的上半平面版本.

下面不加证明地列出 $H^{1}(\mathbb{R})$ 的其他特征刻画.

\begin{theorem}\label{thm:1.12.6}
对 $f \in L^{1}(\mathbb{R})$,下面四者等价.

(i) $f \in H^{1}(\mathbb{R})$.

(ii) $\displaystyle\int_{\mathbb{R}}\sup_{t>0}|P_{t} * f(x)|\dd x < +\infty$.

(iii) $\displaystyle\int_{\mathbb{R}}\sup_{|u - x| < t}|P_{t} * f(u)|\dd x < +\infty$.

(iv)(原子分解)存在数列 $\{\lambda_{j}\}_{j=1}^{\infty}$ 与函数列 $a_{j}(x)$ 满足:

(a) $\displaystyle\sum_{j=1}^{\infty}|\lambda_{j}| < \infty$;

(b) $a_{j}(x)$ 的支集包含在某个区间 $Q_{j}$ 中,且 $\int_{Q_{j}}a_{j}(x)\dd x = 0$,$\| a_{j} \|_{L^{\infty}} \leq |Q_{j}|^{-1}$;

(c) $f(x) = \displaystyle\sum_{j=1}^{\infty}\lambda_{j}a_{j}(x)$ in $L^{1}(\mathbb{R})$.

此时,(ii),(iii) 的积分值,以及 (iv) 中 $f$ 的表示中 $\displaystyle\sum_{j=1}^{\infty}|\lambda_{j}|$ 的下确界,分别与 $\| f \|_{H^{1}}$ 等价.
\end{theorem}

% ============ 第1章 §1.13 BMO空间(原书 pp.114-122) ============
\section{BMO 空间}\label{sec:1.13}
\index{BMO 函数}
\index{有界平均振动函数}
\index{John--Nirenberg 定理}


作为本章的结束,我们讨论 Hardy 空间 $H^{1}(\mathbb{R})$ 的对偶空间,同时也是与锐极大函数有关的函数空间----BMO 空间.

\begin{definition}\label{def:1.13.1}
若 $f \in L_{\mathrm{loc}}^{1}(\mathbb{R})$ 满足
\begin{equation*}
\| f \|_{\mathrm{BMO}} := \sup_{Q}\frac{1}{|Q|}\int_{Q}|f(x) - f_{Q}|\dd x < +\infty,
\end{equation*}
则称 $f$ 为 BMO 函数(有界平均振动函数).这里,上确界对一切区间 $Q$ 取,$f_{Q} = |Q|^{-1}\int_{Q}f(x)\dd x$.
\end{definition}

容易看出,$\| f \|_{\mathrm{BMO}} = \sup_{x \in \mathbb{R}}f^{\#}(x)$.

\begin{lemma}\label{lem:1.13.2}
对 $f \in L_{\mathrm{loc}}^{1}(\mathbb{R})$,
\begin{equation*}
\sup_{Q}\inf_{c \in \mathbb{C}}\frac{1}{|Q|}\int_{Q}|f(x) - c|\dd x \leq \sup_{Q}\frac{1}{|Q|}\int_{Q}|f(x) - f_{Q}|\dd x
\end{equation*}
\begin{equation*}
\leq 2\sup_{Q}\inf_{c \in \mathbb{C}}\frac{1}{|Q|}\int_{Q}|f(x) - c|\dd x.
\end{equation*}
\end{lemma}

\begin{proof}
对区间 $Q$ 与任意 $c \in \mathbb{C}$,
\begin{align*}
\dfrac{1}{|Q|}\displaystyle\int_{Q}|f(x) - f_{Q}|\dd x & \leq \dfrac{1}{|Q|}\int_{Q}|f(x) - c|\dd x + \dfrac{1}{|Q|}\int_{Q}|c - f_{Q}|\dd x \\
& = 2\dfrac{1}{|Q|}\displaystyle\int_{Q}|f(x) - c|\dd x
\end{align*}
成立,利用它即可.
\end{proof}

与 Hilbert 变换相关,有下面的结果.

\begin{theorem}\label{thm:1.13.3}
存在常数 $C$ 使
\begin{equation*}
\| Hf \|_{\mathrm{BMO}} \leq C\| f \|_{L^{\infty}} \quad (f \in L_{c}^{\infty}(\mathbb{R})).
\end{equation*}
\end{theorem}

\begin{proof}
设 $f$ 为支集是有界集合的有界函数,$Q$ 为以 $x_{0}$ 为中心,宽 $2r$ 的区间 $(x_{0} - r, x_{0} + r)$,$2Q = (x_{0} - 2r, x_{0} + 2r)$.令 $f_{1}(x) = f(x)\chi_{2Q}(x)$,$f_{2}(x) = f(x) - f_{1}(x)$.则当 $x \in Q$,$y \in (2Q)^{c}$ 时,$|x - y| \geq |x_{0} - y| - |x_{0} - x| > |x_{0} - y|/2$,故
\begin{align*}
\left|Hf(x) - \displaystyle\int_{(2Q)^{c}}\frac{f(y)}{x_{0} - y}\dd y\right| & = \left|Hf_{1}(x) + \displaystyle\int_{(2Q)^{c}}\left(\frac{1}{x - y} - \frac{1}{x_{0} - y}\right)f(y)\dd y\right| \\
& \leq |Hf_{1}(x)| + \displaystyle\int_{(2Q)^{c}}\frac{2|x - x_{0}||f(y)|}{|x_{0} - y|^{2}}\dd y \\
& \leq |Hf_{1}(x)| + 2\| f \|_{L^{\infty}}.
\end{align*}
因此,也利用 Hilbert 变换的 $L^{2}$ 有界性(定理~\ref{thm:1.3.5}),
\begin{align*}
\dfrac{1}{|Q|}\displaystyle\int_{Q}\left|Hf(x) - \int_{(2Q)^{c}}\frac{f(y)}{x_{0} - y}\dd y\right|\dd x & \leq \dfrac{1}{|Q|}\int_{Q}|Hf_{1}(x)|\dd x + 2\| f \|_{L^{\infty}} \\
& \leq \left(\dfrac{1}{|Q|}\displaystyle\int_{Q}|Hf_{1}(x)|^{2}\dd x\right)^{1/2} + 2\| f \|_{L^{\infty}}
\end{align*}
\begin{align*}
&\leq \left(\dfrac{1}{|Q|}\displaystyle\int_{\mathbb{R}}|f_{1}(x)|^{2}\dd x\right)^{1/2} + 2\| f \|_{L^{\infty}} \\
&\leq (\sqrt{2} + 2)\| f \|_{L^{\infty}}.
\end{align*}
因 $Q$ 是任意区间,由前引理~\ref{lem:1.13.2},
\begin{equation*}
\| Hf \|_{\mathrm{BMO}} \leq 2(\sqrt{2} + 2)\| f \|_{L^{\infty}}.
\end{equation*}
\end{proof}

对 Hilbert 变换稍作修正,可知它把有界函数变成 BMO 函数.

\begin{theorem}\label{thm:1.13.4}
对 $f \in L^{\infty}(\mathbb{R})$,
\begin{equation*}
\tilde{H}f(x) = \lim_{j \to \infty}\left(H\left(f\chi_{\{|x| < j\}}\right)(x) - \frac{1}{\pi}\int_{1 < |y| < j}\frac{f(y)}{-y}\dd y\right)
\end{equation*}
对几乎所有的 $x$ 存在,并在任意有界集上依 $L^{1}$ 范数收敛.进而,存在常数 $C$ 使
\begin{equation*}
\| \tilde{H}f \|_{\mathrm{BMO}} \leq C\| f \|_{L^{\infty}} \quad (f \in L^{\infty}(\mathbb{R})).
\end{equation*}
\end{theorem}

\begin{proof}
任意固定 $A > 0$.令 $j_{0}$ 为满足 $j_{0} \geq 2A$ 的最小整数.对 $j \in \mathbb{N}$,$j > j_{0}$,当 $|x| \leq A$ 时,
\begin{align*}
g_{j}(x) &:= H\left(f\chi_{\{|x| < j\}}\right)(x) - \dfrac{1}{\pi}\int_{1 < |y| < j}\frac{f(y)}{-y}\dd y \\
&= H\left(f\chi_{\{|x| < j_{0}\}}\right)(x) - \dfrac{1}{\pi}\int_{1 < |y| < j_{0}}\frac{f(y)}{-y}\dd y \\
&\quad + \dfrac{1}{\pi}\displaystyle\int_{j_{0} \leq |y| < j}\left(\frac{1}{x - y} - \frac{1}{-y}\right)f(y)\dd y,
\end{align*}
且
\begin{equation*}
\left|\int_{j \leq |y|}\left(\frac{1}{x - y} - \frac{1}{-y}\right)f(y)\dd y\right| \leq \| f \|_{L^{\infty}}\int_{j \leq |y|}\frac{2|x|}{|y|^{2}}\dd y \leq \frac{4A}{j}\| f \|_{L^{\infty}},
\end{equation*}
故当 $j \to \infty$ 时,
\begin{equation*}
g_{j}(x) = H\left(f\chi_{\{|x| < j\}}\right)(x) - \frac{1}{\pi}\int_{1 < |y| < j}\frac{f(y)}{-y}\dd y
\end{equation*}
在 $[-A, A]$ 上一致收敛到
\begin{equation*}
H\left(f\chi_{\{|x| < j_{0}\}}\right)(x) - \frac{1}{\pi}\int_{1 < |y| < j_{0}}\frac{f(y)}{-y}\dd y + \frac{1}{\pi}\int_{j_{0} \leq |y|}\left(\frac{1}{x - y} - \frac{1}{-y}\right)f(y)\dd y,
\end{equation*}
因此依 $L^{1}([-A, A])$ 范数收敛.

对区间 $Q$,$g_{j}$ 在 $Q$ 上依 $L^{1}$ 范数收敛到 $\tilde{H}f$,故也利用定理~\ref{thm:1.13.3},
\begin{align*}
\dfrac{1}{|Q|}\displaystyle\int_{Q}\left|\tilde{H}f(x) - \frac{1}{|Q|}\int_{Q}\tilde{H}f(y)\dd y\right|\dd x&= \displaystyle\lim_{j \to \infty}\frac{1}{|Q|}\int_{Q}\left|g_{j}(x) - \frac{1}{|Q|}\int_{Q}g_{j}(y)\dd y\right|\dd x  \\
&= \displaystyle\lim_{j \to \infty}\frac{1}{|Q|}\int_{Q}\left|H\left(f\chi_{\{|x| < j\}}\right)(x) - \frac{1}{|Q|}\int_{Q}H\left(f\chi_{\{|x| < j\}}\right)(y)\dd y\right|\dd x \\
&\leq \displaystyle\sup_{j}\| H\left(f\chi_{\{|x| < j\}}\right) \|_{\mathrm{BMO}} \leq C\| f \|_{L^{\infty}}.
\end{align*}
由此得
\begin{equation*}
\| \tilde{H}f \|_{\mathrm{BMO}} \leq C\| f \|_{L^{\infty}}.
\end{equation*}
\end{proof}

其次,考察 $\mathrm{BMO}(\mathbb{R})$ 与 $H^{1}(\mathbb{R})$ 的关系.

\begin{theorem}\label{thm:1.13.5}
若 $T$ 是 $H^{1}(\mathbb{R})$ 上的有界线性泛函,则存在有界函数 $g_{1}, g_{2}$ 使
\begin{equation*}
T(f) = \int_{\mathbb{R}}\left(f(x)g_{1}(x) + Hf(x)g_{2}(x)\right)\dd x \quad (f \in H^{1}(\mathbb{R})).
\end{equation*}
此时,$\| g_{1} \|_{L^{\infty}}, \| g_{2} \|_{L^{\infty}} \leq \| T \|_{\mathrm{op}}$.其中 $\| T \|_{\mathrm{op}} = \inf\{C > 0;\ |T(f)| \leq C\| f \|_{H^{1}}\}$.

反之,若 $g_{1}, g_{2} \in L^{\infty}(\mathbb{R})$,则
\begin{align*}
\left|\displaystyle\int_{\mathbb{R}}f(x)\left(g_{1}(x) + \tilde{H}g_{2}(x)\right)\dd x\right| & \leq \max\left(\| g_{1} \|_{L^{\infty}}, \| g_{2} \|_{L^{\infty}}\right)\| f \|_{H^{1}} \\
& \quad (f \in H_{00}^{1}(\mathbb{R})).
\end{align*}
\end{theorem}

\begin{proof}
在 $L^{1}(\mathbb{R}) \times L^{1}(\mathbb{R})$ 上赋予范数 $\| (f_{1}, f_{2}) \| = \| f_{1} \|_{L^{1}} + \| f_{2} \|_{L^{1}}$ 的空间记为 $L^{1}(\mathbb{R}) \oplus L^{1}(\mathbb{R})$,则此空间是 Banach 空间,且 $B := \{(f, Hf) \in L^{1}(\mathbb{R}) \times L^{1}(\mathbb{R});\ f \in H^{1}(\mathbb{R})\}$ 是 $L^{1}(\mathbb{R}) \oplus L^{1}(\mathbb{R})$ 的闭子空间.$H^{1}(\mathbb{R})$ 作为 Banach 空间可与 $B$ 视为同一.若 $|T(f)| \leq \| T \|_{\mathrm{op}}\| f \|_{H^{1}}$,则 $|T(f)| \leq \| T \|_{\mathrm{op}}\left(\| f \|_{L^{1}} + \| Hf \|_{L^{1}}\right)$,故 $T$ 是 $B$ 上的有界线性泛函.因此由 Hahn--Banach 定理可扩张为 $L^{1}(\mathbb{R}) \oplus L^{1}(\mathbb{R})$ 上的有界线性泛函.因 $L^{1}(\mathbb{R}) \oplus L^{1}(\mathbb{R})$ 的对偶空间是 $L^{\infty}(\mathbb{R}) \oplus L^{\infty}(\mathbb{R})$,故存在有界函数 $g_{1}, g_{2}$ 满足
\begin{equation*}
\begin{aligned}
T(f) &= \displaystyle\int_{\mathbb{R}}\left(f(x)g_{1}(x) + Hf(x)g_{2}(x)\right)\dd x, \\
\| g_{1} \|_{L^{\infty}}, \| g_{2} \|_{L^{\infty}} &\leq \| T \|_{\mathrm{op}}.
\end{aligned}
\end{equation*}

反之,设 $g_{1}, g_{2}$ 为有界函数,$f \in H_{00}^{1}(\mathbb{R})$.因 $f \in \mathcal{S}(\mathbb{R})$,$|f(x)| \leq C(1 + |x|)^{-2}$.另外,$\int_{\mathbb{R}}f(x)\dd x = 0$.

若能证明对任意有界函数 $g$,
\begin{equation}\tag{1.13.1}
\int_{\mathbb{R}}f(x)\tilde{H}g(x)\dd x = -\int_{\mathbb{R}}Hf(x)g(x)\dd x,
\label{eq:1.13.1}\end{equation}
则
\begin{align*}
\left|\displaystyle\int_{\mathbb{R}}f(x)\left(g_{1}(x) + \tilde{H}g_{2}(x)\right)\dd x\right| &= \left|\int_{\mathbb{R}}\left(f(x)g_{1}(x) - Hf(x)g_{2}(x)\right)\dd x\right| \\
&\leq \left(\| f \|_{L^{1}}\| g_{1} \|_{L^{\infty}} + \| Hf \|_{L^{1}}\| g_{2} \|_{L^{\infty}}\right) \leq \max\left(\| g_{1} \|_{L^{\infty}}, \| g_{2} \|_{L^{\infty}}\right)\| f \|_{H^{1}}.
\end{align*}
于是来证明 \eqref{eq:1.13.1}.令 $\chi_{a} = \chi_{\{x \in \mathbb{R};\ |x| \leq a\}}$.依次注意 $Hf \in L^{1}(\mathbb{R})$,$g \in L^{\infty}(\mathbb{R})$,$f, g\chi_{2^{k}} \in L^{2}(\mathbb{R})$,$\int_{\mathbb{R}}f(x)\dd x = 0$,
\begin{align*}
-\displaystyle\int_{\mathbb{R}}Hf(x)g(x)\dd x &= -\lim_{k \to \infty}\int_{\mathbb{R}}Hf(x)g(x)\chi_{2^{k}}(x)\dd x \\
&= \displaystyle\lim_{k \to \infty}\int_{\mathbb{R}}f(x)H(g\chi_{2^{k}})(x)\dd x \\
&= \displaystyle\lim_{k \to \infty}\int_{\mathbb{R}}f(x)\left(H(g\chi_{2^{k}})(x) + \frac{1}{\pi}\int_{1 \leq |y| \leq 2^{k}}\frac{g(y)}{y}\dd y\right)\dd x
\end{align*}
\begin{align*}
&= \displaystyle\lim_{k \to \infty}\int_{|x| \leq 2}f(x)\left(H(g\chi_{2^{k}})(x) + \frac{1}{\pi}\int_{1 \leq |y| \leq 2^{k}}\frac{g(y)}{y}\dd y\right)\dd x \\
&+ \displaystyle\lim_{k \to \infty}\sum_{j=1}^{\infty}\int_{2^{j} < |x| \leq 2^{j+1}}f(x)\left(H(g\chi_{2^{k}})(x) + \frac{1}{\pi}\int_{1 \leq |y| \leq 2^{k}}\frac{g(y)}{y}\dd y\right)\dd x.
\end{align*}
令
\begin{equation*}
a_{j}(k) = \int_{2^{j} < |x| \leq 2^{j+1}}f(x)\left(H(g\chi_{2^{k}})(x) + \frac{1}{\pi}\int_{1 \leq |y| \leq 2^{k}}\frac{g(y)}{y}\dd y\right)\dd x,
\end{equation*}
则

(i) 当 $k \leq j + 2$ 时,
\begin{align*}
a_{j}(k)&\leq \displaystyle\sup_{2^{j} < |x| \leq 2^{j+1}}|f(x)|\int_{2^{j} < |x| \leq 2^{j+1}}\left(|H(g\chi_{2^{k}})(x)| + \frac{1}{\pi}\int_{1 \leq |y| \leq 2^{k}}\frac{|g(y)|}{|y|}\dd y\right)\dd x  \\
&\leq C2^{-2j}\left\{\left(\displaystyle\int_{2^{j} < |x| \leq 2^{j+1}}\dd x\right)^{1/2}\left(\int_{2^{j} < |x| \leq 2^{j+1}}|H(g\chi_{2^{k}})(x)|^{2}\dd x\right)^{1/2}\right. \\
&\quad + \left.2^{j+2}k\dfrac{\log 2}{\pi}\| g \|_{L^{\infty}}\right\} \\
&\leq C2^{-2j}\left\{2^{(j+1)/2}\| g\chi_{2^{k}} \|_{L^{2}} + 2^{j+3}(j + 2)\dfrac{\log 2}{\pi}\| g \|_{L^{\infty}}\right\} \\
&\leq Cj2^{-j}\| g \|_{L^{\infty}}.
\end{align*}

(ii) 当 $k > j + 2$ 时,
\begin{align*}
H(g\chi_{2^{k}})(x) + \dfrac{1}{\pi}\displaystyle\int_{1 \leq |y| \leq 2^{k}}\frac{g(y)}{y}\dd y & = H(g\chi_{2^{j+2}})(x) + \dfrac{1}{\pi}\int_{1 \leq |y| \leq 2^{j+2}}\frac{g(y)}{y}\dd y \\
& \quad + \dfrac{1}{\pi}\displaystyle\int_{2^{j+2} \leq |y| \leq 2^{k}}\left(\frac{1}{x - y} + \frac{1}{y}\right)g(y)\dd y,
\end{align*}
且当 $2^{j} < |x| \leq 2^{j+1}$ 时,
\begin{align*}
\left|\displaystyle\int_{2^{j+2} \leq |y| \leq 2^{k}}\left(\frac{1}{x - y} + \frac{1}{y}\right)g(y)\dd y\right| & \leq \int_{2^{j+2} \leq |y| \leq 2^{k}}\frac{|x|}{|x - y||y|}|g(y)|\dd y \\
& \leq \displaystyle\int_{2^{j+2} \leq |y| \leq 2^{k}}\frac{2|x|}{|y|^{2}}|g(y)|\dd y \leq C\| g \|_{L^{\infty}},
\end{align*}
故把 (i) 中 $j + 2 = k$ 的情形也用于此,
\begin{equation*}
a_{j}(k) \leq Cj2^{-j}\| g \|_{L^{\infty}}.
\end{equation*}
(i),(ii) 合并,$a_{j}(k) \leq Cj2^{-j}\| g \|_{L^{\infty}}$($j = 1, 2, \ldots$),故
\begin{equation*}
\lim_{k \to \infty}\sum_{j=1}^{\infty}a_{j}(k) = \sum_{j=1}^{\infty}\lim_{k \to \infty}a_{j}(k).
\end{equation*}
因此,
\begin{align*}
-\displaystyle\int_{\mathbb{R}}Hf(x)g(x)\dd x&= \displaystyle\int_{|x| \leq 2}f(x)\left(\lim_{k \to \infty}H(g\chi_{2^{k}})(x) + \frac{1}{\pi}\int_{1 \leq |y| \leq 2^{k}}\frac{g(y)}{y}\dd y\right)\dd x  \\
&\quad + \displaystyle\sum_{j=1}^{\infty}\int_{2^{j} < |x| \leq 2^{j+1}}f(x)\lim_{k \to \infty}\left(H(g\chi_{2^{k}})(x) + \frac{1}{\pi}\int_{1 \leq |y| \leq 2^{k}}\frac{g(y)}{y}\dd y\right)\dd x \\
&= \displaystyle\int_{\mathbb{R}}f(x)\lim_{k \to \infty}\left(H(g\chi_{2^{k}})(x) + \frac{1}{\pi}\int_{1 \leq |y| \leq 2^{k}}\frac{g(y)}{y}\dd y\right)\dd x \\
&= \displaystyle\int_{\mathbb{R}}f(x)\tilde{H}g(x)\dd x.
\end{align*}
于是式~\eqref{eq:1.13.1} 得证.
\end{proof}

再进一步,$\mathrm{BMO}(\mathbb{R})$ 函数具有 $g_{1} + \tilde{H}g_{2}$($g_{1}, g_{2} \in L^{\infty}(\mathbb{R})$)的形式,但因篇幅关系,这里不加证明地列出下面的定理.

\begin{theorem}\label{thm:1.13.6}
$H^{1}(\mathbb{R})$ 的对偶空间在下面的意义下是 $\mathrm{BMO}(\mathbb{R})$.

(a) 若 $g \in \mathrm{BMO}(\mathbb{R})$,则对 $f \in H_{00}^{1}(\mathbb{R})$,
\begin{equation}\tag{1.13.2}
\ell(f) := \int_{\mathbb{R}}f(x)g(x)\dd x
\label{eq:1.13.2}\end{equation}
存在,且
\begin{equation*}
\left|\int_{\mathbb{R}}f(x)g(x)\dd x\right| \leq \| f \|_{H^{1}}\| g \|_{\mathrm{BMO}},
\end{equation*}
$\ell$ 可作为 $H^{1}(\mathbb{R})$ 上的有界线性泛函唯一扩张.

(b) 反之,若 $\ell$ 是 $H^{1}(\mathbb{R})$ 上的有界线性泛函,则存在对 $f \in H_{00}^{1}(\mathbb{R})$ 满足式~\eqref{eq:1.13.2} 的 $\mathrm{BMO}(\mathbb{R})$ 函数 $g$,且除常数函数之差外唯一确定.这个 $g$ 与定理~\ref{thm:1.13.5} 所得者除常数之差外一致.
\end{theorem}

下面不加证明地列出关于 $\mathrm{BMO}(\mathbb{R})$ 的基本事实.

\begin{lemma}\label{lem:1.13.7}
(i) 若 $f \in L^{\infty}(\mathbb{R})$,则 $\| f \|_{\mathrm{BMO}} \leq 2\| f \|_{L^{\infty}}$.

(ii) 若 $f \in \mathrm{BMO}(\mathbb{R})$,$\alpha > 1$,则
\begin{equation*}
\int_{\mathbb{R}}\frac{|f(x)|}{(1 + |x|)^{\alpha}}\dd x \leq C_{\alpha}\| f \|_{\mathrm{BMO}}.
\end{equation*}

(iii) 若 $|F(x) - F(y)| \leq C_{0}|x - y|$($x, y \in \mathbb{R}$),则对实值 BMO 函数 $f$,
\begin{equation*}
\| F(f) \|_{\mathrm{BMO}} \leq 2C_{0}\| f \|_{\mathrm{BMO}}.
\end{equation*}
由此,若 $f, g \in \mathrm{BMO}(\mathbb{R})$ 且为实值,则
\begin{align*}
\| \max(f, g) \|_{\mathrm{BMO}} &\leq 2\max\left(\| f \|_{\mathrm{BMO}}, \| g \|_{\mathrm{BMO}}\right), \\
\| \min(f, g) \|_{\mathrm{BMO}} &\leq 2\max\left(\| f \|_{\mathrm{BMO}}, \| g \|_{\mathrm{BMO}}\right).
\end{align*}
\end{lemma}

\begin{theorem}[John--Nirenberg]\label{thm:1.13.8}
若 $f \in \mathrm{BMO}(\mathbb{R})$,则对任意区间 $Q$,存在正定数 $C_{0}, C_{1}$ 使
\begin{equation*}
\frac{\left|\{x \in Q;\ |f(x) - f_{Q}| > \lambda\}\right|}{|Q|} \leq C_{0}e^{-C_{1}\lambda/\| f \|_{\mathrm{BMO}}} \quad (\lambda > 0).
\end{equation*}
它们不依赖于 $f, \lambda, Q$.
\end{theorem}

由此,

\begin{lemma}\label{lem:1.13.9}
设 $1 < p < \infty$.$f \in \mathrm{BMO}(\mathbb{R})$ 的充要条件是 $f \in L_{\mathrm{loc}}^{p}(\mathbb{R})$ 且
\begin{equation*}
\| f \|_{\mathrm{BMO}, p} := \sup_{Q:\ \text{区间}}\left(\frac{1}{|Q|}\int_{Q}|f(x) - f_{Q}|^{p}\dd x\right)^{1/p} < \infty.
\end{equation*}
此时,存在正定数 $C_{p}$ 使 $\| f \|_{\mathrm{BMO}} \leq \| f \|_{\mathrm{BMO}, p} \leq C_{p}\| f \|_{\mathrm{BMO}}$.
\end{lemma}

关于在 $\mathrm{BMO}(\mathbb{R})$ 中弱意义下稠密的好的空间,

\begin{lemma}\label{lem:1.13.10}
设 $f \in \mathrm{BMO}(\mathbb{R})$,则存在函数列 $f_{n} \in C_{c}^{\infty}(\mathbb{R})$ 满足
\begin{equation*}
\begin{aligned}
\| f_{n} \|_{\mathrm{BMO}} &\leq 2\| f \|_{\mathrm{BMO}}, \\
\lim_{n \to \infty}\displaystyle\int_{\mathbb{R}}f_{n}(x)\varphi(x)\dd x &= \int_{\mathbb{R}}f(x)\varphi(x)\dd x \quad (\varphi \in \mathcal{S}(\mathbb{R})).
\end{aligned}
\end{equation*}
\end{lemma}


% ============ 第1章 参考文献(原书 pp.122-123) ============
\section*{参考文献}\addcontentsline{toc}{section}{参考文献}
\begin{thebibliography}{99}
% 日文条目使用日文字体(Fandol 字体无日文假名/和制汉字字形;用 TeX Live 自带的 Harano Aji Mincho)
\newfontfamily\jpref{Harano Aji Mincho}
\item[] 这里,列举本章参考的文献以及进一步深化本章内容的文献.
\item[] 单行本有:
\bibitem[1]{ch1:1} {\jpref 猪狩惺, 実解析入門, 岩波書店, 東京, 1996.}
\bibitem[2]{ch1:2} {\jpref 伊藤清三, ルベーグ積分入門, 裳華房, 東京, 1963.}
\bibitem[3]{ch1:3} R. R. Coifman and Y. Meyer, Au-delà des opérateurs pseudodifférentiels, Astérisque \#57, Société Math. France, 1978.
\bibitem[4]{ch1:4} J. García-Cuerva and J. L. Rubio de Francia, Weighted Norm Inequalities and Related Topics, North-Holland Math. Studies 116, Amsterdam, 1985.
\bibitem[5]{ch1:5} L. Grafakos, Classical and Modern Fourier Analysis, Pearson Education, NJ, 2004.
\bibitem[6]{ch1:6} J.-L. Journé, Calderón--Zygmund Operators, Pseudo-Differential Operators and the Cauchy Integral of Calderón, Lecture Notes in Math. \#994, Springer-Verlag, Berlin, 1983.
\bibitem[7]{ch1:7} E. M. Stein, Singular Integrals and Differentiability Properties of Functions, Princeton University Press, Princeton, NJ, 1970.
\bibitem[8]{ch1:8} E. M. Stein and G. Weiss, Introduction to Fourier Analysis on Euclidean Spaces, Princeton University Press, Princeton, NJ, 1971.
\bibitem[9]{ch1:9} E. M. Stein, Harmonic Analysis: Real Variable Methods, Orthogonality, and Oscillatory Integrals, Princeton University Press, Princeton, NJ, 1993.
\bibitem[10]{ch1:10} A. Torchinsky, Real-Variable Methods in Harmonic Analysis, Academic Press, New York, 1986.
\bibitem[11]{ch1:11} A. Uchiyama, Hardy Spaces on the Euclidean Space, Springer-Verlag, Tokyo, 2001.
\bibitem[12]{ch1:12} A. Zygmund, Trigonometric Series, I and II, 2nd Edition, Cambridge University Press, Cambridge, UK, 1959.
\item[] 第 1.1 节用文献 \cite{ch1:1}, \cite{ch1:2}, \cite{ch1:3}.第 1.2,1.3 节用文献 \cite{ch1:8}.第 1.4～1.8 节用文献 \cite{ch1:4}, \cite{ch1:5}, \cite{ch1:6}, \cite{ch1:7}, \cite{ch1:10}.第 1.8 节中 Marcinkiewicz 插值定理的证明参考了文献 \cite{ch1:12}.第 1.9 节用文献 \cite{ch1:3},第 1.10 节用文献 \cite{ch1:5},第 1.11 节用文献 \cite{ch1:4}, \cite{ch1:5}.关于 $H^{1}(\mathbb{R})$,$\mathrm{BMO}(\mathbb{R})$,与论文 C. Fefferman and E. M. Stein, $H^{p}$ spaces of several variables, Acta Math. 129 (1972), 137-193. 一起参考了文献 \cite{ch1:11}.
\item[] 想在实分析领域进一步扩展知识的人,不妨翻阅文献 \cite{ch1:5}, \cite{ch1:7}, \cite{ch1:9}.关于 Hardy 空间,BMO 空间,文献 \cite{ch1:11} 很详细.
\end{thebibliography}

% ============ 第2章 用复函数论与泛函分析方法研究Hardy空间的理论（中路貴彦，原书 pp.124-202 = PDF pp.136-214） ============
\chapter{用复函数论与泛函分析方法研究 Hardy 空间的理论（中路貴彦）}\label{ch:2}

% ============ 第2章 §2.1 Hardy空间的定义(原书 pp.124-126) ============
\section{Hardy 空间的定义}\label{sec:2.1}
\index{Hardy 空间!(单位圆板上的)}


$\mathbb{C}$ 为复平面,$D = \{z \in \mathbb{C};\ |z| < 1\}$,$\partial D = \{z \in \mathbb{C};\ |z| = 1\}$,且 $\overline{D} = \{z \in \mathbb{C};\ |z| \leq 1\}$.$h(D)$ 表示 $D$ 上实值调和函数的全体,$H(D)$ 表示 $D$ 上正则函数的全体.对 $f \in H(D)$,$\operatorname{Re}(f)$ 表示 $f$ 的实部,$\operatorname{Im}(f)$ 表示 $f$ 的虚部.令 $\operatorname{Re} H(D) = \{\operatorname{Re}(f);\ f \in H(D)\}$,则因 $D$ 是单连通的,$\operatorname{Re} H(D) = h(D)$ 成立.又 $\operatorname{Im} H(D) = \{\operatorname{Im}(f);\ f \in H(D)\} = h(D)$.为研究属于 $h(D)$ 与 $H(D)$ 的函数,可以考察多种度量函数在 $D$ 的边界附近绝对值的增长度的量.特别地,当 $F$ 在 $D$ 上连续时,
\begin{equation*}
M_{p}(r, F) = \left\{\frac{1}{2\pi}\int_{-\pi}^{\pi}|F(re^{i\theta})|^{p}\dd\theta\right\}^{1/p} \quad (0 < p < \infty)
\end{equation*}
与
\begin{equation*}
M_{\infty}(r, F) = \max_{-\pi \leq \theta < \pi}|F(re^{i\theta})|
\end{equation*}
是应用上重要的度量 $F$ 的增长度的量.当然,当 $F$ 在 $\overline{D} = \{z \in \mathbb{C};\ |z| \leq 1\}$ 上连续时,对任意 $p$,$\sup_{0 \leq r < 1}M_{p}(r, F) < \infty$.这里对实值调和函数与正则函数定义 Hardy 空间 $h^{p}$ 与 $H^{p}$.对 $0 < p \leq \infty$,
\begin{equation*}
h^{p} = \left\{u \in h(D);\ \sup_{0 \leq r < 1}M_{p}(r, u) < \infty\right\}
\end{equation*}
且
\begin{equation*}
H^{p} = \left\{f \in H(D);\ \sup_{0 \leq r < 1}M_{p}(r, f) < \infty\right\}
\end{equation*}
由定义,$h^{\infty} \subset h^{p} \subset h^{q}$($p > q$)且 $H^{\infty} \subset H^{p} \subset H^{q}$($p > q$)成立.$h^{\infty}$ 与 $H^{\infty}$ 显然是线性空间.当 $0 < p < \infty$ 时,
\begin{equation*}
a^{p} \leq (a + b)^{p} \leq 2^{p}(a^{p} + b^{p}) \quad (a \geq 0, b \geq 0)
\end{equation*}
故 $h^{p}$ 与 $H^{p}$ 是线性空间.虽然 $\operatorname{Re} H(D) = h(D)$,但 $\operatorname{Re} H^{p} = h^{p}$ 是否成立,对一般的 $p$ 是难以证明的.多项式 $\sum_{n=0}^{\ell}a_{n}z^{n}$ 属于 $H^{\infty}$,$\sum_{n=0}^{\ell}a_{n}z^{n} + \sum_{n=0}^{\ell}\overline{a_{n}z^{n}}$ 属于 $h^{\infty}$.当 $0 < p < 1$ 时 $H^{p}$ 作为空间是很大的,$(1-z)^{-1}$ 不属于 $H^{1}$,但可以证明 $(1-z)^{-1} \in \bigcap_{0 < p < 1}H^{p}$.另外,$\operatorname{Re}(1-z)^{-1}$ 与 $\operatorname{Re}(1+z)(1-z)^{-1}$ 属于 $h^{1}$ 这一点将在后面证明.

下一个引理中的不等式在 Hardy 空间 $H^{p}$ 的研究中将在后面用到.

\begin{lemma}\label{lem:2.1.1}
设 $0 < p \leq \infty$.若 $f \in H(D)$ 且 $0 < r < r' < 1$,则 $M_{p}(r, f) \leq M_{p}(r', f)$.
\end{lemma}

\begin{proof}
设 $p \neq \infty$.取 $u(z)$ 为在 $|z| < r'$ 上调和,在 $|z| \leq r'$ 上连续,且在 $|z| = r'$ 上等于 $|f(z)|^{p}$ 的函数.由于 $|f|^{p}$ 是次调和函数,在 $|z| \leq r'$ 上 $|f(z)|^{p} \leq u(z)$.因此,利用调和函数的 Gauss 平均值定理,
\begin{align*}
M_{p}(r, f) & \leq \frac{1}{2\pi}\int_{-\pi}^{\pi}u(re^{i\theta})\dd\theta = u(0) \\
& = \frac{1}{2\pi}\int_{-\pi}^{\pi}u(r'e^{i\theta})\dd\theta = M_{p}(r', f).
\end{align*}
当 $p = \infty$ 时,由正则函数的最大值原理是显然的.
\end{proof}

% ============ 第2章 §2.2 Poisson核与Cauchy核(原书 pp.126-130) ============
\section{Poisson 核与 Cauchy 核}\label{sec:2.2}
\index{Poisson 核!(单位圆板的)}
\index{Cauchy 核}
\index{Poisson 积分}
\index{Cauchy 积分}
\index{Cauchy--Stieltjes 积分}
\index{Poisson--Stieltjes 积分}
\index{Herglotz 积分}
\index{Jensen 公式}
\index{零点集}


2.1 节中出现的 $(1-z)^{-1}$ 与 $\operatorname{Re}\left\{(1+z)(1-z)^{-1}\right\}$ 是非常重要的函数.$(1-z)^{-1}$ 称为 Cauchy 核,$\operatorname{Re}\left\{(1+z)(1-z)^{-1}\right\}$ 称为 Poisson 核.这两个核以如下方式相联系.对 $z \in D$,设 $z^{*}$ 为关于 $\partial D$ 的镜像,则 $z^{*} = z/|z|^{2}$,故
\begin{equation*}
\frac{1}{1 - z} - \frac{1}{1 - z^{*}} = \frac{1 - |z|^{2}}{|1 - z|^{2}} = \operatorname{Re}\left(\frac{1 + z}{1 - z}\right)
\end{equation*}
成立.用极坐标表示,设 $z = re^{i\theta}$($r \neq 1$),记
\begin{equation*}
\frac{1}{1 - re^{i\theta}} = C(r, \theta)
\end{equation*}
且
\begin{equation*}
\operatorname{Re}\left(\frac{1 + re^{i\theta}}{1 - re^{i\theta}}\right) = \frac{1 - r^{2}}{1 - 2r\cos\theta + r^{2}} = P(r, \theta).
\end{equation*}
若 $f$ 在包含 $\overline{D}$ 的开圆盘上正则,则由 Cauchy 积分表示与 Cauchy 积分定理
\begin{align*}
f(z) & = \frac{1}{2\pi i}\int_{|\zeta| = 1}\frac{f(\zeta)}{\zeta - z}\dd\zeta \\
& = \frac{1}{2\pi i}\int_{|\zeta| = 1}f(\zeta)\left(\frac{1}{\zeta - z} - \frac{1}{\zeta - z^{*}}\right)\dd\zeta
\end{align*}
因此
\begin{align*}
f(re^{i\theta}) & = \frac{1}{2\pi}\int_{-\pi}^{\pi}f(e^{it})C(r, \theta - t)\dd t \\
& = \frac{1}{2\pi}\int_{-\pi}^{\pi}f(e^{it})P(r, \theta - t)\dd t
\end{align*}
成立.

$M = M[-\pi, \pi]$ 表示 $[-\pi, \pi]$ 上(有界)复 Borel 测度的全体,对 $\mu \in M$ 令 $\|\mu\| = |\mu|([-\pi, \pi])$.对 $0 < p \leq \infty$,用 $L^{p} = L^{p}[-\pi, \pi]$ 表示 $[-\pi, \pi]$ 上 $p$ 次 Lebesgue 可积函数的全体,则 $L^{1}$ 可以等距嵌入 $M$.事实上,对 $F \in L^{1}$,令 $\mu = F(t)\dd t$,则 $\mu \in M$ 且 $\|\mu\| = \int_{-\pi}^{\pi}|F(t)|\dd t = \|F\|_{1}$.对 $\mu \in M$,
\begin{equation*}
f(re^{i\theta}) = \frac{1}{2\pi}\int_{-\pi}^{\pi}C(r, \theta - t)\dd\mu(t)
\end{equation*}
称为 Cauchy--Stieltjes 积分,
\begin{equation*}
u(re^{i\theta}) = \frac{1}{2\pi}\int_{-\pi}^{\pi}P(r, \theta - t)\dd\mu(t)
\end{equation*}
称为 Poisson--Stieltjes 积分.当 $\mu = F(t)\dd t$ 时,分别称为 Cauchy 积分,Poisson 积分.当 $\mu$ 是 Dirac 测度 $\delta_{t=0}$ 时,其 Cauchy--Stieltjes 积分是 $C(r, \theta)$,Poisson--Stieltjes 积分是 $P(r, \theta)$.

由于 $C(r, \theta - t) = e^{it}/(e^{it} - re^{i\theta})$ 关于 $z = re^{i\theta}$ 正则,Cauchy--Stieltjes 积分属于 $H(D)$.

由于 $P(r, \theta - t) = \operatorname{Re}(e^{it} + re^{i\theta})/(e^{it} - re^{i\theta})$ 关于 $z = re^{i\theta}$ 调和,当 $\mu$ 是实测度时 Poisson--Stieltjes 积分属于 $h(D)$.事实上,Poisson--Stieltjes 积分属于 $h^{1}$ 且 $\mu$ 的 Cauchy--Stieltjes 积分属于 $\bigcap_{0 < p < 1}H^{p}$ 这一点将在 2.3 节证明.

正如 2.1 节所注意到的,对 $u \in h(D)$,存在 $f \in H(D)$ 使得 $\operatorname{Re}(f) = u$.这可以由以下事实证明:当 $u$ 在包含 $\overline{D}$ 的开集上调和时,令
\begin{equation*}
f(z) = \frac{1}{2\pi}\int_{-\pi}^{\pi}\frac{e^{it} + z}{e^{it} - z}u(e^{it})\dd t \quad (z \in D),
\end{equation*}
则 $f$ 在 $D$ 上正则,两边取实部,
\begin{equation*}
\operatorname{Re}\left\{f(z)\right\} = \frac{1}{2\pi}\int_{-\pi}^{\pi}P(r, \theta - t)u(e^{it})\dd t \quad (z \in D)
\end{equation*}
且 $\operatorname{Re}\left\{f(z)\right\} = u(z)$($z \in D$).因此,当 $f$ 在 $\overline{D}$ 上正则且无零点时,$\operatorname{Re}\left\{\log f(z)\right\} = \log|f(z)|$($z \in \overline{D}$),故
\begin{equation*}
\log|f(z)| = \frac{1}{2\pi}\int_{-\pi}^{\pi}P(r, \theta - t)\log|f(e^{it})|\dd t \quad (z \in D).
\end{equation*}
令 $z = 0$,
\begin{equation*}
\log|f(0)| = \frac{1}{2\pi}\int_{-\pi}^{\pi}\log|f(e^{it})|\dd t.
\end{equation*}
将这个结果推广到 $f$ 在 $D$ 内有零点的情况.这称为 Jensen 公式.

设 $\{a_{j}\}$ 为 $f$ 在 $D$ 中的零点集.若为零点阶数为 $\ell$,则连续写出 $\ell$ 次.当在 $\overline{D}$ 上正则的 $f(z)$ 在 $z = a$ 处为零点时,除去该零点时不用 $f(z)/(z - a)$,而令
\begin{equation*}
h(z) = f(z)\Big/\frac{z - a}{1 - \bar{a}z},
\end{equation*}
则 $|h(e^{i\theta})| = |f(e^{i\theta})|$.这个想法经常被使用.

\begin{proposition}\label{prop:2.2.1}
$f$ 在 $\overline{D}$ 上正则,$f(0) \neq 0$ 且在 $\partial D$ 上无零点.设 $\{a_{j}\}$ 为 $f$ 在 $D$ 内的零点集,则 $\{a_{j}\}$ 为有限 $n$ 个,且
\begin{equation*}
\log|f(0)| = \frac{1}{2\pi}\int_{-\pi}^{\pi}\log|f(e^{i\theta})|\dd\theta + \sum_{j = 1}^{n}\log|a_{j}|.
\end{equation*}
\end{proposition}

\begin{proof}
由于 $f$ 在 $\overline{D}$ 上正则,由一致定理,$D$ 中的零点有限.除去 $f$ 在 $D$ 中的零点,令
\begin{equation*}
h(z) = f(z)\Bigg/\prod_{j = 1}^{n}\frac{z - a_{j}}{1 - \bar{a}_{j}z},
\end{equation*}
则如上所注意的,
\begin{equation*}
\log|h(0)| = \frac{1}{2\pi}\int_{-\pi}^{\pi}\log|h(e^{i\theta})|\dd\theta,
\end{equation*}
由于 $h(0) = f(0)/\prod_{j=1}^{n}(-a_{j})$ 且 $|h(e^{i\theta})| = |f(e^{i\theta})|$,本命题得证.
\end{proof}

\begin{proposition}\label{prop:2.2.2}
Poisson 核与 Cauchy 核
\begin{equation*}
P(r, \theta) = \frac{1 - r^{2}}{1 - 2r\cos\theta + r^{2}}, \quad C(r, \theta) = \frac{1}{1 - re^{i\theta}}
\end{equation*}
具有如下性质:

(1) $P(r, \theta) > 0$($0 \leq r < 1$),

(2) $P(r, \theta + 2\pi) = P(r, \theta)$($0 \leq r < 1$),

(3) $\frac{1}{2\pi}\int_{-\pi}^{\pi}P(r, \theta)\dd\theta = 1$($0 \leq r < 1$),

(4) $\operatorname{Re}\left\{C(r, \theta)\right\} = \frac{1 - r\cos\theta}{1 - 2r\cos\theta + r^{2}} > 0$,

(5) $C(r, \theta) - C\left(\frac{1}{r}, \theta\right) = P(r, \theta)$.
\end{proposition}

\begin{proof}
除 (3) 外都显然.(3) 对 $\operatorname{Re}\left(\frac{1 + z}{1 - z}\right)$ 用 Gauss 平均值定理来证明.
\end{proof}

命题~\ref{prop:2.2.1} 若把 $\overline{D}$ 换成 $\{z \in \mathbb{C};\ |z| \leq r\}$,则稍微改变形式后仍然成立,后面将多次用到.命题~\ref{prop:2.2.2} 是 Poisson 核的重要性质,后面将用到.

对 $\mu \in M[-\pi, \pi]$,
\begin{equation*}
f(z) = \frac{1}{2\pi}\int_{-\pi}^{\pi}\frac{e^{it} + z}{e^{it} - z}\dd\mu(t) \quad (z \in D)
\end{equation*}
称为 Herglotz 积分.当 $\mu$ 是实测度时,其实部是 Poisson--Stieltjes 积分.利用这个关系,证明 Poisson--Stieltjes 积分的唯一性.

\begin{proposition}\label{prop:2.2.3}
$\mu \in M[-\pi, \pi]$ 是实测度且满足 $\mu(\{-\pi\}) = \mu(\{\pi\})$.若 $\mu$ 的 Poisson--Stieltjes 积分在 $D$ 上为零,则 $\mu \equiv 0$.
\end{proposition}

\begin{proof}
由上面注意到的 Herglotz 积分与 Poisson 积分的关系,若
\begin{equation*}
u(z) = \int_{-\pi}^{\pi}P(r, \theta - t)\dd\mu(t) = 0 \quad (z \in D),
\end{equation*}
则存在实常数 $\gamma$ 使得
\begin{equation*}
\int_{-\pi}^{\pi}\frac{e^{it} + z}{e^{it} - z}\dd\mu(t) = i\gamma \quad (z \in D).
\end{equation*}
由于
\begin{equation*}
\frac{e^{it} + z}{e^{it} - z} = 1 + 2\sum_{n = 1}^{\infty}e^{-int}z^{n},
\end{equation*}
故 $\int_{-\pi}^{\pi}e^{-int}\dd\mu(t) = 0$($n = 0, \pm 1, \pm 2, \ldots$),由 Fourier 系数的唯一性,$\mu \equiv 0$.
\end{proof}

以后,当涉及 Poisson 积分的唯一性时,设 $[-\pi, \pi]$ 上的测度满足 $\mu(\{-\pi\}) = \mu(\{\pi\})$.另外,同样地,$\mu$ 也可以看作 $\partial D$ 上的测度.

% ============ 第2章 §2.3 径向极限与Fatou定理(原书 pp.130-132) ============
\section{径向极限与 Fatou 定理}\label{sec:2.3}
\index{Fatou 定理}
\index{径向极限}


本节证明:由 Poisson--Stieltjes 积分表示的函数 $u$ 属于 $h^{1}$,且 $\lim_{r \to 1}u(re^{i\theta})$ 关于 $\theta$ 几乎处处存在.这个极限称为径向极限(radial limit).定理~\ref{thm:2.3.2} 称为 Fatou 定理.

\begin{theorem}\label{thm:2.3.1}
设 $\mu$ 为 $[-\pi, \pi]$ 上的(有界)实测度,且
\begin{equation*}
u(re^{i\theta}) = \frac{1}{2\pi}\int_{-\pi}^{\pi}P(r, \theta - t)\dd\mu(t),
\end{equation*}
则 $u$ 属于 $h^{1}$.
\end{theorem}

\begin{proof}
$u \in h(D)$ 已在 2.2 节注意到.为证明 $u \in h^{1}$,由定义只需证明 $\sup_{0 \leq r < 1}M_{1}(r, u) < \infty$.由于 $L^{1}$ 的对偶空间是 $L^{\infty}$,
\begin{equation*}
M_{1}(r, u) = \sup\left\{\left|\frac{1}{2\pi}\int_{-\pi}^{\pi}u(re^{i\theta})g(\theta)\dd\theta\right|;\ g \in L^{\infty} \text{ 且 } \|g\|_{\infty} \leq 1\right\},
\end{equation*}
由下列不等式可证明 $\sup_{0 \leq r < 1}M_{1}(r, u) \leq |\mu|([-\pi, \pi]) < \infty$.设 $g \in L^{\infty}$,由 Fubini 定理与命题~\ref{prop:2.2.1},
\begin{align*}
\left|\int_{-\pi}^{\pi}u(re^{i\theta})g(\theta)\dd\theta\right| & = \left|\frac{1}{2\pi}\int_{-\pi}^{\pi}\dd\mu(t)\int_{-\pi}^{\pi}P(r, \theta - t)g(\theta)\dd\theta\right| \\
& \leq \int_{-\pi}^{\pi}\dd|\mu|(t)\frac{1}{2\pi}\int_{-\pi}^{\pi}P(r, \theta - t)\|g\|_{\infty}\dd\theta \\
& = \|g\|_{\infty}|\mu|([-\pi, \pi]).
\end{align*}
\end{proof}

对实测度 $\mu \in M[-\pi, \pi]$ 令
\begin{equation*}
\mu(\theta) = \int_{0}^{\theta}\dd\mu(t).
\end{equation*}
即 $\mu(\theta)$ 表示区间 $[0, \theta]$ 或区间 $[\theta, 0]$ 关于 $\mu$ 的积分.

\begin{theorem}[Fatou]\label{thm:2.3.2}
$\mu \in M[-\pi, \pi]$ 是实测度且 $\mu'(\theta_{0})$ 存在时,设
\begin{equation*}
u(re^{i\theta}) = \frac{1}{2\pi}\int_{-\pi}^{\pi}P(r, \theta - t)\dd\mu(t),
\end{equation*}
则 $\lim_{r \to 1}u(re^{i\theta_{0}}) = \mu'(\theta_{0})$ 成立.
\end{theorem}

\begin{proof}
为简化记号,设 $\theta_{0} = 0$.令 $d\nu(t) = d\mu(t) - \mu'(0)\dd t$ 且 $d\lambda(t) = (\partial P(r, t)/\partial t)\dd t$,则分部积分得
\begin{align*}
& u(r) - \mu'(0) - \frac{1}{2\pi}\bigl[P(r, t)\nu(t)\bigr]_{-\pi}^{\pi} \\
& \quad = \frac{1}{2\pi}\int_{-\pi}^{\pi}P(r, t)\dd\nu(t) - \frac{1}{2\pi}\bigl[P(r, t)\nu(t)\bigr]_{-\pi}^{\pi} \\
& \quad = -\frac{1}{2\pi}\int_{-\pi}^{\pi}\nu(t)\dd\lambda(t) \\
& \quad = -\frac{1}{2\pi}\int_{-\pi}^{\pi}\bigl(\mu(t) - \mu'(0)t\bigr)\left(\frac{\partial}{\partial t}P(r, t)\right)\dd t \\
& \quad = -\frac{1}{2\pi}\int_{\delta < |t| < \pi}\bigl(\mu(t) - \mu'(0)t\bigr)\left(\frac{\partial}{\partial t}P(r, t)\right)\dd t \\
& \qquad - \frac{1}{2\pi}\int_{|t| \leq \delta}\bigl(\mu(t) - \mu'(0)t\bigr)\left(\frac{\partial}{\partial t}P(r, t)\right)\dd t \\
& \quad = I_{\delta}^{1} + I_{\delta}^{2}.
\end{align*}
此时,
\begin{align*}
|I_{\delta}^{1}| &\leq \sup_{\delta \leq |t| \leq \pi}\left|\frac{\partial}{\partial t}P(r, t)\right| \cdot \frac{1}{2\pi}\int_{\delta \leq |t| \leq \pi}|\mu(t) - \mu'(0)t|\dd t \\
\quad &\leq \frac{2r(1 - r^{2})}{(1 - 2r\cos\delta + r^{2})^{2}} \cdot \frac{1}{2\pi}\int_{\delta \leq |t| \leq \pi}|\mu(t) - \mu'(0)t|\dd t \\
\quad &\rightarrow 0 \quad (r \rightarrow 1).
\end{align*}
这是因为 $\mu$ 是有界实测度,所以 $\mu(t)$ 是 $[-\pi, \pi]$ 上的有界函数.另一方面,由 $\mu(0) = 0$,
\begin{equation*}
|I_{\delta}^{2}| \leq \sup_{|t| \leq \delta}\left|\frac{\mu(t) - \mu(0)}{t} - \mu'(0)\right|\frac{1}{2\pi}\int_{|t| \leq \delta}\left|t\frac{\partial}{\partial t}P(r, t)\right|\dd t.
\end{equation*}
对任意 $\varepsilon > 0$,把 $\delta$ 取得足够小,
\begin{equation*}
\sup_{|t| \leq \delta}\left|\frac{\mu(t) - \mu(0)}{t} - \mu'(0)\right| \leq \varepsilon
\end{equation*}
可以把 $r$ 取得足够接近 1,则
\begin{equation*}
|I_{\delta}^{2}| \leq \frac{\varepsilon}{2\pi}\int_{|t| \leq \delta}\left|t\frac{\partial}{\partial t}P(r, t)\right|\dd t < 2\varepsilon.
\end{equation*}
这里用到了
\begin{align*}
\frac{1}{2\pi}\int_{|t| \leq \delta}\left|t\frac{\partial}{\partial t}P(r, t)\right|\dd t & = \frac{1}{\pi}\int_{0}^{\delta}t\left(-\frac{\partial}{\partial t}P(r, t)\right)\dd t \\
& = \frac{1}{\pi}\left(-\delta P(r, \delta) + \int_{0}^{\delta}P(r, t)\dd t\right).
\end{align*}

另外,$\lim_{r \to 1}\bigl[P(r, t)\nu(t)\bigr]_{-\pi}^{\pi} = 0$,所以对任意 $\varepsilon > 0$,把 $r$ 取得足够接近 1,则可使 $|u(r) - \mu'(0)| \leq |I_{\delta}^{1}| + |I_{\delta}^{2}| < 3\varepsilon$.
\end{proof}

定理~\ref{thm:2.3.2} 中若 $d\mu = f(t)\dd t$,则
\begin{equation*}
\mu(\theta) = \int_{0}^{\theta}\dd\mu(t) = \int_{0}^{\theta}f(t)\dd t,
\end{equation*}
故 $\mu'(\theta) = f(\theta)$ a.e. $\theta \in [-\pi, \pi]$.因此 $\lim_{r \to 1}u(re^{i\theta}) = f(\theta)$ a.e. $\theta \in [-\pi, \pi]$.

% ============ 第2章 §2.4 Poisson-Stieltjes积分表示(原书 pp.132-134) ============
\section{Poisson--Stieltjes 积分表示}\label{sec:2.4}

本节证明属于 $h^{1}$ 的函数可以用 Poisson--Stieltjes 积分表示.因此证明属于 $h^{1}$ 的函数都拥有径向极限.

\begin{theorem}\label{thm:2.4.1}
若 $u \in h^{1}$,则存在唯一的实测度 $\mu \in M[-\pi, \pi]$ 满足 $\mu(\{-\pi\}) = \mu(\{\pi\})$,且
\begin{equation*}
u(re^{i\theta}) = \frac{1}{2\pi}\int_{-\pi}^{\pi}P(r, \theta - t)\dd\mu(t).
\end{equation*}
\end{theorem}

\begin{proof}
令 $r_{n} = 1 - \frac{1}{n}$ 且 $u_{n}(z) = u(r_{n}z)$,则 $u_{n}(z)$ 在 $|z| < 1/(1 - \frac{1}{n})$ 上调和,故如 2.2 节所示,
\begin{equation*}
u_{n}(z) = \frac{1}{2\pi}\int_{-\pi}^{\pi}P(r, \theta - t)u_{n}(e^{it})\dd t,
\end{equation*}
可以用 Poisson 积分表示.由于 $u \in h^{1}$,
\begin{align*}
\|u_{n}\|_{1} & = \int_{-\pi}^{\pi}|u_{n}(e^{it})|\dd t \leq \sup_{0 \leq r < 1}\int_{-\pi}^{\pi}|u(re^{it})|\dd t \\
& = 2\pi\sup_{0 \leq r < 1}M_{1}(r, u) < \infty,
\end{align*}
所以 $\{u_{n}\}$ 属于 $L^{1} = L^{1}[-\pi, \pi]$ 的半径为 $C$ 的球.这里 $C = 2\pi\sup_{0 \leq r < 1}M_{1}(r, u)$.

如 2.2 节所注意到的,$L^{1}$ 可以等距嵌入 $M = M[-\pi, \pi]$.$M$ 是 $[-\pi, \pi]$ 上连续函数的全体 $C[-\pi, \pi]$ 的对偶空间.因 $M$ 的单位球在弱$^{*}$拓扑下紧,故存在 $\{u_{n}\}$ 的子列 $\{u_{n_{j}}\}$ 与 $\mu \in M$,使得 $u_{n_{j}}$ 弱$^{*}$收敛于 $\mu$.此时 $\mu(\{-\pi\}) = \mu(\{\pi\})$.因此,固定 $r$ 与 $\theta$,因 $P(r, \theta - t) \in C[-\pi, \pi]$,
\begin{equation*}
\lim_{j \rightarrow \infty}\int_{-\pi}^{\pi}P(r, \theta - t)u_{n_{j}}(e^{it})\dd t = \int_{-\pi}^{\pi}P(r, \theta - t)\dd\mu(t)
\end{equation*}
成立.对每个 $j$,
\begin{equation*}
u_{n_{j}}(re^{i\theta}) = \frac{1}{2\pi}\int_{-\pi}^{\pi}P(r, \theta - t)u_{n_{j}}(e^{it})\dd t
\end{equation*}
且 $\lim_{j \rightarrow \infty}u_{n_{j}}(re^{i\theta}) = u(re^{i\theta})$,所以
\begin{equation*}
u(re^{i\theta}) = \frac{1}{2\pi}\int_{-\pi}^{\pi}P(r, \theta - t)\dd\mu(t)
\end{equation*}
成立.$\mu$ 的唯一性是命题~\ref{prop:2.2.3} 的结果.
\end{proof}

定理~\ref{thm:2.4.1} 中当 $u \in h^{p}$($1 < p \leq \infty$)时,可以取 $d\mu = f(t)\dd t$ 且 $f \in L^{p}$.这是因为,设 $1/p + 1/q = 1$,则 $L^{p} = (L^{q})^{*}$,可以像定理~\ref{thm:2.4.1} 的证明那样证明.

若 $f \in H^{1}$ 且 $f = u + iv$,则 $u, v \in h^{1}$,所以 $f$ 可以用复测度作 Poisson--Stieltjes 积分.2.12 节将证明该测度关于 Lebesgue 测度绝对连续.

\begin{corollary}\label{cor:2.4.2}
若 $u \in h(D)$ 在 $D$ 上为正,则存在正测度 $\mu \in M[-\pi, \pi]$ 满足 $\mu(\{-\pi\}) = \mu(\{\pi\})$,且
\begin{equation*}
u(re^{i\theta}) = \frac{1}{2\pi}\int_{-\pi}^{\pi}P(r, \theta - t)\dd\mu(t).
\end{equation*}
\end{corollary}

\begin{proof}
$u \in h^{1}$.因为
\begin{equation*}
\int_{-\pi}^{\pi}|u(re^{i\theta})|\dd\theta = \int_{-\pi}^{\pi}u(re^{i\theta})\dd\theta = 2\pi u(0).
\end{equation*}
由定理~\ref{thm:2.4.1},$u$ 可以由 $\mu$ 作积分表示,但 $\mu$ 是由 $u$ 定义的 $u_{n_{j}}$ 的弱$^{*}$极限.由 $u_{n_{j}}(z) > 0$($z \in \overline{D}$)知 $\mu$ 是正测度.
\end{proof}

% ============ 第2章 §2.5 Hardy空间的边界值 (I)(原书 pp.134-137) ============
\section{Hardy 空间的边界值 (I)}\label{sec:2.5}
\index{Nevanlinna 类}


如 2.4 节所注意到的,属于 $h^{1}$ 与 $H^{1}$ 的函数利用 Poisson--Stieltjes 积分表示,可以证明几乎处处存在径向极限.属于 $H^{p}$ 的函数,当 $0 < p < 1$ 时有时不能用 Poisson--Stieltjes 积分表示.但本节证明:属于 $H^{p}$ 的函数都几乎处处存在径向极限.

\begin{lemma}\label{lem:2.5.1}
$f \in H(D)$ 且
\begin{equation*}
\sup_{0 \leq r < 1}\int_{-\pi}^{\pi}\log^{+}|f(re^{i\theta})|\dd\theta < \infty,
\end{equation*}
则存在 $h, k \in H^{\infty}$ 使得 $f = h/k$.
\end{lemma}

\begin{proof}
设 $f \in H(D)$ 且 $f(z) \not\equiv 0$,则存在 $m \geq 0$ 使得 $f(z) = \sum_{n=m}^{\infty}a_{n}z^{n}$ 且 $a_{m} \neq 0$.设 $\{z_{n}\}$ 为 $f$ 的非零零点集,且 $0 < |z_{1}| \leq |z_{2}| \leq \cdots < 1$.若 $\{z_{n}\} \cap \{z;\ |z| = \rho < 1\} = \emptyset$,则
\begin{equation*}
F(z) = \log\left\{f(z)\frac{\rho^{m}}{z^{m}}\prod_{|z_{n}| < \rho}\left(\frac{\rho^{2} - \bar{z}_{n}z}{\rho(z - z_{n})}\right)\right\}
\end{equation*}
在 $|z| \leq \rho$ 上正则,且在 $|z| = \rho$ 上 $\operatorname{Re}\left\{F(z)\right\} = \log|f(z)|$.因此由 2.2 节所注意的事实,
\begin{equation*}
F(z) = \frac{1}{2\pi}\int_{-\pi}^{\pi}\log|f(\rho e^{it})|\frac{\rho e^{it} + z}{\rho e^{it} - z}\dd t + iC.
\end{equation*}
因此,当 $\rho < 1$ 时,令
\begin{equation*}
h_{\rho}(z) = \frac{z^{m}}{\rho^{m}}\prod_{|z_{n}| < \rho}\frac{\rho(z - z_{n})}{\rho^{2} - \bar{z}_{n}z}\exp\left\{-\frac{1}{2\pi}\int_{-\pi}^{\pi}\log^{-}|f(\rho e^{it})|\frac{\rho e^{it} + z}{\rho e^{it} - z}\dd t + iC\right\}
\end{equation*}
且
\begin{equation*}
k_{\rho}(z) = \exp\left\{-\frac{1}{2\pi}\int_{-\pi}^{\pi}\log^{+}|f(\rho e^{it})|\frac{\rho e^{it} + z}{\rho e^{it} - z}\dd t\right\},
\end{equation*}
则 $f(z) = h_{\rho}(z)/k_{\rho}(z)$.下面令 $\rho \to 1$ 来证明本引理.

设 $\{\rho_{\ell}\}$ 为满足 $\rho_{\ell} \nearrow 1$ 且 $\{z_{n}\} \cap \{z;\ |z| = \rho_{\ell}\} = \emptyset$ 的数列.令 $H_{\ell}(z) = h_{\rho_{\ell}}(\rho_{\ell}z)$,$K_{\ell}(z) = k_{\rho_{\ell}}(\rho_{\ell}z)$,则 $f(\rho_{\ell}z) = H_{\ell}(z)/K_{\ell}(z)$($z \in D$).$H_{\ell}$ 与 $K_{\ell}$ 在 $D$ 上正则,且 $|H_{\ell}(z)| \leq 1$,$|K_{\ell}(z)| \leq 1$($z \in D$).因此 $\{H_{\ell}\}$ 与 $\{K_{\ell}\}$ 是正规族,故对任意含于 $D$ 的闭集 $E$,存在数列 $\{\ell_{j}\}$ 使 $H_{\ell_{j}}$ 与 $K_{\ell_{j}}$ 在 $E$ 上一致收敛.设其极限为 $h$ 与 $k$,由一致定理,$h$ 与 $k$ 与 $E$ 无关,在 $D$ 上正则,$|h(z)| \leq 1$ 且 $|k(z)| \leq 1$.因此令 $j \to \infty$,应当有 $f(z) = h(z)/k(z)$($z \in D$),但必须证明 $k(z) \not\equiv 0$.这里才首次用到关于 $f$ 的条件
\begin{equation*}
\sup_{0 \leq r < 1}\int_{-\pi}^{\pi}\log^{+}|f(re^{i\theta})|\dd\theta < \infty.
\end{equation*}
由于 $\log^{+}|f(\rho_{\ell}z)| \geq -\log|K_{\ell}(z)|$,利用 Jensen 公式(命题~\ref{prop:2.2.1})与 $f$ 的条件,得 $|K_{\ell}(0)| \geq \varepsilon > 0$,所以 $k(z) \not\equiv 0$.
\end{proof}

\begin{lemma}\label{lem:2.5.2}
若 $k \in H^{\infty}$ 且 $k(z) \not\equiv 0$,则 $k$ 几乎处处存在径向极限 $k(e^{i\theta})$,且 $\log|k(e^{i\theta})| \in L^{1}$.
\end{lemma}

\begin{proof}
可以设 $\|k\|_{\infty} \leq 1$.由 $k(z) \not\equiv 0$,可以写 $k(z) = z^{m}k_{0}(z)$,$m \geq 0$ 且 $k_{0}(0) \neq 0$,故可以设 $k(0) \neq 0$.由 Jensen 公式(命题~\ref{prop:2.2.1}),
\begin{equation*}
\frac{1}{2\pi}\int_{-\pi}^{\pi}\log|k(re^{i\theta})|\dd\theta = \log|k(0)| + \sum_{|z_{n}| < r}\log\frac{r}{|z_{n}|}
\end{equation*}
成立.这里 $\{z_{n}\}$ 是 $k$ 的零点集.因此
\begin{equation*}
\frac{1}{2\pi}\int_{-\pi}^{\pi}\log|k(re^{i\theta})|\dd\theta
\end{equation*}
关于 $r$ 单调递增.由于 $k$ 属于 $H^{1}$,如 2.4 节所注意到的,几乎处处存在径向极限 $k(e^{i\theta})$,故由 Lebesgue 积分的 Fatou 引理,
\begin{equation*}
\int_{-\pi}^{\pi}\left|\log|k(e^{i\theta})|\right|\dd\theta \leq \varliminf_{r \to 1}\left\{-\int_{-\pi}^{\pi}\log|k(re^{i\theta})|\dd\theta\right\}.
\end{equation*}
上式右边的积分单调递减,故为有限值,$\log|k(e^{i\theta})| \in L^{1}$.
\end{proof}

\begin{theorem}\label{thm:2.5.3}
当 $0 < p \leq \infty$ 时,若 $f \in H^{p}$ 且 $f(z) \not\equiv 0$,则 $f$ 几乎处处存在径向极限 $f(e^{i\theta})$.此时 $\log|f(e^{i\theta})| \in L^{1}$ 且 $f(e^{i\theta}) \in L^{p}$.
\end{theorem}

\begin{proof}
当 $0 < p < \infty$ 时,由 $\lim_{x \to \infty}(\log x)/x^{p} = 0$,若 $f \in H^{p}$ 则
\begin{equation*}
\sup_{0 \leq r < 1}\int_{-\pi}^{\pi}\log^{+}|f(re^{i\theta})|\dd\theta < \infty.
\end{equation*}
因此由引理~\ref{lem:2.5.1},若 $f(z) \not\equiv 0$,则存在 $h, k \in H^{\infty}$ 使得 $f = h/k$.由引理~\ref{lem:2.5.2},$h$ 与 $k$ 几乎处处存在径向极限 $h(e^{i\theta})$ 与 $k(e^{i\theta})$,且 $k(e^{i\theta})$ 在正测度集上不为零,所以径向极限 $f(e^{i\theta})$ 存在.由于 $f \in H^{p}$,由 Fatou 引理知 $f(e^{i\theta}) \in L^{p}$.又由引理~\ref{lem:2.5.2},$\log|h(e^{i\theta})| \in L^{1}$ 且 $\log|k(e^{i\theta})| \in L^{1}$,故 $\log|f(e^{i\theta})| \in L^{1}$.
\end{proof}

引理~\ref{lem:2.5.1} 的逆也成立.用 Jensen 公式证明逆命题是容易的.满足引理~\ref{lem:2.5.1} 的条件的 $H(D)$ 的函数的全体称为 Nevanlinna 类 $N$,且 $N \supset \bigcup_{p > 0}H^{p}$.

% ============ 第2章 §2.6 Blaschke积与H^p的零点集(原书 pp.137-141) ============
\section{\texorpdfstring{Blaschke 积与 $H^p$ 的零点集}{Blaschke 积与 Hp 的零点集}}\label{sec:2.6}
\index{Blaschke 积!(单位圆板上的)}


当 $|a| < 1$ 且 $a \neq 0$ 时,$\frac{|a|}{a}\frac{a - z}{1 - \bar{a}z}$ 是从 $D$ 到 $D$ 上的一次变换,且
\begin{align*}
\left|1 - \frac{|a|}{a}\frac{a - z}{1 - \bar{a}z}\right| & = (1 - |a|)\frac{|a + |a|z|}{|a|\cdot|1 - \bar{a}z|} \\
& \leq (1 - |a|)\cdot\frac{2}{1 - |z|} \quad (z \in D).
\end{align*}
因此,设 $\{a_{j}\}_{j=1}^{\infty} \subset D \setminus \{0\}$,则
\begin{equation*}
\sum_{j = 1}^{\infty}\left|1 - \frac{|a_{j}|}{a_{j}}\frac{a_{j} - z}{1 - \bar{a}_{j}z}\right| \leq \frac{2}{1 - |z|}\sum_{j = 1}^{\infty}(1 - |a_{j}|) \quad (z \in D \setminus \{a_{j}\}_{j=1}^{\infty}),
\end{equation*}
所以
\begin{equation*}
\prod_{j = 1}^{\infty}\frac{|a_{j}|}{a_{j}}\frac{a_{j} - z}{1 - \bar{a}_{j}z}
\end{equation*}
对每个 $z \in D$ 绝对收敛的充要条件是
\begin{equation*}
\sum_{j = 1}^{\infty}(1 - |a_{j}|) < \infty
\end{equation*}
(必要性由引理~\ref{lem:2.6.2} 证明).当 $\sum_{j=1}^{\infty}(1 - |a_{j}|) < \infty$ 时,
\begin{equation*}
B(z) = z^{m}\prod_{j = 1}^{\infty}\frac{|a_{j}|}{a_{j}}\frac{a_{j} - z}{1 - \bar{a}_{j}z}
\end{equation*}
称为 Blaschke 积.这里 $m$ 是非负整数,且允许 $a_{j} = a_{i}$($j \neq i$).

当 $0 < x < 1$ 时 $\lim_{x \to 1}(1 - x)/(-\log x) = 1$,故 Blaschke 条件 $\sum_{j=1}^{\infty}(1 - |a_{j}|) < \infty$ 与 $\sum_{j=1}^{\infty}\log 1/|a_{j}| < \infty$ 等价.

\begin{theorem}\label{thm:2.6.1}
对 $\{a_{j}\}_{j=1}^{\infty}$,设 $B$ 为其 Blaschke 积,则 $B \in H^{\infty}$,$|B(z)| < 1$($z \in D$),且 $B$ 的径向极限满足 $|B(e^{i\theta})| = 1$ a.e. $\theta$.$\{a_{j}\}_{j=1}^{\infty}$ 是 $B$ 的零点,$B$ 在 $D \setminus \{0\}$ 上无其他零点.
\end{theorem}

\begin{proof}
若 $\sum_{j=1}^{\infty}(1 - |a_{j}|) < \infty$,则如上所述,无穷乘积 $B$ 在 $D$ 的任意紧集上一致收敛,故 $B \in H(D)$.此时 $B \in H^{\infty}$ 且 $|B(z)| < 1$($z \in D$)是显然的.又可以把 $\{a_{j}\}_{j=1}^{\infty} \cup \{0\}$ 视为 $B$ 的零点集,故设 $B(0) \neq 0$,则假设 $\sum_{j=1}^{\infty}(1 - |a_{j}|) < \infty$ 蕴含 $\log|B(0)| = \sum_{j=1}^{\infty}\log|a_{j}| > -\infty$.此时,由 Jensen 公式(命题~\ref{prop:2.2.1}),
\begin{align*}
\frac{1}{2\pi}\int_{-\pi}^{\pi}\log|B(re^{i\theta})|\dd\theta & = \sum_{|a_{j}| < r}\log\left(\frac{r}{|a_{j}|}\right) + \log|B(0)| \\
& = \sum_{|a_{j}| < r}\log\left(\frac{r}{|a_{j}|}\right) - \sum_{j = 1}^{\infty}\log\frac{1}{|a_{j}|}.
\end{align*}
对任意 $\varepsilon > 0$,存在 $n$ 使得 $\sum_{j > n}\log 1/|a_{j}| < \varepsilon$.把 $r$ 取得足够接近 1 使 $|a_{j}| < r$($j = 1, 2, \ldots, n$),则
\begin{equation*}
\frac{1}{2\pi}\int_{-\pi}^{\pi}\log|B(re^{i\theta})|\dd\theta \geq \sum_{j = 1}^{n}\log\frac{r}{|a_{j}|} - \sum_{j = 1}^{n}\log\frac{1}{|a_{j}|} - \varepsilon.
\end{equation*}
这里,再把 $r$ 取得足够接近 1,则
\begin{equation*}
\frac{1}{2\pi}\int_{-\pi}^{\pi}\log|B(re^{i\theta})|\dd\theta \geq -2\varepsilon,
\end{equation*}
故令 $\varepsilon \to 0$,
\begin{equation*}
\varlimsup_{r \to 1}\frac{1}{2\pi}\int_{-\pi}^{\pi}\log|B(re^{i\theta})|\dd\theta \geq 0.
\end{equation*}
由引理~\ref{lem:2.5.2},$\lim_{r \to 1}B(re^{i\theta}) = B(e^{i\theta})$ a.e. $\theta$ 存在,且 $\log|B(e^{i\theta})| \leq 0$ a.e. $\theta$.由 Lebesgue 积分的 Fatou 引理,
\begin{align*}
0 &\leq \frac{1}{2\pi}\int_{-\pi}^{\pi}\bigl(-\log|B(e^{i\theta})|\bigr)\dd\theta \\
&\leq \varliminf_{r \to 1}\frac{1}{2\pi}\int_{-\pi}^{\pi}\bigl(-\log|B(re^{i\theta})|\bigr)\dd\theta \\
&= -\varlimsup_{r \to 1}\frac{1}{2\pi}\int_{-\pi}^{\pi}\log|B(re^{i\theta})|\dd\theta \leq 0
\end{align*}
故 $|B(e^{i\theta})| = 1$ a.e. $\theta$.
\end{proof}

若 $\{a_{j}\}_{j=1}^{\infty} \subset D$ 满足 $\sum_{j=1}^{\infty}(1 - |a_{j}|) < \infty$,则定理~\ref{thm:2.6.1} 表明存在 $f \in H^{\infty}$ 以其零点集为 $\{a_{j}\}_{j=1}^{\infty}$.其逆是否成立是自然的问题,下面回答这个问题.

\begin{lemma}\label{lem:2.6.2}
设 $\{a_{j}\}_{j=1}^{\infty} \subset D$ 为 $f \in H(D)$ 的零点集.若 $f(z) \not\equiv 0$ 且
\begin{equation*}
\sup_{0 \leq r < 1}\int_{-\pi}^{\pi}\log^{+}|f(re^{i\theta})|\dd\theta < \infty,
\end{equation*}
则 $\sum_{j=1}^{\infty}(1 - |a_{j}|) < \infty$.
\end{lemma}

\begin{proof}
若 $f(0) = 0$,则对 $f(z)/z^{\ell}$ 证明即可,故设 $f(0) \neq 0$.若 $0 < r < 1$ 且 $|a_{j}| \neq r$($j = 1, 2, \ldots$),则由 Jensen 公式(命题~\ref{prop:2.2.1}),
\begin{equation*}
\log|f(0)| = \frac{1}{2\pi}\int_{-\pi}^{\pi}\log|f(re^{i\theta})|\dd\theta + \sum_{|a_{j}| < r}\log\left|\frac{a_{j}}{r}\right|.
\end{equation*}
令 $M = \sup_{0 \leq r < 1}\frac{1}{2\pi}\int_{-\pi}^{\pi}\log^{+}|f(re^{i\theta})|\dd\theta$,则
\begin{equation*}
\sum_{|a_{j}| < r}\log\left|\frac{r}{a_{j}}\right| \leq M - \log|f(0)|,
\end{equation*}
令 $r \to 1$,对任意 $k$,
\begin{equation*}
\sum_{j = 1}^{k}\log\frac{1}{|a_{j}|} \leq M - \log|f(0)|,
\end{equation*}
故
\begin{equation*}
\sum_{j = 1}^{\infty}\log\frac{1}{|a_{j}|} < \infty.
\end{equation*}
\end{proof}

\begin{theorem}\label{thm:2.6.3}
设 $0 < p \leq \infty$.若 $\{a_{j}\}_{j=1}^{\infty}$ 为 $f \in H^{p}$ 的零点集,则 $\sum_{j=1}^{\infty}(1 - |a_{j}|) < \infty$.反之,若 $\sum_{j=1}^{\infty}(1 - |a_{j}|) < \infty$,则存在以 $\{a_{j}\}_{j=1}^{\infty}$ 为零点集的 $f \in H^{p}$.
\end{theorem}

\begin{proof}
可以设 $f(0) \neq 0$.由定理~\ref{thm:2.5.3} 的证明,若 $f \in H^{p}$,则 $f$ 满足引理~\ref{lem:2.6.2} 的假设,故 $\sum_{j=1}^{\infty}(1 - |a_{j}|) < \infty$.反之,若 $\sum_{j=1}^{\infty}(1 - |a_{j}|) < \infty$,由定理~\ref{thm:2.6.1},取 $\{a_{j}\}_{j=1}^{\infty}$ 的 Blaschke 积 $B$ 作为 $f$ 即可.
\end{proof}

\begin{theorem}\label{thm:2.6.4}
设 $0 < p \leq \infty$.若 $f \in H^{p}$ 且 $f(z) \not\equiv 0$,则 $f$ 可以分解为 $f(z) = B(z)g(z)$.这里 $B$ 是 Blaschke 积,$g$ 在 $D$ 内无零点且属于 $H^{p}$.
\end{theorem}

\begin{proof}
可以设 $f(0) \neq 0$ 且其零点集为无限集.设 $\{a_{j}\}_{j=1}^{\infty}$ 为 $f \in H^{p}$ 的零点集,由定理~\ref{thm:2.6.3} 得 $\sum_{j=1}^{\infty}(1 - |a_{j}|) < \infty$.由定理~\ref{thm:2.6.1} 存在 $\{a_{j}\}_{j=1}^{\infty}$ 的 Blaschke 积 $B$,令 $g(z) = f(z)/B(z)$,则 $g$ 在 $D$ 内无零点.还需证明 $g \in H^{p}$.

令
\begin{equation*}
B_{n}(z) = \prod_{j = 1}^{n}\frac{|a_{j}|}{a_{j}}\frac{a_{j} - z}{1 - \bar{a}_{j}z} \quad \text{且} \quad g_{n}(z) = f(z)/B_{n}(z),
\end{equation*}
则对每个 $n$ 和任意 $\varepsilon > 0$,当 $|z|$ 足够接近 1 时可使 $|B_{n}(z)| > 1 - \varepsilon$.因此,当 $r$ 足够接近 1 时,
\begin{align*}
\frac{1}{2\pi}\int_{-\pi}^{\pi}|g_{n}(re^{i\theta})|^{p}\dd\theta & \leq (1 - \varepsilon)^{-p}\frac{1}{2\pi}\int_{-\pi}^{\pi}|f(re^{i\theta})|^{p}\dd\theta \\
& \leq (1 - \varepsilon)^{-p}\sup_{0 \leq r < 1}M_{p}(r, f)^{p}.
\end{align*}
令 $\varepsilon \to 0$,当 $r$ 足够接近 1 时,
\begin{equation*}
\frac{1}{2\pi}\int_{-\pi}^{\pi}|g_{n}(re^{i\theta})|^{p}\dd\theta \leq \sup_{0 \leq r < 1}M_{p}(r, f)^{p}.
\end{equation*}
由引理~\ref{lem:2.1.1},可证当 $r' < r$ 时 $M_{p}(r', g_{n}) \leq M_{p}(r, g_{n})$,故上面的不等式对所有 $r$ 成立.因此由 Fatou 引理,
\begin{equation*}
\sup_{0 \leq r < 1}M_{p}(r, g) \leq \sup_{0 \leq r < 1}M_{p}(r, f) < \infty
\end{equation*}
得证.
\end{proof}

% ============ 第2章 §2.7 Hardy空间的边界值 (II)(原书 pp.141-144) ============
\section{Hardy 空间的边界值 (II)}\label{sec:2.7}

2.5 节证明了属于 $H^{p}$ 的函数 $f$ 几乎处处存在径向极限 $\lim_{r \to 1}f(re^{i\theta}) = f(e^{i\theta})$,且 $f \in L^{p}$.$H^{p}$ 在 2.1 节是在 $D$ 的内部定义的,本节证明 $H^{p}$ 可以用径向极限来定义.另外,由此证明在 $L^{p}$ 的距离下 $\lim_{r \to 1}f(re^{i\theta}) = f(e^{i\theta})$.

\begin{theorem}\label{thm:2.7.1}
设 $0 < p \leq \infty$.若 $p \neq \infty$,则
\begin{align*}
\sup_{0 \leq r < 1}M_{p}(r, f) & = \lim_{r \to 1}M_{p}(r, f) \\
& = \left\{\frac{1}{2\pi}\int_{-\pi}^{\pi}|f(e^{i\theta})|^{p}\dd\theta\right\}^{1/p} = \|f\|_{p},
\end{align*}
又
\begin{equation*}
\sup_{0 \leq r < 1}M_{\infty}(r, f) = \esssup_{-\pi \leq \theta < \pi}|f(e^{i\theta})| = \|f\|_{\infty}.
\end{equation*}
因此,$H^{p}$ 在 $1 \leq p \leq \infty$ 时是范数空间,在 $0 < p < 1$ 时是度量空间.
\end{theorem}

\begin{proof}
若 $f(z) = \sum_{j=0}^{\infty}a_{j}z^{j} \in H^{2}$,则
\begin{equation*}
M_{2}(r, f) = \left(\sum_{j = 0}^{\infty}|a_{j}|^{2}r^{2j}\right)^{1/2},
\end{equation*}
故由径向极限的存在性与 Lebesgue 积分的 Fatou 引理,
\begin{align*}
\sup_{0 \leq r < 1}M_{2}(r, f) & = \lim_{r \to 1}M_{2}(r, f) \\
& = \left(\sum_{j = 0}^{\infty}|a_{j}|^{2}\right)^{1/2} = \left\{\frac{1}{2\pi}\int_{-\pi}^{\pi}|f(e^{i\theta})|^{2}\dd\theta\right\}^{1/2}.
\end{align*}
若 $f \in H^{p}$($0 < p < \infty$),由定理~\ref{thm:2.6.4} 可分解 $f = Bg$.由于 $g$ 在 $D$ 内无零点,$[g(z)]^{p/2} \in H^{2}$,且 $|B(e^{i\theta})| = 1$ a.e. $\theta$,所以
\begin{align*}
\int_{-\pi}^{\pi}|f(re^{i\theta})|^{p}\dd\theta &\leq \int_{-\pi}^{\pi}|g(re^{i\theta})|^{p}\dd\theta \\
&\rightarrow \int_{-\pi}^{\pi}|g(e^{i\theta})|^{p}\dd\theta = \int_{-\pi}^{\pi}|f(e^{i\theta})|^{p}\dd\theta.
\end{align*}
由引理~\ref{lem:2.1.1},$\sup_{0 \leq r < 1}M_{p}(r, f) = \lim_{r \to 1}M_{p}(r, f)$,由 Fatou 引理,$p \neq \infty$ 的情形得证.当 $p = \infty$ 时,$\|f\|_{\infty} \leq \sup_{0 \leq r < 1}M_{\infty}(r, f)$ 由径向极限的存在性与最大值原理是显然的.若 $f \in H^{\infty}$ 则 $f \in H^{2}$,所以证明它可以作 Poisson 积分表示是容易的,由此导出 $\sup_{0 \leq r < 1}M_{\infty}(r, f) \leq \|f\|_{\infty}$.
\end{proof}

\begin{lemma}\label{lem:2.7.2}
设 $0 < p < \infty$ 且 $k_{n} \in L^{p}$($n = 1, 2, \ldots$).若 $\lim_{n \to \infty}k_{n}(\theta) = k(\theta)$ a.e. $\theta$,且
\begin{equation*}
\lim_{n \to \infty}\int_{-\pi}^{\pi}|k_{n}(\theta)|^{p}\dd\theta = \int_{-\pi}^{\pi}|k(\theta)|^{p}\dd\theta < \infty,
\end{equation*}
则
\begin{equation*}
\lim_{n \to \infty}\int_{-\pi}^{\pi}|k_{n}(\theta) - k(\theta)|^{p}\dd\theta = 0.
\end{equation*}
\end{lemma}

\begin{proof}
设 $E$ 为 $[-\pi, \pi]$ 的 Lebesgue 可测集,令
\begin{equation*}
\Phi_{n}(E) = \int_{E}|k_{n}(\theta)|^{p}\dd\theta, \quad \Phi(E) = \int_{E}|k(\theta)|^{p}\dd\theta.
\end{equation*}
设 $E^{c}$ 为 $E$ 在 $[-\pi, \pi]$ 中的补集,由 Fatou 引理与假设,
\begin{align*}
\Phi(E) & \leq \varliminf_{n \to \infty}\Phi_{n}(E) \leq \varlimsup_{n \to \infty}\Phi_{n}(E) \\
& = \varliminf_{n \to \infty}\Phi_{n}([-\pi, \pi]) - \varliminf_{n \to \infty}\Phi_{n}(E^{c}) \\
& \leq \Phi([-\pi, \pi]) - \Phi(E^{c}) = \Phi(E).
\end{align*}
因此,对每个 $E \subset [-\pi, \pi]$,$\lim_{n \to \infty}\Phi_{n}(E) = \Phi(E)$ 成立.对任意 $\varepsilon > 0$,取 $F \subset [-\pi, \pi]$ 满足 $\Phi(F^{c}) < \varepsilon$.取 $\delta > 0$ 使得若 $Q \subset F$ 且 $d\theta(Q) < \delta$ 则 $\Phi(Q) < \varepsilon$.这由 $\Phi$ 关于 $d\theta$ 绝对连续是可能的.由 Egorov 定理,存在 $Q \subset F$ 使 $d\theta(Q) < \delta$,且 $k_{n}$ 在 $E = F \cap Q^{c}$ 上一致收敛.因此,利用当 $a \geq 0$ 且 $b \geq 0$ 时 $(a + b)^{p} \leq 2^{p}(a^{p} + b^{p})$ 成立,
\begin{align*}
\int_{-\pi}^{\pi}|k_{n}(\theta) - k(\theta)|^{p}\dd\theta&= \int_{F^{c}}|k_{n}(\theta) - k(\theta)|^{p}\dd\theta + \int_{Q}|k_{n}(\theta) - k(\theta)|^{p}\dd\theta  \\
&\quad + \int_{E}|k_{n}(\theta) - k(\theta)|^{p}\dd\theta \\
&\leq 2^{p}\left(\Phi_{n}(F^{c}) + \Phi(F^{c}) + \Phi_{n}(Q) + \Phi(Q)\right) \\
&\quad + \int_{E}|k_{n}(\theta) - k(\theta)|^{p}\dd\theta.
\end{align*}
由 $\lim_{n \to \infty}\Phi_{n}(F^{c}) = \Phi(F^{c}) < \varepsilon$ 且 $\lim_{n \to \infty}\Phi_{n}(Q) = \Phi(Q) < \varepsilon$,以及一致收敛给出
\begin{equation*}
\lim_{n \to \infty}\int_{E}|k_{n}(\theta) - k(\theta)|^{p}\dd\theta = 0,
\end{equation*}
故对任意 $\varepsilon > 0$,
\begin{equation*}
\varlimsup_{n \to \infty}\int_{-\pi}^{\pi}|k_{n}(\theta) - k(\theta)|^{p}\dd\theta \leq 2^{p}\cdot 4\varepsilon.
\end{equation*}
\end{proof}

\begin{theorem}\label{thm:2.7.3}
当 $0 < p < \infty$ 时,若 $f \in H^{p}$,则
\begin{equation*}
\lim_{r \to 1}\int_{-\pi}^{\pi}|f(re^{i\theta}) - f(e^{i\theta})|^{p}\dd\theta = 0.
\end{equation*}
\end{theorem}

\begin{proof}
由定理~\ref{thm:2.5.3} 与定理~\ref{thm:2.7.1},$\lim_{r \to 1}f(re^{i\theta}) = f(e^{i\theta})$ a.e. $\theta$,且
\begin{equation*}
\lim_{r \to 1}\int_{-\pi}^{\pi}|f(re^{i\theta})|^{p}\dd\theta = \int_{-\pi}^{\pi}|f(e^{i\theta})|^{p}\dd\theta.
\end{equation*}
由引理~\ref{lem:2.7.2},
\begin{equation*}
\lim_{r \to 1}\int_{-\pi}^{\pi}|f(re^{i\theta}) - f(e^{i\theta})|^{p}\dd\theta = 0.
\end{equation*}
\end{proof}

\begin{corollary}\label{cor:2.7.4}
当 $0 < p \leq \infty$ 时,若 $f \in H^{p}$,则
\begin{equation*}
\lim_{r \to 1}\int_{-\pi}^{\pi}\left|\log^{+}|f(re^{i\theta})| - \log^{+}|f(e^{i\theta})|\right|\dd\theta = 0.
\end{equation*}
\end{corollary}

\begin{proof}
若 $a \geq 0, b \geq 0$ 且 $0 < p \leq 1$,则
\begin{equation*}
\left|\log^{+}a - \log^{+}b\right| \leq \frac{1}{p}|a - b|^{p},
\end{equation*}
由此与定理~\ref{thm:2.7.3} 可以证明系.上面的不等式由
\begin{equation*}
\log x \leq \frac{1}{p}(x - 1)^{p} \quad (x \geq 1)
\end{equation*}
当 $1 \leq b < a$ 且 $x = a/b$ 时可得.对 $1 < p \leq \infty$,由于若 $p < p'$ 则 $H^{p} \supset H^{p'}$,比 $0 < p \leq 1$ 的情形更显然.
\end{proof}

% ============ 第2章 §2.8 内函数与外函数(原书 pp.144-147) ============
\section{内函数与外函数}\label{sec:2.8}
\index{内函数!(单位圆板上的)}
\index{外函数!(单位圆板上的)}
\index{内函数!(单位圆板上的)奇异内函数}
\index{外函数!(单位圆板上的)奇异内函数}
\index{因式分解定理}


定理~\ref{thm:2.6.4} 给出了属于 $H^{p}$ 的函数 $f$ 的分解 $f = Bg$,其中 $g$ 在 $D$ 内无零点,本节进一步分解 $g$.事实上,利用对某个实测度 $\mu \in M[-\pi, \pi]$ 的 Herglotz 积分表示(见 2.2 节)
\begin{equation*}
\log g(z) = \frac{1}{2\pi}\int_{-\pi}^{\pi}\frac{e^{it} + z}{e^{it} - z}\dd\mu(t) \quad (z \in D).
\end{equation*}
此时,对 $\mu$ 作 Lebesgue 分解 $\mu = \mu_{a} + \mu_{s}$,$\mu_{s}$ 是正测度.这里 $\mu_{a}$ 关于 Lebesgue 测度绝对连续且 $\mu_{s}$ 奇异.定理~\ref{thm:2.8.3} 表明 $\mu$ 的这个分解给出 $g$ 的分解.下面的定理~\ref{thm:2.8.1} 是命题~\ref{prop:2.2.1} 的推广,但对证明因数分解定理~\ref{thm:2.8.3} 很重要.

\begin{theorem}\label{thm:2.8.1}
设 $0 < p \leq \infty$.若 $f \in H^{p}$,则
\begin{equation*}
\log|f(re^{i\theta})| \leq \frac{1}{2\pi}\int_{-\pi}^{\pi}P(r, \theta - t)\log|f(e^{it})|\dd t.
\end{equation*}
\end{theorem}

\begin{proof}
由定理~\ref{thm:2.6.4} 可分解 $f = Bg$,故若对 $g$ 证明,则由 $|B(z)| < 1$($z \in D$)且 $|B(e^{i\theta})| = 1$ a.e. $\theta$ 可对 $f$ 证明不等式.因此可以认为 $f$ 在 $D$ 内无零点.此时 $\log|f(z)|$ 在 $D$ 上调和,故当 $r < \rho < 1$ 时,
\begin{equation*}
\log|f(\rho re^{i\theta})| = \frac{1}{2\pi}\int_{-\pi}^{\pi}P(r, \theta - t)\log|f(\rho e^{it})|\dd t.
\end{equation*}
由系~\ref{cor:2.7.4},
\begin{equation*}
\lim_{\rho \to 1}\int_{-\pi}^{\pi}P(r, \theta - t)\log^{+}|f(\rho e^{it})|\dd t = \int_{-\pi}^{\pi}P(r, \theta - t)\log^{+}|f(e^{it})|\dd t,
\end{equation*}
且由 Lebesgue 积分的 Fatou 引理,
\begin{equation*}
\lim_{\rho \to 1}\int_{-\pi}^{\pi}P(r, \theta - t)\log^{-}|f(\rho e^{it})|\dd t \geq \int_{-\pi}^{\pi}P(r, \theta - t)\log^{-}|f(e^{it})|\dd t,
\end{equation*}
由此可导出本定理.
\end{proof}

$Q \in H^{\infty}$ 称为内函数,是指 $|Q(e^{i\theta})| = 1$ a.e. $\theta$;$F \in H^{p}$ 称为外函数,是指
\begin{equation*}
\log|F(re^{i\theta})| = \frac{1}{2\pi}\int_{-\pi}^{\pi}P(r, \theta - t)\log|F(e^{it})|\dd t.
\end{equation*}
定理~\ref{thm:2.8.1} 表明不等式总是成立.若 $F$ 是外函数,则 $\log|F| \in L^{1}$,故 $F$ 在 $D$ 内无零点.由定理~\ref{thm:2.6.1},Blaschke 积 $B$ 是内函数.下面的定理~\ref{thm:2.8.2} 与在 $D$ 内无零点的内函数和外函数的 Herglotz 积分表示有关.设 $\mu \in M[-\pi, \pi]$ 为关于 Lebesgue 测度奇异的正测度,且
\begin{equation*}
S(z) = \exp\left(-\int_{-\pi}^{\pi}\frac{e^{it} + z}{e^{it} - z}\dd\mu(t)\right),
\end{equation*}
则 $\mu'(\theta) = 0$ a.e. $\theta$,且
\begin{equation*}
\log|S(z)| = -\int_{-\pi}^{\pi}P(r, \theta - t)\dd\mu(t),
\end{equation*}
故由定理~\ref{thm:2.3.2},$|S(e^{i\theta})| = 1$ a.e. $\theta$.因此 $S$ 是内函数,且 $S$ 在 $D$ 内无零点.$S$ 称为奇异内函数.

\begin{theorem}\label{thm:2.8.2}
(1) 若 $Q$ 是在 $D$ 内无零点的内函数,则存在奇异的有界正测度 $\mu$,满足 $\mu(\{-\pi\}) = \mu(\{\pi\})$,使得
\begin{equation*}
Q(z) = e^{i\gamma}\exp\left\{-\int_{-\pi}^{\pi}\frac{e^{it} + z}{e^{it} - z}\dd\mu(t)\right\}
\end{equation*}
可以表示.这里 $\gamma$ 是实常数.

(2) 若 $F$ 是外函数,则
\begin{equation*}
F(z) = e^{i\gamma}\exp\left\{\int_{-\pi}^{\pi}\frac{e^{it} + z}{e^{it} - z}\log|F(e^{it})|\dd t\right\}
\end{equation*}
可以表示.这里 $\gamma$ 是实常数.
\end{theorem}

\begin{proof}
(1) 由最大值原理与 Fatou 定理~\ref{thm:2.3.2},$|Q(z)| < 1$($z \in D$),故 $-\log|Q(z)|$ 是 $D$ 上的正调和函数.由系~\ref{cor:2.4.2},存在正测度 $\mu \in M[-\pi, \pi]$ 使得
\begin{equation*}
-\log|Q(z)| = \frac{1}{2\pi}\int_{-\pi}^{\pi}P(r, \theta - t)\dd\mu(t) \quad (z \in D).
\end{equation*}
由 $|Q(e^{i\theta})| = 1$ a.e. $\theta$ 与定理~\ref{thm:2.3.2},$\mu'(\theta) = 0$ a.e. $\theta$.因此 $\mu$ 是奇异正测度.

令
\begin{equation*}
q(z) = \frac{1}{2\pi}\int_{-\pi}^{\pi}\frac{e^{it} + z}{e^{it} - z}\dd\mu(t),
\end{equation*}
则 $\operatorname{Re}(q(z)) = -\log|Q(z)|$($z \in D$).因此存在实常数 $\gamma$ 使得 $-\log Q(z) = q(z) + i\gamma$.

(2) 当 $F$ 是外函数时,令
\begin{equation*}
G(z) = \frac{1}{2\pi}\int_{-\pi}^{\pi}\frac{e^{it} + z}{e^{it} - z}\log|F(e^{it})|\dd t,
\end{equation*}
则由外函数的定义,$\operatorname{Re}G(z) = \log|F(z)|$($z \in D$).因此存在实常数 $\gamma$ 使得 $\log F(z) = G(z) + i\gamma$.
\end{proof}

\begin{theorem}\label{thm:2.8.3}
设 $0 < p \leq \infty$.若 $f \in H^{p}$ 不恒等于零,则 $f$ 可以唯一分解为 $f = BSF$.这里 $B$ 是 Blaschke 积,$S$ 是奇异内函数,$F$ 是 $H^{p}$ 中的外函数.
\end{theorem}

\begin{proof}
由定理~\ref{thm:2.6.4},可分解 $f = Bg$,$g$ 在 $D$ 内无零点.由定理~\ref{thm:2.5.3},$f \in L^{p}$ 且 $\log|f| \in L^{1}$.

令
\begin{equation*}
F(z) = \exp\left\{\frac{1}{2\pi}\int_{-\pi}^{\pi}\frac{e^{it} + z}{e^{it} - z}\log|f(e^{it})|\dd t\right\},
\end{equation*}
则 $|g(e^{i\theta})| = |F(e^{i\theta})|$ a.e. $\theta$.因此令 $g(z)/F(z) = Q(z)$,$Q$ 是在 $D$ 内无零点的内函数,由定理~\ref{thm:2.8.2} 的 (1),$Q(z) = e^{i\gamma}S(z)$.把 $F$ 换成 $e^{i\gamma}F$,定理得证.
\end{proof}

% ============ 第2章 §2.9 H^1与积分表示(原书 pp.147-149) ============
\section{\texorpdfstring{$H^1$ 与积分表示}{H1 与积分表示}}\label{sec:2.9}

2.4 节证明了属于 $h^{1}$ 的函数可以用 Poisson--Stieltjes 积分表示来刻画.本节证明属于 $\operatorname{Re}H^{1}$ 的函数可以用 Poisson 积分表示,即 $\mu$ 关于 Lebesgue 测度绝对连续.因此 $\operatorname{Re}H^{1} \subset h^{1}$.

\begin{theorem}\label{thm:2.9.1}
若 $f \in H^{1}$,则
\begin{equation*}
f(z) = \frac{1}{2\pi}\int_{-\pi}^{\pi}P(r, \theta - t)f(e^{it})\dd t.
\end{equation*}
\end{theorem}

\begin{proof}
当 $f \in H^{1}$ 时,若 $0 < \rho < 1$,则
\begin{equation*}
f(\rho z) = \frac{1}{2\pi}\int_{-\pi}^{\pi}P(r, \theta - t)f(\rho e^{it})\dd t,
\end{equation*}
且由定理 2.7.2,
\begin{equation*}
\int_{-\pi}^{\pi}|f(\rho e^{it}) - f(e^{it})|\dd t \rightarrow 0 \quad (\rho \rightarrow 1).
\end{equation*}
令
\begin{equation*}
F(z) = \frac{1}{2\pi}\int_{-\pi}^{\pi}P(r, \theta - t)f(e^{it})\dd t,
\end{equation*}
则 $f(\rho z) \to F(z)$,故 $F(z) = f(z)$.
\end{proof}

\begin{corollary}\label{cor:2.9.2}
若 $u \in \operatorname{Re}H^{1}$,则 $u$ 可以用 Poisson 积分表示.
\end{corollary}

对 $\phi \in L^{1}$,
\begin{equation*}
\hat{\phi}(n) = \frac{1}{2\pi}\int_{-\pi}^{\pi}\phi(t)e^{-int}\dd t \quad (n = 0, \pm 1, \pm 2, \ldots)
\end{equation*}
称为 $\phi$ 的 Fourier 系数.$\phi \in L^{1}$ 的 Poisson 积分在 $D$ 上调和,但何时正则?定理~\ref{thm:2.9.1} 表明当 $\phi$ 是 $H^{1}$ 的函数的边界值时正是如此.下面的命题~\ref{prop:2.9.3} 用 $\phi$ 的 Fourier 系数的语言回答这个问题.这将在下面的 2.10 节中用于 $H^{1}$ 的边界值的刻画.

\begin{proposition}\label{prop:2.9.3}
$\phi \in L^{1}$ 且 $\hat{\phi}(n) = 0$($n < 0$)时,令
\begin{equation*}
f(z) = \frac{1}{2\pi}\int_{-\pi}^{\pi}P(r, \theta - t)\phi(e^{it})\dd t \quad (z = re^{i\theta}),
\end{equation*}
则 $f \in H^{1}$ 且 $f(e^{i\theta}) = \phi(e^{i\theta})$ a.e. $\theta$.
\end{proposition}

\begin{proof}
由 $\hat{\phi}(n) = 0$($n < 0$)与 Poisson 核的展开
\begin{equation*}
P(r, \theta - t) = 1 + \sum_{n = 1}^{\infty}\left(z^{n}e^{-int} + \bar{z}^{n}e^{int}\right),
\end{equation*}
可以证明 $f(z) = \sum_{n=0}^{\infty}\hat{\phi}(n)z^{n}$($z \in D$).$f \in H(D)$ 由 $\phi \in L^{1}$ 是显然的,但由命题~\ref{prop:2.2.2} 的 (3),
\begin{equation*}
M_{1}(r, f) \leq \frac{1}{2\pi}\int_{-\pi}^{\pi}|\phi(e^{it})|\dd t,
\end{equation*}
所以 $f \in H^{1}$.由定理~\ref{thm:2.3.2},$f(e^{i\theta}) = \phi(e^{i\theta})$ a.e. $\theta$.
\end{proof}

若 $f \in H^{1}$,则
\begin{equation*}
f(z) = \frac{1}{2\pi}\int_{-\pi}^{\pi}C(r, \theta - t)f(e^{it})\dd t.
\end{equation*}
这可以与定理~\ref{thm:2.9.1} 同样地证明,但 2.12 节将给出另一个证明.

% ============ 第2章 §2.10 Hardy空间的边界值 (III)(原书 pp.149-153) ============
\section{Hardy 空间的边界值 (III)}\label{sec:2.10}

对 $0 < p \leq \infty$,令
\begin{equation*}
H_{*}^{p} = \left\{f(e^{i\theta}) = \lim_{r \to 1}f(re^{i\theta}) \text{ a.e. } \theta;\ f \in H^{p}\right\}.
\end{equation*}
当 $f, g \in H_{*}^{p}$ 且 $f(e^{i\theta}) = g(e^{i\theta})$ a.e. $\theta$ 时,把 $f$ 与 $g$ 视为同一,则由定理~\ref{thm:2.5.3},$H_{*}^{p} \subset L^{p}$.此时,由定理~\ref{thm:2.7.1},$H^{p}$ 与 $H_{*}^{p}$ 可以等同看待.本节证明 $H_{*}^{p}$ 是 $L^{p}$ 的闭子空间,并由此导出 $H^{p}$ 完备.

对 $f \in L^{p}$,令
\begin{equation*}
\|f\|_{p}^{p} = \int_{-\pi}^{\pi}|f(t)|^{p}\frac{\dd t}{2\pi} \quad (p \neq \infty)
\end{equation*}
且 $\|f\|_{\infty} = \esssup_{-\pi \leq t \leq \pi}|f(t)|$.当 $1 \leq p \leq \infty$ 时 $\|f\|_{p}$ 是范数,$L^{p}$ 是 Banach 空间.当 $0 < p < 1$ 时,令 $d(f, g) = \|f - g\|_{p}^{p}$,则是距离,$L^{p}$ 是完备度量空间.

本节证明当 $1 \leq p \leq \infty$ 时 $H_{*}^{p} = \{f \in L^{p};\ \hat{f}(n) = 0 \ (n < 0)\}$.

\begin{lemma}\label{lem:2.10.1}
当 $0 < p \leq \infty$ 时,若 $f \in H^{p}$,则
\begin{equation*}
|f(z)| \leq 2^{1/p}\|f\|_{p}(1 - r)^{-1/p} \quad (r = |z|).
\end{equation*}
\end{lemma}

\begin{proof}
当 $p = 1$ 时,由定理~\ref{thm:2.9.1},
\begin{equation*}
|f(z)| \leq \frac{1}{2\pi}\int_{-\pi}^{\pi}P(r, \theta - t)|f(e^{it})|\dd t \leq \frac{1 + r}{1 - r}\|f\|_{1},
\end{equation*}
故显然.当 $p = \infty$ 时由定理~\ref{thm:2.7.1} 显然,故设 $0 < p < \infty$.由定理~\ref{thm:2.6.4} 可分解 $f = Bg$,$|f(z)| \leq |g(z)|$ 且 $\|f\|_{p} = \|g\|_{p}$.由于 $[g(z)]^{p} \in H^{1}$,由 $p = 1$ 的情形,
\begin{equation*}
|g(z)|^{p} \leq \frac{1 + r}{1 - r}\|g^{p}\|_{1} = \frac{1 + r}{1 - r}\|f\|_{p}^{p},
\end{equation*}
引理得证.
\end{proof}

\begin{theorem}\label{thm:2.10.2}
当 $0 < p \leq \infty$ 时,$H_{*}^{p}$ 是 $L^{p}$ 的闭子空间.当 $p \neq \infty$ 时,$e^{i\theta}$ 的多项式的全体在 $H_{*}^{p}$ 中稠密.
\end{theorem}

\begin{proof}
当 $p \neq \infty$ 时,证明 $e^{i\theta}$ 的多项式的全体在 $H_{*}^{p}$ 中稠密.设 $f \in H_{*}^{p}$ 且 $\varepsilon > 0$,由 $H_{*}^{p}$ 的定义可以把 $f$ 看作属于 $H^{p}$,故由定理~\ref{thm:2.7.3},存在 $0 < \rho < 1$ 使得
\begin{equation*}
\|f_{\rho} - f\|_{p} < \varepsilon.
\end{equation*}
这里 $f_{\rho}(z) = f(\rho z)$.设 $s_{n}(z)$ 为 $f(z)$ 在 $z = 0$ 处 Taylor 展开的第 $n$ 个部分和.此时,由于 $s_{n}(z)$ 在 $|z| = \rho$ 上一致收敛于 $f(z)$,对足够大的 $n$,
\begin{equation*}
\|s_{n\rho} - f_{\rho}\|_{p} < \varepsilon.
\end{equation*}
这里 $s_{n\rho}(z) = s_{n}(\rho z)$.因此,利用当 $0 < p < \infty$ 时
\begin{equation*}
(a + b)^{p} \leq 2^{p}(a^{p} + b^{p}) \quad (a \geq 0, b \geq 0)
\end{equation*}
成立,可以证明 $e^{i\theta}$ 的多项式在 $H_{*}^{p}$ 中稠密.

证明 $H_{*}^{p}$ 是 $L^{p}$ 的闭子空间.设 $f_{n} \in H_{*}^{p}$,$\phi \in L^{p}$ 且 $\|f_{n} - \phi\|_{p} \to 0$($n \to \infty$),则 $\{f_{n}\}$ 在 $L^{p}$ 中有界,故由引理~\ref{lem:2.10.1},对任意 $\rho < 1$,$\{f_{n}(z)\}$ 在 $|z| \leq \rho$ 上一致有界.$\{f_{n}\}$ 在 $|z| \leq \rho$ 上是正规族,故存在其子列 $\{f_{n_{j}}\}$ 使得 $f_{n_{j}}(z)$ 在 $|z| \leq \rho$ 上一致收敛于正则函数 $f(z)$.由一致定理可证 $f$ 与 $\rho$ 无关,且由 $\{f_{n}\}$ 在 $L^{p}$ 中有界知 $f \in H^{p}$.证明 $\phi(\theta) = f(e^{i\theta})$ a.e. $\theta$ 后证明结束.对任意 $\varepsilon > 0$,存在 $N$ 使得当 $n, m \geq N$ 时 $\|f_{n} - f_{m}\|_{p} < \varepsilon$.因此,当 $m \geq N$ 且 $r < 1$ 时,
\begin{align*}
M_{p}(r, f - f_{m}) & = \lim_{j \to \infty}M_{p}(r, f_{n_{j}} - f_{m}) \\
& \leq \lim_{j \to \infty}\|f_{n_{j}} - f_{m}\|_{p} \leq \varepsilon.
\end{align*}
由定理~\ref{thm:2.7.1},当 $m \geq N$ 时 $\|f - f_{m}\|_{p} \leq \varepsilon$.因此 $\|f - f_{n}\|_{p} \to 0$ 且 $\phi(\theta) = f(e^{i\theta})$ a.e. $\theta$.
\end{proof}

下面两个系是定理~\ref{thm:2.7.1} 与上述定理的结果.

\begin{corollary}\label{cor:2.10.3}
当 $1 \leq p \leq \infty$ 时,$H^{p}$ 是 Banach 空间.
\end{corollary}

\begin{corollary}\label{cor:2.10.4}
当 $0 < p < 1$ 时,$H^{p}$ 是完备度量空间.
\end{corollary}

当 $f(z) = \sum_{n=0}^{\infty}a_{n}z^{n} \in H^{2}$ 时,令 $F(e^{i\theta}) = \sum_{n=0}^{\infty}a_{n}e^{in\theta}$,则 $F \in L^{2}$,但更一般地证明 $F(e^{i\theta}) = f(e^{i\theta})$ a.e. $\theta$,即 $\hat{f}(n) = a_{n}$($n \geq 0$).

\begin{theorem}\label{thm:2.10.5}
设 $1 \leq p \leq \infty$.若 $f(z) = \sum_{n=0}^{\infty}a_{n}z^{n} \in H^{p}$,则 $\hat{f}(n) = a_{n}$($n \geq 0$)且 $\hat{f}(n) = 0$($n < 0$).这里 $\hat{f}(n)$ 是 $f(e^{i\theta})$ 的 Fourier 系数.更进一步,
\begin{equation*}
H_{*}^{p} = \left\{\phi \in L^{p};\ \hat{\phi}(n) = 0 \ (n = -1, -2, \ldots)\right\}.
\end{equation*}
\end{theorem}

\begin{proof}
设 $f \in H^{p}$,则当 $0 < r < 1$ 时
\begin{equation*}
\hat{f}_{r}(n) = \frac{1}{2\pi}\int_{-\pi}^{\pi}e^{-int}f(re^{it})\dd t = \begin{cases} r^{n}a_{n} & (n \geq 0) \\ 0 & (n < 0) \end{cases},
\end{equation*}
且 $|\hat{f}_{r}(n) - \hat{f}(n)| \leq \|f_{r} - f\|_{1}$.因此,由定理~\ref{thm:2.7.3},令 $r \to 1$,$\hat{f}(n) = \lim_{r \to 1}\hat{f}_{r}(n)$,定理的前半得证.对后半,由前半只需证明 $H_{*}^{p} \supset \{\phi \in L^{p};\ \hat{\phi}(n) = 0 \ (n < 0)\}$.这由命题~\ref{prop:2.9.3},当 $p = 1$ 时已证明.当 $p \neq 1$ 时也类似.
\end{proof}

在 $H^{p}$ 的边界值的全体 $H_{*}^{p}$ 中,$H_{*}^{1}$ 不包含除常数外的实值函数,由定理~\ref{thm:2.10.5} 是显然的.当 $0 < p < 1$ 时,$H_{*}^{p}$ 中存在除常数外的实值函数.事实上,令 $f(z) = i(1 + z)(1 - z)^{-1}$,则 $f \in \bigcap_{p < 1}H^{p}$,且 $f(e^{i\theta}) = -2\sin\theta|1 - e^{i\theta}|^{-2} \in \bigcap_{p < 1}H_{*}^{p}$.另一方面,$H_{*}^{1/2}$ 不包含除常数外的非负函数 \cite{ch2:1}.下面的命题为了在 2.16 节使用而更一般地叙述 \cite{ch2:2}.

\begin{proposition}\label{prop:2.10.6}
若 $f \in e^{-in\theta}H_{*}^{1/2}$,$n \geq 0$,在 $\partial D$ 上非负,则 $f$ 是次数至多 $n$ 的三角多项式.
\end{proposition}

\begin{proof}
属于 $e^{-in\theta}H_{*}^{1}$ 的非负值函数是次数至多 $n$ 的三角多项式是显然的.设 $f \in e^{-in\theta}H_{*}^{1/2}$ 是非负值函数,为证明命题只需证明 $f \in e^{-in\theta}H_{*}^{1}$.设 $f(e^{i\theta}) = e^{-in\theta}g(e^{i\theta})$ 且 $g \in H_{*}^{1/2}$,由定理~\ref{thm:2.8.3},$g = QG$ 可唯一分解.这里 $Q$ 是内函数,$G$ 是外函数.令 $g_{1} = QG^{1/2}$,$g_{2} = G^{1/2}$,则 $g = g_{1}\cdot g_{2}$,且 $g_{j} \in H_{*}^{1}$($j = 1, 2$).此时 $|g_{1}| = |g_{2}|$.由于 $g_{1}\cdot e^{-in\theta}g_{2} = |g_{1}|^{2}$,故 $e^{-in\theta}g_{2} = \bar{g}_{1}$.因此 $g_{1} + e^{-in\theta}g_{2}$ 与 $i(g_{1} - e^{-in\theta}g_{2})$ 是实值函数.

当 $0 < r < 1$ 时,设 $A_{r} = \{z \in \mathbb{C};\ r < |z| < 1\}$.若 $F$ 在 $A_{r}$ 上正则,
\begin{equation*}
\sup_{r \leq t < 1}\int_{-\pi}^{\pi}|F(te^{i\theta})|\dd\theta < \infty,
\end{equation*}
且 $F$ 在 $\partial D$ 上实值,则 $F$ 可以解析延拓越过 $\partial D$.这个结果超出本书的范围,但众所周知,其证明见 \cite{ch2:3}.

令
\begin{equation*}
F_{1}(z) = g_{1}(z) + \frac{g_{2}(z)}{z^{n}}, \quad F_{2}(z) = i\left(g_{1}(z) - \frac{g_{2}(z)}{z^{n}}\right),
\end{equation*}
则 $F_{1}$ 与 $F_{2}$ 在 $\partial D$ 上正则.因此 $g_{1}$ 在 $\partial D$ 上正则,且由 $g_{1} \in H_{*}^{\infty}$ 得 $g \in H_{*}^{\infty}$.这表明 $f$ 属于 $e^{-in\theta}H_{*}^{1}$.
\end{proof}

% ============ 第2章 §2.11 H^p (0<p<1)与积分表示(原书 pp.153-154) ============
\section{\texorpdfstring{$H^p$($0<p<1$)与积分表示}{Hp(0<p<1)与积分表示}}\label{sec:2.11}
\index{Fourier 系数}


$(1+z)(1-z)^{-1}$ 属于 $\bigcap_{p<1} H^{p}$ 这一事实很容易证明.定理~\ref{thm:2.11.2} 利用引理~\ref{lem:2.11.1} 推广了这一结果.

\begin{lemma}\label{lem:2.11.1}
设 $G(z)$ 在 $D$ 上实值且次调和.若 $\phi \in H^{\infty}$,$\|\phi\|_{\infty} \leq 1$ 且 $\phi(0) = 0$,则
\begin{equation*}
\int_{-\pi}^{\pi} G(\phi(re^{i\theta}))\dd \theta \leq \int_{-\pi}^{\pi} G(re^{i\theta})\dd \theta \quad (0 < r < 1).
\end{equation*}
\end{lemma}

\begin{proof}
设 $U(z)$ 在 $|z| < r$ 上调和且在 $|z| = r$ 上 $U(z) = G(z)$.因 $G(z)$ 次调和,在 $|z| \leq r$ 上有 $G(z) \leq U(z)$.由 Schwarz 引理,$|\phi(z)| \leq |z|$,故 $G(\phi(z)) \leq U(\phi(z))$,因此
\begin{align*}
\frac{1}{2\pi}\int_{-\pi}^{\pi} G(\phi(re^{i\theta}))\dd \theta & \leq \frac{1}{2\pi}\int_{-\pi}^{\pi} U(\phi(re^{i\theta}))\dd \theta \\
& = U(\phi(0)) = U(0) = \frac{1}{2\pi}\int_{-\pi}^{\pi} U(re^{i\theta})\dd \theta \\
& = \frac{1}{2\pi}\int_{-\pi}^{\pi} G(re^{i\theta})\dd \theta .
\end{align*}
\end{proof}

\begin{theorem}\label{thm:2.11.2}
若 $f \in H(D)$ 且 $\operatorname{Re} f(z) > 0$($z \in D$),则 $f \in \bigcap_{p<1} H^{p}$.
\end{theorem}

\begin{proof}
因 $\operatorname{Re} f(z) > 0$($z \in D$),可设 $f(0) = 1$.令 $\phi(z) = (f(z)-1)/(f(z)+1)$,则 $\phi \in H^{\infty}$,$\|\phi\|_{\infty} \leq 1$ 且 $\phi(0) = 0$.$f(z) = (1+\phi(z))/(1-\phi(z))$,令
\begin{equation*}
G(z) = \left| \frac{1+z}{1-z} \right|^{p},
\end{equation*}
则 $G(z)$ 次调和且 $|f(z)|^{p} = G(\phi(z))$,故由引理~\ref{lem:2.11.1} 可证 $f \in \bigcap_{p<1} H^{p}$.
\end{proof}

设 $F$ 为复测度 $\mu \in M[-\pi, \pi]$ 的 Cauchy--Stieltjes 积分,即
\begin{equation*}
F(z) = \frac{1}{2\pi}\int_{-\pi}^{\pi} C(r, \theta - t)\dd \mu(t) = \frac{1}{2\pi}\int_{-\pi}^{\pi} \frac{e^{it}}{e^{it} - z}\dd \mu(t),
\end{equation*}
则 $F(z)$ 在 $|z| \neq 1$ 时正则.

\begin{theorem}\label{thm:2.11.3}
若 $F$ 为复测度 $\mu \in M[-\pi, \pi]$ 的 Cauchy--Stieltjes 积分,则 $F \in \bigcap_{p<1} H^{p}$.
\end{theorem}

\begin{proof}
可设 $\mu$ 为正测度,此时由命题~\ref{prop:2.2.2} 有 $\operatorname{Re} F(z) > 0$.由定理~\ref{thm:2.11.2},$F \in \bigcap_{p<1} H^{p}$.
\end{proof}

已知在定理~\ref{thm:2.11.3} 中,其逆不成立,而且也不能把 $\bigcap_{p<1} H^{p}$ 换成 $H^{1}$.

% ============ 第2章 §2.12 Riesz兄弟定理(原书 pp.154-155) ============
\section{Riesz 兄弟定理}\label{sec:2.12}
\index{Riesz 兄弟定理}


本节考察 Cauchy--Stieltjes 积分何时属于 $H^{1}$.这与测度论意义上令人惊奇的 Riesz 兄弟(F. and M. Riesz)定理密切相关.在 2.11 节中我们注意到,属于 $\bigcap_{p<1} H^{p}$ 的函数未必有 Poisson--Stieltjes 积分表示,而定理 2.12.1 表明对于 $H^{1}$ 这是可能的.本节中 $\hat{\mu}$ 表示 $\mu$ 的 Fourier 系数.

\begin{lemma}\label{lem:2.12.1}
设 $\mu \in M[-\pi,\pi]$ 为复测度且 $\hat{\mu}(n) = 0$($n < 0$),则 $\mu$ 的 Poisson--Stieltjes 积分属于 $H^{1}$,而且实际上是 Poisson 积分.
\end{lemma}

\begin{proof}
设 $F(z)$ 为 $\mu$ 的 Cauchy--Stieltjes 积分,则当 $|z| < 1$ 时
\begin{equation*}
\frac{1}{1 - ze^{-it}} = \sum_{n=0}^{\infty} e^{-int}z^{n}
\end{equation*}
且
\begin{equation*}
\frac{e^{it}}{e^{it} - \frac{1}{\bar{z}}} = -\sum_{n=0}^{\infty} e^{i(n+1)t}\bar{z}^{n+1},
\end{equation*}
所以 $F(z) = \sum_{n=0}^{\infty}\hat{\mu}(n)z^{n}$,并且
\begin{equation*}
F\left(\frac{1}{\bar{z}}\right) = -\sum_{n=0}^{\infty}\hat{\mu}(-(n+1))\left(\frac{1}{\bar{z}}\right)^{n+1} \equiv 0.
\end{equation*}
于是当 $|z| < 1$ 时成立
\begin{equation*}
F(z) = F(z) - F\left(\frac{1}{\bar{z}}\right) = \frac{1}{2\pi}\int_{-\pi}^{\pi} P(r, \theta - t)\dd \mu(t),
\end{equation*}
由定理~\ref{thm:2.3.1} 的证明知 $F \in H^{1}$.
\end{proof}

\begin{theorem}[Riesz 兄弟]\label{thm:2.12.2}
若 $\mu \in M[-\pi, \pi]$ 为满足 $\mu(\{-\pi\}) = \mu(\{\pi\})$ 的复测度且 $\hat{\mu}(n) = 0$($n < 0$),则 $\mu$ 关于 Lebesgue 测度绝对连续.
\end{theorem}

\begin{proof}
由引理~\ref{lem:2.12.1},$\mu$ 的 Poisson--Stieltjes 积分 $f(z)$ 属于 $H^{1}$.由定理~\ref{thm:2.9.1},$f(z)$ 可由 $f(e^{it})$ 作 Poisson 积分表示.由 Poisson 积分表示的唯一性(命题~\ref{prop:2.2.3}),$\dd \mu(t) = f(e^{it})\dd t$.
\end{proof}

\begin{proposition}\label{prop:2.12.3}
$f \in H^{1}$ 可由 $f$ 的边界值 $f(e^{it})$ 作 Cauchy 积分表示.
\end{proposition}

\begin{proof}
设 $F$ 为 $f(e^{it})$ 的 Cauchy 积分.令 $\dd \mu = f(e^{it})\dd t$,则由定理~\ref{thm:2.10.5} 有 $\hat{\mu}(n) = 0$($n < 0$).由引理~\ref{lem:2.12.1},$F \in H^{1}$ 且 $F$ 是 Poisson 积分,并由定理~\ref{thm:2.9.1} 等于 $f$.即
\begin{align*}
F(z) & = \frac{1}{2\pi}\int_{-\pi}^{\pi} \frac{e^{it}}{e^{it} - z} f(e^{it})\dd t \\
& = \frac{1}{2\pi}\int_{-\pi}^{\pi} P(r, \theta - t) f(e^{it})\dd t = f(z).
\end{align*}
\end{proof}

% ============ 第2章 §2.13 有界线性泛函(原书 pp.155-158) ============
\section{有界线性泛函}\label{sec:2.13}
\index{圆板代数}


当 $0 < p \leq \infty$ 时,对固定的 $a \in D$,将 $f(a)$ 对应于 $f \in H^{p}$ 的映射,由引理~\ref{lem:2.10.1} 是有界线性泛函.本节考察 $H^{p}$ 上有界线性泛函的全体 $(H^{p})^{*}$.

由于 $H^{p}$ 与其边界值的全体 $H_{*}^{p}$ 由定理~\ref{thm:2.7.1} 与定理~\ref{thm:2.10.5} 可视为同一,利用 $H^{p}$ 是 $L^{p}$ 的线性子空间这一事实来考察 $(H^{p})^{*}$.当 $1 \leq p < \infty$ 时,设 $1/p + 1/q = 1$,则 $(L^{p})^{*}$ 与 $L^{q}$ 等距同构,即 $(L^{p})^{*} = L^{q}$ 是众所周知的.另外当 $0 < p < 1$ 时,$(L^{p})^{*} = \{0\}$ 也是已知的.但如上所述,一般 $(H^{p})^{*} \neq \{0\}$.

设 $1 \leq p \leq \infty$,则因 $H^{p} \subset L^{p}$,把 $L^{p}$ 上的有界线性泛函限制到 $H^{p}$ 上即得 $(H^{p})^{*}$ 的元素.又由 Hahn--Banach 定理,$H^{p}$ 上的有界线性泛函可扩张到 $L^{p}$ 上,故得 $(L^{p})^{*} = L^{q}$($p \neq \infty$)的元素.将这一事实准确叙述即为定理~\ref{thm:2.13.1}.

\begin{theorem}\label{thm:2.13.1}
设 $1 \leq p < \infty$ 且 $1/p + 1/q = 1$.$(H^{p})^{*}$ 与 $L^{q}/zH^{q}$ 等距同构.即 $(H^{p})^{*} = L^{q}/zH^{q}$.
\end{theorem}

\begin{proof}
对 $g \in L^{q}$,令
\begin{equation*}
\Phi_{g}(k) = \int_{-\pi}^{\pi} k(e^{i\theta}) g(e^{i\theta})\dd \theta / 2\pi \quad (k \in L^{p}),
\end{equation*}
则 $\Phi_{g} \in (L^{p})^{*}$ 且 $\|\Phi_{g}\| = \|g\|_{q}$.令 $\phi = \Phi_{g}|H^{p}$,则 $\phi \in (H^{p})^{*}$.由定理~\ref{thm:2.10.5},
\begin{align*}
(H^{p})^{\perp} & = \left\{g \in L^{q}; \int_{-\pi}^{\pi} f(e^{i\theta}) g(e^{i\theta})\dd \theta / 2\pi = 0 \ (f \in H^{p}) \right\} \\
& = zH^{q},
\end{align*}
因此对任意 $h \in zH^{q}$,$\phi = \Phi_{g+h}|H^{p}$.于是 $L^{q}/zH^{q}$ 可缩小地嵌入 $(H^{p})^{*}$.等距性由下面的逆命题的证明得出.

反过来,设 $\phi \in (H^{p})^{*}$.由 Hahn--Banach 定理,存在 $\Phi \in (L^{p})^{*}$,使得 $\phi = \Phi|H^{p}$ 且 $\|\phi\| = \|\Phi\|$.由 $(L^{p})^{*} = L^{q}$,存在 $g \in L^{q}$,使得 $\Phi = \Phi_{g}$ 且 $\|\Phi\| = \|g\|_{q}$.对任意 $h \in zH^{q}$,$\Phi|H^{p} = \Phi_{g+h}|H^{p}$,故 $\|g\|_{q} = \inf \|g + h\|_{q}$.于是 $(H^{p})^{*}$ 可等距地嵌入 $L^{q}/zH^{q}$.
\end{proof}

当 $p = \infty$ 时,因 $(L^{\infty})^{*} \supsetneq L^{1}$,定理~\ref{thm:2.13.1} 不成立.但类似于定理~\ref{thm:2.13.1} 的证明,可以证明 $(H^{\infty})^{*} = (L^{\infty})^{*} / (H^{\infty})^{\perp}$.

\begin{theorem}\label{thm:2.13.2}
设 $1 < p \leq \infty$ 且 $1/p + 1/q = 1$.$H^{p}$ 与 $(L^{q}/zH^{q})^{*}$ 等距同构.即 $H^{p} = (L^{q}/zH^{q})^{*}$.
\end{theorem}

\begin{proof}
对 $g \in H^{p}$,令
\begin{equation*}
\Phi_{g}(k) = \int_{-\pi}^{\pi} k(e^{i\theta}) g(e^{i\theta})\dd \theta / 2\pi \quad (k \in L^{q}),
\end{equation*}
则 $\Phi_{g} \in (L^{q})^{*}$ 且 $\|\Phi_{g}\| = \|g\|_{p}$.因 $\Phi_{g}$ 在 $zH^{q} = (H^{p})^{\perp}$ 上为零,令 $\phi_{g}(k + zH^{q}) = \Phi_{g}(k)$,则 $\phi_{g} \in (L^{q}/zH^{q})^{*}$ 且 $\|\phi_{g}\| \leq \|g\|_{p}$.于是 $H^{p}$ 可缩小地嵌入 $(L^{q}/zH^{q})^{*}$.

反过来,设 $\phi \in (L^{q}/zH^{q})^{*}$,令
\begin{equation*}
\Phi(k) = \phi(k + zH^{q}) \quad (k \in L^{q}),
\end{equation*}
则 $\Phi \in (L^{q})^{*}$ 且 $\|\Phi\| \leq \|\phi\|$.由 $(L^{q})^{*} = L^{p}$,存在 $g \in L^{p}$,使得 $\Phi = \Phi_{g}$ 且 $\|\Phi\| = \|g\|_{p}$.因 $\Phi_{g}$ 在 $zH^{q}$ 上为零,$g \in (zH^{q})^{\perp} = H^{p}$.于是 $(L^{q}/zH^{q})^{*}$ 可缩小地嵌入 $H^{p}$.
\end{proof}

定理~\ref{thm:2.13.2} 在 p = 1 时不成立,但利用 Riesz 兄弟定理,$H^{1}$ 成为某个 Banach 空间的共轭空间.

令
\begin{equation*}
C = \{f \in C[-\pi, \pi]; f(-\pi) = f(\pi)\},
\end{equation*}
则 $C$ 可与 $D$ 的边界 $\partial D$ 上复值连续函数的全体视为同一.

令
\begin{equation*}
A = \{f \in C; \hat{f}(n) = 0 \ (n < 0)\},
\end{equation*}
则其与 $H(D)$ 中在 $\overline{D}$ 上连续的函数限制到 $\partial D$ 上所得者一致.$A$ 称为圆盘代数.

\begin{theorem}\label{thm:2.13.3}
$H^{1}$ 与 $(C/zA)^{*}$ 等距同构.即 $H^{1} = (C/zA)^{*}$.
\end{theorem}

\begin{proof}
因 $C^{*}$ 与 $\partial D$ 上(有界)复测度的全体 $M(\partial D)$ 一致,由定理~\ref{thm:2.12.2},
\begin{align*}
(zA)^{\perp} & = \{\mu \in M(\partial D); \hat{\mu}(n) = 0 \quad (n < 0)\} \\
& = H^{1},
\end{align*}
因此,定理可以与引理 2.13.2 的证明完全相同地证明.
\end{proof}

% ============ 第2章 §2.14 极值问题(原书 pp.158-161) ============
\section{极值问题}\label{sec:2.14}
\index{极值问题}
\index{极值函数}
\index{极值核}
\index{共轭关系}


当 $\phi \in (H^{p})^{*}$ 时,对
\begin{equation*}
\|\phi\| = \sup \{|\phi(f)|; f \in H^{p}, \|f\|_{p} \leq 1\}
\end{equation*}
的极值问题指的是下面的 (i)--(iv).(i) 上确界的值是什么?(ii) 上确界能否成为最大值?(iii) 给出最大值的函数(极值函数)是否唯一?(iv) 极值函数是什么样的函数?

一般的 $\phi \in (H^{p})^{*}$($1 \leq p < \infty$)可表示为
\begin{align*}
\phi(f) & = \frac{1}{2\pi}\int_{-\pi}^{\pi} f(e^{it}) e^{it} k(e^{it})\dd t \\
& = \frac{1}{2\pi i}\int_{|z|=1} f(z) k(z)\dd z \quad (f \in H^{p}),
\end{align*}
其中 $k \in L^{q}$ 且 $1/p + 1/q = 1$.$\phi \in (H^{\infty})^{*}$ 一般不能用 $k \in L^{1}$ 表示.这是因为 $(L^{\infty})^{*} \supsetneq L^{1}$.但显然用 $k \in L^{1}$ 表示的 $\phi$ 属于 $(H^{\infty})^{*}$.如 2.13 节所示,当 $1 \leq p < \infty$ 时 $(H^{p})^{*} = L^{q}/zH^{q}$,当 $p = \infty$ 时 $(H^{\infty})^{*} = (L^{1}/zH^{1})^{**} \supset L^{1}/zH^{1}$,因此若 $h - k \in H^{q}$,则 $h$ 与 $k$ 在 $(H^{p})^{*}$ 中给出同一元素.于是 $\|\phi\| \leq \|k + g\|_{q}$($g \in H^{q}$),但可以证明下面的共轭关系(定理~\ref{thm:2.14.2}),进而可以解 $1 < p < \infty$ 情形的极值问题 (ii).

\begin{lemma}\label{lem:2.14.1}
设 $X$ 为 Banach 空间,$S$ 为 $X$ 的闭子空间且 $S^{\perp} = \{\phi \in X^{*}; \phi(x) = 0 \ (x \in S)\}$.

(1) 若 $\phi \in X^{*}$,则
\begin{equation*}
\sup \{|\phi(x)|; x \in S, \|x\| \leq 1\} = \min \{\|\phi + \psi\|; \psi \in S^{\perp}\}.
\end{equation*}

(2) 若 $x \in X$,则
\begin{equation*}
\max \{|\psi(x)|; \psi \in S^{\perp}, \|\psi\| \leq 1\} = \inf \{\|x + y\|; y \in S\}.
\end{equation*}
\end{lemma}

\begin{proof}
(1) 定理~\ref{thm:2.13.1} 的证明是通用的,$S^{*} = X^{*}/S^{\perp}$ 成立.其证明即证明了 (1).(2) 定理~\ref{thm:2.13.2} 的证明是通用的,$S^{\perp} = (X/S)^{*}$ 成立.设 $\psi \in S^{\perp}$ 且 $\|\psi\| \leq 1$,则对任意 $x \in X$ 和任意 $y \in S$,
\begin{equation*}
|\psi(x)| = |\psi(x + y)| \leq \|x + y\|,
\end{equation*}
所以
\begin{equation*}
\sup \{|\psi(x)|; \psi \in S^{\perp}, \|\psi\| \leq 1\} \leq \|x + S\|.
\end{equation*}
由 Hahn--Banach 定理的推论 [4,引理 30.1],存在 $\Psi_{0} \in (X/S)^{*}$,使得 $\|\Psi_{0}(x+S)\| = \|x+S\|$ 且 $\|\Psi_{0}\| = 1$.令 $\psi_{0}(x) = \Psi_{0}(x+S)$,则 $\psi_{0} \in S^{\perp}$ 且 $\|\psi_{0}\| = \|\Psi_{0}\| = 1$,故 (2) 得证.
\end{proof}

\begin{theorem}\label{thm:2.14.2}
设 $1 \leq p \leq \infty$ 且 $1/p + 1/q = 1$.若 $k \in L^{q}$,则
\begin{align*}
\sup \left\{\left| \frac{1}{2\pi i}\int_{|z|=1} f(z) k(z)\dd z \right|; f \in H^{p}, \|f\|_{p} \leq 1 \right\}&= \min \{\|k - g\|_{q}; g \in H^{q}\}.
 \end{align*}
$p \neq 1$ 时 sup 可改为 max.
\end{theorem}

\begin{proof}
当 $1 < p \leq \infty$ 时,取 $X = L^{q}$,$S^{\perp} = zH^{p}$ 且 $S = H^{q}$,用引理~\ref{lem:2.14.1} 的 (2) 得
\begin{align*}
\max \left\{\left| \frac{1}{2\pi i}\int_{|z|=1} f(z) k(z)\dd z \right|; f \in H^{p}, \|f\|_{p} \leq 1 \right\}&= \inf \{\|k - g\|_{q}; g \in H^{q}\}.
 \end{align*}
另外,取 $X = L^{p}$,$S^{\perp} = H^{q}$ 且 $S = zH^{p}$,用引理~\ref{lem:2.14.1} 的 (1),inf 可改为 min.当 $p = 1$ 时,取 $X = L^{1}$,$S = zH^{1}$ 且 $S^{\perp} = H^{\infty}$,用引理~\ref{lem:2.14.1} 的 (1) 即可.
\end{proof}

共轭关系中给出最小值的函数称为极值核.

\begin{lemma}\label{lem:2.14.3}
设 $1 \leq p \leq \infty$ 且 $1/p + 1/q = 1$.$F \in H^{p}$ 为 $\phi \in (H^{p})^{*}$ 的极值函数且 $\phi(F) > 0$,其共轭关系的极值核为 $K$ 的充要条件如下:
\begin{equation*}
e^{i\theta} F(e^{i\theta}) K(e^{i\theta}) \geq 0 \quad \text{a.e. } \theta,
\end{equation*}
并且,若 $p = 1$,则 a.e. $\theta$ 有 $|K(e^{i\theta})| = \|K\|_{\infty}$;若 $1 < p < \infty$,则 a.e. $\theta$ 有 $|F(e^{i\theta})|^{p} = \|K\|_{q}^{-q}|K(e^{i\theta})|^{q}$;若 $p = \infty$,则在 $\{\theta; K(e^{i\theta}) \neq 0\}$ 上 a.e. $\theta$ 有 $|F(e^{i\theta})| = 1$.
\end{lemma}

\begin{proof}
由定理~\ref{thm:2.14.2},
\begin{equation*}
\frac{1}{2\pi}\int_{-\pi}^{\pi} F(e^{i\theta}) e^{i\theta} K(e^{i\theta})\dd \theta = \|K\|_{q}\|F\|_{p}.
\end{equation*}
由 Hölder 不等式,
\begin{equation*}
\frac{1}{2\pi}\int_{-\pi}^{\pi} e^{i\theta} F(e^{i\theta}) K(e^{i\theta})\dd \theta = \frac{1}{2\pi}\int_{-\pi}^{\pi} |e^{i\theta} F(e^{i\theta}) K(e^{i\theta})|\dd \theta,
\end{equation*}
故 $e^{i\theta}F(e^{i\theta})K(e^{i\theta}) \geq 0$ a.e. $\theta$.又因 Hölder 不等式中等式成立,后半部分关于 $|K|$ 与 $|F|$ 的关系式得证.
\end{proof}

\begin{theorem}\label{thm:2.14.4}
设 $1 \leq p \leq \infty$ 且 $\phi \in (H^{p})^{*}$.

(1) 当 $1 < p < \infty$ 时,$H^{p}$ 的极值问题的解存在,且除常数倍外唯一.

(2) 当 $p = \infty$ 时,若 $\phi \in L^{1}/zH^{1}$,则 $H^{\infty}$ 的极值问题的解存在,且除常数倍外唯一.

(3) 当 p = 1 时,若 $\phi \in C/zA$,则 $H^{1}$ 的极值问题的解存在,但不一定除常数倍外唯一.
\end{theorem}

\begin{proof}
(1) 与 (2) 中,解的存在性是定理~\ref{thm:2.14.2} 的结果,解的唯一性由引理~\ref{lem:2.14.3} 得出.(3) 中,若 $\phi \in C/zA$,则 $\phi$ 是 *弱连续的,$H^{1}$ 的单位球是 *弱紧的,故解存在.
\end{proof}

若
\begin{equation*}
\phi(f) = \frac{1}{2\pi i}\int_{|z|=1} f(z)\bar{z}^{2}\dd z,
\end{equation*}
则解有无穷多个.

% ============ 第2章 §2.15 端点与暴露点(原书 pp.161-165) ============
\section{端点与暴露点}\label{sec:2.15}
\index{端点}
\index{暴露点}


在 $H^{p}$($1 \leq p \leq \infty$)的单位球 $S^{p}$ 中,作为 $\phi \in (H^{p})^{*}$ 的极值问题除常数倍外唯一解的 $f \in S^{p}$ 称为露点(exposed point).$S^{p}$ 的露点全部属于 $S^{p}$ 的边界.由引理~\ref{lem:2.14.3} 与定理~\ref{thm:2.14.4},当 $1 < p < \infty$ 时,$S^{p}$ 的边界上的函数全部是露点.p = 1,$\infty$ 时,露点是 $S^{p}$ 中非常特殊的函数.在 $S^{p}$ 中,不是 $S^{p}$ 的两个不同函数的真正凸组合的 $f \in S^{p}$ 称为端点(extreme point).$S^{p}$ 的端点全部属于 $S^{p}$ 的边界.当 $1 < p < \infty$ 时逆成立是 Minkowski 不等式的结果.本节考察 p = 1,$\infty$ 情形下的端点与露点.容易证明露点总是端点.在 $1 < p < \infty$ 时 $S^{p}$ 的露点与端点同 $S^{p}$ 的边界一致,但我们将证明在 p = 1,$\infty$ 时存在是端点但不是露点的边界点.

\begin{lemma}\label{lem:2.15.1}
设 p = 1 或 $\infty$.f 不是 $S^{p}$ 的端点的充要条件是,存在非零的 $g \in S^{p}$,使得 $\|f \pm g\|_{p} \leq 1$.
\end{lemma}

\begin{proof}
若 f 不是端点,则存在 0 < a < 1 与 $h, k \in S^{p}$($h \not\equiv k$),使得 $f = ah + (1-a)k$.取 $\varepsilon$ 使 $0 < a \pm \varepsilon < 1$.令 $\ell = f + \varepsilon(h - k)$ 且 $m = f - \varepsilon(h - k)$,则 $\ell$ 与 $m$ 属于 $S^{p}$.令 $g = (\ell - m)/2$,则 $g \not\equiv 0$,$f + g = \ell$ 且 $f - g = m$.逆命题由 $f = (f + g)/2 + (f - g)/2$ 显然.这个证明对 $1 < p < \infty$ 也成立.
\end{proof}

\begin{lemma}\label{lem:2.15.2}
设 $g, f \in H^{1}$ 且 a.e. t 有 $|g(e^{it})| \leq |f(e^{it})|$.若 f 是外函数,则 $|g(z)| \leq |f(z)|$($z \in D$).
\end{lemma}

\begin{proof}
由定理~\ref{thm:2.8.1},
\begin{equation*}
\log|g(re^{i\theta})| \leq \frac{1}{2\pi}\int_{-\pi}^{\pi} P(r, \theta - t)\log|g(e^{it})|\dd t,
\end{equation*}
又因 f 是外函数,
\begin{equation*}
\log|f(re^{i\theta})| = \frac{1}{2\pi}\int_{-\pi}^{\pi} P(r, \theta - t)\log|f(e^{it})|\dd t.
\end{equation*}
这证明了 $|g(re^{i\theta})| \leq |f(re^{i\theta})|$($0 \leq r < 1, -\pi \leq \theta < \pi$).
\end{proof}

\begin{theorem}\label{thm:2.15.3}
设 $f \in H^{1}$.f 是 $S^{1}$ 的端点的充要条件是 f 满足 $\|f\|_{1} = 1$ 且为外函数.
\end{theorem}

\begin{proof}
设 $f \in S^{1}$ 为外函数.若 f 不是端点,由引理~\ref{lem:2.15.1},存在非零 $g \in S^{1}$,使得 $\|f \pm g\|_{1} \leq 1$.由 $\|f\|_{1} = 1$ 得 $\|f + g\|_{1} = \|f - g\|_{1} = 1$.于是令 k = g/f,
\begin{equation*}
\int_{-\pi}^{\pi} \left\{|1 + k(e^{i\theta})| + |1 - k(e^{i\theta})| - 2\right\} |f(e^{i\theta})|\dd \theta \leq 0.
\end{equation*}
因 a.e. $\theta$ 有 $|1 + k(e^{i\theta})| + |1 - k(e^{i\theta})| \geq 2$,故 a.e. $\theta$ 有 $|1 + k(e^{i\theta})| + |1 - k(e^{i\theta})| = 2$,从而 a.e. $\theta$ 有 $-1 \leq k(e^{i\theta}) \leq 1$.于是 a.e. $\theta$ 有 $|g(e^{i\theta})| \leq |f(e^{i\theta})|$,由引理~\ref{lem:2.15.2},k 属于 $H^{\infty}$.由定理~\ref{thm:2.10.5},$\hat{k}(n) = 0$($n = \pm 1, \pm 2, \ldots$),由 Fourier 系数的唯一性,k 为常数,但由 $|1 + k(e^{i\theta})| = |1 - k(e^{i\theta})|$ 得 $k \equiv 0$.于是 $g \equiv 0$.这与假设矛盾,故 f 是端点.

设 f 为 $S^{1}$ 的端点.此时 $\|f\|_{1} = 1$.若 f 不是外函数,由定理~\ref{thm:2.8.3},f = qh 可作因式分解,其中 q 是非常值的内函数,h 是外函数.对实常数 $\alpha$,
\begin{equation*}
g = \frac{1}{2}(e^{i\alpha}q + e^{-i\alpha}\bar{q})f,
\end{equation*}
则 $g \in H^{1}$.因 q 不是常数,$g \not\equiv 0$.选取 $\alpha$ 使得
\begin{equation*}
\int_{-\pi}^{\pi} \mathrm{Re}\left\{e^{i\alpha}q(e^{i\theta})\right\} |f(e^{i\theta})|\dd \theta = 0,
\end{equation*}
则
\begin{equation*}
\|f \pm g\|_{1} = \int_{-\pi}^{\pi} |1 \pm \operatorname{Re}\left\{e^{i\alpha}q(e^{i\theta})\right\}| |f(e^{i\theta})|\dd \theta = 1.
\end{equation*}
这与引理~\ref{lem:2.15.1} 矛盾,故 f 是外函数.
\end{proof}

\begin{theorem}\label{thm:2.15.4}
设 $f \in H^{\infty}$.f 是 $S^{\infty}$ 的端点的充要条件是 $\|f\|_{\infty} = 1$ 且
\begin{equation*}
\int_{-\pi}^{\pi} \log(1 - |f(e^{i\theta})|)\dd \theta = -\infty.
\end{equation*}
\end{theorem}

\begin{proof}
设 $f \in S^{\infty}$ 满足
\begin{equation*}
\int_{-\pi}^{\pi} \log(1 - |f(e^{i\theta})|)\dd \theta = -\infty.
\end{equation*}
设 $g \in S^{\infty}$ 满足 $\|f \pm g\|_{\infty} \leq 1$,则
\begin{equation*}
|g(e^{i\theta})|^{2} \leq 1 - |f(e^{i\theta})|^{2} \leq 2(1 - |f(e^{i\theta})|),
\end{equation*}
故 $\int_{-\pi}^{\pi} \log|g(e^{i\theta})|\dd \theta = -\infty$.由定理~\ref{thm:2.5.3},$g \equiv 0$.于是由引理~\ref{lem:2.15.1},f 是 $S^{\infty}$ 的端点.

反之,若 $\|f\|_{\infty} = 1$ 且
\begin{equation*}
\int_{-\pi}^{\pi} \log(1 - |f(e^{i\theta})|)\dd \theta > -\infty,
\end{equation*}
令
\begin{equation*}
g(z) = \exp\left\{\frac{1}{2\pi}\int_{-\pi}^{\pi} \frac{e^{it} + z}{e^{it} - z}\log(1 - |f(e^{it})|)\dd t\right\},
\end{equation*}
则因 $\log(1 - |f(e^{it})|) \leq 0$,$|g(z)| \leq 1$,且由定理~\ref{thm:2.3.2},a.e. t 有 $|g(e^{it})| = 1 - |f(e^{it})|$.于是 $\|f \pm g\|_{\infty} \leq 1$ 且 $g \in S^{\infty}$,由引理~\ref{lem:2.15.1},f 不是 $S^{\infty}$ 的端点.
\end{proof}

\begin{proposition}\label{prop:2.15.5}
设 $f \in S^{1}$ 且 $\|f\|_{1} = 1$,则 f 是 $S^{1}$ 的露点的充要条件是:满足 $g \in H^{1}$ 且 a.e. $\theta$ 有 $g(e^{i\theta})/f(e^{i\theta}) \geq 0$ 的 g 仅为 f 的标量倍.
\end{proposition}

\begin{proof}
设 f 是 $S^{1}$ 的露点.可设存在 $\phi \in (H^{1})^{*}$ 且 $\phi(f) > 0$.若非零 $g \in H^{1}$ 满足 a.e. $\theta$ 有 $g(e^{i\theta})/f(e^{i\theta}) \geq 0$,则由引理~\ref{lem:2.14.3},$g/\|g\|_{1}$ 也是 $\phi$ 的极值函数.因 f 是 $S^{1}$ 的露点,g 是 f 的标量倍.设 f 不是 $S^{1}$ 的露点.令 K = $|f|/zf$,设 $\phi$ 为由 K 导出的 $(H^{1})^{*}$ 的元素,则 f 是 $\phi$ 的极值函数.因 f 不是露点,存在 $g \in S^{1}$,不是 f 的标量倍,且 $\phi(g) = \phi(f)$.由引理~\ref{lem:2.14.3},a.e. $\theta$ 有 $g(e^{i\theta})/f(e^{i\theta}) \geq 0$.
\end{proof}

\begin{proposition}\label{prop:2.15.6}
设 $f \in S^{\infty}$,若 f 在 $\partial D$ 上正测度的集合上绝对值等于 1,则 f 是 $S^{\infty}$ 的露点.
\end{proposition}

\begin{proof}
令 $E = \{\theta \in [-\pi, \pi]; |f(e^{i\theta})| = 1\}$,则 $\dd \theta(E) > 0$.
\begin{equation*}
K(e^{i\theta}) = \left\{ \begin{array}{c l}
e^{-i\theta}|f(e^{i\theta})|/f(e^{i\theta}) & (\theta \in E), \\
0 & (\theta \notin E)
\end{array} \right.
\end{equation*}
则 $K \in L^{1}$,且 a.e. $\theta$ 有 $e^{i\theta}f(e^{i\theta})K(e^{i\theta}) \geq 0$.由引理~\ref{lem:2.14.3},f 是由极值核 K 定义的 $\phi \in (H^{\infty})^{*}$ 的极值函数,且 $\phi(f) > 0$.设 g 为 $\phi$ 的另一极值函数,则存在实常数 $\alpha$,a.e. $\theta$ 有 $e^{i\alpha}e^{i\theta}g(e^{i\theta})K(e^{i\theta}) \geq 0$,故 $g = e^{-i\alpha}f$.
\end{proof}

\begin{example}\label{exa:2.15.7}
设 $f \in S^{1}$ 且 $\|f\|_{1} = 1$.若存在 $k \in H^{\infty}$($k^{-1} \in H^{\infty}$)使得 a.e. $\theta$ 有 $\operatorname{Re} f(e^{i\theta})k(e^{i\theta}) \geq 0$,则 f 是 $S^{1}$ 的端点.
\end{example}

\begin{proof}
首先证明当 $g = fk$ 时 g 是外函数.由定理~\ref{thm:2.8.3},$g = B_{1}S_{1}h_{1}$ 可作因式分解.因 a.e. $\theta$ 有 $\operatorname{Re} g(e^{i\theta}) \geq 0$,由定理~\ref{thm:2.9.1},$\operatorname{Re} g(z) > 0$($z \in D$),故 $g = S_{1}h_{1}$.因 $\operatorname{Re} g^{-1}(z) > 0$($z \in D$),由定理~\ref{thm:2.11.2},$g^{-1} \in \bigcap_{p<1} H^{p}$.由定理~\ref{thm:2.8.3},$g^{-1} = S_{2}h_{2}$ 可作因式分解.于是 $1 = S_{1}S_{2}h_{1}h_{2}$.由定义,对 j = 1,2,
\begin{equation*}
S_{j}(z) = e^{i\gamma_{j}}\exp\left\{-\int_{-\pi}^{\pi} \frac{e^{it} + z}{e^{it} - z}\dd \mu_{j}(t)\right\},
\end{equation*}
其中 $\mu_{j} \in M[-\pi, \pi]$ 是奇异正测度.又由定义,对 j = 1,2,
\begin{equation*}
h_{j}(z) = \exp\left\{\int_{-\pi}^{\pi} \frac{e^{it} + z}{e^{it} - z}\log|h_{j}(e^{it})|\dd t\right\}.
\end{equation*}
于是由 $1 = S_{1}S_{2}h_{1}h_{2}$,a.e. $\theta$ 有 $|(h_{1}h_{2})(e^{i\theta})| = 1$,故 $(h_{1}h_{2})(z) = 1$($z \in D$).于是 $(S_{1}S_{2})(z) = 1$($z \in D$),故 $\mu_{1} + \mu_{2} = 0$,从而 $\mu_{1} = \mu_{2} = 0$.于是 $S_{1} = e^{i\gamma_{1}}$,$g = e^{i\gamma_{1}}h_{1}$ 是外函数.k 是外函数可同样证明,故 $f = k^{-1}g$ 也是外函数.
\end{proof}

\begin{example}\label{exa:2.15.8}
设 $f \in S^{1}$ 且 $\|f\|_{1} = 1$.若存在 $k \in H^{\infty}$($k^{-1} \in H^{\infty}$)使得 a.e. $\theta$ 有 $\operatorname{Re} f(e^{i\theta})k(e^{i\theta}) \geq 0$,则 f 是 $S^{1}$ 的露点.
\end{example}

\begin{proof}
令 $g = fk/\|fk\|_{1}$.因 $k, k^{-1} \in H^{\infty}$,由命题~\ref{prop:2.15.5},只需证明:若 $h \in S^{1}$ 满足 a.e. $\theta$ 有 $h(e^{i\theta})/g(e^{i\theta}) \geq 0$,则 h = g.由 a.e. $\theta$ 有 $\mathrm{Re}(h(e^{i\theta}) + g(e^{i\theta}))/2 \geq 0$,由例~\ref{exa:2.15.7},$(h + g)/2$ 是外函数.另一方面,若 h 不是 g 的标量倍,则与定理~\ref{thm:2.15.3} 矛盾.
\end{proof}

定理~\ref{thm:2.15.3} 是 de Leeuw--Rudin 的结果 \cite{ch2:5}.已知命题~\ref{prop:2.15.6} 的逆也成立 \cite{ch2:6}.例~\ref{exa:2.15.8} 出自 \cite{ch2:7},其证明依据 \cite{ch2:8}.$S^{1}$ 的露点与端点相比研究得还不充分.关于露点的研究见 \cite{ch2:9} 与 \cite{ch2:10}.

% ============ 第2章 §2.16 极值问题的解(原书 pp.165-173) ============
\section{极值问题的解}\label{sec:2.16}

当 $1 \leq p \leq \infty$ 时,对 $\phi \in (H^{p})^{*}$,其极值问题的解中满足 $\phi(f) > 0$ 者的全体记为 $S_{\phi}^{p}$.即
\begin{equation*}
S_{\phi}^{p} = \{f \in S^{p}; \phi(f) = \|\phi\|\},
\end{equation*}
其中 $S^{p} = \{f \in H^{p}; \|f\|_{p} = 1\}$.$S_{\phi}^{p}$ 是凸集合.设 $\langle S_{\phi}^{p} \rangle$ 为包含 $S_{\phi}^{p}$ 的 $H^{p}$ 的最小线性子空间.由定理~\ref{thm:2.14.4},当 $1 < p < \infty$ 时 $\langle S_{\phi}^{p} \rangle$ 是 1 维的.又当 $p = \infty$ 且 $\phi \in L^{1}/zH^{1}$ 时 $\langle S_{\phi}^{p} \rangle$ 是 1 维的.但当 p = 1 时,即使 $\phi \in C/zA$,$\langle S_{\phi}^{p} \rangle$ 也未必是 1 维的.当 $1/p + 1/q = 1$ 时,对 $k \in L^{q}$,
\begin{equation*}
\phi(f) = \frac{1}{2\pi i}\int_{|z|=1} f(z) k(z)\dd z \quad (f \in H^{p})
\end{equation*}
时记 $S_{\phi}^{p} = S_{k}^{p}$.当 $p \neq \infty$ 时,由定理~\ref{thm:2.13.1},对 $\phi \in (H^{p})^{*}$ 总存在上述那样的 k.

引理~\ref{lem:2.14.3} 给出了 $\phi \in (H^{p})^{*}$ 的极值核与极值解的关系,由此立刻导出下面的命题.

\begin{proposition}\label{prop:2.16.1}
设 $1 \leq p < \infty$,$\phi \in (H^{p})^{*}$ 且 $\|\phi\| = 1$.设 K 为 $\phi$ 的极值核,则
\begin{equation*}
S_{\phi}^{p} = \left\{f \in S^{p}; K(e^{i\theta}) = e^{-i\theta}\frac{|f(e^{i\theta})|^{p}}{f(e^{i\theta})} \text{ a.e. } \theta \right\}.
\end{equation*}
这里,p = 1 时假定 $S_{\phi}^{1} \neq \emptyset$.
\end{proposition}

\begin{proof}
由引理~\ref{lem:2.14.3},若 $f \in S_{\phi}^{p}$,则 $e^{i\theta}K(e^{i\theta})f(e^{i\theta}) = |K(e^{i\theta})||f(e^{i\theta})|$ 且 $|K(e^{i\theta})| = |f(e^{i\theta})|^{p/q}$,由 $1/p + 1/q = 1$ 即明.
\end{proof}

k 的具体例子如下.

\begin{example}\label{exa:2.16.2}
设 k 为 $\phi$ 的积分表示的核.设 $f(z) = \sum_{j=0}^{\infty} a_{j}z^{j} \in H(D)$.

(1) 若 $k(z) = \sum_{j=1}^{n}\frac{c_{j}}{z - \beta_{j}}$($|\beta_{j}| < 1$),则 $\phi(f) = \sum_{j=1}^{n} c_{j}f(\beta_{j})$.

(2) 若 $k(z) = n!(z-\beta)^{-n-1}$($|\beta| < 1$),则 $\phi(f) = f^{(n)}(\beta)$.

(3) 若 $k(z) = \sum_{j=0}^{n} c_{j}z^{-j-1}$,则 $\phi(f) = \sum_{j=0}^{n} c_{j}a_{j}$.
\end{example}

下面的例子描绘了 (3) 中 $c_{0} = \cdots = c_{n-1} = 0$ 且 $c_{n} = 1$ 时的 $S_{\phi}^{p}$.它在 $\phi$ 为连续函数时描绘 $S_{\phi}^{p}$ 是重要的.

\begin{example}\label{exa:2.16.3}
当 $k(z) = z^{-(n+1)}$ 时,若 $p \neq 1$ 则 $S_{k}^{p} = \{z^{n}\}$,p = 1 时
\begin{equation*}
S_{k}^{1} = \left\{\gamma \prod_{j=1}^{n} (z - \alpha_{j})(1 - \bar{\alpha}_{j}z) \in S^{1}; \gamma > 0 \ \text{且} \ \alpha_{j} \in \mathbb{C}, |\alpha_{j}| \leq 1 \right\}.
\end{equation*}
\end{example}

\begin{proof}
p $\neq$ 1 时由定理~\ref{thm:2.14.4} 显然.p = 1 时,$S_{k}^{1}$ 包含右边由命题~\ref{prop:2.16.1} 显然.证明 $S_{k}^{1}$ 包含于右边.对 $f \in S_{k}^{1}$,令 $F(z) = f(z)/z^{n}$,由命题~\ref{prop:2.16.1},$F(z) \geq 0$($z \in \partial D$),故 $F(z) = \sum_{j=-n}^{n} c_{j}z^{j}$,$\bar{c}_{j} = c_{-j}$ 且 $c_{n} \neq 0$.$F(z)$ 是有理函数,且 $F(z) \geq 0$($z \in \partial D$),由镜像原理,若 $F(\alpha) = 0$ 则 $F(1/\bar{\alpha}) = 0$.$f(z) = z^{n}F(z)$ 是 2n 次多项式,有 2n 个零点.$\partial D$ 上的零点为偶重.于是
\begin{equation*}
f(z) = \gamma \prod_{j=1}^{n} (z - \alpha_{j})(1 - \bar{\alpha}_{j}z),
\end{equation*}
其中 $|\alpha_{j}| \leq 1$ 且 $\gamma > 0$.
\end{proof}

\begin{lemma}\label{lem:2.16.4}
设 $q_{j}$($1 \leq j \leq \ell$)为非常值的内函数,$q = \prod_{j=1}^{\ell} q_{j}$ 且 h 为外函数.设 $K = \overline{zq}|h|/h$,则对任意 $\{b_{j}\}_{j=1}^{\ell} \subset D$,存在 $f \in S_{K}^{1}$,使得 $f(b_{1}) = \cdots = f(b_{\ell}) = 0$.
\end{lemma}

\begin{proof}
对 a.e. $\theta$,
\begin{equation*}
\overline{q(e^{i\theta})}\prod_{j=1}^{\ell} (q_{j}(e^{i\theta}) - q_{j}(0))(1 - \overline{q_{j}(0)}q_{j}(e^{i\theta})) \geq 0,
\end{equation*}
故选取常数 $\gamma > 0$,由命题~\ref{prop:2.16.1},
\begin{equation*}
g = \gamma h\prod_{j=1}^{\ell} (q_{j} - q_{j}(0))(1 - \overline{q_{j}(0)}q_{j})
\end{equation*}
属于 $S_{K}^{1}$.因 $g/z^{\ell} \in H^{1}$,对任意 $\{b_{j}\}_{j=1}^{\ell} \subset D$,存在常数 $\delta > 0$,使得
\begin{equation*}
f = \delta \frac{g}{z^{\ell}}\prod_{j=1}^{\ell} (z - b_{j})(1 - \bar{b}_{j}z)
\end{equation*}
属于 $S_{K}^{1}$,故 f 为所求函数.
\end{proof}

\begin{lemma}\label{lem:2.16.5}
设 $K = \bar{z}^{m+1}|g|/g$,m 为非负整数且 g 为 $S^{1}$ 的露点.任意 $f \in S_{K}^{1}$ 的因式分解中的内函数,是在 D 中零点个数不超过 $(m+1)$ 的有限 Blaschke 积.
\end{lemma}

\begin{proof}
若对 $\{a_{j}\}_{j=1}^{m+1} \subset D$ 存在 $f \in S_{K}^{1}$,使得 $f(a_{1}) = \cdots = f(a_{m+1}) = 0$,则存在 $\gamma > 0$ 与 $h \in H^{1}$,使得
\begin{equation*}
f(z) = \gamma \prod_{j=1}^{m+1} (z - a_{j})(1 - \bar{a}_{j}z) h.
\end{equation*}
于是由例~\ref{exa:2.16.3},存在 $\gamma_{0} > 0$,可令 $\gamma_{0}z^{m+1}h \in S_{K}^{1}$.此时由引理~\ref{lem:2.14.3},a.e. $\theta$ 有 $e^{i\theta}h(e^{i\theta})/g(e^{i\theta}) \geq 0$,由露点的定义,g 是 zh 的常数倍.于是 g 不是外函数,与 g 是 $S^{1}$ 的露点矛盾.若 f = qh 可分解且 q 为奇异内函数,则 $q = \prod_{j=1}^{m+1} q_{j}$ 可分解,其中 $q_{j}$($1 \leq j \leq m+1$)不是常数.于是取 $\delta > 0$,$|\alpha_{j}| \leq 1$($1 \leq j \leq m+1$)且
\begin{equation*}
F(z) = \delta \prod_{j=1}^{m+1} (q_{j} - \alpha_{j})(1 - \bar{\alpha}_{j}q_{j}) h,
\end{equation*}
则 F 可属于 $S_{K}$.选取 $q_{j} - \alpha_{j}$ 使在 D 中有零点,则 F 有 $(m+1)$ 个零点.这与上面所证矛盾,故 f 的内函数是所求的 Blaschke 积.
\end{proof}

\begin{lemma}\label{lem:2.16.6}
设 $K = \overline{z}^{m+1}|g|/g$,m 为非负整数且 g 为 $S^{1}$ 的露点.此时
\begin{equation*}
S_{K}^{1} = (\{\gamma\} \times S_{\bar{z}^{m+1}}^{1} \times g) \cap S^{1}.
\end{equation*}
\end{lemma}

\begin{proof}
对 m 用归纳法证明.令 $K_{m} = \overline{z}^{m+1}|g|/g$.若 m = 0,因 g 是 $S^{1}$ 的露点,$S_{K_{0}}^{1} = \{g\}$.设对 $1 \leq m \leq n$ 本引理成立,即 $S_{K_{m}}^{1} = (\{\gamma\} \times S_{\bar{z}^{m+1}}^{1} \times g) \cap S^{1}$ 成立.$S_{K_{n+1}}^{1} \supseteq (\{\gamma\} \times S_{\bar{z}^{n+2}}^{1} \times g) \cap S^{1}$ 由例~\ref{exa:2.16.3} 显然,故证明 $S_{K_{n+1}}^{1} \subseteq (\{\gamma\} \times S_{\bar{z}^{n+2}}^{1} \times g) \cap S^{1}$.若 $f \in S_{K_{n+1}}^{1}$,则
\begin{equation*}
F = (f + z^{n+1}g)/2 \in S_{K_{n+1}}^{1}.
\end{equation*}
因 F 不是 $S^{1}$ 的端点,由定理~\ref{thm:2.15.3},F 不是外函数.由引理~\ref{lem:2.16.5},存在某个 $1 \leq \ell \leq n+1$,$\{a_{j}\}_{j=1}^{\ell} \subset D$ 与外函数 $h \in H^{1}$,使得
\begin{equation*}
F = \prod_{j=1}^{\ell} \frac{z - a_{j}}{1 - \bar{a}_{j}z} h.
\end{equation*}
对每个 j,因 $\mathrm{Re}(1/(1-\bar{a}_{j}z)) \geq 0$,由定理~\ref{thm:2.15.3} 与例~\ref{exa:2.15.7},$(1-\bar{a}_{j}z)^{-1}$ 是外函数.有限个外函数的乘积仍为外函数,故令
\begin{equation*}
k = \frac{1}{\alpha}\prod_{j=1}^{\ell} \frac{h}{(1 - \bar{a}_{j}z)^{2}},
\end{equation*}
并选取 $\alpha > 0$ 使 $\|k\|_{1} = 1$,则 k 是 $H^{1}$ 的外函数,且
\begin{equation*}
F = \alpha \prod_{j=1}^{\ell} (z - a_{j})(1 - \bar{a}_{j}z) k.
\end{equation*}
\begin{sloppypar}
于是 a.e. $\theta$ 有 $e^{i\ell\theta}k(e^{i\theta})/F(e^{i\theta}) \geq 0$.又因 a.e. $\theta$ 有 $F(e^{i\theta})e^{-i(n+1)\theta}/g(e^{i\theta}) \geq 0$,故 a.e. $\theta$ 有 $e^{i\ell\theta}k(e^{i\theta})e^{-i(n+1)\theta}/g(e^{i\theta}) \geq 0$,即 $k \in S_{K_{n-\ell+1}}^{1}$.
\end{sloppypar}由归纳假设,
\begin{equation*}
k = \beta \prod_{j=1}^{n-\ell+1} (z - b_{j})(1 - \bar{b}_{j}z) g,
\end{equation*}
其中 $\beta > 0$ 且 $\{b_{j}\}_{j=1}^{n-\ell+1} \subset \partial D$.令 $a_{\ell+j} = b_{j}$($1 \leq j \leq n - \ell + 1$),则
\begin{equation*}
F = \alpha\beta \prod_{j=1}^{n+1} (z - a_{j})(1 - \bar{a}_{j}z) g.
\end{equation*}
令
\begin{equation*}
p = 2\alpha\beta \prod_{j=1}^{n+1} (z - a_{j})(1 - \bar{a}_{j}z) - z^{n+1},
\end{equation*}
则 $f = 2F - z^{n+1}g = pg$.因 $f \in S_{K_{n+1}}^{1}$,a.e. $\theta$ 有 $e^{-i(n+1)\theta}g(e^{i\theta})^{-1}f(e^{i\theta}) \geq 0$,故 a.e. $\theta$ 有 $e^{-i(n+1)\theta}p(e^{i\theta}) \geq 0$.于是 $p \in \{\gamma\} \times S_{\bar{z}^{n+2}}^{1}$,故 f 属于 $(\{\gamma\} \times S_{\bar{z}^{n+2}}^{1} \times g) \cap S^{1}$.
\end{proof}

\begin{lemma}\label{lem:2.16.7}
设 $k \notin H^{\infty}$ 且 $S_{k}^{1}$ 是非空 $S^{1}$ 的 *弱紧集.则存在非负整数 m,使得任意 $f \in S_{k}^{1}$ 的因式分解中的内函数,是在 D 中零点个数不超过 m 的有限 Blaschke 积.
\end{lemma}

\begin{proof}
若对 D 中任意有限个点,以之为零点的函数存在于 $S_{k}^{1}$,则由引理~\ref{lem:2.16.4},对每个自然数 n,
\begin{equation*}
f_{n}\left(\frac{1}{2}\right) = f_{n}\left(\frac{1}{3}\right) = \dots = f_{n}\left(\frac{1}{n}\right) = 0
\end{equation*}
的 $f_{n}$ 存在于 $S_{k}^{1}$.因 $S_{k}^{1}$ 是 *弱紧的,$\{f_{n}\}$ 的子列在 $H^{1} = (C/zA)^{*}$(参见定理~\ref{thm:2.13.3})的 *弱拓扑下收敛到 $f \in S_{k}^{1}$.对每个 n,点 $z = 1/n$ 处的点赋值由命题 2.9.4 关于 *弱拓扑连续,故 $f(1/n) = 0$($n = 2, 3, \ldots$).由唯一性定理,$f \equiv 0$.这与 $f \in S_{k}^{1}$ 矛盾.因此,由于引理~\ref{lem:2.16.4} 表明任意 $f \in S_{k}^{1}$ 的因式分解中的内函数是有限 Blaschke 积,本引理证明完毕.
\end{proof}

\begin{theorem}\label{thm:2.16.8}
设 $\phi \in C/zA$,即 $k \in C$.此时存在非负整数 m 与 $S^{1}$ 的露点 g,使得
\begin{equation*}
S_{\phi}^{1} = S_{k}^{1} = \left\{\gamma \prod_{j=1}^{m} (z - a_{j})(1 - \bar{a}_{j}z) g \in S^{1}; \gamma > 0, \{a_{j}\}_{j=1}^{m} \subset D \right\}.
\end{equation*}
\end{theorem}

\begin{proof}
由假设与定理~\ref{thm:2.13.3},$S_{\phi}^{1}$ 是 *弱紧的.由引理~\ref{lem:2.16.7},任意 $f \in S_{\phi}^{1}$ 的因式分解中的内函数是有限 Blaschke 积,其零点个数不超过某个常数.设 m 为其零点最大个数,则存在 $f \in S_{\phi}^{1}$,$f = \prod_{j=1}^{m} \frac{z - a_{j}}{1 - \bar{a}_{j}z} h$,其中 $h \in H^{1}$ 是外函数.
\begin{equation*}
g = \gamma h / \prod_{j=1}^{m} (1 - \bar{a}_{j}z)^{2},
\end{equation*}
选取 $\gamma > 0$ 可使 $g \in S^{1}$.
\begin{equation*}
\frac{f(e^{i\theta})}{e^{im\theta}g(e^{i\theta})} = \frac{1}{\gamma}\frac{\prod_{j=1}^{m} (e^{i\theta} - a_{j})(1 - \bar{a}_{j}e^{i\theta})}{e^{im\theta}} \geq 0 \ \text{a.e.} \ \theta,
\end{equation*}
故由命题~\ref{prop:2.16.1},$z^{m}g \in S_{\phi}^{1}$.一旦知道 g 是 $S^{1}$ 的露点,令 $K = \bar{z}^{m+1}|g|/g$ 并用引理~\ref{lem:2.16.6},定理证明完毕.若 g 不是 $S^{1}$ 的露点,由命题~\ref{prop:2.15.5},存在 $G \in S^{1}$($g \neq G$)使得 $z^{m}G \in S_{\phi}^{1}$.此时 $z^{m}(g + G)/2 \in S_{\phi}^{1}$,但由定理~\ref{thm:2.15.3},$(g + G)/2$ 不是外函数,这与 $S_{\phi}^{1}$ 中函数的内函数部分是 Blaschke 积且零点个数不超过 m 相矛盾.
\end{proof}

\begin{theorem}\label{thm:2.16.9}
设 $1 < p < \infty$,$1/p + 1/q = 1$ 且 $\phi \in (H^{p})^{*}$,$\|\phi\| = 1$.设 K 为 $\phi$ 的极值核,G 为满足 a.e. $\theta$ 有 $|K(e^{i\theta})| = |G(e^{i\theta})|$ 的外函数.设 $\phi_{0} \in (H^{1})^{*}$ 由 $K_{0} = K/G$ 定义.若 $\phi_{0} \in C/zA$,则存在 $\{a_{j}\}_{j=1}^{n} \subset \overline{D}$ 与 $\gamma > 0$,使得
$K_{0} = \bar{z}^{n+1}\frac{|g|}{g}$ 且 g 是 $S^{1}$ 的露点,
\begin{equation*}
G = \gamma^{1/q}\prod_{j=1}^{s} (1 - \bar{a}_{j}z)^{2/q}\left[ \prod_{j=s+1}^{n} (z - a_{j})(1 - \bar{a}_{j}z) \right]^{1/q} g^{1/q},
\end{equation*}
且 $S_{\phi}^{p}$ 由形如
\begin{equation*}
\gamma^{1/p}\prod_{j=1}^{s} (z - a_{j})(1 - \bar{a}_{j}z)^{2/p - 1}\left[ \prod_{j=s+1}^{n} (z - a_{j})(1 - \bar{a}_{j}z) \right]^{1/p} g^{1/p}
\end{equation*}
的 $S^{p}$ 的元素组成.其中 $|a_{j}| < 1$($1 \leq j \leq s$),$|a_{j}| = 1$($s + 1 \leq j \leq n$)且 $\gamma > 0$.
\end{theorem}

\begin{proof}
因 $1 < p < \infty$,由定理~\ref{thm:2.14.4},存在 $F \in H^{p}$ 使得 $S_{\phi}^{p} = \{F\}$.此时因 K 是 $\phi$ 的极值核,由引理~\ref{lem:2.14.3},a.e. $\theta$ 有 $|K(e^{i\theta})| = |F(e^{i\theta})|^{p/q}$.由定理~\ref{thm:2.8.3},F = Qh 可作因式分解.令 G = $h^{q/p}$,则 G 是外函数,且 a.e. $\theta$ 有 $|K(e^{i\theta})| = |G(e^{i\theta})|$,且 F = $QG^{p/q}$.因 $K_{0} = K/G$,利用 q/p + 1 = q,由引理~\ref{lem:2.14.3},a.e. $\theta$ 有 $e^{i\theta}Q(e^{i\theta})G^{q}(e^{i\theta})K_{0}(e^{i\theta}) \geq 0$.因 $|QG^{q}| = |F|^{p}$,$QG^{q} \in S^{1}$,故 $QG^{q} \in S_{\phi_{0}}^{1}$.因 $\phi_{0} \in C/zA$,由定理~\ref{thm:2.16.8},$K_{0} = \overline{z}^{n+1}|g|/g$,g 是 $S^{1}$ 的露点,且
\begin{equation*}
QG^{q} = \gamma \prod_{j=1}^{n} (z - a_{j})(1 - \bar{a}_{j}z) g
\end{equation*}
存在 $\gamma > 0$ 与 $\{a_{j}\}_{j=1}^{n} \subset \overline{D}$ 和 g 满足上式.这里 $\{a_{j}\}_{j=1}^{n}$ 可看作满足 $|a_{j}| < 1$($1 \leq j \leq s$)且 $|a_{j}| = 1$($s + 1 \leq j \leq n$).因此
\begin{equation*}
Q = \prod_{j=1}^{s} \frac{z - a_{j}}{1 - \bar{a}_{j}z},
\end{equation*}
且
\begin{equation*}
G^{q} = \gamma \prod_{j=1}^{s} (1 - \bar{a}_{j}z)^{2}\prod_{j=s+1}^{n} (z - a_{j})(1 - \bar{a}_{j}z) g.
\end{equation*}
因 $F = QG^{q/p}$,
\begin{equation*}
F = \gamma^{1/p}\prod_{j=1}^{s} (z - a_{j})(1 - \bar{a}_{j}z)^{2/p - 1}\left[ \prod_{j=s+1}^{n} (z - a_{j})(1 - \bar{a}_{j}z) \right]^{1/p} g^{1/p}.
\end{equation*}
\end{proof}

定义 $\phi$ 的核 k 为有理函数的情形在应用上很重要,被许多人所研究.例如,参见 Duren 的专著 \cite{ch2:11} 的第 8 章.若 k 为有理函数,则 $k \in C$,故可应用定理~\ref{thm:2.16.8}.另外,当 $1 < p < \infty$ 时,已知 k 的极值核 K 又是有理函数,故可应用定理~\ref{thm:2.16.9}.定理~\ref{thm:2.16.8} 在 \cite{ch2:8} 中,定理~\ref{thm:2.16.9} 在 \cite{ch2:12} 中被证明.定理~\ref{thm:2.16.8} 中 $\phi \notin C/zA$ 时 $S_{\phi}^{1}$ 的研究见 \cite{ch2:13}.

% ============ 第2章 §2.17 Pick插值问题(原书 pp.173-179) ============
\section{Pick 插值问题}\label{sec:2.17}
\index{插值问题!最小范数插值问题}


对 D 中相异的有限个点 $\{z_{j}\}_{j=1}^{n}$ 与 $\mathbb{C}$ 中有限个数据 $\{w_{j}\}_{j=1}^{n}$,求属于 $H^{\infty}$ 且满足
\begin{equation*}
f(z_{j}) = w_{j} \quad (1 \leq j \leq n)
\end{equation*}
的函数 f.事实上,对 $1 \leq j \leq n$,令
\begin{equation*}
B_{n}(z) = \prod_{j=1}^{n} \frac{z - z_{j}}{1 - \bar{z}_{j}z} \quad \text{且} \quad B_{nj}(z) = B_{n}(z)\frac{1 - \bar{z}_{j}z}{z - z_{j}}.
\end{equation*}
令
\begin{equation*}
f(z) = \sum_{j=1}^{n} w_{j}\frac{B_{nj}(z)}{B_{nj}(z_{j})},
\end{equation*}
则 $f(z_{j}) = w_{j}$($1 \leq j \leq n$).于是当对 f 没有进一步条件时,有限个点的情形插值问题很容易解决.应用上重要的(有限)插值问题是:何时能选取 $\|f\|_{\infty} \leq 1$.

定理~\ref{thm:2.17.1} 称为最小范数插值问题,将在 2.18 节中使用.定理~\ref{thm:2.17.7} 称为 Pick 定理 \cite{ch2:14}.我们用 Sarason 的新方法,即利用定理~\ref{thm:2.17.5} 证明 Pick 定理.该方法不用复函数论,而用算子论.

\begin{theorem}\label{thm:2.17.1}
设 $1 \leq p \leq \infty$,$\{z_{j}\}_{j=1}^{n} \subset D$ 且 $\{w_{j}\}_{j=1}^{n} \subset \mathbb{C}$.此时存在满足 $f(z_{j}) = w_{j}$($1 \leq j \leq n$)的最小范数的 $f \in H^{p}$.且
\begin{align*}
\min \{\|f\|_{p}; f \in H^{p}, f(z_{j}) = w_{j} (1 \leq j \leq n)\}&= \max \left\{\left| \sum_{j=1}^{n} \frac{w_{j}(1 - |z_{j}|^{2})}{B_{nj}(z_{j})} g(z_{j}) \right|; g \in H^{q}, \|g\|_{q} \leq 1 \right\}.
 \end{align*}
\end{theorem}

\begin{proof}
首先,
\begin{equation*}
F(z) = \sum_{j=1}^{n} w_{j}\frac{B_{nj}(z)}{B_{nj}(z_{j})},
\end{equation*}
则 $F \in H^{p}$.对 $g \in H^{p}$,令
\begin{equation*}
f(z) = F(z) - B_{n}(z)g(z),
\end{equation*}
则 $f(z_{j}) = w_{j}$($1 \leq j \leq n$).对 $\{w_{j}\}_{j=1}^{n}$ 的一般插值函数具有这种形式.令 k = $F/B_{n}$,则 $k \in L^{p}$,且
\begin{equation*}
\inf_{g \in H^{p}} \|F - B_{n}g\|_{p} = \inf_{g \in H^{p}} \|k - g\|_{p}.
\end{equation*}
由定理~\ref{thm:2.14.2},inf 可改为 min,故存在满足 $f(z_{j}) = w_{j}$($1 \leq j \leq n$)的最小范数的 $f \in H^{p}$.

因
\begin{equation*}
k(z) = \sum_{j=1}^{n} \frac{w_{j}}{B_{nj}(z_{j})}\frac{1 - \bar{z}_{j}z}{z - z_{j}},
\end{equation*}
对 $h \in H^{q}$($1/p + 1/q = 1$),令
\begin{equation*}
\phi(f) = \frac{1}{2\pi i}\int_{-\pi}^{\pi} k(z) h(z)\dd z = \sum_{j=1}^{n} \frac{w_{j}(1 - |z_{j}|^{2})}{B_{nj}(z_{j})} h(z_{j}),
\end{equation*}
则 $\phi \in (H^{q})^{*}$.由定理~\ref{thm:2.14.2},当 $p \neq 1$ 时
\begin{align*}
\min \{\|f\|_{p}; f \in H^{p}, f(z_{j}) = w_{j} (1 \leq j \leq n)\}&= \min \{\|k - g\|_{p}; g \in H^{p}\}  \\
\quad &= \max \left\{\left| \sum_{j=1}^{n} \frac{w_{j}(1 - |z_{j}|^{2})}{B_{nj}(z_{j})} h(z_{j}) \right|; h \in H^{q}, \|h\|_{q} \leq 1 \right\}.
\end{align*}
p = 1 时,因 $k \in C$,由定理~\ref{thm:2.14.2} 与定理~\ref{thm:2.14.4} 的 (3),上述等式成立.
\end{proof}

对 $f \in H^{\infty}$,设 $M_{f}$ 为 $H^{2}$ 上的乘法算子,则它是 $H^{2}$ 上的有界线性算子.若在 $H^{2}$ 上定义内积 $\langle\ ,\ \rangle$ 为
\begin{equation*}
\langle f, g \rangle = \frac{1}{2\pi}\int_{-\pi}^{\pi} f(e^{i\theta})\overline{g(e^{i\theta})}\dd \theta \quad (f, g \in H^{2}),
\end{equation*}
则 $H^{2}$ 成为 Hilbert 空间.对内函数 q,设 $N = H^{2} \ominus qH^{2}$,即 $qH^{2}$ 在 $H^{2}$ 中的正交补空间.设 P 为从 $H^{2}$ 到 N 的正交射影.定义 N 上的有界线性算子 $T_{f}$ 为 $T_{f} = PM_{f}|N$.$T_{f}^{*}$ 表示 $T_{f}$ 的共轭算子.定理~\ref{thm:2.17.5} 在 q 为一般内函数时由 Sarason \cite{ch2:15} 证明.

\begin{lemma}\label{lem:2.17.2}
当 $q = B_{n}$ 时,N 由 $\{(1 - \bar{z}_{j}z)^{-1}\}_{j=1}^{n}$ 生成.
\end{lemma}

\begin{proof}
当 $|a| < 1$ 时,由命题 2.9.4,若 $g \in H^{2}$ 则
\begin{equation*}
g(a) = \frac{1}{2\pi}\int_{-\pi}^{\pi} \frac{e^{it}}{e^{it} - a} g(e^{it})\dd t.
\end{equation*}
因
\begin{equation*}
\frac{e^{it}}{e^{it} - a} = \frac{1}{1 - ae^{-it}},
\end{equation*}
$\{(1 - \bar{z}_{j}z)^{-1}\}_{j=1}^{n}$ 与 $B_{n}H^{2}$ 正交.因 $B_{n}H^{2} = \{g \in H^{2}; g(z_{j}) = 0 \ (1 \leq j \leq n)\}$,N 是 n 维空间,故 N 由 $\{(1 - \bar{z}_{j}z)^{-1}\}_{j=1}^{n}$ 生成.
\end{proof}

\begin{lemma}\label{lem:2.17.3}
当 $q = B_{n}$ 时,若 $f \in H^{\infty}$,则
\begin{equation*}
T_{f}^{*}\left(\frac{1}{1 - \bar{z}_{j}z}\right) = \frac{\overline{f(z_{j})}}{1 - \bar{z}_{j}z} \quad (1 \leq j \leq n).
\end{equation*}
\end{lemma}

\begin{proof}
令 $k_{j} = (1 - \bar{z}_{j}z)^{-1}$($1 \leq j \leq n$),则对任意 $g \in N$,
\begin{align*}
\langle T_{f}^{*}k_{j}, g \rangle & = \langle k_{j}, fg \rangle = \overline{\langle fg, k_{j} \rangle} = \overline{f(z_{j})g(z_{j})} \\
& = \langle \overline{f(z_{j})}k_{j}, g \rangle,
\end{align*}
故 $T_{f}^{*}k_{j} = \overline{f(z_{j})}k_{j}$.
\end{proof}

\begin{lemma}\label{lem:2.17.4}
设 $q = B_{n}$.此时对任意 $f \in H^{\infty}$,$T_{z}T_{f} = T_{f}T_{z}$.而且若 A 是 N 上的有界线性算子且 $T_{z}A = AT_{z}$,则对任意 $h \in N$ 有 $T_{h}A = AT_{h}$.
\end{lemma}

\begin{proof}
对任意 $f, g \in H^{\infty}$,
\begin{align*}
T_{f}T_{g} & = PM_{f}PM_{g}|N = PM_{f}(P + I - P)M_{g}|N \\
& = PM_{fg}|N = T_{fg},
\end{align*}
这里用到 $(I-P)H^{2} = B_{n}H^{2}$ 故 $M_{f}(I-P)M_{g}N \subset B_{n}H^{2}$.这证明了前半部分.又当 $T_{z}A = AT_{z}$ 时,若 f 是 z 的多项式则 $T_{f}A = AT_{f}$.若 $h \in N$,由引理~\ref{lem:2.17.2},h 可由 z 的多项式一致逼近.利用这些事实与 $\|T_{f}\| \leq \|f\|_{\infty}$,可以证明 $T_{h}A = AT_{h}$.
\end{proof}

\begin{theorem}[Sarason]\label{thm:2.17.5}
设 $q = B_{n}$.若 A 是 N 上的有界线性算子且 $T_{z}A = AT_{z}$,则存在 $f \in H^{\infty}$,使得 $A = T_{f}$ 且 $\|A\| = \|f\|_{\infty}$.
\end{theorem}

\begin{proof}
由引理~\ref{lem:2.17.2},$N \subset H^{\infty}$,令 F = AP1,则 $F \in H^{\infty}$.由引理~\ref{lem:2.17.4},设 h, g $\in N$,则 $T_{h}T_{F} = T_{F}T_{h}$,故
\begin{align*}
\langle Ah, g \rangle & = \langle AT_{h}P1, g \rangle = \langle T_{h}AP1, g \rangle \\
& = \langle T_{h}T_{F}P1, g \rangle = \langle T_{F}T_{h}P1, g \rangle \\
& = \langle T_{F}h, g \rangle .
\end{align*}
因 $h, g \in N$ 任意,$A = T_{F}$,且 $\|A\| \leq \|F\|_{\infty}$.

若 $F - G \in B_{n}H^{\infty}$ 则 $T_{F} = T_{G}$,故 $\|A\| \leq \|F + B_{n}H^{\infty}\|$.若能证明 $\|A\| = \|F + B_{n}H^{\infty}\|$,则由定理~\ref{thm:2.14.2} 可证明存在 f 使得 $A = T_{f}$ 且 $\|A\| = \|f\|_{\infty}$.设 $h \in H^{2}$,$k \in B_{n}\overline{H}_{0}^{2}$,则
\begin{align*}
\langle Fh, k \rangle & = \langle FPh, k \rangle + \langle F(I - P)h, k \rangle \\
& = \langle FPh, k \rangle = \langle PFPh, k \rangle .
\end{align*}
于是令 $k = B_{n}\bar{g}$,则
\begin{equation*}
\left| \int_{-\pi}^{\pi} \bar{B}_{n}Fhg \dd \theta / 2\pi \right| \leq \|T_{F}\|\|h\|_{2}\|g\|_{2}.
\end{equation*}
任意 $\ell \in e^{i\theta}H^{1}$ 由定理~\ref{thm:2.8.3} 可因式分解为 $\ell = QL$,令 $h = L^{1/2}$ 且 $g = QL^{1/2}$,则 $\|\ell\|_{1} = \|h\|_{2}\|g\|_{2}$.于是
\begin{equation*}
\left| \int_{-\pi}^{\pi} \bar{B}_{n}F\ell \dd \theta / 2\pi \right| \leq \|T_{F}\|\|\ell\|_{1},
\end{equation*}
由定理~\ref{thm:2.14.2},
\begin{equation*}
\|\bar{B}_{n}F + H^{\infty}\| \leq \|T_{F}\|,
\end{equation*}
故可证 $\|T_{F}\| = \|F + B_{n}H^{\infty}\|$.
\end{proof}

\begin{lemma}\label{lem:2.17.6}
设 $q = B_{n}$.若 $B(1 - \bar{z}_{j}z)^{-1} = \bar{w}_{j}(1 - \bar{z}_{j}z)^{-1}$($1 \leq j \leq n$),则 B 可扩张为 N 上的有界线性算子 $\tilde{B}$,且 $\|\tilde{B}\| \leq 1$ 的充要条件是
\begin{equation*}
\sum_{j,k=1}^{n} \frac{1 - w_{j}\bar{w}_{k}}{1 - z_{j}\bar{z}_{k}} \alpha_{j}\bar{\alpha}_{k} \geq 0 \quad (\{\alpha_{j}\}_{j=1}^{n} \subset \mathbb{C}).
\end{equation*}
\end{lemma}

\begin{proof}
由引理~\ref{lem:2.17.2},若 $h \in N$,则 $h = \sum_{j=1}^{n} \alpha_{j}(1 - \bar{z}_{j}z)^{-1}$,令
\begin{equation*}
\tilde{B}h = \sum_{j=1}^{n} \alpha_{j}\bar{w}_{j}(1 - \bar{z}_{j}z)^{-1},
\end{equation*}
则 $\tilde{B}$ 是 N 上的有界线性算子.$\|h\|_{2}^{2} = \sum_{j,k}\alpha_{j}\bar{\alpha}_{k}(1 - \bar{z}_{j}z_{k})^{-1}$,且
\begin{equation*}
\|\tilde{B}h\|_{2}^{2} = \sum_{j,k}\alpha_{j}\bar{\alpha}_{k}\bar{w}_{j}w_{k}(1 - \bar{z}_{j}z_{k})^{-1}.
\end{equation*}
于是 $\|\tilde{B}h\|_{2} \leq \|h\|_{2}$ 的充要条件是,对任意 $\{\alpha_{j}\}_{j=1}^{n}$,
\begin{equation*}
\sum_{j,k=1}^{n} \frac{1 - w_{j}\bar{w}_{k}}{1 - z_{j}\bar{z}_{k}} \bar{\alpha}_{j}\alpha_{k} \geq 0.
\end{equation*}
\end{proof}

\begin{theorem}[Pick]\label{thm:2.17.7}
设 $\{z_{j}\}_{j=1}^{n} \subset D$ 且 $\{w_{j}\}_{j=1}^{n} \subset \mathbb{C}$.存在满足 $f(z_{j}) = w_{j}$($1 \leq j \leq n$)的 $f \in H^{\infty}$,且 $\|f\|_{\infty} \leq 1$ 的充要条件是
\begin{equation*}
\sum_{j,k=1}^{n} \frac{1 - w_{j}\bar{w}_{k}}{1 - z_{j}\bar{z}_{k}} \alpha_{j}\bar{\alpha}_{k} \geq 0 \quad (\{\alpha_{j}\}_{j=1}^{n} \subset \mathbb{C}).
\end{equation*}
\end{theorem}

\begin{proof}
若存在满足 $f(z_{j}) = w_{j}$($1 \leq j \leq n$)且 $\|f\|_{\infty} \leq 1$ 的 $f \in H^{\infty}$,则由引理~\ref{lem:2.17.3} 与 $\|T_{f}^{*}\| = \|T_{f}\|$ 和引理~\ref{lem:2.17.6},
\begin{equation*}
\sum_{j,k=1}^{n} \frac{1 - w_{j}\bar{w}_{k}}{1 - z_{j}\bar{z}_{k}} \alpha_{j}\bar{\alpha}_{k} \geq 0 \quad (\{\alpha_{j}\}_{j=1}^{n} \subset \mathbb{C}).
\end{equation*}
反之,若
\begin{equation*}
\sum_{j,k=1}^{n} \frac{1 - w_{j}\bar{w}_{k}}{1 - z_{j}\bar{z}_{k}} \alpha_{j}\bar{\alpha}_{k} \geq 0,
\end{equation*}
令 $B(1 - \bar{z}_{j}z)^{-1} = \bar{w}_{j}(1 - \bar{z}_{j}z)^{-1}$,由引理~\ref{lem:2.17.6},$\|\tilde{B}\| \leq 1$.因 $\tilde{B}T_{z}^{*} = T_{z}^{*}\tilde{B}$,令 $A = \tilde{B}^{*}$,则 $T_{z}A = AT_{z}$.于是由定理~\ref{thm:2.17.5},存在 $f \in H^{\infty}$,使得 $A = T_{f}$ 且 $\|A\| = \|f\|_{\infty}$.由引理~\ref{lem:2.17.3},$T_{f}^{*}(1 - \bar{z}_{j}z)^{-1} = \overline{f(z_{j})}(1 - \bar{z}_{j}z)^{-1}$,故 $f(z_{j}) = w_{j}$ 且 $\|f\|_{\infty} = \|\tilde{B}\| \leq 1$.
\end{proof}

% ============ 第2章 §2.18 Carleson插值问题 (I)(原书 pp.179-185) ============
\section{Carleson 插值问题 (I)}\label{sec:2.18}
\index{一致分离}
\index{分离}


对 D 中相异的无限个点 $\{z_{j}\}_{j=1}^{\infty}$ 与 $\mathbb{C}$ 中无限个数据 $\{w_{j}\}_{j=1}^{\infty}$,求属于 $H^{p}$ 且满足
\begin{equation*}
f(z_{j}) = w_{j} \quad (1 \leq j < \infty)
\end{equation*}
的函数 f,与 $\{z_{j}\}$ 为有限个点的情况不同,不一定总存在.

如 2.17 节那样,对相异的 $\{z_{j}\}_{j=1}^{n}$ 定义 Blaschke 积 $B_{n}(z)$ 与 $B_{nj}(z)$,令
\begin{equation*}
f_{n}(z) = \sum_{j=1}^{n} w_{j}\frac{B_{nj}(z)}{B_{nj}(z_{j})}.
\end{equation*}
此时 $f_{n}(z_{j}) = w_{j}$($1 \leq j \leq n$).下面的命题~\ref{prop:2.18.1} 给出无限个点的插值问题的一般解答.

\begin{proposition}\label{prop:2.18.1}
设 $1 \leq p \leq \infty$ 且 $1/p + 1/q = 1$.又设 $\{z_{j}\}_{j=1}^{\infty} \subset D$ 且 $\{w_{j}\}_{j=1}^{\infty} \subset \mathbb{C}$.此时存在满足 $f(z_{j}) = w_{j}$($1 \leq j < \infty$)的 $f \in H^{p}$ 的充要条件是,存在常数 $\gamma > 0$,使得对所有 $g \in H^{q}$,
\begin{equation*}
\left| \sum_{j=1}^{n} \frac{w_{j}}{B_{nj}(z_{j})}(1 - |z_{j}|^{2}) g(z_{j}) \right| \leq \gamma\|g\|_{q} \quad (n = 1, 2, \dots).
\end{equation*}
\end{proposition}

\begin{proof}
设 $w = (w_{j})_{j=1}^{\infty}$,对 $n \geq 1$,令
\begin{equation*}
m_{p,n}(w) = \inf \{\|f_{n} + B_{n}h\|_{p}; h \in H^{p}\},
\end{equation*}
由定理~\ref{thm:2.17.1},
\begin{align*}
m_{p,n}(w) &= \sup \left\{\left| \sum_{j=1}^{n} \frac{w_{j}}{B_{nj}(z_{j})}(1 - |z_{j}|^{2}) g(z_{j}) \right|; \right. \\
& \left. g \in H^{q} \ \text{且} \ \|g\|_{q} \leq 1 \right\}.
\end{align*}
设 $\gamma = \sup_{n}m_{p,n}(w) < \infty$.由定理~\ref{thm:2.14.2},存在 $F_{n} \in f_{n} + B_{n}H^{p}$ 且 $\|F_{n}\|_{p} = m_{p,n}(w)$.此时 $\|F_{n}\|_{p} \leq \gamma$($n \geq 1$)且 $F_{n}(z_{j}) = w_{j}$($1 \leq j \leq n$).由定理~\ref{thm:2.13.2} 与定理~\ref{thm:2.13.3},$H^{p}$ 的单位球是 *弱紧的,故 $\{F_{n}\}$ 存在子列 $\{F_{n_{k}}\}$,$F_{n_{k}}$ 在 *弱拓扑下收敛到某个 $f \in H^{p}$.由命题 2.9.4,
\begin{equation*}
\phi_{j}(h) = h(z_{j}) = \frac{1}{2\pi i}\int_{|z|=1} \frac{h(z)}{z - z_{j}}\dd z \quad (h \in H^{p}),
\end{equation*}
因 $(z - z_{j})^{-1} \in L^{q} \cap C$,$\phi_{j}$ 是 *弱连续的.于是 $f(z_{j}) = w_{j}$($1 \leq j < \infty$).反之,若存在满足 $f(z_{j}) = w_{j}$($1 \leq j < \infty$)的 $f \in H^{p}$,则
\begin{equation*}
\inf \{\|f + B_{n}h\|_{p}; h \in H^{p}\} \leq \|f\|_{p},
\end{equation*}
故 $\sup_{n}m_{p,n}(w) < \infty$.定理得证.
\end{proof}

命题~\ref{prop:2.18.1} 作为无限个点的插值问题的解答完全不充分.下面的命题~\ref{prop:2.18.2} 给出对特殊数据 $\{w_{j}\}_{j=1}^{\infty}$ 的解答.设 $\delta_{\ell j} = 0$($\ell \neq j$)且 $\delta_{\ell j} = 1$($\ell = j$).

\begin{proposition}\label{prop:2.18.2}
设 $0 < p \leq \infty$ 且 $(w_{j})_{j=1}^{\infty} = (\delta_{\ell j})_{j=1}^{\infty}$.存在满足
\begin{equation*}
f(z_{j}) = w_{j} \quad (1 \leq j < \infty)
\end{equation*}
的 $f \in H^{p}$ 的充要条件是 $\sum_{j=1}^{\infty}(1 - |z_{j}|) < \infty$.
\end{proposition}

\begin{proof}
若 $\sum_{j=1}^{\infty}(1 - |z_{j}|) < \infty$,由定理~\ref{thm:2.6.1},存在对 $\{z_{j}\}_{j=1}^{\infty}$ 的 Blaschke 积 B.对任意自然数 $\ell$,令
\begin{equation*}
B_{\ell}(z) = B(z)\frac{z_{\ell}}{|z_{\ell}|}\frac{1 - \bar{z}_{\ell}z}{z_{\ell} - z},
\end{equation*}
令
\begin{equation*}
f(z) = B_{\ell}(z)/B_{\ell}(z_{\ell}),
\end{equation*}
则 $f \in H^{p}$ 且 $f(z_{j}) = \delta_{\ell j}$($1 \leq j < \infty$).反之,若存在满足 $f(z_{j}) = \delta_{\ell j}$($1 \leq j < \infty$)的 $f \in H^{p}$,由定理~\ref{thm:2.6.3} 得 $\sum_{j=1}^{\infty}(1 - |z_{j}|) < \infty$.
\end{proof}

令
\begin{equation*}
\ell^{0} = \{(w_{j}); (w_{j}) = (\delta_{\ell j}), 1 \leq \ell < \infty\},
\end{equation*}
若 $\sum_{j=1}^{\infty}(1 - |z_{j}|) < \infty$,由命题~\ref{prop:2.18.2},$\{(f(z_{j})); f \in H^{p}\} \supset \ell^{0}$.设 $\ell^{\infty} = \{(w_{j}); \sup_{j}|w_{j}| < \infty\}$,则
\begin{equation*}
\ell^{\infty} \supset \{(f(z_{j})); f \in H^{\infty}\} \supset \ell^{0}.
\end{equation*}
弄清 $\{(f(z_{j})); f \in H^{\infty}\}$ 是 $\ell^{\infty}$ 的怎样的子空间是最终目标,Carleson \cite{ch2:16} 给出了 $\{(f(z_{j})); f \in H^{\infty}\} = \ell^{\infty}$ 成立的充要条件.这一 Carleson 的结果被 Shapiro 与 Shields \cite{ch2:17} 推广到 $1 \leq p < \infty$,被 Kabaila \cite{ch2:18} 推广到 $0 < p < 1$.但他们研究的不是 $\{(f(z_{j})); f \in H^{p}\}$,而是 $\{((1 - |z_{j}|^{2})^{1/p}f(z_{j})); f \in H^{p}\}$.对 $0 < p < \infty$,设 $\ell^{p} = \{(w_{j}); \sum_{j=1}^{\infty}|w_{j}|^{p} < \infty\}$.

若 $\{z_{j}\}_{j=1}^{\infty} \subset D$ 满足 $\sum_{j=1}^{\infty}(1 - |z_{j}|) < \infty$,由定理~\ref{thm:2.6.1} 存在 Blaschke 积
\begin{equation*}
B(z) = \prod_{j=1}^{\infty} \frac{|z_{j}|}{z_{j}}\frac{z_{j} - z}{1 - \bar{z}_{j}z}.
\end{equation*}
\begin{equation*}
\begin{aligned}
\rho_{k,n} &= \prod_{\substack{j=1\\ j\neq k}}^{n}\left|\frac{z_{j} - z_{k}}{1 - \bar{z}_{j}z_{k}}\right|\\
\rho_{k} &= \prod_{\substack{j=1\\ j\neq k}}^{\infty}\left|\frac{z_{j} - z_{k}}{1 - \bar{z}_{j}z_{k}}\right|
\end{aligned} \quad (1 \leq k \leq n),
\end{equation*}
则 $\rho_{k,n} \geq \rho_{k,n+1}$ 且 $\lim_{n \to \infty}\rho_{k,n} = \rho_{k}$($k \geq 1$).一般 $\rho_{k} > 0$ 未必成立,但当 $\rho_{k} > 0$($k \geq 1$)时称 $\{z_{j}\}_{j=1}^{\infty}$ 是分离的.当 $\inf_{k}\rho_{k} > 0$ 时称 $\{z_{j}\}_{j=1}^{\infty}$ 是一致分离的.

\begin{theorem}\label{thm:2.18.3}
设 $0 < p \leq \infty$.$\{((1 - |z_{j}|^{2})^{1/p}f(z_{j})); f \in H^{p}\} = \ell^{p}$ 成立的充要条件是 $\{z_{j}\}_{j=1}^{\infty}$ 一致分离.
\end{theorem}

设 $\ell^{p}((1 - |z|^{2})^{1/p}) = \{(w_{j}); \sum_{j=1}^{\infty}(1 - |z_{j}|^{2})|w_{j}|^{p} < \infty\}$,则 $\{((1 - |z_{j}|^{2})^{1/p}f(z_{j})); f \in H^{p}\} = \ell^{p}$ 与 $\{(f(z_{j})); f \in H^{p}\} = \ell^{p}((1 - |z|^{2})^{1/p})$ 等价.

下面证明定理~\ref{thm:2.18.3},p = 2 的情形(准确地说充分条件证明的一部分)证明后,$p \neq 2$ 的情形可由之导出.下面的命题~\ref{prop:2.18.4} 的证明将在 2.19 节给出,本节利用它给出目标定理~\ref{thm:2.18.3} 的证明.

\begin{proposition}\label{prop:2.18.4}
$\{((1 - |z_{j}|^{2})^{1/2}f(z_{j})); f \in H^{2}\} = \ell^{2}$ 成立的充要条件是 $\{z_{j}\}_{j=1}^{\infty}$ 一致分离.
\end{proposition}

\begin{proof}
引理~\ref{lem:2.18.5} 证明必要性.若 $\inf_{j}\rho_{j} > 0$,则 $\{((1 - |z_{j}|^{2})^{1/2}f(z_{j})); f \in H^{2}\} = \ell^{2}$ 将在 2.19 节证明.
\end{proof}

\begin{lemma}\label{lem:2.18.5}
设 $0 < p \leq \infty$.若 $\{((1 - |z_{j}|^{2})^{1/p}f(z_{j})); f \in H^{p}\} = \ell^{p}$,则 $\inf_{j}\rho_{j} > 0$.
\end{lemma}

\begin{proof}
令 $T_{p}(f) = ((1 - |z_{j}|^{2})^{1/p}f(z_{j}))$,由假设 $T_{p}(H^{p}) = \ell^{p}$.设 $f_{n} \in H^{p}$,$f_{n} \to f$ 且
\begin{equation*}
T_{p}(f_{n}) = ((1 - |z_{j}|^{2})^{1/p}f_{n}(z_{j})) \to w \in \ell^{p},
\end{equation*}
由 $f_{n}(z_{j}) \to f(z_{j})$($j \geq 1$),$T_{p}(f) = w$.这表明 $T_{p}$ 是闭算子,由闭图像定理,$T_{p}$ 是从 $H^{p}$ 到 $\ell^{p}$ 的有界算子.

令
\begin{equation*}
\ker T_{p} = \{f \in H^{p}; f(z_{j}) = 0 \ (j \geq 1)\},
\end{equation*}
因 $\ker T_{p}$ 是算子 $T_{p}$ 的核,由 $T_{p}(H^{p}) = \ell^{p}$,可从 $T_{p}$ 定义从 $H^{p}/\ker T_{p}$ 到 $\ell^{p}$ 的双射算子 $\underline{T_{p}}$,即令 $\underline{T_{p}}(f + \ker T_{p}) = T_{p}(f)$($f \in H^{p}$).此时 $\underline{T_{p}}$ 可逆,故对任意 $w \in \ell^{p}$,
\begin{equation*}
(1 - |z_{j}|^{2})^{1/p}f(z_{j}) = w_{j} \quad (j \geq 1)
\end{equation*}
且设 $\gamma = \|\underline{T_{p}}^{-1}\|$,则存在满足
\begin{equation*}
\|f\|_{p} \leq \gamma\|w\|_{p}
\end{equation*}
的 $f \in H^{p}$.于是对每个 k,存在 $f_{k} \in H^{p}$,使得 $\|f_{k}\|_{p} \leq \gamma$ 且
\begin{equation*}
f_{k}(z_{j}) = \left\{ \begin{array}{l l}
(1 - |z_{k}|^{2})^{-1/p} & (j = k) \\
0 & (j \neq k)
\end{array} \right..
\end{equation*}
对 n > k,令
\begin{equation*}
F_{nk}(z) = f_{k}(z)\prod_{\substack{j=1\\ j\neq k}}^{n}\frac{1 - \bar{z}_{j}z}{z_{j} - z},
\end{equation*}
则 $F_{nk} \in H^{p}$ 且 $\|F_{nk}\|_{p} = \|f_{k}\|_{p} \leq \gamma$.由引理~\ref{lem:2.10.1},
\begin{equation*}
|F_{nk}(z_{k})| \leq 2^{1/p}\|F_{nk}\|_{p}(1 - |z_{k}|)^{-1/p},
\end{equation*}
故
\begin{equation*}
\left|\prod_{\substack{j=1\\ j\neq k}}^{n}\frac{1 - \bar{z}_{j}z_{k}}{z_{j} - z_{k}}\right| = |f_{k}(z_{k})|^{-1}|F_{nk}(z_{k})| \leq 2^{1/p}\gamma.
\end{equation*}
这证明了 $\inf_{j}\rho_{j} > 0$.
\end{proof}

\begin{lemma}\label{lem:2.18.6}
当 $0 < p \leq \infty$ 时,若 $\inf_{j}\rho_{j} > 0$,则 $\{((1 - |z_{j}|^{2})^{1/p}f(z_{j})); f \in H^{p}\} \subset \ell^{p}$.
\end{lemma}

\begin{proof}
p = $\infty$ 时显然,可设 0 < p < $\infty$.由定理~\ref{thm:2.6.4},$f \in H^{p}$ 可写成 $f(z) = B(z)[F(z)]^{2/p}$,其中 B 是 Blaschke 积,F 是 D 中无零点的 $H^{2}$ 函数.于是由命题~\ref{prop:2.18.4},
\begin{equation*}
\sum_{j=1}^{\infty}(1 - |z_{j}|^{2})|f(z_{j})|^{p} \leq \sum_{j=1}^{\infty}(1 - |z_{j}|^{2})|F(z_{j})|^{2} < \infty.
\end{equation*}
\end{proof}

\begin{lemma}\label{lem:2.18.7}
当 $1 \leq p \leq \infty$ 时,若 $\inf_{j}\rho_{j} > 0$,则 $\{((1 - |z_{j}|^{2})^{1/p}f(z_{j})); f \in H^{p}\} \supset \ell^{p}$.
\end{lemma}

\begin{proof}
由定理~\ref{thm:2.17.1},存在满足
\begin{equation*}
g_{n}(z_{j}) = (1 - |z_{j}|^{2})^{-1/p}w_{j} \quad (1 \leq j \leq n)
\end{equation*}
的最小范数的 $g_{n} \in H^{p}$.又当 $p \neq 1$ 时,设 $1/p + 1/q = 1$,则存在满足 $\|f\|_{q} = 1$ 的 $f \in H^{q}$,使得
\begin{equation*}
\|g_{n}\|_{p} = \left| \sum_{j=1}^{n} \frac{w_{j}(1 - |z_{j}|^{2})^{1/q}}{B_{nj}(z_{j})} f(z_{j}) \right|.
\end{equation*}
将
\begin{equation*}
\frac{|w_{j}|(1 - |z_{j}|^{2})^{1/q}}{|B_{nj}(z_{j})|}|f(z_{j})| \leq \rho_{j}^{-1}\cdot|w_{j}|\cdot(1 - |z_{j}|^{2})^{1/q}|f(z_{j})|
\end{equation*}
代入,由 Hölder 不等式,
\begin{align*}
\|g_{n}\|_{p} &\leq (\inf_{j}\rho_{j})^{-1}\|w\|_{p}\left\{\sum_{j=1}^{n}(1 - |z_{j}|^{2})|f(z_{j})|^{q}\right\}^{1/q} \\
&\leq \gamma(\inf_{j}\rho_{j})^{-1}\|w\|_{p}\|f\|_{q} \\
&\leq \gamma(\inf_{j}\rho_{j})^{-1}\|w\|_{p}.
\end{align*}
这里 $\gamma$ 的存在由引理~\ref{lem:2.18.6} 与闭图像定理保证.

由引理~\ref{lem:2.10.1},$\{g_{n}\}$ 是正规族,故存在 $g \in H(D)$,对每个 R < 1,$\{g_{n}\}$ 的某个子列在 $\{z; |z| < R\}$ 上一致收敛.
\end{proof}

\begin{lemma}\label{lem:2.18.8}
当 $0 < p \leq 1$ 时,若 $\inf_{j}\rho_{j} > 0$,则 $\{((1 - |z_{j}|^{2})^{1/p}f(z_{j})); f \in H^{p}\} \supset \ell^{p}$.
\end{lemma}

\begin{proof}
令 $f_{j}(z) = (1 - \bar{z}_{j}z)^{-2/p}$,
\begin{equation*}
f(z) = \sum_{j=1}^{\infty}(1 - |z_{j}|^{2})^{1/p}w_{j}[B_{j}(z_{j})]^{-1}B_{j}(z)f_{j}(z).
\end{equation*}
设 $(w_{j}) \in \ell^{p}$ 且 $\inf_{j}\rho_{j} > 0$,则对任意 r,f 在 $|z| \leq r < 1$ 上一致收敛,故 $f \in H(D)$.进而由 Minkowski 不等式,
\begin{align*}
M_{p}^{p}(r, f) &\leq (\inf_{j}\rho_{j})^{-p}\sum_{j=1}^{\infty}(1 - |z_{j}|^{2})|w_{j}|^{p}M_{p}^{p}(r, f_{j}) \\
&\leq (\inf_{j}\rho_{j})^{-p}\|w\|_{p}^{p}.
\end{align*}
\end{proof}

定理 2.18.2 的证明 必要条件是引理~\ref{lem:2.18.5}.证明充分性.设 $\inf_{j}\rho_{j} > 0$,由利用命题~\ref{prop:2.18.4}(尚未证明的充分性)证明的引理~\ref{lem:2.18.6},当 0 < p $\leq \infty$ 时
\begin{equation*}
\{((1 - |z_{j}|^{2})^{1/p}f(z_{j})); f \in H^{p}\} \subset \ell^{p}.
\end{equation*}
关于相反的包含关系,当 0 < p $\leq$ 1 时,由引理~\ref{lem:2.18.8},
\begin{equation*}
\{((1 - |z_{j}|^{2})^{1/p}f(z_{j})); f \in H^{p}\} \supset \ell^{p}.
\end{equation*}
当 $1 \leq p \leq \infty$ 时由引理~\ref{lem:2.18.7} 证明.定理 2.18.2 的证明中还剩下的只有命题~\ref{prop:2.18.4} 的充分性.

% ============ 第2章 §2.19 Carleson插值问题 (II)(原书 pp.185-190) ============
\section{Carleson 插值问题 (II)}\label{sec:2.19}

本节证明命题~\ref{prop:2.18.4} 的充分性.此时由引理~\ref{lem:2.18.6} 完成命题~\ref{prop:2.18.4} 的证明.另外给出一致分离的 $\{z_{j}\}_{j=1}^{\infty}$ 的例子.

\begin{lemma}\label{lem:2.19.1}
若 $\rho_{j} > 0$($j \geq 1$),则
\begin{equation*}
\sum_{j=1}^{\infty}\frac{(1 - |z_{j}|^{2})(1 - |z_{k}|^{2})}{|1 - \bar{z}_{j}z_{k}|^{2}} \leq 1 - 2\log\rho_{k} \quad (k \geq 1).
\end{equation*}
\end{lemma}

\begin{proof}
令 $\alpha_{jk} = (1 - |z_{j}|^{2})(1 - |z_{k}|^{2})/|1 - \bar{z}_{j}z_{k}|^{2}$,由简单计算,
\begin{equation*}
\left| \frac{z_{k} - z_{j}}{1 - \bar{z}_{j}z_{k}} \right|^{2} = 1 - \alpha_{jk}.
\end{equation*}
于是
\begin{align*}
\rho_{k}^{2} &= \prod_{j \neq k}\left| \frac{z_{k} - z_{j}}{1 - \bar{z}_{j}z_{k}} \right|^{2} = \prod_{j \neq k}(1 - \alpha_{jk}) \\
&\leq \exp\left\{-\sum_{j \neq k}\alpha_{jk}\right\},
\end{align*}
利用 $0 < \alpha_{jk} \leq 1$ 即证引理.
\end{proof}

\begin{lemma}\label{lem:2.19.2}
设 $a_{jk} \in \mathbb{C}$($1 \leq j, k \leq n$)满足 $a_{kj} = \overline{a_{jk}}$ 且
\begin{equation*}
\sum_{j=1}^{n}|a_{jk}| \leq M \quad (1 \leq k \leq n).
\end{equation*}
此时对任意 $\{x_{j}\}_{j=1}^{n} \subset \mathbb{C}$,
\begin{equation*}
\left| \sum_{j,k=1}^{n} a_{jk}x_{j}\bar{x}_{k} \right| \leq M\sum_{j=1}^{n}|x_{j}|^{2}.
\end{equation*}
\end{lemma}

\begin{proof}
用 Schwartz 不等式,
\begin{align*}
\sum_{j,k=1}^{n}|a_{jk}x_{j}\bar{x}_{k}|&\leq \sum_{k=1}^{n}\left(\sum_{j=1}^{n}|a_{jk}||x_{j}|^{2}\right)^{1/2}\left(\sum_{j=1}^{n}|a_{jk}||x_{k}|^{2}\right)^{1/2}  \\
&\leq M^{1/2}\sum_{k=1}^{n}\left(\sum_{j=1}^{n}|a_{jk}||x_{j}|^{2}\right)^{1/2}|x_{k}| \\
&\leq M^{1/2}\left\{\sum_{k=1}^{n}\left(\sum_{j=1}^{n}|a_{jk}||x_{j}|^{2}\right)\right\}^{1/2}\left(\sum_{k=1}^{n}|x_{k}|^{2}\right)^{1/2} \\
&= M^{1/2}\left\{\sum_{j=1}^{n}\left(\sum_{k=1}^{n}|a_{jk}|\right)|x_{j}|^{2}\right\}^{1/2}\left(\sum_{k=1}^{n}|x_{k}|^{2}\right)^{1/2} \\
&\leq M\left(\sum_{j=1}^{n}|x_{j}|^{2}\right)^{1/2}\left(\sum_{k=1}^{n}|x_{k}|^{2}\right)^{1/2}.
\end{align*}
\end{proof}

如 2.17 节那样在 $H^{2}$ 上定义内积 $\langle\ ,\ \rangle$.即对 $f, g \in H^{2}$,
\begin{equation*}
\langle f, g \rangle = \frac{1}{2\pi}\int_{-\pi}^{\pi} f(e^{i\theta})\overline{g(e^{i\theta})}\dd \theta.
\end{equation*}
设 $B_{n}(z)$ 与 $B_{nk}(z)$($k \leq n$)为 2.17 节定义的 Blaschke 积.证明当 $\inf_{j}\rho_{j} > 0$ 时 $\{((1 - |z_{j}|^{2})^{1/2}f(z_{j})); f \in H^{2}\} = \ell^{2}$.

设 $w = (w_{k}) \in \ell^{2}$ 为任意元素,令
\begin{equation*}
F_{nk}(z) = (1 - |z_{k}|^{2})^{3/2}[B_{n}(z)]^{2}(z - z_{k})^{-2}
\end{equation*}
且
\begin{equation*}
f_{n}(z) = \sum_{k=1}^{n} w_{k}[B_{nk}(z_{k})]^{-2}F_{nk}(z).
\end{equation*}
此时
\begin{equation*}
(1 - |z_{k}|^{2})^{1/2}f_{n}(z_{k}) = w_{k} \quad (1 \leq k \leq n)
\end{equation*}
且
\begin{align*}
& \|f_{n}\|_{2}^{2} = \langle f_{n}, f_{n} \rangle \\
& \qquad = \sum_{j,k=1}^{n} w_{j}[B_{nj}(z_{j})]^{-2}\bar{w}_{k}[\overline{B_{nk}}(z_{k})]^{-2}\langle F_{nj}, F_{nk} \rangle .
\end{align*}
由留数定理,
\begin{equation*}
\langle F_{nj}, F_{nk} \rangle = (1 - |z_{j}|^{2})^{3/2}(1 - |z_{k}|^{2})^{3/2}(1 + z_{j}\bar{z}_{k})(1 - z_{j}\bar{z}_{k})^{-3},
\end{equation*}
由 $(1 - |z_{j}|^{2})(1 - |z_{k}|^{2}) \leq |1 - z_{j}\bar{z}_{k}|^{2}$,
\begin{equation*}
|\langle F_{nj}, F_{nk} \rangle| \leq 2(1 - |z_{j}|^{2})(1 - |z_{k}|^{2})|1 - z_{j}\bar{z}_{k}|^{-2}.
\end{equation*}
由引理~\ref{lem:2.19.1},
\begin{equation*}
\sum_{j=1}^{n}|\langle F_{nj}, F_{nk} \rangle| \leq 2(1 - 2\log\rho_{k}) \quad (1 \leq k \leq n),
\end{equation*}
由引理~\ref{lem:2.19.2},
\begin{align*}
\|f_{n}\|_{2}^{2} &\leq \sum_{j,k=1}^{n}|w_{j}|\rho_{j,n}^{-2}|w_{k}|\rho_{k,n}^{-2}|\langle F_{nj}, F_{nk} \rangle| \\
&\leq 2(1 - 2\log\inf_{1 \leq k \leq n}\rho_{k})\sum_{j=1}^{n}|w_{j}|^{2}\rho_{j,n}^{-4} \\
&\leq 2(1 - 2\log\inf_{1 \leq k \leq n}\rho_{k})\sup_{1 \leq j \leq n}\rho_{j}^{-4}\sum_{j=1}^{n}|w_{j}|^{2} \\
&\leq 2(1 - 2\log\inf_{j}\rho_{j})(\inf_{j}\rho_{j})^{-4}\sum_{j=1}^{\infty}|w_{j}|^{2}.
\end{align*}
如此由引理~\ref{lem:2.10.1} 可证 $\{f_{n}\}$ 是正规族.于是存在 $f \in H^{2}$,对每个 R < 1,$\{f_{n}\}$ 的某个子列在 $\{z; |z| < R\}$ 上一致收敛.于是
\begin{equation*}
(1 - |z_{j}|^{2})^{1/2}f(z_{j}) = w_{j} \quad (j \geq 1)
\end{equation*}
且
\begin{equation*}
\|f\|_{2} \leq 2(1 - 2\log\inf_{j}\rho_{j})(\inf_{j}\rho_{j})^{-2}\|w\|_{2}.
\end{equation*}
如此证明了 $\{((1 - |z_{j}|^{2})^{1/2}f(z_{j})); f \in H^{2}\} \supset \ell^{2}$.

还必须证明 $\{((1 - |z_{j}|^{2})^{1/2}f(z_{j})); f \in H^{2}\} \subset \ell^{2}$.对任意 $w = (w_{j}) \in \ell^{2}$,由前半的证明,设 $\gamma = 2(1 - 2\log\inf_{j}\rho_{j})(\inf_{j}\rho_{j})^{-2}$,则
\begin{equation*}
\min \{\|g\|_{2}; g \in H^{2}, (1 - |z_{j}|^{2})^{1/2}g(z_{j}) = w_{j} \ (1 \leq j \leq n)\} \leq \gamma\|w\|_{2}.
\end{equation*}
于是由定理~\ref{thm:2.17.1},任意 $f \in H^{2}$ 满足下面的不等式:
\begin{equation*}
\left| \sum_{j=1}^{n} \frac{w_{j}(1 - |z_{j}|^{2})^{1/2}}{B_{nj}(z_{j})} f(z_{j}) \right| \leq \gamma\|w\|_{2}\|f\|_{2}.
\end{equation*}
这里利用 $(\ell^{2})^{*} = \ell^{2}$,对任意 n,
\begin{equation*}
\left(\sum_{j=1}^{n}\frac{(1 - |z_{j}|^{2})}{\rho_{j,n}^{2}}|f(z_{j})|^{2}\right)^{1/2} \leq \gamma\|f\|_{2}.
\end{equation*}
于是设 $\gamma_{0} = 2(1 - 2\log\inf_{j}\rho_{j})(\inf_{j}\rho_{j})^{-1}$,则
\begin{equation*}
\left(\sum_{j=1}^{\infty}(1 - |z_{j}|^{2})|f(z_{j})|^{2}\right)^{1/2} \leq \gamma_{0}\|f\|_{2},
\end{equation*}
故导出 $\{((1 - |z_{j}|^{2})^{1/2}f(z_{j})); f \in H^{2}\} \subset \ell^{2}$.于是与前半合在一起 $\{((1 - |z_{j}|^{2})^{1/2}f(z_{j})); f \in H^{2}\} = \ell^{2}$.

接下来的 Hayman \cite{ch2:19} 与 Newman \cite{ch2:20} 的例~\ref{exa:2.19.4} 回答了是否存在一致分离的 $\{z_{j}\}_{j=1}^{\infty}$ 这一问题.

\begin{lemma}\label{lem:2.19.3}
若 $|\alpha| < 1$ 且 $|\beta| < 1$,则
\begin{equation*}
\left| \frac{\alpha - \beta}{1 - \bar{\alpha}\beta} \right| \geq \frac{|\alpha| - |\beta|}{1 - |\alpha\beta|}.
\end{equation*}
\end{lemma}

\begin{proof}
当 $1 \leq A < B$ 时,
\begin{equation*}
\frac{A + x}{B + x} \geq \frac{A - 1}{B - 1} \quad (-1 \leq x \leq 1),
\end{equation*}
令 a > 0,$A = \frac{a^{2} + r^{2}}{2ar}$,$B = \frac{1 + a^{2}r^{2}}{2ar}$ 且 $x = -\cos\theta$,则
\begin{equation*}
\left| \frac{a - re^{i\theta}}{1 - are^{i\theta}} \right|^{2} = \frac{A + x}{B + x} \geq \frac{A - 1}{B - 1} = \left(\frac{a - r}{1 - ar}\right)^{2}.
\end{equation*}
当 $a = \alpha$,$re^{i\theta} = \beta$ 时即证引理,一般 $\alpha$ 的情形容易导出.
\end{proof}

\begin{example}\label{exa:2.19.4}
当 $\{z_{j}\}_{j=1}^{\infty} \subset D$ 时,若存在 0 < c < 1 使得
\begin{equation*}
1 - |z_{n}| \leq c(1 - |z_{n-1}|) \quad (1 \leq n < \infty),
\end{equation*}
则 $\{z_{j}\}_{j=1}^{\infty}$ 是一致分离的.
\end{example}

\begin{proof}
当 j > k 时,由假设,
\begin{equation*}
1 - |z_{j}| \leq c^{j-k}(1 - |z_{k}|) \quad (1 \leq k < \infty).
\end{equation*}
于是
\begin{equation*}
|z_{j}| - |z_{k}| \geq (1 - c^{j-k})(1 - |z_{k}|)
\end{equation*}
且
\begin{align*}
1 - |z_{j}z_{k}| & = 1 - |z_{j}| + |z_{j}|(1 - |z_{k}|) \\
& \leq (1 + c^{j-k})(1 - |z_{k}|).
\end{align*}
因此由引理~\ref{lem:2.19.3},当 j > k 时,
\begin{equation*}
\left| \frac{z_{k} - z_{j}}{1 - \bar{z}_{j}z_{k}} \right| \geq \frac{|z_{j}| - |z_{k}|}{1 - |z_{j}z_{k}|} \geq \frac{1 - c^{j-k}}{1 + c^{j-k}}.
\end{equation*}
当 j < k 时,上面的不等式为
\begin{equation*}
\left| \frac{z_{k} - z_{j}}{1 - \bar{z}_{j}z_{k}} \right| \geq \frac{1 - c^{k-j}}{1 + c^{k-j}}.
\end{equation*}
如此
\begin{align*}
\prod_{\substack{j=1\\ j\neq k}}^{\infty}\left| \frac{z_{k} - z_{j}}{1 - \bar{z}_{j}z_{k}} \right| &= \prod_{j=1}^{k-1}\left| \frac{z_{k} - z_{j}}{1 - \bar{z}_{j}z_{k}} \right| \times \prod_{j=k+1}^{\infty}\left| \frac{z_{k} - z_{j}}{1 - \bar{z}_{j}z_{k}} \right| \\
&\geq \prod_{n=1}^{k-1}\frac{1 - c^{n}}{1 + c^{n}}\prod_{n=1}^{\infty}\frac{1 - c^{n}}{1 + c^{n}}.
\end{align*}
因 $\sum_{n=1}^{\infty}2c^{n}/(1 + c^{n}) < \infty$,上面的无限乘积 $\prod_{n=1}^{\infty}(1 - c^{n})/(1 + c^{n}) > 0$,故 $\{z_{j}\}_{j=1}^{\infty}$ 是一致分离的.
\end{proof}

Carleson 插值问题针对的是正则函数,关于调和函数的研究见 \cite{ch2:21}.

% ============ 第2章 §2.20 半平面上的Hardy空间(原书 pp.190-200) ============
\section{半平面上的 Hardy 空间}\label{sec:2.20}
\index{Hardy 空间!(上半平面上的)}
\index{Blaschke 积!(上半平面上的)}
\index{内函数!(上半平面上的)}
\index{外函数!(上半平面上的)}
\index{Poisson 核!(上半平面的)}


$\mathcal{H}(\mathcal{D})$ 表示上半平面 $\mathcal{D} = \{z \in \mathbb{C}; \operatorname{Im}(z) > 0\}$ 上正则函数的全体.对 $f \in \mathcal{H}(\mathcal{D})$,当 $0 < y < \infty$ 时令
\begin{equation*}
\mathcal{M}_{p}(y, f) = \left\{\int_{-\infty}^{\infty}|f(x + iy)|^{p}\dd x\right\}^{1/p} \quad (0 < p < \infty).
\end{equation*}
$\mathcal{D}$ 上的 Hardy 空间 $\mathcal{H}^{p}$ 定义如下:
\begin{equation*}
\mathcal{H}^{p} = \left\{f \in \mathcal{H}(\mathcal{D}); \sup_{0 < y < \infty}\mathcal{M}_{p}(y, f) < \infty \right\},
\end{equation*}
且 $\mathcal{H}^{\infty}$ 表示 $\mathcal{D}$ 上有界正则函数的全体.

本节作为单位圆盘 $D$ 上的 Hardy 空间 $H^{p}$ 的结果的应用,叙述上半平面 $\mathcal{D}$ 上的 Hardy 空间 $\mathcal{H}^{p}$.令 $\overline{\mathcal{D}} = \{z \in \mathbb{C}; \operatorname{Im}(z) \geq 0\}$ 且
\begin{equation*}
z = \phi(w) = \frac{i - w}{i + w} \quad (w \in \overline{\mathcal{D}}),
\end{equation*}
则 $\phi$ 给出从 $\overline{\mathcal{D}}$ 到 $\overline{D}$ 的一对一正则映射.
\begin{equation*}
w = \psi(z) = i\frac{1 - z}{1 + z},
\end{equation*}
则 $\psi \circ \phi(w) = w$ 且 $\phi \circ \psi(z) = z$,$\psi$ 是从 $\overline{D}$ 到 $\overline{\mathcal{D}}$ 的一对一正则映射.对 $g \in H(D)$,$g \circ \phi \in \mathcal{H}(\mathcal{D})$;反之,对 $f \in \mathcal{H}(\mathcal{D})$,$f \circ \psi \in H(D)$.但对 $g \in H^{p}$,$g \circ \phi \in \mathcal{H}^{p}$ 未必成立.实际上,对某个 $0 < \alpha < 1$,$g(z) = 1/(1-z)^{\alpha}$ 属于 $H^{1}$,但容易证明 $g \notin \bigcup_{p \geq 1}\mathcal{H}^{p}$.另一方面,对 $f \in \mathcal{H}^{p}$ 可以证明 $f \circ \psi \in H^{p}$,故 2.19 节为止关于 $H^{p}$ 的结果可用于 $\mathcal{H}^{p}$ 的研究.另外,若 g 在 $D$ 上调和,则 $g \circ \phi$ 在 $\mathcal{D}$ 上调和;反之,若 f 在 $\mathcal{D}$ 上调和,则 $f \circ \psi$ 在 $D$ 上调和.

在单位圆盘 $D$ 上的 Hardy 空间的研究中,Cauchy 核,Poisson 核与 Herglotz 核起了重要作用.把这 3 个核变换到上半平面 $\mathcal{D}$ 上.写 $\phi(t) = e^{i\tau}$($t \in \partial \mathcal{D}$),Cauchy 核变换为
\begin{equation*}
\frac{1}{1 - z} = \frac{1}{2}\left(1 + \frac{i}{w}\right), \quad \frac{e^{i\tau}}{e^{i\tau} - z} = \frac{1}{2}\left(1 + i\frac{1 + tw}{w - t}\right),
\end{equation*}
Poisson 核变换为
\begin{equation*}
\operatorname{Re}\left(\frac{1 + z}{1 - z}\right) = \operatorname{Re}\left(\frac{i}{w}\right), \quad \operatorname{Re}\left(\frac{e^{i\tau} + z}{e^{i\tau} - z}\right) = \operatorname{Re}\left(i\frac{1 + tw}{w - t}\right),
\end{equation*}
Herglotz 核变换为
\begin{equation*}
\frac{1 + z}{1 - z} = \frac{i}{w}, \quad \frac{e^{i\tau} + z}{e^{i\tau} - z} = i\frac{1 + tw}{w - t}.
\end{equation*}
这些变换后的核在 $\mathcal{D}$ 上再次变得重要.

在 $D$ 上重要的一阶变换如下变换.设 $a = \phi(b)$,则
\begin{equation*}
\frac{a - z}{1 - \bar{a}z} = \frac{\phi(b) - \phi(w)}{1 - \overline{\phi(b)}\phi(w)} = \frac{i - \bar{b}}{i + b}\frac{w - b}{w - \bar{b}},
\end{equation*}
且
\begin{equation*}
\frac{|a|}{a}\frac{a - z}{1 - \bar{a}z} = \frac{|\phi(b)|}{\phi(b)}\frac{\phi(b) - \phi(w)}{1 - \overline{\phi(b)}\phi(w)} = \frac{|1 + b^{2}|}{1 + b^{2}}\frac{w - b}{w - \bar{b}}.
\end{equation*}
若设 $\phi(t) = e^{i\tau}$($-\infty < t < \infty$),则对 $g \in L^{1}[-\pi, \pi]$,令 $f(t) = g \circ \phi(t)$,
\begin{equation*}
\frac{1}{2\pi}\int_{-\pi}^{\pi} g(e^{i\tau})\dd \tau = \frac{1}{\pi}\int_{-\infty}^{\infty} f(t)\frac{1}{1 + t^{2}}\dd t.
\end{equation*}
设 $z = re^{i\theta}$,$w = x + iy$ 且 $z = \phi(w)$,$0 \leq r < 1$,则
\begin{equation*}
\frac{1}{2\pi}\int_{-\pi}^{\pi} g(e^{i\tau})P(r, \theta - \tau)\dd \tau = \frac{1}{\pi}\int_{-\infty}^{\infty} f(t)P(y, x - t)\dd t.
\end{equation*}
其中 $P(y, x - t) = y/((x-t)^{2} + y^{2})$.$P(y, x) = \mathrm{Re}(i/w) = y/(x^{2} + y^{2})$ 称为 $\mathcal{D}$ 上的 Poisson 核.

\begin{theorem}\label{thm:2.20.1}
设 F 在 $(-\infty, \infty)$ 上关于 $(1 + t^{2})^{-1}\dd t$ 可积.当 y > 0 时,
\begin{equation*}
f(x + iy) = \frac{1}{\pi}\int_{-\infty}^{\infty} F(t)P(y, x - t)\dd t,
\end{equation*}
则 f 在 $\mathcal{D}$ 上调和.进而 $\lim_{y \to 0}f(x + iy) = F(x)$ a.e. $x \in (-\infty, \infty)$.这里 $y \to 0$ 指的是 $y = (1 - r^{2})/(1 + r^{2} + 2r\cos\theta) \to 0$($r \to 1$).
\end{theorem}

\begin{proof}
令 $G(e^{i\tau}) = (F \circ \psi)(e^{i\tau}) = F(t)$,因 F 关于 $(1 + t^{2})^{-1}\dd t$ 可积,$G \in L^{1}[-\pi, \pi]$.于是设 $z = re^{i\theta}$,$\psi(z) = w = x + iy$,则
\begin{align*}
g(re^{i\theta}) & = \frac{1}{2\pi}\int_{-\pi}^{\pi} G(e^{i\tau})P(r, \theta - \tau)\dd \tau \\
& = \frac{1}{\pi}\int_{-\infty}^{\infty} F(t)P(y, x - t)\dd t = f(x + iy).
\end{align*}
由
\begin{equation*}
y = \frac{1 - r^{2}}{1 + r^{2} + 2r\cos\theta}, \quad x = \frac{2r\sin\theta}{1 + r^{2} + 2r\cos\theta},
\end{equation*}
由定理~\ref{thm:2.3.2} 及其后的注意,$\lim_{r \to 1} g(re^{i\theta}) = G(e^{i\theta})$ a.e. $\theta \in [-\pi, \pi]$,故可证明 $\lim_{y \to 0}f(x + iy) = F(x)$ a.e. $x \in (-\infty, \infty)$.
\end{proof}

已知沿与 x 轴正交的直线取 $y \to 0$ 时定理~\ref{thm:2.20.1} 也成立.

\begin{lemma}\label{lem:2.20.2}
设 $p \geq 1$ 且 $f \in \mathcal{H}^{p}$.设 $w = \zeta + i\eta$ 且 $y_{0} > 0$,则以下成立.
\begin{equation*}
f(w) = \frac{1}{2\pi i}\int_{-\infty}^{\infty}\frac{f(t + iy_{0})}{t + iy_{0} - w}\dd t \quad (\eta > y_{0}), \qquad (1)
\end{equation*}
\begin{equation*}
\frac{1}{2\pi i}\int_{-\infty}^{\infty}\frac{f(t + iy_{0})}{t + iy_{0} - w}\dd t = 0 \quad (\eta < y_{0}), \qquad (2)
\end{equation*}
\begin{equation*}
f(w) = \frac{1}{\pi}\int_{-\infty}^{\infty} f(t + iy_{0})\frac{\eta - y_{0}}{(t - \zeta)^{2} + (\eta - y_{0})^{2}}\dd t \quad (\eta > y_{0}). \qquad (3)
\end{equation*}
\end{lemma}

\begin{proof}
(1) 固定 w,设 $\operatorname{Im}(w) > 0$.选取 x, y 使得 $-x < \mathrm{Re}(w) < x$,$y_{0} < \mathrm{Im}(w) < y$.由对矩形 $-x + iy_{0}$,$x + iy_{0}$,$x + iy$,$-x + iy$ 的 Cauchy 积分公式,
\begin{align*}
2\pi i f(w) & = \int_{-x}^{x}\frac{f(t + iy_{0})}{t + iy_{0} - w}\dd t - \int_{-x}^{x}\frac{f(t + iy)}{t + iy - w}\dd t \\
& \quad + i\int_{y_{0}}^{y}\frac{f(x + is)}{x + is - w}\dd s - i\int_{y_{0}}^{y}\frac{f(-x + is)}{-x + is - w}\dd s \\
& = I_{1}(x) + I_{2}(x) + I_{3}(x) + I_{4}(x),
\end{align*}
取 $\alpha$ 使 $\alpha > 2|\mathrm{Re}(w)|$,将上式两边在区间 $[\alpha, 2\alpha]$ 上平均,得
\begin{equation*}
2\pi i f(w) = \sum_{j=1}^{4}\frac{1}{\alpha}\int_{\alpha}^{2\alpha}I_{j}(x)\dd x.
\end{equation*}
对每个 j 令
\begin{equation*}
I_{j}(\infty) = \lim_{\alpha \to \infty}\frac{1}{\alpha}\int_{\alpha}^{2\alpha}I_{j}(x)\dd x,
\end{equation*}
则 $I_{3}(\infty) = I_{4}(\infty) = 0$(后文证明),
\begin{equation*}
I_{1}(\infty) = \int_{-\infty}^{\infty}\frac{f(t + iy_{0})}{t + iy_{0} - w}\dd t
\end{equation*}
且
\begin{equation*}
I_{2}(\infty) = -\int_{-\infty}^{\infty}\frac{f(t + iy)}{t + iy - w}\dd t
\end{equation*}
可以证明(后文证明).于是对任意 $y > \operatorname{Im}(w)$,
\begin{equation*}
2\pi i f(w) = \int_{-\infty}^{\infty}\frac{f(t + iy_{0})}{t + iy_{0} - w}\dd t - \int_{-\infty}^{\infty}\frac{f(t + iy)}{t + iy - w}\dd t.
\end{equation*}
令 $y \to \infty$,则可证明
\begin{equation*}
-I_{2}(\infty) = \int_{-\infty}^{\infty}\frac{f(t + iy)}{t + iy - w}\dd t \longrightarrow 0 \ \text{(后文证明)},
\end{equation*}
故
\begin{equation*}
2\pi i f(w) = \int_{-\infty}^{\infty}\frac{f(t + iy_{0})}{t + iy_{0} - w}\dd t \quad (\operatorname{Im}(w) > y_{0}),
\end{equation*}
(1) 的证明结束.

下面证明上述假定成立之处.$I_{1}(\infty)$ 与 $I_{2}(\infty)$ 的积分表示,由于 s = $y_{0}$ 或 s = y 时 $y_{0} < \mathrm{Im}(w) < y$,
\begin{equation*}
\int_{-\infty}^{\infty}\frac{|f(t + is)|}{|t + is - w|}\dd t < \infty,
\end{equation*}
由 $\lim_{x \to \infty}I_{j}(x)$(j = 1,2)存在即明.$I_{2}(\infty)$ 的积分表示在 $y \to \infty$ 时收敛到 0,分 p = 1 与 p > 1 两种情形如下证明.p = 1 时,
\begin{align*}
\left| \int_{-\infty}^{\infty}\frac{f(t + iy)}{t + iy - w}\dd t \right| & \leq \int_{-\infty}^{\infty}\frac{|f(t + iy)|}{y - \mathrm{Im}(w)}\dd t \\
& \leq \frac{1}{y - \mathrm{Im}(w)}\sup_{0 < y < \infty}\mathcal{M}_{1}(y, f) = O\left(\frac{1}{y}\right),
\end{align*}
p > 1 时,设 $1/p + 1/q = 1$,由 Hölder 不等式,
\begin{align*}
\left| \int_{-\infty}^{\infty}\frac{f(t + iy)}{t + iy - w}\dd t \right|&\leq \sup_{0 < y < \infty}\mathcal{M}_{p}(y, f)\left(\int_{-\infty}^{\infty}\frac{\dd t}{|t + iy - w|^{q}}\right)^{1/q} = O(y^{-1 + \frac{1}{q}}).
 \end{align*}
至此,还剩下证明假定的 $I_{3}(\infty) = I_{4}(\infty) = 0$.$I_{4}(\infty) = 0$ 类似,只证明 $I_{3}(\infty) = 0$.因 $\alpha > 2|\mathrm{Re}(w)|$,$\alpha < x < 2\alpha$,故 $|x + is - w| \geq |x - \mathrm{Re}(w)| \geq \alpha/2$,所以
\begin{align*}
\left| \frac{1}{\alpha}\int_{\alpha}^{2\alpha}I_{3}(x)\dd x \right| & \leq \frac{1}{\alpha}\int_{\alpha}^{2\alpha}\dd x\int_{y_{0}}^{y}\left| \frac{f(x + is)}{x + is - w} \right|\dd s \\
& \leq \frac{2}{\alpha^{2}}\int_{y_{0}}^{y}\dd s\int_{\alpha}^{2\alpha}|f(x + is)|\dd x.
\end{align*}
设 $1/p + 1/q = 1$,用 Hölder 不等式,
\begin{equation*}
\int_{\alpha}^{2\alpha}|f(x + is)|\dd x \leq \alpha^{1/q}\sup_{0 < y < \infty}\mathcal{M}_{p}(y, f),
\end{equation*}
代入上式得
\begin{equation*}
\left| \frac{1}{\alpha}\int_{\alpha}^{2\alpha}I_{3}(x)\dd x \right| \leq 2\alpha^{\frac{1}{q} - 2}(y - y_{0})\sup_{0 < y < \infty}\mathcal{M}_{p}(y, f).
\end{equation*}
因 $f \in \mathcal{H}^{p}$,令 $\alpha \to \infty$,
\begin{equation*}
I_{3}(\infty) = \lim_{\alpha \to \infty}\frac{1}{\alpha}\int_{\alpha}^{2\alpha}I_{3}(x)\dd x = 0.
\end{equation*}

(2) 因 $\operatorname{Im}(w) < y_{0}$,对矩形 $-x + iy_{0}$,$x + iy_{0}$,$x + iy$,$-x + iy$ 用 Cauchy 定理,
\begin{align*}
0 &= \int_{-x}^{x}\frac{f(t + iy_{0})}{t + iy_{0} - w}\dd t - \int_{-x}^{x}\frac{f(t + iy)}{t + iy - w}\dd t \\
&\quad + i\int_{y_{0}}^{y}\frac{f(x + is)}{x + is - w}\dd s - i\int_{y_{0}}^{y}\frac{f(-x + is)}{-x + is - w}\dd s \\
\quad &= I_{1}(x) + I_{2}(x) + I_{3}(x) + I_{4}(x).
\end{align*}
与 (1) 相同,可证明 $\lim_{y \to \infty}I_{2}(\infty) = I_{3}(\infty) = I_{4}(\infty) = 0$,故
\begin{equation*}
\frac{1}{2\pi i}\int_{-\infty}^{\infty}\frac{f(t + iy_{0})}{t + iy_{0} - w}\dd t = 0 \quad (\operatorname{Im}(w) < y_{0}).
\end{equation*}

(3) 设 $w = \zeta + i\eta$,$\eta > y_{0}$ 且 $w' = i2y_{0} + \zeta - i\eta$.因 $\operatorname{Im}(w') < y_{0}$,由 (2),
\begin{equation*}
0 = \frac{1}{2\pi i}\int_{-\infty}^{\infty}\frac{f(t + iy_{0})}{t + iy_{0} - w'}\dd t,
\end{equation*}
从 (1) 中 $f(w)$ 的积分表示减去它,
\begin{equation*}
f(w) = \frac{1}{\pi}\int_{-\infty}^{\infty} f(t + iy_{0})\frac{\eta - y_{0}}{(t - \zeta)^{2} + (\eta - y_{0})^{2}}\dd t.
\end{equation*}
\end{proof}

\begin{lemma}\label{lem:2.20.3}
设 $1 \leq p < \infty$.设 $\Gamma = \{z; |z| = R < 1\}$ 且 $\gamma$ 为 $\Gamma$ 内部任意圆周,若 $g \in H(D)$,则
\begin{equation*}
\int_{\gamma}|g(z)|^{p}|\dd z| \leq 2\int_{\Gamma}|g(z)|^{p}|\dd z|.
\end{equation*}
\end{lemma}

\begin{proof}
当 $z \in \gamma$ 时,
\begin{equation*}
g(z) = \frac{1}{2\pi}\int_{-\pi}^{\pi} g(Re^{i\theta})\frac{R^{2} - |z|^{2}}{|Re^{i\theta} - z|^{2}}\dd \theta,
\end{equation*}
用 Hölder 不等式,
\begin{align*}
\int_{\gamma}|g(z)|^{p}|\dd z|&\leq \int_{\gamma}\left(\frac{1}{2\pi}\int_{-\pi}^{\pi}|g(Re^{i\theta})|^{p}\frac{R^{2} - |z|^{2}}{|Re^{i\theta} - z|^{2}}\dd \theta\right)|\dd z|  \\
&\leq \int_{-\pi}^{\pi}|g(Re^{i\theta})|^{p}\left(\frac{1}{2\pi}\int_{\gamma}\frac{R^{2} - |z|^{2}}{|Re^{i\theta} - z|^{2}}|\dd z|\right)\dd \theta.
\end{align*}
设圆 $\gamma$ 的中心为 a,半径为 r,则 $r \leq |Re^{-i\theta} - a|$,由 Gauss 平均值定理,
\begin{equation*}
\frac{1}{2\pi}\int_{\gamma}\frac{R^{2} - |z|^{2}}{|Re^{i\theta} - z|^{2}}|\dd z| = r \times \frac{R^{2} - |a|^{2}}{|Re^{i\theta} - a|^{2}} \leq 2.
\end{equation*}
引理得证.
\end{proof}

\begin{theorem}\label{thm:2.20.4}
设 $1 \leq p \leq \infty$.对 $f \in \mathcal{H}^{p}$ 令 $g = f \circ \psi$,则 $g \in H^{p}$.其中 $\psi(z) = i(1 - z)/(1 + z)$.
\end{theorem}

\begin{proof}
p = $\infty$ 时显然,设 $1 \leq p < \infty$.设 $g = f \circ \psi$,若能证明
\begin{equation*}
\sup_{r}\int_{-\pi}^{\pi}|g(re^{i\theta})|^{p}\dd \theta/2\pi \leq 2\sup_{y}\int_{-\infty}^{\infty}|f(x + iy)|^{p}\dd x/\pi,
\end{equation*}
则由假设 $f \in \mathcal{H}^{p}$,g 属于 $H^{p}$.

对 0 < r < 1 令 $\gamma_{r} = \{z; |z| = r\}$,对 $0 < \delta$ 令
\begin{equation*}
\Gamma_{\delta} = \left\{z; z = \phi(w) = \frac{i - w}{i + w}, \operatorname{Im}(w) = \delta \right\},
\end{equation*}
则 $\Gamma_{\delta}$ 是中心 $-\delta/(1+\delta)$,半径 $1/(1+\delta)$ 的圆.对任意 r < 1,存在 $\delta > 0$,使得 $\gamma_{r}$ 包含在 $\Gamma_{\delta}$ 的内部.由引理~\ref{lem:2.20.2} 的 (3),f 在 $\{w; \operatorname{Im}(w) \geq \delta\} \cup \{\text{无穷远点}\}$ 上连续.因为 $w \to \infty$ 时,在 $\{w; \operatorname{Im}(w) \geq \delta\}$ 上 $f(w) \to 0$.于是 g 在圆盘 $\{z; z = \phi(w), \operatorname{Im}(w) \geq \delta\}$ 上连续且在其内部正则.因 $\gamma_{r}$ 包含于该圆盘内部,由引理~\ref{lem:2.20.3} 定理得证.
\end{proof}

\begin{theorem}\label{thm:2.20.5}
设 $1 \leq p < \infty$ 且 $f \in \mathcal{H}^{p}$.$y \to 0$ 的含义与定理~\ref{thm:2.20.1} 相同.

(1) $\lim_{y \to 0}f(x + iy)$ 在 a.e. $x \in (-\infty, \infty)$ 存在.

(2) $f(x) = \lim_{y \to 0}f(x + iy)$ 属于 $L^{p}(-\infty, \infty)$,且
\begin{equation*}
f(x + iy) = \frac{1}{\pi}\int_{-\infty}^{\infty} f(t)P(y, x - t)\dd t \quad (y > 0).
\end{equation*}
\end{theorem}

\begin{proof}
(1) 令 $g = f \circ \psi$,由定理~\ref{thm:2.20.4},$g \in H^{p}$,由定理~\ref{thm:2.5.3},$\lim_{r \to 1}g(re^{i\theta})$ 在 a.e. $\theta \in [-\pi, \pi]$ 存在.于是定理~\ref{thm:2.20.1} 的证明证明了 (1).

(2) 再令 $g = f \circ \psi$,因 $p \geq 1$,由定理~\ref{thm:2.9.1},
\begin{equation*}
g(z) = \frac{1}{2\pi}\int_{-\pi}^{\pi} g(e^{i\tau})P(r, \theta - \tau)\dd \tau .
\end{equation*}
由定理~\ref{thm:2.20.1} 前的注意,设 $\psi(z) = x + iy$,则
\begin{equation*}
f(x + iy) = \frac{1}{\pi}\int_{-\infty}^{\infty} f(t)P(y, x - t)\dd t.
\end{equation*}
\end{proof}

利用定理~\ref{thm:2.20.4},与定理~\ref{thm:2.20.5} 一样,$p \geq 1$ 时在 $H^{p}$ 中成立的许多结果在 $\mathcal{H}^{p}$ 中也成立.这里叙述属于 $\mathcal{H}^{p}$ 的函数的零点.

设
\begin{equation*}
\{b_{j}\}_{j=1}^{\infty} \subset \mathcal{D} \quad \text{且} \quad \sum_{j=1}^{\infty}\frac{\mathrm{Im}(b_{j})}{|i + b_{j}|^{2}} < \infty,
\end{equation*}
则
\begin{equation*}
\mathcal{B}(w) = \left(\frac{w - i}{w + i}\right)^{m}\prod_{j=1}^{\infty}\frac{|b_{j}^{2} + 1|}{b_{j}^{2} + 1}\cdot\frac{w - b_{j}}{w - \bar{b}_{j}}
\end{equation*}
称为上半平面 $\mathcal{D}$ 上的 Blaschke 积.这里 m 是非负整数,且允许 $b_{j} = b_{i}$($j \neq i$).设 $a_{j} = \phi(b_{j})$,$z = \phi(w)$,则
\begin{equation*}
\frac{|a_{j}|}{a_{j}}\frac{a_{j} - z}{1 - \bar{a}_{j}z} = \frac{|b_{j}^{2} + 1|}{b_{j}^{2} + 1}\cdot\frac{w - b_{j}}{w - \bar{b}_{j}},
\end{equation*}
故令 $B(z) = \mathcal{B} \circ \psi(z)$,即 $B \circ \phi(w) = \mathcal{B}(w)$,则 B(z) 是 D 上 $\{a_{j}\}_{j=1}^{\infty}$ 的 Blaschke 积.实际上
\begin{equation*}
1 - |a_{j}|^{2} = 1 - \left| \frac{i - b_{j}}{i + b_{j}} \right|^{2} = \frac{4\mathrm{Im}(b_{j})}{|i + b_{j}|^{2}},
\end{equation*}
且由
\begin{equation*}
\sum_{j=1}^{\infty}\frac{\operatorname{Im}(b_{j})}{|i + b_{j}|^{2}} < \infty \quad \text{得} \quad \sum_{j=1}^{\infty}(1 - |a_{j}|^{2}) < \infty.
\end{equation*}
由定理~\ref{thm:2.6.1},$B \in H^{\infty}$,故 $\mathcal{B} \in \mathcal{H}^{\infty}$,$|\mathcal{B}(w)| < 1$($w \in \mathcal{D}$)且 $|\mathcal{B}(x)| = 1$ a.e. $x \in (-\infty, \infty)$.$\{b_{j}\}_{j=1}^{\infty}$ 是 $\mathcal{B}$ 的零点,且 $\mathcal{B}$ 在 $\mathcal{D}$ 上没有其他零点.

\begin{theorem}\label{thm:2.20.6}
设 $1 \leq p \leq \infty$.设 $\{b_{j}\}_{j=1}^{\infty}$ 为 $f \in \mathcal{H}^{p}$ 的零点集合,则
\begin{equation*}
\sum_{j=1}^{\infty}\frac{\operatorname{Im}(b_{j})}{|i + b_{j}|^{2}} < \infty.
\end{equation*}
反之,若 $\sum_{j=1}^{\infty}\mathrm{Im}(b_{j})/|i + b_{j}|^{2} < \infty$,则存在以 $\{b_{j}\}_{j=1}^{\infty}$ 为零点集合的 $f \in \mathcal{H}^{p}$.
\end{theorem}

\begin{proof}
令 $g = f \circ \psi$,由定理~\ref{thm:2.20.4},$g \in H^{p}$,设 $a_{j} = \phi(b_{j})$,则 $\{a_{j}\}_{j=1}^{\infty}$ 是 g 的零点集合.由定理~\ref{thm:2.6.3},$\sum_{j=1}^{\infty}(1 - |a_{j}|) < \infty$,故 $\sum_{j=1}^{\infty}\mathrm{Im}(b_{j})/|i + b_{j}|^{2} < \infty$.反之,由本定理上面的注意,设 $\mathcal{B}$ 为 $\{b_{j}\}_{j=1}^{\infty}$ 的 Blaschke 积,则 $\mathcal{B} \in \mathcal{H}^{\infty}$.令 $h(w) = (w + i)^{-2/p}$,则 h 属于 $\mathcal{H}^{p}$ 且在 $\mathcal{D}$ 上没有零点.f = $\mathcal{B}h$ 即为所求.
\end{proof}

\begin{theorem}\label{thm:2.20.7}
设 $1 \leq p \leq \infty$.若 $f \in \mathcal{H}^{p}$ 且 $f(w) \neq 0$,则 $f(w) = \mathcal{B}(w)F(w)$ 可分解.其中 $\mathcal{B}$ 是 Blaschke 积,F 属于 $\mathcal{H}^{p}$,且在 $\mathcal{D}$ 上没有零点.
\end{theorem}

\begin{proof}
因 $g = f \circ \psi \in H^{p}$,由定理~\ref{thm:2.6.4},g = Bh 可分解,其中 B 是 Blaschke 积,$h \in H^{p}$ 在 D 上没有零点.于是由定理~\ref{thm:2.9.1},设 $z = re^{i\theta}$,则
\begin{equation*}
h^{p}(z) = \frac{1}{2\pi}\int_{-\pi}^{\pi} h^{p}(e^{i\tau})P(r, \theta - \tau)\dd \tau
\end{equation*}
可作 Poisson 积分表示.$B \circ \phi$ 是 $\mathcal{D}$ 上的 Blaschke 积,$F = h \circ \phi$ 在 $\mathcal{D}$ 上没有零点.剩下的问题是证明 $F \in \mathcal{H}^{p}$.由定理~\ref{thm:2.20.1} 前的注意,
\begin{equation*}
F^{p}(x + iy) = \frac{1}{\pi}\int_{-\infty}^{\infty} F^{p}(t)P(y, x - t)\dd t \quad (y > 0),
\end{equation*}
利用
\begin{equation*}
\frac{1}{\pi}\int_{-\infty}^{\infty} P(y, x - t)\dd x = \frac{1}{2\pi}\int_{-\pi}^{\pi} P(r, \theta - \tau)\dd \tau = 1,
\end{equation*}
由 Fubini 定理,
\begin{equation*}
\sup_{y > 0}\int_{-\infty}^{\infty}|F(x + iy)|^{p}\dd x < \infty,
\end{equation*}
故可证明 $F \in \mathcal{H}^{p}$.这里用到了 $|F(t)| = |f(t)|$ a.e. t.
\end{proof}

对关于 $(-\infty, \infty)$ 上 Lebesgue 测度奇异的(有界)Borel 正测度 $\nu$,
\begin{equation*}
\mathcal{S}(w) = \exp\left\{i\int_{-\infty}^{\infty}\frac{1 + tw}{t - w}\dd \nu(t)\right\}
\end{equation*}
称为奇异内函数.已知 $\mathcal{D}$ 上的任意内函数(即属于 $\mathcal{H}^{\infty}$,其边界值几乎处处绝对值等于 1 的函数)存在常数 $\alpha \geq 0$,可写成
\begin{equation*}
e^{i\alpha w}\mathcal{B}(w)\mathcal{S}(w).
\end{equation*}
若存在实常数 $\gamma$,$G \in \mathcal{H}^{p}$ 可写成
\begin{equation*}
G(w) = e^{i\gamma}\exp\left\{\frac{1}{\pi i}\int_{-\infty}^{\infty}\frac{1 + tw}{t - w}\frac{\log|G(t)|}{1 + t^{2}}\dd t \right\},
\end{equation*}
则称为 $\mathcal{D}$ 上的外函数.已知任意 $\mathcal{H}^{p}$ 的函数可写成内函数与外函数的乘积.


% ============ 第2章 参考文献(原书 pp.200-202) ============
\section*{参考文献}\addcontentsline{toc}{section}{参考文献}
\begin{thebibliography}{99}
\item[] 注意．本章主要参考 Duren 的专著 \cite{ch2:11} 写成.另外 Hoffman 的书 \cite{ch2:22} 与 Koosis 的专著 \cite{ch2:23} 也被部分参考.
\bibitem[1]{ch2:1} J. Neuwirth and D. J. Newman, Positive $H^{1/2}$ functions are constant, Proc. Amer. Math. Soc. 18(1967), 958.
\bibitem[2]{ch2:2} H. Helson and D. Sarason, Past and future, Math. Scand. 21(1967), 5--16.
\bibitem[3]{ch2:3} H. L. Royden, The boundary values of analytic and harmonic functions, Math. Zeitschr. 78(1962), 1--24.
\bibitem[4]{ch2:4} 竹之内脩, 函数解析, 近代数学講座 13, 朝倉書店, 東京, 1967.
\bibitem[5]{ch2:5} K. deLeeuw and W. Rudin, Extreme points and extremum problems in $H^{1}$, Pacific J. Math. 8(1958), 467--485.
\bibitem[6]{ch2:6} K. deLeeuw, W. Rudin and J. Wermer, The isometries of some function spaces, Proc. Amer. Math. Soc. 11(1960), 694--698.
\bibitem[7]{ch2:7} K. Yabuta, Some uniqueness theorems for $H^{p}(U^{n})$ functions, Tôhoku Math. J. 24(1972), 353--357.
\bibitem[8]{ch2:8} T. Nakazi, Exposed points and extremal problems in $H^{1}$, J. Funct. Anal. 53(1983), 224--230.
\bibitem[9]{ch2:9} D. Sarason, Exposed points in $H^{1}$, I, in Operator Theory: Advances and Applications, 41(1989), 485--496.
\bibitem[10]{ch2:10} J. Inoue, An example of a non-exposed extreme function on the unit ball of $H^{1}$, Proc. Edinburgh Math. Soc. 37(1993), 47--51.
\bibitem[11]{ch2:11} P. L. K. Duren, Theory of $H^{p}$ Spaces, Academic Press, New York, 1970.
\bibitem[12]{ch2:12} T. Nakazi, Extremal problems in $H^{p}$, J. Austral. Math. Soc. 52(1992), 103--110.
\bibitem[13]{ch2:13} E. Hayashi, The solution sets of extremal problems in $H^{1}$, Proc. Amer. Math. Soc. 93(1985), 690--696.
\bibitem[14]{ch2:14} G. Pick, Über die Beschränkungen analytische Funktionen, welche durch vorgegebene Funktionswerte bewirkt werden, Math. Ann. 77(1916), 7--23.
\bibitem[15]{ch2:15} D. Sarason, Generalized interpolation in $H^{\infty}$, Trans. Amer. Math. Soc. 127(1967), 179--203.
\bibitem[16]{ch2:16} L. Carleson, An interpolations problem by bounded analytic functions, Amer. J. Math. 80(1958), 921--930.
\bibitem[17]{ch2:17} H. S. Shapiro and A. Shields, On some interpolation problems for analytic functions, Amer. J. Math. 83(1961), 513--532.
\bibitem[18]{ch2:18} V. Kabaila, Interpolation sequences for the $H_{p}$ classes in the case p < 1, Litovsk. Mat. Sb. 3(1963), 141--147 (in Russian).
\bibitem[19]{ch2:19} W. K. Hayman, Interpolation by bounded functions, Ann. Inst. Fourier (Grenoble) 8(1959), 277--290.
\bibitem[20]{ch2:20} D. J. Newman, Interpolation in $H^{\infty}$, Trans. Amer. Math. Soc. 92(1959), 501--507.
\bibitem[21]{ch2:21} J. Garnett, Harmonic interpolating sequences, $L^{p}$ and BMO, Ann. Inst. Fourier (Grenoble) 28, 215--228.
\bibitem[22]{ch2:22} K. Hoffman, Banach Spaces of Analytic Functions, Prentice-Hall, Englewood Cliffs, New Jersey, 1962.
\bibitem[23]{ch2:23} P. Koosis, Introduction to $H_{p}$ Spaces, London Math. Society Lecture Notes Series, 40. Cambridge Univ. Press, London and New York, 1980.
\end{thebibliography}

% ============ 第3章 Fourier分析中的交换Banach代数（佐藤圓治，原书 pp.203-309 = PDF pp.215-321） ============
\chapter{Fourier 分析中的交换 Banach 代数（佐藤圓治）}\label{ch:3}

% ============ 第3章 §3.1 交换Banach代数(原书 pp.203-226) ============
\section{交换 Banach 代数}\label{sec:3.1}
\index{Banach 代数}
\index{Banach 代数!交换----}
\index{Banach 代数!添加单位元的交换----}
\index{交换 Banach 代数}
\index{谱}
\index{预解式}
\index{谱半径}
\index{理想}
\index{理想!极大----}
\index{理想!正则----}
\index{理想!闭----}
\index{正则!(理想)}
\index{Gelfand 变换}
\index{Gelfand 表示}
\index{极大理想空间}
\index{范数拓扑}
\index{*-弱拓扑}
\index{hk 拓扑}
\index{谱综合}
\index{谱综合集}
\index{谱综合!谱综合可能}
\index{谱综合!谱综合不可能}
\index{半单}
\index{闭图像定理}
\index{特征}
\index{Ditkin 条件}
\index{作用函数!(半单交换 Banach 代的)}


本节的目的在于叙述交换 Banach 代数的一般理论.Fourier 分析中常出现的交换 Banach 代数,其极大理想空间的结构往往复杂,但这里的例子处理的是简单的情形.下一节将处理 Fourier 分析中出现的主要交换 Banach 代数.另外,函数分析中众所周知的基础知识视为已知.

\subsection{交换 Banach 代数的定义}\label{sec:3.1.1}

\begin{definition}\label{def:3.1.1}
集合 $A$ 具有如下性质.

(1) $A$ 是复数域 $\mathbf{C}$ 上的环 $(+, \cdot)$.即,

(i) $A$ 是 $\mathbf{C}$ 上的向量空间.

(ii) $x(yz) = (xy)z$,$x(y + z) = xy + xz$,$\alpha(xy) = (\alpha x)y = x(\alpha y)$($x, y, z \in A$,$\alpha \in \mathbf{C}$).

(2) 对 $x \in A$,赋予非负值 $\|x\|$ 使其满足下列条件(把满足下述 (i), (ii), (iii) 的 $\|\cdot\|$ 称为范数).

(i) $\|x\| \geq 0$,且 $\|x\| = 0$ 的情形仅限于 $x = 0$.

(ii) $\|\alpha x\| = |\alpha|\,\|x\|$.

(iii) $\|x + y\| \leq \|x\| + \|y\|$.

(vi) $\|xy\| \leq \|x\| \cdot \|y\|$.

(3) $A$ 关于由范数导出的距离 $d(x, y) = \|x - y\|$ 完备.

此时,称 $A$ 为 Banach 代数.特别地,当对 $x, y \in A$ 有 $xy = yx$ 时,称 $A$ 为交换 Banach 代数.又,当存在满足 $ex = xe = x$ 的 $e \in A$ 时,称 $e$ 为 $A$ 的单位元.以后,$A$ 中的拓扑均按由范数 $\|\cdot\|$ 导出的距离 $d(x, y)$ 来考虑.
\end{definition}

在交换 Banach 代数中,若存在单位元 $e$,则设 $\|e\| = 1$.若不存在单位元,令
\begin{equation*}
A_{e} = \{(x, \lambda):\ x \in A,\ \lambda \in \mathbf{C}\}.
\end{equation*}
此时,定义
和 $(x, \lambda) + (y, \mu) = (x + y, \lambda + \mu)$,
数乘 $\alpha(x, \lambda) = (\alpha x, \alpha\lambda)$,
积 $(x, \lambda) \cdot (y, \mu) = (xy + \mu x + \lambda y, \lambda\mu)$,
\begin{align*}
\text{范数} & \|(x, \lambda)\| = \|x\| + |\lambda|, \\
& (x, y \in A,\ \lambda, \mu, \alpha \in \mathbf{C})
\end{align*}
则 $A_{e}$ 成为以 $e = (0, 1)$ 为单位元的交换 Banach 代数.此时,在 $A_{e}$ 中,$\{(x, 0):\ x \in A\}$ 作为交换 Banach 代数可与 $A$ 等同.称这个 $A_{e}$ 为对 $A$ 添加单位元而得到的交换 Banach 代数.

例．(1) 若 $\mathbf{T} = \{|z| = 1:\ z \in \mathbf{C}\}$,则 $\mathbf{T}$ 是紧集,且关于乘法构成交换群.又,若以 $C(\mathbf{T})$ 表示 $\mathbf{T}$ 上复值连续函数的全体,则按通常的函数运算与范数 $\|f\|_{\infty} = \sup\{|f(z)|:\ z \in \mathbf{T}\}$,$C(\mathbf{T})$ 成为交换 Banach 代数.

(2) 令 $\ell^{1} = \{a = \{a_{n}\}:\ \sum_{n \in \mathbf{Z}} |a_{n}| < \infty,\ a_{n} \in \mathbf{C}\ (n \in \mathbf{Z})\}$.$\mathbf{Z}$ 表示整数集.又,对 $a = \{a_{n}\}$,$b = \{b_{n}\} \in \ell^{1}$,$\alpha, \beta \in \mathbf{C}$,有 $\alpha a + \beta b \in \ell^{1}$.进一步,令 $\|a\|_{\ell^{1}} = \sum_{n \in \mathbf{Z}} |a_{n}| < \infty$,$a * b = \{\sum_{k \in \mathbf{Z}} a_{n-k} b_{k}\}$,则 $\ell^{1}$ 关于范数 $\|\cdot\|_{\ell^{1}}$ 与运算 $*$ 成为交换 Banach 代数.

此后,对紧 Hausdorff 空间 $X$,用 $C(X)$ 表示以 $\|f\|_{\infty} = \sup_{x \in X} |f(x)|$ 为范数的 Banach 空间.

\subsection{逆元,预解式与谱}\label{sec:3.1.2}

\begin{definition}\label{def:3.1.2}
设 $A$ 为具有单位元 $e$ 的交换 Banach 代数,令
\begin{equation*}
A^{-1} = \{x \in A:\ \text{存在 } x \text{ 的逆元 } x^{-1} \text{ 使得 } xx^{-1} = e\}.
\end{equation*}
\end{definition}

\begin{theorem}\label{thm:3.1.3}
设 $A$ 为具有单位元 $e$ 的交换 Banach 代数.

(1) $A^{-1} \supset \{x:\ \|e - x\| < 1\}$,且 $A^{-1}$ 是开集.

(2) 由 $x$ 到 $x^{-1}$ 的对应是从 $A^{-1}$ 到 $A^{-1}$ 的连续映射.
\end{theorem}

\begin{proof}
(1) 当 $\|x\| < 1$ 时,令 $x^{0} = e$,则 $\sum_{n=0}^{\infty} x^{n}$ 在 $A$ 中绝对收敛.因此 $\sum_{n=0}^{\infty} x^{n}$ 收敛.此时
\begin{equation*}
(e - x)\sum_{n=0}^{\infty} x^{n} = \sum_{n=0}^{\infty} x^{n} - \sum_{n=1}^{\infty} x^{n} = e
\end{equation*}
故 $\{e - x:\ \|x\| < 1\} \subset A^{-1}$.于是 $\{x:\ \|e - x\| < 1\} \subset A^{-1}$.

进一步,当 $x \in A^{-1}$ 时,若 $\|y\| < \|x^{-1}\|^{-1}$,则
\begin{equation*}
\|x^{-1} y\| \leq \|x^{-1}\| \cdot \|y\| < 1
\end{equation*}
故由前半可知 $e - (-x^{-1}y) \in A^{-1}$.从而若 $\|y\| < \|x^{-1}\|^{-1}$,则 $x + y \in A^{-1}$,可见 $A^{-1}$ 是开集.

(2) 利用 (1),
\begin{align*}
(x + y)^{-1} - x^{-1} & = x^{-1}\bigl((e + x^{-1}y)^{-1} - e\bigr) \\
& = x^{-1}\Bigl(\sum_{n=0}^{\infty} (-x^{-1}y)^{n} - e\Bigr) \\
& = x^{-2}y\sum_{n=0}^{\infty} (-1)^{n+1}(x^{-1}y)^{n},
\end{align*}
故
\begin{align*}
\|(x + y)^{-1} - x^{-1}\| & \leq \|x^{-2}y\|\sum_{n=0}^{\infty}\|x^{-1}y\|^{n} \\
& \leq \|x^{-2}\| \cdot \|y\|\sum_{n=0}^{\infty}\bigl(\|x^{-1}\| \cdot \|y\|\bigr)^{n} \\
& = \frac{\|x^{-2}\| \cdot \|y\|}{1 - \|x^{-1}\| \cdot \|y\|} \to 0 \quad (\|y\| \to 0)
\end{align*}
由此可知,当 $\|y\| \to 0$ 时,$\|(x + y)^{-1} - x^{-1}\| \to 0$.
\end{proof}

\begin{definition}\label{def:3.1.4}

(1) 设 $A$ 为具有单位元 $e$ 的交换 Banach 代数.

(i) 对 $x \in A$,
\begin{equation*}
\Sp_{A}(x) = \{\lambda \in \mathbf{C}:\ x - \lambda e \notin A^{-1}\}
\end{equation*}
称为 $x$ 的(在 $A$ 中的)谱.$\Sp_{A}(x)$ 也记作 $\Sp(x, A)$.

(ii) 对满足 $\lambda \notin \Sp_{A}(x)$ 的 $\lambda \in \mathbf{C}$,
\begin{equation*}
R(x, \lambda) = (x - \lambda e)^{-1}
\end{equation*}
称为 $x$ 的预解式.

(2) 当 $A$ 为没有单位元的交换 Banach 代数时,令 $\Sp_{A}(x) = \Sp_{A_{e}}(x)$.此时容易看出,对 $\lambda \neq 0$,$\lambda \notin \Sp_{A}(x)$ 的充要条件是存在 $y \in A$ 使得 $xy - \lambda^{-1}x - \lambda y = 0$.
\end{definition}

\begin{definition}\label{def:3.1.5}
设 $\mathbf{C} \supset D$ 为开集,对交换 Banach 代数 $A$,若对任意 $\lambda_{0} \in D$,
\begin{equation*}
\lim_{\lambda \to \lambda_{0}} \frac{F(\lambda) - F(\lambda_{0})}{\lambda - \lambda_{0}}
\end{equation*}
在 $A$ 中存在,则称 $F(\lambda):\ D \to A$ 在 $D$ 中正则.
\end{definition}

\begin{theorem}\label{thm:3.1.6}
设 $A$ 为交换 Banach 代数.此时,对 $x \in A$,下述 (i)～(iii) 成立.

(i) $\Sp_{A}(x)$ 是闭集.

(ii) $\Sp_{A}(x) \subset \{\lambda:\ |\lambda| \leq \|x\|\}$.

(iii) $\Sp_{A}(x) \neq \emptyset$.

特别地,由 (i), (ii) 可知 $\Sp_{A}(x)$ 是复平面上的紧集.
\end{theorem}

\begin{proof}
首先设 $A$ 具有单位元 $e$.(i) 令 $F(\lambda) = x - \lambda e$,则 $F(\lambda):\ \mathbf{C} \to A$ 是连续映射.事实上,由
\begin{equation*}
\|F(\lambda) - F(\lambda_{0})\| = \|(x - \lambda e) - (x - \lambda_{0}e)\| = |\lambda - \lambda_{0}|\,\|e\|
\end{equation*}
可知.这里,由定理~\ref{thm:3.1.3},$A^{-1}$ 是开集,故 $(A^{-1})^{c}$ 是闭集.于是因 $F$ 连续,$F^{-1}((A^{-1})^{c})$ 是复平面的闭子集.从而 $\Sp_{A}(x) = F^{-1}((A^{-1})^{c})$ 是闭集.

(ii) 若 $|\lambda| > \|x\|$,则由定理~\ref{thm:3.1.3},$x - \lambda e = -\lambda(e - \lambda^{-1}x) \in A^{-1}$,故 $\lambda \notin \Sp_{A}(x)$.因此 $\Sp_{A}(x) \subset \{\lambda:\ |\lambda| \leq \|x\|\}$.从而 $\Sp_{A}(x)$ 是 $\mathbf{C}$ 的紧子集.

(iii) 证明 $\Sp_{A}(x) \neq \emptyset$.当 $\lambda_{0} \notin \Sp_{A}(x)$ 时,$G(\lambda) = R(x, \lambda)$ 在 $\lambda = \lambda_{0}$ 处正则,且当 $|\lambda| > \|x\|$ 时,
\begin{equation*}
\lim_{\lambda \to \lambda_{0}} \frac{R(x, \lambda) - R(x, \lambda_{0})}{\lambda - \lambda_{0}} = \bigl(R(x, \lambda_{0})\bigr)^{2},
\end{equation*}
\begin{equation*}
R(x, \lambda) = -\frac{1}{\lambda}\sum_{n=0}^{\infty}\left(\frac{x}{\lambda}\right)^{n}.
\end{equation*}
事实上,由 $R(x, \lambda) - R(x, \lambda_{0}) = (\lambda - \lambda_{0})R(x, \lambda)R(x, \lambda_{0})$,以及当 $|\lambda| > \|x\|$ 时
\begin{equation*}
(x - \lambda e)^{-1} = \left(-\frac{1}{\lambda}\right)\left(e - \frac{x}{\lambda}\right)^{-1} = \left(-\frac{1}{\lambda}\right)\sum_{n=0}^{\infty}\left(\frac{x}{\lambda}\right)^{n}
\end{equation*}
即得.于是,考虑 $G(\lambda)$,若 $\Sp_{A}(x) = \emptyset$,则作为向量值函数 $G$ 在 $\mathbf{C}$ 上正则且有界.从而由向量值函数论(cf.\ \cite{ch3:3}),存在 $y \in A$ 使得 $G(\lambda) = y$($\lambda \in \mathbf{C}$).于是,$\lambda = 0$ 时 $y = G(0) = x^{-1}$,$\lambda = 1$ 时 $y = G(1) = (x - e)^{-1}$.这与 $y = G(\lambda) \in A^{-1}$ 相矛盾.因此 $\Sp_{A}(x) \neq \emptyset$.

当 $A$ 没有单位元时,由 $\Sp_{A}(x) = \Sp_{A_{e}}(x)$,把 $x$ 作为 $A_{e}$ 的元来考虑,由前半的论证可知 $\Sp_{A}(x)$ 是 $\mathbf{C}$ 的非空紧子集.
\end{proof}

\begin{corollary}\label{cor:3.1.7}
具有单位元 $e$ 的交换 Banach 代数 $A$ 若为域,则(作为环)与复数域同构.
\end{corollary}

\begin{proof}
设 $A$ 为域.若 $x \in A$,则由定理~\ref{thm:3.1.6},$\Sp_{A}(x) \neq \emptyset$.因此存在 $\lambda \in \mathbf{C}$ 使得 $x - \lambda e \notin A^{-1}$.此时因 $A$ 是域,$x - \lambda e = 0$.最终得到 $x = \lambda e$,由这一对应 $x \mapsto \lambda$,$A$ 与 $\mathbf{C}$ 作为环同构.
\end{proof}

\begin{definition}\label{def:3.1.8}
对 $x \in A$,定义
\begin{equation*}
\rho(x) = \sup\{|\lambda|:\ \lambda \in \Sp_{A}(x)\},
\end{equation*}
称 $\rho(x)$ 为 $x$ 的谱半径.
\end{definition}

\begin{theorem}\label{thm:3.1.9}
下式成立:
\begin{equation*}
\rho(x) = \lim_{n \to \infty}\|x^{n}\|^{1/n}.
\end{equation*}
\end{theorem}

\begin{proof}
设 $A$ 具有单位元 $e$.若 $\lambda \in \Sp_{A}(x)$,则 $\lambda^{n} \in \Sp_{A}(x^{n})$.由定理~\ref{thm:3.1.6},$|\lambda^{n}| \leq \|x^{n}\|$,故 $|\lambda| \leq \varliminf_{n \to \infty}\|x^{n}\|^{1/n}$.因此
\begin{equation*}
\rho(x) \leq \varliminf_{n \to \infty}\|x^{n}\|^{1/n}.
\end{equation*}
现在,考虑
\begin{equation*}
F(\lambda) = \left(-\frac{1}{\lambda}\right)\sum_{n=0}^{\infty}\left(\frac{x}{\lambda}\right)^{n},
\end{equation*}
当 $|\lambda| > \varlimsup_{n \to \infty}\sqrt[n]{\|x^{n}\|}$ 时,$F(\lambda)$ 的右边在 $A$ 中收敛,且 $F(\lambda) \in A$(cf.\ \cite{ch3:3}, \cite{ch3:8}).又,当 $\|\lambda^{-1}x\| < 1$ 时,由定理~\ref{thm:3.1.6},
\begin{equation*}
(x - \lambda e)^{-1} = F(\lambda),
\end{equation*}
故注意 $(x - \lambda e)F(1/\lambda) = e$.进一步,若 $|\lambda| > \rho(x)$,则由 $\rho(x)$ 的定义 $\lambda \notin \Sp_{A}(x)$,故存在正则的 $R(x, \lambda) = (x - \lambda e)^{-1}$.综上,
\begin{equation*}
\frac{1}{\rho(x)} \leq \frac{1}{\varlimsup_{n \to \infty}\sqrt[n]{\|x^{n}\|}}.
\end{equation*}
于是
\begin{equation*}
\varlimsup_{n \to \infty}\|x^{n}\|^{1/n} \leq \rho(x)
\end{equation*}
与前半合并即得
\begin{equation*}
\rho(x) = \lim_{n \to \infty}\|x^{n}\|^{1/n}.
\end{equation*}
当 $A$ 没有单位元时,把 $x \in A$ 放到 $A_{e}$ 中考虑,由 $\Sp_{A}(x) = \Sp_{A_{e}}(x)$ 归结为具有单位元的情形.
\end{proof}

\subsection{理想}\label{sec:3.1.3}

以下假设 $A$ 为交换 Banach 代数.

\begin{definition}\label{def:3.1.10}

(1) 若 $A$ 的子集 $I$ 是 $\mathbf{C}$ 上的环,且当 $x \in I$,$y \in A$ 时 $xy \in I$ 成立,则称 $I$ 为 $A$ 的理想.

(2) 若 $I$ 是理想,且不存在满足 $I \subsetneq J \subsetneq A$ 的理想 $J$,则称 $I$ 为极大理想.
\end{definition}

下面容易证明.

\begin{proposition}\label{prop:3.1.11}
当 $I$ 为 $A$ 的理想时,$I$ 的等价类 $A/I = \{x + I:\ x \in A\}$ 通过
(i) $(x + I) + (y + I) = x + y + I$,
(ii) $\lambda(x + I) = \lambda x + I$,
(iii) $(x + I)(y + I) = xy + I$
成为 $\mathbf{C}$ 上的环.$A/I$ 称为商环.特别地,当 $I$ 为闭集时,$A/I$ 由商拓扑 $\|x + I\|_{A/I} = \inf_{z \in I}\|x + z\|$ 成为交换 Banach 代数.
\end{proposition}

\begin{definition}\label{def:3.1.12}
$A$ 的理想 $I$ 称为正则理想,是指商环 $A/I$ 具有单位元.此时,$A/I$ 具有单位元的充要条件是存在 $u \in A$ 使得对任意 $x \in A$ 有 $ux - x \in I$.称这个 $u$ 为 $A$ 模 $I$ 的单位元.
\end{definition}

\begin{lemma}\label{lem:3.1.13}
当 $I \subsetneq A$ 且 $I$ 为正则理想时,存在正则极大理想 $J$ 满足 $I \subset J \subsetneq A$.
\end{lemma}

\begin{proof}
用 Zorn 引理证明.设 $I$ 为与 $A$ 不同的正则理想,则 $A$ 模 $I$ 的单位元 $u$ 显然不属于 $I$.事实上,若 $u \in I$,则由 $ux - x \in I$($x \in A$)得 $I = A$.现在,令 $\mathcal{J}$ 为满足 $I \subset J$,$u \notin J$ 的 $A$ 的理想 $J$ 的全体,则 $\mathcal{J}$ 关于包含关系构成半序集.取 $\mathcal{J}$ 的任意全序子集 $\mathcal{K}$,则 $L = \bigcup_{J \in \mathcal{K}} J$ 是包含 $I$ 的 $A$ 的理想,且 $u \notin L$.因此 $L \in \mathcal{J}$.从而由 Zorn 引理,$\mathcal{J}$ 中存在极大元 $J$.即满足 $I \subset J \subsetneq A$ 的理想.此时,对任意 $x \in A$,$ux - x \in I \subset J$,故 $J$ 是包含 $I$ 的正则极大理想.
\end{proof}

\begin{proposition}\label{prop:3.1.14}

(i) 若 $A$ 具有单位元 $e$,则 $A$ 的理想 $I$ 是正则理想,且 $A$ 的极大理想是闭集(是闭集的理想称为闭理想).

(ii) 设 $A$ 没有单位元.若 $I \subset A$ 是正则极大理想,则 $I$ 是闭理想.
\end{proposition}

\begin{proof}
(i) 当 $x + I \in A/I$ 时,由 $(e + I)(x + I) = x + I$,$I$ 是正则理想.设 $I$ 为极大理想.若 $\|e - x\| < 1$,则由定理~\ref{thm:3.1.3},$x \notin I$.因此若 $x \in I$,则 $\|e - x\| \geq 1$.于是用 $\overline{I}$ 表示 $I$ 的闭包,任取 $y \in \overline{I}$,则 $\|e - y\| \geq 1$.从而 $\overline{I} \subsetneqq A$,且 $\overline{I}$ 是理想.又因 $I$ 是极大理想,$I = \overline{I}$.因此是闭集.

(ii) 设 $u$ 为 $A$ 模 $I$ 的单位元.此时令 $\|u - I\| = \inf_{x \in I}\|u - x\|$,则 $\|u - I\| \geq 1$.事实上,假设存在 $v \in I$ 使得 $\|u - v\| < 1$.固定任意 $x \in A$,令
\begin{equation*}
y = x + \Bigl(\sum_{n=1}^{\infty}(u - v)^{n}\Bigr)x,
\end{equation*}
则
\begin{align*}
y - uy + vy & = x + \Bigl(\sum_{n=1}^{\infty}(u - v)^{n}\Bigr)x \\
& \quad - u\Bigl(x + \Bigl(\sum_{n=1}^{\infty}(u - v)^{n}\Bigr)x\Bigr) + v\Bigl(x + \Bigl(\sum_{n=1}^{\infty}(u - v)^{n}\Bigr)x\Bigr) \\
& = x + \Bigl(\sum_{n=1}^{\infty}(u - v)^{n}\Bigr)x - (u - v)x - \Bigl(\sum_{n=2}^{\infty}(u - v)^{n}\Bigr)x \\
& = x
\end{align*}
即,由 $v \in I$ 与 $u$ 的定义,$x \in I$.于是 $A = I$,矛盾.因此 $\|u - I\| \geq 1$ 成立.$\overline{I}$ 是 $I$ 的闭包,由 $I \subset \overline{I}$,$\overline{I}$ 是正则理想.因 $\|u - x\| \geq 1$($x \in I$),故 $\|u - x\| \geq 1$($x \in \overline{I}$).于是 $\overline{I} \subsetneqq A$.从而由 $I$ 是极大理想,$I = \overline{I}$.
\end{proof}

\begin{lemma}\label{lem:3.1.15}
当 $I$ 为与 $A$ 不同的正则理想时,$I$ 是正则极大理想的充要条件是 $A/I$ 与 $\mathbf{C}$ 同构.
\end{lemma}

\begin{proof}
设 $A/I$ 与 $\mathbf{C}$ 同构.若 $I$ 不是极大理想,则由补题~\ref{lem:3.1.13},存在极大理想 $J$ 使得 $I \subsetneq J$.取 $w \in J \cap I^{c}$,则与 $w + I$ 对应的是非零的 $\lambda \in \mathbf{C}$.因此 $w + I$ 在 $A/I$ 中具有逆元.于是 $J = A$,产生矛盾,故 $I$ 是极大理想.

反之,设 $I$ 为正则极大理想.由命题~\ref{prop:3.1.11},命题~\ref{prop:3.1.14},$A/I$ 成为具有单位元 $u + I$ 的交换 Banach 代数.此时 $A/I$ 是域.事实上,设 $x_{0} + I \in A/I$ 且 $x_{0} \notin I$,则由 $I$ 是极大理想,$\{x_{0}x + I:\ x \in A\} = A/I$.因为 $\{xx_{0} + y:\ x \in A,\ y \in I\}$ 是包含 $I$ 的理想.因此存在 $z \in A$ 使得 $zx_{0} + I = (z + I)(x_{0} + I) = u + I$,故 $x_{0} + I$ 具有逆元,$A/I$ 是域.从而由系~\ref{cor:3.1.7},$A/I$ 与 $\mathbf{C}$ 同构.
\end{proof}

\begin{theorem}\label{thm:3.1.16}
设 $A$ 为没有单位元的交换 Banach 代数,$A_{e}$ 为对 $A$ 添加单位元 $e$ 而得到的交换 Banach 代数.进一步,令 $M_{1}$ 为 $A_{e}$ 的正则极大理想的全体,$M_{0}$ 为 $A$ 的正则极大理想的全体,则 $M_{1} \setminus \{A\} \ni I \mapsto I \cap A \in M_{0}$ 是双射.
\end{theorem}

\begin{proof}
首先证明若 $I \in \mathfrak{M}_{1} \setminus \{A\}$,则 $I \cap A \in \mathfrak{M}_{0}$.首先,$I \cap A$ 是 $A$ 的理想,由假设容易看出 $I \not\subset A$.进一步,存在 $u \in A$ 使得 $u - e \in I$.事实上,因 $I \not\subset A$,存在非零复数 $\lambda_{0}$ 与 $y \in A$ 使得 $y + \lambda_{0}e \in I$.即 $y/\lambda_{0} + e \in I$.令 $-u = y/\lambda_{0}$,则 $e - u \in I$.由此得 $ux - x \in I$($x \in A$).因此 $ux - x \in I \cap A$($x \in A$),故 $I \cap A$ 是 $A$ 的正则理想.进一步,$I \cap A$ 是 $A$ 的极大理想.因为,若存在 $A$ 的理想 $J$ 满足 $A \cap I \subsetneq J \subset A$,则存在 $v_{0} \in J$ 使得 $v_{0} \notin I \cap A$.此时,$v_{0} + I \in A_{e}/I$ 由补题~\ref{lem:3.1.15} 对应于适当的非零复数 $\lambda' \in \mathbf{C}$,故 $v_{0}/\lambda' - e \in I$.于是由 $e - u \in I$,$v_{0}/\lambda' - e \in I$ 得 $v_{0}/\lambda' - u \in I \cap A \subset J$.再考虑 $v_{0} \in J$,得 $u \in J$.从而 $J = A$.综上得 $I \cap A \in \mathfrak{M}_{0}$.

其次证明上述映射是单射.事实上,设 $I_{1}, I_{2} \in \mathfrak{M}_{1} \setminus \{A\}$,$I_{1} \cap A = I_{2} \cap A$.此时,由前述,存在 $u \in A$ 使得 $ux - x \in I_{1} \cap A$($x \in A$).于是令
\begin{equation*}
J = \{x \in A_{e}:\ xu \in I_{1} \cap A\},
\end{equation*}
则 $J$ 是 $A_{e}$ 的理想且 $J \supset I_{1}$.这里,因 $I_{1} \cap A \subsetneq A$,$e \notin J$,得 $J \subsetneq A_{e}$.从而由 $I_{1}$ 是极大理想,$J = I_{1}$.同理 $J = I_{2}$.因此 $I_{1} = I_{2}$.

进一步证明该映射是满射.事实上,固定任意 $J \in M_{0}$,取适当的 $u \in A$ 使 $ux - x \in J$($x \in A$),令
\begin{equation*}
I' = \{x \in A_{e}:\ ux \in J\},
\end{equation*}
则 $I'$ 是 $A_{e}$ 的理想,且 $u \notin I'$.事实上,若 $u \in I'$,则 $u^{2} \in J$,由 $u^{2} - u \in J$ 得 $u \in J$,矛盾.于是取满足 $I \supset I'$ 且 $u \notin I$ 的 $A_{e}$ 的极大理想 $I$,由 $u \notin I$ 得 $I \neq A$,故 $I \in \mathfrak{M}_{1} \setminus \{A\}$.此时,若 $x \in J$,则 $ux \in J$,故 $x \in I'$.因此 $J \subset I$,即 $J \subset I \cap A$.因 $J$ 是极大理想,得 $J = I \cap A$.
\end{proof}

例．对任意 $x_{0} \in \mathbf{T}$,
\begin{equation*}
I(x_{0}) = \{f \in C(\mathbf{T}):\ f(x_{0}) = 0\}
\end{equation*}
是交换 Banach 代数 $C(\mathbf{T})$ 的极大理想.

\subsection{Gelfand 变换}\label{sec:3.1.4}

\begin{definition}\label{def:3.1.17}
取交换 Banach 代数 $A$,若 $\chi:\ A \to \mathbf{C}$ 具有如下性质,则称 $\chi$ 为特征.

(i) $\chi(\alpha x + \beta y) = \alpha\chi(x) + \beta\chi(y)$,\quad (ii) $\chi(xy) = \chi(x)\chi(y)$ \quad($x, y \in A$,$\alpha, \beta \in \mathbf{C}$)

用 $\Delta(A)$ 表示 $A$ 的满足 $\chi \neq 0$ 的特征 $\chi$ 的全体.
\end{definition}

\begin{theorem}\label{thm:3.1.18}
$\Delta(A)$ 与 $A$ 的正则极大理想全体 $\mathfrak{M}(A)$ 通过对应 $\Delta(A) \ni \chi \mapsto \chi^{-1}(0) \in \mathfrak{M}(A)$ 构成双射.
\end{theorem}

\begin{proof}
首先证明上述对应是满射.事实上,设 $I \in \mathfrak{M}(A)$,则由补题~\ref{lem:3.1.15},$A/I$ 与 $\mathbf{C}$ 同构,此时对应
\begin{equation*}
A/I \ni x + I \mapsto \chi(x + I) \in \mathbf{C}
\end{equation*}
是同构.进一步,令 $\chi(x) = \chi(x + I)$,则 $\chi$ 是特征.事实上,可定义性由若 $x - x' \in I$ 则按 $\chi$ 的定义 $\chi(x) = \chi(x + I) = \chi(x' + I)$ 可知.此时 $\chi^{-1}(0) = I$.反之,对 $\Delta(A) \ni \chi$,令 $I = \chi^{-1}(0)$,因 $\chi$ 是同态,$A/I$ 与 $\mathbf{C}$ 同构.

其次证明单射.事实上,若 $I = \chi^{-1}(0) = (\chi')^{-1}(0) = I'$,则 $\chi(x + I) = \chi'(x + I)$,故 $\chi = \chi'$.
\end{proof}

\begin{corollary}\label{cor:3.1.19}
设 $A$ 为交换 Banach 代数,则对 $x \in A$ 下式成立.

(1) 当 $A$ 具有单位元时,$\Sp_{A}(x) = \{\chi(x):\ \chi \in \Delta(A)\}$.

(2) 当 $A$ 没有单位元时,$\Sp_{A}(x) = \{\chi(x):\ \chi \in \Delta(A)\} \cup \{0\}$.
\end{corollary}

\begin{proof}
(1) 设 $A$ 的单位元为 $e$.当 $\lambda \in \Sp_{A}(x)$ 时,$x - \lambda e \notin A^{-1}$.此时存在 $I \in \mathfrak{M}(A)$ 使得 $x - \lambda e \in I$.从而由定理~\ref{thm:3.1.18},存在 $\chi \in \Delta(A)$ 使 $\chi^{-1}(0) = I$.即 $\chi(x - \lambda e) = 0$.于是 $\chi(x) = \lambda$.因此 $\lambda \in \{\chi(x):\ \chi \in \Delta(A)\}$.反向是显然的.

(2) 设 $A_{e}$ 为对 $A$ 添加单位元 $e$ 而得到的交换 Banach 代数,由 $\Sp_{A}(x) = \Sp_{A_{e}}(x)$,
\begin{equation*}
\Sp_{A_{e}}(x) = \{\chi(x):\ \chi \in \Delta(A_{e})\} = \{\chi(x):\ \chi \in \Delta(A)\} \cup \{0\}
\end{equation*}
即得.
\end{proof}

\begin{corollary}\label{cor:3.1.20}
当 $\chi \in \Delta(A)$ 时,$|\chi(x)| \leq \|x\|$($x \in A$).特别地,$\chi$ 是 $A$ 上的复值连续函数.
\end{corollary}

\begin{proof}
由系~\ref{cor:3.1.19},$\Sp_{A}(x) = \{\chi(x):\ \chi \in \Delta(A)\}$,或 $\Sp_{A}(x) = \{\chi(x):\ \chi \in \Delta(A)\} \cup \{0\}$.由于 $\rho(x) = \sup\{|\lambda|:\ \lambda \in \Sp_{A}(x)\}$,故 $|\chi(x)| \leq \|x\|$.特别地,当 $\|x_{n} - x\| \to 0$ 时,
\begin{equation*}
|\chi(x_{n}) - \chi(x)| = |\chi(x_{n} - x)| \leq \|x_{n} - x\| \to 0
\end{equation*}
故 $\chi:\ A \to \mathbf{C}$ 连续.
\end{proof}

\begin{definition}\label{def:3.1.21}
对 Banach 空间 $A$,把 $A$ 上的连续线性复值函数(称为 $A$ 上的连续线性泛函)的全体定义为 $A^{*}$.此时,$u \in A^{*}$ 在 $A$ 上连续的充要条件是存在常数 $C > 0$ 使得 $|u(x)| \leq C\|x\|$($x \in A$).$A^{*}$ 由通常的运算成为向量空间.
\end{definition}

\begin{definition}\label{def:3.1.22}

(1) 令 $u \in A^{*}$,$\|u\| = \sup\{|u(x)|:\ x \in A,\ \|x\| = 1\}$,则 $A^{*}$ 关于 $\|\cdot\|$ 成为 Banach 空间.称这个拓扑为 $A^{*}$ 的范数拓扑.

(2) 取 $u_{0} \in A^{*}$,把 $u_{0}$ 的邻域 $U$ 定义如下.
\begin{align*}
U &= U(u_{0};\ x_{1}, \ldots, x_{n}, \varepsilon) \\
&= \{u \in A^{*}:\ |u(x_{j}) - u_{0}(x_{j})| < \varepsilon\ (j = 1, \ldots, n)\}
\end{align*}
其中 $\varepsilon > 0$,$x_{1}, \ldots, x_{n} \in A$,$n \in \mathbf{N}$($\mathbf{N}$ 表示自然数集).此时,
\begin{equation*}
W(u_{0}) = \{U = U(u_{0};\ x_{1}, \ldots, x_{n}, \varepsilon):\ \varepsilon > 0,\ x_{1}, \ldots, x_{n} \in A,\ n \in \mathbf{N}\}
\end{equation*}
构成 $u_{0}$ 的基本邻域系.由它导出的 $A^{*}$ 的拓扑称为 $*$-弱拓扑.

(3) 在 $\Delta(A)$ 上引入 $A^{*}$ 的 $*$-弱拓扑的相对拓扑.
\end{definition}

\begin{lemma}\label{lem:3.1.23}

(1) $\Delta(A)$ 是 Hausdorff 空间.

(2) $\Delta(A)$ 在 $A^{*}$ 中关于 $*$-弱拓扑的闭包 $\overline{\Delta(A)}$ 为:当 $A$ 具有单位元时,$\overline{\Delta(A)} = \Delta(A)$;当 $A$ 没有单位元时,$\overline{\Delta(A)} = \Delta(A) \cup \{0\}$.其中 $0 \in A^{*}$ 为 $0(x) \equiv 0$($x \in A$).
\end{lemma}

\begin{proof}
(1) 设 $\chi_{1} \neq \chi_{2}$,$\chi_{1}, \chi_{2} \in \Delta(A)$.由 $\chi_{1} \neq \chi_{2}$,存在 $x \in A$ 使 $\chi_{1}(x) \neq \chi_{2}(x)$.于是令
\begin{equation*}
\varepsilon = \frac{|\chi_{1}(x) - \chi_{2}(x)|}{3} > 0.
\end{equation*}
进一步令 $U_{j} = \{\chi \in \Delta(A):\ |\chi(x) - \chi_{j}(x)| < \varepsilon\}$($j = 1, 2$),则 $U_{1} \cap U_{2} = \emptyset$.因为若 $U_{1} \cap U_{2} \ni \chi_{0}$,则 $|\chi_{0}(x) - \chi_{j}(x)| < \varepsilon$.于是 $|\chi_{1}(x) - \chi_{2}(x)| < 2\varepsilon = 2|\chi_{1}(x) - \chi_{2}(x)|/3$,矛盾.

(2) 取 $u_{0} \in \overline{\Delta(A)}$.对 $x, y \in A$,$\varepsilon > 0$,
\begin{equation*}
U = \{u \in A^{*}:\ |u(x) - u_{0}(x)| < \varepsilon,\ |u(y) - u_{0}(y)| < \varepsilon,\ |u(xy) - u_{0}(xy)| < \varepsilon\}
\end{equation*}
这是 $u_{0}$ 的邻域,故 $U \cap \Delta(A) \neq \emptyset$.因此存在 $\chi \in U \cap \Delta(A)$ 满足 $|\chi(x) - u_{0}(x)| < \varepsilon$,$|\chi(y) - u_{0}(y)| < \varepsilon$,$|\chi(xy) - u_{0}(xy)| < \varepsilon$.从而
\begin{align*}
& |u_{0}(xy) - u_{0}(x)u_{0}(y)| \leq |u_{0}(xy) - \chi(xy)| + |\chi(xy) - \chi(x)\chi(y)| \\
& \qquad + |\chi(y) - u_{0}(y)||\chi(x)| + |\chi(x) - u_{0}(x)||u_{0}(y)| \\
& \qquad \leq \varepsilon(1 + \|x\| + |u_{0}(y)|).
\end{align*}
因此证得 $u_{0}(xy) = u_{0}(x)u_{0}(y)$,故 $u_{0} \in \Delta(A) \cup \{0\}$.若 $A$ 具有单位元,则 $u_{0} \in \Delta(A)$,故 $\overline{\Delta(A)} = \Delta(A)$.没有单位元时,$\overline{\Delta(A)} = \Delta(A) \cup \{0\}$.
\end{proof}

下面是在函数分析中众所周知的定理.

\begin{theorem}\label{thm:3.1.24}
令 $B = \{u \in A^{*}:\ \|u\| \leq 1\}$,则 $B$ 在 $*$-弱拓扑下紧.
\end{theorem}

利用定理~\ref{thm:3.1.24} 可得下面结论.

\begin{corollary}\label{cor:3.1.25}
对交换 Banach 代数 $A$ 下式成立:

(1) 若 $A$ 具有单位元,则 $\Delta(A)$ 紧.

(2) 若 $A$ 没有单位元,则 $\Delta(A)$ 的每一点具有紧邻域.
\end{corollary}

\begin{proof}
(1) 由 $\overline{\Delta(A)} = \Delta(A) \subset B = \{u \in A^{*}:\ \|u\| \leq 1\}$,利用定理~\ref{thm:3.1.24} 可知 $\Delta(A)$ 紧.

(2) 由 $\overline{\Delta(A)} = \Delta(A) \cup \{0\}$($\subset B$),利用定理~\ref{thm:3.1.24} 可知 $\Delta(A) \cup \{0\}$ 紧.因此 $\Delta(A)$ 的每一点具有紧邻域.
\end{proof}

\begin{definition}\label{def:3.1.26}
设 $A$ 为交换 Banach 代数,$\mathfrak{M}(A)$ 为 $A$ 的正则极大理想的全体.按如下方式在 $\mathfrak{M}(A)$ 上引入拓扑:由定理~\ref{thm:3.1.18},$\mathfrak{M}(A)$ 与 $\Delta(A)$ 一一对应,故把 $\mathfrak{M}(A) \ni M$ 与 $\chi_{M} \in \Delta(A)$ 对应来考虑.取 $M_{0} \in \mathfrak{M}(A)$ 的邻域
\begin{equation*}
\begin{aligned}
U(M_{0}) &= \{M \in \mathfrak{M}(A):\ |\chi_{M}(x_{j}) - \chi_{M_{0}}(x_{j})| < \varepsilon\ (j = 1, \ldots, n)\} \\
&\bigl(= U(M_{0};\ \{x_{j}\}_{1}^{n}, \varepsilon)\bigr)
\end{aligned}
\end{equation*}
(其中 $x_{j} \in A$,$\varepsilon > 0$,$n \in \mathbf{N}$),令
\begin{equation*}
\mathcal{W}(M_{0}) = \{U(M_{0};\ \{x_{j}\}_{1}^{n}, \varepsilon):\ x_{j} \in A\ (j = 1, \ldots, n),\ \varepsilon > 0,\ n \in \mathbf{N}\}
\end{equation*}
为 $M_{0}$ 的基本邻域系.这样通过 $\Delta(A)$ 的拓扑而引入拓扑的 $\mathfrak{M}(A)$ 称为 $A$ 的极大理想空间.
\end{definition}

\begin{definition}\label{def:3.1.27}
设 $A$ 为交换 Banach 代数,$\mathfrak{M}(A)$ 为 $A$ 的极大理想空间.对 $x \in A$,令 $\hat{x}(M) = \chi_{M}(x)$($\chi_{M} \in \Delta(A)$),则 $\hat{x}$ 是 $\mathfrak{M}(A)$ 上的复值函数.称其为 $x$ 的 Gelfand 变换(Gelfand 表示).在不会引起混淆时,把 $\Delta(A)$ 与 $\mathfrak{M}(A)$ 等同,把 $\hat{x}(M)$ 表为 $\hat{x}(\chi_{M})$.
\end{definition}

\begin{proposition}\label{prop:3.1.28}

(1) $\hat{x}$ 是 $\mathfrak{M}(A)$ 上的连续函数.

(2) $\mathfrak{M}(A)$ 的拓扑是使 $\hat{x}$ 连续的最弱拓扑.

(3) 对任意 $\varepsilon > 0$,$\{M \in \mathfrak{M}(A):\ |\hat{x}(M)| \geq \varepsilon\}$ 是 $\mathfrak{M}(A)$ 的紧子集.
\end{proposition}

\begin{proof}
(1) 取 $x \in A$,$M_{0} \in \mathfrak{M}(A)$,$\varepsilon > 0$.令 $U = \{M \in \mathfrak{M}(A):\ |\hat{x}(M) - \hat{x}(M_{0})| < \varepsilon\}$,则由 $\mathfrak{M}(A)$ 的拓扑定义,$U$ 是 $M_{0} \in \mathfrak{M}(A)$ 的邻域.若 $M \in U$,则 $|\hat{x}(M) - \hat{x}(M_{0})| < \varepsilon$,故 $\hat{x}$ 在 $M_{0}$ 连续.因此 $\hat{x}$ 是 $\mathfrak{M}(A)$ 上的连续函数.

(2) 取 $x \in A$,$\varepsilon > 0$,$M_{0} \in \mathfrak{M}(A)$.令 $U = \{M \in \mathfrak{M}(A):\ |\hat{x}(M) - \hat{x}(M_{0})| < \varepsilon\}$.设 $\mathcal{O}$ 为集合 $\mathfrak{M}(A)$ 上使 $\hat{x}$ 连续的拓扑的开集系.此时,因 $\hat{x}$ 关于 $\mathcal{O}$ 的拓扑连续,$U \in \mathcal{O}$.因此 $A$ 的极大理想空间 $\mathfrak{M}(A)$ 的开集的全体包含于 $\mathcal{O}$.从而极大理想空间 $\mathfrak{M}(A)$ 的拓扑是使 $\hat{x}$ 连续的 $\mathfrak{M}(A)$ 的最弱拓扑.

(3) 由系~\ref{cor:3.1.25} 证明.
\end{proof}

\begin{definition}\label{def:3.1.29}
设 $A$ 为交换 Banach 代数.此时称 $A$ 为半单的,是指若 $x \in \bigcap_{M \in \mathfrak{M}(A)} M$,则 $x = 0$.
\end{definition}

下面的定理在函数分析中众所周知.

\textbf{闭图像定理}．X, Y 为 Banach 空间,$T:\ X \to Y$ 为线性算子.进一步,若当
\begin{equation*}
\lim_{n \to \infty} x_{n} = x, \quad \lim_{n \to \infty} Tx_{n} = y
\end{equation*}
时有 $Tx = y$,则 $T$ 连续.

\begin{theorem}\label{thm:3.1.30}
设 $A, B$ 为交换 Banach 代数,$B$ 半单.此时,若 $T:\ A \to B$ 是线性映射且 $T(xy) = (Tx)(Ty)$($x, y \in A$),则 $T$ 连续.
\end{theorem}

\begin{proof}
由闭图像定理,只需证明当 $x, x_{n} \in A$ 且 $z \in B$,
\begin{equation*}
\lim_{n \to \infty} x_{n} = x \quad \text{且} \quad \lim_{n \to \infty} Tx_{n} = z
\end{equation*}
时,$Tx = z$.对任意 $\chi \in \Delta(B)$,$\chi(T(xy)) = \chi(TxTy) = \chi(Tx)\chi(Ty) = (\chi \circ T)(x)(\chi \circ T)(y)$,故 $\chi \circ T$ 是 $A$ 的特征.因此 $\chi \circ T \in \Delta(A) \cup \{0\}$.从而 $\chi \circ T:\ A \to \mathbf{C}$ 连续.因此当 $\lim_{n \to \infty} x_{n} = x$ 时,
\begin{equation*}
\lim_{n \to \infty}\chi \circ T(x_{n}) = \chi \circ T(x).
\end{equation*}
另一方面,由 $\lim_{n \to \infty} Tx_{n} = z$,$\chi \in \Delta(B)$,
\begin{equation*}
\lim_{n \to \infty}\chi(Tx_{n}) = \chi(z),
\end{equation*}
故 $\chi \circ T(x) = \chi(z)$.即 $\chi_{M}(Tx) = \chi_{M}(z)$($M \in \mathfrak{M}(B)$).于是 $0 = \chi_{M}(Tx - z) = (Tx - z)^{\wedge}(M)$($\chi \in \Delta(B)$).因 $B$ 半单,$Tx = z$.
\end{proof}

\begin{corollary}\label{cor:3.1.31}
若 $A$ 是半单的交换 Banach 代数,则 $A$ 的范数拓扑唯一确定.
\end{corollary}

\begin{proof}
设 $(A, \|\cdot\|_{1})$,$(A, \|\cdot\|_{2})$ 为半单的交换 Banach 代数.在定理~\ref{thm:3.1.30} 中令 $T$ 为恒等映射.$T$ 线性且 $T(xy) = T(x)T(y)$,故 $T$ 满足定理~\ref{thm:3.1.30} 的假设.因此存在 $C_{1} > 0$ 使得 $\|x\|_{2} \leq C_{1}\|x\|_{1}$.反之同理,故存在 $C > 0$ 使得 $C^{-1}\|x\|_{1} \leq \|x\|_{2} \leq C\|x\|_{1}$($x \in A$).
\end{proof}

例．(1) 交换 Banach 代数 $C(\mathbf{T})$ 的 Gelfand 变换是恒等映射.事实上,用定理~\ref{thm:3.1.16} 后的例中的记号,极大理想具有 $I(x_{0})$($x_{0} \in \mathbf{T}$)的形式,且由 $x_{0} \in \mathbf{T}$ 到 $I(x_{0}) \in \mathfrak{M}(C(\mathbf{T}))$ 的映射连续,$\mathbf{T}$ 与 $\mathfrak{M}(C(\mathbf{T}))$ 同胚.因此成为同胚映射.

(2) 交换 Banach 代数 $\ell^{1}$ 的 Gelfand 变换对 $a = \{a_{n}\}$ 为 $\hat{a}(x) = \sum_{n \in \mathbf{Z}} a_{n}e^{inx}$.极大理想空间是 $\mathbf{T}$.事实上,由 $x \in \mathbf{T}$ 到 $M_{x} = \{a = \{a_{n}\} \in \ell^{1}:\ \hat{a}(x) = 0\} \in \mathfrak{M}(\ell^{1})$ 的映射是双射且连续.因此 $\mathbf{T}$ 与 $\mathfrak{M}(\ell^{1})$ 同胚(参照 3.2 节 $A(\mathbf{T})$ 一节).

\subsection{正则交换 Banach 代数}\label{sec:3.1.5}

\begin{definition}\label{def:3.1.32}
设 $A$ 为交换 Banach 代数.

(1) 对 $D \subset A$,令 $h(D) = \{M \in \mathfrak{M}(A):\ M \supset D\}$.此时利用 Gelfand 变换,$h(D) = \{\chi \in \Delta(A):\ \hat{x} = 0\ (x \in D)\}$.

(2) 对 $E \subset \mathfrak{M}(A)$,令 $k(E) = \bigcap\{M:\ M \in E\}$.此时利用 Gelfand 变换,也可认为 $k(E) = \{x \in A:\ \hat{x}(\chi) = 0\ (\chi \in E)\}$.其中,对空集 $\emptyset$,令 $k(\emptyset) = A$.
\end{definition}

\begin{theorem}\label{thm:3.1.33}
设 $E, E_{1}, E_{2} \subset \mathfrak{M}(A)$,则下式成立.

(1) $hk(E) \supset E$.

(2) $hk(E) = hk(hk(E))$.

(3) $hk(\emptyset) = \emptyset$.

(4) $hk(E_{1} \cup E_{2}) = hk(E_{1}) \cup hk(E_{2})$.
\end{theorem}

\begin{proof}
(1) 若 $\chi \in E$,则当 $x \in k(E)$ 时 $\hat{x}(\chi) = 0$,故 $\chi \in hk(E)$.

(2) 由 (1),$khk(E) \subset k(E)$.另一方面,对 $\chi \in hk(E)$ 取 $x \in k(E)$,由定义~\ref{def:3.1.32},$\hat{x}(\chi) = 0$,故 $k(E) \subset khk(E)$.因此 $khk(E) = k(E)$.从而得 $kh(hk(E)) = hk(E)$.

(3) 由 $k(\emptyset) = A$,$h(A) = \emptyset$,$hk(\emptyset) = \emptyset$.

(4) 由 $E_{1} \cup E_{2} \supset E_{j}$($j = 1, 2$),$hk(E_{1} \cup E_{2}) \supset hk(E_{j})$($j = 1, 2$).因此 $hk(E_{1} \cup E_{2}) \supset hk(E_{1}) \cup hk(E_{2})$.反之,设 $\chi \notin hk(E_{1}) \cup hk(E_{2})$,则存在 $x_{j} \in k(E_{j})$,$\chi(x_{j}) \neq 0$ 的 $x_{j} \in A$($j = 1, 2$).此时 $x_{1}x_{2} \in k(E_{1} \cup E_{2})$.事实上,若 $\chi' \in E_{1} \cup E_{2}$,当 $\chi' \in E_{1}$ 时,由 $x_{1} \in k(E_{1})$,$\chi'(x_{1}) = 0$,故 $\chi'(x_{1}x_{2}) = 0$.同理,$\chi' \in E_{2}$ 的情形也有 $\chi'(x_{1}x_{2}) = 0$.因此由 $\chi(x_{1}x_{2}) \neq 0$,$\chi \notin hk(E_{1} \cup E_{2})$.于是若 $\chi \notin hk(E_{1}) \cup hk(E_{2})$,则 $\chi \notin hk(E_{1} \cup E_{2})$,故 $hk(E_{1} \cup E_{2}) \subset hk(E_{1}) \cup hk(E_{2})$.
\end{proof}

\begin{definition}\label{def:3.1.34}
对 $E \subset \mathfrak{M}(A)$ 对应 $hk(E)$,由定理~\ref{thm:3.1.33} 在 $\mathfrak{M}(A)$ 上引入了闭包拓扑(把满足 $E = hk(E)$ 的 $E$ 作为闭集来引入拓扑).称其为 $\mathfrak{M}(A)$ 上的 hk 拓扑.
\end{definition}

\begin{definition}\label{def:3.1.35}
对交换 Banach 代数 $A$,令 $\hat{A} = \{\hat{x}:\ x \in A\}$.设 $F$ 为 $\mathfrak{M}(A)$ 的闭集,$\chi_{0} \notin F$,$\chi_{0} \in \mathfrak{M}(A)$,若存在 $x \in A$ 使得 $\hat{x}(\chi) = 0$($\chi \in F$)且 $\hat{x}(\chi_{0}) \neq 0$,则称交换 Banach 代数 $A$ 为正则的.
\end{definition}

\begin{theorem}\label{thm:3.1.36}
下列命题等价.

(1) 交换 Banach 代数 $A$ 正则.

(2) $\mathfrak{M}(A)$ 的拓扑与 hk 拓扑一致.

特别地,当 $A$ 具有单位元时,下列也等价.

(3) 设 $F$ 为 $\mathfrak{M}(A)$ 的闭集,$K$ 为 $\mathfrak{M}(A)$ 的紧集且 $K \cap F = \emptyset$,则存在 $x \in A$ 使得 $\hat{x}(\chi) = 0$($\chi \in F$),$\hat{x}(\chi) = 1$($\chi \in K$).
\end{theorem}

\begin{proof}
首先证明 (1) 与 (2) 等价.设 $F$ 为 $\mathfrak{M}(A)$ 的闭集.此时 $hk(F) \subset F$.事实上,由 (1) 的假设,若 $\chi \notin F$,则存在 $x \in A$ 使 $\hat{x}(F) = 0$,$\hat{x}(\chi) \neq 0$.此时由 $x \in k(F)$,$\chi \notin hk(F)$,故得 $hk(F) \subset F$.从而由定理~\ref{thm:3.1.33},$F \subset hk(F)$,故 $hk(F) = F$.因此当 $F$($\subset \mathfrak{M}(A)$)为闭集时,在 hk 拓扑下也是闭集.反之,设 $F$($\subset \mathfrak{M}(A)$)在 hk 拓扑下为闭集.此时 $hk(F) = F$,且 $h(k(F))$ 是 $\mathfrak{M}(A)$ 的闭集,故 $F$ 是 $\mathfrak{M}(A)$ 的闭集.假设 (3),则 (1) 显然.反之假设 (1) 证明 (3).为此考虑下列引理.
\end{proof}

\begin{lemma}\label{lem:3.1.37}
设 $A$ 为具有单位元 $e$ 的正则交换 Banach 代数,$I$ 为 $A$ 的理想,$E$ 为 $\mathfrak{M}(A)$ 的非空紧集且 $h(I) \cap E = \emptyset$.此时存在 $x \in I$ 使得 $\hat{x}(\chi) = 1$($\chi \in E$).
\end{lemma}

\begin{proof}
商环 $A/k(E)$ 的极大理想空间 $\mathfrak{M}(A/k(E))$ 与 $E$ 同胚.事实上,若 $\chi \in E$,令 $\widetilde{\chi}(x + k(E)) = \chi(x)$,则可定义 $A/k(E)$ 上的非零特征.进一步,把 $\chi \in E$ 对应到 $\widetilde{\chi}$ 的映射是从 $E$ 到 $\Delta(A/k(E))$ 的双射.因为,对 $\widetilde{\chi_{0}} \in \Delta(A/k(E))$,令 $\chi_{0}(x) = \widetilde{\chi_{0}}(x + k(E))$,若 $x \in k(E)$ 则 $\widetilde{\chi_{0}}(x) = 0$,由 (1),$hk(E) = E$,故 $\chi_{0} \in E$.因此该对应是满射.单射显然.又该映射连续,故 $E$ 与 $\mathfrak{M}(A/k(E))$ 同胚.现在,由假设 $h(I \cup k(E)) \subset h(I) \cap hk(E) = h(I) \cap E = \emptyset$,故 $h(I \cup k(E)) = \emptyset$.因此不存在同时包含 $I$,$k(E)$ 的 $A$ 的极大理想.即同时包含 $I$,$k(E)$ 的理想是 $A$.于是设 $e$ 为 $A$ 的单位元,存在 $u \in I$,$u \notin k(E)$ 使 $u + k(E) = e + k(E)$.因此存在 $v \in k(E)$ 使 $u = e + v$.即 $u \in I$,且 $\hat{u}(\chi) = \hat{e}(\chi) = 1$($\chi \in E$).
\end{proof}

((1) 蕴含 (3) 的证明)

在补题~\ref{lem:3.1.37} 中令 $I = k(F)$,$E = K$,则存在 $x \in A$ 使 $x \in k(F)$,$\hat{x}(\chi) = 1$($\chi \in K$).即 $x \in A$,$\hat{x}(\chi) = 0$($\chi \in F$),$\hat{x}(\chi) = 1$($\chi \in K$).

\begin{theorem}\label{thm:3.1.38}
设 $A$ 为具有单位元 $e$ 的正则交换 Banach 代数.设 $U_{1}, \ldots, U_{n}$ 为 $\mathfrak{M}(A)$ 的开覆盖.此时存在 $x_{j} \in A$($j = 1, \ldots, n$)使得 $\supp \widehat{x}_{j} \subset U_{j}$($j = 1, \ldots, n$),$\sum_{j=1}^{n} \widehat{x}_{j}(\chi) = 1$($\chi \in \mathfrak{M}(A)$).
\end{theorem}

\begin{proof}
$n = 1$ 时取 $x_{1} = e$ 即可.假设 $n - 1$ 时成立.设 $V, V_{3}, V_{4}, \ldots, V_{n}$ 为 $\mathfrak{M}(A)$ 的开覆盖且 $V \subset \overline{V} \subset U_{1} \cup U_{2}$,$V_{j} \subset \overline{V_{j}} \subset U_{j}$($j = 3, \ldots, n$).此时由 $n - 1$ 的假设,存在 $x, x_{3}, \ldots, x_{n} \in A$ 满足 $\supp \hat{x} \subset V$,$\supp \widehat{x_{j}} \subset V_{j}$,$\hat{x}(\chi) + \widehat{x_{3}}(\chi) + \cdots + \widehat{x_{n}}(\chi) = 1$($\chi \in \mathfrak{M}(A)$).此时令 $E = \supp \hat{x}$,$E_{1} = E \setminus U_{2}$,$E_{2} = E \setminus U_{1}$,由定理~\ref{thm:3.1.36} 存在 $y \in A$ 使 $\supp \hat{y} \subset V \cap U_{1}$,$\hat{y}(\chi) = 1$($\chi \in E_{1}$).事实上,$(V \cap U_{1})^{c} \cap E_{1} = \emptyset$.此时令 $x_{1} = xy$,$x_{2} = x(e - y)$,则 $x_{1} + x_{2} = x$,且 $\supp \widehat{x_{1}} \subset V \cap U_{1} \subset U_{1}$,$\supp \widehat{x_{2}} \subset E \cap U_{2} \subset U_{2}$.因此满足 $\widehat{x_{1}} + \cdots + \widehat{x_{n}} = 1$,$\supp \widehat{x_{j}} \subset U_{j}$($j = 1, \ldots, n$).
\end{proof}

\begin{corollary}\label{cor:3.1.39}
设 $A$ 为具有单位元 $e$ 的半单正则交换 Banach 代数.设 $I$ 为 $A$ 的理想,$f$ 为 $A$ 的极大理想空间 $\mathfrak{M}(A)$ 上的连续函数.此时,若对任意 $\chi_{0} \in \mathfrak{M}(A)$ 存在 $x_{0} \in I$ 以及 $\chi_{0}$ 的某邻域 $U_{0}$,使得 $f(\chi) = \widehat{x_{0}}(\chi)$($\chi \in U_{0}$),则存在 $x \in I$ 使 $\hat{x}(\chi) = f(\chi)$($\chi \in \mathfrak{M}(A)$).
\end{corollary}

\begin{proof}
对每个 $\chi_{0} \in \mathfrak{M}(A)$ 按假设对应 $\chi_{0}$ 的邻域 $U_{0}$,则 $\{U_{0}\}$ 是 $\mathfrak{M}(A)$ 的开覆盖,因 $\mathfrak{M}(A)$ 紧,从 $\{U_{0}\}$ 中可取 $\mathfrak{M}(A)$ 的有限开覆盖 $U_{1}, \ldots, U_{n}$.对此,由定理~\ref{thm:3.1.38} 选取 $z_{1}, \ldots, z_{n} \in A$,$\supp \widehat{z}_{j} \subset U_{j}$,$\widehat{z}_{1} + \cdots + \widehat{z}_{n} = 1$.此时
\begin{equation*}
f(\chi)\widehat{z_{1}}(\chi) + \cdots + f(\chi)\widehat{z_{n}}(\chi) = f(\chi) \quad (\chi \in \mathfrak{M}(A))
\end{equation*}
成立.这里,由假设,不妨设存在 $x_{j} \in I$ 使 $f(\chi) = \hat{x}_{j}(\chi)$($\chi \in U_{j}$),故
\begin{equation*}
\widehat{x_{1}}(\chi)\widehat{z_{1}}(\chi) + \cdots + \widehat{x_{n}}(\chi)\widehat{z_{n}}(\chi) = f(\chi) \quad (\chi \in \mathfrak{M}(A))
\end{equation*}
成立.因 $A$ 半单且 $I$ 是理想,令 $x = x_{1}z_{1} + \cdots + x_{n}z_{n}$,则 $x \in I$,$\hat{x} = f$.
\end{proof}

例．(1) 交换 Banach 代数 $C(\mathbf{T})$ 的 Gelfand 变换是恒等映射,故正则.

(2) 交换 Banach 代数 $\ell^{1}$ 的极大理想空间是 $\mathbf{T}$,且 $\mathbf{T}$ 上的 2 次连续可微函数 $f(x)$ 可表为 $f(x) = \sum_{n \in \mathbf{Z}} a_{n}e^{inx}$,$\sum_{n \in \mathbf{Z}} |a_{n}| < \infty$,故 $\ell^{1}$ 正则(参照 3.2.1 小节).

\subsection{谱综合}\label{sec:3.1.6}

\begin{definition}\label{def:3.1.40}
设 $A$ 为正则交换 Banach 代数.

(1) 设 $E \subset \mathfrak{M}(A)$ 为闭集,$I$ 为正则闭理想.当满足 $h(I) = E$ 的 $I$ 唯一确定时,称 $E$ 为交换 Banach 代数 $A$ 的谱综合集.

(2) 当 $\mathfrak{M}(A)$ 中存在不是谱综合集的闭集时,称交换 Banach 代数 $A$ 为谱综合不可能的.

(3) 当 $\mathfrak{M}(A)$ 的所有闭集都是谱综合集时,称交换 Banach 代数 $A$ 为谱综合可能的.
\end{definition}

以下在本节中,设 $A$ 为具有单位元 $e$ 的正则交换 Banach 代数.

\begin{definition}\label{def:3.1.41}
对 $E \subset \mathfrak{M}(A)$,令 $I(E) = \{x \in A:\ \hat{x}(\chi) = 0\ (\chi \in E)\}$,则 $I(E)$ 是闭理想.又令 $j(E) = \{x \in A:\ \supp \hat{x} \text{ 是紧集,且 } \supp \hat{x} \cap E = \emptyset\}$,$J(E) = \overline{j(E)}$($j(E)$ 在 $A$ 中的闭包),则 $j(E), J(E)$ 都是理想.
\end{definition}

\begin{lemma}\label{lem:3.1.42}
设 $A$ 半单.若 $I \subset \mathfrak{M}(A)$ 为闭集,则 $I(E)$ 是满足 $h(I) = E$ 的理想 $I$ 中最大的理想,$j(E)$ 是满足 $h(I) = E$ 的理想 $I$ 中最小的理想.
\end{lemma}

\begin{proof}
当 $h(I) = E$ 时,若 $x \in I$,则 $\widehat{x}(\chi) = 0$($\chi \in E$),故 $x \in I(E)$.又 $h(j(E)) \supset E$ 显然.若 $\chi \notin E$,则因 $A$ 正则,存在 $x_{0} \in A$ 使得 $\widehat{x_{0}}(\chi) = 1$,$\supp \widehat{x_{0}}$ 是紧集,$\supp \widehat{x_{0}} \cap E = \emptyset$.此时 $x_{0} \in j(E)$,故 $\chi \notin h(j(E))$,得 $h(j(E)) = E$.又取满足 $h(I) = E$ 的理想 $I$.此时,对满足 $K \cap E = \emptyset$ 的紧集 $K \subset \mathfrak{M}(A)$,由补题~\ref{lem:3.1.37} 存在 $a \in I$ 使 $\widehat{a}(\chi) = 1$($\chi \in K$),$\widehat{a}(\chi) = 0$($\chi \in E$).因此,对满足 $\supp \widehat{x} \subset K$ 的 $x \in A$,$\widehat{a}\widehat{x} = \widehat{x}$,因 $A$ 半单,$x = ax \in I$.从而 $j(E) \subset I$.
\end{proof}

\begin{note}\label{note:3.1.43}
由补题~\ref{lem:3.1.42},闭集 $E \subset \mathfrak{M}(A)$ 是 $A$ 的谱综合集的充要条件是 $I(E) = J(E)$.
\end{note}

\begin{definition}\label{def:3.1.44}
交换 Banach 代数 $A$ 满足 Ditkin 条件,是指当 $\chi \in \mathfrak{M}(A)$ 与 $x \in A$ 满足 $\hat{x}(\chi) = 0$ 时,存在 $\{x_{n}\} \subset A$ 使得 $\widehat{x_{n}} = 0$ 在 $\chi$ 的某邻域内成立,且 $\lim_{n \to \infty} x_{n}x = x$.
\end{definition}

\begin{theorem}\label{thm:3.1.45}
设 $A$ 为具有单位元 $e$ 的半单正则交换 Banach 代数,且满足 Ditkin 条件.设 $I$ 为 $A$ 的闭理想,$x_{0} \in kh(I)$.此时,若 $h(\{x_{0}\}) \cap h(I)$ 的拓扑边界不含非空完全集,则 $x_{0} \in I$.
\end{theorem}

\begin{proof}
对给定的 $x_{0} \in kh(I)$,令
\begin{equation*}
E = \{\chi \in \Delta(A):\ \text{对 } I \text{ 的任意元 } x \text{,不存在 } \chi \text{ 的邻域使 } \widehat{x_{0}} = \widehat{x} \text{ 成立}\},
\end{equation*}
若能证明 $E = \emptyset$,则由系~\ref{cor:3.1.39} 即得定理~\ref{thm:3.1.45}.现在,若 $\chi \notin h(I)$,由补题~\ref{lem:3.1.37},存在 $a \in I$ 及 $\chi$ 的邻域使 $\widehat{a} = 1$,故 $\chi \notin E$.又,若 $\chi \in \mathfrak{M}(A)$ 不在 $h(\{x_{0}\})$ 的拓扑边界上,则 $\chi$ 是 $h(\{x_{0}\})$ 的内点或 $\chi \notin h(\{x_{0}\})$.若 $\chi$ 是 $h(\{x_{0}\})$ 的内点,则 $\widehat{x_{0}} = 0$ 在 $\chi$ 的某邻域内成立,故 $\chi \notin E$.当 $\chi \notin h(\{x_{0}\})$ 时,由 $x_{0}$($\in kh(I)$)的假设,$h(\{x_{0}\}) \supset hkh(I) = h(I)$,故 $\chi \notin h(I)$.因此 $\chi \notin E$.综上,证明了 $E$ 包含于 $h(\{x_{0}\}) \cap h(I)$ 的拓扑边界.进一步,$E$ 没有孤立点.事实上,若 $\chi \in E$ 是孤立点,则存在 $\chi$ 的邻域 $U$,使得 $U$ 在 $\mathfrak{M}(A)$ 中的闭包 $\overline{U}$ 满足 $\overline{U} \cap E = \{\chi\}$.因 $\widehat{x_{0}}(\chi) = 0$,由 Ditkin 条件,存在 $\{z_{n}\} \subset A$ 使得 $\widehat{z_{n}} = 0$ 在 $\chi$ 的邻域内成立,且 $\lim_{n \to \infty} z_{n}x_{0} = x_{0}$.此时,取在 $\chi$ 的邻域内 $\widehat{y} = 1$ 且 $\supp \widehat{y} \subset U$ 的 $y \in A$,由系~\ref{cor:3.1.39},$z_{n}yx_{0} \in I$.因此由 $\lim_{n \to \infty} z_{n}yx_{0} = yx_{0} \in I$,$yx_{0} \in I$.于是由 $E$ 的定义,$\chi \notin E$,产生矛盾.综上,$E$ 是完全集.由假设,$h(\{x_{0}\}) \cap h(I)$ 的拓扑边界不含非空完全集,故 $E = \emptyset$.
\end{proof}

\begin{corollary}\label{cor:3.1.46}
设 $A$ 为具有单位元 $e$ 的半单正则且满足 Ditkin 条件的交换 Banach 代数,若闭集 $E \subset \mathfrak{M}(A)$ 的拓扑边界不含非空完全集,则 $E$ 是($A$ 的)谱综合集.
\end{corollary}

\begin{proof}
设 $x_{0} \in I(E)$.此时,由 $h(J(E)) = E$,$x_{0} \in k(E) = kh(J(E))$.进一步,利用假设,$h(\{x_{0}\}) \cap E = E$,故 $h(\{x_{0}\}) \cap h(J(E))$ 的拓扑边界不含完全集.因此由定理~\ref{thm:3.1.45},$x_{0} \in J(E)$,故证得 $I(E) = J(E)$.
\end{proof}

\subsection{半单交换 Banach 代数的作用函数}\label{sec:3.1.7}

\begin{definition}\label{def:3.1.47}
设 $A$ 为半单交换 Banach 代数.设 $G \subset \mathbf{C}$,$F:\ G \to \mathbf{C}$ 为函数.对 $x \in A$,$\Sp_{A}(x) \subset G$,若存在 $y \in A$ 使得
\begin{equation*}
F(\widehat{x}(M)) = \widehat{y}(M) \quad (M \in \mathfrak{M}(A))
\end{equation*}
则称 $F$ 为 $A$ 的作用函数.
\end{definition}

\begin{theorem}\label{thm:3.1.48}
设 $A$ 为具有单位元 $e$ 的半单交换 Banach 代数.若 $F$ 是单连通区域 $G$($\subset \mathbf{C}$)上的正则函数,则 $F$ 是 $A$ 的作用函数.即,设 $\Gamma$ 为 $G$ 中具有长度的闭曲线,且 $\supp \hat{x}$ 包含于 $\Gamma$ 的内部,则
\begin{equation*}
F(x) = \frac{1}{2\pi i}\int_{\Gamma}\frac{F(z)}{z - x}\dd z \in A
\end{equation*}
可以定义,且 $\widehat{F(x)}(M) = F(\widehat{x}(M))$($M \in \mathfrak{M}(A)$)成立.
\end{theorem}

\begin{proof}
由向量值函数的函数论证明(cf.\ \cite{ch3:3}, \cite{ch3:8}).细节从略.
\end{proof}

关于解析函数作用于 Banach 代数,在 \cite{ch3:1}, \cite{ch3:2} 等中有论述.

% ============ 第3章 §3.2 若干交换Banach代数的Gelfand表示(原书 pp.226-237) ============
\section{若干交换 Banach 代数的 Gelfand 表示}\label{sec:3.2}
\index{Fourier 级数}
\index{Fourier 系数}
\index{测度代数}
\index{Fourier--Stieltjes 变换}
\index{乘子!(单位圆周上的函数的)}
\index{乘子!乘子算子(单位圆周上的函数的)}
\index{Calderón--Zygmund 定理}
\index{Katznelson 问题}
\index{Parreau 定理}
\index{Wiener--Pitt 现象}
\index{Williamson 定理}


本节作为 Fourier 分析中出现的交换 Banach 代数的例子,考察具有绝对收敛 Fourier 级数的连续函数空间,Fourier 变换的空间等,并叙述它们的极大理想空间.另外,Riemann-Lebesgue 引理等初等 Fourier 分析的知识视为已知.

\subsection{\texorpdfstring{关于 $A(\mathbf{T})$}{关于 A(T)}}\label{sec:3.2.1}

\begin{definition}\label{def:3.2.1}
定义
\begin{equation*}
A(\mathbf{T}) = \left\{f \in C(\mathbf{T}):\ \|f\|_{A(\mathbf{T})} = \sum_{n=-\infty}^{\infty}|\hat{f}(n)| < \infty\right\}.
\end{equation*}
其中,
\begin{equation*}
\hat{f}(n) = \frac{1}{2\pi}\int_{0}^{2\pi}f(x)e^{-inx}\dd x
\end{equation*}
(一般地,对可积函数 $f$,称 $\hat{f}(n)$ 为 $f$ 的 Fourier 系数,称形式级数 $\sum_{n \in \mathbf{Z}}\hat{f}(n)e^{inx}$ 为 $f$ 的 Fourier 级数).又,和与数乘定义为
\begin{equation*}
(\alpha f + \beta g)(x) = \alpha f(x) + \beta g(x) \quad (\alpha, \beta \in \mathbf{C},\ f, g \in A(\mathbf{T})),
\end{equation*}
积定义为
\begin{equation*}
(fg)(x) = f(x)g(x) \quad (f, g \in A(\mathbf{T})),
\end{equation*}
则 $A(\mathbf{T})$ 成为环.进一步,$\|\cdot\|_{A(\mathbf{T})}$ 满足范数的性质,且关于这个范数完备.又,$\|fg\|_{A(\mathbf{T})} \leq \|f\|_{A(\mathbf{T})}\|g\|_{A(\mathbf{T})}$ 成立,故由具有绝对收敛 Fourier 级数的连续函数构成的空间 $A(\mathbf{T})$ 是以常数函数 $1$ 为单位元的交换 Banach 代数.
\end{definition}

\begin{theorem}\label{thm:3.2.2}
$A(\mathbf{T})$ 的极大理想空间 $\mathfrak{M}(A(\mathbf{T}))$ 与 $\mathbf{T}$ 等同.Gelfand 变换是恒等映射.又,$A(\mathbf{T})$ 半单.
\end{theorem}

\begin{proof}
由 $\mathfrak{M}(A(\mathbf{T})) = \Delta(A(\mathbf{T}))$,考察 $\chi \in \Delta(A(\mathbf{T}))$.因 $e^{i\theta} \in A(\mathbf{T})$,
\begin{equation*}
\chi(e^{i\theta})\chi(e^{-i\theta}) = \chi(e^{i\theta}e^{-i\theta}) = \chi(1) = 1
\end{equation*}
成立.因此 $\chi(e^{i\theta})\chi(e^{-i\theta}) = 1$.进一步,$|\chi(e^{i\theta})| \leq \|e^{i\theta}\|_{A(\mathbf{T})} = 1$,故 $|\chi(e^{i\theta})| = 1$.从而可令 $\chi(e^{i\theta}) = e^{i\xi}$($\xi \in [0, 2\pi)$).于是,对 $f \in A(\mathbf{T})$,$f(\theta) = \sum_{n \in \mathbf{Z}}\hat{f}(n)e^{in\theta}$,令 $f_{N}(\theta) = \sum_{|n| \leq N}\hat{f}(n)e^{in\theta}$,由 $\|f_{N} - f\|_{A(\mathbf{T})} \to 0$($N \to \infty$),
\begin{equation*}
|\chi(f) - \chi(f_{N})| = |\chi(f - f_{N})| \leq \|f_{N} - f\|_{A(\mathbf{T})} \to 0 \quad (N \to \infty)
\end{equation*}
成立.因此
\begin{equation*}
\chi(f) = \sum_{n \in \mathbf{Z}}\hat{f}(n)e^{in\xi}.
\end{equation*}
即,对 $\chi \in \Delta(A(\mathbf{T}))$ 对应 $\xi \in [0, 2\pi)$.反之,对 $\xi \in [0, 2\pi)$ 考虑 $\chi_{\xi}:\ f \mapsto f(\xi)$,则 $\chi_{\xi} \in \Delta(A(\mathbf{T}))$.因此,由 $\xi \in [0, 2\pi)$ 到 $\chi_{\xi} \in \Delta(A(\mathbf{T}))$ 的对应是满射.由 $\chi_{\xi}(f) = \chi_{\xi'}(f)$($f \in A(\mathbf{T})$)得 $e^{i\xi} = e^{i\xi'}$($\xi, \xi' \in [0, 2\pi)$),故 $\xi = \xi'$,是单射.又,当 $\lim_{n \to \infty}\xi_{n} = \xi$ 时,
\begin{equation*}
\lim_{n \to \infty}\chi_{\xi_{n}}(f) = \lim_{n \to \infty}f(\xi_{n}) = f(\xi)
\end{equation*}
故
\begin{equation*}
\lim_{n \to \infty}\chi_{\xi_{n}}(f) = \chi_{\xi}(f)
\end{equation*}
成立.因 $\mathbf{T}$ 是紧 Hausdorff 空间,$\Delta(A(\mathbf{T}))$ 与 $\mathbf{T}$ 同胚.因此 Gelfand 变换是恒等映射.由此,$A(\mathbf{T})$ 半单.
\end{proof}

\begin{note}\label{note:3.2.3}
$n \in \mathbf{N}$,令 $\mathbf{T}^{n} = \{\theta = (\theta_{1}, \ldots, \theta_{n}):\ \theta_{j} \in [0, 2\pi)\ (j = 1, 2, \ldots, n)\}$.定义~\ref{def:3.2.1} 的多维版本定义为
\begin{equation*}
A(\mathbf{T}^{n}) = \left\{f \in C(\mathbf{T}^{n}):\ f(\theta) = \sum_{m = (m_{1}, \ldots, m_{n}) \in \mathbf{Z}^{n}}\hat{f}(m)e^{im \cdot \theta},\ \|f\|_{A(\mathbf{T}^{n})} = \sum_{m \in \mathbf{Z}^{n}}|\hat{f}(m)| < \infty\right\}.
\end{equation*}
其中 $m \cdot \theta = \sum_{j=1}^{n}m_{j}\theta_{j}$,
\begin{equation*}
\hat{f}(m) = \frac{1}{(2\pi)^{n}}\int_{0}^{2\pi}\int_{0}^{2\pi}\cdots\int_{0}^{2\pi}f(\theta_{1}, \ldots, \theta_{n})e^{-im \cdot \theta}\dd\theta.
\end{equation*}
此时,对 $A(\mathbf{T}^{n})$ 也有与定理~\ref{thm:3.2.2} 相同的结论.
\end{note}

\subsection{\texorpdfstring{Fourier 变换构成的交换 Banach 代数 $A(\mathbb{R}^n)$}{Fourier 变换构成的交换 Banach 代数 A(R^n)}}\label{sec:3.2.2}

\begin{definition}\label{def:3.2.4}
$n$ 维 Euclid 空间 $\mathbf{R}^{n}$ 上的 Lebesgue 可积函数空间定义为
\begin{equation*}
L^{1}(\mathbf{R}^{n}) = \left\{f:\ \|f\|_{1} = \frac{1}{(2\pi)^{n}}\int_{\mathbf{R}^{n}}|f(x)|\dd x < \infty\right\}.
\end{equation*}
进一步,由
\begin{equation*}
\hat{f}(\xi) = \frac{1}{(2\pi)^{n}}\int_{\mathbf{R}^{n}}f(x)e^{-i\xi \cdot x}\dd x
\end{equation*}
定义 $f$ 的 Fourier 变换.其中 $\xi \cdot x$ 是内积.
\end{definition}

对 $f, g \in L^{1}(\mathbf{R}^{n})$,$\alpha, \beta \in \mathbf{C}$,由 $(\alpha f + \beta g)(x) = \alpha f(x) + \beta g(x)$,以及 $f * g(x) = \frac{1}{(2\pi)^{n}}\int_{\mathbf{R}^{n}}f(x - y)g(y)\dd y$,$L^{1}(\mathbf{R}^{n})$ 关于和,数乘,积成为环.又,$L^{1}(\mathbf{R}^{n})$ 以 $\|\cdot\|_{1}$ 为范数成为 Banach 空间,且满足 $\|f * g\|_{1} \leq \|f\|_{1}\|g\|_{1}$.综上,$L^{1}(\mathbf{R}^{n})$ 成为交换 Banach 代数.

又,对测度空间 $(X, \mu)$,用 $L^{\infty}(X)$ 表示以 $\|f\|_{L^{\infty}} = \inf\{M:\ |f(x)| < M\ (\mu\text{-a.e.})\}$ 为范数的 Banach 空间.以后,$L^{1}(\mathbf{R})^{*} = L^{\infty}(\mathbf{R})$,$L^{1}(\mathbf{T})^{*} = L^{\infty}(\mathbf{T})$ 等视为已知.

\begin{theorem}\label{thm:3.2.5}
下列事实成立.

(1) $L^{1}(\mathbf{R}^{n})$ 是没有单位元的交换 Banach 代数.

(2) $\mathfrak{M}(L^{1}(\mathbf{R}^{n})) = \mathbf{R}^{n}$.$L^{1}(\mathbf{R}^{n})$ 的 Gelfand 变换是 Fourier 变换.

(3) 令 $A(\mathbf{R}^{n}) = \{\hat{f}:\ f \in L^{1}(\mathbf{R}^{n})\}$,$\|\hat{f}\|_{A(\mathbf{R}^{n})} = \|f\|_{1}$,则 $A(\mathbf{R}^{n})$ 关于函数通常运算的和,数乘,积与范数 $\|\cdot\|_{A(\mathbf{R}^{n})}$ 成为没有单位元的交换 Banach 代数.

(4) $\mathfrak{M}(A(\mathbf{R}^{n})) = \mathbf{R}^{n}$.$A(\mathbf{R}^{n})$ 的 Gelfand 变换是恒等映射.
\end{theorem}

\begin{proof}
(1) 若 $L^{1}(\mathbf{R}^{n})$ 具有单位元 $e$,则 $e * e = e$.由此 $\hat{e} \cdot \hat{e} = \hat{e}$.因此 $\hat{e}(\xi) = 0, 1$($\xi \in \mathbf{R}^{n}$).由 Riemann-Lebesgue 引理,$\hat{e}$ 是在无穷远点取 $0$ 的连续函数,故 $\hat{e} \equiv 0$.从而由 Fourier 变换的性质,$e \equiv 0$.这与 $e$ 是单位元相矛盾.

(2) 以 $n = 1$ 证明($n > 1$ 时同理).对 $\xi \in \mathbf{R}$,令 $\chi_{\xi}(f) = \hat{f}(\xi)$,由 Fourier 变换的性质,$\chi_{\xi} \in \Delta(L^{1}(\mathbf{R}))$.又,若 $\chi_{\xi}(f) = \chi_{\xi'}(f)$($f \in L^{1}(\mathbf{R})$),则 $\xi = \xi'$,故此对应是单射.进一步,若 $\lim_{n \to \infty}\xi_{n} = \xi$,由 $\lim_{n \to \infty}\hat{f}(\xi_{n}) = \hat{f}(\xi)$,$\xi \mapsto \chi_{\xi}$ 连续.而且此对应是满射.事实上,设 $\chi \in \Delta(L^{1}(\mathbf{R}))$,则 $\chi:\ L^{1}(\mathbf{R}) \to \mathbf{C}$ 是同态.这里,$L^{1}(\mathbf{R})$ 上连续线性泛函的空间是 $L^{\infty}(\mathbf{R})$,故由 $|\chi(f)| \leq \|f\|_{1}$,存在 $\phi \in L^{\infty}(\mathbf{R})$ 使得
\begin{equation*}
\chi(f) = \int f(x)\phi(x)\frac{\dd x}{2\pi}.
\end{equation*}
进一步,$\|\phi\|_{L^{\infty}} = 1$.此时,若 $f, g \in L^{1}(\mathbf{R})$,则
\begin{align*}
\int \chi(f)g(y)\phi(y)\frac{\dd y}{2\pi} & = \chi(f)\chi(g) \\
& = \chi(f * g) \\
& = \int (f * g)(x)\phi(x)\frac{\dd x}{2\pi} \\
& = \iint g(y)f(x - y)\frac{\dd y}{2\pi}\phi(x)\frac{\dd x}{2\pi} \\
& = \int g(y)\int f(x - y)\phi(x)\frac{\dd x}{2\pi}\frac{\dd y}{2\pi} \\
& = \int g(y)\chi(f_{y})\frac{\dd y}{2\pi}
\end{align*}
成立.其中 $f_{y}(x) = f(x - y)$.因此
\begin{equation*}
\int \chi(f)\phi(y)g(y)\frac{\dd y}{2\pi} = \int \chi(f_{y})g(y)\frac{\dd y}{2\pi} \quad (g \in L^{1}(\mathbf{R}))
\end{equation*}
故得 $\chi(f)\phi(y) = \chi(f_{y})$.这里固定满足 $\chi(f) \neq 0$ 的 $f$.当 $y \to 0$ 时,$\|f_{y} - f\|_{1} \to 0$,故由 $\phi(y) = \chi(f_{y})/\chi(f)$(a.e.\ $y$),$\phi$ 可视为连续的.把 $y$ 换成 $x + y$,则
\begin{equation*}
\chi(f)\phi(x + y) = \chi(f_{x + y}) = \chi((f_{x})_{y}) = \chi(f_{x})\phi(y) = \chi(f)\phi(x)\phi(y)
\end{equation*}
故 $\phi(x + y) = \phi(x)\phi(y)$,$|\phi(x)| \leq 1$.因此 $\phi(0) = \phi(0)^{2}$.由此,$\phi(0) = 1$.从而由 $1 = \phi(0) = \phi(x)\phi(-x)$,$|\phi(x)| \leq 1$,$|\phi(x)| = 1$($x \in \mathbf{R}$).此时,存在 $\delta > 0$ 使
\begin{equation*}
\int_{0}^{\delta}\phi(x)\frac{\dd x}{2\pi} = \alpha \neq 0.
\end{equation*}
又,
\begin{equation*}
\alpha\phi(x) = \int_{0}^{\delta}\phi(y)\phi(x)\frac{\dd y}{2\pi} = \int_{0}^{\delta}\phi(x + y)\frac{\dd y}{2\pi} = \int_{x}^{x + \delta}\phi(y)\frac{\dd y}{2\pi}
\end{equation*}
故 $\phi(x)$ 可微.因此,利用
\begin{equation*}
\frac{\phi(x + h) - \phi(x)}{h} = \frac{\phi(x)\phi(h) - \phi(x)}{h} = \phi(x)\frac{\phi(h) - 1}{h}
\end{equation*}
得 $\phi'(x) = \phi'(0)\phi(x)$,$|\phi(x)| = 1$,$\phi(0) = 1$.从而存在 $\xi \in \mathbf{R}$ 使 $\phi(x) = e^{-i\xi x}$.因此
\begin{equation*}
\chi(f) = \int f(x)e^{-i\xi x}\frac{\dd x}{2\pi} = \hat{f}(\xi).
\end{equation*}
综上,$\mathbf{R} \ni \xi \mapsto \chi_{\xi} \in \Delta(L^{1}(\mathbf{R}))$ 是满射.

反之,取 $\xi_{0} \in \mathbf{R}$,$\varepsilon > 0$,设 $|\xi - \xi_{0}| < \varepsilon$.此时,取 $f \in L^{1}(\mathbf{R})$ 使 $|\hat{f}(\xi)| < |\hat{f}(\xi_{0})| = 1$($\xi \in \mathbf{R}$,$\xi \neq \xi_{0}$).这里,取适当的 $\delta > 0$ 考虑 $V = \{\chi_{\xi}:\ |\chi_{\xi}(f) - \chi_{\xi_{0}}(f)| < \delta\}$,则 $V$ 是 $\Delta(L^{1}(\mathbf{R}))$ 的开集,若 $\chi_{\xi} \in V$,可令 $|\xi - \xi_{0}| < \varepsilon$.因此 $\xi \mapsto \chi_{\xi}$ 的逆映射也连续.因此 $\mathbf{R}$ 与 $\Delta(L^{1}(\mathbf{R}))$ 同胚.从而 $f$ 的 Gelfand 变换是 $f$ 的 Fourier 变换.

(3), (4) 由 (1), (2) 容易证明.
\end{proof}

\subsection{\texorpdfstring{测度构成的交换 Banach 代数(测度代数)$M(\mathbf{T})$}{测度构成的交换 Banach 代数(测度代数)M(T)}}\label{sec:3.2.3}

这里,叙述 $\mathbf{T}$ 上有界正则 Borel 测度构成的交换 Banach 代数 $M(\mathbf{T})$ 的若干性质.$M(\mathbf{T})$ 作为交换 Banach 代数具有复杂的结构,有许多研究(\cite{ch3:2}, \cite{ch3:12}, \cite{ch3:13} 等).

定义
\begin{equation*}
L^{1}(\mathbf{T}) = \left\{f:\ \|f\|_{1}\ \bigl(= \|f\|_{L^{1}(\mathbf{T})}\bigr) = \frac{1}{2\pi}\int_{0}^{2\pi}|f(x)|\dd x < \infty\right\}.
\end{equation*}

\begin{definition}\label{def:3.2.6}

(1) 用 $M(\mathbf{T})$ 表示 $\mathbf{T}$ 上有界正则复值 Borel 测度的全体.设 $\mathcal{B}$ 为 $\mathbf{T}$ 上 Borel 集的全体,对 $\mu \in M(\mathbf{T})$,令
\begin{align*}
& |\mu|(E) = \\
& \sup\Bigl\{\sum_{j=1}^{n}|\mu(E_{j})|:\ E = \bigcup_{j=1}^{n}E_{j},\ E_{j} \in \mathcal{B},\ \{E_{j}\}_{j} \text{ 互不相交},\ n \in \mathbf{N}\Bigr\}, \\
& \|\mu\| = \|\mu\|_{M(\mathbf{T})} = |\mu|(\mathbf{T})
\end{align*}
又,当 $\alpha, \beta \in \mathbf{C}$,$\mu, \nu \in M(\mathbf{T})$ 时,定义
\begin{equation*}
\begin{aligned}
(\alpha\mu + \beta\nu)(E) &= \alpha \cdot \mu(E) + \beta \cdot \nu(E), \\
\mu * \nu(E) &= \int \mu(E - x)\dd\nu(x) \quad (E \in \mathcal{B})
\end{aligned}
\end{equation*}
则 $\alpha\mu + \beta\nu$,$\mu * \nu \in M(\mathbf{T})$.此时,$M(\mathbf{T})$ 关于范数 $\|\cdot\|$,积 $*$,和,数乘成为交换 Banach 代数.进一步,令 $\delta_{0}(E) = 1$($0 \in E$),$0$($0 \notin E$),则 $\delta_{0}$ 是 $M(\mathbf{T})$ 的单位元(称 $\delta_{0}$ 为 Dirac 测度).因此 $M(\mathbf{T})$ 是具有单位元的交换 Banach 代数.称其为 $\mathbf{T}$ 上的测度代数.

(2) 对 $\mu \in M(\mathbf{T})$,定义 Fourier-Stieltjes 变换为
\begin{equation*}
\hat{\mu}(n) = \int_{0}^{2\pi}e^{-inx}\dd\mu(x) \quad (n \in \mathbf{Z}).
\end{equation*}
\end{definition}

对 $\mu \in M(\mathbf{T})$,令 $T_{\mu}f = \int f(x)\dd\mu(x)$,则 $T_{\mu}$ 是 $C(\mathbf{T})$ 上的连续线性泛函.此时,取 $C(\mathbf{T})$ 上连续线性泛函的空间 $C(\mathbf{T})^{*}$,由 Riesz 表示定理,通过 $\mu \in M(\mathbf{T})$ 到 $T_{\mu} \in C(\mathbf{T})^{*}$ 的对应,$M(\mathbf{T})$ 与 $C(\mathbf{T})^{*}$ 作为 Banach 空间等同.以后,Riesz 表示定理视为已知(cf.\ \cite{ch3:9}, \cite{ch3:11}).

\begin{lemma}\label{lem:3.2.7}

(1) $L^{1}(\mathbf{T})$ 是 $M(\mathbf{T})$ 的闭理想.

(2) $L^{1}(\mathbf{T})$ 的极大理想空间是 $\mathbf{Z}$,且可视为 $M(\mathbf{T})$ 的极大理想空间 $\Delta(M(\mathbf{T}))$ 的开集.
\end{lemma}

\begin{proof}
(1) 对 $f \in L^{1}(\mathbf{T})$,令 $\mu_{f}(E) = \int_{E}f(x)\frac{\dd x}{2\pi}$,则 $\mu_{f} \in M(\mathbf{T})$ 且 $\|\mu_{f}\| = \|f\|_{1}$,故 $L^{1}(\mathbf{T})$ 是 $M(\mathbf{T})$ 的闭子空间.又,设 $\mathbf{T}$ 上的 Lebesgue 测度为 $dm(x) = \frac{\dd x}{2\pi}$,设 $m(E) = 0$.此时,若 $f \in L^{1}(\mathbf{T})$,$\mu \in M(\mathbf{T})$,则
\begin{equation*}
\mu_{f} * \mu(E) = \int_{0}^{2\pi}\mu_{f}(E - x)\dd\mu(x) = 0
\end{equation*}
故得 $f * \mu = \mu_{f} * \mu \in L^{1}(\mathbf{T})$.因此 $L^{1}(\mathbf{T})$ 是 $M(\mathbf{T})$ 的理想.

(2) $L^{1}(\mathbf{T})$ 的极大理想空间可与定理~\ref{thm:3.2.5} 同样地证明与 $\mathbf{Z}$ 等同(利用 $\phi$ 具有 $2\pi$ 周期性).因此 $\Delta(L^{1}(\mathbf{T})) = \mathbf{Z}$.由 $L^{1}(\mathbf{T})$ 是 $M(\mathbf{T})$ 的理想,在 $\Delta(M(\mathbf{T}))$ 中 $\mathbf{Z}$ 是开集.
\end{proof}

\begin{theorem}[Williamson \cite{ch3:55}]\label{thm:3.2.8}
存在 $\mu \in M(\mathbf{T})$ 使 $\Sp(\mu, M(\mathbf{T}))$ 真包含 $\hat{\mu}(\mathbf{Z})$ 的闭包.由此,$\overline{\mathbf{Z}}$ 是 $\Delta(M(\mathbf{T}))$ 的真子集.
\end{theorem}

\begin{proof}
由后述的注意~\ref{note:3.3.12},存在 $\mu \in M(\mathbf{T})$ 使得 $\mu(\mathbf{T}) = 1$,$\mu \geq 0$(称这样的 $\mu$ 为概率测度),$\hat{\mu}$ 取实数值,且任意一点的测度为 $0$.进一步,此测度可取为:当令 $\mu^{n} = \mu^{n-1} * \mu$($n > 1$)时,$\delta_{0}, \mu, \mu^{2}, \mu^{3}, \ldots$ 具有落在互不相交的集合上的测度.此时,对 $\lambda \in \mathbf{C}$,$|\lambda| = 1$,
\begin{align*}
& \|(\lambda\delta_{0} + \mu)^{n}\| = \left\|\sum_{k=0}^{n}\binom{n}{k}\lambda^{k}\mu^{n-k}\right\| \\
& \qquad = \sum_{k=0}^{n}\binom{n}{k}|\lambda|^{k}\|\mu\|^{n-k} \\
& \qquad = \sum_{k=0}^{n}\binom{n}{k} = 2^{n}
\end{align*}
成立.因此 $\lim_{n \to \infty}\|(\lambda\delta_{0} + \mu)^{n}\|^{\frac{1}{n}} = 2$,故存在 $\chi \in \Delta(M(\mathbf{T}))$ 使 $|\chi(\lambda\delta_{0} + \mu)| = 2$.从而得 $\chi(\mu) = \lambda$.因此,按 $\mu$ 的取法,$\overline{\hat{\mu}(\mathbf{Z})} \subsetneq \Sp(\mu, M(\mathbf{T}))$.由此与补题~\ref{lem:3.2.7} 得 $\overline{\mathbf{Z}} \subsetneq \Delta(M(\mathbf{T}))$.
\end{proof}

下面,用定理~\ref{thm:3.2.8} 的 $\mu$,令 $\nu = \mu^{2} + \delta_{0}$,即可容易得到.

\begin{corollary}\label{cor:3.2.9}
存在满足 $\hat{\nu} \geq 1$ 的 $\nu \in M(\mathbf{T})$,且 $1/\hat{\nu}$ 不是 Fourier-Stieltjes 变换.
\end{corollary}

这些结果本质上是由 Wiener-Pitt \cite{ch3:54} 于 1938 年得到的,把这种现像称为 Wiener-Pitt 现象.在这个方向上,F.\ Parreau 得到了更精密的结果,是对下述 Katznelson 问题 \cite{ch3:6} 的否定解决.

\textbf{Katznelson 问题}．对 $\mu \in M(\mathbf{T})$,若 $\Sp(\mu, M(\mathbf{T}))$ 是实数的子集,则 $\hat{\mu}(\mathbf{Z})$ 在 $\Sp(\mu, M(\mathbf{T}))$ 中稠密吗?

\begin{theorem}[Parreau \cite{ch3:36}]\label{thm:3.2.10}
存在 $\mu \in M(\mathbf{T})$ 满足:
\begin{equation*}
\mathbf{R} \supset \Sp(\mu, M(\mathbf{T})) \supsetneqq \overline{\hat{\mu}(\mathbf{Z})}.
\end{equation*}
\end{theorem}

\subsection{\texorpdfstring{交换 Banach 代数 $\mathfrak{M}_p(\mathbf{T})$}{交换 Banach 代数 M_p(T)}}\label{sec:3.2.4}

设 $1 \leq p < \infty$.令 $L^{p}(\mathbf{T}) = \{f:\ \|f\|_{p} = \left(\frac{1}{2\pi}\int_{0}^{2\pi}|f(x)|^{p}\dd x\right)^{1/p} < \infty\}$,则 $L^{p}(\mathbf{T})$ 由通常的运算成为 Banach 空间.对 $\mathbf{T}$ 上的函数 $f(x)$,令 $\tau_{y}f(x) = f(x - y)$,称满足 $T\tau_{y} = \tau_{y}T$($y \in \mathbf{T}$)的 $L^{p}(\mathbf{T})$ 上的连续线性算子 $T$ 为乘子算子,其全体记为 $\mathfrak{M}_{p}(\mathbf{T})$.此时,$\mathfrak{M}_{p}(\mathbf{T})$ 以算子范数 $\|T\|_{\mathfrak{M}_{p}(\mathbf{T})} = \sup\{\|Tf\|_{p}:\ \|f\|_{p} \leq 1\}$ 与自然的运算成为 Banach 空间.又,利用对偶性可证明 $T(f * g) = Tf * g$($f, g \in C(\mathbf{T})$).进一步,具有下列性质(cf.\ \cite{ch3:9}).

\begin{lemma}\label{lem:3.2.11}
设 $1 \leq p \leq 2$.

(1) 当 $T \in \mathfrak{M}_{p}(\mathbf{T})$ 时,存在 $\hat{T} \in \ell^{\infty}$ 使 $(Tf)^{\wedge}(n) = \hat{T}(n)\hat{f}(n)$($f \in C(\mathbf{T})$).这里,$\ell^{\infty}$ 表示由 $\mathbf{Z}$ 上的有界数列 $\{a_{n}\}$ 构成的,以 $\|\{a_{n}\}\|_{\ell^{\infty}} = \sup_{n}|a_{n}|$ 为范数的 Banach 空间.又,记 $M_{p}(\mathbf{Z}) = \{\hat{T}:\ T \in \mathfrak{M}_{p}(\mathbf{T})\}$,称 $M_{p}(\mathbf{Z})$ 的元为乘子.$M_{p}(\mathbf{Z})$ 的范数由 $\mathfrak{M}_{p}(\mathbf{T})$ 自然引入.

(2) $\mathfrak{M}_{1}(\mathbf{T})$ 作为 Banach 空间与 $M(\mathbf{T})$ 等同.

(3) $M_{2}(\mathbf{Z})$ 作为 Banach 空间与 $\ell^{\infty}$ 等同.
\end{lemma}

\begin{proof}
(1) 当 $f(x) = e^{-inx}$ 时,令 $\hat{T}(n) = \int_{0}^{2\pi}Tf(x)e^{-inx}\frac{\dd x}{2\pi}$,由 Hölder 不等式,$|\hat{T}(n)| \leq \|T\|_{\mathfrak{M}_{p}}$.此时,$T(e^{-inx}) = \hat{T}(n)e^{inx}$,故对 $f \in C(\mathbf{T})$,
\begin{equation*}
\widehat{Tf}(n) = \hat{T}(n)\hat{f}(n) \quad (n \in \mathbf{Z}).
\end{equation*}
又,由三角多项式全体在 $L^{p}(\mathbf{T})$ 中稠密,$\{\hat{T}(n)\}$ 唯一确定.

(2) 若 $\mu \in M(\mathbf{T})$,由 $T_{\mu}f = \mu * f$,$T_{\mu}f \in L^{1}(\mathbf{T})$($f \in L^{1}(\mathbf{T})$).又,由 Fubini 定理,$\|T_{\mu}f\|_{1} = \int\left|\int f(x - y)\dd\mu(y)\right|\frac{\dd x}{2\pi} \leq \|\mu\|\|f\|_{1}$,得 $\|T_{\mu}\|_{\mathfrak{M}_{1}} \leq \|\mu\|$.反之,当 $T \in \mathfrak{M}_{1}(\mathbf{T})$ 时,对 $\mathbf{T}$ 上的三角多项式 $g$,对应 $\sum_{n}\hat{g}(n)\hat{T}(n) = \int T(e^{-inx})g(-x)\frac{\dd x}{2\pi}$,由 (1),$\left|\sum_{n}\hat{g}(n)\hat{T}(n)\right| \leq \|T\|_{\mathfrak{M}_{1}}\|g\|_{\infty}$,故由 Riesz 表示定理,存在 $\mu \in M(\mathbf{T})$ 使 $\int g(-x)\dd\mu(x) = \sum_{n}\hat{g}(n)\hat{T}(n)$.特别地,当 $g(x) = e^{inx}$ 时,$\hat{\mu}(n) = \hat{T}(n)$,故可写为 $T = T_{\mu}$.因此 $\|\mu\| \leq \|T\|_{\mathfrak{M}_{1}}$.综上,$\|\mu\| = \|T_{\mu}\|_{\mathfrak{M}_{1}}$,$\mathfrak{M}_{1}(\mathbf{T})$ 作为 Banach 空间与 $M(\mathbf{T})$ 等同.

(3) 若 $\{a_{n}\} \in \ell^{\infty}$,则当 $f \in L^{2}(\mathbf{T})$ 时,由 Parseval 等式,
\begin{equation*}
\sum\left|a_{n}\hat{f}(n)\right|^{2} \leq \|\{a_{n}\}\|_{\ell^{\infty}}^{2}\sum|\hat{f}(n)|^{2} = \|\{a_{n}\}\|_{\ell^{\infty}}^{2}\|f\|_{2}^{2} < \infty
\end{equation*}
故存在 $g \in L^{2}(\mathbf{T})$ 使 $\hat{g}(n) = a_{n}\hat{f}(n)$($n \in \mathbf{Z}$).因此,由 $\widehat{Tf}(n) = a_{n}\hat{f}(n)$ 确定 $T \in \mathfrak{M}_{2}(\mathbf{T})$,且可知 $\|T\|_{\mathfrak{M}_{2}} \leq \|\{a_{n}\}\|_{\ell^{\infty}}$.反之,若 $T \in \mathfrak{M}_{2}(\mathbf{T})$,像 (2) 那样令 $\hat{T}(n) = \int T(e^{-inx})e^{-inx}\frac{\dd x}{2\pi}$,由 Schwarz 不等式可证明 $\hat{T} \in \ell^{\infty}$ 且 $\|\{\hat{T}(n)\}\|_{\ell^{\infty}} \leq \|T\|_{\mathfrak{M}_{2}}$.
\end{proof}

下面的证明由补题~\ref{lem:3.2.11} 容易得到.

\begin{theorem}\label{thm:3.2.12}
设 $1 \leq p \leq 2$.$\mathfrak{M}_{p}(\mathbf{T})$ 关于算子的复合,成为以 $I = \delta_{0}$ 为单位元的交换 Banach 代数.
\end{theorem}

\subsection{奇异积分算子构成的交换 Banach 代数}\label{sec:3.2.5}

这里,对 $n \geq 2$,只介绍奇异积分算子全体构成交换 Banach 代数这一 Calderón-Zygmund \cite{ch3:14} 的结果.证明等细节请参照 \cite{ch3:14}.

\begin{definition}\label{def:3.2.13}
$\Omega(x)$ 是 $\mathbf{R}^{n} \setminus \{0\}$ 上的函数,设 $S^{n-1}$ 为 $\mathbf{R}^{n}$ 的单位球面,对 $1 < q < \infty$ 满足下列性质:

(1) $\Omega(rx) = \Omega(x)$($x \in S^{n-1}$,$r > 0$)

(2) $\Omega \in L^{q}(S^{n-1})$

(3) $\int_{S^{n-1}}\Omega(x')\dd\sigma(x') = 0$.其中 $\dd\sigma$ 表示 $S^{n-1}$ 上的面积元素,设 $L^{q}(S^{n-1}) = \{f:\ \|f\|_{L^{q}(S^{n-1})} = \left(\int_{S^{n-1}}|f(x')|^{q}\dd\sigma(x')\right)^{1/q} < \infty\}$.
\end{definition}

令 $K(x) = \Omega(x)/|x|^{n}$($x \neq 0$)时,奇异积分算子
\begin{equation*}
K * f(x) = \lim_{\varepsilon \to 0}\int_{|x - y| > \varepsilon}K(x - y)f(y)\dd y
\end{equation*}
是 $L^{p}(\mathbf{R}^{n})$($1 < p < \infty$)上的有界线性算子,其 Fourier 变换为
\begin{equation*}
\mathcal{F}(K * f) = \mathcal{F}(K)\mathcal{F}(f) \quad \left(f \in C_{c}^{\infty}(\mathbf{R}^{n})\right),
\end{equation*}
且已知在几乎所有点上 $|\mathcal{F}(K)| \leq C\|K\|_{L^{q}(S^{n-1})}$ 成立(cf.\ \cite{ch3:5}).其中 $C_{c}^{\infty}(\mathbf{R}^{n})$ 是具有紧支集的无限次连续可偏导函数的全体.把上述 $K$ 的全体记为 $\mathcal{K}$.此时,对 $1 < q < \infty$,定义下列奇异积分算子的族
\begin{equation*}
A_{q} = \{T:\ Tf = \alpha f + K * f,\ \alpha \in \mathbf{C},\ K \in \mathcal{K}\}.
\end{equation*}
进一步,对 $T \in A_{q}$,$Tf = \alpha f + K * f$,令
\begin{equation*}
|T|_{q} = |\alpha| + \|K\|_{L^{q}(S^{n-1})}.
\end{equation*}
此时,下式成立:

\begin{theorem}[Calderón-Zygmund]\label{thm:3.2.14}
$A_{q}$ 关于算子的通常运算与范数 $|\cdot|_{q}$,成为具有单位元的交换 Banach 代数.
\end{theorem}

又,关于 $A_{q}$ 的极大理想空间,已知下列结论.

\begin{theorem}\label{thm:3.2.15}
具有单位元的交换 Banach 代数 $A_{q}$ 的极大理想空间 $\mathfrak{M}(A_{q})$ 与 $S^{n-1}$ 等同.
\end{theorem}

% ============ 第3章 §3.3 A(T)中的谱综合(原书 pp.237-249) ============
\section{\texorpdfstring{$A(\mathbf{T})$ 中的谱综合}{A(T) 中的谱综合}}\label{sec:3.3}
\index{Malliavin 定理}
\index{Kronecker 集}
\index{Kronecker 定理}
\index{Cantor 集}
\index{在 Q 上线性无关}


这里,首先叙述谱综合集的若干例子,后半部分叙述证明 $A(\mathbf{T})$ 谱综合不可能的 Malliavin 的结果 \cite{ch3:34}.其后,叙述作为稀疏集(thin set)的 Kronecker 集.

\subsection{谱综合集的例子}\label{sec:3.3.1}

\begin{theorem}\label{thm:3.3.1}
$\mathbf{T}$ 上的一点是 $A(\mathbf{T})$ 的谱综合集.
\end{theorem}

\begin{proof}
考虑平移,只需证明 $\{0\} \subset \mathbf{T} = [-\pi, \pi)$ 是谱综合集.因为,设 $\{0\}$ 是谱综合集.取 $x_{0} \in \mathbf{T}$,$f \in A(\mathbf{T})$,$f(x_{0}) = 0$.令 $f_{-x_{0}}(x) = f(x + x_{0})$,则 $f_{-x_{0}} \in A(\mathbf{T})$,$f_{-x_{0}}(0) = f(x_{0}) = 0$.因 $\{0\}$ 是谱综合集,对任意 $\varepsilon > 0$ 存在 $g \in A(\mathbf{T})$ 使得 $g$ 在原点某邻域内为 $0$ 且 $\|f_{-x_{0}} - g\|_{A(\mathbf{T})} < \varepsilon$.令 $h(x) = g(x - x_{0})$,则 $h \in A(\mathbf{T})$,$h$ 在 $x_{0}$ 的某邻域内为 $0$.这里,由 $f(x + x_{0}) - h(x + x_{0}) = f_{-x_{0}}(x) - g(x)$,$\|f - h\|_{A(\mathbf{T})} = \|f_{-x_{0}} - g\|_{A(\mathbf{T})} < \varepsilon$.于是下面证明 $\{0\}$ 是 $A(\mathbf{T})$ 的谱综合集.设 $f(0) = 0$,$f \in A(\mathbf{T})$.现在,令 $0 < \varepsilon < 1/2$,
\begin{equation*}
\Delta_{\varepsilon}(x) = \max\left\{0, 1 - \frac{|x|}{\varepsilon}\right\}, \quad V_{\varepsilon} = 2\Delta_{2\varepsilon} - \Delta_{\varepsilon},
\end{equation*}
则
\begin{equation*}
\lim_{\varepsilon \to 0}\|V_{\varepsilon}f\|_{A(\mathbf{T})} = 0.
\end{equation*}
事实上,对连续可微函数 $f \in C^{1}(\mathbf{T})$,
\begin{equation*}
\|V_{\varepsilon}f\|_{A(\mathbf{T})} = \sum_{n \in \mathbf{Z}}\left|\widehat{V_{\varepsilon}f}(n)\right| = \left|\widehat{V_{\varepsilon}f}(0)\right| + \sum_{n \neq 0}\left|\widehat{V_{\varepsilon}f}(n)\right|.
\end{equation*}
这里考虑
\begin{equation*}
\left|\widehat{V_{\varepsilon}f}(0)\right| = \left|\frac{1}{2\pi}\int_{-\pi}^{\pi}V_{\varepsilon}(x)f(x)\dd x\right| = \left|\frac{1}{2\pi}\int_{-2\varepsilon}^{2\varepsilon}V_{\varepsilon}(x)f(x)\dd x\right|.
\end{equation*}
由
\begin{equation*}
\int_{-2\varepsilon}^{2\varepsilon}V_{\varepsilon}(x)f(x)\frac{\dd x}{2\pi} = \left[x(V_{\varepsilon}f)(x)\frac{1}{2\pi}\right]_{-2\varepsilon}^{2\varepsilon} - \frac{1}{2\pi}\int_{-2\varepsilon}^{2\varepsilon}x(V_{\varepsilon}f)'(x)\dd x
\end{equation*}
对右边用 Schwarz 不等式,则
\begin{align*}
\left|\int_{-2\varepsilon}^{2\varepsilon}V_{\varepsilon}(x)f(x)\frac{\dd x}{2\pi}\right| & \leq \left|\int_{-2\varepsilon}^{2\varepsilon}x(V_{\varepsilon}f)'(x)\frac{\dd x}{2\pi}\right| \\
& \leq \left(\int_{-2\varepsilon}^{2\varepsilon}x^{2}\frac{\dd x}{2\pi}\right)^{1/2}\left(\int_{-\pi}^{\pi}\left|(V_{\varepsilon}f)'(x)\right|^{2}\frac{\dd x}{2\pi}\right)^{1/2} \\
& = C\varepsilon^{3/2}\left(\int_{-\pi}^{\pi}\left|(V_{\varepsilon}f)'(x)\right|^{2}\frac{\dd x}{2\pi}\right)^{1/2}
\end{align*}
其次,当 $n \neq 0$ 时,
\begin{align*}
\widehat{V_{\varepsilon}f}(n) & = \frac{1}{2\pi}\int_{-\pi}^{\pi}(V_{\varepsilon}f)(x)e^{-inx}\dd x \\
& = \left[-\frac{e^{-inx}}{2\pi in}(V_{\varepsilon}f)(x)\right]_{-\pi}^{\pi} + \frac{1}{2\pi in}\int_{-\pi}^{\pi}(V_{\varepsilon}f)'(x)e^{-inx}\dd x
\end{align*}
故
\begin{equation*}
\widehat{V_{\varepsilon}f}(n) = \frac{1}{2\pi in}\int_{-\pi}^{\pi}(V_{\varepsilon}f)'(x)e^{inx}\dd x.
\end{equation*}
又,
\begin{align*}
\sum_{n \neq 0}\left|\widehat{V_{\varepsilon}f}(n)\right| & = \sum_{n \neq 0}\frac{1}{n}\left|\int_{-\pi}^{\pi}(V_{\varepsilon}f)'(x)e^{-inx}\frac{\dd x}{2\pi}\right| \\
& \leq \left(\sum_{n \neq 0}\frac{1}{n^{2}}\right)^{1/2}\left(\sum_{n \neq 0}\left|\widehat{(V_{\varepsilon}f)'}(n)\right|^{2}\right)^{1/2} \\
& \leq C\left(\frac{1}{2\pi}\int_{-\pi}^{\pi}\left|(V_{\varepsilon}f)'(x)\right|^{2}\dd x\right)^{1/2}.
\end{align*}
这里,
\begin{equation*}
\left(\int_{-\pi}^{\pi}\left|(V_{\varepsilon}f)'(x)\right|^{2}\frac{\dd x}{2\pi}\right)^{1/2} \leq C\varepsilon^{1/2}.
\end{equation*}
事实上,在几乎所有点上 $(V_{\varepsilon}f)'(x) = V_{\varepsilon}'(x)f(x) + V_{\varepsilon}(x)f'(x)$,由
\begin{equation*}
V_{\varepsilon}'(x) = \left\{\begin{array}{l l}
\frac{1}{2\varepsilon} & (-2\varepsilon < x < -\varepsilon) \\
-\frac{1}{2\varepsilon} & (\varepsilon < x < 2\varepsilon) \\
0 & (|x| \leq \varepsilon,\ |x| \geq 2\varepsilon)
\end{array}\right.
\end{equation*}
与 $f(0) = 0$ 即得.综上,
\begin{equation*}
\lim_{\varepsilon \to 0}\|V_{\varepsilon}f\|_{A(\mathbf{T})} = 0.
\end{equation*}
对一般的 $f \in A(\mathbf{T})$,由连续可微函数全体在 $A(\mathbf{T})$ 中稠密与 $\|V_{\varepsilon}\|_{A(\mathbf{T})} \leq 3$ 即得.$\|V_{\varepsilon}\|_{A(\mathbf{T})} \leq 3$ 由
\begin{equation*}
\widehat{\Delta_{\varepsilon}}(n) = \frac{2}{\pi} \cdot \frac{\sin^{2}\left(\frac{\varepsilon n}{2}\right)}{\varepsilon n^{2}} \geq 0, \quad \|\Delta_{\varepsilon}\|_{A(\mathbf{T})} = \Delta_{\varepsilon}(0) = 1
\end{equation*}
可知.于是设 $f \in A(\mathbf{T})$,$f(0) = 0$.令 $g_{\varepsilon} = (1 - V_{\varepsilon})f$,则 $g_{\varepsilon}$ 在原点某邻域内为 $0$,且 $g_{\varepsilon} \in A(\mathbf{T})$.又,
\begin{equation*}
\|g_{\varepsilon} - f\|_{A(\mathbf{T})} = \|V_{\varepsilon}f\|_{A(\mathbf{T})} \to 0 \quad (\varepsilon \to 0)
\end{equation*}
故
\begin{equation*}
\lim_{\varepsilon \to 0}\|g_{\varepsilon} - f\|_{A(\mathbf{T})} = 0.
\end{equation*}
因此证明了 $\{0\}$ 是 $A(\mathbf{T})$ 的谱综合集.
\end{proof}

\begin{corollary}\label{cor:3.3.2}
闭区间 $[a, b] \subset \mathbf{T}$ 是 $A(\mathbf{T})$ 的谱综合集.
\end{corollary}

\begin{proof}
不妨设 $[a, b] \subset (-\pi, \pi)$.令 $V_{\varepsilon}^{a}(x) = V_{\varepsilon}(x - a)$,$V_{\varepsilon}^{b}(x) = V_{\varepsilon}(x - b)$.设 $f \in A(\mathbf{T})$ 满足 $f(x) = 0$($x \in [a, b]$).令 $g_{\varepsilon} = (1 - V_{\varepsilon}^{a})(1 - V_{\varepsilon}^{b})f$,则 $g_{\varepsilon} \in A(\mathbf{T})$ 在 $[a, b]$ 的某邻域内为 $0$.又
\begin{align*}
& \|f - g_{\varepsilon}\|_{A(\mathbf{T})} \leq \|f - (1 - V_{\varepsilon}^{a})f\|_{A(\mathbf{T})} \\
& \qquad + \|(1 - V_{\varepsilon}^{a})f - (1 - V_{\varepsilon}^{a})(1 - V_{\varepsilon}^{b})f\|_{A(\mathbf{T})} \\
& \qquad \leq \|V_{\varepsilon}^{a}f\|_{A(\mathbf{T})} + \|1 - V_{\varepsilon}^{a}\|_{A(\mathbf{T})}\|V_{\varepsilon}^{b}f\|_{A(\mathbf{T})} \\
& \qquad \leq \|V_{\varepsilon}(f_{-a})\|_{A(\mathbf{T})} + 4\|V_{\varepsilon}(f_{-b})\|_{A(\mathbf{T})} \to 0 \quad (\varepsilon \to 0)
\end{align*}
故 $\|f - g_{\varepsilon}\|_{A(\mathbf{T})} \to 0$($\varepsilon \to 0$).
\end{proof}

\subsection{Malliavin 定理}\label{sec:3.3.2}

这里叙述交换 Banach 代数 $A(\mathbf{T})$ 谱综合不可能.今天,这一结果被称为 Malliavin 定理(cf.\ \cite{ch3:6}, \cite{ch3:7}, \cite{ch3:8}, \cite{ch3:11}, \cite{ch3:34}, \cite{ch3:39}).

\begin{theorem}[Malliavin]\label{thm:3.3.3}
$A(\mathbf{T})$ 谱综合不可能.
\end{theorem}

下面叙述的补题~\ref{lem:3.3.4},补题~\ref{lem:3.3.5} 证明上述定理.

\begin{lemma}\label{lem:3.3.4}
设实值函数 $f \in A(\mathbf{T})$ 具有下列性质:$t^{p}\omega(t) \in L^{1}(\mathbf{R})$($p = 0, 1, \ldots$),其中
\begin{equation*}
\omega(t) = \sup_{\|g\|_{A(\mathbf{T})} \leq 1}\left|\int e^{itf(x)}g(x)\frac{\dd x}{2\pi}\right|.
\end{equation*}
此时,存在 $a \in \mathbf{R}$ 使得由 $(f - a)^{p}$($p = 0, 1, \ldots$)生成的 $A(\mathbf{T})$ 的闭理想 $I_{p}$($p = 0, 1, \ldots$)全都不同.
\end{lemma}

\begin{proof}
令 $K = f(\mathbf{T})$,固定 $\tau \in C(\mathbf{T})$.考虑从 $\varphi \in C(K)$ 到 $\int \varphi(f(x))\tau(x)\frac{\dd x}{2\pi}$ 的复值函数,它是线性的,且
\begin{equation*}
\left|\int \varphi(f(x))\tau(x)\frac{\dd x}{2\pi}\right| \leq \|\varphi\|_{\infty}\|\tau\|_{\infty}\|\mu\|
\end{equation*}
故这是连续线性泛函.因此由 Riesz 表示定理,存在 $\lambda_{\tau} \in M(K)$($K$ 上有界正则 Borel 测度的全体)使得
\begin{equation*}
\int \varphi(f(x))\tau(x)\frac{\dd x}{2\pi} = \int_{K}\varphi(t)\dd\lambda_{\tau}(t).
\end{equation*}
这里,考察由 $\tau \in C(\mathbf{T})$ 到 $\lambda_{\tau} \in M(K)$ 的映射,线性是显然的.进一步,当 $\psi \in C(K)$ 时,$d\lambda_{\psi(f)\tau} = \psi(x)\dd\lambda_{\tau}$ 成立.事实上,
\begin{equation*}
\int_{K}\varphi(t)\psi(t)\dd\lambda_{\tau}(t) = \int \varphi(f(x))\psi(f(x))\tau(x)\frac{\dd x}{2\pi} = \int_{K}\varphi(t)\dd\lambda_{\psi(f)\tau}(t).
\end{equation*}
于是令 $\varphi(t) = e^{ist}$,$\tau = g(x)$,则
\begin{equation*}
\int e^{isf(x)}g(x)\frac{\dd x}{2\pi} = \int_{K}e^{ist}\dd\lambda_{g}(t).
\end{equation*}
此时,利用假设,
\begin{equation*}
\left|\int_{K}e^{ist}\dd\lambda_{g}\right| \leq \omega(s)\|g\|_{A(\mathbf{T})}, \quad s^{p}\omega(s) \in L^{1}(\mathbf{R}) \ (p = 0, 1, \ldots)
\end{equation*}
故可写为 $d\lambda_{g}(t) = m_{g}(t)\dd t$,$m_{g} \in C^{\infty}(\mathbf{R})$($\mathbf{R}$ 上无限次可微函数的全体).因此
\begin{equation*}
m_{g}(t) = \frac{1}{2\pi}\int_{-\infty}^{\infty}e^{its}\widehat{\lambda_{g}}(s)\dd s
\end{equation*}
得
\begin{align*}
\frac{\dd^{p}}{\dd t^{p}}m_{g}(t) & = \frac{1}{2\pi}\int_{-\infty}^{\infty}(is)^{p}e^{its}\widehat{\lambda_{g}}(s)\dd s \\
& = \frac{i^{p}}{2\pi}\int_{-\infty}^{\infty}s^{p}e^{its}\widehat{\lambda_{g}}(s)\dd s.
\end{align*}
从而
\begin{equation*}
\left|\frac{\dd^{p}}{\dd t^{p}}m_{g}(t)\right| \leq \frac{1}{2\pi}\int_{-\infty}^{\infty}|s|^{p}\omega(s)\dd s\,\|g\|_{A(\mathbf{T})}
\end{equation*}
得
\begin{equation*}
\left|\frac{\dd^{p}}{\dd t^{p}}m_{g}(t)\right| \leq C_{p}\|g\|_{A(\mathbf{T})}.
\end{equation*}
因此,对 $a \in \mathbf{R}$,
\begin{equation*}
\Phi_{p}:\ g \mapsto \left(\frac{\dd^{p}}{\dd t^{p}}m_{g}\right)(a)
\end{equation*}
是 $A(\mathbf{T})$ 上的连续线性泛函.于是令 $g \equiv 1$,则
\begin{equation*}
\int_{-\infty}^{\infty}m_{1}(t)\dd t = \widehat{\lambda_{1}}(0) = \int \frac{\dd x}{2\pi} = 1 \neq 0
\end{equation*}
故取 $a \in \mathbf{R}$ 使 $m_{1}(a) \neq 0$,则对 $p, q = 1, 2, \ldots$,由 $\lambda_{\tau}$ 的性质,
\begin{align*}
\Phi_{p}((f - a)^{q}) & = \left(\frac{\dd^{p}}{\dd t^{p}}m_{(f-a)^{q}}\right)(a) \\
& = \left(\frac{\dd^{p}}{\dd t^{p}}(t - a)^{q}m_{1}\right)(a).
\end{align*}
因此,若 $q > p$,则 $\Phi_{p}((f - a)^{q}) = 0$;若 $q = p$,则 $\Phi_{p}((f - a)^{p} \cdot 1) = p!\,m_{1}(a) \neq 0$,故对 $p = 1, 2, \ldots$,由 $(f - a)^{p}$ 生成的各个理想全都不同.
\end{proof}

其次,证明满足补题~\ref{lem:3.3.4} 假设的 $f$ 的存在性.

\begin{lemma}\label{lem:3.3.5}
存在实值函数 $f \in A(\mathbf{T})$ 与 $\delta > 0$ 满足
\begin{equation*}
\omega(t) = \sup_{\|g\|_{A(\mathbf{T})} \leq 1}\left|\int_{0}^{2\pi}e^{itf(x)}g(x)\dd x\right| = O\left(e^{-\delta|t|^{1/2}}\right).
\end{equation*}
由此可证明 $f$ 满足补题~\ref{lem:3.3.4}.
\end{lemma}

\begin{proof}
第 1 阶段.设 $0 \neq g_{\alpha}$,$0 \neq h_{\alpha} \in A(\mathbf{T})$($\alpha \in \mathbf{R}$)连续依赖于 $\alpha$.此时,对任意 $M > 0$,$\varepsilon > 0$,存在 $k \in \mathbf{N}$ 使得
\begin{equation*}
\|g_{\alpha}(x)h_{\alpha}(kx)\|_{A^{\infty}} \leq (1 + \varepsilon)\|g_{\alpha}\|_{A^{\infty}}\|h_{\alpha}\|_{A^{\infty}} \quad (|\alpha| \leq M).
\end{equation*}
其中 $\|f\|_{A^{\infty}} = \sup_{n \in \mathbf{Z}}\left|\hat{f}(n)\right|$.

第 1 阶段的证明.固定 $\alpha$,设 $k$ 为自然数.
\begin{equation*}
\widehat{g_{\alpha}(x)h_{\alpha}(kx)}(l) = \sum_{j}\widehat{h_{\alpha}}(j)\widehat{g_{\alpha}}(l - kj)
\end{equation*}
取 $l'$,$j'$ 满足 $l' = l - kj'$,$-k/2 \leq l' \leq (k/2) + 1$,则
\begin{equation*}
\left|\sum_{j}\widehat{g_{\alpha}}(l - kj)\widehat{h_{\alpha}}(j) - \widehat{g_{\alpha}}(l')\widehat{h_{\alpha}}(j')\right| \leq \sum_{|n| \geq \frac{k}{2} - 1}\left|\widehat{g_{\alpha}}(n)\right|\|h_{\alpha}\|_{A^{\infty}}.
\end{equation*}
因此
\begin{equation*}
\|g_{\alpha}(x)h_{\alpha}(kx)\|_{A^{\infty}} \leq \sum_{|n| \geq \frac{k}{2} - 1}\left|\widehat{g_{\alpha}}(n)\right|\|h_{\alpha}\|_{A^{\infty}} + \|g_{\alpha}\|_{A^{\infty}}\|h_{\alpha}\|_{A^{\infty}}
\end{equation*}
以及
\begin{equation*}
\left|\widehat{g_{\alpha}}(l')\widehat{h_{\alpha}}(j')\right| \leq \sum_{|n| \geq \frac{k}{2} - 1}\left|\widehat{g_{\alpha}}(n)\right|\|h_{\alpha}\|_{A^{\infty}} + \|g_{\alpha}(x)h_{\alpha}(kx)\|_{A^{\infty}}.
\end{equation*}
由后一式得
\begin{equation*}
\|g_{\alpha}\|_{A^{\infty}}\|h_{\alpha}\|_{A^{\infty}} \leq \sum_{|n| \geq \frac{k}{2} - 1}\left|\widehat{g_{\alpha}}(n)\right|\|h_{\alpha}\|_{A^{\infty}} + \|g_{\alpha}(x)h_{\alpha}(kx)\|_{A^{\infty}}.
\end{equation*}
因此,固定 $\alpha$ 时,对 $\varepsilon > 0$ 取充分大的 $k$,
\begin{equation*}
\left|\|g_{\alpha}\|_{A^{\infty}}\|h_{\alpha}\|_{A^{\infty}} - \|g_{\alpha}(x)h_{\alpha}(kx)\|_{A^{\infty}}\right| < \varepsilon.
\end{equation*}
这里,利用 $[-M, M]$ 是紧集,对 $\eta > 0$ 存在 $\delta > 0$ 与 $\alpha_{1}, \ldots, \alpha_{n} \in [-M, M]$,使得对任意 $|\alpha| \leq M$,存在满足 $|\alpha - \alpha_{j}| < \delta$ 的 $\alpha_{j}$,且 $\|g_{\alpha} - g_{\alpha_{j}}\|_{A(\mathbf{T})} < \eta$,$\|h_{\alpha} - h_{\alpha_{j}}\|_{A(\mathbf{T})} < \eta$.此时,取充分大的 $k$,
\begin{equation*}
\left|\|g_{\alpha_{j}}\|_{A^{\infty}}\|h_{\alpha_{j}}\|_{A^{\infty}} - \|g_{\alpha_{j}}(x)h_{\alpha_{j}}(kx)\|_{A^{\infty}}\right| < \eta \quad (j = 1, \ldots, n),
\end{equation*}
令
\begin{equation*}
C = \sup_{\alpha \in [-M, M]}\|g_{\alpha}\|_{A^{\infty}} < \infty, \quad D = \sup_{\alpha \in [-M, M]}\|h_{\alpha}\|_{A^{\infty}} < \infty,
\end{equation*}
则对 $\alpha \in [-M, M]$,存在 $\alpha_{j}$,
\begin{align*}
& \left|\|g_{\alpha}\|_{A^{\infty}}\|h_{\alpha}\|_{A^{\infty}} - \|g_{\alpha}(x)h_{\alpha}(kx)\|_{A^{\infty}}\right| \\
& \leq \left|\|g_{\alpha}\|_{A^{\infty}}\|h_{\alpha}\|_{A^{\infty}} - \|g_{\alpha_{j}}\|_{A^{\infty}}\|h_{\alpha_{j}}\|_{A^{\infty}}\right| \\
& \qquad + \left|\|g_{\alpha_{j}}\|_{A^{\infty}}\|h_{\alpha_{j}}\|_{A^{\infty}} - \|g_{\alpha_{j}}(x)h_{\alpha_{j}}(kx)\|_{A^{\infty}}\right| \\
& \qquad + \|g_{\alpha_{j}}(x)h_{\alpha_{j}}(kx) - g_{\alpha}(x)h_{\alpha_{j}}(kx)\|_{A^{\infty}} \\
& \qquad + \|g_{\alpha}(x)h_{\alpha_{j}}(kx) - g_{\alpha}(x)h_{\alpha}(kx)\|_{A^{\infty}} \\
& \leq \|g_{\alpha} - g_{\alpha_{j}}\|_{A(\mathbf{T})}\|h_{\alpha}\|_{A^{\infty}} + \|g_{\alpha_{j}}\|_{A^{\infty}}\|h_{\alpha} - h_{\alpha_{j}}\|_{A(\mathbf{T})} \\
& \qquad + \|g_{\alpha} - g_{\alpha_{j}}\|_{A(\mathbf{T})}\|h_{\alpha_{j}}\|_{A^{\infty}} + \|g_{\alpha}\|_{A^{\infty}}\|h_{\alpha} - h_{\alpha_{j}}\|_{A(\mathbf{T})} \\
& \leq 2(C + D)\eta.
\end{align*}
这里,把 $\eta > 0$ 取得足够小,使右边小于 $\varepsilon$ 即可.

第 2 阶段.对不含原点的闭区间 $[a, b]$,存在 $\rho$ 使得
\begin{equation*}
\left\|e^{i\alpha \cos kx}\right\|_{A^{\infty}} \leq \rho < 1 \quad (k \in \mathbf{Z},\ \alpha \in [a, b]).
\end{equation*}
由 $\|f(kx)\|_{A^{\infty}} = \|f(x)\|_{A^{\infty}}$,$\left\|e^{i\alpha \cos kx}\right\|_{A^{\infty}} = \|e^{i\alpha \cos x}\|_{A^{\infty}}$.又,令 $f_{\alpha}(x) = e^{i\alpha \cos x}$,利用 Parseval 等式,
\begin{equation*}
1 = \frac{1}{2\pi}\int_{0}^{2\pi}\left|e^{i\alpha \cos x}\right|^{2}\dd x = \sum_{n=-\infty}^{\infty}\left|\widehat{f_{\alpha}}(n)\right|^{2}
\end{equation*}
故对 $\alpha \in [a, b]$,由 $\alpha \neq 0$,$\left|\widehat{f_{\alpha}}(n)\right| < 1$($n \in \mathbf{Z}$).进一步,若 $\sup_{\alpha \in [a, b]}\sup_{n}\left|\widehat{f_{\alpha}}(n)\right| = 1$,则存在 $\{\alpha_{j}\}$,$\{n_{\alpha_{j}}\}$ 使 $\left|\widehat{f_{\alpha_{j}}}(n_{\alpha_{j}})\right| \to 1$($j \to \infty$).因 $[a, b]$ 是紧集,不妨设 $\alpha_{j} \to \alpha_{0}$($j \to \infty$).此时
\begin{equation*}
\left|\widehat{f_{\alpha_{j}}}(n) - \widehat{f_{\alpha_{0}}}(n)\right| \leq \int\left|e^{i\alpha_{j}\cos x} - e^{i\alpha_{0}\cos x}\right|\dd x \to 0 \quad (\alpha_{j} \to \alpha_{0})
\end{equation*}
且若 $|n_{\alpha_{j}}| \to \infty$,则 $\widehat{f_{\alpha_{0}}}(n_{\alpha_{j}}) \to 0$,故不妨设 $\{n_{\alpha_{j}}\}$ 有界.进一步,由 $\left|\widehat{f_{\alpha_{j}}}(n_{\alpha_{j}})\right| \to 1$($j \to \infty$),不妨设 $n_{\alpha_{j}} = n_{\alpha_{0}}$.即 $f_{\alpha_{0}}(x) = e^{in_{\alpha_{0}}x}$.这意味著 $e^{in_{\alpha_{0}}x} = e^{i\alpha\cos x}$($x \in \mathbf{T}$),产生矛盾.

第 3 阶段.令
\begin{equation*}
f(x) = \sum_{k=1}^{\infty}\frac{1}{k^{2}}\cos p_{k}x
\end{equation*}
时,适当选取 $\{p_{k}\} \subset \mathbf{N}$,则 $f(x)$ 即为所求.

第 3 阶段的证明.选取满足 $\prod_{k=1}^{\infty}(1 + \varepsilon_{k}) < 2$ 的 $\varepsilon_{k} > 0$,$p_{k}$ 如下选取.令 $p_{1} = 1$,设 $p_{n}$ 以前已选取.利用第 1 阶段选取 $p_{n+1}$:
\begin{align*}
& \left\|\prod_{k=1}^{n+1}e^{i(u/k^{2})\cos p_{k}x}\right\|_{A^{\infty}} \\
& = \left\|e^{i\sum_{k=1}^{n+1}(u/k^{2})\cos p_{k}x}\right\|_{A^{\infty}} \\
& \leq (1 + \varepsilon_{n})\left\|\prod_{k=1}^{n}e^{i(u/k^{2})\cos p_{k}x}\right\|_{A^{\infty}}\left\|e^{i(u/(n+1)^{2})\cos p_{n+1}x}\right\|_{A^{\infty}} \quad (|u| \leq n^{2}).
\end{align*}
这样选取 $\{p_{k}\} \subset \mathbf{N}$.此时,
\begin{equation*}
f(x) = \sum_{k=1}^{\infty}\frac{1}{k^{2}}\cos p_{k}x
\end{equation*}
即为所求.事实上,固定 $u \neq 0$.不妨设 $n \in \mathbf{N}$ 满足 $4^{n-1} \leq |u| < 4^{n}$.此时,由 $\left\|e^{i(u/k^{2})\cos p_{k}x}\right\|_{A^{\infty}} \leq 1$,
\begin{align*}
\|e^{iuf}\|_{A^{\infty}} & = \left\|e^{iu\sum_{k=1}^{\infty}(1/k^{2})\cos p_{k}x}\right\|_{A^{\infty}} \\
& \leq \left\|\prod_{k=1}^{2^{n+1}}e^{i(u/k^{2})\cos p_{k}x}\right\|_{A^{\infty}}\left\|\prod_{k=2^{n+1}+1}^{\infty}e^{i(u/k^{2})\cos p_{k}x}\right\|_{A^{\infty}} \\
& \leq 2\prod_{k=2^{n}}^{2^{n+1}}(1 + \varepsilon_{k})\left\|e^{i(u/k^{2})\cos p_{k}x}\right\|_{A^{\infty}}.
\end{align*}
这里,
\begin{equation*}
\frac{1}{16} = \frac{4^{n-1}}{4^{n+1}} \leq \frac{|u|}{k^{2}} < \frac{4^{n}}{4^{n}} = 1 \quad (2^{n} \leq k \leq 2^{n+1})
\end{equation*}
故 $1/16 \leq |u|/k^{2} < 1$.从而利用第 2 阶段($a = 1/16$,$b = 1$,$\alpha = |u|/k^{2}$),
\begin{equation*}
\|e^{iuf}\|_{A^{\infty}} \leq 2\left(\prod_{k=2^{n}}^{2^{n+1}}(1 + \varepsilon_{k})\right)\rho^{2^{n}} \quad (0 < \rho < 1).
\end{equation*}
因此,令 $\delta = -\log\rho$,
\begin{equation*}
\|e^{iuf}\|_{A^{\infty}} \leq 4e^{|u|^{1/2}\log\rho} = 4e^{-\delta|u|^{1/2}}.
\end{equation*}
于是,现在设 $g \in A(\mathbf{T})$,$\|g\|_{A(\mathbf{T})} \leq 1$,则
\begin{align*}
\left|\int_{0}^{2\pi}e^{itf(x)}g(x)\dd x\right| & = \left|\int_{0}^{2\pi}e^{itf(x)}\sum_{n \in \mathbf{Z}}\hat{g}(n)e^{inx}\dd x\right| \\
& \leq \sum|\hat{g}(n)|\left|\int_{0}^{2\pi}e^{itf(x)}e^{inx}\dd x\right| \\
& \leq 2\pi\|g\|_{A(\mathbf{T})}\|e^{itf}\|_{A^{\infty}}
\end{align*}
得 $\omega(t) = O\left(e^{-\delta|t|^{1/2}}\right)$.至此,第 3 阶段证明完毕.
\end{proof}

综上,命题得证.

\subsection{关于 Kronecker 集}\label{sec:3.3.3}

\begin{definition}\label{def:3.3.6}
$\mathbf{T}$ 的闭子集 $K$ 称为 Kronecker 集,是指对 $f \in S(K) = \{f \in C(K):\ |f(x)| = 1\ (x \in K)\}$,任意取 $\varepsilon > 0$ 时,存在 $n \in \mathbf{Z}$ 使得
\begin{equation*}
\max_{x \in K}|f(x) - e^{inx}| < \varepsilon.
\end{equation*}
\end{definition}

\begin{definition}\label{def:3.3.7}
$E \subset \mathbf{T}$ 在 $\mathbf{Q}$ 上线性无关,是指对 $\{x_{1}, \ldots, x_{k}\} \subset E \subset [-\pi, \pi) = \mathbf{T}$,$\{q_{0}, q_{1}, \ldots, q_{k}\} \subset \mathbf{Z}$,若 $q_{1}x_{1} + \cdots + q_{k}x_{k} = \pi q_{0}$,则 $q_{0} = q_{1} = \cdots = q_{k} = 0$.
\end{definition}

\textbf{Kronecker 定理}(\cite{ch3:8})．若 $\{x_{1}, x_{2}, \ldots, x_{n}\} \subset \mathbf{T}$ 在 $\mathbf{Q}$ 上线性无关,则对任意 $\{\beta_{1}, \beta_{2}, \ldots, \beta_{n}\} \subset \mathbf{R}$ 与任意 $\varepsilon > 0$,存在 $k \in \mathbf{Z}$ 使得
\begin{equation*}
\left|e^{ikx_{j}} - e^{i\beta_{j}}\right| < \varepsilon \quad (j = 1, \ldots, n).
\end{equation*}

下面,利用 Kronecker 定理容易证明.

\begin{theorem}\label{thm:3.3.8}
$\mathbf{T}$ 的线性无关的有限子集 $E$ 是 Kronecker 集.
\end{theorem}

\begin{theorem}\label{thm:3.3.9}
Kronecker 集 $E$ 是线性无关的集合.
\end{theorem}

\begin{proof}
对 $\{x_{1}, \ldots, x_{k}\} \subset E$,$q_{j} \in \mathbf{Q}$($j = 0, 1, \ldots, k$),设
\begin{equation*}
\sum_{j=1}^{k}q_{j}x_{j} = 2\pi q_{0}.
\end{equation*}
此时,对 $n \in \mathbf{Z}$,
\begin{equation*}
\prod_{j=1}^{k}(e^{inx_{j}})^{q_{j}} = e^{in\sum_{j=1}^{k}q_{j}x_{j}} = e^{i2\pi q_{0}n}
\end{equation*}
因此 $\prod_{j=1}^{k}(e^{inx_{j}})^{q_{j}} = 1$.这里,由 Kronecker 集的定义,取 $\alpha_{j} \in \mathbf{R}$($j = 1, \ldots, k$),则对 $\varepsilon > 0$ 存在 $m \in \mathbf{Z}$ 使得 $\left|e^{imx_{j}} - e^{i\alpha_{j}}\right| < \varepsilon$($j = 1, \ldots, k$).因此 $e^{i\sum_{j=1}^{k}\alpha_{j}q_{j}} = 1$($\alpha_{j} \in \mathbf{R}$).因此 $q_{j} = 0$($j = 1, \ldots, k$).从而 $q_{0} = 0$.综上,$E$ 线性无关.
\end{proof}

下面,由 Rudin \cite{ch3:11} 证明非平凡完全不连通完全集(称为 Cantor 集)构成的 Kronecker 集的存在性.

\begin{lemma}\label{lem:3.3.10}
设 $U_{1}, U_{2}, \ldots, U_{k}$ 为 $(0, \pi)$ 的非空互不相交开区间列,则存在 $x_{j} \in U_{j}$($j = 1, \ldots, k$)使 $\{x_{1}, \ldots, x_{k}\}$ 是 Kronecker 集.
\end{lemma}

\begin{proof}
由于 $U_{j}$ 是不可数集,可取 $x_{j} \in U_{j}$($j = 1, \ldots, k$)使得 $\{x_{1}, \ldots, x_{k}\}$ 线性无关.此时,由定理~\ref{thm:3.3.8} 成立.
\end{proof}

\begin{theorem}\label{thm:3.3.11}
$\mathbf{T}$ 中存在是 Cantor 集的 Kronecker 集.
\end{theorem}

\begin{proof}
令 $P_{1}^{0} = \left[\frac{1}{3}, \frac{1}{2}\right] \subset (-\pi, \pi)$.设对 $r \geq 1$,已作出 $s = 2^{r-1}$ 个互不相交的闭区间 $P_{1}^{r-1}, P_{2}^{r-1}, \ldots, P_{s}^{r-1}$.此时,在 $P_{j}^{r-1}$ 中取互不相交的开区间 $W_{2j-1}, W_{2j}$($1 \leq j \leq s$).这样,由补题~\ref{lem:3.3.10},存在 Kronecker 集 $\{x_{1}^{r}, \ldots, x_{2s}^{r}\}$($x_{j}^{r} \in W_{j}$,$1 \leq j \leq 2s$).这里,存在 $A_{r} \subset \mathbf{Z}$(有限集)使得对任意 $\{\alpha_{1}, \ldots, \alpha_{2^{r}}\} \subset \mathbf{R}$,取适当的 $n \in A_{r}$,
\begin{equation*}
\left|e^{i\alpha_{j}} - e^{inx_{j}^{r}}\right| < \frac{1}{r} \quad (j = 1, \ldots, 2^{r}).
\end{equation*}
此时,存在闭区间 $P_{j}^{r} \subset W_{j}$ 使得
\begin{equation*}
\left|e^{inx} - e^{inx_{j}^{r}}\right| < \frac{1}{r} \quad (x \in P_{j}^{r},\ n \in A_{r}).
\end{equation*}
进一步,设 $P_{j}^{r}$ 的长度 $\left|P_{j}^{r}\right| < 1/r$.如此归纳继续,令
\begin{equation*}
P = \bigcap_{r=1}^{\infty}\bigcup_{j=1}^{2^{r}}P_{j}^{r}.
\end{equation*}
由构造方法,$P$ 是 Cantor 集.设 $f \in S(P)$,$\varepsilon > 0$.由 $f \in C(P)$,$f$ 在 $P$ 上一致连续.因此存在足够大的 $r_{0}$ 使得 $f(P \cap P_{j}^{r_{0}})$($1 \leq j \leq 2^{r_{0}}$)不成为 $\{z \in \mathbf{C}:\ |z| = 1\}$ 的全部.此时,由 $P \cap P_{j}^{r_{0}} \subset (-\pi, \pi)$,存在 $\widetilde{f} \in S\left(\bigcup_{j}P_{j}^{r_{0}}\right)$ 使 $\widetilde{f}\Big|_{P \cap \left(\bigcup_{j}P_{j}^{r_{0}}\right)} = f$.进一步,由 $P_{j}^{r}$ 的构造方法,对 $\varepsilon > 0$,$r > r_{0}$,$r > 3/\varepsilon$,可设 $\left|\widetilde{f}(x) - \widetilde{f}(x_{j}^{r})\right| < \varepsilon/3$($x \in P_{j}^{r}$,$1 \leq j \leq 2^{r}$).又由构造方法,存在 $n \in A_{r}$ 使 $\left|\widetilde{f}(x_{j}^{r}) - e^{inx_{j}^{r}}\right| < 1/r$($1 \leq j \leq 2^{r}$).因此,对 $x \in \bigcup_{j}P_{j}^{r}$,
\begin{align*}
\left|\widetilde{f}(x) - e^{inx}\right| & \leq \left|\widetilde{f}(x) - \widetilde{f}(x_{j}^{r})\right| + \left|\widetilde{f}(x_{j}^{r}) - e^{inx_{j}^{r}}\right| + \left|e^{inx_{j}^{r}} - e^{inx}\right| \\
& < \frac{\varepsilon}{3} + \frac{1}{r} + \frac{1}{r} < \varepsilon.
\end{align*}
因此,$\left|f(x) - e^{inx}\right| < \varepsilon$($x \in P$).
\end{proof}

\begin{note}\label{note:3.3.12}
(cf.\ \cite{ch3:11})上述中令 $Q = P \cup (-P)$,取满足 $Q$ 上任意一点测度为 $0$ 的 $\mu \in M(\mathbf{T})$ 时,测度 $\delta_{0}, \mu, \mu^{2}, \mu^{3}, \ldots, \mu^{n} = \mu^{n-1} * \mu, \ldots$ 聚集在互不相交的集合上.这一点可由 $P$ 是线性无关集与测度乘积的定义证明.事实上,取 $n, m \geq 1$,令 $Q_{m} = \{x_{1} + \cdots + x_{m}:\ x_{j} \in Q\}$,则 $\mu^{n}(Q_{m}) = 0$($m < n$).例如,$n = 3$,$m = 1$,$\mu \geq 0$ 时,
\begin{equation}\tag{3.3.1}\begin{aligned}
\mu^{3}(Q_{m}) & = \displaystyle\int \mu^{2}(Q - x)\dd\mu(x) \\
& = \displaystyle\int\!\!\int \chi_{Q-x}(w)\dd\mu^{2}(w)\dd\mu(x) \\
& = \mu \times \mu \times \mu\bigl(\{(x, y, z) \in Q \times Q \times Q:\ x + y + z \in Q\}\bigr) \\
& \leq \sum_{i \neq j,\ 1 \leq i, j \leq 3}\mu \times \mu \times \mu\bigl(\{(x_{1}, x_{2}, x_{3}) \in Q \times Q \times Q:\ x_{i} \pm x_{j} = 0\}\bigr).
\end{aligned}\label{eq:3.3.1}\end{equation}
其中,$\chi_{Q-x}(w)$ 是 $Q - x$ 的特征函数.这里,用 Fubini 定理与 $\mu(\{x\}) = 0$($x \in \mathbf{T}$),\eqref{eq:3.3.1} 的右边为 $0$.用这样的想法可得 $\mu^{n}(Q_{m}) = 0$($m < n$).
\end{note}

由定义与 Riemann-Lebesgue 引理可知,Kronecker 集的 Lebesgue 测度为 $0$,且是完全不连通集.Varopoulos \cite{ch3:50} 利用这一性质对谱综合证明了下列结论.其证明是考察 $A(\mathbf{T})$ 上的连续线性泛函,充分利用 Kronecker 集的性质而进行的函数分析方法的证明,饶有趣味,但此处从略.证明请参照论文 \cite{ch3:50}.

\begin{theorem}\label{thm:3.3.13}
$\mathbf{T}$ 的 Kronecker 集是 $A(\mathbf{T})$ 的谱综合集.
\end{theorem}

即使在非离散局部紧 Abel 群上,也可用特征定义 Kronecker 集,但众所周知,一般未必完全不连通.Saeki \cite{ch3:44} 发展了 Varopoulos \cite{ch3:50} 的方法,得到了推广到非离散局部紧 Abel 群的结果.

% ============ 第3章 §3.4 关于谱综合----Varopoulos 的方法(原书 pp.249-265) ============
\section{关于谱综合----Varopoulos 的方法}\label{sec:3.4}
\index{谱综合!非谱综合集}
\index{Helson 集}
\index{相对零化}
\index{受限锥性质}
\index{曲率}
\index{第二基本矩阵}
\index{Katznelson 定理}
\index{Wiener--Lévy 定理}
\index{作用函数!(函数空间的)}
\index{作用!(函数空间的情形)}


Schwartz \cite{ch3:49} 证明了 $\mathbf{R}^{3}$ 的单位球面 $S^{2}$ 不是 $L^{1}(\mathbf{R}^{3})$ 中的谱综合集.这里,我们叙述利用这一结果,关于交换 Banach 代数谱综合的 Varopoulos 的研究 \cite{ch3:52}(以及其参考文献 cf. \cite{ch3:23}).此外,也对 L. Schwartz 这一研究在几何方向上的发展作一补充.

\subsection{Schwartz 的结果}\label{sec:3.4.1}

下面叙述的是 Schwartz \cite{ch3:49} 的结果.

\begin{theorem}\label{thm:3.4.1}
当 $n \geq 3$ 时,$S^{n-1}$ 不是 $A(\mathbf{R}^{n})$ 的谱综合集.
\end{theorem}

\begin{proof}
令
\begin{equation*}
I(S^{n-1}) = \{f \in A(\mathbf{R}^{n}): f(y) = 0 \ (y \in S^{n-1})\},
\end{equation*}
\begin{equation*}
I_{1}(S^{n-1}) = \left\{f \in I(S^{n-1}) \cap C^{\infty}(\mathbf{R}^{n}): \frac{\partial f}{\partial y_{1}}(y) = 0 \ (y \in S^{n-1})\right\}.
\end{equation*}
以 $\mu$ 表示 $S^{n-1}$ 上的面积元素,并设 $\mu(S^{n-1}) = 1$.此时
\begin{equation*}
\widehat{\mu}(y) = \int_{S^{n-1}} e^{-iy \cdot x}\dd\mu(x) = \Gamma\left(\frac{n}{2}\right)\left(\frac{2}{|y|}\right)^{\frac{n-2}{2}} J_{\frac{n-2}{2}}(|y|) \quad (n \geq 3,\ y \neq 0)
\end{equation*}
(cf. \cite{ch3:11}).其中,$\Gamma$ 是 Gamma 函数,$J_{\alpha}$ 是第一类 Bessel 函数.此时,对 $y = (y_{1}, \ldots, y_{n})$,
\begin{equation*}
|y_{1}||\widehat{\mu}(y)| \leq C\left(\frac{1}{|y|}\right)^{\frac{n-2}{2}}\left|J_{\frac{n-2}{2}}(|y|)\right||y|
\end{equation*}
由 Bessel 函数的性质($C$ 是适当的常数)
\begin{equation*}
\left\{\begin{array}{ll}
\left|J_{\frac{n-2}{2}}(|y|)\right| \leq C'|y|^{\frac{n-2}{2}} & (|y| \to 0),\\
\left|J_{\frac{n-2}{2}}(|y|)\right| \leq C'|y|^{-\frac{1}{2}} & (|y| \to \infty)
\end{array}\right.
\end{equation*}
可知当 $n \geq 3$ 时,对适当的常数 $C$ 有 $|y_{1}\widehat{\mu}(y)| \leq C$($y \in \mathbf{R}^{n}$).于是,令
\begin{align*}
Tf &= -i\int_{\mathbf{R}^{n}}\tilde{f}(y)y_{1}\widehat{\mu}(y)\frac{\dd y}{2\pi}\\
&\quad \left(\tilde{f} \in C_{c}^{\infty}(\mathbf{R}^{n}) \subset A(\mathbf{R}^{n}),\ f \text{ 是 } \tilde{f} \text{ 的 Fourier 变换}\right)
\end{align*}
则
\begin{equation*}
|Tf| \leq C\int_{\mathbf{R}^{n}}|\tilde{f}(y)|\dd y = \|f\|_{A(\mathbf{R}^{n})}.
\end{equation*}
又因 $C_{c}^{\infty}(\mathbf{R}^{n})$ 在 $A(\mathbf{R}^{n})$ 中稠密,故可得 $|Tf| \leq C\|f\|_{A(\mathbf{R}^{n})}$.进一步,
\begin{align*}
Tf &= -i\int_{\mathbf{R}^{n}}\tilde{f}(y)y_{1}\widehat{\mu}(y)\frac{\dd y}{2\pi}\\
&= \iint (-iy_{1})\tilde{f}(y)e^{-iy \cdot x}\frac{\dd y}{2\pi}\dd\mu(x)\\
&= \int \lim_{h \to 0}\frac{\int \tilde{f}(y)e^{-i(x+h_{1})\cdot y}\frac{\dd y}{2\pi} - \int \tilde{f}(y)e^{-ix \cdot y}\frac{\dd y}{2\pi}}{h}\dd\mu(x)
\end{align*}
\begin{equation*}
= \int \frac{\partial f}{\partial x_{1}}(x)\dd\mu(x)
\end{equation*}
其中设 $h_{1} = (h, 0, \ldots, 0)$.此时,若 $f \in I_{1}(S^{n-1})$,则 $Tf = 0$.可是,在 $1/2 < |x| < 3/2$ 上取满足 $g(x) = x_{1}(|x|^{2} - 1)$ 的 $g \in C_{c}^{\infty}(\mathbf{R}^{n})$,则 $g \in I(S^{n-1})$,而且
\begin{equation*}
Tg = \int_{S^{n-1}} 2x_{1}^{2}\dd\mu(x) > 0.
\end{equation*}
于是由泛函分析的基础知识
\begin{equation*}
I(S^{n-1}) \supsetneqq \overline{I_{1}(S^{n-1})} \supseteq J(S^{n-1})
\end{equation*}
成立.因此,考虑到 $\overline{I_{1}(S^{n-1})}$ 是 $A(\mathbf{R}^{n})$ 的闭理想,$S^{n-1}$ 不是谱综合集.
\end{proof}

\subsection{\texorpdfstring{$A(\mathbf{T}^3)$ 的谱综合}{A(T^3) 的谱综合}}\label{sec:3.4.2}

由 3.4.1 小节可知 $A(\mathbf{R}^{3})$ 是不可谱综合的交换 Banach 代数.利用这一点,可如下证明 $A(\mathbf{T}^{3})$ 不可谱综合.

\begin{proposition}\label{prop:3.4.2}
设
\begin{equation*}
E = \left\{(x_{1}, x_{2}, x_{3}) \in \mathbf{R}^{3}: |x_{1}| \geq \frac{\pi}{2} \text{ 或 } |x_{2}| \geq \frac{\pi}{2} \text{ 或 } |x_{3}| \geq \frac{\pi}{2}\right\},
\end{equation*}
则 $E$ 是 $A(\mathbf{R}^{3})$ 的谱综合集.
\end{proposition}

\begin{proof}
设 $f \in A(\mathbf{R}^{3})$,$f = 0$ on $E$,并令 $f_{\alpha}(x) = \alpha^{3}f(\alpha x)$,$0 < \alpha < 1$,则
\begin{equation*}
\|f_{\alpha} - f\|_{A(\mathbf{R}^{3})} \to 0 \quad (\alpha \to 1),
\end{equation*}
且 $f_{\alpha}$ 在 $E$ 的邻域上为 0.因此 $E$ 是谱综合集.
\end{proof}

\begin{proposition}\label{prop:3.4.3}
在命题~\ref{prop:3.4.2} 的记号下,$E \cup S^{2}$ 不是 $A(\mathbf{R}^{3})$ 的谱综合集.
\end{proposition}

\begin{proof}
假设 $E \cup S^{2}$ 是谱综合集,取 $f \in I(S^{2})$.对 $\varepsilon > 0$,取 $h \in A(\mathbf{R}^{3})$ 使得 $\supp h$ 是紧集且 $\|f - hf\|_{A(\mathbf{R}^{3})} < \varepsilon$.再取 $g \in A(\mathbf{R}^{3})$ 使得 $g$ 在 $S^{2}$ 的邻域上为 0,在 $\supp h \cap E$ 上为 1.此时因 $fh - ghf \in I(E \cup S^{2})$,由 $E \cup S^{2}$ 是谱综合集的假设,存在在 $E \cup S^{2}$ 的邻域上为 0 的 $k \in A(\mathbf{R}^{3})$ 使得 $\|k - (fh - fhg)\|_{A(\mathbf{R}^{3})} < \varepsilon$.从而
\begin{equation*}
\|f - (k + fhg)\|_{A(\mathbf{R}^{3})} \leq \|f - fh\|_{A(\mathbf{R}^{3})} + \|fh - (k + fhg)\|_{A(\mathbf{R}^{3})} < 2\varepsilon,
\end{equation*}
于是 $S^{2}$ 成为 $A(\mathbf{R}^{3})$ 的谱综合集,产生矛盾.
\end{proof}

\begin{proposition}\label{prop:3.4.4}
设 $f \in A(\mathbf{T}^{3})$ 满足 $f(x_{1}, x_{2}, x_{3}) = 0$($\frac{\pi}{2} \leq |x_{1}| \leq \pi$ 或 $\frac{\pi}{2} \leq |x_{2}| \leq \pi$ 或 $\frac{\pi}{2} \leq |x_{3}| \leq \pi$).此时,令
\begin{equation*}
g(x_{1}, x_{2}, x_{3}) = \left\{\begin{array}{ll}
f(x_{1}, x_{2}, x_{3}) & (|x_{j}| \leq \frac{\pi}{2},\ j = 1, 2, 3),\\
0 & (\text{其他})
\end{array}\right.
\end{equation*}
则 $g \in A(\mathbf{R}^{3})$,且存在常数 $C_{1}, C_{2} > 0$ 使得
\begin{equation*}
C_{1}\|g\|_{A(\mathbf{R}^{3})} \leq \|f\|_{A(\mathbf{T}^{3})} \leq C_{2}\|g\|_{A(\mathbf{R}^{3})}.
\end{equation*}
\end{proposition}

\begin{proof}
设 $E' = \{(x_{1}, x_{2}, x_{3}): |x_{j}| \leq \frac{\pi}{2},\ j = 1, 2, 3\}$.取 $h \in C_{c}^{\infty}(\mathbf{R}^{3})$ 使得 $h$ 在 $E'$ 的邻域上为 1 且 $\supp h \subset \{(x_{1}, x_{2}, x_{3}): |x_{j}| \leq \frac{2\pi}{3},\ j = 1, 2, 3\}$.对 $f \in A(\mathbf{T}^{3})$,由 $g(x) = f(x)h(x)$($x = (x_{1}, x_{2}, x_{3})$)得 $g \in A(\mathbf{R}^{3})$ 且
\begin{equation*}
\|g\|_{A(\mathbf{R}^{3})} \leq \|h\|_{A(\mathbf{R}^{3})}\|f\|_{A(\mathbf{T}^{3})}.
\end{equation*}
反之,若 $g \in A(\mathbf{R}^{3})$,则当 $f = gh$ 时,
\begin{equation*}
\widehat{f}(n) = \frac{1}{(2\pi)^{3}}\int_{\mathbf{R}^{3}}g(x)h(x)e^{-ix \cdot n}\dd x \quad (n \in \mathbf{Z}^{3}).
\end{equation*}
由于 $\widehat{h} \in L^{1}(\mathbf{R}^{3})$,
\begin{equation*}
\widehat{f}(n) = \int_{\mathbf{R}^{3}}\widehat{g}(y)\widehat{h}(n - y)\dd y.
\end{equation*}
此时由 $h \in C_{c}^{\infty}(\mathbf{R}^{3})$,存在 $b > 0$ 使得
\begin{equation*}
\sum_{n \in \mathbf{Z}^{3}}|\widehat{h}(n - y)| < b \quad (y \in \mathbf{R}^{3}),
\end{equation*}
故得 $\|f\|_{A(\mathbf{T}^{3})} \leq b\|g\|_{A(\mathbf{R}^{3})}$.
\end{proof}

\begin{theorem}\label{thm:3.4.5}
$A(\mathbf{T}^{3})$ 不可谱综合.
\end{theorem}

\begin{proof}
利用命题~\ref{prop:3.4.2}～3.4.4 的记号.令
\begin{equation*}
\widetilde{E} = \left\{x \in \mathbf{T}^{3}: |x_{1}| \geq \frac{\pi}{2} \text{ 或 } |x_{2}| \geq \frac{\pi}{2} \text{ 或 } |x_{3}| \geq \frac{\pi}{2}\right\}.
\end{equation*}
此时 $\widetilde{E}$ 是 $A(\mathbf{T}^{3})$ 的谱综合集.事实上,若 $h \in I(\widetilde{E}) = \{f \in A(\mathbf{T}^{3}): f = 0 \text{ on } \widetilde{E}\}$,则由命题~\ref{prop:3.4.4},取 $h$ 所对应的 $g \in A(\mathbf{R}^{3})$,则 $g \in I(E)$.由命题~\ref{prop:3.4.2},对 $\varepsilon > 0$,存在在 $E$ 的邻域上为 0 且 $\|g - g_{0}\|_{A(\mathbf{R}^{3})} < \varepsilon$ 的 $g_{0} \in A(\mathbf{R}^{3})$.此时,设 $h_{0}$ 为命题~\ref{prop:3.4.4} 中对应于 $g_{0}$ 的函数,则 $h_{0}$ 在 $\widetilde{E}$ 的邻域上为 0 且满足 $\|h - h_{0}\|_{A(\mathbf{T}^{3})} < C\varepsilon$,所以 $\widetilde{E}$ 是 $A(\mathbf{T}^{3})$ 的谱综合集.另一方面,设 $\widetilde{S^{2}} = \{x = (x_{1}, x_{2}, x_{3}) \in \mathbf{T}^{3}: x_{1}^{2} + x_{2}^{2} + x_{3}^{2} = 1\}$,$\widetilde{K} = \widetilde{E} \cup \widetilde{S^{2}}$,若 $\widetilde{K}$ 成为 $A(\mathbf{T}^{3})$ 的谱综合集,则由命题~\ref{prop:3.4.4},$K = E \cup S^{2}$ 成为 $A(\mathbf{R}^{3})$ 的谱综合集,与命题~\ref{prop:3.4.3} 矛盾.因此 $\widetilde{K}$ 不是 $A(\mathbf{T}^{3})$ 的谱综合集.故 $A(\mathbf{T}^{3})$ 不可谱综合.
\end{proof}

\subsection{\texorpdfstring{交换 Banach 代数 $V(K)$}{交换 Banach 代数 V(K)}}\label{sec:3.4.3}

基于 Varopoulos \cite{ch3:52},设 $K_{1}, K_{2}$ 为紧 Hausdorff 空间,如下定义 $V(K_{1} \times K_{2})$.

\begin{definition}\label{def:3.4.6}
称 $\varphi \in V(K_{1} \times K_{2})$,是指 $\varphi(x, y) = \sum_{j=1}^{\infty} f_{j}(x)g_{j}(y)$,其中 $f_{j} \in C(K_{1})$,$g_{j} \in C(K_{2})$ 满足 $\sum_{j=1}^{\infty}\|f_{j}\|_{\infty}\|g_{j}\|_{\infty} < \infty$.不过,$\|\cdot\|_{\infty}$ 是作为连续函数空间的范数.此时令
\begin{equation*}
\begin{aligned}
&\|\varphi\|_{V} = \inf\left\{\sum_{j=1}^{\infty}\|f_{j}\|_{\infty}\|g_{j}\|_{\infty}: \varphi(x, y) = \sum_{j=1}^{\infty}f_{j}(x)g_{j}(y),\right.\\
&\left.\sum_{j=1}^{\infty}\|f_{j}\|_{\infty}\|g_{j}\|_{\infty} < \infty\right\}.
\end{aligned}
\end{equation*}
$V(K_{1} \times K_{2}) = C(K_{1}) \widehat{\otimes}\, C(K_{2})$ 也记为 $V(K_{1} \times K_{2})$.
\end{definition}

\begin{proposition}\label{prop:3.4.7}
$V(K_{1} \times K_{2})$ 关于 $\|\cdot\|_{V}$ 是半单,正则的交换 Banach 代数,其极大理想空间为 $K_{1} \times K_{2}$.
\end{proposition}

\begin{proof}
$V(K_{1} \times K_{2})$ 关于 $\|\cdot\|_{V}$ 成为交换 Banach 代数这一点是容易证明的.又,若 $(x, y) \in K_{1} \times K_{2}$,则由 $V(K_{1} \times K_{2}) \ni \varphi \mapsto \varphi(x, y) \in \mathbf{C}$ 可把 $K_{1} \times K_{2}$ 视为 $\Delta(V(K_{1} \times K_{2}))$ 的子集.反之,若 $h \in \Delta(V(K_{1} \times K_{2}))$,则存在 $\varphi_{0} \in V(K_{1} \times K_{2})$ 使得 $h(\varphi_{0}) \neq 0$.此时,对 $f \in C(K_{1})$,$g \in C(K_{2})$ 令
\begin{equation*}
h_{1}(f) = \frac{h(f\varphi_{0})}{h(\varphi_{0})},\quad h_{2}(g) = \frac{h(g\varphi_{0})}{h(\varphi_{0})},
\end{equation*}
则 $h_{1} \in \Delta(C(K_{1}))$,$h_{2} \in \Delta(C(K_{2}))$.事实上,若 $h(\varphi') \neq 0$,则由
\begin{equation*}
\frac{h(f\varphi_{0})}{h(\varphi_{0})} = \frac{h(f\varphi_{0})h(\varphi')}{h(\varphi_{0}\varphi')} = \frac{h(f\varphi_{0}\varphi')}{h(\varphi_{0}\varphi')} = \frac{h(f\varphi')}{h(\varphi')}
\end{equation*}
成立,可得
\begin{align*}
h_{1}(f_{1}f_{2}) &= \frac{h(f_{1}f_{2}\varphi_{0}^{2})}{h(\varphi_{0}^{2})} = \frac{h(f_{1}\varphi_{0})h(f_{2}\varphi_{0})}{h(\varphi_{0})h(\varphi_{0})} = h_{1}(f_{1})h_{1}(f_{2})\\
&\quad (f_{1}, f_{2} \in C(K_{1})).
\end{align*}
因此,由 $\mathfrak{M}(C(K_{1})) = K_{1}$,存在 $x_{1} \in K_{1}$ 使得 $h_{1}(f) = f(x_{1})$($f \in C(K_{1})$).对 $h_{2}$ 也同样.于是
\begin{equation*}
\begin{aligned}
i: K_{1} \times K_{2} &\ni (x_{1}, x_{2}) \to h_{(x_{1}, x_{2})} \in \Delta(V(K_{1} \times K_{2}))\\
&(h_{(x_{1}, x_{2})}(\varphi) = \varphi(x_{1}, x_{2}) \in V(K_{1} \times K_{2}))
\end{aligned}
\end{equation*}
是双射.而且由于它连续,故是同胚映射,$K_{1} \times K_{2}$ 与 $\mathfrak{M}(V(K_{1} \times K_{2}))$ 在自然对应下等同.此外,半单,正则也是容易看出的.
\end{proof}

这里令 $V(\mathbf{T}^{3}) = V(\mathbf{T}^{3} \times \mathbf{T}^{3})$,关于 $V(\mathbf{T}^{3})$ 的谱综合,下式成立.

\begin{theorem}\label{thm:3.4.8}
$V(\mathbf{T}^{3})$ 不可谱综合.
\end{theorem}

\begin{proof}
从映射 $M$,$P$ 的定义开始.如下定义 $M: C(\mathbf{T}^{3}) \to C(\mathbf{T}^{3} \times \mathbf{T}^{3})$,$P: C(\mathbf{T}^{3} \times \mathbf{T}^{3}) \to C(\mathbf{T}^{3})$.其中 $dm(y)$ 是 $\mathbf{T}^{3}$ 上的 Lebesgue 测度,且 $m(\mathbf{T}^{3}) = 1$.
\begin{equation*}
(Mf)(x, y) = f(x + y), \quad (P\varphi)(x) = \int_{\mathbf{T}^{3}}\varphi(x - y, y)\dd m(y).
\end{equation*}
此时,若 $f \in C(\mathbf{T}^{3})$ 则 $Mf(x, y) \in C(\mathbf{T}^{3} \times \mathbf{T}^{3})$;若 $\varphi \in C(\mathbf{T}^{3} \times \mathbf{T}^{3})$ 则 $P\varphi \in C(\mathbf{T}^{3})$.又由
\begin{equation*}
(PMf)(x) = \int_{\mathbf{T}^{3}}(Mf)(x - y, y)\dd m(y) = \int_{\mathbf{T}^{3}}f(x)\dd m(y) = f(x)
\end{equation*}
$PM$ 是恒等映射.利用这一点来证明 $V(\mathbf{T}^{3})$ 不可谱综合.

由定理~\ref{thm:3.4.5},$A(\mathbf{T}^{3})$ 不可谱综合,故存在非谱综合集 $E$ 和 $f \in A(\mathbf{T}^{3})$ 使得 $f(x) = 0$($x \in E$),且 $f$ 不能被在 $E$ 的邻域上为零的 $g \in A(\mathbf{T}^{3})$ 所逼近.这里令 $E^{*} = \{(x, y): x + y \in E\}$,由 $f \in A(\mathbf{T}^{3})$ 可写成 $f(x) = \sum_{n \in \mathbf{Z}^{3}} a_{n}e^{inx}$,$\sum_{n \in \mathbf{Z}^{3}}|a_{n}| < \infty$.此时
\begin{equation*}
Mf(x, y) = f(x + y) = \sum_{n \in \mathbf{Z}^{3}}a_{n}e^{inx}e^{iny} = \sum_{n \in \mathbf{Z}^{3}}(a_{n}e^{inx})e^{iny}
\end{equation*}
所以 $Mf \in V(\mathbf{T}^{3})$.而且 $Mf(z) = 0$($z \in E^{*}$).现在,设 $\psi(x, y) = \sum f_{j}(x)g_{j}(y) \in V(\mathbf{T}^{3})$ 在 $E^{*}$ 的邻域上为零,则
\begin{equation*}
P\psi(x) = \int \psi(x - y, y)\dd m(y) = \sum_{j} f_{j} * g_{j}(x),
\end{equation*}
即得 $P\psi(x) = \sum f_{j} * g_{j}(x)$.这里由 $f_{j}, g_{j} \in L^{2}(\mathbf{T}^{3})$,
\begin{equation*}
\|f_{j} * g_{j}\|_{A(\mathbf{T}^{3})} \leq \|\widehat{f_{j}}\|_{2}\|\widehat{g_{j}}\|_{2} \leq \|f_{j}\|_{\infty}\|g_{j}\|_{\infty},
\end{equation*}
且 $\sum\|f_{j} * g_{j}\|_{A(\mathbf{T}^{3})} < \infty$,故 $P\psi \in A(\mathbf{T}^{3})$,且 $\sup_{\|\psi\|_{V} \leq 1}\|P\psi\|_{A(\mathbf{T}^{3})} \leq 1$.因此
\begin{equation*}
\|f - P\psi\|_{A(\mathbf{T}^{3})} = \|P(Mf - \psi)\|_{A(\mathbf{T}^{3})} \leq \|Mf - \psi\|_{V(\mathbf{T}^{3})}.
\end{equation*}
又,$P\psi$ 在 $E$ 的邻域上为 0.从而 $\inf\{\|Mf - \psi\|_{V(\mathbf{T}^{3})}: \psi \in V(\mathbf{T}^{3}),\ \psi \text{ 在 } E^{*} \text{ 的邻域上为零}\} > 0$,由 $E$ 的性质,$V(\mathbf{T}^{3})$ 不可谱综合.
\end{proof}

\subsection{\texorpdfstring{$A(\mathbf{T})$ 的谱综合}{A(T) 的谱综合}}\label{sec:3.4.4}

这里,用 Varopoulos 的方法 \cite{ch3:52}(cf. \cite{ch3:23})证明 $A(\mathbf{T})$ 不可谱综合.如后所述,该方法并不限于 $\mathbf{T}$,对非离散紧 Abel 群也可作类似的讨论.首先,设 $D_{2}$ 为循环群 $\{0, 1\}$ 的可数直积所构成的紧 Abel 群.把 $D_{2}$ 表示为 $D_{2} = \prod_{j=1}^{\infty}\{0, 1\}$,拓扑取直积拓扑.此时,因 $D_{2}$ 是 Cantor 集,利用这一点如下推进.如 3.3.3 小节所述,$\mathbf{T}$ 中存在作为 Cantor 集的 Kronecker 集,记之为 $K$.把 $K$ 分解为 $K = K_{1} \cup K_{2}$,$K_{1} \cap K_{2} = \emptyset$,其中 $K_{1}, K_{2}$ 是作为 Cantor 集的 Kronecker 集.本小节参考了 \cite{ch3:23}.

\begin{definition}\label{def:3.4.9}
如下定义:
(1) $\widetilde{K} = K_{1} + K_{2} = \{k_{1} + k_{2}: k_{j} \in K_{j},\ j = 1, 2\}$.此时,由 Kronecker 集的线性无关性,$k \in K$ 可唯一表示为 $k = k_{1} + k_{2}$($k_{j} \in K_{j}$,$j = 1, 2$).

(2) 注意:因 $D_{2}$ 是 Cantor 集,$D_{2}$ 与 $K_{j}$($j = 1, 2$)可等同,利用这一点可把 $D_{2} \times D_{2}$ 视为 $K_{1} + K_{2} = \widetilde{K}$.

(3) 当 $E \subset \mathbf{T}$ 为紧集时,把 $A(\mathbf{T})$ 到 $E$ 的限制 $A(E)$ 定义为商环 $A(E) = A(\mathbf{T})/I(E)$.
\end{definition}

\begin{proposition}\label{prop:3.4.10}
若 $\mathbf{T}$ 的紧集 $E$ 是 Kronecker 集,则 $A(E) = A(\mathbf{T})/I(E)$ 与 $E$ 上的连续函数空间 $C(E)$ 作为 Banach 空间同构.
\end{proposition}

\begin{proof}
把 $A(E)$ 看作 $A(\mathbf{T})$ 中函数在 $E$ 上的限制,只需证明 $A(E) = C(E)$.为证明这一点,由 Riesz 定理 $C(E)^{*} = M(E)$,只需对 $\mu \in M(E)$ 证明
\begin{equation*}
\|\mu\|_{M(E)} \leq C\|\widehat{\mu}\|_{\infty}.
\end{equation*}
不过这里令 $\|\widehat{\mu}\|_{\infty} = \sup_{m \in \mathbf{Z}}|\widehat{\mu}(m)|$.若这已证明,则因 $A(E)$ 与 $C(E)$ 都是 Banach 空间,由泛函分析中的值域定理可得 $A(E) = C(E)$.于是,对 $\mu \in M(E)$,取适当的可测函数 $f(x)$ 使 $|f(x)| = 1$($x \in E$),有 $\|\mu\| = \left|\int f(x)\dd\mu(x)\right|$;而由 $\mu \in M(E)$,$f(x)$ 可取为连续函数.又因 $E$ 是 Kronecker 集,$f(x)$ 可由适当的 $e^{imx}$ 在 $E$ 上一致逼近,故 $\|\mu\| \leq \|\widehat{\mu}\|_{\infty}$ 成立.
\end{proof}

\begin{proposition}\label{prop:3.4.11}
$A(\widetilde{K})$ 作为 Banach 代数与 $V(D_{2}) = C(D_{2}) \widehat{\otimes}\, C(D_{2})$ 同构.
\end{proposition}

\begin{proof}
注意:由定义,$\varphi \in V(D_{2})$ 可表示为 $D_{2} \times D_{2}$ 上的连续函数.

首先证明 $A(\widetilde{K}) \subset V(D_{2})$.$f \in A(\widetilde{K})$ 可写成 $f(z) = \sum_{j \in \mathbf{Z}} a_{j}e^{ijz}$($z \in \widetilde{K}$),$\sum_{j \in \mathbf{Z}}|a_{j}| < \infty$.此时,若设 $z = x + y$,$x \in K_{1}$,$y \in K_{2}$,则
\begin{equation*}
\varphi(x, y) = f(x + y) = \sum_{j \in \mathbf{Z}}a_{j}e^{ijx}e^{ijy}, \quad \sum_{j \in \mathbf{Z}}\|a_{j}e^{ijx}\|_{\infty}\|e^{ijy}\|_{\infty} < \infty,
\end{equation*}
所以 $\varphi(x, y) \in V(K_{1} \times K_{2}) = C(K_{1}) \widehat{\otimes}\, C(K_{2}) = C(D_{2}) \widehat{\otimes}\, C(D_{2}) = V(D_{2})$.于是证明了 $A(\widetilde{K}) \subset V(D_{2})$.

反之,若 $\varphi \in V(D_{2})$,则 $\varphi(x, y) = \sum_{j=1}^{\infty}f_{j}(x)g_{j}(y)$,$\sum_{j}\|f_{j}\|_{\infty}\|g_{j}\|_{\infty} \leq 2\|\varphi\|_{V(D_{2})}$.此时把 $K_{1} = D_{1}, K_{2} = D_{2}$ 看待,则 $\varphi(x, y) \in C(K_{1} \times K_{2})$.又因 $K_{1}, K_{2}$ 是 Kronecker 集,由命题~\ref{prop:3.4.10} 的证明,存在常数 $C > 0$ 使得 $\|f_{j}\|_{A(K_{1})} \leq C\|f_{j}\|_{C(K_{1})}$,$\|g_{j}\|_{A(K_{2})} \leq C\|g_{j}\|_{C(K_{2})}$.这里令
\begin{equation*}
f_{j}(x) = \sum_{l \in \mathbf{Z}}a_{jl}e^{ilx}, \quad \sum_{l}|a_{jl}| \leq \left(\frac{3}{2}\right)\|f_{j}\|_{A(K_{1})},
\end{equation*}
\begin{equation*}
g_{j}(x) = \sum_{m \in \mathbf{Z}}b_{jm}e^{imx}, \quad \sum_{m}|b_{jm}| \leq \left(\frac{3}{2}\right)\|g_{j}\|_{A(K_{2})}
\end{equation*}
($j = 1, 2, \ldots$),则
\begin{align*}
\sum_{j,l,m}|a_{jl}||b_{jm}| &\leq \frac{9}{4}\sum_{j}\|f_{j}\|_{A(K_{1})}\|g_{j}\|_{A(K_{2})}\\
&\leq \frac{9}{4}C^{2}\sum_{j}\|f_{j}\|_{\infty}\|g_{j}\|_{\infty} \leq 5C^{2}\|\varphi\|_{V(D_{2})}.
\end{align*}
进一步,令
\begin{equation*}
f(z) = \left\{\begin{array}{ll}
e^{ilz} & (z \in K_{1})\\
e^{imz} & (z \in K_{2})
\end{array}\right.
\end{equation*}
因 $K = K_{1} \cup K_{2}$ 是 Kronecker 集,对 $\varepsilon > 0$ 存在 $k_{lm} \in \mathbf{Z}$ 使得 $|f(z) - e^{ik_{lm}z}| < \varepsilon$($z \in K$).现在令
\begin{equation*}
T_{\varepsilon}\varphi(z) = \sum_{j,l,m}a_{jl}b_{jm}e^{ik_{lm}z},
\end{equation*}
则
\begin{equation*}
\|T_{\varepsilon}\varphi\|_{A(\widetilde{K})} \leq \sum_{j,l,m}|a_{jl}||b_{jm}| \leq 5C^{2}\|\varphi\|_{V(D_{2})}.
\end{equation*}
又由 $|e^{ilx}e^{imy} - e^{ik_{lm}(x+y)}| < 2\varepsilon$($x \in K_{1}$,$y \in K_{2}$),
\begin{equation*}
\|T_{\varepsilon}\varphi - \varphi\|_{V(D_{2})} \leq 2\varepsilon\sum_{j,l,m}|a_{jl}||b_{jm}| \leq 10C^{2}\varepsilon\|\varphi\|_{V(D_{2})}.
\end{equation*}
此时,考虑级数
\begin{equation}\tag{3.4.1}
T_{\varepsilon}\varphi + T_{\varepsilon}(I - T_{\varepsilon})\varphi + \cdots + T_{\varepsilon}(I - T_{\varepsilon})^{m}\varphi + \cdots
\label{eq:3.4.1}\end{equation}
其中 $I$ 是恒等映射.在这种情况下,因
\begin{equation*}
\|T_{\varepsilon}(I - T_{\varepsilon})^{m}\varphi\|_{A(\widetilde{K})} \leq 5C^{2}\|(I - T_{\varepsilon})^{m}\varphi\|_{V(D_{2})} \leq 5C^{2}(10C^{2}\varepsilon)^{m}\|\varphi\|_{V(D_{2})},
\end{equation*}
当 $\varepsilon$ 充分小时\eqref{eq:3.4.1}在 $A(\widetilde{K})$ 中收敛.又由 $\varphi - T_{\varepsilon}\varphi = \sum_{m=1}^{N}T_{\varepsilon}(I - T_{\varepsilon})^{m}\varphi + (I - T_{\varepsilon})^{N+1}\varphi$,$\|(I - T_{\varepsilon})^{N+1}\varphi\|_{V(D_{2})} \leq (10C^{2}\varepsilon)^{N+1}\|\varphi\|_{V(D_{2})} \to 0$($N \to \infty$),得 $\varphi = \sum_{m=0}^{\infty}T_{\varepsilon}(I - T_{\varepsilon})^{m}\varphi$,所以 $\varphi \in A(\widetilde{K})$.
\end{proof}

\begin{theorem}\label{thm:3.4.12}
$A(\widetilde{K})$ 不可谱综合.
\end{theorem}

\begin{proof}
设 $T = \{e^{i2\pi x}: 0 \leq x < 1\}$,如下定义从 $D_{2}$ 到 $T^{3}$ 的映射 $\tau$:
\begin{equation*}
\tau: (\varepsilon_{1}, \varepsilon_{2}, \ldots) \to (e^{2\pi i(0.\varepsilon_{1}\varepsilon_{4}\cdots)}, e^{2\pi i(0.\varepsilon_{2}\varepsilon_{5}\cdots)}, e^{2\pi i(0.\varepsilon_{3}\varepsilon_{6}\cdots)}).
\end{equation*}
这里 $\varepsilon_{j} = 0, 1$,$0.\varepsilon_{j_{1}}\varepsilon_{j_{2}}\cdots$ 表示由二进制展开得到的区间 $[0, 1)$ 的元素.这一从 $D_{2}$ 到 $T^{3}$ 的对应连续,且除 $D_{2}$ 的可数个集合外是一一的.由此在函数空间之间诱导的映射一律用 $\widetilde{\tau}$(即 $\widetilde{\tau}F = F(\tau)$)表示.又由 $\tau$ 的对应,$C(D_{2})$ 自然被视为 $L^{\infty}(\mathbf{T}^{3})$ 的子集.把这个嵌入记为 $\widetilde{i}$.另外,由 $\tau$ 诱导的从 $V(D_{2})$ 到 $L^{\infty}(\mathbf{T}^{3}) \widehat{\otimes}\, L^{\infty}(\mathbf{T}^{3})$ 的映射也用 $\widetilde{i}$ 表示.进一步如下定义交换 Banach 代数 $L^{\infty}(\mathbf{T}^{3}) \widehat{\otimes}\, L^{\infty}(\mathbf{T}^{3})$:称 $\phi \in L^{\infty}(\mathbf{T}^{3}) \widehat{\otimes}\, L^{\infty}(\mathbf{T}^{3})$,是指
\begin{align*}
\phi(x, y) &= \sum_{j=1}^{\infty}f_{j}(x)g_{j}(y) \quad (f_{j}(x), g_{j}(x) \in L^{\infty}(\mathbf{T}^{3})),\\
&\sum_{j}\|f_{j}\|_{L^{\infty}}\|g_{j}\|_{L^{\infty}} < \infty.
\end{align*}
此时,由
\begin{equation*}
\|\phi\|_{L^{\infty}(\mathbf{T}^{3}) \widehat{\otimes} L^{\infty}(\mathbf{T}^{3})} = \inf\left\{\sum_{j=1}^{\infty}\|f_{j}\|_{L^{\infty}}\|g_{j}\|_{L^{\infty}}: \text{上述表示}\right\}
\end{equation*}
$L^{\infty}(\mathbf{T}^{3}) \widehat{\otimes}\, L^{\infty}(\mathbf{T}^{3})$ 如 3.4.3 小节那样成为交换 Banach 代数.

于是,利用 $\widetilde{\tau}, \widetilde{i}$,
\begin{align*}
V(\mathbf{T}^{3}) = C(\mathbf{T}^{3}) \widehat{\otimes}\, C(\mathbf{T}^{3}) &\xrightarrow{\widetilde{\tau}} C(D_{2}) \widehat{\otimes}\, C(D_{2}) = V(D_{2})\\
&\xrightarrow{\widetilde{i}} L^{\infty}(\mathbf{T}^{3}) \widehat{\otimes}\, L^{\infty}(\mathbf{T}^{3})
\end{align*}
中各算子作用素范数不超过 1.此时,由定理~\ref{thm:3.4.8},$V(\mathbf{T}^{3})$ 不可谱综合.设闭集 $E \subset \mathbf{T}^{3} \times \mathbf{T}^{3}$ 是 $V(\mathbf{T}^{3})$ 的非谱综合集,则存在在 $E$ 上为零,且不能被在 $E$ 的邻域上为零的 $V(\mathbf{T}^{3})$ 的元逼近的 $f \in V(\mathbf{T}^{3})$.此时,$(\tau \times \tau)^{-1}(E)$ 是 $D_{2} \times D_{2}$ 的闭集,由 $\widetilde{\tau}f \in V(D_{2})$ 及 $f$ 的取法,$\widetilde{\tau}f(z) = 0$($z \in (\tau \times \tau)^{-1}(E)$).若设 $h \in V(D_{2})$ 在 $(\tau \times \tau)^{-1}(E)$ 的邻域上为零,且逼近了 $\widetilde{\tau}f$.此时,令 $\tilde{i}h(x, y) = \sum_{j}h_{1j}(x)h_{2j}(y)$,则由 $\tilde{i}h(x, y) = \sum(\tilde{i}h_{1j})(x)(\tilde{i}h_{2j})(y)$,$\tilde{i}h(x, y)$ 可视为 $L^{\infty}(\mathbf{T}^{3}) \widehat{\otimes}\, L^{\infty}(\mathbf{T}^{3})$ 的元.于是,定义函数列 $\{e_{n}(x, y)\}_{n}$ 为 $e_{n}(x, y) = e_{1n}(x)e_{2n}(y)$.其中,函数列 $\{e_{jn}\}_{n}$($j = 1, 2$)满足 $e_{jn} \in C(\mathbf{T}^{3})$,$e_{jn} \geq 0$,$\{\supp e_{jn}\}_{n}$ 是递减列,其交集为 $\{0\}$,而且 $\int e_{jn}(x)\dd x = 1$.这里,考虑
\begin{equation*}
(e_{n} * \widetilde{i}h)(x, y) = \sum_{j}(e_{1n} * (\widetilde{i}h_{1j}))(x)(e_{2n} * (\widetilde{i}h_{2j}))(y),
\end{equation*}
则 $e_{n} * \widetilde{i}h$ 在 $E$ 的邻域上为零,而且由构造方法,$e_{n} * \widetilde{i}h \in V(\mathbf{T}^{3})$,$\|e_{n} * \widetilde{i}h - f\|_{V(\mathbf{T}^{3})} \to 0$($n \to \infty$).这与 $E$ 的假设矛盾.因此 $(\tau \times \tau)^{-1}(E)$ 是 $V(D_{2})$ 的非谱综合集,$V(D_{2})$ 不可谱综合.从而由命题~\ref{prop:3.4.11},$A(\widetilde{K})$ 也不可谱综合.
\end{proof}

\begin{theorem}\label{thm:3.4.13}
$A(\mathbf{T})$ 不可谱综合.
\end{theorem}

\begin{proof}
由定理~\ref{thm:3.4.12},$A(\widetilde{K})$ 不可谱综合.因此,存在 $A(\widetilde{K})$ 的非谱综合集 $E_{0} \subset \widetilde{K} = K_{1} + K_{2}$.此时,假设 $E_{0}$ 是 $A(\mathbf{T})$ 的谱综合集.若 $f(x) = 0$($x \in E_{0}$),$f \in A(\widetilde{K})$,则存在 $\widetilde{f} \in A(\mathbf{T})$ 使得 $\widetilde{f}(x) = f(x)$($x \in \widetilde{K}$).由 $\widetilde{f}(x) = 0$($x \in E_{0}$)且 $E_{0}$ 是 $A(\mathbf{T})$ 的谱综合集,对 $\varepsilon > 0$ 存在 $\widetilde{g} \in A(\mathbf{T})$,$\widetilde{g}$ 在 $E_{0}$ 的 $\mathbf{T}$ 中的邻域上为零,且 $\|\widetilde{f} - \widetilde{g}\|_{A(\mathbf{T})} < \varepsilon$.因此,若令 $g(x) = \widetilde{g}(x)$($x \in \widetilde{K}$),则 $g \in A(\widetilde{K})$,$g$ 在 $\widetilde{K}$ 的相对拓扑下于 $E_{0}$ 的邻域上为零,且 $\|f - g\|_{A(\widetilde{K})} < \varepsilon$,这与 $E_{0}$ 是 $A(\widetilde{K})$ 的非谱综合集矛盾.故 $E_{0}$ 是 $A(\mathbf{T})$ 的非谱综合集,从而证明了 $A(\mathbf{T})$ 不可谱综合.
\end{proof}

\begin{theorem}\label{thm:3.4.14}
$A(\mathbf{R})$ 不可谱综合.
\end{theorem}

\begin{proof}
设 $E_{0} \subset T$ 为定理~\ref{thm:3.4.13} 中的 $A(\mathbf{T})$ 的非谱综合集,$E_{0}, \widetilde{K}$ 可取为 $\{|x| < \frac{\pi}{2}\}$ 的子集.假设 $E_{0}$ 是 $A(\mathbf{R})$ 的谱综合集.此时,设 $f \in A(\widetilde{K})$,$f(x) = 0$($x \in E_{0}$).$f$ 可用同一记号看作 $A(\mathbf{T})$ 到 $\widetilde{K}$ 的限制.现在,取 $h \in A(\mathbf{R})$ 使得 $h$ 在 $\widetilde{K}$ 的邻域上为 1 且 $\supp h \subset \{|x| < \frac{\pi}{2}\}$.此时,由命题~\ref{prop:3.4.4} 的想法,$fh \in A(\mathbf{R})$.因 $(fh)(x) = 0$($x \in E_{0}$),由 $E_{0}$ 是 $A(\mathbf{R})$ 的谱综合集,对 $\varepsilon > 0$ 存在 $\widetilde{g} \in A(\mathbf{R})$,$\widetilde{g}$ 在 $E_{0}$ 的邻域上为零且 $\|fh - \widetilde{g}\|_{A(\mathbf{R})} < \varepsilon$.于是取满足 $k(x) = 1$($x \in \supp h$),$\supp k \subset \{|x| < \frac{\pi}{2}\}$,$\|k\|_{A(\mathbf{R})} \leq 2$ 的 $k$,则由命题~\ref{prop:3.4.4} 的想法,$k\widetilde{g}$ 可视为 $A(\mathbf{T})$ 的元,故
\begin{align*}
\|k(\widetilde{g} - fh)\|_{A(\widetilde{K})} &\leq \|k(\widetilde{g} - fh)\|_{A(\mathbf{T})}\\
&\leq C\|k(\widetilde{g} - fh)\|_{A(\mathbf{R})} \leq 2C\|\widetilde{g} - fh\|_{A(\mathbf{R})} < 2C\varepsilon.
\end{align*}
其中 $C$ 是适当的常数.这里 $(k\widetilde{g} - fh)(x) = \widetilde{g}(x) - f(x)$($x \in \widetilde{K}$),且 $\widetilde{g}(x)$ 在 $E_{0}$ 的邻域上为零.综上,$\|\widetilde{g} - f\|_{A(\widetilde{K})} < 2C\varepsilon$,$\widetilde{g}$ 在 $E_{0}$ 的邻域上为零,与 $E_{0}$ 是 $A(\widetilde{K})$ 的非谱综合集矛盾.因此 $E_{0}$ 是 $A(\mathbf{R})$ 的非谱综合集.
\end{proof}

\begin{note}\label{note:3.4.15}
$A(\mathbf{T})$,$A(\mathbf{R})$ 不可谱综合的 Varopoulos 证明方法,依赖于作为 Cantor 集的 Kronecker 集的存在.因此,对于 $\mathbf{T}$ 以外的,可用这一方法进行讨论的非离散紧 Abel 群,也可作类似的讨论(cf. \cite{ch3:52}).
\end{note}

\subsection{\texorpdfstring{$A(\mathbf{T})$ 的非谱综合集}{A(T) 的非谱综合集}}\label{sec:3.4.5}

\begin{definition}\label{def:3.4.16}
紧集 $E \subset \mathbf{T}$ 称为 Helson 集,是指把 $A(\mathbf{T})/I(E)$ 看作 $A(\mathbf{T})$ 到 $E$ 的限制 $A(E)$ 时,$A(E) = C(E)$ 成立.
\end{definition}

即,$E$ 上任意连续函数都是具有绝对收敛 Fourier 级数的连续函数在 $E$ 上的限制.

从而,由命题~\ref{prop:3.4.10},Kronecker 集是 Helson 集.

如 3.3.3 小节所述,Varopoulos \cite{ch3:50} 证明了 Kronecker 集是 $A(\mathbf{T})$ 的谱综合集.其次,Helson 集的情形成为问题.Körner \cite{ch3:33} 于 1972 年,对长期未解决的``Helson 集是否为谱综合集''的问题,证明了``存在既是 Helson 集又是非谱综合集的集合'',给予了否定的解答.其后不久,R. Kaufman 对 T. W. Körner 的证明给出了简洁的别证.进而,S. Saeki 推广了该证明.这里只介绍结果.详情请参阅 R. Kaufman 的论文或 Saeki \cite{ch3:45}(cf. \cite{ch3:46}).

下面是由 Rudin \cite{ch3:11} 给出的(cf. \cite{ch3:33}, \cite{ch3:42}).

\begin{proposition}\label{prop:3.4.17}
$\mathbf{T}$ 存在具有如下性质的完全不连通的紧子集 $E$:

(1) $E$ 的 Lebesgue 测度为零.

(2) $\{\sigma \in M(E): \lim_{|n| \to \infty}\widehat{\sigma}(n) = 0,\ \sigma \neq 0\} \neq \emptyset$.
\end{proposition}

关于这一点,Kaufman \cite{ch3:32}(cf. \cite{ch3:45})证明了下面的事实.

\begin{theorem}[{$\;$\cite{ch3:32}(cf. \cite{ch3:45})}]\label{thm:3.4.18}
在命题~\ref{prop:3.4.17} 的记号下,存在满足 $E_{0} \subset \operatorname{supp}\sigma$ 的 Helson 集 $E_{0}$,使得 $\{T \in PM(E_{0}): \lim_{|n| \to \infty}\widehat{T}(n) = 0,\ T \neq 0\} \neq \emptyset$.其中,设 $PM(E_{0}) = \{s \in PM(\mathbf{T}): \supp s \subset E_{0}\}$,$\supp s$ 是指满足下列条件的开集 $O$ 的补集的最小者:对所有支集含于 $O$ 的 $f \in A(\mathbf{T})$,$(s, f) = 0$.
\end{theorem}

\begin{corollary}[{$\;$\cite{ch3:32},cf. \cite{ch3:45}}]\label{cor:3.4.19}
$\mathbf{T}$ 中存在既是 Helson 集又是非谱综合集的集合.
\end{corollary}

由定理~\ref{thm:3.4.18},系~\ref{cor:3.4.19} 的证明如下进行.

\begin{proof}[系~\ref{cor:3.4.19} 的证明]
设 $E_{0}$ 为定理~\ref{thm:3.4.18} 中的 Helson 集.此时,作为 Helson 集的性质,已知
\begin{equation*}
\left\{\mu \in M(\mathbf{T}): \lim_{|n| \to \infty}\widehat{\mu}(n) = 0,\ \mu \neq 0\right\} = \emptyset
\end{equation*}
成立(cf. \cite{ch3:11}).若 $E_{0}$ 是 $A(\mathbf{T})$ 的谱综合集,则定理~\ref{thm:3.4.18} 中的 $T$ 是 $(A(\mathbf{T})/J(E_{0}))^{*}$ 的元,由 $I(E_{0}) = J(E_{0})$ 得 $T \in A(E_{0})^{*}$.因 $E_{0}$ 是 Helson 集,$A(E_{0}) = C(E_{0})$,由 Riesz 表示定理得 $T \in M(E_{0})$.这与证明开头所述相矛盾.因此 $E_{0}$ 是非谱综合集.
\end{proof}

\subsection{Schwartz 结果的发展}\label{sec:3.4.6}

如前所述,L. Schwartz 证明了 $S^{n-1}$($n \geq 3$)是交换 Banach 代数 $A(\mathbf{R}^{n})$ 中的非谱综合集.将其发展的研究见于 N. Th. Varopoulos,Y. Domar,D. Müller 等人.这里简单介绍 Müller \cite{ch3:35}.详情请连同文献参阅 \cite{ch3:35}.以下使用 3.1 节的记号.

\begin{definition}\label{def:3.4.20}
(1) 设 $M$ 是 $\mathbf{R}^{n}$ 的 $k$ 维子流形,对 $\alpha \in M$ 设 $(X, \Omega)$ 为坐标邻域.即 $\Omega$ 是 $\mathbf{R}^{k}$ 的开集,$X \in C^{\infty}(\Omega, \mathbf{R}^{n})$,$X(x_{0}) = \alpha$($0 \leq k \leq n - 1$).此时,设 $N$ 为 $\alpha$ 处的法向量,令
\begin{equation*}
h_{ij}(\alpha; X, N) = N \cdot \frac{\partial^{2}}{\partial x_{i}\partial x_{j}}X(x_{1}, \ldots, x_{k}) = \frac{\partial^{2}}{\partial x_{i}\partial x_{j}}(N \cdot X),
\end{equation*}
由此定义第二基本矩阵 $(h_{ij}(\alpha; X, N))_{i,j}$.

(2) 令 $\pi(\alpha; X, N) = \{y \in \mathbf{R}^{k}: \sum_{j=1}^{k}h_{ij}y^{j} = 0,\ i = 1, \ldots, k\}$,$\pi(y) = (h_{ij})(y_{i})$,$y = (y_{i})$,则 $\pi$ 可看作从 $\alpha$ 处的切平面 $M_{\alpha}$ 到自身的映射,$\pi(\alpha; X, N)$ 可看作切平面 $M_{\alpha}$ 的子集.而且 $\dim\pi(\alpha; X, N) = k - \operatorname{rank}(h_{ij})$.另请注意 $\pi(\alpha; X, N)$ 不依赖于 $X$ 的取法.$y \in \pi(\alpha; X, N)$ 意味着:对过 $\alpha$ 的 $M$ 上的曲线 $p(t)$($p(0) = \alpha$,$\dot{p}(0) = y$,$-1 \leq t \leq 1$),沿 $p(t)$ 的法向量 $\{N(x)\}$ 不``摆动''.

(3) 令
\begin{equation*}
\pi(\alpha; X) = \bigcap_{N: \alpha \text{ 处的法向量}}\pi(\alpha; X, N).
\end{equation*}
由 (2) 的注意,$\pi(\alpha; X)$ 不依赖于 $X$ 的取法,所以 $\dim(\pi(\alpha; X))$ 也不依赖于坐标邻域.称 $\dim\pi(\alpha; X)$ 为 $M$ 在 $\alpha$ 处的 relative nullity(相对零化度),记为 $\nu(\alpha) = \dim\pi(\alpha; X)$.
\end{definition}

\begin{example}\label{exa:3.4.21}
当 $n = 3$ 时,即在 $\mathbf{R}^{3}$ 中如下:

(i) 当 $M = S^{2}$ 时,由 $\pi(\alpha; X, N) = \{0\}$ 得 $\nu(\alpha) = 0$($\alpha \in M$).

(ii) 当 $M = S^{1}$(平面的单位圆)时,$\nu(\alpha) = 0$($\alpha \in M$).

(iii) 当 $M$ 为圆柱面时,$\nu(\alpha) = 1$($\alpha \in M$).
\end{example}

\begin{definition}\label{def:3.4.22}
设 $M$ 是 $\mathbf{R}^{n}$ 的 $k$ 维流形,对 $\alpha \in M$ 设 $(X, \Omega)$ 为 $\alpha$ 处的坐标邻域系,$N_{0}$ 为 $\alpha$ 处的法向量.此时如下定义 $\alpha$ 处的曲率 $K(\alpha)$:
\begin{equation*}
K(\alpha) = \frac{\det(h_{ij}(\alpha; X, N_{0}))}{\det\left(\frac{\partial X}{\partial x_{i}}, \frac{\partial X}{\partial x_{j}}\right)}.
\end{equation*}
这里注意:因 $M$ 是 $k$ 维流形,分母不为零.又,$K(\alpha) = 0$ 的充要条件是分子为零,即 $\nu(\alpha) = k - \operatorname{rank}(h_{ij}) > 0$.
\end{definition}

以下设 $n = 3$,在 $\mathbf{R}^{3}$ 的空间中讨论.

\begin{definition}\label{def:3.4.23}
$F \subset \mathbf{R}^{2}$ 在 $y_{0} \in \mathbf{R}^{2}$ 具有 restricted cone property(限制锥性质),是指:对 $y_{0}$ 的邻域 $V_{0}$,某个 $0 < \sigma < 1$,$y_{1} \in \mathbf{R}^{2}$($\|y_{1}\| = 1$),令 $K = \{y \in \mathbf{R}^{2}: (1 - \sigma)\|y\| \leq y \cdot y_{1} \leq \sigma\}$ 时,$F \cap V_{0} - K \subset F$ 成立.其中 $\|z\|$ 表示 $z$ 与原点的距离.
\end{definition}

\begin{definition}\label{def:3.4.24}
(1) 紧集 $E(\subset M)$ 具有 restricted cone property(限制锥性质),是指:对任意 $\alpha \in E$,存在 $\alpha$ 的邻域 $V$,使得 $E \cap V$ 向 $\alpha$ 处切平面的投影在 $\alpha$ 具有 restricted cone property(限制锥性质).

(2) 令 $J_{1}(E) = \{\varphi \in C_{c}^{\infty}(\mathbf{R}^{3}): \varphi(x) = 0 \ (x \in E)\}$,对 $l = 1, 2, \ldots$,把 $J_{1}(E)^{l}$ 定义为由 $J_{1}(E)$ 的任意 $l$ 个元 $\varphi_{1}\cdots\varphi_{l}$ 生成的 $A(\mathbf{R}^{3})$ 的理想.
\end{definition}

于是,在下文中,关于 $\mathbf{R}^{3}$ 的子流形 $M$ 的具有 restricted cone property(限制锥性质)的紧子集是否为谱综合集,成立如下结果.

\begin{theorem}\label{thm:3.4.25}
设 $M$ 是 $\mathbf{R}^{3}$ 的 $k$ 维子流形($0 \leq k \leq 2$),relative nullity(相对零化度)恒为 $\nu$($0 \leq \nu \leq 2$).又设 $E \subset M$ 是具有 restricted cone property(限制锥性质)的紧集.此时,令 $d(E) = \min\{l \in \mathbb{N}: J_{1}(E)^{l} = J(E)\}$,则 $d(E) \leq \left[k - \frac{\nu}{2} + 1\right]$.特别地,当 $M$ 是超曲面时,$d(E) = \left[k - \frac{\nu}{2} + 1\right]$.
\end{theorem}

关于谱综合,利用定理~\ref{thm:3.4.25} 可证明作为 L. Schwartz 结果的推广的下面的事实.另外,这类讨论在 D. Müller 的论文 \cite{ch3:35} 中对高维情形作了研究.

\begin{theorem}\label{thm:3.4.26}
在定理~\ref{thm:3.4.25} 的条件下,进一步设 $M$ 是 2 维流形.此时,$E$ 是 $A(\mathbf{R}^{3})$ 的谱综合集的充要条件是 $\nu \geq 1$.特别地,$S^{2}$ 不是 $A(\mathbf{R}^{3})$ 的谱综合集.
\end{theorem}

% ============ 第3章 §3.5 关于作用函数(原书 pp.265-305)+ 结语(原书 p.305) ============
\section{关于作用函数}\label{sec:3.5}
\index{作用函数}
\index{乘子空间!乘子空间 Mp(Z) 的作用函数}
\index{作用!(乘子空间的情形)}
\index{Rudin--Shapiro 多项式}
\index{Riesz--Thorin 定理}
\index{插值定理!Stein 的}
\index{Hörmander 问题}
\index{Zafran 定理}


在 3.1 节中已给出交换 Banach 代数作用函数的定义,但在那种情形下,是把它作为极大理想空间上的函数空间上的作用来考虑的.本节中的作用函数,例如对交换 Banach 代数 $M(\mathbf{T})$,是通过把它看作真含于极大理想空间的 $\mathbf{Z}$ 上的函数空间来考虑的.这样,本节所处理的交换 Banach 代数的作用函数,均作为对 Fourier 变换上的函数空间的作用来考虑.通过研究函数空间的作用函数,可以看出交换 Banach 代数的特征.

\subsection{作用函数的定义}\label{sec:3.5.1}

\begin{definition}\label{def:3.5.1}
设 $X$ 是局部紧 Hausdorff 空间,$A$,$B$ 是 $X$ 上的函数空间.设 $G$ 为 $\mathbf{C}$ 的子集,$\Phi$ 为 $G$ 上的复值连续函数.称 $\Phi$ 从 $A$ 作用于 $B$,是指:若 $f \in A$,$f(X) \subset G$,则 $\Phi \circ f \in B$.此时称 $\Phi$ 为从 $A$ 到 $B$ 的作用函数.当 $A = B$ 时,简称为 $\Phi$ 作用于 $A$,或 $\Phi$ 是 $A$ 的作用函数.
\end{definition}

一般地,设 $\mathfrak{A}$ 为半单交换 Banach 代数,如 3.1 节所述,通过 Gelfand 表示,$\mathfrak{A}$ 可表示为某个局部紧 Hausdorff 空间 $X$($\mathfrak{A}$ 的极大理想空间)上的函数空间 $A = \widehat{\mathfrak{A}}$.因此,作用函数的研究也可在半单交换 Banach 代数上进行.本节所处理的交换 Banach 代数,是 Fourier 变换恰为 Gelfand 变换的情形.于是,基于 Fourier 变换的性质考虑作为 Fourier 变换的函数空间,并处理其上的作用函数.在这种情况下,借用 $A = \widehat{\mathfrak{A}}$ 的作用函数的说法,也称为 $\mathfrak{A}$ 的作用函数.

\subsection{\texorpdfstring{$A(\mathbf{T})$ 的作用函数}{A(T) 的作用函数}}\label{sec:3.5.2}

下面是 3.1 节定理~\ref{thm:3.1.48} 的特殊情形.但这里利用 $A(\mathbf{T})$ 的性质直接证明.

\begin{theorem}[Wiener--Lévy]\label{thm:3.5.2}
若 $f \in A(\mathbf{T})$,$\Phi$ 在 $f(\mathbf{T})$ 的邻域上解析,则 $\Phi(f) \in A(\mathbf{T})$.
\end{theorem}

\begin{proof}
设 $f$ 为 $\mathbf{T}$ 上的函数.若 $f$ 在 $\mathbf{T}$ 的任意点 $x_{0} \in T$ 的邻域上与 $A(\mathbf{T})$ 的元一致,则 $f \in A(\mathbf{T})$.因为由 $C^{1}(\mathbf{T}) \subset A(\mathbf{T})$,利用 3.1 节系~\ref{cor:3.1.39} 的想法即可容易地证明.其次,设 $\Phi$ 在 $\mathbf{C}$ 的开集 $G$ 上解析.此时,若 $f(\mathbf{T}) \subset G$,$f \in A(\mathbf{T})$,则 $\Phi(f) \in A(\mathbf{T})$.事实上,由前述事实,对 $a \in T = [-\pi, \pi)$,只要在 $a$ 的邻域上存在满足 $F_{a}(x) = \Phi(f(x))$ 的 $F_{a} \in A(\mathbf{T})$ 即可.设 $f(a) = c$.由 $\Phi$ 的解析性,在 $c$ 的邻域上,对适当的 $\varepsilon_{0} > 0$ 可作 Taylor 展开 $\Phi(z) = \sum_{n=0}^{\infty}a_{n}(z - c)^{n}$($|z - c| < \varepsilon_{0}$).在 3.3 节的记号下,把 $\varepsilon > 0$ 取得充分小,则在 $a$ 的适当邻域 $U(a)$ 上,$V_{\varepsilon}^{a}(x)(f(x) - c) = f(x) - c$,$\|V_{\varepsilon}^{a}(x)(f(x) - c)\|_{A(\mathbf{T})} < \varepsilon_{0}$.综上,$\Phi(f(x)) = \sum_{n=0}^{\infty}a_{n}(f(x) - c)^{n}$($x \in U(a)$),$V_{\varepsilon}^{a}(x)(f(x) - c) = f(x) - c$($x \in U(a)$).这里,令 $F_{a}(x) = \sum_{n=0}^{\infty}a_{n}(V_{\varepsilon}^{a}(x)(f(x) - c))^{n}$,则
\begin{equation*}
\sum_{n=0}^{\infty}|a_{n}|\|V_{\varepsilon}^{a}(x)(f(x) - c)\|_{A(\mathbf{T})}^{n} \leq \sum_{n=0}^{\infty}|a_{n}|\varepsilon_{0}^{n} < \infty,
\end{equation*}
故 $\Phi(f(x)) = F_{a}(x)$($x \in U(a)$),$F_{a} \in A(\mathbf{T})$.从而定理得证.
\end{proof}

1958 年,Katznelson \cite{ch3:31} 考虑了上述定理的逆命题,得到如下结果(cf. \cite{ch3:7}, \cite{ch3:8}, \cite{ch3:11}).今天,它被称为 Katznelson 定理.

\begin{theorem}[Katznelson]\label{thm:3.5.3}
设 $\Phi$ 为 $[-1, 1]$ 上的复值函数.假设对任意 $f \in A(\mathbf{T})$,$f(\mathbf{T}) \subset [-1, 1]$,$\Phi(f) \in A(\mathbf{T})$ 成立.此时,$\Phi$ 可延拓为 $[-1, 1]$ 的邻域上的解析函数.
\end{theorem}

证明分若干引理进行.在叙述引理之前,先注意作用函数 $\Phi$ 在 $[-1, 1]$ 上连续.事实上,只需证明在 $x = 0$ 处连续即可,设 $\Phi(0) = 0$ 并假设在 0 处不连续.此时,存在 $\eta, \delta > 0$ 和 $\{a_{n}\}_{n} \subset [-1, 1]$ 使得 $\sum_{n}|a_{n}| < \eta$,$|\Phi(a_{n})| > \delta$.这里,利用 3.3 节定义的函数 $\Delta_{\varepsilon}$,可知存在 $f \in A(\mathbf{T})$,$f(\frac{1}{n}) = a_{n}$($n = 2, 3, \ldots$),$f(\mathbf{T}) \subset [-1, 1]$.此时,因 $\Phi(f) \in A(\mathbf{T})$,$\lim_{n \to \infty}\Phi(f(\frac{1}{n})) = \Phi(f(0))$,可导出与 $\delta$,$\{a_{n}\}_{n}$ 的取法相矛盾.

\begin{lemma}\label{lem:3.5.4}
当 $[-1, 1]$ 上的连续函数 $\Phi$ 作用于 $A(\mathbf{T})$ 时,存在 $\varepsilon, C > 0$,使得实值函数 $f \in A(\mathbf{T})$ 若满足 $\|f\|_{A(\mathbf{T})} < \varepsilon$,则 $\|\Phi(f)\|_{A(\mathbf{T})} \leq C$.
\end{lemma}

\begin{proof}
\begin{sloppypar}
用反证法证明.若命题不成立,则存在实值函数的序列 $\{f_{n}\} \subset A(\mathbf{T})$ 使得 $\|f_{n}\|_{A(\mathbf{T})} \to 0$,$\|\Phi(f_{n})\|_{A(\mathbf{T})} \to \infty$($n \to \infty$).此时,若对任意点 $x \in \mathbf{T}$,存在含 $x$ 的某个闭区间 $I_{x}$ 使得 $\|\Phi(f_{n})\|_{A(I_{x})} \leq C$($C$ 是与 $x, f_{n}$ 无关的常数)(cf. 定义~\ref{def:3.4.9}),则由定理~\ref{thm:3.1.38} 与 $\|\Phi(f_{n})\|_{A(\mathbf{T})} \to \infty$($n \to \infty$)矛盾,故不妨设存在某点 $x_{0} \in \mathbf{T}$,使得对含 $x_{0}$ 的任意闭区间 $I$,$\|\Phi(f_{n})\|_{A(I)} \to \infty$($n \to \infty$).进而,不妨设 $x_{0} = 0$.又由定理~\ref{thm:3.3.1},把闭区间 $I$ 的长度记为 $|I|$,若 $|I| \to 0$,则 $\|\Phi(f_{n})\|_{A(I)} \to |\Phi(f_{n})(0)|$,所以固定 $n$ 时,$\{\|\Phi(f_{n})\|_{A(I)}\}_{I}$ 有界.从而,再次利用定理~\ref{thm:3.1.38} 的想法,对 $\|\Phi(f_{n})\|_{A(I)}$,把含 0 的闭区间 $J \subset I$ 取得足够小,令 $L$ 为 $I \setminus J$ 的闭包,则可做到
\end{sloppypar}
\begin{equation*}
\|\Phi(f_{n})\|_{A(L)} \geq \frac{1}{5}\|\Phi(f_{n})\|_{A(I)} - \|\Phi\|_{\infty}.
\end{equation*}
由此,取适当的数列 $\{n_{k}\}$,$\{n_{k}'\}$,令 $I_{k} = [-2^{-n_{k}}, 2^{-n_{k}}]$,$J_{k} = [-2^{-n_{k}'}, 2^{-n_{k}'}]$,$L_{k} = \overline{I_{k} \setminus J_{k}}$,则可做到 $\|\Phi(f_{n_{k}})\|_{A(L_{k})} \geq k$,$\|f_{n_{k}}\|_{A(\mathbf{T})} \leq 2^{-k}$.进而,令 $V_{k} = 2\Delta_{2^{-n_{k}+1}} - \Delta_{2^{-n_{k}}}$,在 $L_{k}$ 上 $V_{k} - V_{k+1} = 1$;若 $j \neq k$,则可取 $V_{j} - V_{j+1} = 0$.于是令 $f = \sum_{k=5}^{\infty}(V_{k} - V_{k+1})f_{n_{k}}$,由 $\|V_{k}\|_{A(\mathbf{T})} \leq 3$,$\|f\|_{A(\mathbf{T})} \leq 1$.又由 $f(x) = f_{n_{k}}(x)$($x \in L_{k}$),$\|\Phi(f)\|_{A(\mathbf{T})} \geq \|\Phi(f_{n_{k}})\|_{A(L_{k})} \geq k$($k = 5, 6, \ldots$),与 $\Phi(f) \in A(\mathbf{T})$ 矛盾.
\end{proof}

\begin{lemma}\label{lem:3.5.5}
若 $g \in A(\mathbf{T})$,$h \in A(\mathbf{T})$,则
\begin{equation*}
\lim_{k \to \infty}\|g(x)h(kx)\|_{A(\mathbf{T})} = \|g\|_{A(\mathbf{T})}\|h\|_{A(\mathbf{T})}
\end{equation*}
成立.
\end{lemma}

\begin{proof}
\begin{sloppypar}
$\lim_{k \to \infty}\|g(x)h(kx)\|_{A(\mathbf{T})} \leq \|g\|_{A(\mathbf{T})}\|h\|_{A(\mathbf{T})}$ 是显然的,故证明反向不等式.设 $h_{n}(x) = \sum_{|j| \leq n}\widehat{h}(j)e^{ijx}$,则对任意 $\varepsilon > 0$,存在 $N$ 使得 $\|h\|_{A(\mathbf{T})} - \varepsilon < \|h_{N}\|_{A(\mathbf{T})}$.此时,由 $g(x)h(kx) = g(x)h_{N}(kx) + g(x)(h(kx) - h_{N}(kx))$,
$\|g(x)h(kx)\|_{A(\mathbf{T})} \geq \|g(x)h_{N}(kx)\|_{A(\mathbf{T})} - \|g(x)(h(kx) - h_{N}(kx))\|_{A(\mathbf{T})}$,且
\end{sloppypar}
\begin{align*}
\|g(x)(h(kx) - h_{N}(kx))\|_{A(\mathbf{T})} &\leq \|g\|_{A(\mathbf{T})}\|h(kx) - h_{N}(kx)\|_{A(\mathbf{T})}\\
&\leq \varepsilon\|g\|_{A(\mathbf{T})}
\end{align*}
成立.因此
\begin{equation*}
\|g(x)h(kx)\|_{A(\mathbf{T})} \geq \|g(x)h_{N}(kx)\|_{A(\mathbf{T})} - \varepsilon\|g\|_{A(\mathbf{T})}
\end{equation*}
成立.又与 $\{h_{n}\}_{n}$ 同样地定义 $\{g_{n}\}_{n}$,则因 $g(x)h_{N}(kx) = g_{M}(x)h_{N}(kx) + (g(x) - g_{M}(x))h_{N}(kx)$,把 $M$ 取得足够大时,
\begin{equation*}
\|g(x)h_{N}(kx)\|_{A(\mathbf{T})} \geq \|g_{M}(x)h_{N}(kx)\|_{A(\mathbf{T})} - \varepsilon\|h\|_{A(\mathbf{T})}
\end{equation*}
成立.这里把 $k$ 取得足够大时,可证明
\begin{equation*}
\|g_{M}(x)h_{N}(kx)\|_{A(\mathbf{T})} \geq (\|g\|_{A(\mathbf{T})} - 2\varepsilon)(\|h\|_{A(\mathbf{T})} - 2\varepsilon),
\end{equation*}
故对 $\eta > 0$,只要 $k$ 足够大,
\begin{equation*}
\|g(x)h(kx)\|_{A(\mathbf{T})} \geq \|g\|_{A(\mathbf{T})}\|h\|_{A(\mathbf{T})} - \eta,
\end{equation*}
命题得证.
\end{proof}

\begin{lemma}\label{lem:3.5.6}
对 $\varepsilon > 0$,
\begin{equation*}
\|e^{i\varepsilon\cos x}\|_{A(\mathbf{T})} = 1 + \varepsilon + O(\varepsilon^{2})
\end{equation*}
成立.这里,$g(x) = O(x^{2})$($x \to 0$)表示 $g(x)/x^{2}$ 在 $x \to 0$ 时有界.
\end{lemma}

\begin{proof}
在
\begin{equation*}
e^{i\varepsilon\cos x} = \sum_{n=0}^{\infty}\frac{(i\varepsilon\cos x)^{n}}{n!}
\end{equation*}
中,由 $\|\cos x\|_{A(\mathbf{T})} \leq 1$ 得 $\|e^{i\varepsilon\cos x}\|_{A(\mathbf{T})} \leq 1 + \varepsilon + O(\varepsilon^{2})$.反之,由 $\|e^{i\varepsilon\cos x}\|_{A(\mathbf{T})} \geq 1 + \varepsilon - \left\|\sum_{n=2}^{\infty}\frac{(i\varepsilon\cos x)^{n}}{n!}\right\|_{A(\mathbf{T})}$ 及 $\|\cos x\|_{A(\mathbf{T})} \leq 1$ 得 $\|e^{i\varepsilon\cos x}\|_{A(\mathbf{T})} \geq 1 + \varepsilon - O(\varepsilon^{2})$.
\end{proof}

\begin{lemma}\label{lem:3.5.7}
对 $r > 0$,
\begin{equation*}
\sup\left\{\|e^{if}\|_{A(\mathbf{T})}: f \text{ 是实值函数, } \|f\|_{A(\mathbf{T})} \leq r\right\} = e^{r}
\end{equation*}
成立.
\end{lemma}

\begin{proof}
若 $\|f\|_{A(\mathbf{T})} \leq r$,则 $\|e^{if}\|_{A(\mathbf{T})} \leq e^{r}$,故证明反向不等式.令 $k_{1} = 1$,设 $f(x) = \frac{r}{n}(\cos k_{1}x + \cdots + \cos k_{n}x)$($k_{2}, \ldots, k_{n}$ 是自然数),由引理~\ref{lem:3.5.5},把 $k_{2}, \ldots, k_{n}$ 取得足够大,则
\begin{equation*}
\left\|\prod_{j=1}^{n}e^{i\frac{r}{n}\cos k_{j}x}\right\|_{A(\mathbf{T})} > \left(1 - \frac{1}{n}\right)\prod_{j=1}^{n}\left\|e^{i\frac{r}{n}\cos k_{j}x}\right\|_{A(\mathbf{T})}.
\end{equation*}
进而利用引理~\ref{lem:3.5.6},对任意 $\varepsilon > 0$,把 $n$ 取得足够大,则 $\left(1 - \frac{1}{n}\right)\left(1 + \frac{r}{n} + O\left(\frac{r^{2}}{n^{2}}\right)\right)^{n} > (1 - \varepsilon)e^{r}$.此时,得 $\|f\|_{A(\mathbf{T})} \leq r$ 且 $\|e^{if}\|_{A(\mathbf{T})} > (1 - \varepsilon)e^{r}$.
\end{proof}

\begin{proof}[定理~\ref{thm:3.5.3} 的证明]
如前所述,先注意 $\Phi$ 在 $[-1, 1]$ 上连续.首先证明 $\Phi(x)$ 在 $x = 0$ 的邻域上解析.由引理~\ref{lem:3.5.4},只要 $\alpha > 0$ 足够小,就存在常数 $C > 0$,使得当 $\|f\|_{A(\mathbf{T})} \leq 1$,$f$ 为实值函数,$-\pi \leq x < \pi$ 时,$\|\Phi(\alpha\sin(f(t) + x))\|_{A(\mathbf{T})} \leq C$.现在,令 $\Phi_{1}(x) = \Phi(\alpha\sin x)$,则由向量值函数的积分,对实值函数 $f \in A(\mathbf{T})$,
\begin{equation*}
\widehat{\Phi_{1}}(n)e^{inf(t)} = \frac{1}{2\pi}\int_{-\pi}^{\pi}\Phi_{1}(f(t) + x)e^{-inx}\dd x
\end{equation*}
成立.从而得 $\left|\widehat{\Phi_{1}}(n)\right|\|e^{inf}\|_{A(\mathbf{T})} \leq C$.此时,由引理~\ref{lem:3.5.7},$\left|\widehat{\Phi_{1}}(n)\right| \leq Ce^{-|n|}$,所以由众所周知的事实,$\Phi_{1}(z)$ 在 $|\mathrm{Im}z| < 1$ 上解析,利用 $\Phi(x) = \Phi_{1}\left(\arcsin\frac{x}{\alpha}\right)$,$\Phi$ 在 $x = 0$ 的邻域上解析.若 $a \in (-1, 1)$,令 $\Psi(x) = \Phi(x + a)$,取适当的 $\varepsilon > 0$,则 $\Psi$ 是在 $[-\varepsilon, \varepsilon]$ 上定义的 $A(\mathbf{T})$ 的作用函数.因此在 $x = 0$ 的邻域上解析,所以 $\Phi$ 在 $x = a$ 的邻域上解析.最后证明在 $x = \pm 1$(端点)处解析.令 $\Psi(x) = \Phi(1 - x^{2})$($x \in [-1, 1]$),因 $\Psi$ 是在 $[-1, 1]$ 上定义的 $A(\mathbf{T})$ 的作用函数,在 $x = 0$ 的邻域上解析.因 $\Psi$ 是偶函数,
\begin{equation*}
\Psi(x) = \sum_{n=0}^{\infty}a_{n}x^{2n}.
\end{equation*}
因此,$\Phi(1 - x) = \sum_{n=0}^{\infty}a_{n}x^{n}$,故对适当的 $\varepsilon_{0} > 0$,
\begin{equation*}
\Phi(x) = \sum_{n=0}^{\infty}a_{n}(1 - x)^{n} \quad (|x| < \varepsilon_{0})
\end{equation*}
可作 Taylor 展开.从而在 $x = 1$ 的邻域上解析.同样可证明在 $x = -1$ 的邻域上解析.
\end{proof}

\subsection{\texorpdfstring{$L^1(\mathbf{T})$ 与 $M(\mathbf{T})$ 的作用函数}{L1(T) 与 M(T) 的作用函数}}\label{sec:3.5.3}

这里参照 Rudin \cite{ch3:11}(cf. \cite{ch3:11} 的参考文献)叙述 $L^{1}(\mathbf{T})$ 与 $M(\mathbf{T})$ 的作用函数.

\begin{definition}\label{def:3.5.8}
$[-1, 1]$ 上的函数 $\Phi(x)$ 是 $L^{1}(\mathbf{T})$ 的作用函数,是指:令 $A(\mathbf{Z}) = \{\{\widehat{f}(n)\}_{n \in \mathbf{Z}}: f \in L^{1}(\mathbf{T})\}$,$\Phi$ 是 $A(\mathbf{Z})$ 上的作用函数.即,当 $\widehat{f} \in A(\mathbf{Z})$,$\widehat{f}(\mathbf{Z}) \subset [-1, 1]$ 时,存在 $\Phi \circ f \in L^{1}(\mathbf{T})$ 使得 $\widehat{\Phi \circ f} = \Phi(\widehat{f}) \in A(\mathbf{Z})$.又,$[-1, 1]$ 上的函数 $\Phi(x)$ 是 $M(\mathbf{T})$ 的作用函数,是指:令 $B(\mathbf{Z}) = \{\{\widehat{\mu}(n)\}_{n \in \mathbf{Z}}: \mu \in M(\mathbf{T})\}$,$\Phi$ 是 $B(\mathbf{Z})$ 上的作用函数.即,当 $\widehat{\mu} \in B(\mathbf{Z})$,$\widehat{\mu}(\mathbf{Z}) \subset [-1, 1]$ 时,存在 $\Phi \circ \mu \in M(\mathbf{T})$ 使得 $\widehat{\Phi \circ \mu} = \Phi(\widehat{\mu}) \in B(\mathbf{Z})$.把 $A(\mathbf{Z})$ 看作由 $L^{1}(\mathbf{T})$ 自然引入的范数及作为函数的通常运算所构成的交换 Banach 代数,$B(\mathbf{Z})$ 同样看作由 $M(\mathbf{T})$ 引入的交换 Banach 代数.
\end{definition}

\begin{theorem}\label{thm:3.5.9}
设 $\Phi$ 为在 $[-1, 1]$ 上定义的函数.此时,$\Phi$ 是 $L^{1}(\mathbf{T})$ 的作用函数的充要条件是:$\Phi(0) = 0$ 且 $\Phi(x)$ 在 0 的邻域上解析.
\end{theorem}

该定理的充分条件如下证明.设 $\Phi$ 在 $|x| < \varepsilon$($\varepsilon > 0$)上解析.取 $f \in L^{1}(\mathbf{T})$,$\hat{f}(\mathbf{Z}) \subset [-1, 1]$,则对任意 $\varepsilon > 0$,取适当的三角多项式 $P$,$Q$,可使 $\Phi(\widehat{f}(n)) = \widehat{Q}(n) + \Phi(\widehat{f}(n) - \widehat{P}(n))$($n \in \mathbf{Z}$),$\|f - P\|_{1} < \varepsilon$.因此由 $\Phi$ 的假设,$\Phi(\widehat{f}(n) - \widehat{P}(n)) \in A(\mathbf{Z})$,故 $\Phi$ 是 $L^{1}(\mathbf{T})$ 的作用函数.

必要性的证明如下进行.首先,因常数函数是 $L^{1}(\mathbf{T})$ 的元,若 $\Phi$ 作用于 $L^{1}(\mathbf{T})$,由 Riemann--Lebesgue 引理得 $\Phi(0) = 0$.

其次,叙述两个引理.

\begin{lemma}\label{lem:3.5.10}
$\Phi(x)$ 在 $x = 0$ 处连续.
\end{lemma}

\begin{proof}
设 $0 < \eta < 1$.若 $\Phi$ 在 0 处不连续,则存在 $\{a_{n}\} \subset R$,$\delta > 0$ 使得 $\sum_{n=1}^{\infty}2|a_{n}| < \eta$,$|\Phi(a_{n})| > \delta$.此时,令 $k_{n} = 5^{n}$,$f(x) = \sum_{n=1}^{\infty}2a_{n}\cos k_{n}x$,则 $\Phi(\widehat{f}) \in A(\mathbf{Z}) \subset B(\mathbf{Z})$.因此存在 $\mu \in M(\mathbf{T})$ 使得
\begin{equation*}
\widehat{\mu}(l) = \left\{\begin{array}{ll}
\Phi(a_{n}) & (l = \pm k_{n})\\
0 & (l \neq \pm k_{n})
\end{array}\right.
\end{equation*}
于是,令 $d\mu_{n}(x) = e^{-ik_{n}x}\dd\mu(x)$,设 $\sigma \in M(\mathbf{T})$ 为 $\{\mu_{n}\} \subset M(\mathbf{T})$ 关于 $*$-弱拓扑的一个聚点.由 $|\widehat{\mu_{n}}(0)| = |\Phi(a_{n})| > \delta$,$\sigma \neq 0$.又对 $l \neq 0$,$l \in \mathbf{Z}$,$\widehat{\mu_{n}}(l) = \widehat{\mu}(k_{n} + l)$,故注意 $\widehat{\mu_{n}}(l) \neq 0$ 只在 $l = -2k_{n}$ 时发生.因此 $\widehat{\sigma}(l) = 0$($l \neq 0$,$l \in \mathbf{Z}$).从而 $\sigma$ 是非零的常数函数.可是,由 Riemann--Lebesgue 引理,对任意三角多项式 $f(x)$,
\begin{equation*}
\int_{0}^{2\pi}f(x)\dd\sigma(x) = \lim_{n \to \infty}\int_{0}^{2\pi}f(x)\dd\mu_{n}(x) = \lim_{n \to \infty}\int_{0}^{2\pi}f(x)e^{-ik_{n}x}\dd\nu(x)
\end{equation*}
成立.其中,$\nu$ 是 $\mu$ 集中于 Lebesgue 测度为零的集合上的部分(奇异测度部分).由此,利用 $\sigma, \nu \in M(\mathbf{T})$,对可测集 $E$ 有 $|\sigma|(E) \leq |\nu|(E)$.这与 $\sigma$ 是非零常数函数矛盾.
\end{proof}

\begin{lemma}\label{lem:3.5.11}
设 $\Phi$ 为 $[-\eta, \eta]$($0 < \eta < 1$)上的,作用于 $L^{1}(\mathbf{T})$ 的函数.此时,存在 $0 < \varepsilon_{0} < 1$,$C > 0$,使得若 $\|\sigma\|_{M(\mathbf{T})} < \varepsilon_{0}$,$\widehat{\sigma}(\mathbf{Z}) \subset \mathbf{R}$,则 $\Phi \circ \widehat{\sigma} \in B(\mathbf{Z})$,$\|\Phi \circ \sigma\|_{M(\mathbf{T})} \leq C$.其中设 $\widehat{\Phi \circ \sigma}(n) = \Phi(\widehat{\sigma}(n))$.
\end{lemma}

\begin{proof}
只需对三角多项式 $f$(代替引理中的 $\sigma$)证明即可.事实上,设对三角多项式已证明.设 $\|\sigma\|_{M(\mathbf{T})} < \varepsilon_{0}$,$\widehat{\sigma}(\mathbf{Z}) \subset \mathbf{R}$.此时,取三角多项式列 $\{k_{n}\}$ 满足 $\hat{k_{n}}(m) = 1$($|m| \leq n$),$\|k_{n}\|_{1} = 1$,则 $\|\Phi \circ (\sigma * k_{n})\|_{M(\mathbf{T})} \leq C$.$\{\Phi \circ (\sigma * k_{n})\}$ 关于 $*$-弱拓扑的极限是 $\Phi \circ \sigma$,故得结果.于是,假设存在三角多项式列 $\{f_{n}\}$ 使得 $\widehat{f_{n}}(\mathbf{Z}) \subset [-1, 1]$,$\|f_{n}\|_{1} < 2^{-n}$,$\|\Phi \circ f_{n}\|_{M(\mathbf{T})} > n$($n \in \mathbf{N}$).这里,选取 $\mathbf{T}$ 上的三角多项式 $h_{n}$ 作为三角多项式列 $\{h_{n}\}$:$\|h_{n}\|_{1} < 2$,在 $\operatorname{supp}\widehat{f_{n}}$ 上为 1,且 $0 \leq \widehat{h_{n}} \leq 1$.此时,注意 Fourier 系数的平移不改变函数的范数以及 $\Phi(0) = 0$,不妨设 $\operatorname{supp}\widehat{h_{n}} \cap \operatorname{supp}\widehat{h_{m}} = \emptyset$($n \neq m$),并令 $g = \sum_{n=1}^{\infty}f_{n}$,则 $g \in L^{1}(\mathbf{T})$,且 $h_{n} * (\Phi \circ g) = h_{n} * (\Phi \circ f_{n}) = \Phi \circ f_{n}$,$\|\Phi \circ f_{n}\|_{1} = \|h_{n} * (\Phi \circ g)\|_{1} \leq 2\|\Phi \circ g\|_{M(\mathbf{T})} < \infty$.这与 $\|\Phi \circ f_{n}\|_{M(\mathbf{T})} \to \infty$($n \to \infty$)矛盾.
\end{proof}

\begin{proof}[定理~\ref{thm:3.5.9} 的证明]
先注意 $\Phi(0) = 0$.设 $\Phi$ 作用于 $A(\mathbf{Z})$.对 $\varepsilon_{0} > 0$ 和 $a$($|a| < \varepsilon_{0}$),令 $\Phi_{1}(x) = \Phi(x + a) - \Phi(a)$,由引理~\ref{lem:3.5.11},当 $\widehat{f} \in A(\mathbf{Z})$,$\widehat{f}(\mathbf{Z}) \subset [-\eta, \eta]$($0 < \eta < 1$),$\|f\|_{1} < \varepsilon_{0} - |a|$ 时,$\Phi_{1}(\widehat{f}) \in B(\mathbf{Z})$.此时,利用引理~\ref{lem:3.5.10} 的想法,$\Phi_{1}$ 在 $x = 0$ 处连续.因此 $\Phi$ 在 $x = a$ 处连续,得 $\Phi$ 在 $(-\varepsilon_{0}, \varepsilon_{0})$ 上连续.于是,取 $0 < r < \varepsilon_{0}/e$,令 $\Phi_{2}(x) = \Phi(r\sin x)$($x \in \mathbf{R}$),则 $\Phi_{2}$ 是 $\mathbf{R}$ 上的 $2\pi$ 周期连续函数.现在,取 $\mu \in M(\mathbf{T})$ 满足 $\|\mu\|_{M(\mathbf{T})} \leq 1$,$\widehat{\mu}(\mathbf{Z}) \subset \mathbf{R}$,则对 $a \in \mathbf{R}$,
\begin{align*}
&\|\sin(\widehat{\mu} + a)\|_{B(\mathbf{Z})} \leq |\cos a|\|\sin\widehat{\mu}\|_{B(\mathbf{Z})} + |\sin a|\|\cos\widehat{\mu}\|_{B(\mathbf{Z})}\\
&\qquad \leq e^{\|\widehat{\mu}\|_{B(\mathbf{Z})}} = e^{\|\mu\|_{M(\mathbf{T})}} \leq e
\end{align*}
成立.因此,若 $a \in R$,$\widehat{\mu}(\mathbf{Z}) \subset \mathbf{R}$,$\|\mu\|_{M(\mathbf{T})} \leq 1$,则由 $\|r\sin(\widehat{\mu} + a)\|_{B(\mathbf{Z})} \leq re \leq \varepsilon_{0}$,利用引理~\ref{lem:3.5.11} 得 $\|\Phi_{2}(\widehat{\mu} + a)\|_{B(\mathbf{Z})} \leq \|\Phi(r\sin(\widehat{\mu} + a))\|_{B(\mathbf{Z})} \leq C$.这里,因 $\Phi_{2}$ 是连续函数,
\begin{equation*}
\Phi_{2}(x) \sim \sum_{-\infty}^{\infty}c_{n}e^{inx}
\end{equation*}
作 Fourier 级数展开.取注意~\ref{note:3.3.12} 所述的,作为 Cantor 集的 Kronecker 集 $K$,设 $\mu \in M(K \cup (-K))$ 满足任意 1 点的测度为 0,$\widehat{\mu}(\mathbf{Z}) \subset \mathbf{R}$,$0 < \|\mu\|_{M(\mathbf{T})} \leq 1$.因 $\Phi_{2}$ 连续,利用向量值函数的积分(cf. \cite{ch3:3}, \cite{ch3:8}),估计
\begin{equation*}
c_{n}e^{in\widehat{\mu}(j)} = \lim_{N \to \infty}\frac{1}{N}\sum_{k=1}^{N}\Phi_{2}\left(\frac{2\pi k}{N} + \widehat{\mu}(j)\right)e^{-i2\pi k/N},
\end{equation*}
得 $|c_{n}|\|e^{in\mu}\|_{M(\mathbf{T})} \leq C$($n \in \mathbf{Z}$).由 $\mu$ 的性质,$\|e^{in\mu}\|_{M(\mathbf{T})} = e^{|n|\|\mu\|_{M(\mathbf{T})}}$($n \in \mathbf{Z}$),故得 $|c_{n}| \leq Ce^{-|n|}$($n \in \mathbf{Z}$).从而 $\Phi_{2}(z)$ 在 $|\mathrm{Im}z| < 1$ 上解析.因此 $\Phi(x) = \Phi_{2}\left(\arcsin\frac{x}{r}\right)$ 在 0 的邻域上解析.
\end{proof}

关于 $M(\mathbf{T})$ 的作用函数,下面的事实成立.

\begin{theorem}\label{thm:3.5.12}
$[-1, 1]$ 上的复值函数 $\Phi$ 是 $M(\mathbf{T})$ 的作用函数的充要条件是:$\Phi$ 可延拓为整函数.
\end{theorem}

首先证明下面的事实.

\begin{lemma}\label{lem:3.5.13}
若 $\Phi_{1}$ 是 $\mathbf{R}$ 上的 $2\pi$ 周期连续函数且作用于 $B(\mathbf{Z})$,则当 $\widehat{\mu} \in B(\mathbf{Z})$,$\widehat{\mu}(\mathbf{Z}) \subset \mathbf{R}$ 时,存在 $\delta > 0$,$C < \infty$ 使得
\begin{equation*}
\|\Phi_{1}(\widehat{\mu} + \widehat{\sigma})\|_{B(\mathbf{Z})} \leq C \quad (\sigma \in M(\mathbf{T}),\ \widehat{\sigma} \in \mathbf{R},\ \|\sigma\|_{M(\mathbf{T})} < \delta).
\end{equation*}
\end{lemma}

\begin{proof}
不妨设 $\Phi_{1}(0) = 0$.只需证明 $\|\Phi_{1}(\widehat{\mu} + \widehat{f})\|_{B(\mathbf{Z})} \leq C$($f \in L^{1}(\mathbf{T})$,$\widehat{f} \in \mathbf{R}$,$\|f\|_{1} < \delta$).因为,由 $L^{1}(\mathbf{T})$ 是 $M(\mathbf{T})$ 的理想以及 $M(\mathbf{T})$ 的闭单位球关于 $*$-弱拓扑是紧集即可证明.现在,假设上述事实不成立.此时,存在 $\mathbf{T}$ 上取实值的三角多项式列 $\{f_{n}\}$ 使得 $\|f_{n}\|_{1} \to 0$,$\|\Phi_{1} \circ (\widehat{\mu} + \widehat{f_{n}})\|_{B(\mathbf{Z})} \to \infty$($n \to \infty$).为导出矛盾,如下构造 $\{\lambda_{j}\} \subset M(\mathbf{T})$.令 $n_{1} = 1$,设 $n_{1}, \ldots, n_{j}$ 与三角多项式列 $f_{n_{1}}, \ldots, f_{n_{j}}$ 已选出,令 $\lambda_{j} = \mu + f_{n_{1}} + \cdots + f_{n_{j}}$.此时,取次数足够高的 Fejér 核 $k_{j}(x) = \sum_{|m| \leq n_{j}}\left(1 - \frac{|m|}{n_{j} + 1}\right)e^{imx}$,则
\begin{equation*}
\|k_{j} * (\Phi_{1} \circ \lambda_{j})\|_{M(\mathbf{T})} > \frac{1}{2}\|\Phi_{1} \circ \lambda_{j}\|_{M(\mathbf{T})}.
\end{equation*}
因 $\Phi_{1}$ 连续,把 $n_{j+1}$ 取得足够大,可做到
\begin{align*}
\|f_{n_{j+1}}\|_{1} &< 2^{-j}, \quad \|\Phi_{1} \circ (\lambda_{j} + f_{n_{j+1}})\|_{M(\mathbf{T})} \geq j + 1,\\
\left\|k_{i} * \left((\Phi_{1} \circ \lambda_{j}) - \Phi_{1} \circ (\lambda_{j} + f_{n_{j+1}})\right)\right\|_{1} &< 2^{-j} \quad (1 \leq i \leq j).
\end{align*}
这样归纳地取 $\{n_{j}\}$ 与 $\{k_{j}\}$ 的序列,令 $\lambda = \mu + \sum_{j=1}^{\infty}f_{n_{j}}$,$\tau = \Phi_{1} \circ \lambda$.这里注意 $\widehat{\tau}(\gamma) = \lim_{N \to \infty}\Phi_{1}(\widehat{\lambda_{N}}(\gamma))$,考虑到 $k_{s}$ 是 Fejér 核,
\begin{equation*}
\lim_{N \to \infty}\|k_{s} * (\tau - \Phi_{1} \circ \lambda_{N})\|_{M(\mathbf{T})} = 0 \quad (s = 1, 2, \ldots)
\end{equation*}
成立.因此
\begin{equation*}
k_{s} * (\tau - \Phi_{1} \circ \lambda_{s}) = \lim_{N \to \infty}\sum_{j=s}^{N-1}k_{s} * (\Phi_{1} \circ \lambda_{j+1} - \Phi_{1} \circ \lambda_{j}),
\end{equation*}
所以
\begin{equation*}
\|k_{s} * (\tau - \Phi_{1} \circ \lambda_{s})\|_{M(\mathbf{T})} \leq \sum_{j=s}^{\infty}2^{-j} \leq 1 \quad (s \in \mathbf{N}).
\end{equation*}
从而
\begin{align*}
2\|\tau\|_{M(\mathbf{T})} &\geq \|k_{s} * \tau\|_{M(\mathbf{T})}\\
&\geq \|k_{s} * (\Phi_{1} \circ \lambda_{s})\|_{M(\mathbf{T})} - 1\\
&\geq \frac{1}{2}\|\Phi_{1} \circ \lambda_{s}\|_{M(\mathbf{T})} - 1 \geq \frac{1}{2}s - 1.
\end{align*}
这是矛盾.
\end{proof}

\begin{proof}[定理~\ref{thm:3.5.12} 的证明]
只需证明必要条件.取 $0 < r < 1$,令 $\Phi_{3}(x) = \Phi(r\sin x)$,则 $\Phi_{3}$ 是在 $\mathbf{R}$ 上定义的 $2\pi$ 周期函数,且是 $M(\mathbf{T})$ 的作用函数.此时,利用引理~\ref{lem:3.5.10} 可证明 $\Phi_{3}$ 是 $\mathbf{R}$ 上的连续函数.又,把 $\Phi_{3}$ 的 Fourier 级数记为 $\Phi_{3} \sim \sum_{n}c_{n}e^{inx}$.此时,与定理~\ref{thm:3.5.9} 的证明一样利用引理~\ref{lem:3.5.13},取适当的 $\mu \in M(\mathbf{T})$(cf. 注意~\ref{note:3.3.12}),则
\begin{equation*}
|c_{n}| \leq Ce^{-|n|\|\mu\|_{M(\mathbf{T})}} \quad (n \in \mathbf{Z})
\end{equation*}
成立.由此,$\Phi_{3}(z)$ 在 $|\mathrm{Im}z| < \|\mu\|_{M(\mathbf{T})}$ 上解析,且 $\|\mu\|_{M(\mathbf{T})}$ 可以取得任意大,故 $\Phi_{3}$ 与某个整函数一致.现在,令 $\Psi_{j}(x) = \Phi(r_{j}\sin x)$($j = 1, 2$,$0 < r_{1} < r_{2} < 1$),如前所述,$\Psi_{j}$ 在 $\mathbf{R}$ 上与某个整函数一致.由此,对 $j = 1, 2$,
\begin{equation*}
\Phi(x) = \Psi_{j}\left(\arcsin\frac{x}{r_{j}}\right) \quad (-r_{j} < x < r_{j}),
\end{equation*}
$\Phi(x)$ 在 0 的邻域上与解析函数一致.又,利用多值函数 $\arcsin$ 作解析延拓,可证明 $\Phi$ 在 $[-1, 1]$ 上与某个整函数一致.
\end{proof}

\begin{note}\label{note:3.5.14}
由定理~\ref{thm:3.5.12} 也可导出 $M(\mathbf{T})$ 的 Wiener--Pitt 现象.事实上,令 $\Phi(z) = \frac{1}{1 + z^{2}}$,因该函数不是整函数,故不作用于 $B(\mathbf{Z})$.因此,可知存在 $\mu \in M(\mathbf{T})$ 使得 $\hat{\mu}(n) \in \mathbf{R}$($n \in \mathbf{Z}$),$i \in \mathrm{Sp}(\mu, M(\mathbf{T}))$.
\end{note}

\subsection{\texorpdfstring{乘子空间 $M_p(\mathbb{Z})$ 的作用函数}{乘子空间 M_p(Z) 的作用函数}}\label{sec:3.5.4}

这里叙述乘子空间 $M_{p}(\mathbf{Z})$ 的作用函数.$[-1, 1]$ 上的函数 $\Phi(x)$ 是 $M_{p}(\mathbf{Z})$ 上的作用函数,是指:当 $\hat{T} \in M_{p}(\mathbf{Z})$,$\hat{T}(\mathbf{Z}) \subset [-1, 1]$ 时,存在 $\Phi \circ T \in \mathfrak{M}_{p}(\mathbf{T})$ 使得 $\widehat{\Phi \circ T} = \Phi(\hat{T}) \in M_{p}(\mathbf{Z})$.如 3.2 节所见,$\mathfrak{M}_{p}(\mathbf{T})$ 关于算子范数成为交换 Banach 代数.因此,$M_{p}(\mathbf{Z}) = \{\hat{T} \in \ell^{\infty}: T \in \mathfrak{M}_{p}(\mathbf{T})\}$ 关于由 $\mathfrak{M}_{p}(\mathbf{T})$ 引入的范数与函数的运算成为含单位元的交换 Banach 代数.如前所见,$\mathfrak{M}_{1}(\mathbf{T}) = M(\mathbf{T})$(或 $M_{1}(\mathbf{Z}) = B(\mathbf{Z})$),$M_{2}(\mathbf{Z}) = \ell^{\infty}$,$\mathfrak{M}_{1}(\mathbf{T}) \subset \mathfrak{M}_{p}(\mathbf{T}) \subset \mathfrak{M}_{2}(\mathbf{T})$($1 \leq p \leq 2$)成立.从而,当 $T \in \mathfrak{M}_{p}(\mathbf{T})$ 时,$\overline{\hat{T}(\mathbf{Z})} \subset \mathrm{Sp}(T, \mathfrak{M}_{p}(\mathbf{T}))$.特别是,若 $f \in L^{1}(\mathbf{T})$,则关于算子 $T_{f}(g) = f * g$ 有 $\mathrm{Sp}(T_{f}, \mathfrak{M}_{p}(\mathbf{T})) = \hat{f}(\mathbf{Z}) \cup \{0\}$($1 \leq p \leq 2$).

1960 年,Hörmander \cite{ch3:22} 把 $m_p$ 取为由 $\{T_f : f \in L^1(\mathbf{T})\}$ 在 $\mathfrak{M}_p(\mathbf{T})$ 中的闭包构成的交换 Banach 代数,证明了下面的事实:

\begin{proposition}\label{prop:3.5.15}
(1) $\Delta(m_p) = \mathbf{Z}$.

(2) 在 0 的邻域上解析的函数 $\Phi$ 若满足 $\phi(0) = 0$,则作用于 $m_{p}$.

(3) 对 $1 < p < 2$,令 $C_{0}\mathfrak{M}_{p}(\mathbf{T}) = \{T \in \mathfrak{M}_{p}(\mathbf{T}): \hat{T}(n) \to 0 \ (|n| \to \infty)\}$.设 $1 < p < q < 2$,若满足 $\Phi(0) = 0$ 的函数 $\Phi$ 在 0 的邻域上解析,则 $\Phi$ 从 $C_{0}\mathfrak{M}_{p}(\mathbf{T})$ 作用于 $\mathfrak{M}_{q}(\mathbf{T})$.
\end{proposition}

又,他还提出了如下问题:``当 $1 < p < 2$ 时,$C_{0}\mathfrak{M}_{p}(\mathbf{T}) = m_{p}$ 是否成立?''Igari \cite{ch3:25} 对命题~\ref{prop:3.5.15}(3) 的 $p = q$ 的情形作了相应研究,决定了 $\mathfrak{M}_{p}(\mathbf{T})$($= M_{p}(\mathbf{Z})$)的作用函数.这里叙述 S. Igari 的结果.另外,Hörmander 问题由 M. Zafran 解决一事,将在 3.5.6 小节补充.

下面的插值定理省略证明(cf. \cite{ch3:4}).

\begin{theorem}[Riesz--Thorin]\label{thm:3.5.16}
设 $(X, \mu)$,$(Y, \nu)$ 为 $\sigma$-有限的测度空间.设 $1 \leq p_{0}, q_{0}, p_{1}, q_{1} \leq \infty$($p_{0} \neq p_{1}$,$q_{0} \neq q_{1}$),$T$ 为从 $X$ 上简单函数的空间到 $Y$ 上可测函数空间的线性算子,且满足:存在常数 $M_{0}, M_{1} > 0$,对 $X$ 上的简单函数 $f(x)$ 有
\begin{equation*}
\|Tf\|_{q_{j}} \leq M_{j}\|f\|_{p_{j}} \quad (j = 0, 1).
\end{equation*}
此时,对 $0 < \theta < 1$,$\|Tf\|_{q} \leq M_{0}^{1-\theta}M_{1}^{\theta}\|f\|_{p}$ 成立.其中 $p, q$ 由
\begin{equation*}
\frac{1}{p} = \frac{1 - \theta}{p_{0}} + \frac{\theta}{p_{1}}, \quad \frac{1}{q} = \frac{1 - \theta}{q_{0}} + \frac{\theta}{q_{1}}
\end{equation*}
给出.
\end{theorem}

下面的定理是 de Leeuw \cite{ch3:17}(cf. \cite{ch3:27})的结果.

\begin{theorem}\label{thm:3.5.17}
设 $1 < p < 2$.设 $m$ 是 $\mathbf{R}$ 上具有 $2\pi$ 周期的函数,$m \in L^{\infty}(\mathbf{R})$.此时,对任意 $\{a_{n}\} \in \ell^{p}$,令 $\hat{f}(x) = \sum a_{n}e^{inx}$,存在某个 $\{b_{n}\} \in \ell^{p}$ 使 $m(x)\hat{f}(x) = \sum_{n \in \mathbf{Z}}b_{n}e^{inx}$ 成立的充要条件是:对任意 $L^{p}(\mathbf{R})$ 的元 $g$ 的 Fourier 变换 $\hat{g}$,$m\hat{g}$ 是 $L^{p}(\mathbf{R})$ 的 Fourier 变换.其中,设 $\ell^{p} = \{a = \{a_{n}\}_{n=-\infty}^{\infty}: \|a\|_{\ell^{p}} = (\sum_{n \in \mathbf{Z}}|a_{n}|^{p})^{1/p} < \infty\}$,$\hat{a}(x) = \sum_{n}a_{n}e^{inx}$.
\end{theorem}

\begin{proof}
证明条件是必要的.对 $a = \{a_{n}\} \in \ell^{p}$,令 $\widehat{T_{m}}a(x) = m(x)\widehat{a}(x)$,由闭图像定理,$T_{m}$ 是从 $\ell^{p}$ 到 $\ell^{p}$ 的有界线性算子.把这样的 $T_{m}$ 的全体记为 $\mathfrak{M}_{p}(\mathbf{Z})$.又把 $T_{m}$ 的算子范数记为 $\|T_{m}\|_{\mathfrak{M}_{p}(\mathbf{Z})}$(或 $\|m\|_{M_{p}(\mathbf{T})}$).于是令 $m_{n} = \frac{1}{2\pi}\int_{-\pi}^{\pi}m(x)e^{-inx}\dd x$($n \in \mathbf{Z}$),进一步对 $f \in L^{p}(\mathbf{R})$ 令 $Kf(x) = \sum_{n}m_{n}f(x + n)$.这里,
\begin{equation*}
\infty > \int_{-\infty}^{\infty}|f(x)|^{p}\dd x = \int_{-1/2}^{1/2}\sum_{n \in \mathbf{Z}}|f(x + n)|^{p}\dd x
\end{equation*}
故 $\{f(x + n)\} \in \ell^{p}$(a.e. $x \in [-\frac{1}{2}, \frac{1}{2}]$).又由 $T_{m}$ 是 $\ell^{p}$ 上的有界算子以及 Banach 空间 $\ell^{p}$ 上的连续线性泛函的空间是 Banach 空间 $\ell^{p'}$($1/p + 1/p' = 1$),得 $\{m_{n}\} \in \ell^{p'}$.因此由 Hölder 不等式,
\begin{align*}
|Kf(x)| &\leq \sum_{n}|m_{n}||f(x + n)| \leq \left(\sum|m_{n}|^{p'}\right)^{1/p'}\left(\sum|f(x + n)|^{p}\right)^{1/p}\\
&< \infty \quad \left(\text{a.e. } x \in \left[-\frac{1}{2}, \frac{1}{2}\right]\right).
\end{align*}
又
\begin{equation*}
\begin{aligned}
\int_{-\infty}^{\infty}|Kf(x)|^{p}\dd x &= \int_{[-\frac{1}{2}, \frac{1}{2}]}\sum_{n \in \mathbf{Z}}|Kf(x + n)|^{p}\dd x\\
&= \int_{[-\frac{1}{2}, \frac{1}{2}]}\sum_{n \in \mathbf{Z}}\left|\sum_{k}m_{k}f(x + n + k)\right|^{p}\dd x.
\end{aligned}
\end{equation*}
这里利用 $T_{m}$ 是从 $\ell^{p}$ 到 $\ell^{p}$ 的有界线性算子,
\begin{align*}
&\int_{-\infty}^{\infty}|Kf(x)|^{p}\dd x \leq \|T_{m}\|_{\mathfrak{M}_{p}(\mathbf{Z})}^{p}\int_{[-\frac{1}{2}, \frac{1}{2}]}\left(\sum_{n}|f(x + n)|^{p}\right)^{p/p}\dd x\\
&= 2\pi\|T_{m}\|_{\mathfrak{M}_{p}(\mathbf{Z})}^{p}\|f\|_{L^{p}(\mathbf{R})}^{p}.
\end{align*}
因此 $K$ 是 $L^{p}(\mathbf{R})$ 上的有界线性算子.其次证明 $\widehat{Kf}(\xi) = m(\xi)\widehat{f}(\xi)$($f \in C_{c}^{\infty}(\mathbf{R})$).设 $j$ 为自然数,令 $s_{j}(x) = 2\pi j\Delta_{1}(jx)$($\Delta_{1}$ 是定理~\ref{thm:3.3.1} 的记号),$\sigma_{j}(x) = \sum_{n \in \mathbf{Z}}s_{j}(x + 2n\pi)$.此时,因 $m$ 是 $2\pi$ 周期函数,
\begin{align*}
\sigma_{j} * m(x) &= \frac{1}{2\pi}\int_{-\pi}^{\pi}\sigma_{j}(y)m(x - y)\dd y\\
&= \frac{1}{2\pi}\int_{-\pi}^{\pi}\sum_{n}s_{j}(y + 2n\pi)m(x - y)\dd y\\
&= \frac{1}{2\pi}\int_{-\infty}^{\infty}s_{j}(y)m(x - y)\dd y = s_{j} * m(x).
\end{align*}
此时,不妨设 $s_{j} * m(x) \to m(x)$ a.e..事实上,对 $0 < a < b < \infty$,因 $m \in L^{1}([a, b])$,
\begin{equation*}
\int_{-\infty}^{\infty}|m(x - y) - m(x)|\chi_{[a, b]}(x)\dd x \to 0 \quad (y \to 0).
\end{equation*}
因此
\begin{align*}
\int_{-\infty}^{\infty}\left|s_{j} * m(x)\chi_{[a, b]}(x) - m(x)\chi_{[a, b]}(x)\right|\dd x&\leq \int_{-\infty}^{\infty}\int_{-\infty}^{\infty}s_{j}(y)|m(x - y) - m(x)|\chi_{[a, b]}(x)\dd y\dd x\\
&= \int_{-\infty}^{\infty}s_{j}(y)\left(\int_{-\infty}^{\infty}|m(x - y) - m(x)|\chi_{[a, b]}(x)\dd x\right)\dd y \to 0 \quad (j \to \infty).
\end{align*}
从而,存在子列 $\{j_{k}\}$ 使得 $s_{j_{k}} * m(x) \to m(x)$ a.e..这里利用 $\|\sigma_{j}\|_{L^{1}(\mathbf{T})} = 1$ 计算,由
\begin{equation*}
\|T_{\sigma_{j} * m}\|_{\mathfrak{M}_{p}(\mathbf{Z})} = \|\sigma_{j} * m\|_{M_{p}(\mathbf{T})} \leq \|\sigma_{j}\|_{L^{1}(\mathbf{T})}\|m\|_{M_{p}(\mathbf{T})} = \|m\|_{M_{p}(\mathbf{T})},
\end{equation*}
注意 $T_{\sigma_{j} * m}$ 是从 $\ell^{p}$ 到 $\ell^{p}$ 的有界线性算子.由此,设 $s_{j, n} = \widehat{\sigma}_{j}(n)$,
\begin{equation*}
K_{j}f(x) = \sum_{n}s_{j, n}m_{n}f(x + n),
\end{equation*}
则 $K_{j}$ 与 $K$ 一样是从 $L^{p}(\mathbf{R})$ 到 $L^{p}(\mathbf{R})$ 的有界线性算子.此时
\begin{equation*}
|K_{j}f(x) - Kf(x)| \leq \sum_{n}|s_{j, n} - 1||m_{n}f(x + n)|.
\end{equation*}
这里,由
\begin{equation*}
s_{j, n} = \left(\sin\frac{n}{2j}\right)^{2}\left(\frac{n}{2j}\right)^{-2}, \quad \sum_{n}|m_{n}f(x + n)| < \infty \ \text{ a.e. } x,
\end{equation*}
$K_{j}f(x) \to Kf(x)$ a.e.($j \to \infty$).又
\begin{align*}
\widehat{K_{j}f}(x) &= \frac{1}{2\pi}\int_{-\infty}^{\infty}\sum_{n}s_{j, n}m_{n}f(y + n)e^{-ixy}\dd y\\
&= \sum_{n}s_{j, n}m_{n}e^{-inx}\widehat{f}(x),
\end{align*}
利用 $\widehat{s}_{j}(n) = s_{j, n}$,
\begin{equation*}
\widehat{K_{j}f}(x) = (s_{j} * m)(x)\widehat{f}(x).
\end{equation*}
因此,由 Parseval 等式,
\begin{align*}
&\|K_{j_{k}}f - T_{m}f\|_{L^{2}(\mathbf{R})} = \|\widehat{K_{j_{k}}f} - m\widehat{f}\|_{L^{2}(\mathbf{R})}\\
&\qquad = \|s_{j_{k}} * m(x)\widehat{f}(x) - m(x)\widehat{f}(x)\|_{L^{2}(\mathbf{R})}^{2}\\
&\qquad = \|(s_{j_{k}} * m - m)\widehat{f}\|_{L^{2}(\mathbf{R})}^{2} \to 0 \quad (k \to \infty),
\end{align*}
得 $\widehat{Kf} = m\widehat{f}$ a.e..

其次证明逆命题.设满足 $\widehat{Kf} = m\widehat{f}$ 的算子 $K$ 是从 $L^{p}(\mathbf{R})$ 到 $L^{p}(\mathbf{R})$ 的有界线性算子.此时,把算子范数记为 $\|K\|_{\mathfrak{M}_{p}(\mathbf{R})}$ 或 $\|m\|_{M_{p}(\mathbf{R})}$.又对 $c = \{c_{n}\}$,$d = \{d_{n}\} \in \ell^{p}$(于是 $c_{n}$,$d_{n}$ 除有限个外为 0),令 $f(x) = \sum_{m}c_{m}\chi_{(-\frac{1}{2}, \frac{1}{2})}(x + m)$,$g(x) = \sum_{n}d_{n}\chi_{(-\frac{1}{2}, \frac{1}{2})}(x + n)$,则 $\|f\|_{L^{p}(\mathbf{R})}^{p} = \frac{1}{2\pi}\|c\|_{\ell^{p}}^{p}$,$\|g\|_{L^{p'}(\mathbf{R})}^{p'} = \frac{1}{2\pi}\|d\|_{\ell^{p'}}^{p'}$.进一步,由 Parseval 等式,
\begin{align*}
\left|\int_{-\infty}^{\infty}\widehat{Kf}(x)\overline{\widehat{g}(x)}\dd x\right| &= \left|\frac{1}{2\pi}\int Kf(y)\overline{g(y)}\dd y\right|\\
&\leq \|m\|_{M_{p}(\mathbf{R})}\|f\|_{p}\|g\|_{p'}\\
&= \frac{1}{2\pi}\|m\|_{M_{p}(\mathbf{R})}\|c\|_{\ell^{p}}\|d\|_{\ell^{p'}}.
\end{align*}
另一方面,
\begin{align*}
\int_{-\infty}^{\infty}\widehat{Kf}(x)\overline{\widehat{g}(x)}\dd x&= \int_{-\infty}^{\infty}m(x)\widehat{f}(x)\overline{\widehat{g}(x)}\dd x\\
&= \int_{-\infty}^{\infty}m(x)\left(\sum_{k}c_{k}e^{-ikx}\widehat{\chi_{(-\frac{1}{2}, \frac{1}{2})}}(x)\right)\overline{\left(\sum_{n}d_{n}e^{-inx}\widehat{\chi_{(-\frac{1}{2}, \frac{1}{2})}}(x)\right)}\dd x\\
&= 2\pi\sum_{k, n}c_{k}\overline{d_{n}}\int_{0}^{2\pi}m(x)\sum_{j}\left|\widehat{\chi_{(-\frac{1}{2}, \frac{1}{2})}}(x + 2j\pi)\right|^{2}e^{-i(n-k)x}\dd x\\
&= 2\pi\sum_{k, n}c_{k}\overline{d_{n}}m_{n - k}
\end{align*}
成立(利用 Parseval 等式,$\sum_{j}|\widehat{\chi_{(-\frac{1}{2}, \frac{1}{2})}}(x + 2j\pi)|^{2} = \frac{1}{4\pi^{2}}$).从而
\begin{equation*}
\sup_{\|d\|_{\ell^{p'}} \leq 1}\left|\sum_{k, n}c_{k}\overline{d_{n}}m_{n - k}\right| \leq \|m\|_{M_{p}(\mathbf{R})}\|c\|_{\ell^{p}},
\end{equation*}
故满足 $\widehat{T_{m}f} = m\widehat{f}$ 的算子 $T_{m}$ 是 $\ell^{p}$ 上的有界线性算子,且 $\|T_{m}\|_{\mathfrak{M}_{p}(\mathbf{Z})} = \|m\|_{M_{p}(\mathbf{T})} \leq \|m\|_{M_{p}(\mathbf{R})}$.
\end{proof}

关于 $M_{p}(\mathbf{Z})$ 的元,下面的事实成立.该定理出自 K. de Leeuw \cite{ch3:17},Igari \cite{ch3:25},但这里按 Wozniakowski \cite{ch3:56},Kaneko and Sato \cite{ch3:28} 给出证明.

\begin{theorem}\label{thm:3.5.18}
设 $1 \leq p < \infty$.设 $\mathbf{R}$ 上的有界函数 $\lambda$ 在每个点连续.此时,令 $\widehat{Tf}(\xi) = \lambda(\xi)\widehat{f}(\xi)$,$T$ 是从 $L^{p}(\mathbf{R})$ 到 $L^{p}(\mathbf{R})$ 的有界线性算子的充要条件是:对任意 $\varepsilon > 0$,由 $\widehat{T_{\varepsilon}}f(n) = \lambda(\varepsilon n)\widehat{f}(n)$ 定义的 $\widetilde{T_{\varepsilon}}$ 是 $L^{p}(\mathbf{T})$ 上的有界线性算子,且令 $\lambda_{\varepsilon}(n) = \lambda(\varepsilon n)$,关于其算子范数 $\|T_{\varepsilon}\|_{\mathfrak{M}_{p}(\mathbf{T})} = \|\lambda_{\varepsilon}\|_{M_{p}(\mathbf{Z})}$,$\sup_{\varepsilon > 0}\|\lambda_{\varepsilon}\|_{M_{p}(\mathbf{Z})} < \infty$ 成立.此时,$\|T\|_{M_{p}(\mathbf{R})} = \sup_{\varepsilon > 0}\|\lambda_{\varepsilon}\|_{M_{p}(\mathbf{Z})}$ 成立.其中,$\|T\|_{\mathfrak{M}_{p}(\mathbf{R})}$ 表示 $T$ 在 $L^{p}(\mathbf{R})$ 上的算子范数.
\end{theorem}

\begin{note}\label{note:3.5.19}
把关于 $\lambda$ 的条件``在每个点连续''放宽为``在 $\mathbf{Z}$ 上的点连续'',通过审视证明也可证明 $\lambda|_{\mathbf{Z}} \in M_{p}(\mathbf{Z})$.
\end{note}

\begin{proof}
设 $\lambda$ 是 $\mathbf{R}$ 上的有界连续函数,对任意 $\varepsilon > 0$,令 $\widehat{T_{\varepsilon}f}(\xi) = \lambda(\varepsilon\xi)\widehat{f}(\xi)$($f \in C_c(\mathbf{R})$),设 $T = T_1$ 是 $L^p(\mathbf{R})$ 上的有界线性算子.又由 $\widehat{T_{\varepsilon}g}(n) = \lambda(\varepsilon n)\widehat{g}(n)$($g \in C(\mathbf{T})$)定义算子 $\widetilde{T_{\varepsilon}}$.设 $F$ 为 $\mathbf{T}$ 上的三角多项式.取 $\chi_k \in C^\infty(\mathbf{R})$,$\chi_k \geq 0$,满足 $\chi_k(\tau) = 1$($|\tau| \leq 2k\pi$),$\chi_k(\tau) = 0$($|\tau| > 2(k+1)\pi$),令 $\chi_k^\eta(x) = \chi_k(\eta x)$($\eta > 0$).进一步令 $\gamma_{\eta,k}(x) = \chi_k^\eta(x)(\widetilde{T_{\varepsilon}}F)(x) - T_{\varepsilon}(\chi_k^\eta F)(x)$,并令
\begin{equation*}
\widetilde{\gamma_{\eta,k}}(\xi) = \frac{1}{2\pi}\int_{-\infty}^{\infty}\chi_k^\eta(x)(\widetilde{T_{\varepsilon}}F)(x)e^{i\xi x}\dd x - \frac{1}{2\pi}\int_{-\infty}^{\infty}T_{\varepsilon}(\chi_k^\eta F)(x)e^{i\xi x}\dd x = I - II.
\end{equation*}

于是计算 I,II.
\begin{align*}
\mathrm{I} &= \frac{1}{2\pi}\int_{-\infty}^{\infty}\chi_k^\eta(x)\sum_j\lambda(j\varepsilon)\widehat{F}(j)e^{ijx}e^{i\xi x}\dd x\\
&= \frac{1}{2\pi}\sum_j\lambda(j\varepsilon)\widehat{F}(j)\int_{-\infty}^{\infty}\chi_k^\eta(x)e^{i(j+\xi)x}\dd x\\
&= \sum_j\lambda(j\varepsilon)\widehat{F}(j)\widehat{\chi_k^\eta}(-j-\xi),
\end{align*}
其中,
\begin{equation*}
\widehat{\chi_k^\eta}(-j-\xi) = \frac{1}{2\pi}\int_{-\infty}^{\infty}\chi_k(x)e^{i(j+\xi)x/\eta}\frac{1}{\eta}\dd x,
\end{equation*}
所以
\begin{equation*}
\mathrm{I} = \sum_j\lambda(j\varepsilon)\widehat{F}(j)\widehat{\chi_k}\left(\frac{-j-\xi}{\eta}\right)\frac{1}{\eta}.
\end{equation*}
又,
\begin{align*}
\mathrm{II} &= \frac{1}{2\pi}\int_{-\infty}^{\infty}T_{\varepsilon}(\chi_k^\eta F)(x)e^{i\xi x}\dd x\\
&= \lambda(-\varepsilon\xi)\frac{1}{2\pi}\int_{-\infty}^{\infty}\chi_k^\eta(x)F(x)e^{i\xi x}\dd x\\
&= \lambda(-\varepsilon\xi)\sum_j\widehat{F}(j)\frac{1}{2\pi}\int_{-\infty}^{\infty}\chi_k^\eta(x)e^{i(j+\xi)x}\dd x
\end{align*}
\begin{equation*}
= \lambda(-\varepsilon\xi)\sum_j\widehat{F}(j)\eta^{-1}\widehat{\chi_k}\left(\frac{-j-\xi}{\eta}\right),
\end{equation*}
所以,最终
\begin{equation*}
\mathrm{I} - \mathrm{II} = \sum_j\widehat{F}(j)\eta^{-1}\widehat{\chi_k}\left(\frac{-j-\xi}{\eta}\right)(\lambda(j\varepsilon) - \lambda(-\xi\varepsilon)).
\end{equation*}
因此,
\begin{equation*}
\widetilde{\gamma_{\eta,k}}(\xi) = \sum_j\widehat{F}(j)\eta^{-1}\widehat{\chi_k}\left(\frac{\xi - j}{\eta}\right)(\lambda(j\varepsilon) - \lambda(\xi\varepsilon)),
\end{equation*}
所以
\begin{align*}
\left\|\widetilde{\gamma_{\eta,k}}\right\|_{1} &\leq \sum\left|\widehat{F}(j)\right|\int_{-\infty}^{\infty}\eta^{-1}\left|\widehat{\chi_k}\left(\frac{\xi - j}{\eta}\right)\right||\lambda(j\varepsilon) - \lambda(\xi\varepsilon)|\dd\xi\\
&\leq \sum\left|\widehat{F}(j)\right|\int_{-\infty}^{\infty}|\widehat{\chi_k}(\xi)||\lambda(j\varepsilon) - \lambda(j\varepsilon + \eta\xi\varepsilon)|\dd\xi.
\end{align*}
因此
\begin{equation*}
|\gamma_{\eta,k}(x)| \leq \|\widetilde{\gamma_{\eta,k}}\|_{1} \leq \sum_j\left|\widehat{F}(j)\right|\int_{-\infty}^{\infty}|\widehat{\chi_k}(\xi)||\lambda(j\varepsilon) - \lambda(j\varepsilon + \eta\xi\varepsilon)|\dd\xi.
\end{equation*}
因 $F$ 是三角多项式,固定 $k$ 时,若 $\eta \to 0$,则 $\|\widetilde{\gamma_{\eta,k}}\|_{1} \to 0$.于是,对 $a > 0$,取 $\eta > 0$ 使 $\|\widetilde{\gamma_{\eta,k}}\|_{1} < a$,设 $t > a$,则
\begin{equation*}
\left|\{x: |T_{\varepsilon}(\chi_k^\eta F)(x)| > t - \|\widetilde{\gamma_{\eta,k}}\|_{1}\}\right| \geq \left|\{x: |\chi_k^\eta(\widetilde{T_{\varepsilon}}F)(x)| > t\}\right|
\end{equation*}
(对集合 $E$,$|E|$ 表示 $E$ 的 Lebesgue 测度).从而
\begin{align*}
\mathrm{III} &= \int_{a}^{\infty}p\left|\left\{\left|\chi_k^\eta(\widetilde{T_{\varepsilon}}F)(x)\right| > t\right\}\right|t^{p-1}\dd t\\
&\leq \int_{a}^{\infty}p\left|\{|T_{\varepsilon}(\chi_k^\eta F)(x)| > t - \|\widetilde{\gamma_{\eta,k}}\|_{1}\}\right|t^{p-1}\dd t\\
&\leq \int_{a - \|\widetilde{\gamma_{\eta,k}}\|_{1}}^{\infty}p\left|\{|T_{\varepsilon}(\chi_k^\eta F)(x)| > t\}\right|(t + \|\widetilde{\gamma_{\eta,k}}\|_{1})^{p-1}\dd t
\end{align*}
\begin{align*}
&= \int_{a - \|\widetilde{\gamma_{\eta,k}}\|_{1}}^{\infty}p\left(t|\{|T_{\varepsilon}(\chi_k^\eta F)(x)| > t\}|^{1/p}\right)^{p-1}\left(1 + \frac{\|\widetilde{\gamma_{\eta,k}}\|_{1}}{t}\right)^{p-1}\\
&\leq \left(1 + \frac{\|\widetilde{\gamma_{\eta,k}}\|_{1}}{a - \|\widetilde{\gamma_{\eta,k}}\|_{1}}\right)^{p-1}\|T_{\varepsilon}(\chi_k^\eta F)\|_{p}^{p}.
\end{align*}
这里,因 $T_{\varepsilon}(\chi_{k}^{\eta}F)(x) = (T(\chi_{k}^{\eta}F)^{\varepsilon})(x/\varepsilon)$(其中 $(\chi_{k}^{\eta}F)^{\varepsilon}(x) = (\chi_{k}^{\eta}F)(\varepsilon x)$),
\begin{align*}
&\mathrm{III} \leq \left(1 + \frac{\|\widetilde{\gamma_{\eta,k}}\|_{1}}{a - \|\widetilde{\gamma_{\eta,k}}\|_{1}}\right)^{p}\left\|(T(\chi_{k}^{\eta}F)^{\varepsilon})\left(\frac{\cdot}{\varepsilon}\right)\right\|_{p}^{p}\\
&\quad \leq \left(1 + \frac{\|\widetilde{\gamma_{\eta,k}}\|_{1}}{a - \|\widetilde{\gamma_{\eta,k}}\|_{1}}\right)^{p}\varepsilon\|T\|_{\mathfrak{M}_{p}(\mathbf{R})}^{p}\|(\chi_{k}^{\eta}F)^{\varepsilon}\|_{p}^{p}\\
&\quad = \left(1 + \frac{\|\widetilde{\gamma_{\eta,k}}\|_{1}}{a - \|\widetilde{\gamma_{\eta,k}}\|_{1}}\right)^{p}\|T\|_{\mathfrak{M}_{p}(\mathbf{R})}^{p}\|\chi_{k}^{\eta}F\|_{p}^{p}.
\end{align*}
这里,由 $0 \leq \chi_{k}^{\eta} \leq 1$,$\supp\chi_{k}^{\eta} \subset \left[\frac{-2(k+1)\pi}{\eta}, \frac{2(k+1)\pi}{\eta}\right]$,对 $Q_{l} = [-\pi, \pi] + 2l\pi$($l \in Z$),满足 $Q_{l} \cap \left[\frac{-2(k+1)\pi}{\eta}, \frac{2(k+1)\pi}{\eta}\right] \neq \emptyset$ 的 $l$ 的个数至多为 $\left[2\left(\frac{k+1}{\eta}\right)\right]$.因此
\begin{equation*}
\|\chi_{k}^{\eta}F\|_{L^{p}(\mathbf{R})} \leq \left\{2\left(\frac{k+1}{\eta} + 1\right)\right\}^{1/p}\|F\|_{L^{p}(\mathbf{T})}.
\end{equation*}
另一方面,由 $\chi_{k}^{\eta}(x) = 1$($x \in \left[\frac{-2k\pi}{\eta}, \frac{2k\pi}{\eta}\right]$),满足 $Q_{j} \subset \left[\frac{-2k\pi}{\eta}, \frac{2k\pi}{\eta}\right]$ 的 $Q_{j}$ 的个数至少为 $\left[2\left(\frac{k}{\eta} - 1\right)\right]$,所以
\begin{align*}
\left|\left\{\left|\chi_{k}^{\eta}(\widetilde{T_{\varepsilon}}F)\right| > t\right\}\right| &\geq \left|\bigcup_{j}\left\{Q_{j} \cap \left\{\left|\widetilde{T_{\varepsilon}}F\right| > t\right\}\right\}: Q_{j} \subset \left[\frac{-2k\pi}{\eta}, \frac{2k\pi}{\eta}\right]\right|\\
&\geq \left\{2\left(\frac{k}{\eta} - 1\right)\right\}\left|\left\{\left|\widetilde{T_{\varepsilon}}F\right| > t\right\}\right|.
\end{align*}
综合以上,
\begin{align*}
&\left\{2\left(\frac{k}{\eta} - 1\right)\right\}^{1/p}\left\{\int_{a}^{\infty}p\left|\left\{\left|\widetilde{T_{\varepsilon}}F\right| > t\right\}\right|^{1/p}t^{p-1}\dd t\right\}^{1/p}\\
&\qquad \leq \left\{2\left(\frac{k+1}{\eta} + 1\right)\right\}^{1/p}\left\{1 + \frac{\|\widetilde{\gamma_{\eta,k}}\|_{1}}{a - \|\widetilde{\gamma_{\eta,k}}\|_{1}}\right\}\|T\|_{\mathfrak{M}_{p}(\mathbf{R})}\|F\|_{p}.
\end{align*}
因此
\begin{align*}
\{2(k - \eta)\}^{1/p}\left\{\int_{a}^{\infty}p\left|\left\{\left|\widetilde{T_{\varepsilon}}F\right| > t\right\}\right|^{1/p}t^{p-1}\dd t\right\}^{1/p}&\leq \{2(k + 1 + \eta)\}^{1/p}\left\{1 + \frac{\|\widetilde{\gamma_{\eta,k}}\|_{1}}{a - \|\widetilde{\gamma_{\eta,k}}\|_{1}}\right\}\|T\|_{\mathfrak{M}_{p}(\mathbf{R})}\|F\|_{p}.
\end{align*}
这里令 $\eta \to 0$,得
\begin{align*}
(2k)^{1/p}\left\{\int_{a}^{\infty}p\left|\left\{\left|\widetilde{T_{\varepsilon}}F\right| > t\right\}\right|^{1/p}t^{p-1}\dd t\right\}^{1/p}&\leq \{2(k + 1)\}^{1/p}\|T\|_{\mathfrak{M}_{p}(\mathbf{R})}\|F\|_{p}.
\end{align*}
除以 $k$ 后,令 $k \to \infty$,再在 $a \to 0$ 时利用 Fatou 引理,得
\begin{equation*}
\|\widetilde{T_{\varepsilon}}F\|_{p} \leq \|T\|_{\mathfrak{M}_{p}(\mathbf{R})}\|F\|_{p}.
\end{equation*}

反之,设 $\sup_{\varepsilon>0}\|\widetilde{T_{\varepsilon}}\|_{M_{p}(\mathbf{T})}<\infty$.取 $f\in C_{c}^{\infty}(\mathbf{R})$,令 $F_{\varepsilon}(x)=\varepsilon^{-1}\sum_{k\in\mathbf{Z}}f(\varepsilon^{-1}(x+2k\pi))$,则 $F_{\varepsilon}\in L^{1}(\mathbf{T})$.于是,
\begin{align*}
\widehat{F_{\varepsilon}}(j) &= \frac{1}{2\pi}\int_{-\pi}^{\pi}F_{\varepsilon}(x)e^{-ijx}\dd x\\
&= \frac{1}{2\pi}\sum_{k}\varepsilon^{-1}\int_{-\pi}^{\pi}f(\varepsilon^{-1}(x+2k\pi))e^{-ijx}\dd x\\
&= \frac{1}{2\pi}\int_{-\infty}^{\infty}\varepsilon^{-1}f(\varepsilon^{-1}x)e^{-ijx}\dd x = \widehat{f}(\varepsilon j),
\end{align*}
得 $\widehat{F_{\varepsilon}}(j)=\widehat{f}(\varepsilon j)$.因此 $F_{\varepsilon}(x)=\sum_{j}\widehat{f}(\varepsilon j)e^{ijx}$.这里,考虑 Riemann 积分的定义,
\begin{align*}
|Tf(x)| &= \left|\int_{-\infty}^{\infty}\lambda(\xi)\widehat{f}(\xi)e^{ix\xi}\dd\xi\right|\\
&= \lim_{\varepsilon \to 0}\left|\sum_{j}\lambda(\varepsilon j)\widehat{f}(\varepsilon j)e^{ix\varepsilon j}\varepsilon\right|
\end{align*}
\begin{align*}
&= \lim_{\varepsilon \to 0}\left|\sum_{j}\lambda(\varepsilon j)\widehat{F_{\varepsilon}}(j)e^{i\varepsilon xj}\varepsilon\right|\\
&= \lim_{\varepsilon \to 0}\left|\widetilde{T_{\varepsilon}}(\varepsilon F_{\varepsilon})(\varepsilon x)\right|.
\end{align*}
进一步,由 $\lim_{\varepsilon\to0}\chi_{[-\pi,\pi)}(\varepsilon x)=1$($x\in\mathbf{R}$),
\begin{equation*}
|Tf(x)| = \lim_{\varepsilon \to 0}\left|\widetilde{T_{\varepsilon}}(\varepsilon F_{\varepsilon})(\varepsilon x)\right|\chi_{[-\pi,\pi)}(\varepsilon x),
\end{equation*}
所以
\begin{align*}
\|Tf\|_{L^{p}(\mathbf{R})} &= \left(\int_{-\infty}^{\infty}|Tf|^{p}\frac{\dd x}{2\pi}\right)^{1/p}\\
&\leq \liminf_{\varepsilon \to 0}\left(\frac{1}{2\pi}\int_{-\pi}^{\pi}\left|\widetilde{T_{\varepsilon}}(\varepsilon F_{\varepsilon})(x)\right|^{p}\varepsilon^{-1}\dd x\right)^{1/p}\\
&\leq \liminf_{\varepsilon \to 0}\varepsilon^{-1/p}\|\widetilde{T_{\varepsilon}}(\varepsilon F_{\varepsilon})\|_{L^{p}(\mathbf{T})}\\
&\leq \left(\sup_{\varepsilon}\|\widetilde{T_{\varepsilon}}\|_{\mathfrak{M}_{p}(\mathbf{T})}\right)\liminf_{\varepsilon \to 0}\varepsilon^{-1/p}\|\varepsilon F_{\varepsilon}\|_{L^{p}(\mathbf{T})}\\
&\leq \left(\sup_{\varepsilon}\|\widetilde{T_{\varepsilon}}\|_{\mathfrak{M}_{p}(\mathbf{T})}\right)\|f\|_{L^{p}(\mathbf{R})}.
\end{align*}
因此
\begin{equation*}
\|Tf\|_{L^{p}(\mathbf{R})} \leq \left(\sup_{\varepsilon > 0}\|\widetilde{T_{\varepsilon}}\|_{\mathfrak{M}_{p}(\mathbf{T})}\right)\|f\|_{L^{p}(\mathbf{R})}.
\end{equation*}
\end{proof}

下面研究乘子空间 $M_{p}(\mathbf{Z})$($1 < p < 2$)的作用函数.为此,先叙述若干引理(\cite{ch3:24},cf. \cite{ch3:43}, \cite{ch3:48}).

\begin{lemma}\label{lem:3.5.20}
设 $1 < p < 2$.存在 $K_{p}(>1)$,使得对每个自然数 $j$,存在满足下列条件的 $\mathbf{T}$ 上的三角多项式 $\phi = \phi_{p,j}$:

(i) $\widehat{\phi}(n) \geq 0$($n \in \mathbf{Z}$),$\|\phi\|_{L^1(\mathbf{T})} = 1$.

(ii) $\|\exp(ij\phi)\|_{M_p(\mathbf{T})} > K_p^j$.

这里,$\exp(ij\phi) \in M(\mathbf{T})$ 由 $\exp(ij\phi)^{\wedge}(n) = \exp(ij\widehat{\phi}(n))$($n \in \mathbf{Z}$)定义.
\end{lemma}

\begin{proof}
分几个阶段叙述.

第 1 阶段.令 $g(x)=e^{i\cos x}$,因 $g\in A(\mathbf{T})$,
\begin{equation*}
K_{p}^{\prime} = \left(\sum|\widehat{g}(n)|^{p}\right)^{1/p} < \sum|\widehat{g}(n)| < \infty.
\end{equation*}
又由 $1 < p < 2$,$K_{p}^{\prime} > \left(\sum|g(n)|^{2}\right)^{1/2} = \|g\|_{2} = 1$,故 $K_{p}^{\prime} > 1$.现在取 $0 < \mu < 1$ 使 $K_{p} = \mu K_{p}^{\prime} > 1$.

第 2 阶段.令 $Q(x) = \sum_{j=1}^{r}\cos n_{j}x$,适当选取 $n_{1}, \ldots, n_{r}$,令 $g = e^{iQ}$,则 $A_{p}(g) = \left(\sum|\widehat{g}(n)|^{p}\right)^{1/p} > K_{p}^{r}$($r = 1, 2, \ldots$).事实上,令 $n_{1} = 1$,由第 1 阶段即成立.设 $n = k < r$ 时成立.此时,对 $f = e^{i\sum_{j=1}^{k}\cos n_{j}x}$,$g = e^{i\cos n_{k+1}x}$,由 $A_{p}(fg) \leq \|f\|_{A(\mathbf{T})}A_{p}(g)$,$f$ 不妨取为三角多项式.此时,把 $n_{k+1}$ 取得足够大,$A_{p}(fg) = A_{p}(f)A_{p}(g)$ 成立,故可做到 $A_{p}(fg) > K_{p}^{k}\mu A_{p}(g) > K_{p}^{k+1}$.从而 $A_{p}\left(e^{i\sum_{j=1}^{k+1}\cos n_{j}x}\right) > K_{p}^{k+1}$.

第 3 阶段.对上述 $K_{p} > 1$,存在 $\mathbf{T}$ 上的概率测度 $\mu$ 满足 $\widehat{\mu}(n) \in \mathbf{R}$($n \in \mathbf{Z}$)且 $\|\exp(ij\mu/2)\|_{\mathfrak{M}_{p}(\mathbf{T})} > K_{p}^{j}$.事实上,由第 2 阶段,存在 $\mathbf{T}$ 上的三角多项式 $Q$ 满足 $A_{p}(e^{iQ}) > K_{p}^{j}$($\|Q\|_{A(\mathbf{T})} \leq j$).把 $T_{\varphi}$ 看作由 $\widehat{T_{\varphi}}f(x) = \varphi(x)\widehat{f}(x)$($\widehat{f}(x) = \sum f_n e^{inx}$)给出的 $\ell^p$ 上的算子.其算子范数为 $\|T_{\varphi}\|_{\mathfrak{M}_{p}(\mathbf{Z})} = \|\varphi\|_{M_p(\mathbf{T})} \geq A_p(\varphi)$,且 $\|Q\|_{A(\mathbf{T})} = \|Q\|_{M_1(\mathbf{T})}$(cf. \cite{ch3:9},引理~\ref{lem:3.2.11}),故令 $\varphi = \frac{1}{j}Q$,得 $\|e^{ij\varphi}\|_{M_p(\mathbf{T})} > K_p^j$($j = 1, 2, \ldots$).此时,对该 $\varphi$ 应用定理~\ref{thm:3.5.17},再对其应用定理~\ref{thm:3.5.18},适当取 $\varepsilon > 0$,则存在概率测度 $\mu$ 满足 $\widehat{\mu}(n) = \varphi(\varepsilon n)$,$\widehat{\mu}(n) \in \mathbf{R}$($n \in \mathbf{Z}$)且 $\|e^{ij\mu}\|_{\mathfrak{M}_{p}(\mathbf{T})} > K_p^j$($j = 1, 2, \ldots$).进一步,
\begin{equation*}
\|e^{ij\mu/2}\|_{\mathfrak{M}_{p}(\mathbf{T})}^{2} \geq \|e^{ij\mu}\|_{\mathfrak{M}_{p}(\mathbf{T})} > K_{p}^{j},
\end{equation*}
所以 $\|e^{ij\mu/2}\|_{\mathfrak{M}_{p}(\mathbf{T})} > \left(\sqrt{K_{p}}\right)^{j}$.这里令 $\nu = (\delta_{0} + \mu)/2$,则 $\widehat{\nu}(n) \geq 0$($n \in \mathbf{Z}$),$\nu$ 是概率测度,且 $\|\exp(ij\nu)\|_{\mathfrak{M}_{p}(\mathbf{T})} > \sqrt{K_{p}}^{j}$($j = 1, 2, \ldots$).从而,取次数足够高的 Fejér 核 $\sigma_{n}$,令 $\phi = \nu * \sigma_{n}$,把 $\sqrt{K_{p}}$ 重新记为 $K_{p}$,即得结果.
\end{proof}

\begin{theorem}\label{thm:3.5.21}
设 $1 < p < 2$,则存在 $\mathbf{T}$ 上的概率测度 $\mu$($\hat{\mu} \geq 0$)和常数 $K_{p} > 1$,使得令
\begin{equation*}
Tf = e^{i\mu} * f \quad (f \in L^{p}(\mathbf{T})),
\end{equation*}
算子 $T$ 在 $\mathfrak{M}_{p}(\mathbf{T})$ 中的谱半径 $\rho(T)$ 满足 $\rho(T) \geq K_{p}$.
\end{theorem}

\begin{proof}
设 $Q$ 为有理数的集合.设 $\{p_{n}\}$ 为使任意 $p \in \mathbf{Q} \cap (1, 2)$ 出现无穷多次的数列.由引理~\ref{lem:3.5.20},对自然数 $j$ 存在 $\mathbf{T}$ 上的三角多项式 $\phi_{j}$ 使得 $\widehat{\phi_{j}} \geq 0$,$\|\phi_{j}\|_{L^{1}(\mathbf{T})} = 1$,$\|\exp(ij\phi_{j})\|_{M_{p}(\mathbf{T})} > K_{p}^{j}$($p = p_{j}$).于是,对 $p = p_{j}$,取三角多项式 $f_{j}$ 使得
\begin{equation*}
\|f_{j}\|_{p} = 1, \quad \left\|\sum_{k}\exp(ij\widehat{\phi_{j}}(k))\widehat{f_{j}}(k)\exp(ikt)\right\|_{p} > K_{p}^{j}.
\end{equation*}
此时,存在自然数 $m_{j}$ 使得
\begin{equation*}
\left(\supp\widehat{\phi_{j}}\right) \cup \left(\supp\widehat{f_{j}}\right) \subset \{-m_{j}, -m_{j} + 1, \ldots, m_{j} - 1, m_{j}\}.
\end{equation*}
现在,设 $\{r_{j}\} \subset \mathbb{N}$ 为急剧增大的序列,则满足
\begin{equation*}
d\mu_{n}(t) = \phi_{1}(r_{1}t)\phi_{2}(r_{2}t)\cdots\phi_{n}(r_{n}t)\frac{\dd t}{2\pi}
\end{equation*}
的 $\mu_{n}\in M(\mathbf{T})$ 关于 $*$-弱拓扑收敛于概率测度 $\mu\in M(\mathbf{T})$,且满足
\begin{equation*}
\widehat{\mu}(k_{1}r_{1} + \cdots + k_{n}r_{n}) = \prod_{j=1}^{n}\widehat{\phi_{j}}(k_{j}) \quad (|k_{j}| \leq m_{j},\ j = 1, 2, \ldots, n).
\end{equation*}
进一步,令 $g_{j}(t)=f_{j}(r_{j}t)$,
\begin{align*}
\left\|\sum_{k}\exp(ij\widehat{\mu}(k))\widehat{g_{j}}(k)\exp(ikt)\right\|_{p} &= \left\|\sum_{k}\exp(ij\widehat{\phi_{j}}(k))\widehat{f_{j}}(k)\exp(ikt)\right\|_{p}\\
&> K_{p}^{j}.
\end{align*}
因此 $\|\exp(ij\mu)\|_{M_{p}(\mathbf{T})}>K_{p}^{j}$($p=p_{j}$),由 Riesz--Thorin 插值定理,对所有 $p\in(1,2)$ 有 $\|\exp(ij\mu)\|_{M_{p}(\mathbf{T})}>K_{p}^{j}$.从而由谱半径的定义即得结果.
\end{proof}

\begin{note}\label{note:3.5.22}
对上述 $\mu$,由 Riesz--Thorin 插值定理,下式成立:
\begin{equation*}
\sup\left\{\|e^{ij\mu}\|_{\mathfrak{M}_{p}(\mathbf{T})}: \mu \in M(\mathbf{T}),\ \hat{\mu}(\mathbf{Z}) \subset \mathbf{R},\ \|\mu\|_{M(\mathbf{T})} \leq 1\right\} \leq e^{|j|(2/p - 1)} \quad (j \in \mathbf{Z}).
\end{equation*}
\end{note}

\begin{lemma}\label{lem:3.5.23}
在 $[-1, 1]$ 上定义的函数 $\Phi$ 若作用于 $M_{p}(\mathbf{Z})$,则连续.
\end{lemma}

\begin{proof}
设 $\Phi$ 在 $a \in [-1, 1]$ 处不连续.此时,存在 $\{a_j\} \subset [-1, 1]$ 和 $B \in \mathbf{R}$ 使得
\begin{equation*}
\Phi(a) \neq B, \quad a_i \neq a_j \ (i \neq j), \quad \sum_{j=-\infty}^{\infty}|a_j - a| < \infty, \quad \sum_{j=-\infty}^{\infty}|\Phi(a_j) - B| < \infty.
\end{equation*}
此时不妨设 $\Phi(a)=0$.又,适当选取 $\{\varepsilon_{j}\}_{j=-\infty}^{\infty}$($\varepsilon_{j}=\pm1$),取 $f \in L^{p}(\mathbf{T})$ 使得
\begin{equation*}
\sum_{j=-\infty}^{\infty}\varepsilon_{j}\widehat{f}(j)e^{ijx} \notin L^{p}(\mathbf{T})
\end{equation*}
($f$ 的存在性参见 \cite{ch3:24}).现在,令 $\eta_{j} = \Phi(a_{j})$($\varepsilon_{j} = 1$),$\eta_{j} = 0$($\varepsilon_{j} = -1$),则 $\sum_{-\infty}^{\infty}\widehat{f}(j)\eta_{j}e^{ijx} \notin L^{p}(\mathbf{T})$.事实上,由
\begin{align*}
\sum_{j=-\infty}^{\infty}\widehat{f}(j)\eta_{j}e^{ijx}&= \frac{B}{2}\sum_{j=-\infty}^{\infty}\widehat{f}(j)e^{ijx} + \frac{B}{2}\sum_{j=-\infty}^{\infty}\widehat{f}(j)\varepsilon_{j}e^{ijx} + \sum_{\varepsilon_{j}=1}\widehat{f}(j)(\Phi(a_{j}) - B)e^{ijx},
\end{align*}
右边第 1 项,第 3 项是 $L^{p}(\mathbf{T})$ 的元,而第 2 项不是.于是,令 $\varphi(j)=a_{j}$($\varepsilon_{j}=1$),$\varphi(j)=a$($\varepsilon_{j}\neq1$),则对 $g\in L^{1}(\mathbf{T})$,
\begin{equation*}
\sum_{j}\widehat{g}(j)\varphi(j)e^{ijx} = \sum_{j}\widehat{g}(j)(\varphi(j) - a)e^{ijx} + a\sum_{j}\widehat{g}(j)e^{ijx}.
\end{equation*}
因此右边是 $L^{1}(\mathbf{T})$ 的元,故 $\varphi \in M_{1}(\mathbf{Z})$.从而由 $\Phi$ 是 $M_{p}(\mathbf{Z})$ 的作用函数,$\Phi(\varphi) \in M_{p}(\mathbf{Z})$.可是,$\Phi(\varphi(j)) = \Phi(a_j) = \eta_j$,$\{\eta_j\} \notin M_p(\mathbf{Z})$,故 $\Phi$ 不作用于 $M_p(\mathbf{Z})$,矛盾.
\end{proof}

$M_{p}(\mathbf{Z})$ 的作用函数如下确定.

\begin{theorem}\label{thm:3.5.24}
$[-1, 1]$ 上的函数 $\Phi$ 作用于 $M_{p}(\mathbf{Z})$ 的充要条件是:与某个整函数一致.
\end{theorem}

\begin{proof}
设 $\Phi$ 是作用于 $[-1, 1]$ 上的 $M_p(\mathbf{Z})$ 的函数.考虑函数 $\Phi(\varepsilon\sin t)$($0 < \varepsilon < 1$),它是 $\mathbf{R}$ 上的 $2\pi$ 周期函数且作用于 $M_p(\mathbf{Z})$.因此,不妨从一开始假设 $\Phi$ 是 $\mathbf{R}$ 上的 $2\pi$ 周期函数,只需证明它与某个整函数一致(此后的讨论与 $M(\mathbf{T})$ 的作用函数的确定情形相同).此时,由引理~\ref{lem:3.5.23} 不妨设 $\Phi$ 连续.又,像引理~\ref{lem:3.5.13} 的证明那样考虑,对任意 $\lambda \in \mathbf{R}$ 存在 $C_\lambda > 0$ 使得 $\|\Phi(\lambda\hat{\mu} + a)\|_{M_p(\mathbf{Z})} \leq C_\lambda$($-\pi \leq a \leq \pi$).于是,把 $\Phi$ 的 Fourier 级数记为 $\Phi \sim \sum a_n e^{int}$.设 $j$ 为任意自然数,$n \in \mathbf{Z}$,则
\begin{align*}
a_{n}\exp(inj\widehat{\mu}) &= \frac{1}{2\pi}\int_{-\pi}^{\pi}\Phi(x)e^{-in(x-j\widehat{\mu})}\dd x\\
&= \frac{1}{2\pi}\int_{-\pi}^{\pi}\Phi(x + j\widehat{\mu})e^{-ijnx}\dd x.
\end{align*}
由前半部分,
\begin{equation*}
\|a_{n}\exp(inj\widehat{\mu})\|_{M_{p}(\mathbf{Z})} \leq \frac{1}{2\pi}\int\|\Phi(x + j\widehat{\mu})\|_{M_{p}(\mathbf{Z})}\dd x \leq C_{j}.
\end{equation*}
因此
\begin{equation*}
|a_{n}| \leq C_{j}\|\exp(inj\widehat{\mu})\|_{M_{p}(\mathbf{Z})}^{-1}.
\end{equation*}
这里,由引理~\ref{lem:3.5.20} 的证明,
\begin{equation*}
\|\exp(inj\widehat{\mu})\|_{M_{p}(\mathbf{Z})}^{-1} \leq C\exp(-|n|j\log K_{p}),
\end{equation*}
所以
\begin{equation*}
|a_{n}| \leq C_{j}'\exp(-|n|j\log K_{p}) \quad (n \in \mathbf{Z})
\end{equation*}
被证明.从而 $\Phi$ 在 $[-1, 1]$ 上与整函数一致.
\end{proof}

\subsection{\texorpdfstring{$M(p,q)$($1<p<q<\infty$)的作用函数}{M(p,q)(1<p<q<∞)的作用函数}}\label{sec:3.5.5}

设 $1 \leq p < q < \infty$.把从 $L^{p}(\mathbf{T})$ 到 $L^{q}(\mathbf{T})$ 的平移不变有界线性算子的全体记为 $\mathfrak{M}_{pq}$($= \{T \in \mathfrak{M}_{p}(\mathbf{T}): L^{p}(\mathbf{T})$ 到 $L^{q}(\mathbf{T})$ 的有界线性算子$\}$),它关于算子范数成为交换 Banach 代数,且 $\mathfrak{M}_{pq} \subset \mathfrak{M}_{p}(\mathbf{T})$.又定义 $M(p,q) = \{\hat{T} \in \ell^{\infty}: T \in \mathfrak{M}_{pq}(\mathbf{T})\}$($\subset M_{p}(\mathbf{Z})$).当 $p = 1$ 时,$\mathfrak{M}_{pq} = L^{q}(\mathbf{T})$(cf. \cite{ch3:9}).Rudin \cite{ch3:41},Rider \cite{ch3:40} 研究了 $M(1, q)$ 的作用函数.这里叙述推广了他们结果的 Igari--Sato \cite{ch3:26} 的结果.另外,$M(p,q)$ 的范数由 $\mathfrak{M}_{pq}$ 引入.对函数 $\phi$ 与算子 $T$,算子 $\phi \circ T$ 由 $\widehat{\phi \circ T}(n) = \phi(\hat{T}(n))$($n \in \mathbf{Z}$)定义,下文使用.

关于 $M(p,q)$ 的作用函数,下面的事实成立.

\begin{theorem}\label{thm:3.5.25}
(1) 设 $1 < p < q \leq 2$.设 $\beta_0 = (1/q - 1/2)/(1/p - 1/q)$,$n_0$ 为满足 $n_0 \geq \beta_0$ 的最小整数.此时,设 $\alpha_1, \ldots, \alpha_{n_0}$ 为任意复数,取 $[-1, 1]$ 上的有界复值函数 $\Phi_0(t)$,则
\begin{equation*}
\Phi(t) = \alpha_{1}t + \cdots + \alpha_{n_0}t^{n_0} + |t|^{\beta_0 + 1}\Phi_0(t)
\end{equation*}
作用于 $M(p,q)$.

(2) 设 $1 < p < 2 < q < \infty$.$[-1, 1]$ 上的复值函数 $\Phi$ 对复数 $\alpha$ 与
\begin{equation*}
\beta_{1} = \min\left\{\frac{\frac{1}{2} - \frac{1}{q}}{\frac{1}{p} - \frac{1}{2}}, \frac{\frac{1}{p} - \frac{1}{2}}{\frac{1}{2} - \frac{1}{q}}\right\}
\end{equation*}
以及有界函数 $\Phi_{0}$,若写成 $\Phi(t)=\alpha t+|t|^{\beta_{1}+1}\Phi_{0}(t)$,则 $\Phi$ 作用于 $M(p,q)$.
\end{theorem}

由此,关于 $M(p,q)$ 的元的谱,经若干计算,下面的事实成立(cf. \cite{ch3:26}).

\begin{corollary}\label{cor:3.5.26}
\begin{sloppypar}
设 $1 < p < q < \infty$.$M(p,q)$ 是不含单位元的交换 Banach 代数,对 $M(p,q) \ni f$,$\mathrm{Sp}(f, M(p,q)) = \overline{f(\mathbf{Z})}$.
\end{sloppypar}
\end{corollary}

\begin{proof}
由 Riesz--Thorin 插值定理和 Riesz 表示定理,不妨设 $q < 2$.设 $0 \neq \lambda \notin \widehat{T}(\mathbf{Z})$,$\lambda \in \mathrm{Sp}(T, \mathfrak{M}_{pq})$.此时,取足够大的自然数 $N$,由 Riesz--Thorin 插值定理,$T^{N} \in \mathfrak{M}_{p2}(\mathbf{T})$(cf. 引理~\ref{lem:3.5.29}).于是,令
\begin{equation*}
\Phi(z) = \left(\frac{1}{z - \lambda} + \frac{1}{\lambda}\right)^{N} = z^{N}\Phi_{0}(z),
\end{equation*}
则 $\Phi_{0}(z)$ 在 $\widehat{T}(\mathbf{Z})$ 的邻域上有界,故作用于 $M_{22}$.因此,令 $\widehat{S}(n)=\overline{\Phi(\widehat{T}(n))}$($n\in\mathbf{Z}$),则 $S\in M_{p2}\subset M_{pq}$.又由 $\widehat{S}(n)(\widehat{T}(n)-\lambda)^{N}\lambda^{N}=\widehat{T}^{N}(n)$($n\in\mathbf{Z}$),$S(T-\lambda)^{N}\lambda^{N}=T^{N}$.从而,取满足 $h(T)=\lambda$ 的 $h\in\Delta(\mathfrak{M}_{pq})$,与 $\lambda\neq0$ 矛盾.设 $0\notin\overline{\widehat{T}(\mathbf{Z})}$.此时,存在 $\varepsilon_{0}>0$ 使得 $\left|\widehat{T}(n)\right|\geq\varepsilon_{0}>0$($n\in\mathbf{Z}$).因此 $\left|\widehat{T^{N}}(n)\right|\geq\varepsilon_{0}^{N}$($n\in\mathbf{Z}$),对任意三角多项式 $f$,
\begin{equation*}
\varepsilon_{0}^{N}\|f\|_{2} \leq \|T^{N}f\|_{2} \leq C\|f\|_{p}.
\end{equation*}
这与 $L^{2}(\mathbf{T})\neq L^{p}(\mathbf{T})$ 矛盾.
\end{proof}

\begin{note}\label{note:3.5.27}
(1) 设 $2 \leq p < q < \infty$.关于 $M(p, q)$ 的作用函数,可证明与定理~\ref{thm:3.5.25} (1) 类似的事实.

(2) 当 $1 < p < q < 2$ 时,定理~\ref{thm:3.5.25} (1) 的条件是否也是必要条件,是未解决的问题.

(3) 设 $1 < q < p < \infty$.由 Doss \cite{ch3:16} 的结果,与 3.5.4 小节的结果合并可知,$M(p, q)$ 的在 $[-1, 1]$ 上定义的作用函数与整函数一致.
\end{note}

其次转入定理~\ref{thm:3.5.25} 的证明.利用若干引理,分若干阶段证明.下面是省略证明,作为 Stein 插值定理广为人知的事实(cf. \cite{ch3:4} 及其参考文献).

\begin{lemma}[Stein]\label{lem:3.5.28}
设 $1 \leq p_{0}, q_{0}, p_{1}, q_{1} \leq \infty$,$\{T_{z}\}$ 为满足下列性质的算子族.

(1) $z \mapsto \int_{-\pi}^{\pi}T_z(f)g\dd x$($f, g \in L^1(\mathbf{T})$)在 $0 < \operatorname{Re} z < 1$ 上解析,在 $0 \leq \operatorname{Re} z \leq 1$ 上连续.

(2) 存在某个 $0 < a < \pi$,使得
\begin{equation*}
e^{-a|\operatorname{Im}z|}\log\left|\int_{-\pi}^{\pi}T_{z}(f)g\dd x\right|
\end{equation*}
在 $0 \leq \operatorname{Re}z \leq 1$ 上一致有界.

(3) 存在常数 $M_0, M_1$,对任意实数 $y$,$\|T_{iy}f\|_{q_0} \leq M_0\|f\|_{p_0}$($f \in L^{p_0}(\mathbf{T})$),$\|T_{1+iy}f\|_{q_1} \leq M_1\|f\|_{p_1}$($f \in L^{p_1}(\mathbf{T})$)($y$ 是实数)成立.

此时,设 $0 < \theta < 1$,$\operatorname{Re}z = \theta$,$1/p = (1 - \theta)/p_0 + \theta/p_1$,$1/q = (1 - \theta)/q_0 + \theta/q_1$,则 $\|T_\theta f\|_q \leq M_0^{1-\theta}M_1^\theta\|f\|_p$($f \in L^p(\mathbf{T})$)成立.
\end{lemma}

\begin{lemma}\label{lem:3.5.29}
设 $1 < p < q < 2$.设 $T \in \mathfrak{M}_{pq}$ 且 $\widehat{T}(n)$ 取实数值.进一步设 $1/q_{n} = 1/p - n(1/p - 1/q)$,$1/s_{n} = 1/2 + n(1/p - 1/q)$.此时,取定理~\ref{thm:3.5.25}(1) 的 $n_{0}$,对 $n = 1, 2, \ldots, n_{0}$,$T^{n} \in \mathfrak{M}_{s_{n}2}$,且 $\|T^{n}\|_{\mathfrak{M}_{s_{n}2}} \leq \|T\|_{\mathfrak{M}_{pq}}^{n}$ 成立.
\end{lemma}

\begin{proof}
设 $T \in \mathfrak{M}_{pq}$,$\widehat{T}(n)$ 取实数值.此时 $\|T\|_{\mathfrak{M}_{pq}} = \|T\|_{\mathfrak{M}_{q'p'}}$.利用这一点与 Riesz--Thorin 插值定理,令 $1/u = (1-\theta)/p + \theta/q'$,$1/v = (1-\theta)/q + \theta/p'$($0 \leq \theta \leq 1$),则 $T \in M_{uv}$ 且 $\|T\|_{M_{uv}} \leq \|T\|_{M_{pq}}$ 成立.令 $\theta = n(1/p - 1/q)/(1/p + 1/q - 1)$,则 $T \in M_{q_n q_{n+1}}$.由此,$T^n \in M_{q_0 q_1} M_{q_1 q_2} \cdots M_{q_{n+1} q_n}$(其中 $q_0 = p$),且 $\|T^n\|_{\mathfrak{M}_{pq_n}} = \|T\|_{\mathfrak{M}_{pq}}^n$($n = 1, \ldots, n_0$),所以由 $q_{n_0 + 1} \geq 2 > q_{n_0}$ 以及 $\|T^n\|_{\mathfrak{M}_{pq_n}} = \|T^n\|_{\mathfrak{M}_{q_n'p'}}$,$q_n < 2$,令 $1/s_n = (1 - \tau)/p + \tau/q_n'$($1/2 = (1 - \tau)/q_n + \tau/p'$),得 $T^n \in M_{s_n2}$.
\end{proof}

\begin{lemma}\label{lem:3.5.30}
设 $1 < p \leq 2 \leq q \leq p'$.取 $\{a(n)\} \in \ell^{\infty}$,$T \in M_{pq}$,令 $\widehat{T}a(n) = \widehat{T}(n)a(n)$,则下式成立:

(1) $\|Ta\|_{\mathfrak{M}_{p2}} \leq \|\{a(n)\}\|_{\ell^{\infty}}\|T\|_{\mathfrak{M}_{pq}}$.

(2) $\|T^2a\|_{\mathfrak{M}_{pp'}} \leq \|\{a(n)\}\|_{\ell^{\infty}}\|T\|_{\mathfrak{M}_{pq}}^{2}$.
\end{lemma}

\begin{proof}
(1) 设 $f \in L^{p}(\mathbf{T})$,则
\begin{equation*}
\|Taf\|_{2} \leq \|\{a(\cdot)\}\|_{\ell^{\infty}}\|Tf\|_{2} \leq \|\{a(\cdot)\}\|_{\ell^{\infty}}\|T\|_{\mathfrak{M}_{p2}}\|f\|_{p}.
\end{equation*}
这里,由 $\|T\|_{M_{p2}} \leq \|T\|_{M_{pq}}$ 得 $\|Ta\|_{M_{p2}} \leq \|\{a(n)\}\|_{\ell^{\infty}}\|T\|_{M_{pq}}$.

(2) 取 $f \in L^{p}(\mathbf{T})$,
\begin{equation*}
\|T^{2}af\|_{p'} \leq \|T\|_{\mathfrak{M}_{2p'}}\|Taf\|_{2} \leq \|T\|_{\mathfrak{M}_{2p'}}\|Ta\|_{\mathfrak{M}_{p2}}\|f\|_{p}.
\end{equation*}
这里,利用 $\|T\|_{M_{2p^{\prime}}}=\|T\|_{M_{p2}}$ 与 (1),得 $\|T^{2}af\|_{p'}\leq\|\{a(n)\}\|_{\ell^{\infty}}\|T\|_{M_{pq}}^{2}$.
\end{proof}

\begin{proof}[定理~\ref{thm:3.5.25} 的证明]
(1) 不妨设 $\Phi_0$ 是有界函数.设 $1 < p < q < 2$ 且 $n_0 \geq 2$.此时不妨设 $\alpha_1 = \cdots = \alpha_{n_0} = 0$,故设 $\Phi(t) = |t|^{\beta_0 + 1}\Phi_0(t)$.现在,对 $0 \leq \operatorname{Re} z \leq 1$ 定义算子 $S^z$ 为
\begin{equation*}
\widehat{S^z}(n) = \operatorname{sgn}\widehat{T}(n)\left|\widehat{T}(n)\right|^{(n_0-1)z+1}\Phi_{0}(\widehat{T}(n)).
\end{equation*}
此时,
\begin{equation*}
\|S^{z}\|_{\mathfrak{M}_{s_{1}2}} \leq \|\Phi_{0}\|_{\infty}\|T\|_{\mathfrak{M}_{s_{1}2}} \leq \|\Phi_{0}\|_{\infty}\|T\|_{\mathfrak{M}_{pq}} \quad (\operatorname{Re}z = 0)
\end{equation*}
以及
\begin{equation*}
\|S^{z}\|_{\mathfrak{M}_{s_{n_0}2}} \leq \|\Phi_{0}\|_{\infty}\|T^{n_{0}}\|_{\mathfrak{M}_{s_{n_0}2}} \leq \|\Phi_{0}\|_{\infty}\|T\|_{\mathfrak{M}_{pq}}^{n_{0}} \quad (\operatorname{Re}z = 1)
\end{equation*}
成立.由此利用 Stein 插值定理,令 $\theta$ 满足 $1/q = (1-\theta)/s_{1} + \theta/s_{n_{0}}$,则
\begin{equation*}
\|S^{\theta}\|_{\mathfrak{M}_{q2}} \leq \|\Phi_{0}\|_{\infty}\|T\|_{\mathfrak{M}_{pq}}^{1+\theta(n_0-1)}.
\end{equation*}
注意这里的 $\theta$ 取法满足 $\theta(n_{0}-1)+1=\beta_{0}$.换言之,$\widehat{S^{\theta+1}}(n)=\left|\widehat{T}(n)\right|^{\beta_{0}+1}\Phi_{0}(\widehat{T}(n))$,
\begin{equation*}
\|S^{\theta+1}\|_{\mathfrak{M}_{pq}} \leq \|S^{\theta+1}\|_{\mathfrak{M}_{p2}} \leq \|T\|_{\mathfrak{M}_{pq}}\|S^{\theta}\|_{\mathfrak{M}_{q2}} \leq \|\Phi_{0}\|_{\infty}\|T\|_{\mathfrak{M}_{pq}}^{\beta_{0}+1}.
\end{equation*}
这样就证明了 $n_{0} \geq 2$ 的情形.$n_{0} = 1$ 的情形也同样证明.事实上,令 $\widehat{R^{z}}(n) = \operatorname{sgn}\widehat{T}(n)|\widehat{T}(n)|^{z}\Phi_{0}(\widehat{T}(n))$,由 $\|R^{iy}\|_{\mathfrak{M}_{22}} \leq \|\Phi_{0}\|_{\infty}$,$\|R^{1+iy}\|_{\mathfrak{M}_{s_{1}2}} \leq \|\Phi_{0}\|_{\infty}\|T\|_{\mathfrak{M}_{pq}}$,令 $1/q = (1-\theta)/2 + \theta/s_{1}$,则 $\|R^{\theta}\|_{\mathfrak{M}_{q2}} \leq \|\Phi_{0}\|_{\infty}\|T\|_{\mathfrak{M}_{pq}}^{\theta}$.此时注意 $\theta = \beta_{0}$.其余与 $n_{0} \geq 2$ 的情形同样证明.

(2) 设 $T \in \mathfrak{M}_{pq}$,$\widehat{T}(\mathbf{Z}) \subset [-1, 1]$.因 $\|T\|_{\mathfrak{M}_{pq}} = \|T\|_{\mathfrak{M}_{q'p'}}$(cf. \cite{ch3:9}),只需考虑 $p' \geq q$ 的情形.设 $\Phi(t) = |t|^{\beta_1 + 1}\Phi_0(t)$.定义算子 $R^z$ 为 $\widehat{R^z}(n) = \operatorname{sgn}\widehat{T}(n)|\widehat{T}(n)|^z\Phi_0(\widehat{T}(n))$,把引理~\ref{lem:3.5.30} 用于 $a(n) = \widehat{R^z}(n)$,得
\begin{equation*}
\|R^{z}\|_{\mathfrak{M}_{p2}} \leq \|\Phi_{0}\|_{\infty}\|T\|_{\mathfrak{M}_{pq}} \quad (\operatorname{Re}z = 0)
\end{equation*}
以及
\begin{equation*}
\|R^{z}\|_{\mathfrak{M}_{pp'}} \leq \|\Phi_{0}\|_{\infty}\|T\|_{\mathfrak{M}_{pq}}^{2} \quad (\operatorname{Re}z = 1).
\end{equation*}
利用 Stein 插值定理,令 $1/q = (1 - \theta)/2 + \theta/p'$($\theta = \beta_1$),则得 $\|TR^{\theta}\|_{\mathfrak{M}_{pq}} \leq \|\Phi_0\|_\infty\|T\|_{\mathfrak{M}_{pq}}^{\theta+1}$.这样就证明了 (2).
\end{proof}

其次证明定理~\ref{thm:3.5.25} 的 (2) 是必要条件.此时,先注意:由 $\Phi$ 作用于 $M(p,q)$ 自然可导出 $\Phi(0)=0$.

\begin{theorem}\label{thm:3.5.31}
设 $1 < p < 2 \leq q < \infty$,$\Phi$ 为 $[-1, 1]$ 上的函数.若 $\Phi$ 作用于 $M(p,q)$,则 $\Phi(t) = \alpha t + |t|^{\beta_{1}+1}\Phi_{0}(t)$.其中 $\alpha$ 是复数,$\Phi_{0}$ 是有界函数,$\beta_{1}$ 是定理~\ref{thm:3.5.25} (2) 中的量.
\end{theorem}

\begin{proof}
如引理~\ref{lem:3.5.30} 那样,不妨设 $1/2 - 1/q \leq 1/p - 1/2$.如下分若干阶段证明.

第 1 阶段.若 $\Phi$ 作用于 $M(p,q)$,则存在两个正数 $C$ 和 $\eta$,使得当 $\|T\|_{\mathfrak{M}_{pq}} < \eta$ 时 $\|\Phi \circ T\|_{\mathfrak{M}_{pq}} \leq C$ 成立.

事实上,假设对任意自然数 $m$,存在 $T_{m} \in \mathfrak{M}_{pq}$,$\widehat{T_{m}}(\mathbf{Z}) \subset [-1, 1]$,使得 $\|T_{m}\|_{\mathfrak{M}_{pq}} < 1/m^{2}$,$\|\Phi \circ T_{m}\|_{\mathfrak{M}_{pq}} > m$.此时不妨设 $\supp\widehat{T_{m}}$ 是有限集,即 $T_{m}$ 是三角多项式.设 $T_{m}$ 为 $N_{m}$ 次三角多项式.此时,令 $T(x) = \sum_{m=1}^{\infty}e^{in_{m}x}T_{m}(x)$,取 $\{n_{m}\}_{m=1}^{\infty}$ 满足 $n_{m} + 3N_{m} < n_{m+1} - 3N_{m+1}$.由 $T_{m}$ 的构造,$T \in M_{pq}$,故 $\Phi(T) \in M_{pq}$.于是,与 3.3 节的 $V_{\varepsilon}$ 同样地,把 $V_{N_{m}}$ 定义为以全体实数为定义域的函数.利用这个 $V_{N_{m}}$,令 $H_{m}(x) = e^{in_{m}x}V_{N_{m}}(x)$,则 $((\Phi \circ T)H_{m})(x) = e^{in_{m}x}(\Phi \circ T_{m})(x)$.由此得 $\|\Phi \circ T_{m}\|_{M_{pq}} \leq \|H_{m}\|_{1}\|\Phi \circ T\|_{M_{pq}} \leq 3\|\Phi(T)\|_{M_{pq}}$.从而这与 $\{T_{m}\}$ 的取法矛盾.

这里,若 $\Phi$ 作用于 $M(p,q)$,则奇函数 $\Phi(t)-\Phi(-t)$ 与偶函数 $\Phi(t)+\Phi(-t)$ 也作用.因 $\Phi$ 是这两个函数的和,下面讨论这两种情形.

第 2 阶段.证明 $\Phi$ 为偶函数时满足定理的条件.为此,先设如下引理.

\begin{lemma}\label{lem:3.5.32}
对素数 $r(\geq 2)$,存在取值 $\delta^{(r)}(n)=r-1$ 或 $-1$ 的数列 $\{\delta^{(r)}(n)\}$,使得
\begin{equation*}
\left|\sum_{n=1}^{N}\delta^{(r)}(n)e^{inx}\right| < (r-1)r(1+\sqrt{r})\sqrt{N} \quad (0 \leq x \leq 2\pi,\ N = 1, 2, \ldots).
\end{equation*}
此时令 $\Delta_{N}^{(r)}(x)=\sum_{n=1}^{N}\delta^{(r)}(n)e^{inx}$.$\Delta_{N}^{(2)}(x)$ 被称为 Rudin--Shapiro 多项式.
\end{lemma}

\begin{proof}
该引理分两步证明.

步骤 1(cf. \cite{ch3:41}).设 $r$ 为满足 $r \geq 2$ 的自然数,$\alpha = e^{2\pi i/r}$,适当取 $\varepsilon_{r}(n)$($\varepsilon_{r}(n) = 1, \alpha, \ldots, \alpha^{r-1}$),则
\begin{equation*}
\left|\sum_{n=1}^{N}\varepsilon_{r}(n)e^{in\theta}\right| \leq r(1+\sqrt{r})N^{1/2} \quad (0 \leq \theta \leq 2\pi,\ N \geq 1).
\end{equation*}
事实上,设 $A_{0}, A_{1}, \ldots, A_{r-1}$ 为任意复数,利用 $\alpha^{sr} = 1$($s = 0, 1, \ldots, r-1$),
\begin{equation}\tag{3.5.1}
\sum_{s=0}^{r-1}\left|\sum_{j=0}^{r-1}\alpha^{sj}A_{j}\right|^{2} = r\sum_{j=0}^{r-1}|A_{j}|^{2}
\label{eq:3.5.1}\end{equation}
成立.其次,把多项式 $P_{k}^{s}(z)$($k = 0, 1, \ldots$;$s = 0, 1, \ldots, r-1$)定义为 $P_{0}^{0}(z) = P_{0}^{1}(z) = \cdots = P_{0}^{r-1}(z) = z$,并对 $k = 0, 1, 2, \ldots$ 令 $P_{k+1}^{s}(z) = \sum_{j=0}^{r-1}z^{jr^{k}}\alpha^{sj}P_{k}^{j}(z)$($s = 0, 1, \ldots, r-1$),则 $P_{k}^{s}(z)$ 是 $r^{k}$ 次多项式.此时,把 $\varepsilon_{r}(n)$ 取为 $P_{k}^{0}(e^{ix})$ 的第 $n$ 个 Fourier 系数,则知它具有所求的性质.事实上,如下证明.由 $P_{k}^{s}(z)$ 的构造,当 $z = e^{ix}$ 时,利用 \eqref{eq:3.5.1},
\begin{align*}
\sum_{s=0}^{r-1}\left|P_{k+1}^{s}(z)\right|^{2} &= \sum_{s=0}^{r-1}\left|\sum_{j=0}^{r-1}z^{jr^{k}}\alpha^{sj}P_{k}^{j}(z)\right|^{2}\\
&= r\sum_{j=0}^{r-1}\left|z^{jr^{k}}P_{k}^{j}(z)\right|^{2}\\
&\vdots\\
&= r^{k+2}.
\end{align*}
因此
\begin{equation}\tag{3.5.2}
\left|P_{k}^{s}(e^{ix})\right| \leq r^{1/2}r^{k/2} \quad (k = 0, 1, \ldots;\ s = 0, 1, \ldots, r-1)
\label{eq:3.5.2}\end{equation}
得证.特别地,当 $s=0$ 时,$\sum_{n=1}^{r^{k}}\varepsilon_{r}(n)e^{inx}=P_{k}^{0}(e^{ix})$,若 $N=r^{k}$,则步骤 1 的结果得证.

其次,当 $1 \leq n \leq r^{k}$ 时,令 $S_{n}(P_{k}^{s})(x) = \sum_{|m| \leq n}\widehat{P_{k}^{s}}(m)e^{imx}$,则
\begin{equation}\tag{3.5.3}
\left|S_{n}(P_{k}^{s})\right| \leq \sqrt{r}(1+\sqrt{r})r^{k/2} \quad (s = 0, 1, \ldots, r-1)
\label{eq:3.5.3}\end{equation}
可证明.这是通过对 $k$ 的数学归纳法,利用 \eqref{eq:3.5.2} 证明的.从而对 $N \geq 1$,取满足 $r^{k-1} \leq N \leq r^{k}$ 的 $k$,由 $\left|\sum_{n=1}^{N}\varepsilon_{r}(n)e^{inx}\right| = \left|S_{n}(P_{k}^{0})\right|$,从 \eqref{eq:3.5.3} 得
\begin{equation*}
\left|\sum_{n=1}^{N}\varepsilon_{r}(n)e^{in\theta}\right| \leq r(1+\sqrt{r})N^{1/2}.
\end{equation*}

步骤 2.叙述引理~\ref{lem:3.5.32} 的证明.在步骤 1 中,取 $\alpha^{t}$($t=1,2,\ldots,r-1$)代替 $\alpha$,因 $r$ 是素数,$\alpha^{t}$ 是 1 的 $r$ 次根.把 $\alpha^{t}$ 作为 $\alpha$ 按步骤 1 那样构造多项式 $P_{k}^{0}$,则 $\widehat{P_{k}^{0}}(n)=(\varepsilon_{r}(n))^{t}$.这里,令
\begin{equation*}
\delta^{(r)}(n) = \sum_{t=1}^{r-1}(\varepsilon_{r}(n))^{t},
\end{equation*}
则 $\delta^{(r)}(n)=r-1$ 或 $-1$.从而
\begin{align*}
\left|\sum_{n=1}^{N}\delta^{(r)}(n)e^{inx}\right| &= \left|\sum_{n=1}^{N}\left(\sum_{t=1}^{r-1}(\varepsilon_{r}(n))^{t}\right)e^{inx}\right|\\
&\leq \sum_{t=1}^{r-1}\left|\sum_{n=1}^{N}(\varepsilon_{r}(n))^{t}e^{inx}\right|\\
&\leq (r-1)r(1+\sqrt{r})N^{1/2}.
\end{align*}
\end{proof}

\begin{lemma}\label{lem:3.5.33}
利用引理~\ref{lem:3.5.32} 的记号,
\begin{equation*}
\|\Delta_{N}^{(r)}\|_{\mathfrak{M}_{pq}} \leq (r-1)r(1+\sqrt{r})N^{1/p - 1/2}
\end{equation*}
成立.
\end{lemma}

\begin{proof}
由乘子的性质,$\|\Delta_{N}^{(r)}\|_{\mathfrak{M}_{22}}=\|\{\delta^{(r)}(n)\}\|_{l^{\infty}}\leq r-1$,$\|\Delta_{N}^{(r)}\|_{\mathfrak{M}_{1s}}\leq\|\Delta_{N}^{(r)}\|_{s}\leq\|\Delta_{N}^{(r)}\|_{\infty}\leq(r-1)r(1+\sqrt{r})\sqrt{N}$ 成立,故取 $0<\theta<1$ 使 $1/p=\theta/1+(1-\theta)/2$,再取 $s$ 使 $1/q=\theta/s+(1-\theta)/2$,利用 Riesz--Thorin 插值定理即可.
\end{proof}

第 2 阶段的证明.利用引理~\ref{lem:3.5.33},取 $r = 2$,$\Delta_N^{(2)}N^{-(1/p - 1/2)}$ 在 $\mathfrak{M}_{pq}$ 中(在 $N$ 无关的意义下)一致有界.从而,存在常数 $C > 0$ 使得
\begin{equation*}
\left\|\Phi \circ \left(\frac{\Delta_{N}^{(2)}}{C N^{1/p - 1/2}}\right)\right\|_{\mathfrak{M}_{pq}} \leq C.
\end{equation*}
由此,因 $\Phi$ 是偶函数,
\begin{equation*}
\left\|\sum_{n=1}^{N}\Phi\left(\frac{1}{C N^{1/p - 1/2}}\right)e^{inx}\right\|_{q} \leq C\left\|\sum_{n=1}^{N}e^{inx}\right\|_{q}.
\end{equation*}
因此
\begin{equation*}
\left|\Phi\left(\frac{1}{C N^{1/p - 1/2}}\right)\right|N^{1 - 1/q} \leq C N^{1 - 1/p}.
\end{equation*}
此时,把 $\Phi(t)$ 换成 $\Phi(\eta t)$($1/2 < \eta < 1$)也同理可行,所以在 0 的邻域上存在常数 $C > 0$ 使得
\begin{equation*}
|\Phi(t)t^{-1}| \leq C|t|^{\beta_{1}}.
\end{equation*}
这里,令 $\Phi_{0}(t)=\Phi(t)t^{-(1+\beta_{1})}$($t\neq0$),由 $p>1$,$\Phi_{0}$ 有界.事实上,若 $t_{n}\in[-1,1]$,$|\Phi_{0}(t_{n})|\to\infty$,取 $m(k)=t_{n}$($k=2^{n}$),$m(k)=0$($k\neq2^{n}$)的 $m=\{m(k)\}$,由 $\Phi(m(\cdot))\in M(p,q)\subset M(2,2)$,$\Phi(m(\cdot))\in\ell^{\infty}$,故得矛盾(cf. \cite{ch3:11}, \cite{ch3:26}).又,若 $\Phi(t)$ 是奇函数,用同样的方法可证明存在常数 $C>0$ 使得 $|\Phi(t)|\leq C|t|$.这个结果后面使用.以上证明了偶函数时 $\Phi$ 作用于 $M(p,q)$ 的必要条件.

第 3 阶段.若 $\Phi$ 是奇函数,则
\begin{equation*}
\left|\Phi(t) - 2\Phi\left(\frac{t}{2}\right)\right| \leq C|t|^{1 + \beta_{1}}.
\end{equation*}
事实上,如下证明.令
\begin{align*}
\Phi\left(\frac{t}{2}\delta^{(3)}(n)\right) &= \frac{1}{3}\left[\Phi(t) - 2\Phi\left(\frac{t}{2}\right)\right] + \frac{1}{3}\left[\Phi(t) + \Phi\left(\frac{t}{2}\right)\right]\delta^{(3)}(n)\\
&= \Psi_{1}(t) + \Psi_{2}(t)\delta^{(3)}(n).
\end{align*}
此时,由第 1 阶段,引理~\ref{lem:3.5.33} 与第 2 阶段后半的注意,不妨设
\begin{equation*}
\left\|\Phi \circ \left(\frac{\Delta_{N}^{(3)}}{C N^{1/p - 1/2}}\right)\right\|_{\mathfrak{M}_{pq}} \leq C, \quad |\Phi(t)| \leq C|t|.
\end{equation*}
从而,注意 $\Psi_{2}$ 也是作用函数,利用引理~\ref{lem:3.5.33},
\begin{equation*}
\begin{aligned}
\left\|\sum_{n=1}^{N}\Psi_{2}\left(\frac{1}{C N^{1/p - 1/2}}\right)\delta^{(3)}(n)e^{inx}\right\|_{\mathfrak{M}_{pq}} &= \left|\Psi_{2}\left(\frac{1}{C N^{1/p - 1/2}}\right)\right|\|\Delta_{N}^{(3)}\|_{\mathfrak{M}_{pq}} \leq C.
\end{aligned}
\end{equation*}
因此
\begin{equation*}
\left\|\sum_{n=1}^{N}\Psi_{1}\left(\frac{1}{C N^{1/p - 1/2}}\right)e^{inx}\right\|_{M(p,q)} \leq C \quad (N = 1, 2, \ldots).
\end{equation*}
换言之,
\begin{equation*}
\left\|\sum_{n=1}^{N}\Psi_{1}\left(\frac{1}{C N^{1/p - 1/2}}\right)e^{inx}\right\|_{q} \leq C\left\|\sum_{n=1}^{N}e^{inx}\right\|_{p}.
\end{equation*}
从而
\begin{equation*}
\left|\Psi_{1}\left(\frac{1}{C N^{1/p - 1/2}}\right)\right|N^{1 - 1/q} \leq C N^{1 - 1/p},
\end{equation*}
所以
\begin{equation*}
|\Psi_{1}(t)| \leq C t^{\frac{1/p - 1/q}{1/p - 1/2}}.
\end{equation*}
综上所述,第 3 阶段得证.

这里,注意第 3 阶段表明:对 $t \in [-1, 1]$,$\{2^{n}\Phi\left(\frac{t}{2^{n}}\right)\}_{n}$ 是 Cauchy 列.于是令
\begin{equation*}
\Phi_{1}(t) = \lim_{n \to \infty}2^{n}\Phi\left(\frac{t}{2^{n}}\right), \quad \Phi_{2}(t) = \Phi(t) - \Phi_{1}(t).
\end{equation*}

第 4 阶段.$|\Phi_2(t)|\leq C|t|^{1 + \beta_1}$.

第 4 阶段如下证明.利用第 3 阶段的结果,
\begin{align*}
|\Phi_{2}(t)| &= \left|\lim_{n \to \infty}2^{n}\Phi\left(\frac{t}{2^{n}}\right) - \Phi(t)\right|\\
&\leq \sum_{n=1}^{\infty}\left|2^{n}\Phi\left(\frac{t}{2^{n}}\right) - 2^{n-1}\Phi\left(\frac{t}{2^{n-1}}\right)\right|\\
&\leq \sum_{n=1}^{\infty}C'2^{n-1}\left|\frac{t}{2^{n-1}}\right|^{1+\beta_{1}}\\
&\leq C|t|^{1+\beta_{1}}.
\end{align*}
即可证明.这样第 4 阶段得证.

第 5 阶段.$\Phi_{1}(t)$ 是 $[-1, 1]$ 上的连续函数.

事实上,如下证明.必要时把 $\Phi_{1}(t)$ 换成 $\Phi_{1}(\eta t)$,只需证明 $\Phi_{1}(t)$ 在 $t=1$ 处连续.设 $\{t_{n}\}$ 为满足 $|1-t_{m}|<2^{-m}$ 的数列,令
\begin{equation*}
N_{m} = \left[2^{m/(\frac{1}{p} - \frac{1}{2})}\right], \quad T_{m} = \frac{1}{2^{m}}\sum_{n=1}^{N_{m}}\left\{\delta^{(2)}(n) + \frac{1 - t_{m}}{2}(1 - \delta^{(2)}(n))\right\}e^{inx}.
\end{equation*}
此时
\begin{equation*}
\|T_{m}\|_{\mathfrak{M}_{pq}} \leq 2^{-m}\left\|\sum_{n=1}^{N_{m}}\delta^{(2)}(n)e^{inx}\right\|_{\mathfrak{M}_{pq}} + C2^{-2m}\left\|\sum_{n=1}^{N_{m}}e^{inx}\right\|_{\mathfrak{M}_{pq}}.
\end{equation*}
右边第 1 项由 $N_{m}$ 的取法与引理~\ref{lem:3.5.33} 关于 $m$ 一致有界.右边第 2 项可用 $2^{-2m}2^{m/(\frac{1}{p}-\frac{1}{2})\cdot(\frac{1}{p}-\frac{1}{q})}$ 控制,由 $1/p+1/q\geq1$ 关于 $m$ 一致有界(cf. \cite{ch3:40}).可是,因 $\Phi_{1}$ 是奇函数,
\begin{align*}
&\Phi_{1}\left(\delta^{(2)}(n) + \frac{1 - t_{m}}{2}(1 - \delta^{(2)}(n))\right)\\
&\qquad = \frac{1}{2}\left[\Phi_{1}(1) - \Phi_{1}(t_{m})\right] + \frac{1}{2}\left[\Phi_{1}(1) + \Phi_{1}(t_{m})\right]\delta^{(2)}(n).
\end{align*}
利用 $2\Phi_{1}\left(\frac{t}{2}\right)=\Phi_{1}(t)$,若取 $\|T_{m}\|_{\mathfrak{M}_{pq}}\leq2^{m_{0}}$,则从上式
\begin{align*}
C &\geq \|\Phi_{1}(2^{-m_{0}}T_{m})\|_{\mathfrak{M}_{pq}}\\
&= 2^{-m-m_{0}}\|\Phi_{1}(2^{m}T_{m})\|_{\mathfrak{M}_{pq}}\\
&\geq C|\Phi_{1}(1) - \Phi_{1}(t_{m})|2^{-m}N_{m}^{1/p - 1/q} - C2^{-m}N_{m}^{1/p - 1/2},
\end{align*}
且因 $2^{-m}N_{m}^{1/p-1/q}\to\infty$($m\to\infty$),必须 $\Phi_{1}(t_{m})\to\Phi_{1}(1)$($m\to\infty$).因此 $\Phi_{1}(t)$ 在 $t=1$ 处连续.以上第 5 阶段得证.

第 6 阶段.$\Phi_1(t) = \alpha t$($t\in [-1,1]$)成立.

第 6 阶段的证明如下进行.

设 $0 < t < 1$,$r$ 为素数.由 $\Phi_{1}(t)$ 是奇函数,
\begin{equation}\tag{3.5.4}\begin{aligned}
\Phi_{1}\left(\frac{t\delta^{(r)}(n)}{2^{m}(r-1)}\right) &= \frac{1}{2^{m}r}\left[\Phi_{1}(t) - (r-1)\Phi_{1}\left(\frac{t}{r-1}\right)\right]\\
&\quad + \frac{1}{2^{m}r}\left[\Phi_{1}(t) + \Phi_{1}\left(\frac{t}{r-1}\right)\right]\delta^{(r)}(n)
\end{aligned}\label{eq:3.5.4}\end{equation}
成立.可是,由引理~\ref{lem:3.5.33},
\begin{equation*}
\left\|t2^{-m}(r-1)^{-1}\Delta_{N_{m}}^{(r)}\right\|_{M(p,q)} \leq Ct2^{-m}(r-1)^{-1}N_{m}^{1/p - 1/2},
\end{equation*}
故令 $N_{m} = [2^{m/(1/p-1/2)}]$,把自然数 $m_{0}$ 取得足够大,\eqref{eq:3.5.4} 右边第 2 项的算子范数即可用 $2^{m_{0}}$ 控制,
\begin{equation*}
\left\|\Phi_{1}\left(\frac{t\Delta_{N}^{(r)}}{2^{m+m_{0}}(r-1)}\right)\right\|_{\mathfrak{M}_{pq}} \leq C
\end{equation*}
($C$ 可以依赖于 $r$).从而
\begin{align*}
2^{m_{0}}C &\geq \frac{1}{2^{m}r}\left|\Phi_{1}(t) - (r-1)\Phi_{1}\left(\frac{t}{r-1}\right)\right|N_{m}^{1/p - 1/q}\\
&\quad - \frac{C}{2^{m}r}\left|\Phi_{1}(t) + \Phi_{1}\left(\frac{t}{r-1}\right)\right|N_{m}^{1/p - 1/2}.
\end{align*}
因 $2^{-m}N_{m}^{1/p-1/2}$ 有界且 $2^{-m}N_{m}^{1/p-1/q}\to\infty$($m\to\infty$),必须
\begin{equation*}
\Phi_{1}(t) = (r-1)\Phi_{1}\left(\frac{t}{r-1}\right).
\end{equation*}
反复利用这一点,设 $q, r$ 为素数,则
\begin{equation*}
\Phi_{1}\left(\left(\frac{r-1}{q-1}\right)^{n}\right) = \left(\frac{r-1}{q-1}\right)^{n}\Phi_{1}(1).
\end{equation*}
这里,因 $\left\{(\frac{r-1}{q-1})^{n}: r, q \text{ 是素数, } n \text{ 是自然数}\right\}$ 在 $(0, 1)$ 中稠密且 $\Phi_{1}(t)$ 连续,$\Phi_{1}(t)=t\Phi_{1}(1)$($0\leq t\leq1$).以上第 6 阶段得证.

综上所述,全部证明完毕.
\end{proof}

\begin{note}\label{note:3.5.34}
定理~\ref{thm:3.5.31} 的 $q = 2$ 的情形,由 Hatori \cite{ch3:19} 推广到了无限紧 Abel 群上.另外,$M_{pq}$ 的极大理想空间在 Hatori \cite{ch3:20} 中作了研究.
\end{note}

\subsection{\texorpdfstring{$C_0 M_p(\mathbb{Z})$($1<p<2$)的作用函数}{C0 M_p(Z)(1<p<2)的作用函数}}\label{sec:3.5.6}

这里,关于 $\mathfrak{M}_{p}(\mathbf{T})$($1 < p < 2$)的子空间 $C_{0}\mathfrak{M}_{p}(\mathbf{T})$ 的作用函数,只叙述 M. Zafran 的研究 \cite{ch3:59}(cf. \cite{ch3:58})的结果.详情请参阅论文.

\begin{definition}\label{def:3.5.35}
设 $1 < p < 2$.$T \in C_{0}\mathfrak{M}_{p}(\mathbf{T})$ 对应于 $\{\widehat{T}(n)\}_{n \in \mathbf{Z}} \in \ell^{\infty}$(即 $\widehat{T}f(n) = \widehat{T}(n)\widehat{f}(n)$,$f \in L^{2}(\mathbf{T})$),但因不会产生混乱,把 $C_{0}\mathfrak{M}_{p}(\mathbf{T})$ 等同于 $C_{0}M_{p}(\mathbf{Z}) = \{\widehat{T}: T \in C_{0}\mathfrak{M}_{p}(\mathbf{T})\}$.设 $\Phi$ 为 $[-1, 1]$ 上的函数.当 $\widehat{T} \in C_{0}M_{p}(\mathbf{Z})$,$\widehat{T}(\mathbf{Z}) \subset [-1, 1]$ 时,若 $\Phi(\widehat{T}) \in C_{0}M_{p}(\mathbf{Z})$,则称 $\Phi$ 作用于 $C_{0}\mathfrak{M}_{p}(\mathbf{T})$.
\end{definition}

现在,在 3.5.4 小节中证明了满足 $\operatorname{Sp}(\mu, \mathfrak{M}_p(\mathbf{T})) \supsetneqq \overline{\widehat{\mu}(\mathbf{Z})}$ 的 $\mu \in M(\mathbf{T})$ 的存在,而 M. Zafran 注意到:对此 $\mu$,$\hat{\mu}$ 在无穷远点不趋于零.事实上,下式成立.

\begin{theorem}[{$\;$\cite{ch3:22}}]\label{thm:3.5.36}
设 $1 \leq q < p \leq 2$.若 $T \in C_{0}\mathfrak{M}_{q}(\mathbf{T})$,则
\begin{equation*}
\operatorname{Sp}(T, \mathfrak{M}_{p}) = \hat{T}(\mathbf{Z}) \cup \{0\}
\end{equation*}
成立.
\end{theorem}

这是由 Riesz--Thorin 插值定理证明的.关于 3.5.4 小节所述的 Hörmander 问题``$m_p = C_0\mathfrak{M}_p(\mathbf{T})$ ?'',Figa-Talamanca--Gaudry \cite{ch3:18} 给予了否定解决.他们证明:当 $1 < p < 2$ 时,存在 $T_0 \in C_0\mathfrak{M}_p(\mathbf{T})$ 使得 $\|T_0 - T_f\|_{\mathfrak{M}_p} \geq c(>0)$($f \in L^1(\mathbf{T})$).但对这个 $T_0$,Zafran \cite{ch3:57} 证明了 $T_0^2 \in m_p$ 成立,指出 $T_0$ 距 $m_p$ 并不太远.进而,M. Zafran 指出了对应于 $p = 1$ 的情形的 $C_0\mathfrak{M}_1(\mathbf{T}) = M_0(\mathbf{T}) = \{\mu \in M(\mathbf{T}): \hat{\mu} \to 0 \ (|n| \to \infty)\}$ 的作用函数的 Varopoulos 研究 \cite{ch3:51} 的重要性.Varopoulos \cite{ch3:51} 注意到:由 $L^1(\mathbf{T}) \subset M_0(\mathbf{T})$,$M_0(\mathbf{T})$ 的作用函数 $\Phi$ 作用于 $L^1(\mathbf{T})$,故可在 0 的邻域上级数展开;并利用发现 $M_0(\mathbf{T})$ 的特别测度的方法,证明了当 $\Phi(z) = \sum_{n=0}^{\infty}a_nz^n$($|z| < \varepsilon_0$,$\varepsilon_0 > 0$)时,右边级数的收敛半径可以取得任意大.M. Zafran 把 Varopoulos \cite{ch3:51} 的这一方法用于 $C_0\mathfrak{M}_p(\mathbf{T})$ 的作用函数的研究,证明了下面的结果.

\begin{theorem}\label{thm:3.5.37}
设 $1 < p < 2$.$[-1, 1]$ 上的函数 $\Phi$ 作用于 $C_{0}\mathfrak{M}_{p}(\mathbf{T})$ 的充要条件是:$\Phi$ 可延拓为整函数.
\end{theorem}

由这个结果可知,在 $C_0\mathfrak{M}_p(\mathbf{T})$ 中也会发生 Wiener--Pitt 现象.即,

\begin{corollary}\label{cor:3.5.38}
存在 $T \in C_{0}\mathfrak{M}_{p}(\mathbf{T})$ 使得
\begin{equation*}
\operatorname{Sp}\left(T, C_{0}\mathfrak{M}_{p}(\mathbf{T})\right) \supsetneqq \overline{\widehat{T}(\mathbf{Z})}
\end{equation*}
成立.
\end{corollary}

关于定理~\ref{thm:3.5.37} 的证明,经 Brown \cite{ch3:13} 的指摘有所改进,但为得到好的估计,需要使用测度的构造,向量值函数的插值定理,Littlewood--Paley 理论等,比较复杂.最后,关于 $\|e^{i\mu}\|_{\mathfrak{M}_{p}(\mathbf{T})}$ 的估计,叙述 M. Zafran 的结果 \cite{ch3:60}.

\begin{theorem}[Zafran]\label{thm:3.5.39}
设 $1 < p < 2$.此时,对 $j \in Z$,
\begin{equation*}
\sup\left\{\|e^{ij\mu}\|_{\mathfrak{M}_{p}(\mathbf{T})}: \mu \in M(\mathbf{T}),\ \hat{\mu} \in \mathbf{R},\ \|\mu\|_{\mathfrak{M}_{p}(\mathbf{T})} \leq 1\right\} \geq c\beta^{|j|}
\end{equation*}
成立.其中,$c>0$,$\beta>1$ 是不依赖于 $p$ 的常数.
\end{theorem}

与定理~\ref{thm:3.5.21} 或注意~\ref{note:3.5.22} 比较可知,这是体现出 $\mathfrak{M}_{p}(\mathbf{T})$ 与 $M(\mathbf{T})$ 的本质大小差异的估计,是非常有趣的结果.

\section*{结语}

如前言所述,第\ref{ch:3}章主要叙述了据称以 N. Wiener 的定理为开端的谱综合与作用函数的话题.关于谱综合,以 P. Malliavin,N. Th. Varopoulos,S. Saeki,T. W. Körner,D. Müller 等人的研究为中心;关于作用函数,则以 J. P. Kahane,W. Rudin,Y. Katznelson,S. Igari,M. Zafran 等人的研究为中心.本章参考了他们的文献.除这里所选取的内容之外,在与交换 Banach 代数的调和分析相关方面,尚有 \cite{ch3:6},\cite{ch3:15},\cite{ch3:21},\cite{ch3:25},\cite{ch3:30},\cite{ch3:36},\cite{ch3:37},\cite{ch3:38},\cite{ch3:46},\cite{ch3:58} 等大量有趣的研究.另外,参考文献的排列顺序为:先书,后论文.

此番,谨向给予我执笔机会的东京女子大学宮地晶彦教授致以衷心的感谢.


% ============ 第3章 参考文献(原书 pp.305-309) ============
\section*{参考文献}\addcontentsline{toc}{section}{参考文献}
\begin{thebibliography}{99}
% 条目照录原文(书名/作者不翻译).页码范围用 en dash.
\bibitem[1]{ch3:1} ブルバキ(倉田令二朗訳),数学原論 スペクトル論 I, 東京図書, 東京, 1971.
\bibitem[2]{ch3:2} 日本数学会 編,関数環とその関連分野(特集),数学 28(1976), no.1.
\bibitem[3]{ch3:3} 荷見守助,調和解析学,槙書店,東京, 1978.
\begin{sloppypar}
\bibitem[4]{ch3:4} J. Bergh and J. Löfström, Interpolation Spaces: An introduction, Springer-Verlag, Berlin-Heidelberg-New York, 1976.
\end{sloppypar}
\bibitem[5]{ch3:5} J. Duoandikoetxea, Fourier analysis, Graduate Studies in Mathematics, A.M.S, 2001.
\bibitem[6]{ch3:6} C. C. Graham, O. C. McGehee, Essays in commutative harmonic analysis, Springer-Verlag, New York-Heidelberg-Berlin, 1979.
\bibitem[7]{ch3:7} J. P. Kahane, Séries de Fourier absolument convergentes, Springer-Verlag, Berlin, 1970.
\bibitem[8]{ch3:8} Y. Katznelson, An introduction to harmonic analysis, Dover Publications, New York, 1976.
\bibitem[9]{ch3:9} R. Larsen, An Introduction to the Theory of Multipliers, Springer-Verlag, Berlin-Heidelberg-New York, 1971.
\bibitem[10]{ch3:10} L. H. Loomis, An introduction to abstract harmonic analysis, Princeton, Van Nostrand, 1953.
\bibitem[11]{ch3:11} W. Rudin, Fourier analysis on groups, Interscience, New York-London, 1962.
\bibitem[12]{ch3:12} J. L. Taylor, Measure algebras, CBMS Conference Report, No. 16, A.M.S, 1973.
\bibitem[13]{ch3:13} G. Brown, Construction of Fourier multipliers, Bull. Austral. Math. Soc. 16(1977), 463--472.
\bibitem[14]{ch3:14} A. P. Calderón and A. Zygmund, Algebras of certain singular operators, Amer. J. Math. 78(1956), 310--320.
\bibitem[15]{ch3:15} Y. Domar, On the spectral synthesis problem for $(n-1)$-dimensional subsets of $R^{n}$, $n \geq 2$, Ark. Mat. 9(1971), 23--37.
\bibitem[16]{ch3:16} R. Doss, Some inclusions in multipliers, Pacific J. Math. 32(1970), 643--646.
\bibitem[17]{ch3:17} K. de Leeuw, On $L_{p}$ multipliers, Ann. of Math. (2) 81(1965), 364--379.
\bibitem[18]{ch3:18} A. Figà-Talamanca and G. I. Gaudry, Multipliers of $L^{p}$ which vanish at infinity, J. Functional Analysis 7(1971), 475--486.
\bibitem[19]{ch3:19} O. Hatori, Functions which operate on algebras of Fourier multipliers, Tohoku Math. J. (2) 47(1995), no. 1, 105--115.
\bibitem[20]{ch3:20} O. Hatori, A characterization of lacunary set and spectral properties of Fourier multipliers, Function spaces (Edwardsville, IL, 1994), 183--203, Lecture Notes in Pure and Appl. Math., 172, Dekker, New York, 1995.
\bibitem[21]{ch3:21} C. S. Herz, Spectral synthesis for the circle, Ann. of Math. 68(1958), 709--712.
\bibitem[22]{ch3:22} L. Hörmander, Estimates for translation invariant operators in $L^{p}$ spaces, Acta Math. 104(1960), 93--140.
\bibitem[23]{ch3:23} S. Igari, スペクトル合成について, 数学 18(1966/1967), 215--222.
\bibitem[24]{ch3:24} S. Igari, Functions of $L^{p}$ multipliers, Tôhoku Math. J. 21(1969), 304--320.
\bibitem[25]{ch3:25} S. Igari, Functions of $L^{p}$-multipliers II, Tôhoku Math. J. (2) 26(1974), 555--561.
\bibitem[26]{ch3:26} S. Igari, E. Sato, Operating functions on Fourier multipliers, Tohoku Math. J. (2) 46(1994), no. 3, 357--366.
\bibitem[27]{ch3:27} M. Jr. Jodeit, Restrictions and extensions of Fourier multipliers, Studia Math. 34(1970), 215--226.
\bibitem[28]{ch3:28} M. Kaneko and E. Sato, Notes on transference of continuity from maximal Fourier multiplier operators on $R^{n}$ to those on $T^{n}$, Interdiscip. Inform. Sci. 4(1998), no. 1, 97--107.
\bibitem[29]{ch3:29} Y. Kanjin, A convolution measure algebra on the unit disc, Tôhoku Math. J. 28(1976), 105--115.
\bibitem[30]{ch3:30} Y. Kanjin, On algebras with convolution structures for Laguerre polynomials, Trans. Amer. Math. Soc. 295(1986), no. 2, 783--794.
\bibitem[31]{ch3:31} Y. Katznelson, Sur les fonctions opérant sur l'algèbre des séries de Fourier absolument convergentes, C. R. Acad. Sci. Paris 247(1958), 404--406.
\bibitem[32]{ch3:32} R. Kaufman, M-sets and distributions, Astérisque 5(1973), 225--230.
\bibitem[33]{ch3:33} T. Körner, A pseudofunction on a Helson set. I, Astérisque 5(1973), 3--224.
\bibitem[34]{ch3:34} P. Malliavin, Impossibilité de la synthèse spectrale sur les groupes abéliens non compacts, Publ. Math. Inst. Hautes Études Sci. 1959, 85--92.
\bibitem[35]{ch3:35} D. Müller, On the spectral synthesis problem for hypersurfaces of $\mathbf{R}^N$, J. Funct. Anal. 47(1982), no. 2, 247--280.
\bibitem[36]{ch3:36} F. Parreau, Measures with real spectra, Invent. Math. 98(1989), no. 2, 311--330.
\bibitem[37]{ch3:37} T. V. Pedersen, Some properties of the Pisier algebra, Math. Proc. Cambridge Philos. Soc. 128(2000), 343--354.
\bibitem[38]{ch3:38} G. Pisier, A remarkable homogeneous Banach algebra, Israel J. Math. 34(1979), no. 1-2, 38--44.
\bibitem[39]{ch3:39} I. Richards, On Malliavin's counterexample to spectral synthesis, Bull. Amer. Math. Soc. 72(1966), 698--700.
\bibitem[40]{ch3:40} D. Rider, Transformations of Fourier coefficients, Pacific J. Math. 19(1966), 347--355.
\bibitem[41]{ch3:41} W. Rudin, Some theorems on Fourier coefficients, Proc. Amer. Math. Soc. 10(1959), 855--859.
\bibitem[42]{ch3:42} W. Rudin, Fourier-Stieltjes transforms of measures on independent sets, Bull. Amer. Math. Soc. 66(1960), 199--202.
\bibitem[43]{ch3:43} W. Rudin, A strong converse of the Wiener--Lévy theorem, Canad. J. Math. 14(1962) 694--701.
\bibitem[44]{ch3:44} S. Saeki, Spectral synthesis for the Kronecker sets, J. Math. Soc. Japan 21(1969), 549--563.
\bibitem[45]{ch3:45} S. Saeki, Helson sets which disobey spectral synthesis, Proc. Amer. Math. Soc. 47(1975), 371--377.
\bibitem[46]{ch3:46} S. Saeki, Tensor products in Harmonic Analysis, KSU Lecture Note.
\bibitem[47]{ch3:47} P. Sarnak, Spectra of singular measures as multipliers on $L_{p}$, J. Functional Analysis 37(1980), 302--317.
\bibitem[48]{ch3:48} E. Sato, Spectra of measures as $L_{p}$ multipliers, Tohoku Math. J. (2) 37(1985), no. 3, 333--342.
\bibitem[49]{ch3:49} L. Schwartz, Sur une propriété de synthèse spectrale dans les groupes non compacts, C. R. Acad. Sci. Paris 227(1948), 424--426.
\bibitem[50]{ch3:50} N. Th. Varopoulos, Sur les ensembles parfaits et les séries trigonométriques, C. R. Acad. Sci. Paris 260(1965), 3831--3834.
\bibitem[51]{ch3:51} N. Th. Varopoulos, The functions that operate on $B_{0}(\Gamma)$ of a discrete group $\Gamma$, Bull. Soc. Math. France 93(1965), 301--321.
\bibitem[52]{ch3:52} N. Th. Varopoulos, Tensor algebras and harmonic analysis, Acta Math. 119(1967), 51--112.
\bibitem[53]{ch3:53} N. Wiener, Tauberian theorems, Ann. of Math. 33(1932), 1--100.
\bibitem[54]{ch3:54} N. Wiener and H. R. Pitt, On absolutely convergent Fourier-Stieltjes transforms, Duke Math. J. 4(1938), no. 2, 420--436.
\bibitem[55]{ch3:55} J. H. Williamson, A theorem on algebras of measures on topological groups, Proc. Edinburgh Math. Soc. 11(1958/1959), 195--206.
\bibitem[56]{ch3:56} K. Wozniakowski, A new proof of the restriction theorem for weak type (1,1) multipliers on $\mathbf{R}^n$, Illinois J. Math. 40(1996), 483.
\bibitem[57]{ch3:57} M. Zafran, On the spectra of multipliers, Pacific J. Math. 47(1973), 609--626.
\bibitem[58]{ch3:58} M. Zafran, The spectra of multiplier transformations on the $L_{p}$ spaces, Ann. of Math. 103(1976), no. 2, 355--374.
\bibitem[59]{ch3:59} M. Zafran, The functions operating on multiplier algebras, J. Functional Analysis 26(1977), no. 3, 289--314.
\bibitem[60]{ch3:60} M. Zafran, Exponential estimates in multiplier algebras, Conference on harmonic analysis in honor of Antoni Zygmund, Vol. I, 114--121, 1983.
\end{thebibliography}

% ============ 第4章 振荡积分与Kakeya问题（田中仁，原书 pp.310-380 = PDF pp.322-392） ============
\chapter{振荡积分与 Kakeya 问题（田中仁）}\label{ch:4}

% ============ 第4章 §4.1 Hardy-Littlewood极大函数与微分定理(原书 pp.310-319) ============
\section{Hardy--Littlewood 极大函数与微分定理}\label{sec:4.1}
\index{Hardy--Littlewood 极大函数}
\index{Young 不等式}
\index{Lebesgue 微分定理}
\index{分布函数}
\index{覆盖引理}


本节中,定义作为调和分析基本工具的 Hardy--Littlewood 极大函数,并考察其性质.

\noindent\textbf{极大函数}\quad 用 $B(x,t)\subset\mathbb{R}^{n}$ 表示以 $x$ 为中心,半径为 $t$ 的开球,用 $|B(x,t)|$ 表示其体积.此时,对局部可积函数 $f \in L_{\mathrm{loc}}^{1}(\mathbb{R}^{n})$,Lebesgue 微分定理
\begin{equation}
\lim_{t \to 0}\frac{1}{|B(x,t)|}\int_{B(x,t)}f(y)\dd y = f(x) \quad (\text{a.e. } x \in \mathbb{R}^{n})\tag{4.1.1}
\label{eq:4.1.1}\end{equation}
成立.为证明此定理,希望考虑 $\lim_{t\to 0}$ 的代替 $\sup_{t>0}$.在此估计中,无需考虑相消(cancellation),Hardy--Littlewood 极大函数
\begin{equation}
\mathcal{M}f(x) := \sup_{t > 0}\frac{1}{|B(x,t)|}\int_{B(x,t)}|f(y)|\dd y\tag{4.1.2}
\label{eq:4.1.2}\end{equation}
就是应当研究的目标.

极大函数 $\mathcal{M}f$ 由其定义立即看出满足如下的次线性性:
\begin{equation}
\begin{cases}
\mathcal{M}(f + g)(x) \leq \mathcal{M}f(x) + \mathcal{M}g(x), \\
\mathcal{M}(cf)(x) = |c|\,\mathcal{M}f(x) \qquad (f, g \in L_{\mathrm{loc}}^{1}(\mathbb{R}^{n}),\ c \in \mathbb{C}).
\end{cases}\tag{4.1.3}
\label{eq:4.1.3}\end{equation}

\noindent\textbf{极大函数的可测性}\quad 设 $f \in L_{\mathrm{loc}}^{1}(\mathbb{R}^{n})$.令\footnote{$\chi_{A}(x)$ 表示集合 $A \subset \mathbb{R}^{n}$ 的特征函数,即 $\chi_{A}(x) = 1$($x \in A$),$\chi_{A}(x) = 0$($x \notin A$).}
\begin{equation*}
\chi_{t}(x) := \frac{\chi_{B(0,t)}(x)}{|B(0,t)|} \quad (t > 0),
\end{equation*}
\begin{equation*}
f_{t}(x) := (\chi_{t} * |f|)(x) = \int_{\mathbb{R}^{n}}\chi_{t}(x - y)|f(y)|\dd y,
\end{equation*}
此时,
\begin{equation*}
\mathcal{M}f(x) = \sup_{t > 0}f_{t}(x).
\end{equation*}

固定 $x$ 时,$f_{t}(x)$ 是 $t > 0$ 的连续函数,故上式右端的上限只要对有理数 $0 < t \in \mathbb{Q}$ 取即可.又固定 $t$ 时,$f_{t}(x)$ 是 $x \in \mathbb{R}^{n}$ 的连续函数.因此,上式右端可由连续函数的可数多个上限表示,从而得到 $\mathcal{M}f$ 的可测性.

\noindent\textbf{分布函数}\quad 为知道函数的相对大小,定义其分布函数.对 $\mathbb{R}^{n}$ 上的可测函数 $g(x)$,其分布函数定义为\footnote{以后,把可测集合 $A \subset \mathbb{R}^{n}$ 的大小(Lebesgue 测度)简记为 $|A|$.}
\begin{equation}
\alpha(\lambda) := |\{x \in \mathbb{R}^{n}: |g(x)| > \lambda\}| \quad (\lambda > 0).\tag{4.1.4}
\label{eq:4.1.4}\end{equation}

用分布函数可以表示 $g \in L^{p}(\mathbb{R}^{n})$ 的范数.

$g \in L^{\infty}$ 时,
\begin{equation*}
\|g\|_{\infty} = \inf\{\lambda : \alpha(\lambda) = 0\}.
\end{equation*}

$g \in L^{p}(\mathbb{R}^{n})$($1 \leq p < \infty$)时,需要稍作计算.首先,将
\begin{align*}
\int_{\mathbb{R}^{n}}|g(y)|^{p}\dd y & = p\int_{\mathbb{R}^{n}}\int_{0}^{|g(y)|}\lambda^{p-1}\dd\lambda\dd y \\
& = p\int_{\mathbb{R}^{n}}\int_{0}^{\infty}\chi_{[0,|g(y)|]}(\lambda)\lambda^{p-1}\dd\lambda\dd y
\end{align*}
改写,用 Fubini 定理交换积分顺序,即得
\begin{equation}
\int_{\mathbb{R}^{n}}|g(y)|^{p}\dd y = p\int_{0}^{\infty}\alpha(\lambda)\lambda^{p-1}\dd\lambda.\tag{4.1.5}
\label{eq:4.1.5}\end{equation}

$g \in L^{p}(\mathbb{R}^{n})$ 时,$\alpha(\lambda)$ 可由 $\|g\|_{p}$ 从上估计,即
\begin{equation}
\alpha(\lambda) \leq \left(\frac{\|g\|_{p}}{\lambda}\right)^{p}.\tag{4.1.6}
\label{eq:4.1.6}\end{equation}
实际上,
\begin{equation*}
\lambda^{p}\alpha(\lambda) \leq \int_{\{x: |g(x)| > \lambda\}}|g(y)|^{p}\dd y \leq \|g\|_{p}^{p}.
\end{equation*}

下面是与极大函数相关的最基本,最重要的定理.作为其推论,得到微分定理.

\begin{theorem}\label{thm:4.1.1}
以下成立.

(a) 若 $f \in L^{p}(\mathbb{R}^{n})$($1 \leq p \leq \infty$),则 $\mathcal{M}f$ 几乎处处有限.

(b) 若 $f \in L^{1}(\mathbb{R}^{n})$,则
\begin{equation*}
|\{x: \mathcal{M}f(x) > \lambda\}| \leq \frac{3^{n}}{\lambda}\|f\|_{L^{1}(\mathbb{R}^{n})} \quad (\lambda > 0).
\end{equation*}

(c) 若 $f \in L^{p}(\mathbb{R}^{n})$($1 < p \leq \infty$),则
\begin{equation*}
\|\mathcal{M}f\|_{L^{p}(\mathbb{R}^{n})} \leq A_{p}\|f\|_{L^{p}(\mathbb{R}^{n})}, \quad A_{p} = 2\left(\frac{3^{n}p}{p-1}\right)^{1/p}.
\end{equation*}

(a) 的成立在 (b),(c) 成立时是显然的.
\end{theorem}

\begin{corollary}\label{cor:4.1.2}
若 $f \in L_{\mathrm{loc}}^{1}(\mathbb{R}^{n})$,则
\begin{equation*}
\lim_{t \to 0}\frac{1}{|B(x,t)|}\int_{B(x,t)}f(y)\dd y = f(x) \quad (\text{a.e. } x \in \mathbb{R}^{n})
\end{equation*}
成立.
\end{corollary}

若用 $d\mu(x)$ 表示只在原点放质量 1 的测度,则在极大函数的定义中用 $d\mu(x)$ 替换 $|f(x)|\dd x$ 来定义 $\mathcal{M}\mu(x)$,则
\begin{equation*}
\mathcal{M}\mu(x) = \frac{C}{|x|^{n}} \notin L^{1}(\mathbb{R}^{n}),
\end{equation*}
\begin{equation*}
|\{x: \mathcal{M}\mu(x) > \lambda\}| = \frac{C}{\lambda}
\end{equation*}
可以算出.因此,在这种情况下 (b) 的估计是最优的.

由于可以选择 $L^{1}(\mathbb{R}^{n})$ 中的函数列 $\{f_{n}\}$ 使得
\begin{equation*}
\lim_{n \to \infty}f_{n}(x)\dd x = d\mu(x),
\end{equation*}
可知 (b) 的估计是最优的.

\noindent\textbf{覆盖引理}\quad 定理的核心是 (b),其证明基于如下的覆盖引理.

\begin{lemma}\label{lem:4.1.3}
设 $B_{1} = B(x_{1},t_{1})$,$B_{2} = B(x_{2},t_{2})$,$\ldots$,$B_{N} = B(x_{N},t_{N})$ 是有限个具有有限大小的球构成的族,令
\begin{equation*}
E := \bigcup_{j=1}^{N}B_{j}.
\end{equation*}
此时,可以从 $\{B_{j}\}$ 中选出互不相交的子族 $B_{j_{1}}, B_{j_{2}}, \ldots, B_{j_{K}}$,使得
\begin{equation*}
|E| \leq 3^{n}\sum_{k=1}^{K}|B_{j_{k}}|.
\end{equation*}
\end{lemma}

\begin{proof}
首先选取 $B_{j_{1}}$ 为满足
\begin{equation*}
t_{j_{1}} = \max\{t_{j}: j = 1, \ldots, N\}
\end{equation*}
的球.其次选取 $B_{j_{2}}$ 为满足
\begin{equation*}
t_{j_{2}} = \max\{t_{j}: B_{j_{1}} \cap B_{j} = \emptyset\}
\end{equation*}
的球.以下归纳地选取 $B_{j_{k}}$ 为满足
\begin{equation*}
t_{j_{k}} = \max\{t_{j}: B_{j} \cap B_{j_{l}} = \emptyset \ (l = 1, \ldots, k-1)\}
\end{equation*}
的球.

对如此得到的 $B_{j_{1}}, B_{j_{2}}, \ldots, B_{j_{K}}$,任意 $B_{j}$ 都可以由某个 $B_{j_{k}}$ 使得
\begin{equation}
B_{j} \subset B(x_{j_{k}}, 3t_{j_{k}})\tag{4.1.7}
\label{eq:4.1.7}\end{equation}
成立.实际上,取 $B_{j_{1}}, B_{j_{2}}, \ldots, B_{j_{K}}$ 中第一个与 $B_{j}$ 相交的球为 $B_{j_{k}}$,则
\begin{equation*}
t_{j} \leq t_{j_{k}}.
\end{equation*}
否则的话,本应选取这个 $B_{j}$ 来代替 $B_{j_{k}}$.因此,由简单的几何可知 \eqref{eq:4.1.7} 成立.

由 \eqref{eq:4.1.7} 得
\begin{equation*}
E \subset \bigcup_{k=1}^{K}B(x_{j_{k}}, 3t_{j_{k}}),
\end{equation*}
从而
\begin{equation*}
|E| \leq \sum_{k=1}^{K}|B(x_{j_{k}}, 3t_{j_{k}})| = 3^{n}\sum_{k=1}^{K}|B_{j_{k}}|.
\end{equation*}
\end{proof}

\begin{proof}[定理~\ref{thm:4.1.1} (b) 的证明]
设 $f \in L^{1}(\mathbb{R}^{n})$,令
\begin{equation*}
E_{\lambda} := \{x \in \mathbb{R}^{n}: \mathcal{M}f(x) > \lambda\} \quad (\lambda > 0).
\end{equation*}
此时,设 $K \subset E_{\lambda}$ 为任意紧集,对每个 $x \in K$,由极大函数的定义选取球 $B(x, t_{x})$,使得
\begin{equation}
\lambda|B(x, t_{x})| < \int_{B(x, t_{x})}|f(y)|\dd y\tag{4.1.8}
\label{eq:4.1.8}\end{equation}
成立.

由 $K$ 的紧性,得到满足
\begin{equation*}
K \subset \bigcup_{j=1}^{N}B_{j}, \quad B_{j} := B(x_{j}, t_{x_{j}})
\end{equation*}
的有限个球构成的族.对此族应用引理~\ref{lem:4.1.3},得
\begin{equation*}
|K| \leq 3^{n}\sum_{k=1}^{K}|B_{j_{k}}|, \quad B_{j_{k}} \cap B_{j_{l}} = \emptyset \quad (k \neq l).
\end{equation*}
用 \eqref{eq:4.1.8} 估计此右端,得
\begin{equation*}
|K| \leq \frac{3^{n}}{\lambda}\sum_{k=1}^{K}\int_{B_{j_{k}}}|f(y)|\dd y \leq \frac{3^{n}}{\lambda}\int_{\mathbb{R}^{n}}|f(y)|\dd y.
\end{equation*}
最后一个不等号是由于 $B_{j_{k}}$ 互不相交.

由于 $K \subset E_{\lambda}$ 可以任意选取,
\begin{equation*}
|E_{\lambda}| \leq \frac{3^{n}}{\lambda}\|f\|_{L^{1}(\mathbb{R}^{n})}.
\end{equation*}
(b) 成立.
\end{proof}

\begin{proof}[定理~\ref{thm:4.1.1} (c) 的证明]
显然,
\begin{equation}
\|\mathcal{M}f\|_{L^{\infty}(\mathbb{R}^{n})} \leq \|f\|_{L^{\infty}(\mathbb{R}^{n})},\tag{4.1.9}
\label{eq:4.1.9}\end{equation}
故设 $f \in L^{p}(\mathbb{R}^{n})$,$1 < p < \infty$ 来证明.

对 $\lambda > 0$,令
\begin{equation*}
f_{\lambda}(x) := |f(x)|\chi_{\{y: |f(y)| > \lambda\}}(x),
\end{equation*}
则分解为
\begin{equation*}
|f(x)| = f_{\lambda}(x) + \{|f(x)| - f_{\lambda}(x)\}.
\end{equation*}

考虑 $\mathcal{M}$ 的次线性性,可知
\begin{equation*}
\{x: \mathcal{M}f(x) > 2\lambda\} \subset \{x: \mathcal{M}f_{\lambda}(x) > \lambda\}.
\end{equation*}
注意到若 $f \in L^{p}$ 则 $f_{\lambda} \in L^{1}$,利用这一点,
\begin{equation}
\begin{aligned}
|E_{\lambda}| & = |\{x: \mathcal{M}f(x) > \lambda\}| \\
& \leq |\{x: \mathcal{M}f_{\lambda/2}(x) > \lambda/2\}| \\
& \leq \frac{2\cdot 3^{n}}{\lambda}\|f_{\lambda/2}\|_{1} \\
& = \frac{2\cdot 3^{n}}{\lambda}\int_{\mathbb{R}^{n}}|f(y)|\chi_{\{|f| > \lambda/2\}}(y)\dd y.
\end{aligned}\tag{4.1.10}
\label{eq:4.1.10}\end{equation}

想起 \eqref{eq:4.1.5},用 \eqref{eq:4.1.10} 得
\begin{equation}
\begin{aligned}
\int_{\mathbb{R}^{n}}\mathcal{M}f(x)^{p}\dd x & = p\int_{0}^{\infty}\lambda^{p-1}|\{x: \mathcal{M}f(x) > \lambda\}|\dd\lambda \\
& \leq 2\cdot 3^{n}p\int_{0}^{\infty}\lambda^{p-2}\int_{\mathbb{R}^{n}}|f(y)|\chi_{\{|f| > \lambda/2\}}(y)\dd y\dd\lambda.
\end{aligned}
\end{equation}
交换积分顺序,上式右端的二重积分等于
\begin{equation*}
\int_{\mathbb{R}^{n}}|f(y)|\int_{0}^{2|f(y)|}\lambda^{p-2}\dd\lambda\dd y = \frac{2^{p-1}}{p-1}\int_{\mathbb{R}^{n}}|f(y)|^{p}\dd y.
\end{equation*}
综合以上全部结果,
\begin{equation*}
\|\mathcal{M}f\|_{L^{p}(\mathbb{R}^{n})} \leq 2\left(\frac{3^{n}p}{p-1}\right)^{1/p}\|f\|_{L^{p}(\mathbb{R}^{n})}
\end{equation*}
得证.
\end{proof}

系~\ref{cor:4.1.2} 的证明需要下面的命题,其证明在本节最后给出.

\begin{proposition}\label{prop:4.1.4}
设 $f \in L^{p}(\mathbb{R}^{n})$($1 \leq p < \infty$),令
\begin{equation*}
f_{t}(x) := \frac{1}{|B(x,t)|}\int_{B(x,t)}f(y)\dd y \quad (t > 0).
\end{equation*}
此时,
\begin{equation*}
\lim_{t \to 0}\|f - f_{t}\|_{L^{p}(\mathbb{R}^{n})} = 0
\end{equation*}
成立.
\end{proposition}

\begin{proof}[系~\ref{cor:4.1.2} 的证明]
微分是局部操作,故只需对 $f \in L^{1}(\mathbb{R}^{n})$ 证明系即可.此时,由上面的命题,
\begin{equation*}
\lim_{t \to 0}\|f - f_{t}\|_{1} = 0,
\end{equation*}
故选取子列 $t_{k} \to 0$($k \to \infty$),可以做到
\begin{equation*}
\lim_{k \to \infty}f_{t_{k}}(x) = f(x) \quad (\text{a.e. } x \in \mathbb{R}^{n}).
\end{equation*}
因此,问题归结为证明对几乎处处 $x \in \mathbb{R}^{n}$,$\lim_{t \to 0}f_{t}(x)$ 存在.

对 $g \in L^{1}(\mathbb{R}^{n})$,其振幅定义为
\begin{equation}
\Omega g(x) := \varlimsup_{t \to 0}g_{t}(x) - \varliminf_{t \to 0}g_{t}(x).\tag{4.1.11}
\label{eq:4.1.11}\end{equation}
若 $g$ 是支集紧的连续函数,则 $\Omega g(x) \equiv 0$.

容易看出,逐点有
\begin{equation*}
\Omega g(x) \leq 2\mathcal{M}g(x).
\end{equation*}
因此,$g \in L^{1}(\mathbb{R}^{n})$ 时,由定理~\ref{thm:4.1.1} (b),
\begin{equation}
|\{x: \Omega g(x) > \epsilon\}| \leq |\{x: 2\mathcal{M}g(x) > \epsilon\}| \leq \frac{2\cdot 3^{n}}{\epsilon}\|g\|_{1} \quad (\epsilon > 0).\tag{4.1.12}
\label{eq:4.1.12}\end{equation}

任意 $f \in L^{1}(\mathbb{R}^{n})$ 可以分解为 $f = g + h$,其中 $g$ 是支集紧的连续函数,$h$ 的 $L^{1}$ 范数可以自由地取得任意小.显然,
\begin{equation*}
\Omega f(x) \leq \Omega g(x) + \Omega h(x),
\end{equation*}
故由上面所确认的事实,
\begin{equation*}
\Omega f(x) \leq \Omega h(x).
\end{equation*}
由此与 \eqref{eq:4.1.12},
\begin{equation*}
|\{x: \Omega f(x) > \epsilon\}| \leq |\{x: \Omega h(x) > \epsilon\}| \leq \frac{2\cdot 3^{n}}{\epsilon}\|h\|_{1} \quad (\epsilon > 0).
\end{equation*}
由于 $\|h\|_{1}$ 可以自由设置,上式左端等于 0.因此,对几乎处处 $x \in \mathbb{R}^{n}$ 有 $\Omega f(x) = 0$,即 $\lim_{t \to 0}f_{t}(x)$ 存在,目标系成立.
\end{proof}

\begin{proof}[命题~\ref{prop:4.1.4} 的证明]
设 $f \in L^{p}(\mathbb{R}^{n})$($1 \leq p < \infty$).对任意 $\epsilon > 0$,选取支集紧的连续函数 $g$ 和满足
\begin{equation*}
\|h\|_{L^{p}(\mathbb{R}^{n})} < \frac{\epsilon}{4}
\end{equation*}
的 $h$,分解为 $f = g + h$.考虑线性性,变形为
\begin{equation*}
f(x) - f_{t}(x) = \left(g(x) - g_{t}(x)\right) + \left(h(x) - h_{t}(x)\right),
\end{equation*}
得
\begin{equation*}
\|f - f_{t}\|_{p} \leq \|g - g_{t}\|_{p} + \|h - h_{t}\|_{p}.
\end{equation*}

由于 $g$ 一致连续,对任意 $D > 0$,选取充分小的 $\delta > 0$,可以做到
\begin{equation*}
|g(x) - g(y)| \leq D \quad (|x - y| \leq \delta).
\end{equation*}
注意
\begin{equation*}
\int_{B(x,t)}\frac{\dd y}{|B(x,t)|} = 1,
\end{equation*}
则
\begin{equation}
\begin{aligned}
|g(x) - g_{t}(x)| & = \left|\int_{B(x,t)}\frac{g(x)}{|B(x,t)|}\dd y - \int_{B(x,t)}\frac{g(y)}{|B(x,t)|}\dd y\right| \\
& \leq \int_{B(x,t)}\frac{|g(x) - g(y)|}{|B(x,t)|}\dd y \\
& \leq D \quad (t < \delta),
\end{aligned}
\end{equation}
故只要 $D > 0$ 取得充分小,
\begin{equation*}
\|g - g_{t}\|_{p} \leq D|\supp(g - g_{t})|^{1/p} \leq \frac{\epsilon}{2} \quad (t < \delta).
\end{equation*}

又
\begin{equation*}
\|h - h_{t}\|_{p} \leq \|h\|_{p} + \|h_{t}\|_{p},
\end{equation*}
把 $h_{t}$ 看作卷积,应用 Young 不等式,\footnote{关于 Young 不等式将在下一节说明.}得
\begin{equation*}
\|h - h_{t}\|_{p} \leq 2\|h\|_{p} \leq \frac{\epsilon}{2}.
\end{equation*}

综合以上,
\begin{equation*}
\|f - f_{t}\|_{p} \leq \epsilon \quad (t < \delta).
\end{equation*}
\end{proof}

% ============ 第4章 §4.2 Hardy-Littlewood-Sobolev不等式(原书 pp.319-324) ============
\section{Hardy--Littlewood--Sobolev 不等式}\label{sec:4.2}
\index{Hardy--Littlewood--Sobolev 不等式}
\index{分数阶积分算子}


先复习 Young 不等式.

\noindent\textbf{Young 不等式}\quad 首先准备一个引理.

\begin{lemma}\label{lem:4.2.1}
设 $f(x)$ 为 $\mathbb{R}^{n}$ 上的可测函数.$f \in L^{p}(\mathbb{R}^{n})$($1 \leq p \leq \infty$)的充分必要条件为:设 $\frac{1}{p} + \frac{1}{p'} = 1$,存在满足
\begin{equation*}
\left|\int_{\mathbb{R}^{n}}f(x)g(x)\dd x\right| \leq M\|g\|_{L^{p'}(\mathbb{R}^{n})} \quad (g \in L^{p'}(\mathbb{R}^{n}))
\end{equation*}
的常数 $M$.特别地,此时
\begin{equation*}
\|f\|_{L^{p}(\mathbb{R}^{n})} \leq M
\end{equation*}
成立.\footnote{以后,对 $1 \leq p \leq \infty$,用 $p'$ 表示满足 $1/p + 1/p' = 1$ 的数,即 $p' = p/(p-1)$.}
\end{lemma}

\begin{proof}
只要设 $f, g$ 均为非负来证明即可.必要性由 Hölder 不等式显然,故证明充分性.

若 $p = 1$,令 $g \equiv 1 \in L^{\infty}$,则 $\|f\|_{1} \leq M$.

设 $1 < p < \infty$.令
\begin{equation*}
f_{t}(x) := f(x)\chi_{B(0,t)}(x)\chi_{[0,t]}(f(x)) \quad (t > 0),
\end{equation*}
则
\begin{equation*}
f_{t}(x)^{p-1} \in L^{p'}(\mathbb{R}^{n}) \quad (t > 0).
\end{equation*}
利用这一点,并注意 $p' = p/(p-1)$,
\begin{equation*}
\int_{\mathbb{R}^{n}}f_{t}(x)^{p}\dd x = \int_{\mathbb{R}^{n}}f_{t}(x)f_{t}(x)^{p-1}\dd x \leq M\left(\int_{\mathbb{R}^{n}}f_{t}(x)^{p}\dd x\right)^{1/p'}.
\end{equation*}
由于 $1/p + 1/p' = 1$,将两边除以
\begin{equation*}
\left(\int_{\mathbb{R}^{n}}f_{t}(x)^{p}\dd x\right)^{1/p'},
\end{equation*}
得
\begin{equation*}
\|f_{t}\|_{p} \leq M.
\end{equation*}
令 $t \to \infty$,由单调收敛定理得 $\|f\|_{p} \leq M$.

$p = \infty$ 时,令
\begin{equation*}
g_{y,t}(x) := \frac{\chi_{B(y,t)}(x)}{|B(y,t)|},
\end{equation*}
则
\begin{equation*}
\|g_{y,t}\|_{1} = 1,
\end{equation*}
故
\begin{equation*}
\int_{\mathbb{R}^{n}}f(x)g_{y,t}(x)\dd x \leq M.
\end{equation*}
令 $t \to 0$,左端由 Lebesgue 微分定理在几乎处处 $y \in \mathbb{R}^{n}$ 收敛到 $f(y)$.因此
\begin{equation*}
f(y) \leq M \quad (\mathrm{a.e.} \ y \in \mathbb{R}^{n}),
\end{equation*}
得 $\|f\|_{\infty} \leq M$.
\end{proof}

进入正题.

\begin{proposition}[Young 不等式]\label{prop:4.2.2}
设 $p, q, r$ 满足
\begin{equation}
\frac{1}{r} = 1 - \left(\frac{1}{p} - \frac{1}{q}\right), \quad 1 \leq p \leq q \leq \infty, \quad 1 \leq r \leq \infty.\tag{4.2.1}
\label{eq:4.2.1}\end{equation}
此时,令
\begin{equation*}
Tf(x) := \int_{\mathbb{R}^{n}}K(x, y)f(y)\dd y,
\end{equation*}
若假设
\begin{equation*}
\sup_{x \in \mathbb{R}^{n}}\left(\int|K(x, y)|^{r}\dd y\right)^{1/r} \leq C, \quad \sup_{y \in \mathbb{R}^{n}}\left(\int|K(x, y)|^{r}\dd x\right)^{1/r} \leq C,
\end{equation*}
则
\begin{equation}
\|Tf\|_{L^{q}(\mathbb{R}^{n})} \leq C\|f\|_{L^{p}(\mathbb{R}^{n})}\tag{4.2.2}
\label{eq:4.2.2}\end{equation}
成立.
\end{proposition}

\begin{proof}
用 M. Riesz 插值定理证明.即插值
\begin{equation*}
\|Tf\|_{L^{\infty}(\mathbb{R}^{n})} \leq C\|f\|_{L^{r'}(\mathbb{R}^{n})}
\end{equation*}
与
\begin{equation*}
\|Tf\|_{L^{r}(\mathbb{R}^{n})} \leq C\|f\|_{L^{1}(\mathbb{R}^{n})},
\end{equation*}
得到结论.

前一不等式由 Hölder 不等式可知成立.

后一不等式:设 $f \in L^{1}(\mathbb{R}^{n})$,$g \in L^{r'}(\mathbb{R}^{n})$,
\begin{equation}
\begin{aligned}
\left|\int Tf(x)g(x)\dd x\right| & = \left|\int\int K(x, y)f(y)g(x)\dd y\dd x\right| \\
& = \left|\int f(y)\int K(x, y)g(x)\dd x\dd y\right| \leq C\|f\|_{1}\times\|g\|_{r'},
\end{aligned}
\end{equation}
可以算出,由引理~\ref{lem:4.2.1} 可知成立.
\end{proof}

\noindent\textbf{分数阶积分算子}\quad 用 Hardy--Littlewood 极大函数的 $L^{p}(\mathbb{R}^{n})$ 估计来证明分数阶积分算子的 $L^{p}(\mathbb{R}^{n}) \to L^{q}(\mathbb{R}^{n})$ 有界性.

设 $r > 1$,分数阶积分算子 $I_{r}$ 定义为
\begin{equation*}
I_{r}f(x) := \int_{\mathbb{R}^{n}}\frac{f(y)}{|x-y|^{n/r}}\dd y \quad (r > 1).
\end{equation*}
关于分数阶积分算子,下面的定理成立.

\begin{theorem}[Hardy--Littlewood--Sobolev 不等式]\label{thm:4.2.3}
设 $p, q, r$ 满足
\begin{equation}
\frac{1}{r} = 1 - \left(\frac{1}{p} - \frac{1}{q}\right), \quad 1 < p < q < \infty, \quad 1 < r < \infty.\tag{4.2.3}
\label{eq:4.2.3}\end{equation}
此时,若 $f \in L^{p}(\mathbb{R}^{n})$,则
\begin{equation}
\|I_{r}f\|_{L^{q}(\mathbb{R}^{n})} \leq C_{n,p,q}\|f\|_{L^{p}(\mathbb{R}^{n})}\tag{4.2.4}
\label{eq:4.2.4}\end{equation}
成立.
\end{theorem}

Hardy--Littlewood--Sobolev 不等式与 Young 不等式的区别在于,$I_{r}$ 的核 $|y|^{-n/r}$ 不属于 $L^{r}(\mathbb{R}^{n})$.

\begin{proof}[定理~\ref{thm:4.2.3} 的证明]
固定 $f \in L^{p}(\mathbb{R}^{n})$ 与 $x \in \mathbb{R}^{n}$,估计 $I_{r}f(x)$.

定义
\begin{equation*}
f_{t}(x) := \frac{1}{t^{n}}\int_{B(x,t)}|f(y)|\dd y \quad (t > 0).
\end{equation*}
此时,对 $y \in \mathbb{R}^{n}$($y \neq x$),
\begin{equation*}
\int_{0}^{\infty}\frac{\chi_{B(x,t)}(y)}{t^{n}}t^{n(1-1/r)-1}\dd t = \int_{|x-y|}^{\infty}t^{(-n/r)-1}\dd t = \frac{r}{n}|x-y|^{-n/r},
\end{equation*}
故考虑交换积分顺序,
\begin{equation}
\begin{aligned}
|I_{r}f(x)| & \leq \int_{\mathbb{R}^{n}}|x-y|^{-n/r}|f(y)|\dd y \\
& = \frac{n}{r}\int_{\mathbb{R}^{n}}\int_{0}^{\infty}\frac{\chi_{B(x,t)}(y)}{t^{n}}t^{n(1-1/r)-1}\dd t\,|f(y)|\dd y \\
& = \frac{n}{r}\int_{0}^{\infty}f_{t}(x)t^{n(1-1/r)-1}\dd t.
\end{aligned}
\end{equation}

设 $\epsilon > 0$,在完成必要的计算之后再确定其取值,将上式右端的积分分为
\begin{equation*}
\mathrm{I} := \int_{0}^{\epsilon}, \quad \mathrm{II} := \int_{\epsilon}^{\infty}
\end{equation*}
两部分分别估计.

注意 $f_{t}(x) \leq C\mathcal{M}f(x)$,
\begin{equation*}
\mathrm{I} = \int_{0}^{\epsilon}f_{t}(x)t^{n(1-1/r)-1}\dd t \leq C\mathcal{M}f(x)\epsilon^{n(1-1/r)}.
\end{equation*}

由 Hölder 不等式与 $f \in L^{p}(\mathbb{R}^{n})$,
\begin{equation}
\begin{aligned}
f_{t}(x) & = \frac{1}{t^{n}}\int_{B(x,t)}|f(y)|\dd y \\
& \leq \frac{1}{t^{n}}|B(x,t)|^{1-1/p}\cdot\left(\int_{B(x,t)}|f(y)|^{p}\dd y\right)^{1/p} \leq Ct^{-n/p}\|f\|_{p}.
\end{aligned}
\end{equation}

由 \eqref{eq:4.2.3} 推出的
\begin{equation*}
1 - \frac{1}{r} - \frac{1}{p} = -\frac{1}{q}
\end{equation*}
与上式,
\begin{equation}
\begin{aligned}
\mathrm{II} &= \int_{\epsilon}^{\infty}f_{t}(x)t^{n(1-1/r)-1}\dd t \\
&\leq C\|f\|_{p}\int_{\epsilon}^{\infty}t^{n(1-1/r-1/p)-1}\dd t \\
&= C\|f\|_{p}\epsilon^{-n/q}.
\end{aligned}
\end{equation}

因此
\begin{equation}
|I_{r}f(x)| \leq C\left(\mathcal{M}f(x)\epsilon^{n(1-1/r)} + \|f\|_{p}\epsilon^{-n/q}\right).\tag{4.2.5}
\label{eq:4.2.5}\end{equation}
此右端在满足
\begin{equation*}
\mathcal{M}f(x)\epsilon^{n(1-1/r)} = \|f\|_{p}\epsilon^{-n/q}
\end{equation*}
的 $\epsilon$ 处几乎取到最小.即,注意由 \eqref{eq:4.2.3} 推出的
\begin{equation*}
1 - \frac{1}{r} + \frac{1}{q} = \frac{1}{p}, \quad 1 - \frac{p}{q} = p - \frac{p}{r},
\end{equation*}
取
\begin{equation*}
\epsilon = \left(\frac{\|f\|_{p}}{\mathcal{M}f(x)}\right)^{p/n},
\end{equation*}
得
\begin{equation}
|I_{r}f(x)| \leq C\|f\|_{p}^{1-p/q}\mathcal{M}f(x)^{p/q}.\tag{4.2.6}
\label{eq:4.2.6}\end{equation}

由极大函数的 $L^{p}(\mathbb{R}^{n})$ 有界性,
\begin{equation*}
\int_{\mathbb{R}^{n}}(\mathcal{M}f(x))^{p}\dd x \leq C\|f\|_{p}^{p},
\end{equation*}
故
\begin{equation*}
\int_{\mathbb{R}^{n}}|I_{r}f(x)|^{q}\dd x \leq C\|f\|_{p}^{q-p}\int_{\mathbb{R}^{n}}(\mathcal{M}f(x))^{p}\dd x \leq C\|f\|_{p}^{q}.
\end{equation*}
两边开 $q$ 次方,
\begin{equation*}
\|I_{r}f\|_{q} \leq C\|f\|_{p}.
\end{equation*}
\end{proof}

% ============ 第4章 §4.3 Fourier变换(原书 pp.324-330) ============
\section{Fourier 变换}\label{sec:4.3}
\index{Fourier 变换}
\index{Fourier 反演公式}
\index{Fourier 变换!(缓增广义函数的)}
\index{缓增广义函数}
\index{缓增广义函数族}
\index{速降函数}
\index{速降函数族}
\index{Plancherel 定理}
\index{Parseval 等式}
\index{Riemann--Lebesgue 定理}
\index{卷积}


本节中,定义 Fourier 变换,并确认其基本性质.对可积函数 $f \in L^{1}(\mathbb{R}^{n})$,其 Fourier 变换定义为
\begin{equation*}
\hat{f}(\xi) := \int_{\mathbb{R}^{n}}e^{-2\pi i x\cdot\xi}f(x)\dd x \quad (\xi \in \mathbb{R}^{n}).
\end{equation*}
这里,$x = (x_{1}, \ldots, x_{n}) \in \mathbb{R}^{n}$,$\xi = (\xi_{1}, \ldots, \xi_{n}) \in \mathbb{R}^{n}$ 的内积表为
\begin{equation*}
x\cdot\xi := x_{1}\xi_{1} + \ldots + x_{n}\xi_{n}.
\end{equation*}

\begin{proposition}\label{prop:4.3.1}
设 $f \in L^{1}(\mathbb{R}^{n})$.

(i)
\begin{equation}
\|\hat{f}\|_{\infty} \leq \|f\|_{1}\tag{4.3.1}
\label{eq:4.3.1}\end{equation}
成立.

(ii) $\hat{f}(\xi)$ 是一致连续的函数.

(iii) $\lim_{|\xi| \to \infty}\hat{f}(\xi) = 0$ 成立.称此为 Riemann--Lebesgue 定理.
\end{proposition}

\begin{proof}
(i) 由定义显然.(ii) 来证明.固定 $\epsilon > 0$.首先选取足够大的 $R > 0$,使得
\begin{equation*}
\int_{|x| > R}|f(x)|\dd x \leq \frac{\epsilon}{4},
\end{equation*}
其次选取足够小的 $\delta > 0$,使得当 $|h| \leq \delta$($h \in \mathbb{R}^{n}$)时
\begin{equation*}
\left|e^{-2\pi i x\cdot h} - 1\right| \leq \frac{\epsilon}{2\|f\|_{1}} \quad (|x| \leq R).
\end{equation*}
利用这些,对 $\xi \in \mathbb{R}^{n}$,$|h| \leq \delta$,
\begin{equation}
\begin{aligned}
|\hat{f}(\xi + h) - \hat{f}(\xi)| & = \left|\int_{\mathbb{R}^{n}}\left(e^{-2\pi i x\cdot h} - 1\right)e^{-2\pi i x\cdot\xi}f(x)\dd x\right| \\
& \leq \int_{|x| \leq R}\left|e^{-2\pi i x\cdot h} - 1\right||f(x)|\dd x + 2\int_{|x| \geq R}|f(x)|\dd x \leq \epsilon.
\end{aligned}
\end{equation}
故 $\hat{f}(\xi)$ 一致连续.(iii) 来证明.

令
\begin{equation*}
I(\alpha) := \int_{a}^{b}e^{-2\pi i\theta\alpha}\dd\theta,
\end{equation*}
可以算出
\begin{equation*}
\lim_{|\alpha| \to \infty}I(\alpha) = 0.
\end{equation*}
由此可知,对任意立方体\footnote{所谓立方体,指各边与各坐标轴平行的立方体.}$Q \subset \mathbb{R}^{n}$ 的特征函数,(iii) 成立.用立方体特征函数的线性组合逼近 $f$,即得所证.
\end{proof}

\begin{definition}[速降函数]\label{def:4.3.2}
无限次连续可微函数 $\phi \in C^{\infty}(\mathbb{R}^{n})$ 在满足
\begin{equation}
\sup_{x \in \mathbb{R}^{n}}\left|x^{\gamma}\partial^{\alpha}\phi(x)\right| < \infty\tag{4.3.2}
\label{eq:4.3.2}\end{equation}
(对任意多重指标 $\alpha, \gamma \in \mathbb{N}_{0}^{n}$)时,称为速降函数.速降函数全体构成的线性空间记为 $\mathcal{S}(\mathbb{R}^{n})$,称为速降函数族.这里,对
\begin{equation*}
\alpha = (\alpha_{1}, \ldots, \alpha_{n}), \quad \gamma = (\gamma_{1}, \ldots, \gamma_{n})
\end{equation*}
表为
\begin{equation*}
x^{\gamma} := (x_{1})^{\gamma_{1}}\ldots(x_{n})^{\gamma_{n}}, \quad \partial^{\alpha} := (\partial/\partial x_{1})^{\alpha_{1}}\ldots(\partial/\partial x_{n})^{\alpha_{n}}.
\end{equation*}
\end{definition}

速降函数是 $L^{p}(\mathbb{R}^{n})$($1 \leq p \leq \infty$)中的函数.

速降函数族 $\mathcal{S}(\mathbb{R}^{n})$ 由 \eqref{eq:4.3.2} 定义的半范数族成为 Fréchet 空间.对 $k \in \mathbb{N}$,$\phi \in \mathcal{S}(\mathbb{R}^{n})$,令
\begin{equation*}
p_{k}(\phi) := \sup\left\{\left|x^{\gamma}\partial^{\alpha}\phi(x)\right|: |\alpha|, |\gamma| \leq k, x \in \mathbb{R}^{n}\right\},
\end{equation*}
定义距离
\begin{equation*}
d(\phi, \psi) := \sum_{k=1}^{\infty}2^{-k}\min\{p_{k}(\phi - \psi), 1\},
\end{equation*}
则 $\mathcal{S}(\mathbb{R}^{n})$ 可以看作度量空间.

令 $D_{j} := \frac{1}{2\pi i}(\partial/\partial x_{j})$.对 $\phi \in \mathcal{S}(\mathbb{R}^{n})$,由分部积分,
\begin{equation*}
\xi_{j}\hat{\phi}(\xi) = \int\left(-D_{j}e^{-2\pi i x\cdot\xi}\right)\phi(x)\dd x = \int e^{-2\pi i x\cdot\xi}D_{j}\phi(x)\dd x,
\end{equation*}
由积分与微分的交换,
\begin{equation*}
D_{j}\hat{\phi}(\xi) = \int e^{-2\pi i x\cdot\xi}(-x_{j})\phi(x)\dd x
\end{equation*}
成立.利用这些,
\begin{equation*}
\xi^{\alpha}D^{\gamma}\hat{\phi}(\xi) = \int e^{-2\pi i x\cdot\xi}D^{\alpha}\left((-x)^{\gamma}\phi(x)\right)\dd x.
\end{equation*}
进而
\begin{equation}
\begin{aligned}
|\xi^{\alpha}D^{\gamma}\hat{\phi}(\xi)| & \leq \int\left|D^{\alpha}\left((-x)^{\gamma}\phi(x)\right)\right|\dd x \\
& = \int\frac{(1 + |x|)^{n+1}}{(1 + |x|)^{n+1}}\left|D^{\alpha}\left((-x)^{\gamma}\phi(x)\right)\right|\dd x \\
& \leq \sup_{y}\left\{(1 + |y|)^{n+1}\left|D^{\alpha}\left((-y)^{\gamma}\phi(y)\right)\right|\right\}\int\frac{1}{(1 + |x|)^{n+1}}\dd x \\
& = C\sup_{y}\left\{(1 + |y|)^{n+1}\left|D^{\alpha}\left((-y)^{\gamma}\phi(y)\right)\right|\right\}.
\end{aligned}
\end{equation}
对左端取上限,可知 Fourier 变换在 $\mathcal{S} \to \mathcal{S}$ 中连续.进而,由下面的定理可知 Fourier 变换是 $\mathcal{S} \to \mathcal{S}$ 的同胚.

\begin{theorem}\label{thm:4.3.3}
设 $\phi \in \mathcal{S}(\mathbb{R}^{n})$.对所有 $x \in \mathbb{R}^{n}$,Fourier 反演公式
\begin{equation}
\int_{\mathbb{R}^{n}}e^{2\pi i x\cdot\xi}\hat{\phi}(\xi)\dd\xi = \phi(x)\tag{4.3.3}
\label{eq:4.3.3}\end{equation}
成立.
\end{theorem}

为证明准备如下两个引理.

\begin{lemma}\label{lem:4.3.4}
若 $f, g \in L^{1}$,则
\begin{equation*}
\int_{\mathbb{R}^{n}}\hat{f}g\dd x = \int_{\mathbb{R}^{n}}f\hat{g}\dd x
\end{equation*}
成立.
\end{lemma}

\begin{lemma}\label{lem:4.3.5}
设 $\epsilon > 0$,则
\begin{equation*}
\int_{\mathbb{R}^{n}}e^{-2\pi i x\cdot\xi}e^{-\pi\epsilon|x|^{2}}\dd x = \epsilon^{-n/2}e^{-\pi|\xi|^{2}/\epsilon}
\end{equation*}
成立.
\end{lemma}

引理~\ref{lem:4.3.4} 由 Fubini 定理显然.引理~\ref{lem:4.3.5} 由变量替换与下面的计算可知:
\begin{equation*}
\begin{aligned}
\int_{-\infty}^{\infty}e^{-2\pi i t\tau}e^{-\pi t^{2}}\dd t & = e^{-\pi\tau^{2}}\int_{-\infty}^{\infty}e^{-\pi(t + i\tau)^{2}}\dd t \\
& = e^{-\pi\tau^{2}}\int_{-\infty}^{\infty}e^{-\pi t^{2}}\dd t = e^{-\pi\tau^{2}}.
\end{aligned}
\end{equation*}

\begin{proof}[定理~\ref{thm:4.3.3} 的证明]
首先注意
\begin{equation*}
e^{2\pi i x\cdot\xi}\hat{\phi}(\xi) = \int e^{-2\pi i (y - x)\cdot\xi}\phi(y)\dd y = \int e^{-2\pi i y\cdot\xi}\phi(y + x)\dd y.
\end{equation*}
由这个等式与上面的两个引理,
\begin{equation*}
\int e^{2\pi i x\cdot\xi}\hat{\phi}(\xi)e^{-\pi\epsilon|\xi|^{2}}\dd\xi = \epsilon^{-n/2}\int\phi(y + x)e^{-\pi|y|^{2}/\epsilon}\dd y.
\end{equation*}
令 $\epsilon \to 0$,上式左端由 Lebesgue 收敛定理收敛到 \eqref{eq:4.3.3} 的左端;上式右端,由于
\begin{equation*}
\epsilon^{-n/2}\int e^{-\pi|y|^{2}/\epsilon}\dd y = 1,
\end{equation*}
可知 $\epsilon^{-n/2}e^{-\pi|y|^{2}/\epsilon}\dd y$ 是恒等逼近,故收敛到 $\phi(x)$.
\end{proof}

$f, g \in L^{1}$ 的卷积定义为
\begin{equation*}
(f * g)(x) := \int f(x - y)g(y)\dd y.
\end{equation*}

\begin{theorem}\label{thm:4.3.6}
若 $\phi, \psi \in \mathcal{S}$,则
\begin{equation}
\int\hat{\phi}\,\overline{\widehat{\psi}}\dd x = \int\phi\,\overline{\psi}\dd x,\tag{4.3.4}
\label{eq:4.3.4}\end{equation}
以及
\begin{equation}
(\phi * \psi)^{\wedge}(\xi) = \hat{\phi}(\xi)\hat{\psi}(\xi)\tag{4.3.5}
\label{eq:4.3.5}\end{equation}
成立.
\end{theorem}

\begin{proof}
令 $\eta = \overline{\hat{\psi}}$,由 Fourier 反演公式可知 $\hat{\eta} = \overline{\psi}$.应用引理~\ref{lem:4.3.4},即得 \eqref{eq:4.3.4}.其次,
\begin{equation}
\begin{aligned}
(\phi * \psi)^{\wedge}(\xi) & = \int e^{-2\pi i x\cdot\xi}\int\phi(x - y)\psi(y)\dd y\dd x \\
& = \int\int e^{-2\pi i (x - y)\cdot\xi}\phi(x - y)e^{-2\pi i y\cdot\xi}\psi(y)\dd y\dd x = \hat{\phi}(\xi)\hat{\psi}(\xi).
\end{aligned}
\end{equation}
得 \eqref{eq:4.3.5}.
\end{proof}

为对更一般的函数考虑 Fourier 变换,定义缓增广义函数.

\begin{definition}[缓增广义函数]\label{def:4.3.7}
速降函数族 $\mathcal{S}(\mathbb{R}^{n})$ 的对偶空间记为 $\mathcal{S}'(\mathbb{R}^{n})$,称为缓增广义函数族.$\mathcal{S}'(\mathbb{R}^{n})$ 的元称为缓增广义函数.
\end{definition}

\begin{definition}\label{def:4.3.8}
对 $u \in \mathcal{S}'(\mathbb{R}^{n})$,其 Fourier 变换定义为
\begin{equation*}
\hat{u}(\phi) = u(\hat{\phi}).
\end{equation*}
\end{definition}

由引理~\ref{lem:4.3.4},若 $u \in L^{1}$,则 $u$ 作为广义函数的 Fourier 变换与作为函数的 Fourier 变换一致.由 Fourier 反演公式,Fourier 变换是 $\mathcal{S}' \to \mathcal{S}'$ 的同胚.特别地,若 $u \in L^{1}$ 且 $\hat{u} \in L^{1}$,则 Fourier 反演公式在 a.e. $x$ 成立.

\begin{theorem}\label{thm:4.3.9}
若 $f \in L^{2}$,则 $\hat{f} \in L^{2}$ 且
\begin{equation}
\|\hat{f}\|_{2} = \|f\|_{2}\tag{4.3.6}
\label{eq:4.3.6}\end{equation}
成立.称此为 Plancherel 定理.进而,若 $f, g \in L^{2}$,则
\begin{equation}
\int\widehat{f}\,\overline{\widehat{g}}\dd x = \int f\,\overline{g}\dd x\tag{4.3.7}
\label{eq:4.3.7}\end{equation}
成立.称此为 Parseval 等式.
\end{theorem}

\begin{proof}
选取 $f_{j}, g_{j} \in \mathcal{S}$ 使得
\begin{equation*}
\|f_{j} - f\|_{2}, \quad \|g_{j} - g\|_{2} \to 0 \quad (j \to \infty).
\end{equation*}
此时,在 $\mathcal{S}'$ 的拓扑中 $f_{j} \to f$,故由 Fourier 变换的连续性,在 $\mathcal{S}'$ 的拓扑中 $\hat{f}_{j} \to \hat{f}$.由 \eqref{eq:4.3.4},
\begin{equation*}
\|\hat{f}_{j} - \hat{f}_{k}\|_{2} = \|f_{j} - f_{k}\|_{2} \to 0 \quad (j, k \to \infty),
\end{equation*}
故存在 $h \in L^{2}$ 使得在 $L^{2}$ 的拓扑中 $\hat{f}_{j} \to h$.这意味着在 $\mathcal{S}'$ 的拓扑中 $\hat{f}_{j} \to h$,从而 $\hat{f} = h \in L^{2}$.

由
\begin{equation*}
\|\hat{f}_{j}\|_{2} = \|f_{j}\|_{2}
\end{equation*}
令 $j \to \infty$,即得 \eqref{eq:4.3.6}.

\eqref{eq:4.3.7} 由 \eqref{eq:4.3.4} 得
\begin{equation*}
\int\hat{f}_{j}\,\overline{\widehat{g}_{j}}\dd x = \int f_{j}\,\overline{g_{j}}\dd x,
\end{equation*}
且已证明 $\hat{f}, \hat{g} \in L^{2}$,故显然.
\end{proof}

由 M. Riesz 插值定理插值 \eqref{eq:4.3.1} 与 \eqref{eq:4.3.6},得到下面的系.

\begin{corollary}\label{cor:4.3.10}
设 $1 \leq p \leq 2$,$p' = p/(p-1)$,则
\begin{equation}
\|\hat{f}\|_{L^{p'}(\mathbb{R}^{n})} \leq \|f\|_{L^{p}(\mathbb{R}^{n})}\tag{4.3.8}
\label{eq:4.3.8}\end{equation}
成立.
\end{corollary}

% ============ 第4章 §4.4 驻相法(原书 pp.330-336) ============
\section{驻相法}\label{sec:4.4}
\index{驻相法}
\index{Van der Corput 引理}
\index{Morse 引理}


本节所讨论的驻相法,是数量地把握 Fourier 变换的振荡衰减的一种手段.该方法重要且应用范围广泛,但由于是强烈依赖于函数光滑性信息的方法,在关于函数的定量估计问题中,还需要进一步的手段.

设 $a \in C_{0}^{\infty}(\mathbb{R}^{n})$.对实值函数 $\phi \in C^{\infty}(\mathbb{R}^{n})$,令
\begin{equation*}
I(\lambda) := \left|\int_{\mathbb{R}^{n}}e^{i\lambda\phi(x)}a(x)\dd x\right| \quad (\lambda > 1).
\end{equation*}
估计 $I(\lambda)$ 是本节的目标.

若 $\phi(x) \not\equiv C$(常数),则当 $\lambda \to \infty$ 时上式右端的被积函数剧烈振荡,其积分值变小.问题在于对其大小的估计.

\begin{theorem}\label{thm:4.4.1}
若 $\phi$ 满足
\begin{equation*}
|\nabla\phi(x)| \geq c > 0 \quad (x \in \supp a),
\end{equation*}
则对任意 $\lambda > 1$,$N = 1, 2, \ldots$,成立如下估计:
\begin{equation}
I(\lambda) = \left|\int_{\mathbb{R}^{n}}e^{i\lambda\phi(x)}a(x)\dd x\right| \leq C_{N}\lambda^{-N}.\tag{4.4.1}
\label{eq:4.4.1}\end{equation}
\end{theorem}

\begin{proof}
对任意 $x \in \supp a$,选取单位向量 $\nu_{x} \in S^{n-1}$ 与 $r_{x} > 0$,使得
\begin{equation*}
\nu_{x}\cdot\nabla\phi(y) \geq c/2 \quad (y \in B(x, r_{x})).
\end{equation*}
由于 $\supp a$ 紧,由此可由有限个 $\alpha_{j} \in C_{0}^{\infty}$,$\nu_{j} \in S^{n-1}$ 做到
\begin{equation*}
\sum_{j=1}^{J}\alpha_{j}(y) = 1 \quad (y \in \supp a), \quad \nu_{j}\cdot\nabla\phi(y) \geq c/2 \quad (y \in \supp\alpha_{j}).
\end{equation*}
为证明,令 $a_{j}(x) = \alpha_{j}(x)a(x)$,只要证明
\begin{equation*}
\left|\int e^{i\lambda\phi(x)}a_{j}(x)\dd x\right| \leq C_{N}\lambda^{-N}
\end{equation*}
即可.必要时变换坐标,设 $\nu_{j} = (1, 0, \ldots, 0)$,可假设
\begin{equation*}
|(\partial/\partial x_{1})\phi(x)| \geq c/2 \quad (x \in \supp\alpha_{j}).
\end{equation*}

由于
\begin{equation*}
(\partial/\partial x_{1})e^{i\lambda\phi(x)} = e^{i\lambda\phi(x)}i\lambda(\partial/\partial x_{1})\phi(x),
\end{equation*}
有
\begin{equation}
\begin{aligned}
\int e^{i\lambda\phi(x)}a_{j}(x)\dd x & = \int(\partial/\partial x_{1})e^{i\lambda\phi(x)}\cdot\frac{a_{j}(x)}{i\lambda(\partial/\partial x_{1})\phi(x)}\dd x \\
& = -\frac{1}{i\lambda}\int e^{i\lambda\phi(x)}(\partial/\partial x_{1})\frac{a_{j}(x)}{(\partial/\partial x_{1})\phi(x)}\dd x.
\end{aligned}
\end{equation}
这里,令
\begin{equation*}
L^{*}(x, D) := (\partial/\partial x_{1})\frac{1}{(\partial/\partial x_{1})\phi(x)},
\end{equation*}
则进一步
\begin{equation*}
\int e^{i\lambda\phi(x)}L^{*}(x, D)a_{j}(x)\dd x = -\frac{1}{i\lambda}\int e^{i\lambda\phi(x)}L^{*}(x, D)^{2}a_{j}(x)\dd x.
\end{equation*}
这样依次重复 $N$ 次,
\begin{equation*}
\int e^{i\lambda\phi(x)}a_{j}(x)\dd x = \left(-\frac{1}{i\lambda}\right)^{N}\int e^{i\lambda\phi(x)}L^{*}(x, D)^{N}a_{j}(x)\dd x.
\end{equation*}
对两边取绝对值,到达目标.
\end{proof}

即使在某个 $x \in \supp a$ 处 $\nabla\phi(x) = 0$,若还有关于 $\phi$ 的导数的信息,仍可估计 $I(\lambda)$.称此为驻相法.以下,将
\begin{equation*}
\left((\partial^{2}/\partial x_{j}\partial x_{k})\phi(x)\right) := \left(\begin{array}{ccc}
(\partial^{2}/\partial x_{1}\partial x_{1})\phi(x) & \ldots & (\partial^{2}/\partial x_{1}\partial x_{n})\phi(x) \\
\vdots & \ddots & \vdots \\
(\partial^{2}/\partial x_{n}\partial x_{1})\phi(x) & \ldots & (\partial^{2}/\partial x_{n}\partial x_{n})\phi(x)
\end{array}\right)
\end{equation*}
表示.

\begin{theorem}\label{thm:4.4.2}
若 $\phi$ 满足
\begin{equation*}
\phi(0) = 0, \quad \nabla\phi(0) = 0, \quad \nabla\phi(x) \neq 0 \quad (x \in \supp a \setminus \{0\}),
\end{equation*}
并且满足
\begin{equation*}
\det\left((\partial^{2}/\partial x_{j}\partial x_{k})\phi(x)\right) \neq 0 \quad (x = 0),
\end{equation*}
则成立如下估计:
\begin{equation}
I(\lambda) = \left|\int_{\mathbb{R}^{n}}e^{i\lambda\phi(x)}a(x)\dd x\right| \leq C\lambda^{-n/2}.\tag{4.4.2}
\label{eq:4.4.2}\end{equation}

以下贯穿本节证明此定理.
\end{theorem}

\noindent\textbf{先从一维情形开始}\quad 以一维情形更一般的形式来证明.

\begin{lemma}[Van der Corput]\label{lem:4.4.3}
若 $\phi$ 满足
\begin{equation*}
\phi(0) = 0, \quad (\dd/\dd x)\phi(0) = 0, \quad (\dd/\dd x)\phi(x) \neq 0 \quad (x \in \supp a \setminus \{0\}),
\end{equation*}
并且满足
\begin{equation*}
(\dd^{2}/\dd x^{2})\phi(0) \neq 0,
\end{equation*}
则成立如下估计:
\begin{equation}
I_{k}(\lambda) := \left|\int_{0}^{\infty}e^{i\lambda\phi(x)}x^{k}a(x)\dd x\right| \leq C\lambda^{-1/2-k/2} \quad (k = 0, 1, \ldots).\tag{4.4.3}
\label{eq:4.4.3}\end{equation}
\end{lemma}

\begin{proof}
由条件,$\phi$ 表为
\begin{equation*}
\phi(x) = x^{2}\phi_{1}(x), \quad \phi_{1}(0) \neq 0.
\end{equation*}
选取 $\rho \in C^{\infty}$ 满足
\begin{equation*}
\rho(x) = \begin{cases} 1 & (x < 1) \\ 0 & (x > 2) \end{cases}
\end{equation*}
利用它,对 $\delta > 0$,
\begin{equation}
\begin{aligned}
I_{k}(\lambda) & = \left|\int_{0}^{\infty}e^{i\lambda\phi(x)}x^{k}a(x)\dd x\right| \\
& \leq \left|\int_{0}^{\infty}e^{i\lambda\phi(x)}x^{k}a(x)\rho(x/\delta)\dd x\right| + \left|\int_{0}^{\infty}e^{i\lambda\phi(x)}x^{k}a(x)(1 - \rho(x/\delta))\dd x\right| \\
& =: \mathrm{I} + \mathrm{II}.
\end{aligned}
\end{equation}

对 $\mathrm{I}$ 取被积函数的绝对值,可估计为
\begin{equation*}
\mathrm{I} \leq C\int_{0}^{2\delta}x^{k}\dd x = C'\delta^{1+k}.
\end{equation*}

对 $\mathrm{II}$,与定理~\ref{thm:4.4.1} 的证明同样,令 $L^{*}(x, D) := (\dd/\dd x)\frac{1}{\phi'(x)}$,
\begin{equation*}
\mathrm{II} = \frac{1}{\lambda^{N}}\left|\int_{0}^{\infty}e^{i\lambda\phi(x)}L^{*}(x, D)^{N}\left(x^{k}a(x)(1 - \rho(x/\delta))\right)\dd x\right|
\end{equation*}
可以算出.由 $|\phi'(x)| \geq Cx$ 以及乘积的微分公式,上式右端被积函数的绝对值可由 $C(x^{k-2N} + x^{k-N}\delta^{-N})$ 估计.因此,若 $N \geq k + 2$,
\begin{equation*}
\mathrm{II} \leq C\lambda^{-N}\delta^{1+k-2N}.
\end{equation*}

合并两个估计,
\begin{equation*}
I_{k}(\lambda) \leq C\left(\delta^{1+k} + \lambda^{-N}\delta^{1+k-2N}\right).
\end{equation*}
取使右端两项相等的 $\delta = \lambda^{-1/2}$,
\begin{equation*}
I_{k}(\lambda) \leq C\lambda^{-1/2-k/2}.
\end{equation*}
\end{proof}

\noindent\textbf{高维情形}\quad 首先令
\begin{equation*}
Q(x) := \frac{1}{2}\left(x_{1}^{2} + \ldots + x_{j}^{2} - x_{j+1}^{2} - \ldots - x_{n}^{2}\right),
\end{equation*}
试探性地考虑 $\phi(x) = Q(x)$ 的情形.$Q(x)$ 显然满足定理的条件.于是用归纳法估计
\begin{equation*}
\left|\int_{\mathbb{R}^{n}}e^{i\lambda Q(x)}a(x)\dd x\right|.
\end{equation*}

令
\begin{equation*}
x = (x_{1}, x'), \quad Q'(x') = Q(x) - x_{1}^{2}/2,
\end{equation*}
改写为
\begin{equation*}
\int_{\mathbb{R}^{n}}e^{i\lambda Q(x)}a(x)\dd x = \int_{-\infty}^{\infty}e^{i\lambda x_{1}^{2}/2}\left\{\int_{\mathbb{R}^{n-1}}e^{i\lambda Q'(x')}a(x)\dd x'\right\}\dd x_{1}.
\end{equation*}
此右端 $\{\ \}$ 的部分,关于 $x_{1}$ 的导数,其大小总可以由 $C\lambda^{-(n-1)/2}$ 从上估计(使用了归纳法假设).因此,想起一维情形的证明,可知
\begin{equation*}
\left|\int_{\mathbb{R}^{n}}e^{i\lambda Q(x)}a(x)\dd x\right| \leq C\lambda^{-n/2}
\end{equation*}
成立.

一般情形通过适当的坐标变换归结为 $\phi(x) = Q(x)$ 的形式.保证这一点的就是下面的 Morse 引理.

\begin{lemma}[Morse]\label{lem:4.4.4}
若 $\phi$ 满足
\begin{equation*}
\phi(0) = 0, \quad \nabla\phi(0) = 0, \quad \det\left((\partial^{2}/\partial x_{j}\partial x_{k})\phi(x)\right) \neq 0 \quad (x = 0),
\end{equation*}
则在 $x = 0$ 的某个邻域内存在光滑坐标变换 $x \to \tilde{x}$,使得
\begin{equation*}
\phi(x) = \frac{1}{2}\left(\tilde{x}_{1}^{2} + \ldots + \tilde{x}_{j}^{2} - \tilde{x}_{j+1}^{2} - \tilde{x}_{n}^{2}\right).
\end{equation*}
\end{lemma}

\begin{proof}
必要时用线性变换变换坐标,在 $x = 0$ 处可假设
\begin{equation*}
\left((\partial^{2}/\partial x_{j}\partial x_{k})\phi(x)\right) = \left(\begin{array}{cccccc}
1 & & & & & 0 \\
& \ddots & & & & \\
& & 1 & & & \\
& & & -1 & & \\
& & & & \ddots & \\
0 & & & & & -1
\end{array}\right).
\end{equation*}
此时,在 $x = 0$ 处
\begin{equation*}
(\partial/\partial x_{1})\phi(x) = 0, \quad (\partial^{2}/\partial x_{1}^{2})\phi(x) = 1,
\end{equation*}
故由隐函数定理,在 $x = 0$ 的某个邻域内存在光滑函数 $X_{1}(x')$ 使得
\begin{equation}
(\partial/\partial x_{1})\phi(X_{1}(x'), x') = 0\tag{4.4.4}
\label{eq:4.4.4}\end{equation}
成立.现在令 $y_{1} = x_{1} - X_{1}(x')$,变换
\begin{equation*}
x = (x_{1}, x') \to (y_{1}, x'),
\end{equation*}
则这是光滑坐标变换.将 $\phi(x) = \phi(y_{1} + X_{1}(x'), x')$ 的右端在 $y_{1} = 0$ 处 Taylor 展开,考虑 \eqref{eq:4.4.4},
\begin{equation}
\begin{aligned}
\phi(x) & = \phi(X_{1}(x'), x') + (\partial/\partial x_{1})\phi(X_{1}(x'), x')y_{1} + c(y_{1}, x')y_{1}^{2}/2 \\
& = \phi(X_{1}(x'), x') + c(y_{1}, x')y_{1}^{2}/2 \quad (c \in C^{\infty},\ c(0) = 1).
\end{aligned}
\end{equation}
令 $\tilde{x}_{1} = y_{1}\sqrt{c(y_{1}, x')}$,变换
\begin{equation*}
(y_{1}, x') \to (\tilde{x}_{1}, x'),
\end{equation*}
这也是光滑坐标变换,且
\begin{equation*}
\phi(x) = \tilde{x}_{1}^{2}/2 + \phi(X_{1}(x'), x').
\end{equation*}
再考虑 \eqref{eq:4.4.4},对 $\phi(X_{1}(x'), x')$ 同样继续,由归纳法到达结论.
\end{proof}

% ============ 第4章 §4.5 非退化振荡积分算子(原书 pp.336-338) ============
\section{非退化振荡积分算子}\label{sec:4.5}
\index{振荡积分算子}
\index{非退化条件}


作为上一节的自然推广,本节讨论非退化振荡积分算子.

设 $a(x, y) \in C_{0}^{\infty}(\mathbb{R}^{n}\times\mathbb{R}^{n})$.对实值函数 $\phi(x, y) \in C^{\infty}(\mathbb{R}^{n}\times\mathbb{R}^{n})$,考虑振荡积分算子
\begin{equation*}
T_{\lambda}f(x) := \int_{\mathbb{R}^{n}}e^{i\lambda\phi(x, y)}a(x, y)f(y)\dd y \quad (\lambda > 1)
\end{equation*}
的估计.当然,目标是数量地把握 $\lambda \to \infty$ 时的衰减.以下对 $\phi(x, y)$ 假设非退化条件
\begin{equation}
\det\left((\partial^{2}/\partial x_{j}\partial y_{k})\phi(x, y)\right) \neq 0.\tag{4.5.1}
\label{eq:4.5.1}\end{equation}

\begin{theorem}\label{thm:4.5.1}
若 $\phi(x, y)$ 在 $a(x, y)$ 的支集上满足 \eqref{eq:4.5.1},则
\begin{equation}
\left\|\int_{\mathbb{R}^{n}}e^{i\lambda\phi(x, y)}a(x, y)f(y)\dd y\right\|_{L^{2}(\mathbb{R}^{n})} \leq C\lambda^{-n/2}\|f\|_{L^{2}(\mathbb{R}^{n})}\tag{4.5.2}
\label{eq:4.5.2}\end{equation}
成立.
\end{theorem}

显然,
\begin{equation}
\left\|\int_{\mathbb{R}^{n}}e^{i\lambda\phi(x, y)}a(x, y)f(y)\dd y\right\|_{L^{\infty}(\mathbb{R}^{n})} \leq C\|f\|_{L^{1}(\mathbb{R}^{n})}.\tag{4.5.3}
\label{eq:4.5.3}\end{equation}
由 M. Riesz 插值定理插值 \eqref{eq:4.5.2},\eqref{eq:4.5.3},得到下面的系.

\begin{corollary}\label{cor:4.5.2}
设 $1 \leq p \leq 2$,$p' = p/(p-1)$,则
\begin{equation}
\|T_{\lambda}f\|_{L^{p'}(\mathbb{R}^{n})} \leq C\lambda^{-n/p'}\|f\|_{L^{p}(\mathbb{R}^{n})}\tag{4.5.4}
\label{eq:4.5.4}\end{equation}
成立.
\end{corollary}

若 $\phi(x, y) = -2\pi ix\cdot y$,则这满足 \eqref{eq:4.5.1} 的条件.变换变量,由 \eqref{eq:4.5.2} 可得
\begin{equation*}
\left\|\int_{\mathbb{R}^{n}}e^{-2\pi i x\cdot y}a(x/\sqrt{\lambda}, y/\sqrt{\lambda})f(y)\dd y\right\|_{L^{2}(\mathbb{R}^{n})} \leq C\|f\|_{L^{2}(\mathbb{R}^{n})}.
\end{equation*}
选取 $a(x, y)$ 在 $(x, y) = (0, 0)$ 处取 1,令 $\lambda \to \infty$,即得 Fourier 变换的 $L^{2}$ 有界性.

\begin{proof}[定理~\ref{thm:4.5.1} 的证明]
由 \eqref{eq:4.5.1} 的条件与 Taylor 定理,可以选取 $c > 0$,使得当 $|y - z|$ 充分小时
\begin{equation}
|\nabla_{x}[\phi(x, y) - \phi(x, z)]| \geq c|y - z|\tag{4.5.5}
\label{eq:4.5.5}\end{equation}
成立.必要时将 $a(x, y)$ 的支集用光滑的有限个单位分解分解,可假设在 $a(x, y)$ 的支集上 \eqref{eq:4.5.5} 成立.

首先改写为
\begin{equation}
\begin{aligned}
\|T_{\lambda}f\|_{2}^{2} & = \int\int K_{\lambda}(y, z)f(y)\overline{f(z)}\dd y\dd z, \\
K_{\lambda}(y, z) & := \int_{\mathbb{R}^{n}}e^{i\lambda[\phi(x, y) - \phi(x, z)]}a(x, y)\overline{a(x, z)}\dd x.
\end{aligned}
\end{equation}
考虑 \eqref{eq:4.5.5},与定理~\ref{thm:4.4.1} 的证明同样,可估计为
\begin{equation*}
|K_{\lambda}(y, z)| \leq C_{N}(1 + \lambda|y - z|)^{-N} \quad (N = 1, 2, \ldots).
\end{equation*}
取 $N = n + 1$,
\begin{equation*}
\int_{\mathbb{R}^{n}}(1 + \lambda|y|)^{-(n+1)}\dd y = C\lambda^{-n},
\end{equation*}
由 Young 不等式,
\begin{equation*}
\left\|\int(1 + \lambda|y - z|)^{-(n+1)}|f(y)|\dd y\right\|_{2} \leq C\lambda^{-n}\|f\|_{2}.
\end{equation*}
故由 Schwarz 不等式,
\begin{equation*}
\|T_{\lambda}f\|_{2}^{2} \leq \int\int(1 + \lambda|y - z|)^{-(n+1)}|f(y)||f(z)|\dd y\dd z \leq C\lambda^{-n}\|f\|_{2}^{2}.
\end{equation*}
到达目标.
\end{proof}

% ============ 第4章 §4.6 Fourier限制问题(Tomas-Stein定理)(原书 pp.338-346) ============
\section{Fourier 限制问题(Tomas--Stein 定理)}\label{sec:4.6}
\index{Fourier 限制问题}
\index{Tomas--Stein 定理}


在上一节中,利用在 $L^{2}$ 估计中可以直接写出其范数这一点,对振荡积分算子进行了估计.本节将进一步深入讨论,考察基于 $L^{2}$ 理论的 Fourier 限制问题.

由系~\ref{cor:4.3.10},$f \in L^{p}(\mathbb{R}^{n})$($1 \leq p \leq 2$)的 Fourier 变换 $\hat{f}$ 是函数,且满足
\begin{equation}
\|\hat{f}\|_{p'} \leq \|f\|_{p} \quad \left(p' = \frac{p}{p-1}\right).\tag{4.6.1}
\label{eq:4.6.1}\end{equation}

$L^{1}(\mathbb{R}^{n})$ 的 Fourier 像是既有界又连续的函数,故将 Fourier 像整体限制在 $\mathbb{R}^{n}$ 的任意子集上都是有意义的.另一方面,$L^{2}(\mathbb{R}^{n})$ 的 Fourier 像是整个 $L^{2}(\mathbb{R}^{n})$,故将 Fourier 像整体限制在 Lebesgue 测度为 0 的集合上没有意义.那么,处于这两者之间的 $L^{p}(\mathbb{R}^{n})$($1 < p < 2$)是如何反映这一极端差异的呢?

令 $\phi \in \mathcal{S}(\mathbb{R}^{n-1})$,
\begin{equation*}
f(x_{1}, x_{2}, \ldots, x_{n}) := \frac{1}{1 + |x_{1}|}\phi(x_{2}, \ldots, x_{n}),
\end{equation*}
则 $f \in L^{p}(\mathbb{R}^{n})$($p > 1$),且在超平面 $\{\xi_{1} = 0\}$ 上 $\hat{f}(\xi) = \infty$.但是,1967 年 E. M. Stein 证明了:在某个 $p$ 的范围内,可以将 $L^{p}(\mathbb{R}^{n})$ 的 Fourier 像限制到具有某种曲率的超曲面上.当然,由 \eqref{eq:4.6.1},限制到正测度集合上总是可能的.

这里讨论单位球面
\begin{equation*}
S^{n-1} = \{x \in \mathbb{R}^{n}: |x| = 1\}
\end{equation*}
上的 Fourier 限制问题.

证明下面的定理是本节的目标.以下用 $d\sigma$ 表示定义在 $S$ 上的 Lebesgue 测度.

\begin{theorem}[Tomas--Stein]\label{thm:4.6.1}
\begin{equation}
\left(\int_{S}|\hat{f}(\xi)|^{q}\dd\sigma\right)^{1/q} \leq C_{p}\|f\|_{L^{p}(\mathbb{R}^{n})} \quad (f \in \mathcal{S}(\mathbb{R}^{n}))\tag{4.6.2}
\label{eq:4.6.2}\end{equation}
当
\begin{equation*}
q = \frac{n-1}{n+1}\frac{p}{p-1}
\end{equation*}
时,在
(1) $1 \leq p \leq \frac{2(n+1)}{n+3}$,$n \geq 3$,
(2) $1 \leq p < \frac{4}{3}$,$n = 2$
的情形下成立.
\end{theorem}

令
\begin{equation*}
\begin{cases}
D := \{y \in \mathbb{R}^{n-1}: |y| \leq 1/2\}, \\
h(y) := \sqrt{1 - |y|^{2}} \quad (y \in D), \\
\tilde{D} := \{(y, h(y)) \in S: y \in D\},
\end{cases}
\end{equation*}
则在 $\tilde{D}$ 上
\begin{equation*}
\dd\sigma = \sqrt{1 + |\nabla h(y)|^{2}}\dd y = \frac{\dd y}{h(y)}.
\end{equation*}
考虑对称性与光滑的单位分解,为证明 \eqref{eq:4.6.2},固定 $0 \leq b(y) \in C_{0}^{\infty}(D)$,只要证明
\begin{equation*}
\left(\int_{\mathbb{R}^{n-1}}|\hat{f}(y, h(y))|^{q}a(y)\dd y\right)^{1/q} \leq C\|f\|_{L^{p}(\mathbb{R}^{n})}, \quad a(y) := \frac{b(y)}{h(y)}
\end{equation*}
即可.再考虑对偶性,证明
\begin{equation}
\left|\int_{\mathbb{R}^{n-1}}\hat{f}(y, h(y))g(y)a(y)\dd y\right| \leq C\|f\|_{L^{p}(\mathbb{R}^{n})}\times\|g\|_{L^{q'}(\mathbb{R}^{n-1})}.\tag{4.6.3}
\label{eq:4.6.3}\end{equation}

由于
\begin{equation}
\begin{aligned}
& \int_{\mathbb{R}^{n-1}}\hat{f}(y, h(y))g(y)a(y)\dd y \\
& = \int_{\mathbb{R}^{n-1}}\int_{\mathbb{R}^{n}}e^{-2\pi i(z_{1}y_{1} + \ldots + z_{n-1}y_{n-1} + z_{n}h(y))}f(z)g(y)a(y)\dd z\dd y,
\end{aligned}
\end{equation}
\eqref{eq:4.6.3} 归结为证明
\begin{equation}
\left\|\int_{\mathbb{R}^{n-1}}e^{-2\pi i(z_{1}y_{1} + \ldots + z_{n-1}y_{n-1} + z_{n}h(y))}g(y)a(y)\dd y\right\|_{L^{p'}(\mathbb{R}^{n})} \leq C\|g\|_{L^{q'}(\mathbb{R}^{n-1})}.\tag{4.6.4}
\label{eq:4.6.4}\end{equation}

用上一节的振荡积分算子方法分析 \eqref{eq:4.6.4}.即设 $a(z, y) \in C_{0}^{\infty}(\mathbb{R}^{n}\times\mathbb{R}^{n-1})$,证明
\begin{equation}
\begin{aligned}
& \left\|\int_{\mathbb{R}^{n-1}}e^{-2\pi i\lambda(z_{1}y_{1} + \ldots + z_{n-1}y_{n-1} + z_{n}h(y))}a(z, y)g(y)\dd y\right\|_{L^{p'}(\mathbb{R}^{n})} \\
& \leq C\lambda^{-n/p'}\|g\|_{L^{q'}(\mathbb{R}^{n-1})}
\end{aligned}\tag{4.6.5}
\label{eq:4.6.5}\end{equation}
成立.只要有这个式子,变形为
\begin{equation}
\begin{aligned}
& \left\|\int_{\mathbb{R}^{n-1}}e^{-2\pi i(z_{1}y_{1} + \ldots + z_{n-1}y_{n-1} + z_{n}h(y))}a(z/\lambda, y)g(y)\dd y\right\|_{L^{p'}(\mathbb{R}^{n})} \\
& \leq C\|g\|_{L^{q'}(\mathbb{R}^{n-1})},
\end{aligned}
\end{equation}
令 $\lambda \to \infty$,即得 \eqref{eq:4.6.4}.以下,令
\begin{equation*}
\phi(z, y) := z_{1}y_{1} + \ldots + z_{n-1}y_{n-1} + z_{n}h(y),
\end{equation*}
考察振荡积分算子
\begin{equation*}
T_{\lambda}g(z) := \int_{\mathbb{R}^{n-1}}e^{i\lambda\phi(z, y)}a(z, y)g(y)\dd y.
\end{equation*}

整理要点,证明下面的命题.

\begin{proposition}\label{prop:4.6.2}
\begin{equation}
\|T_{\lambda}g\|_{L^{p'}(\mathbb{R}^{n})} \leq C\lambda^{-n/p'}\|g\|_{L^{q'}(\mathbb{R}^{n-1})}\tag{4.6.6}
\label{eq:4.6.6}\end{equation}
当 $p' = \frac{n+1}{n-1}\frac{q'}{q'-1}$ 时,在
(1) $1 \leq q' \leq 2$,$n \geq 3$,
(2) $1 \leq q' < 4$,$n = 2$
的情形下成立.
\end{proposition}

\begin{proof}[命题~\ref{prop:4.6.2} (1) 的证明]
$(q', p') = (1, \infty)$ 的情形显然,故证明 $(q', p') = \left(2, \frac{2(n+1)}{n-1}\right)$ 的情形:
\begin{equation}
\|T_{\lambda}g\|_{L^{\frac{2(n+1)}{n-1}}(\mathbb{R}^{n})} \leq C\lambda^{-\frac{n(n-1)}{2(n+1)}}\|g\|_{L^{2}(\mathbb{R}^{n-1})},\tag{4.6.7}
\label{eq:4.6.7}\end{equation}
再用插值法导出结论.

考虑对偶性,为证明 \eqref{eq:4.6.7},令
\begin{equation*}
T_{\lambda}^{*}f(y) := \int_{\mathbb{R}^{n}}e^{i\lambda\phi(z, y)}a(z, y)f(z)\dd z,
\end{equation*}
证明
\begin{equation}
\|T_{\lambda}^{*}f\|_{L^{2}(\mathbb{R}^{n-1})} \leq C\lambda^{-\frac{n(n-1)}{2(n+1)}}\|f\|_{L^{\frac{2(n+1)}{n+3}}(\mathbb{R}^{n})}.\tag{4.6.8}
\label{eq:4.6.8}\end{equation}
此左端可以直接写出:
\begin{equation}
\begin{aligned}
\|T_{\lambda}^{*}f\|_{L^{2}(\mathbb{R}^{n-1})}^{2} & = \int_{\mathbb{R}^{n-1}}\int_{\mathbb{R}^{n}}\int_{\mathbb{R}^{n}}e^{i\lambda[\phi(z, y) - \phi(w, y)]}a(z, y)\overline{a(w, y)}f(z)\overline{f(w)}\dd z\dd w\dd y \\
& \leq \|\tilde{K}_{\lambda}f\|_{L^{\frac{2(n+1)}{n-1}}(\mathbb{R}^{n})}\times\|f\|_{L^{\frac{2(n+1)}{n+3}}(\mathbb{R}^{n})}.
\end{aligned}
\end{equation}
这里
\begin{equation}
\begin{aligned}
\tilde{K}_{\lambda}f(z) & := T_{\lambda}\overline{T_{\lambda}^{*}f}(z) = \int_{\mathbb{R}^{n}}K_{\lambda}(z, w)\overline{f(w)}\dd w, \\
K_{\lambda}(z, w) & := \int_{\mathbb{R}^{n-1}}e^{i\lambda[\phi(z, y) - \phi(w, y)]}a(z, y)\overline{a(w, y)}\dd y.
\end{aligned}
\end{equation}
因此,为证明 \eqref{eq:4.6.8},令
\begin{equation*}
K_{\lambda}f(z) := \int_{\mathbb{R}^{n}}K_{\lambda}(z, w)f(w)\dd w,
\end{equation*}
当 $r = \frac{2(n+1)}{n+3}$ 时证明
\begin{equation}
\|K_{\lambda}f\|_{L^{r'}(\mathbb{R}^{n})} \leq C\lambda^{-\frac{n(n-1)}{n+1}}\|f\|_{L^{r}(\mathbb{R}^{n})}.\tag{4.6.9}
\label{eq:4.6.9}\end{equation}
这个式子的有用之处在于,两边都能估计同一维数的函数.

将 $z = (z', u)$($z' \in \mathbb{R}^{n-1}$),$w = (w', v)$($w' \in \mathbb{R}^{n-1}$)分别分割,固定 $u, v$,令
\begin{equation}
\begin{aligned}
K_{u,v}^{\lambda}f(z') & := \int_{\mathbb{R}^{n-1}}K_{u,v}^{\lambda}(z', w')f(w', v)\dd w', \\
K_{u,v}^{\lambda}(z', w') & := K_{\lambda}(z', u; w', v),
\end{aligned}
\end{equation}
若证明
\begin{equation}
\|K_{u,v}^{\lambda}f\|_{L^{r'}(\mathbb{R}^{n-1})} \leq C\lambda^{-\frac{n(n-1)}{n+1}}|u - v|^{-1+(1/r - 1/r')}\|f(\cdot, v)\|_{L^{r}(\mathbb{R}^{n-1})},\tag{4.6.10}
\label{eq:4.6.10}\end{equation}
则经由一维 Hardy--Littlewood--Sobolev 不等式到达 \eqref{eq:4.6.9}.

实际上,由于
\begin{equation*}
K_{u}^{\lambda}f(z') := \int_{\mathbb{R}^{n}}K_{\lambda}(z', u; w)f(w)\dd w = \int_{-\infty}^{\infty}K_{u,v}^{\lambda}f(z')\dd v,
\end{equation*}
由 Minkowski 积分不等式,
\begin{equation}
\begin{aligned}
\|K_{u}^{\lambda}f\|_{L^{r'}(\mathbb{R}^{n-1})} & \leq \int_{-\infty}^{\infty}\|K_{u,v}^{\lambda}f\|_{L^{r'}(\mathbb{R}^{n-1})}\dd v \\
& \leq C\lambda^{-\frac{n(n-1)}{n+1}}\int_{-\infty}^{\infty}|u - v|^{-1+(1/r - 1/r')}\|f(\cdot, v)\|_{L^{r}(\mathbb{R}^{n-1})}\dd v.
\end{aligned}
\end{equation}
上式左端关于 $u$ 的 $L^{r'}(\mathbb{R})$ 范数等于 \eqref{eq:4.6.9} 的左端,故应用 Hardy--Littlewood--Sobolev 不等式(定理~\ref{thm:4.2.3}),
\begin{equation}
\begin{aligned}
\|K_{\lambda}f\|_{L^{r'}(\mathbb{R}^{n})} & \leq C\lambda^{-\frac{n(n-1)}{n+1}}\left\|\int_{-\infty}^{\infty}|u - v|^{-1+(1/r - 1/r')}\|f(\cdot, v)\|_{L^{r}(\mathbb{R}^{n-1})}\dd v\right\|_{L^{r'}(\mathbb{R})} \\
& \leq C\lambda^{-\frac{n(n-1)}{n+1}}\left\|\|f(\cdot, v)\|_{L^{r}(\mathbb{R}^{n-1})}\right\|_{L^{r}(\mathbb{R})} \\
& = C\lambda^{-\frac{n(n-1)}{n+1}}\|f\|_{L^{r}(\mathbb{R}^{n})}.
\end{aligned}
\end{equation}

\eqref{eq:4.6.10} 由 $L^{2} \to L^{2}$,$L^{1} \to L^{\infty}$ 的插值法得到.

前者,将
\begin{equation}
K_{u,v}^{\lambda}f(z') = \int_{\mathbb{R}^{n-1}}e^{i\lambda\phi(z, y)}a(z, y)\int_{\mathbb{R}^{n-1}}\overline{e^{i\lambda\phi(w, y)}a(w, y)}f(w)\dd w'\dd y
\end{equation}
还原,并想起
\begin{equation*}
\phi(w, y) = w_{1}y_{1} + \ldots + w_{n-1}y_{n-1} + uh(y),
\end{equation*}
则
\begin{equation*}
\det\left(\left(\partial^{2}/\partial w_{j}\partial y_{k}\right)\phi(w, y)\right)_{j, k \leq n-1} = 1 \neq 0,
\end{equation*}
故应用定理~\ref{thm:4.5.1},
\begin{equation}
\begin{aligned}
& \left\|\int_{\mathbb{R}^{n-1}}\overline{e^{i\lambda\phi(w, y)}a(w, y)}f(w)\dd w'\right\|_{L^{2}(\mathbb{R}^{n-1})} \\
& \leq C\lambda^{-(n-1)/2}\|f(\cdot, v)\|_{L^{2}(\mathbb{R}^{n-1})}.
\end{aligned}
\end{equation}
同样考虑,
\begin{equation}
\|K_{u,v}^{\lambda}f\|_{L^{2}(\mathbb{R}^{n-1})} \leq C\lambda^{-(n-1)}\|f(\cdot, v)\|_{L^{2}(\mathbb{R}^{n-1})}.\tag{4.6.11}
\label{eq:4.6.11}\end{equation}

用驻相法证明后者.需要估计的核
\begin{equation}
K_{u,v}^{\lambda}(z', w') = \int_{\mathbb{R}^{n-1}}e^{i\lambda\Phi(y, z, w)}a(z, y)\overline{a(w, y)}\dd y,
\end{equation}
\begin{equation*}
\Phi(y, z, w) := (z_{1} - w_{1})y_{1} + \ldots + (z_{n-1} - w_{n-1})y_{n-1} + (u - v)h(y),
\end{equation*}
验证
\begin{equation}
|K_{u,v}^{\lambda}(z', w')| \leq C(\lambda|u - v|)^{-(n-1)/2}.\tag{4.6.12}
\label{eq:4.6.12}\end{equation}

将
\begin{equation}
\begin{aligned}
& (z_{1} - w_{1})y_{1} + \ldots + (z_{n-1} - w_{n-1})y_{n-1} + (u - v)h(y) \\
& = (u - v)\left(\frac{z_{1} - w_{1}}{u - v}y_{1} + \ldots + \frac{z_{n-1} - w_{n-1}}{u - v}y_{n-1} + h(y)\right)
\end{aligned}
\end{equation}
中的 $\mathbb{R}^{n}$ 上的点设为
\begin{equation*}
A(z, w) := \frac{1}{u - v}(z_{1} - w_{1}, \ldots, z_{n-1} - w_{n-1}, 1),
\end{equation*}
则可知,仅当连接原点与 $A(z, w)$ 的直线与 $\tilde{D}$ 有交点 $(y_{0}, h(y_{0}))$ 时,
\begin{equation*}
\nabla_{y}\Phi(y_{0}, w, z) = 0.
\end{equation*}
不失一般性,可设
\begin{equation*}
(y_{0}, h(y_{0})) = A(z_{0}, w_{0}) = (0, \ldots, 0, 1)
\end{equation*}
继续验证.此时
\begin{equation*}
\Phi(y, z_{0}, w_{0}) = h(y),
\end{equation*}
故
\begin{equation*}
\det\left(\left(\partial^{2}/\partial y_{j}\partial y_{k}\right)\Phi(y, z_{0}, w_{0})\right) = 1 \neq 0 \quad (y = 0),
\end{equation*}
由定理~\ref{thm:4.4.2} 可估计为
\begin{equation*}
\left|K_{u,v}^{\lambda}(0, z_{0}, w_{0})\right| \leq C(\lambda|u - v|)^{-(n-1)/2}.
\end{equation*}
其他情形得到比这更好的估计,作为结论 \eqref{eq:4.6.12} 得证.由此得到
\begin{equation}
\|K_{u,v}^{\lambda}f\|_{L^{\infty}(\mathbb{R}^{n-1})} \leq C\lambda^{-(n-1)/2}|u - v|^{-(n-1)/2}\|f(\cdot, v)\|_{L^{1}(\mathbb{R}^{n-1})}.\tag{4.6.13}
\label{eq:4.6.13}\end{equation}

利用
\begin{equation*}
\begin{cases}
\frac{1}{r} = \frac{1}{2}\left(1 - \frac{2}{n+1}\right) + \frac{2}{n+1}, \\
-\frac{n(n-1)}{n+1} = -(n-1)\left(1 - \frac{2}{n+1}\right) - \frac{n-1}{2}\frac{2}{n+1}, \\
-1 + \left(1/r - 1/r'\right) = -\frac{n-1}{n+1} = -\frac{n-1}{2}\frac{2}{n+1},
\end{cases}
\end{equation*}
插值 \eqref{eq:4.6.11},\eqref{eq:4.6.13},\eqref{eq:4.6.10} 得证.
\end{proof}

\begin{proof}[命题~\ref{prop:4.6.2} (2) 的证明]
这里稍微改变形式,令
\begin{equation}
\begin{aligned}
T_{\lambda}g(z) & := \int_{-\infty}^{\infty}e^{i\lambda\phi(z, \theta)}a(z, \theta)g(\theta)\dd\theta, \\
\phi(z, \theta) & := z_{1}\cos\theta + z_{2}\sin\theta,
\end{aligned}
\end{equation}
在 $p' = 3\frac{q'}{q'-1}$,$1 \leq q' < 4$ 中证明
\begin{equation}
\|T_{\lambda}g\|_{L^{p'}(\mathbb{R}^{2})} \leq C\lambda^{-2/p'}\|g\|_{L^{q'}(\mathbb{R})}.\tag{4.6.14}
\label{eq:4.6.14}\end{equation}
证明的关键在于 $p' > 4$.在这种情况下,可以估计
\begin{equation}
(T_{\lambda}g(z))^{2} = \int_{\mathbb{R}^{2}}e^{i\lambda\Phi(z, y)}b(z, y)g(y_{1})g(y_{2})\dd y,
\end{equation}
\begin{equation*}
\Phi(z, y) := z_{1}(\cos y_{1} + \cos y_{2}) + z_{2}(\sin y_{1} + \sin y_{2}), \quad b(z, y) := a(z, y_{1})a(z, y_{2}).
\end{equation*}

想对
\begin{equation*}
\int_{\mathbb{R}^{2}}e^{i\lambda\Phi(z, y)}b(z, y)f(y)\dd y
\end{equation*}
应用定理~\ref{thm:4.5.1}.但是,由于
\begin{equation*}
\left|\begin{array}{cc}
(\partial^{2}/\partial z_{1}\partial y_{1})\Phi(z, y) & (\partial^{2}/\partial z_{1}\partial y_{2})\Phi(z, y) \\
(\partial^{2}/\partial z_{2}\partial y_{1})\Phi(z, y) & (\partial^{2}/\partial z_{2}\partial y_{2})\Phi(z, y)
\end{array}\right| = -\sin(y_{1} - y_{2}),
\end{equation*}
不能这样直接应用.于是,变换坐标
\begin{equation*}
(u_{1}, u_{2}) = (y_{1} - y_{2}, y_{1} + y_{2}),
\end{equation*}
则 $|du/dy| = 2$,故必要时将 $b$ 的支集分解小,可得
\begin{equation*}
|\det(\partial^{2}\Phi/\partial z\partial u)| \geq c|u_{1}|.
\end{equation*}
进而
\begin{equation}
\begin{aligned}
\Phi(z, u) & = z_{1}\left(\cos\frac{u_{2} + u_{1}}{2} + \cos\frac{u_{2} - u_{1}}{2}\right) \\
& \quad + z_{2}\left(\sin\frac{u_{2} + u_{1}}{2} + \sin\frac{u_{2} - u_{1}}{2}\right),
\end{aligned}
\end{equation}
\begin{equation*}
b(z, u) = a\left(z, \frac{u_{2} + u_{1}}{2}\right)a\left(z, \frac{u_{2} - u_{1}}{2}\right),
\end{equation*}
故 $\Phi(z, u)$,$b(z, u)$ 都是 $u_{1}$ 的偶函数,可以看作 $u_{1}^{2}$ 的连续函数.于是,变换坐标
\begin{equation*}
(v_{1}, v_{2}) = \left(\frac{1}{2}u_{1}^{2}, u_{2}\right),
\end{equation*}
则 $|dv/du| = |u_{1}|$,故
\begin{equation*}
|\det(\partial^{2}\Phi/\partial z\partial v)| \geq c
\end{equation*}
成立.这里,对
\begin{equation*}
\tilde{T}_{\lambda}f(z) := \int_{\mathbb{R}^{2}}e^{i\lambda\Phi(z, v)}b(z, v)f(v)\dd v
\end{equation*}
应用定理~\ref{thm:4.5.1},特别地系~\ref{cor:4.5.2} 成立.即
\begin{equation}
\|\tilde{T}_{\lambda}f\|_{L^{r'}(\mathbb{R}^{2})} \leq C\lambda^{-2/r'}\|f\|_{L^{r}(\mathbb{R}^{2})} \quad (1 \leq r \leq 2).\tag{4.6.15}
\label{eq:4.6.15}\end{equation}

由定义可计算为
\begin{equation}
\begin{aligned}
\|T_{\lambda}g\|_{L^{p'}(\mathbb{R}^{2})}^{2} & = \|(T_{\lambda}g)^{2}\|_{L^{p'/2}(\mathbb{R}^{2})} \\
& = \|\tilde{T}_{\lambda}g(\cdot)g(\cdot)|dv/dy|^{-1}\|_{L^{p'/2}(\mathbb{R}^{2})},
\end{aligned}
\end{equation}
对右端应用 \eqref{eq:4.6.15},
\begin{equation}
\begin{aligned}
\|T_{\lambda}g\|_{L^{p'}(\mathbb{R}^{2})}^{2} & \leq C\lambda^{-4/p'}\left(\int_{\mathbb{R}^{2}}g(y_{1})^{r}g(y_{2})^{r}|y_{1} - y_{2}|^{1-r}\dd y\right)^{1/r}, \\
& \qquad r = p'/(p' - 2).
\end{aligned}
\end{equation}
这里,可计算出
\begin{equation*}
1 - r = -1 + (2 - r), \quad 2 - r = [q'/r]^{-1} - [(q'/r)']^{-1},
\end{equation*}
用一维 Hardy--Littlewood--Sobolev 不等式到达结论.
\end{proof}

% ============ 第4章 §4.7 Nikodym极大函数(Wolff定理)(原书 pp.346-367) ============
\section{Nikodym 极大函数(Wolff 定理)}\label{sec:4.7}
\index{Nikodym 极大函数!小 Nikodym 极大函数}
\index{Nikodym 极大函数!大 Nikodym 极大函数}
\index{Wolff 定理}
\index{Córdoba 筛}
\index{bush 论证}
\index{tube}
\index{tube!$\delta$-tube}
\index{三点插值法}


首先,让我们回顾 Hardy--Littlewood 极大函数.

对局部可积函数 $f \in L_{\mathrm{loc}}^{1}(\mathbb{R}^{n})$,Hardy--Littlewood 极大函数定义为
\begin{equation*}
\mathcal{M}f(x) := \sup_{t>0}\frac{1}{|B(x,t)|}\int_{B(x,t)}|f(y)|\dd y.
\end{equation*}
这里,$B(x,t) \subset \mathbb{R}^{n}$ 表示以 $x$ 为中心,半径 $t$ 的开球,$|B(x,t)|$ 表示其体积.

该极大函数在 $L^{p}(\mathbb{R}^{n})$($p > 1$)上有界这一事实的用途,已通过如下方式得到确认:由它导出 Hardy--Littlewood--Sobolev 不等式,并在上一节的最后本质性地使用了该不等式.也请回忆,有界性成立的核心在于覆盖引理(引理~\ref{lem:4.1.3}).

在 Fourier 分析中,或许是由于测不准原理这一远因,人们会研究一种别具一格的极大函数的估计.

设 $N \gg 1$ 为非常大的参数.对 $a > 0$,(小的)Nikodym 极大函数定义为
\begin{equation}\tag{4.7.1}
\mathcal{K}_{N,a}f(x) := \sup_{T \ni x}\frac{1}{|T|}\int_{T}|f(y)|\dd y.
\label{eq:4.7.1}\end{equation}
这里,上确界对一切与具有形状
\begin{align*}
&a \times \cdots \times a \times (Na) \\
\intertext{即}
&\left\{(x',x_n) \in \mathbb{R}^{n}:\ |x'| < \frac{a}{2},\ |x_n| < \frac{Na}{2}\right\}
\end{align*}
的 tube 合同的,包含 $x$ 且方向(长轴方向)任意的 tube $T$ 取.进而,(大的)Nikodym 极大函数定义为
\begin{equation}\tag{4.7.2}
\mathcal{K}_{N}f(x) := \sup_{a>0}\mathcal{K}_{N,a}f(x).
\label{eq:4.7.2}\end{equation}

Nikodym 极大函数 $\mathcal{K}_{N}$ 的 $L^{p}(\mathbb{R}^{n})$($p > 1$)有界性,用于控制函数的相对大小关于方向的分布.对出现在该极大函数中的 tube 族,回忆覆盖引理(引理~\ref{lem:4.1.3})的证明并加以应用:若长轴方向全部固定,则可原样应用;但若容许方向具有自由度,则不再是以 3 倍包含其余 tube,而是需要在 $(n-1)$ 维方向上放大 $3N$ 倍,在剩下的 1 维方向上放大 3 倍,于是只能得到如下朴素的估计.
\begin{equation*}
\|\mathcal{K}_{N}f\|_{L^{p}(\mathbb{R}^{n})} \leq C_{p,N}\|f\|_{L^{p}(\mathbb{R}^{n})}, \quad C_{p,N} = 2\left(\frac{(3N)^{n-1}3p}{p-1}\right)^{1/p}.
\end{equation*}
这里的论点是,求出右端出现的常数 $C_{p,N}$ 关于 $N$ 的最优估计.为了改进上面的估计应当使用怎样的方法呢?我们将在本节中通篇思考这个问题.

由于是极大函数的话题,以下在本节中假定出现的所有函数都是非负的.

\subsection{小 Nikodym 极大函数}\label{sec:4.7.1}

此后设 $N \gg 1$ 为大的整数参数,并令 $\delta = 1/N$.所谓 $\delta$-tube,是指与
\[
\{(x',x_n) \in \mathbb{R}^{n}:\ |x'| < \delta/2,\ |x_n| < 1/2\}
\]
合同的 tube,$T$ 的方向指长轴的方向.用 $l(T)$ 表示其轴所在的直线,用 $\omega(T) \in S^{n-1}$ 表示其一单位方向向量.暂时,我们将在 $\mathbb{R}^{n}$ 的原点附近进行细致的讨论.

令 $Q_{1} := (-1,1)^{n}$,$Q_{2} := (-2,2)^{n}$.对 $f \in L^{1}(Q_{1})$,小的 Nikodym 极大函数定义为
\begin{equation*}
K_{\delta}f(x) := \sup_{T \ni x}\frac{1}{|T|}\int_{T}f(y)\dd y.
\end{equation*}
这里,上确界对包含 $x$ 且方向任意的 $\delta$-tube $T$ 取.

关于该极大函数,如下猜想问题至今仍未解决.尽管这一问题已由许多优秀的数学家进行研究,但长期以来仍悬而未决.

\textbf{猜想问题.} 存在常数 $C$,$\alpha$,使得
\begin{equation}\tag{4.7.3}
\|K_{\delta}f\|_{L^{n}(Q_{2})} \leq C(\log N)^{\alpha}\|f\|_{L^{n}(Q_{1})}
\label{eq:4.7.3}\end{equation}
成立.

现在,若取
\begin{align*}
f(x) &= \frac{\chi_{[\delta,1]}(|x|)}{|x|}, \\
\intertext{则可得估计}
K_{\delta}f(x) &\geq C^{-1}\log N \quad (|x| < 1), \\
\intertext{从而}
\|K_{\delta}f\|_{L^{n}(Q_{2})} &\geq C^{-1}\log N. \\
\intertext{另一方面,由于}
\|f\|_{L^{n}(Q_{1})} &= C(\log N)^{1/n},
\end{align*}
\eqref{eq:4.7.3}的右端需要出现 $\log N$ 的项.

正如后文将看到的,在 $n = 2$ 时已知此猜想成立.

若 $f \in L^{1}(Q_{1})$,则
\[
K_{\delta}f(x) \leq C N^{n-1}\|f\|_{1},
\]
从而
\begin{equation}\tag{4.7.4}
\|K_{\delta}f\|_{L^{\infty}(Q_{2})} \leq C N^{n-1}\|f\|_{L^{1}(Q_{1})}.
\label{eq:4.7.4}\end{equation}
于是,对\eqref{eq:4.7.3},\eqref{eq:4.7.4}进行插值,作为猜想问题的一个版本,我们来考虑如下问题.

\textbf{猜想问题.} 在 $1 \leq p \leq n$ 时
\begin{equation}\tag{4.7.5}
\left\|K_{\delta}f\right\|_{L^{q}(Q_{2})} \leq C N^{n/p-1}(\log N)^{\alpha n'/p'}\|f\|_{L^{p}(Q_{1})} \quad \left(q = (n-1)p'\right)
\label{eq:4.7.5}\end{equation}
成立.

该猜想问题在 $p = 1$ 时当然成立.因此,论点是使成立的 $p$ 的范围向 $p = n$ 扩展.($\log N$)的指数大小在应用上不成问题.

证明如下定理是本节的目标.

\begin{theorem}[T. Wolff, 1995]\label{thm:4.7.1}
在 $1 \leq p \leq (n+2)/2$ 时\eqref{eq:4.7.5}成立.特别地,
\begin{equation}\tag{4.7.6}\begin{aligned}
\|K_{\delta}f\|_{L^{q}(Q_{2})} & \leq C N^{n/p-1}(\log N)^{\alpha}\|f\|_{L^{p}(Q_{1})} \\
& \quad (p = (n+2)/2,\ q = (n-1)p')
\end{aligned}\label{eq:4.7.6}\end{equation}
这里 $\alpha$ 表示仅依赖于维数 $n$ 的常数.
\end{theorem}

正如后文将确认的,该定理由如下使用分布函数的命题推出.

\begin{proposition}\label{prop:4.7.2}
设 $p = (n+2)/2$,$0 < \lambda < 1$,且 $\|f\|_{L^{\infty}(Q_{1})} = 1$,则
\begin{equation}\tag{4.7.7}\begin{aligned}
\lambda^{p}\delta^{n-p}|\{x \in Q_{2}:\ K_{\delta}f(x) > \lambda\}|^{(p-1)/(n-1)}&\leq C(\log N)^{(n+5)/2}\int_{Q_{1}}f(y)^{2}\dd y
 \end{aligned}\label{eq:4.7.7}\end{equation}
成立.
\end{proposition}

\subsection{命题~\ref{prop:4.7.2} 的证明}\label{sec:4.7.2}

将 $Q_{2}$ 分割为边长为 $\delta$ 的立方体,把由各立方体中心点组成的集合记为 $I$.$I$ 的个数 $\#I$ 为 $\#I = (2N)^{n}$.对 $i \in I$,用 $Q_{i}$ 表示以 $i$ 为中心,边长为 $\delta$ 的分割立方体.

固定 $0 < \lambda < 1$,令
\begin{align*}
\tilde{\mathcal{I}} &:= \mathbf{I} \cap \{x \in Q_{2}:\ K_{\delta}f(x) > \lambda\}, \\
\intertext{则}
\{x \in Q_{2}:\ K_{\delta}f(x) > \lambda\} &\subset \bigcup_{i \in \tilde{\mathcal{I}}}Q_{i}, \\
\intertext{因此我们来估计}
\sum_{i \in \tilde{\mathcal{I}}}|Q_{i}| &= \delta^{n}\#\tilde{\mathcal{I}}.
\end{align*}

由极大函数的定义,对每个 $i \in \tilde{I}$ 各选取一个 $\delta$-tube $T_{i} \ni i$,使得
\begin{equation}\tag{4.7.8}
\lambda|T_{i}| < \int_{T_{i}}f(y)\dd y, \quad i \in T_{i} \quad (i \in \tilde{\mathcal{I}})
\label{eq:4.7.8}\end{equation}
成立.

令
\[
\tilde{D} := \left\{(\omega',\omega_{n}) \in S^{n-1}:\ \omega_{n} > \frac{1}{\sqrt{n}}\right\},
\]
并(在必要时让依赖于维数的常数进入常数中)假定 $\omega(T_{i}) \in \tilde{D}$.

设集合 $\tilde{I}$ 被超平面 $\{x_{n} = t\}$ 截得的断面中出现的点的个数
\begin{align*}
\#(\tilde{\mathcal{I}} \cap \{x_{n} = t\})& \\
\intertext{在 $t = t_{0}$ 处取最大值,令}
\mathbf{I}_{0} &:= \mathbf{I} \cap \{x_{n} = t_{0}\}, \quad \mathcal{I}_{0} := \tilde{\mathcal{I}} \cap \{x_{n} = t_{0}\}, \\
\intertext{则}
\#\tilde{\mathcal{I}} &\leq 2N\#\mathcal{I}_{0}.
\end{align*}
综合以上,为证明\eqref{eq:4.7.7},只需证明
\begin{equation}\tag{4.7.9}
\lambda^{(n+2)/2}\delta^{(n-2)/2}\left(\delta^{n-1}\#\mathcal{I}_{0}\right)^{n/2(n-1)} \leq C(\log N)^{(n+5)/2}\int_{Q_{1}}f(y)^{2}\dd y
\label{eq:4.7.9}\end{equation}
成立即可.

下面的几何引理将被基本地使用.

\begin{lemma}\label{lem:4.7.3}
设 $i, j \in \mathcal{I}_{0}$,且 $T_{i} \cap T_{j} \neq \emptyset$,则
\begin{align*}
(1)\quad |T_{i} \cap T_{j}| &\leq C\frac{\delta^{n}}{|\omega(T_{i}) - \omega(T_{j})| + \delta}, \\
(2)\quad |i - j| &\leq |\omega(T_{i}) - \omega(T_{j})|
\end{align*}
成立.
\end{lemma}

两者的证明都很容易.(1)由 $\omega(T_{i}), \omega(T_{j}) \in \tilde{D}$ 推出.继续讨论.使用\eqref{eq:4.7.8}计算可得
\begin{align*}
\lambda^{2}\delta^{2n-2}(\#\mathcal{I}_{0})^{2} & = \left(\lambda\delta^{n-1}\#\mathcal{I}_{0}\right)^{2} \\
& \leq C\left(\sum_{i \in \mathcal{I}_{0}}\int_{T_{i}}f(y)\dd y\right)^{2} \\
&= C\left(\int\sum_{i \in \mathcal{I}_{0}}\chi_{T_{i}}(y)f(y)\dd y\right)^{2} \\
&\leq C\int\left(\sum_{i \in \mathcal{I}_{0}}\chi_{T_{i}}(y)\right)^{2}\dd y \cdot \int f(y)^{2}\dd y.
\end{align*}
下面我们来估计
\begin{equation}\tag{4.7.10}
\int\left(\sum_{i \in \mathcal{I}_{0}}\chi_{T_{i}}(y)\right)^{2}\dd y.
\label{eq:4.7.10}\end{equation}

\textbf{情形 $n = 2$.} 若 $n = 2$,则由上面的引理和上式,讨论即告完成.实际上,
\[
\eqref{eq:4.7.10} = \sum_{i \in \mathcal{I}_{0}}\sum_{j \in \mathcal{I}_{0}}|T_{i} \cap T_{j}| \leq C\sum_{i \in \mathcal{I}_{0}}\sum_{j \in \mathcal{I}_{0}}\frac{\delta^{2}}{|i-j|+\delta} \leq C\#\mathcal{I}_{0}\delta\log N.
\]
与\eqref{eq:4.7.10}合起来即得\eqref{eq:4.7.9}.

\textbf{估计\eqref{eq:4.7.10}.} 由于在估计中容许出现 $\log N$ 的项,我们充分利用这一点来估计\eqref{eq:4.7.10}.首先改写如下进行估计.
\begin{align*}
\eqref{eq:4.7.10} &= \sum_{i \in \mathcal{I}_{0}}\sum_{j \in \mathcal{I}_{0}}|T_{i} \cap T_{j}| \\
&\leq c\delta^{n-1}\#\mathcal{I}_{0} \\
&\quad + \sum_{i \in \mathcal{I}_{0}}\sum_{k=1}^{\log N}\sum\left\{j \in \mathcal{I}_{0}:\ \delta2^{k-1} \leq |\omega(T_{i}) - \omega(T_{j})| < \delta2^{k}\right\}|T_{i} \cap T_{j}| \\
&\leq c\delta^{n-1}\#\mathcal{I}_{0} \\
&\quad + \log N\sum_{i \in \mathcal{I}_{0}}\sum\left\{j \in \mathcal{I}_{0}:\ \sigma/2 \leq |\omega(T_{i}) - \omega(T_{j})| < \sigma\right\}|T_{i} \cap T_{j}|, \\
\intertext{其中 $\sigma := \delta2^{k_{0}}$,$k_{0}$ 是使级数 $\sum_{k}$ 的项取最大的 $k$ 的值. 如果}
\delta^{n-1}\#\mathcal{I}_{0} &\geq \sum_{i \in \mathcal{I}_{0}}\sum\left\{j \in \mathcal{I}_{0}:\ \sigma/2 \leq |\omega(T_{i}) - \omega(T_{j})| < \sigma\right\}|T_{i} \cap T_{j}| \\
\intertext{成立,则}
\eqref{eq:4.7.10} &\leq C(\log N)\delta^{n-1}\#\mathcal{I}_{0},
\end{align*}
经由\eqref{eq:4.7.10}得\eqref{eq:4.7.9}成立.因此,我们假定
\[
\eqref{eq:4.7.10} \leq C\log N\sum_{i \in \mathcal{I}_{0}}\sum\left\{j \in \mathcal{I}_{0}:\ \sigma/2 \leq |\omega(T_{i}) - \omega(T_{j})| < \sigma\right\}|T_{i} \cap T_{j}|
\]
来继续讨论.

像上面那样把待估计的量分配到 $O(\log N)$ 个部分,并关注其中最大的量进行估计的方法,此后将被基本地使用.

将 $I_{0}$ 分割为边长为 $\sigma$ 的立方体,把其中一个记为 $Q$,考虑 $Q \cap \mathcal{I}_{0}$.对每个 $Q \cap \mathcal{I}_{0}$,在满足 $i \in Q \cap \mathcal{I}_{0}$ 的 $i$ 中,各选取一个使
\[
\sum\left\{j \in \mathcal{I}_{0}:\ \sigma/2 \leq |\omega(T_{i}) - \omega(T_{j})| < \sigma\right\}|T_{i} \cap T_{j}|
\]
取最大的 $i$,把所选取的全体 $i$ 的集合记为 $I_{1}$.当 $i \in I_{1}$ 且 $i \in Q \cap \mathcal{I}_{0}$ 时,记
\begin{align*}
I_{i} &:= Q \cap \mathcal{I}_{0}. \\
\intertext{于是,}
I_{i} \cap I_{j} &= \emptyset \quad (i \neq j,\ i, j \in \mathcal{I}_{1}), \\
\mathcal{I}_{0} &= \bigcup_{i \in \mathcal{I}_{1}}I_{i}, \\
\intertext{并且}
\eqref{eq:4.7.10} &\leq C\log N\sum_{i \in \mathcal{I}_{1}}\#I_{i}\sum\left\{j \in \mathcal{I}_{0}:\ \sigma/2 \leq |\omega(T_{i}) - \omega(T_{j})| < \sigma\right\}|T_{i} \cap T_{j}|. \\
\intertext{对 $i \in I_{1}$,令}
\tilde{J}_{i} &:= \{j \in \mathcal{I}_{0}:\ \sigma/2 \leq |\omega(T_{i}) - \omega(T_{j})| < \sigma,\ T_{i} \cap T_{j} \neq \emptyset\}, \\
\intertext{则由以上和引理~\ref{lem:4.7.3} 知下式成立:}
\eqref{eq:4.7.10} &\leq C(\log N)\frac{\delta^{n}}{\sigma}\sum_{i \in \mathcal{I}_{1}}\#I_{i}\#\tilde{J}_{i}.
\end{align*}
对每个 $j \in \tilde{J}_{i}$,由简单的几何考察可知
\begin{align*}
T_{j} &\subset \{x:\ \operatorname{dist}(x, l(T_{i})) < \sigma\}. \\
\intertext{因此由\eqref{eq:4.7.8},令}
&\begin{cases}
f_{0}(y) := \chi\{\operatorname{dist}(x, l(T_{i})) < \delta\}(y)f(y), \\
f_{k}(y) := \chi\{\delta2^{k-1} \leq \operatorname{dist}(x, l(T_{i})) < \delta2^{k}\}(y)f(y),
\end{cases} \\
\intertext{则}
\lambda|T_{j}| &< \sum_{k=0}^{k_{0}}\int_{T_{j}}f_{k}(y)\dd y.
\end{align*}
利用此式静心思考,可以找到 $\delta < \rho \leq \sigma$ 与 $J_{i} \subset \tilde{J}_{i}$,使得
\begin{equation}\tag{4.7.11}
\frac{\lambda}{\log N}|T_{j}| < \int_{T_{j}}\chi\{\rho/2 \leq \operatorname{dist}(x, l(T_{i})) < \rho\}(y)f(y)\dd y \quad (j \in J_{i},\ i \in \mathcal{I}_{1})
\label{eq:4.7.11}\end{equation}
和
\[
\eqref{eq:4.7.10} \leq C(\log N)^{2}\frac{\delta^{n}}{\sigma}\sum_{i \in \mathcal{I}_{1}}\#I_{i}\#J_{i}
\]
同时成立.

这里给出一条之后会变得重要的注意.对\eqref{eq:4.7.11},利用 $\|f\|_{\infty} = 1$ 作简单的几何考察即可知
\begin{equation}\tag{4.7.12}
\frac{\sigma}{\rho} \leq C\frac{\log N}{\lambda}
\label{eq:4.7.12}\end{equation}
成立.

再增加一个 $\log N$ 的项,然后继续前进.

利用 $\#J_{i} \leq (2N)^{n-1}$,可以找到自然数 $\nu$ 和 $I_{2} \subset I_{1}$,使得
\begin{equation}\tag{4.7.13}
\nu \leq \#J_{i} < 2\nu \quad (i \in \mathcal{I}_{2})
\label{eq:4.7.13}\end{equation}
和
\begin{equation}\tag{4.7.14}
\eqref{eq:4.7.10} \leq C\nu(\log N)^{3}\frac{\delta^{n}}{\sigma}\sum_{i \in \mathcal{I}_{2}}\#I_{i}
\label{eq:4.7.14}\end{equation}
同时成立.

\textbf{辅助极大函数.} 设直线 $L \subset \mathbb{R}^{n}$ 的单位方向向量属于 $\tilde{D}_{0}$.对 $f \in L^{1}(Q_{1})$,在集合 $\mathbf{I}_{0}$ 上定义的辅助极大函数为
\[
K_{\delta}^{L}f(i) := \sup_{T \ni i,\ T \cap L \neq \emptyset}\frac{1}{|T|}\int_{T}f(y)\dd y \quad (i \in \mathbf{I}_{0}).
\]
这里,上确界对包含 $i$ 且与 $L$ 相交的 $\delta$-tube $T$ 取.

关于该辅助极大函数的如下引理,在以后的证明推进中起重要作用.

\begin{lemma}\label{lem:4.7.4}
设 $\delta < \rho \leq \sigma < 1$,且
\begin{align*}
\operatorname{supp}f &\subset \{x:\ \rho/2 \leq \operatorname{dist}(x, L) < \rho\}, \quad \mathcal{I} = \mathbf{I}_{0} \cap \{x:\ \operatorname{dist}(x, L) < \sigma\}, \\
\intertext{则}
\lambda^{2}\delta^{n-1}\#\left\{i \in \mathcal{I}:\ K_{\delta}^{L}f(i) > \lambda\right\}&\leq C(\log N)^{2}(\sigma/\rho)^{n-2}\int f(y)^{2}\dd y \quad (\lambda > 0)
 \end{align*}
成立.
\end{lemma}

\begin{proof}
证明容易,故仅记其概略.令
\begin{align*}
E_{1} &:= \{x:\ \rho/2 \leq \operatorname{dist}(x, L) < \rho\}, \quad E_{2} := \{x:\ \sigma/2 \leq \operatorname{dist}(x, L) < \sigma\}, \\
\intertext{考虑二进制分解,来估计最困难的部分}
I_{0} &:= \left\{i \in \mathbf{I}_{0} \cap E_{2}:\ K_{\delta}^{L}f(i) > \lambda\right\}
\end{align*}
的个数.

与\eqref{eq:4.7.10}相同,此时需要估计 $\left|T_{i} \cap T_{j} \cap E_{1}\right|$ 的和.固定 $T_{i}$ 后,由辅助极大函数的定义可知,与 $T_{i}$ 相交的 $T_{j}$ 本质上只限于包含在由 $L$ 与 $l(T_{i})$ 决定的二维平面的 $\delta$ 邻域中的那些.因此这里可以应用二维的结果.具有如此厚度的平面,随 $E_{1}$,$I_{0}$ 的位置关系,在 $E_{1}$ 上具有 $(\sigma/\rho)^{n-2}$ 的重数.利用这些即证得引理.
\end{proof}

\textbf{估计 $\nu$ 与 $\nu\#\mathcal{I}_{2}$.} 为了从\eqref{eq:4.7.14}中消去 $\nu$,我们估计 $\nu$ 以及 $\nu\#\mathcal{I}_{2}$.

对 $i \in I_{2}$,回忆\eqref{eq:4.7.11},对 $L = l(T_{i})$ 应用引理~\ref{lem:4.7.4},令
\begin{align*}
\overline{T_{i}} &:= \{x:\ \operatorname{dist}(x, l(T_{i})) < \rho\}, \\
\intertext{则可知}
\left(\frac{\lambda}{\log N}\right)^{2}\delta^{n-1}\#J_{i} &\leq C(\log N)^{2}(\sigma/\rho)^{n-2}\int_{\overline{T_{i}}}f(y)^{2}\dd y. \\
\intertext{于是,记}
\gamma_{1} &:= (\log N)^{4}(\sigma/\rho)^{n-2}\lambda^{-2}\delta^{1-n},
\end{align*}
则
\begin{equation}\tag{4.7.15}
\nu \leq \#J_{i} \leq C\gamma_{1}\int_{\overline{T_{i}}}f(y)^{2}\dd y.
\label{eq:4.7.15}\end{equation}
接下来估计
\begin{align*}
\nu\#\mathcal{I}_{2} &\leq C\gamma_{1}\sum_{i \in \mathcal{I}_{2}}\int_{\overline{T_{i}}}f(y)^{2}\dd y. \\
\intertext{以下是被称作 bush-argument 的手法.把}
\sum_{i \in \mathcal{I}_{2}}\int_{\overline{T_{i}}}f(y)^{2}\dd y &= \int f(y)^{2}\left(\sum_{i \in I_{2}}\chi_{\overline{T_{i}}}(y)\right)\dd y =: \int f(y)^{2}G(y)\dd y
\end{align*}
改写,设 $G(y)$ 在 $y = a$ 处取最大值 $\mu$,来估计 $G(y)$.

由于与 $y = a$ 相交的 $\overline{T_{i}}$ 的个数为 $\mu$,令
\[
\mathcal{I}_{3} := \{i \in \mathcal{I}_{2}:\ \overline{T_{i}} \ni a\},
\]
来考察 $\nu\#\mathcal{I}_{3}$.

对 $i \in \mathcal{I}_{3}$,由\eqref{eq:4.7.11}得
\begin{align*}
\frac{\lambda}{\log N}|T_{j}| &< \int_{T_{j}}\chi_{\overline{T_{i}}}(y)f(y)\dd y \quad (j \in J_{i}), \\
\intertext{故利用 $\|f\|_{\infty} = 1$,}
&\begin{aligned}
\frac{\lambda}{\log N}|T_{j}|& \\
&< \int_{T_{j}}\chi\left\{\overline{T_{i}} \setminus B\left(a, \frac{\lambda}{3\log N}\right)\right\}(y)f(y)\dd y + \left|T_{j} \cap B\left(a, \frac{\lambda}{3\log N}\right)\right|.
\end{aligned} \\
\intertext{因此可设}
\frac{\lambda}{3\log N}|T_{j}| &< \int_{T_{j}}\chi\left\{\overline{T_{i}} \setminus B\left(a, \frac{\lambda}{3\log N}\right)\right\}(y)f(y)\dd y \quad (j \in J_{i}). \\
\intertext{用此式再次应用引理~\ref{lem:4.7.4},得}
\nu &\leq \#J_{i} \leq C\gamma_{1}\int\chi\left\{\overline{T_{i}} \setminus B\left(a, \frac{\lambda}{3\log N}\right)\right\}(y)f(y)^{2}\dd y \quad (i \in \mathcal{I}_{3}). \\
\intertext{因此}
\nu\#\mathcal{I}_{3} &\leq C\gamma_{1}\sum_{i \in \mathcal{I}_{3}}\int\chi\left\{\overline{T_{i}} \setminus B\left(a, \frac{\lambda}{3\log N}\right)\right\}(y)f(y)^{2}\dd y.
\end{align*}
利用 $\mathcal{I}_{3}$ 至少分离 $\sigma$ 这一事实作简单的几何考察,右端的和中重数由
\begin{align*}
\gamma_{2} &:= \left(\frac{\rho}{(\lambda/\log N)\sigma}\right)^{n-1} \\
\intertext{的常数倍估计,得}
\nu\#\mathcal{I}_{3} &\leq C\gamma_{1}\gamma_{2}\int f(y)^{2}\dd y.
\end{align*}
因此
\begin{equation}\tag{4.7.16}
\nu\#\mathcal{I}_{2} \leq C\nu^{-1}\gamma_{1}^{2}\gamma_{2}\left(\int f(y)^{2}\dd y\right)^{2}.
\label{eq:4.7.16}\end{equation}

\textbf{\eqref{eq:4.7.9}的证明.} 首先假定
\[
\sigma \geq \left(\delta^{n-1}\#\mathcal{I}_{0}\right)^{1/(n-1)}.
\]
此时,由\eqref{eq:4.7.10},\eqref{eq:4.7.14},\eqref{eq:4.7.15}以及\eqref{eq:4.7.12}计算出
\begin{align*}
\sigma\lambda^{n+2}\delta^{n-2}\left(\delta^{n-1}\#\mathcal{I}_{0}\right) &\leq C(\log N)^{n+5}\left(\int f(y)^{2}\dd y\right)^{2}, \\
\intertext{利用假定得}
\lambda^{n+2}\delta^{n-2}\left(\delta^{n-1}\#\mathcal{I}_{0}\right)^{n/(n-1)} &\leq C(\log N)^{n+5}\left(\int f(y)^{2}\dd y\right)^{2}. \\
\intertext{取正平方根,则\eqref{eq:4.7.9}成立. 其次假定}
\sigma &< \left(\delta^{n-1}\#\mathcal{I}_{0}\right)^{1/(n-1)}.
\end{align*}
第一,由\eqref{eq:4.7.14}得
\begin{equation}\tag{4.7.17}
\eqref{eq:4.7.10} \leq C\nu(\log N)^{3}\frac{\delta^{n}}{\sigma}\#\mathcal{I}_{0}.
\label{eq:4.7.17}\end{equation}
第二,由 $\#I_{i} \leq \left(\frac{\sigma}{\delta}\right)^{n-1}$ 以及\eqref{eq:4.7.14},\eqref{eq:4.7.16}计算出
\[
\eqref{eq:4.7.10} \leq C\nu^{-1}\lambda^{-n-3}\delta^{-2n+3}(\log N)^{n+10}\sigma^{n-2}\left(\frac{\sigma}{\rho}\right)^{n-3}\left(\int f(y)^{2}\dd y\right)^{2},
\]
用\eqref{eq:4.7.12}得
\begin{equation}\tag{4.7.18}
\eqref{eq:4.7.10} \leq C\nu^{-1}\lambda^{-2n}\delta^{-2n+3}(\log N)^{2n+7}\sigma^{n-2}\left(\int f(y)^{2}\dd y\right)^{2}.
\label{eq:4.7.18}\end{equation}
把\eqref{eq:4.7.17},\eqref{eq:4.7.18}的两边相乘并取正平方根,则
\begin{align*}
&\begin{aligned}
\eqref{eq:4.7.10}&\leq C\lambda^{-n}\delta^{-(n-3)/2}(\log N)^{n+5}\sigma^{(n-3)/2}(\#I_{0})^{1/2}\int f(y)^{2}\dd y  \\
&= C\lambda^{-n}\delta^{-n+2}(\log N)^{n+5}\sigma^{(n-3)/2}\left(\delta^{n-1}\#I_{0}\right)^{1/2}\int f(y)^{2}\dd y.
\end{aligned} \\
\intertext{由\eqref{eq:4.7.10}得}
\lambda^{n+2}\delta^{n-2}\left(\delta^{n-1}\#\mathcal{I}_{0}\right)^{3/2} &\leq C(\log N)^{n+5}\sigma^{(n-3)/2}\left(\int f(y)^{2}\dd y\right)^{2}. \\
\intertext{利用假定,此情形也得}
\lambda^{n+2}\delta^{n-2}\left(\delta^{n-1}\#\mathcal{I}_{0}\right)^{n/(n-1)} &\leq C(\log N)^{n+5}\left(\int f(y)^{2}\dd y\right)^{2}.
\end{align*}
取正平方根,则\eqref{eq:4.7.9}成立.

由以上,命题~\ref{prop:4.7.2} 得证.

\subsection{定理~\ref{thm:4.7.1} 的证明}\label{sec:4.7.3}

定理~\ref{thm:4.7.1} 可由在命题~\ref{prop:4.7.2} 中取 $f = \chi_{E}$($E \subset Q_{1}$)所得的结果与\eqref{eq:4.7.4},用 Marcinkiewicz 插值定理证明.这里,虽然计算稍复杂,但我们用二进制分解给出直接证明.用被称为三点插值法的方法来证明\eqref{eq:4.7.6}.由此定理显然成立.

设 $\lambda > 0$,$\epsilon \gg 1$ 在所需的计算之后确定,分解
\begin{align*}
f(x) & = f(x)\chi_{\{f < \lambda\}}(x) + f(x)\chi_{\{\lambda \leq f < \epsilon\lambda\}}(x) + f(x)\chi_{\{\epsilon\lambda \leq f\}}(x) \\
& =: f_{1}(x) + f_{2}(x) + f_{3}(x), \\
\intertext{则由次线性性得}
K_{\delta}f(x) &\leq K_{\delta}f_{1}(x) + K_{\delta}f_{2}(x) + K_{\delta}f_{3}(x).
\end{align*}
显然 $K_{\delta}f_{1}(x) \leq \lambda$,故由上式得
\begin{align*}
&\begin{aligned}
\{x \in Q_{2}:\ K_{\delta}f(x) > 3\lambda\}&\subset \{x \in Q_{2}:\ K_{\delta}f_{2}(x) > \lambda\} \cup \{x \in Q_{2}:\ K_{\delta}f_{3}(x) > \lambda\}.
\end{aligned}  \\
\intertext{因此}
&\begin{aligned}
|\{x \in Q_{2}:\ K_{\delta}f(x) > 3\lambda\}|&\leq |\{x \in Q_{2}:\ K_{\delta}f_{2}(x) > \lambda\}| + |\{x \in Q_{2}:\ K_{\delta}f_{3}(x) > \lambda\}|.
\end{aligned}  \\
\intertext{利用此式,}
&\begin{aligned}
\int_{Q_{2}}K_{\delta}f(x)^{q}\dd x &= 3^{q}q\int_{0}^{\infty}|\{x \in Q_{2}:\ K_{\delta}f(x) > 3\lambda\}|\lambda^{q-1}\dd\lambda \\
\qquad &\leq 3^{q}q\int_{0}^{\infty}|\{x \in Q_{2}:\ K_{\delta}f_{2}(x) > \lambda\}|\lambda^{q-1}\dd\lambda
\end{aligned} \\
&\begin{aligned}
+ 3^{q}q\int_{0}^{\infty}|\{x \in Q_{2}:\ K_{\delta}f_{3}(x) > \lambda\}|\lambda^{q-1}\dd\lambda&=: 3^{q}q(\mathrm{I} + \mathrm{II}).
\end{aligned}
 \end{align*}

\textbf{估计 I.} 首先用命题~\ref{prop:4.7.2} 证明
\begin{equation}\tag{4.7.19}
|\{x \in Q_{2}:\ K_{\delta}f_{2}(x) > \lambda\}|^{p/q} \leq C N^{n-p}(\log N)^{\alpha}(\log\epsilon)^{p}\lambda^{-p}\int f_{2}(y)^{p}\dd y.
\label{eq:4.7.19}\end{equation}

令
\begin{align*}
g_{k} &:= f(x)\chi\{2^{k-1}\lambda \leq f < 2^{k}\lambda\}(x) \quad (k > 0), \\
\intertext{再分解}
f_{2}(x) &= \sum_{k=1}^{\log\epsilon}g_{k}(x), \\
\intertext{则}
|\{x:\ K_{\delta}f_{2}(x) > \lambda\}| &\leq \sum_{k=1}^{\log\epsilon}\left|\left\{x:\ K_{\delta}g_{k}(x) > \frac{\lambda}{\log\epsilon}\right\}\right|. \\
\intertext{改写为}
\left|\left\{x:\ K_{\delta}g_{k}(x) > \frac{\lambda}{\log\epsilon}\right\}\right| &= \left|\left\{x:\ K_{\delta}\left(\frac{g_{k}}{2^{k}\lambda}\right)(x) > \frac{1}{2^{k}\log\epsilon}\right\}\right|, \\
\intertext{即可对该式右端应用命题~\ref{prop:4.7.2}:}
&\begin{aligned}
\left|\left\{x:\ K_{\delta}\left(\frac{g_{k}}{2^{k}\lambda}\right)(x) > \frac{1}{2^{k}\log\epsilon}\right\}\right|^{p/q}&\leq C N^{n-p}(\log N)^{\alpha}(2^{k}\log\epsilon)^{p}\int\left(\frac{g_{k}(y)}{2^{k}\lambda}\right)^{2}\dd y  \\
&\leq C N^{n-p}(\log N)^{\alpha}(\log\epsilon)^{p}\lambda^{-p}\int g_{k}(y)^{p}\dd y.
\end{aligned}
\end{align*}
这里用到了 $(p-1)/(n-1) = p/q$ 与 $\frac{1}{2} \leq \frac{g_{k}(x)}{2^{k}\lambda}$.注意到 $p/q < 1$,由此得
\begin{align*}
&\begin{aligned}
|\{x \in Q_{2}:\ K_{\delta}f_{2}(x) > \lambda\}|^{p/q}&\leq \sum_{k=1}^{\log\epsilon}\left|\left\{x:\ K_{\delta}g_{k}(x) > \frac{\lambda}{\log\epsilon}\right\}\right|^{p/q}  \\
&\leq C N^{n-p}(\log N)^{\alpha}(\log\epsilon)^{p}\lambda^{-p}\int\sum_{k=1}^{\log\epsilon}g_{k}(y)^{p}\dd y \\
&= C N^{n-p}(\log N)^{\alpha}(\log\epsilon)^{p}\lambda^{-p}\int f_{2}(y)^{p}\dd y.
\end{aligned} \\
\intertext{即知\eqref{eq:4.7.19}成立. 用此来估计 I:}
&\begin{aligned}
\mathrm{I} &= \int_{0}^{\infty}|\{x \in Q_{2}:\ K_{\delta}f_{2}(x) > \lambda\}|\lambda^{q-1}\dd\lambda \\
&\leq C N^{qn/p-q}(\log N)^{\alpha q/p}(\log\epsilon)^{q}\int_{0}^{\infty}\left(\int_{Q_{1}}f(y)^{p}\chi_{\{\lambda \leq f < \epsilon\lambda\}}(y)\dd y\right)^{q/p}\frac{\dd\lambda}{\lambda}.
\end{aligned} \\
\intertext{利用 $p/q < 1$,再次用二进制分解,}
&\begin{aligned}
\left(\int_{0}^{\infty}\left(\int f(y)^{p}\chi_{\{\lambda \leq f < \epsilon\lambda\}}(y)\dd y\right)^{q/p}\frac{\dd\lambda}{\lambda}\right)^{p/q}&= \left(\sum_{k \in \mathbf{Z}}\int_{2^{k-1}}^{2^{k}}\left(\int f(y)^{p}\chi_{\{\lambda \leq f < \epsilon\lambda\}}(y)\dd y\right)^{q/p}\frac{\dd\lambda}{\lambda}\right)^{p/q}  \\
&\leq \left(\sum_{k \in \mathbf{Z}}\left(\int f(y)^{p}\chi_{\{2^{k-1} \leq f < \epsilon2^{k}\}}(y)\dd y\right)^{q/p}\int_{2^{k-1}}^{2^{k}}\frac{\dd\lambda}{\lambda}\right)^{p/q} \\
&\leq C\sum_{k \in \mathbf{Z}}\int f(y)^{p}\chi_{\{2^{k-1} \leq f < \epsilon2^{k}\}}(y)\dd y \\
&\leq C\int_{0}^{\infty}\int f(y)^{p}\chi_{\{\lambda \leq f < 8\epsilon\lambda\}}(y)\dd y\frac{\dd\lambda}{\lambda}.
\end{aligned} \\
\intertext{用此得}
\mathrm{I} &\leq C N^{qn/p-q}(\log N)^{\alpha q/p}(\log\epsilon)^{q}\left(\int_{0}^{\infty}\int f(y)^{p}\chi_{\{\lambda \leq f < 8\epsilon\lambda\}}(y)\dd y\frac{\dd\lambda}{\lambda}\right)^{q/p} \\
&\begin{aligned}
&\leq C N^{qn/p-q}(\log N)^{\alpha q/p}(\log\epsilon)^{q}\left(\int f(y)^{p}\int_{f(y)/8\epsilon}^{f(y)}\frac{\dd\lambda}{\lambda}\dd y\right)^{q/p} \\
&= C N^{qn/p-q}(\log N)^{\alpha q/p}(\log\epsilon)^{q+q/p}\left(\int f(y)^{p}\dd y\right)^{q/p}.
\end{aligned}
\end{align*}

\textbf{估计 II.} 由简单的几何考察可知
\begin{align*}
K_{\delta}f(x) &\leq C N^{n-1}\int f(y)\chi_{B(x,1)}(y)\dd y. \\
\intertext{由 Young 不等式,由此得}
\|K_{\delta}f\|_{L^{1}(Q_{2})} &\leq C N^{n-1}\|f\|_{L^{1}(Q_{1})}. \\
\intertext{特别地,关于分布函数有}
|\{x \in Q_{2}:\ K_{\delta}f(x) > \lambda\}| &\leq C\frac{N^{n-1}}{\lambda}\|f\|_{L^{1}(Q_{1})} \quad (\lambda > 0).
\end{align*}
又,若 $f \in L^{p}(Q_{1})$,则 Hölder 不等式给出
\begin{align*}
K_{\delta}f(x) &\leq C N^{(n-1)/p}\|f\|_{L^{p}(Q_{1})}. \\
\intertext{利用这些,}
&\begin{aligned}
\mathrm{II} &= \int_{0}^{\infty}|\{x \in Q_{2}:\ K_{\delta}f_{3}(x) > \lambda\}|\lambda^{q-1}\dd\lambda \\
&\leq C\left(N^{(n-1)/p}\|f\|_{p}\right)^{q-p}\int_{0}^{\infty}|\{x:\ K_{\delta}f_{3}(x) > \lambda\}|\lambda^{p-1}\dd\lambda \\
&\leq C\left(N^{(n-1)/p}\|f\|_{p}\right)^{q-p}N^{n-1}\int_{0}^{\infty}\int f(y)\chi_{\{\epsilon\lambda \leq f\}}(y)\dd y\,\lambda^{p-2}\dd\lambda \\
&\leq C N^{q(n-1)/p}\|f\|_{p}^{q-p}\int f(y)\int_{0}^{f(y)/\epsilon}\lambda^{p-2}\dd\lambda\dd y \\
&\leq C N^{q(n-1)/p}\|f\|_{p}^{q-p}\epsilon^{1-p}\int f(y)^{p}\dd y \\
&= C N^{q(n-1)/p}\epsilon^{1-p}\left(\int f(y)^{p}\dd y\right)^{q/p}.
\end{aligned}
\end{align*}
若选取 $\epsilon = N^{(q-q/p)/(p-1)}$,则合 I 与 II 的估计得
\[
\|K_{\delta}f\|_{L^{q}(Q_{2})} \leq C N^{n/p-1}(\log N)^{(\alpha+1)/p+1}\|f\|_{L^{p}(Q_{1})},
\]
于是\eqref{eq:4.7.6}成立,定理得证.

\subsection{大 Nikodym 极大函数(C\'ordoba 的筛)}\label{sec:4.7.4}

这里,我们想利用定理~\ref{thm:4.7.1} 来考虑对由\eqref{eq:4.7.1},\eqref{eq:4.7.2}定义的大 Nikodym 极大函数的估计.C\'ordoba 的筛论证是其核心.

首先,由\eqref{eq:4.7.6}与 Hölder 不等式可知
\[
\|K_{\delta}f\|_{L^{p}(Q_{2})} \leq C N^{n/p-1}(\log N)^{\alpha}\|f\|_{L^{p}(Q_{1})} \quad (p = (n+2)/2).
\]
回忆 $N = 1/\delta$,将坐标按 $1 \to Na$ 作尺度变换,则由上式得
\begin{align*}
\|\mathcal{K}_{N,a}f\|_{L^{p}(Q_{2Na})} &\leq C N^{n/p-1}(\log N)^{\alpha}\|f\|_{L^{p}(Q_{Na})}, \\
Q_{Na} &:= (-Na, Na)^{n}, \quad Q_{2Na} := (-2Na, 2Na)^{n}
\end{align*}
成立.由此,通过适当分解 $f$ 的支集,得
\begin{equation}\tag{4.7.20}
\|\mathcal{K}_{N,a}f\|_{L^{p}(\mathbb{R}^{n})} \leq C N^{n/p-1}(\log N)^{\alpha}\|f\|_{L^{p}(\mathbb{R}^{n})} \quad (p = (n+2)/2).
\label{eq:4.7.20}\end{equation}
特别地,对一切 $\lambda > 0$,
\[
\begin{aligned}
|\{x \in \mathbb{R}^{n}:\ \mathcal{K}_{N,a}f(x) > \lambda\}| &\leq C N^{n-p}(\log N)^{p\alpha}\lambda^{-p}\|f\|_{p}^{p} \\
&(p = (n+2)/2)
\end{aligned}
\]

下面给出筛引理.

\begin{lemma}\label{lem:4.7.5}
设 $1 \leq p < \infty$.此时,对依赖于 $N$ 的常数 $C_{N}$,若
\begin{align*}
|\{x \in \mathbb{R}^{n}:\ \mathcal{K}_{N,a}f(x) > \lambda\}| &\leq C_{N}\lambda^{-p}\|f\|_{p}^{p} \quad (\lambda > 0) \\
\intertext{成立,则}
|\{x \in \mathbb{R}^{n}:\ \mathcal{K}_{N}f(x) > \lambda\}| &\leq C C_{N}(\log N)\lambda^{-p}\|f\|_{p}^{p} \quad (\lambda > 0)
\end{align*}
成立.其中 $C$ 是仅依赖于维数 $n$ 与 $p$ 的常数.
\end{lemma}

由该引理与\eqref{eq:4.7.20}可确认
\begin{align*}
&\begin{aligned}
|\{x \in \mathbb{R}^{n}:\ \mathcal{K}_{N}f(x) > \lambda\}| &\leq C N^{n-p}(\log N)^{p\alpha+1}\lambda^{-p}\|f\|_{p}^{p} \\
&(p = (n+2)/2)
\end{aligned} \\
\intertext{成立.另一方面,}
|\{x \in \mathbb{R}^{n}:\ \mathcal{K}_{N}f(x) > \lambda\}| &\leq C N^{n-1}\lambda^{-1}\|f\|_{1}
\end{align*}
的成立是显然的,因此利用这两式,通过与定理~\ref{thm:4.7.1} 的证明类似的三点插值法的讨论,可证明如下定理.

\begin{theorem}\label{thm:4.7.6}
在 $1 < p \leq (n+2)/2$ 时
\begin{equation}\tag{4.7.21}
\|\mathcal{K}_{N}f\|_{L^{p}(\mathbb{R}^{n})} \leq C N^{n/p-1}(\log N)^{\alpha}\|f\|_{L^{p}(\mathbb{R}^{n})}
\label{eq:4.7.21}\end{equation}
成立.这里 $\alpha$ 表示依赖于维数 $n$ 与 $p$ 的常数.
\end{theorem}

\textbf{引理~\ref{lem:4.7.5} 的证明.}

\begin{proof}
采用仅在此处使用的记号:与 tube
\[
\{(x',x_n) \in \mathbb{R}^{n}:\ |x'| < a/2,\ |x_n| < b/2\}
\]
合同的 tube 的维数记为 $(a,b)$.另外,对 tube $T$,用 $cT$ 表示保持 $T$ 的中心与轴所在直线不变,将各边的长度放大 $c$ 倍所得的 tube.

设 $f \in L^{p}(\mathbb{R}^{n})$,令
\[
E_{\lambda} := \{x \in \mathbb{R}^{n}:\ \mathcal{K}_{N}f(x) > \lambda\}.
\]
此时,设 $A \subset E_{\lambda}$ 为任意紧集,对每个 $x \in A$,由极大函数的定义,选取包含 $x$,维数为 $(a_{x}, Na_{x})$ 的 tube $T_{x}$,使得
\begin{equation}\tag{4.7.22}
\lambda < \frac{1}{|T_{x}|}\int_{T_{x}}f(y)\dd y
\label{eq:4.7.22}\end{equation}
成立.由 $A$ 的紧性,得到满足
\[
\mathcal{A} \subset \bigcup_{j=1}^{\nu}T_{x_{j}}
\]
的有限个 tube 的族.把 $T_{x_{j}}$ 记为 $T_{j}$,设 $T_{j}$ 的维数为 $(a_{j}, Na_{j})$.

选取满足
\begin{align*}
N^{\kappa} &\leq \max\{a_{j}:\ j = 1, \dots, \nu\} < N^{\kappa+1} \\
\intertext{的 $\kappa \in \mathbf{Z}$.首先令}
I_{0} &:= \{j:\ N^{\kappa} \leq a_{j} < N^{\kappa+1}\}, \\
\intertext{并令}
J_{0} &:= \{j:\ N^{\kappa-1} \leq a_{j} < N^{\kappa}\}, \\
\intertext{设}
\mathcal{A}_{0} &:= \bigcup_{j \in I_{0}}T_{j}, \quad \mathcal{B}_{0} := \bigcup_{j \in J_{0}}T_{j}. \\
\intertext{其次令}
I_{1} &:= \{j:\ N^{\kappa-2} \leq a_{j} < N^{\kappa-1},\ T_{j} \cap \mathcal{A}_{0} = \emptyset\}, \\
\intertext{并令}
J_{1} &:= \{j:\ N^{\kappa-3} \leq a_{j} < N^{\kappa-2},\ T_{j} \cap \mathcal{B}_{0} = \emptyset\}, \\
\intertext{设}
\mathcal{A}_{1} &:= \bigcup_{j \in I_{1}}T_{j}, \quad \mathcal{B}_{1} := \bigcup_{j \in J_{1}}T_{j}. \\
\intertext{以下归纳地令}
I_{k} &:= \{j:\ N^{\kappa-2k} \leq a_{j} < N^{\kappa-(2k-1)},\ T_{j} \cap \mathcal{A}_{l} = \emptyset\ (l = 0, \dots, k-1)\}, \\
\intertext{并令}
J_{k} &:= \{j:\ N^{\kappa-(2k+1)} \leq a_{j} < N^{\kappa-2k},\ T_{j} \cap \mathcal{B}_{l} = \emptyset\ (l = 0, \dots, k-1)\}, \\
\intertext{设}
\mathcal{A}_{k} &:= \bigcup_{j \in I_{k}}T_{j}, \quad \mathcal{B}_{k} := \bigcup_{j \in J_{k}}T_{j}.
\end{align*}
该操作当然在有限次($K$ 次)后终止.

由构造方法可知,$\mathcal{A}_{0}, \ldots, \mathcal{A}_{K}$ 以及 $\mathcal{B}_{0}, \ldots, \mathcal{B}_{K}$ 各自互不相交.

任取 $T_{j}$.若
\[
N^{\kappa-2k} \leq a_{j} < N^{\kappa-(2k-1)}, \quad j \not\in I_{k},
\]
则存在某个 $j' \in I_{k'}$($k' = 0, \ldots, k-1$)使得
\begin{align*}
T_{j} \cap T_{j'} &\neq \emptyset. \\
\intertext{由此考虑 tube 的维数,得}
T_{j} &\subset 3T_{j'}. \\
\intertext{另一方面,若}
N^{\kappa-(2l+1)} &\leq a_{j} < N^{\kappa-2l}, \quad j \not\in J_{l},
\end{align*}
则存在某个 $j' \in J_{l'}$($l' = 0, \ldots, l-1$)使得
\begin{align*}
T_{j} \cap T_{j'} &\neq \emptyset, \\
\intertext{故同理}
T_{j} &\subset 3T_{j'}. \\
\intertext{因此}
\mathcal{A} &\subset \bigcup_{k=0}^{K}\bigcup_{j \in I_{k} \cup J_{k}}3T_{j}. \\
\intertext{由此得}
|\mathcal{A}| &\leq \sum_{k=0}^{K}\left|\bigcup_{j \in I_{k} \cup J_{k}}3T_{j}\right|. \\
\intertext{把该右端分成 $O(\log N)$ 个元素,来应用假设.分解}
&\begin{aligned}
I_{k} &= \bigcup_{m=0}^{\log N}\left\{j \in I_{k}:\ 2^{m}N^{\kappa-2k} \leq a_{j} < 2^{m+1}N^{\kappa-2k}\right\}, \\
J_{k} &= \bigcup_{m=0}^{\log N}\left\{j \in J_{k}:\ 2^{m}N^{\kappa-(2k+1)} \leq a_{j} < 2^{m+1}N^{\kappa-(2k+1)}\right\},
\end{aligned} \\
\intertext{则存在 $m_{0}$,使得}
|\mathcal{A}| &\leq \log N\sum_{k=0}^{K}\left|\bigcup_{j \in I'_{k} \cup J'_{k}}3T_{j}\right|, \\
&\begin{aligned}
I'_{k} &:= \{j \in I_{k}:\ 2^{m_{0}}N^{\kappa-2k} \leq a_{j} < 2^{m_{0}+1}N^{\kappa-2k}\}, \\
J'_{k} &:= \{j \in J_{k}:\ 2^{m_{0}}N^{\kappa-(2k+1)} \leq a_{j} < 2^{m_{0}+1}N^{\kappa-(2k+1)}\}.
\end{aligned} \\
\intertext{设}
\mathcal{A}'_{k} &:= \bigcup_{j \in I'_{k}}T_{j}, \quad \mathcal{B}'_{k} := \bigcup_{j \in J'_{k}}T_{j},
\end{align*}
则 $\mathcal{A}'_{0}, \ldots, \mathcal{A}'_{K}$ 以及 $\mathcal{B}'_{0}, \ldots, \mathcal{B}'_{K}$ 各自互不相交.因此,把它们合在一起的重数不超过 2.

把引理的假设用于
\begin{align*}
&\begin{aligned}
g_{k}(x) &:= f(x)\chi_{\mathcal{A}'_{k}}(x), \quad a = 2^{m_{0}}N^{\kappa-2k}, \\
h_{k}(x) &:= f(x)\chi_{\mathcal{B}'_{k}}(x), \quad a = 2^{m_{0}}N^{\kappa-(2k+1)},
\end{aligned} \\
\intertext{则由\eqref{eq:4.7.22}与假设得}
\left|\bigcup_{j \in I'_{k}}3T_{j}\right| &\leq C C_{N}\lambda^{-p}\|g_{k}\|_{p}^{p}, \\
\left|\bigcup_{j \in J'_{k}}3T_{j}\right| &\leq C C_{N}\lambda^{-p}\|h_{k}\|_{p}^{p}. \\
\intertext{综合以上,得}
|\mathcal{A}| &\leq C C_{N}(\log N)\lambda^{-p}\|f\|_{p}^{p}.
\end{align*}
由于 $A \subset E_{\lambda}$ 是任意的,引理得证.
\end{proof}

% ============ 第4章 §4.8 Kakeya集的几何维数(原书 pp.367-371) ============
\section{Kakeya 集的几何维数}\label{sec:4.8}
\index{Kakeya 集}
\index{Besicovitch 集}
\index{Hausdorff 外测度}
\index{Hausdorff 上维数}
\index{Minkowski 外测度}
\index{Minkowski 上维数}
\index{tube!tube 的维数}


本节作为前一节定理~\ref{thm:4.7.1} 的一个应用,讨论 Kakeya 集的几何维数的估计.首先定义 Kakeya 集.

\begin{definition}\label{def:4.8.1}
$\mathbb{R}^{n}$($n > 1$)中 $n$ 维 Lebesgue 测度为 0 的紧致子集,若其内部包含一切方向的单位线段,则称之为 Kakeya 集或 Besicovitch 集.
\end{definition}

关于这种令人着迷的集合的构造方法,请参考其他书籍.这里,我们只考虑该集合的几何维数的估计.

来定义几何维数.

\begin{definition}\label{def:4.8.2}
对 $E \subset \mathbb{R}^{n}$,其 Hausdorff 外测度定义为
\begin{equation*}
\mathcal{H}_{s}(E) := \sup_{\delta > 0}\inf\left\{\sum_{j}t_{j}^{s}:\ E \subset \bigcup_{j}B(x_{j},t_{j}),\ t_{j} < \delta\right\} \quad (s > 0).
\end{equation*}
\end{definition}

对有界的无限集合 $E \subset \mathbb{R}^{n}$,下式成立:
\[
\mathcal{H}_{n+\epsilon}(E) = 0 \quad (\epsilon > 0), \qquad \mathcal{H}_{0}(E) = \infty.
\]

\begin{definition}\label{def:4.8.3}
对 $E \subset \mathbb{R}^{n}$,其上 Hausdorff 维数定义为
\begin{equation*}
\operatorname{H-dim}(E) := \inf\{s:\ \mathcal{H}_{s}(E) = 0\}.
\end{equation*}
\end{definition}

\begin{definition}\label{def:4.8.4}
对 $E \subset \mathbb{R}^{n}$,其 Minkowski 外测度定义为
\begin{equation*}
\mathcal{M}_{s}(E) := \limsup_{\delta \to 0}\inf\left\{\sum_{j}\delta^{s}:\ E \subset \bigcup_{j}B(x_{j},\delta)\right\} \quad (s > 0).
\end{equation*}
\end{definition}

由定义得
\begin{equation}\tag{4.8.1}
\mathcal{M}_{s}(E) \geq \mathcal{H}_{s}(E).
\label{eq:4.8.1}\end{equation}

\begin{definition}\label{def:4.8.5}
对 $E \subset \mathbb{R}^{n}$,其上 Minkowski 维数定义为
\begin{equation*}
\operatorname{M-dim}(E) := \inf\{s:\ \mathcal{M}_{s}(E) = 0\}.
\end{equation*}
\end{definition}

对 $1 \leq p \leq n$,把"$\mathbb{R}^{n}$ 上所有 Kakeya 集具有至少 $p$ 的 Hausdorff 维数"这一主张记为 $K_{H}(p,n)$,把"$\mathbb{R}^{n}$ 上所有 Kakeya 集具有至少 $p$ 的 Minkowski 维数"这一主张记为 $K_{M}(p,n)$.此时,如下猜想成立.

\textbf{猜想问题.} $K_{H}(n,n)$ 以及 $K_{M}(n,n)$ 成立.

由\eqref{eq:4.8.1}可知 $K_{H}(p,n) \Rightarrow K_{M}(p,n)$ 的关系.

下面说明上一节讨论的小 Nikodym 极大函数 $K_{\delta}f$ 的猜想问题与上面的猜想问题之间的联系.

对 $1 < p \leq n$,把"上一节的\eqref{eq:4.7.5}在 $p$ 时成立"这一主张记为 $K_{N}(p,n)$.重复上一节的主张可知 $K_{N}(n,n)$ 被猜想.特别地,由定理~\ref{thm:4.7.1} 知 $K_{N}((n+2)/2, n)$ 成立.

\begin{lemma}\label{lem:4.8.6}
设 $1 < p \leq n$,则
\[
K_{N}(p,n) \Rightarrow K_{M}(p,n), \quad K_{N}(p,n) \Rightarrow K_{H}(p,n)
\]
的关系成立.
\end{lemma}

\textbf{$K_{N}(p,n) \Rightarrow K_{M}(p,n)$.} 对 Kakeya 集 $E \subset \mathbb{R}^{n}$,当 $s < p$ 时证明 $\mathcal{M}_{s}(E) > 0$.

设 $0 < \delta \ll 1$ 为小参数,把 $E$ 的 $\delta$ 邻域记为 $\mathcal{N}_{\delta}(E)$.此时,若把 $\mathcal{N}_{\delta}(E)$ 的特征函数记为 $f$,则由简单的几何考察可知
\begin{align*}
\|K_{\delta}f\|_{q} &> c > 0. \\
\intertext{用此式应用假设,则对某个 $\alpha$,}
\delta^{n-p}(\log(1/\delta))^{-\alpha} &\leq C|\mathcal{N}_{\delta}(E)|.
\end{align*}
因此,覆盖 $E$ 的球 $B(x_{j}, \delta)$ 的个数可由下界 $C^{-1}\delta^{\epsilon-p}$($\epsilon > 0$)估计,当 $s < p$ 时
\[
\mathcal{M}_{s}(E) \geq \limsup_{\delta \to 0}C^{-1}\delta^{-(p-s)/2} = \infty.
\]

\textbf{$K_{N}(p,n) \Rightarrow K_{H}(p,n)$.} 这次对 Kakeya 集 $E$,当 $s < p$ 时证明 $\mathcal{H}_{s}(E) > 0$.设 $E \subset \bigcup_{j}B(x_{j},t_{j})$,$t_{j} < 1$,令
\[
E_{k} := \bigcup_{j:\ 2^{-k} \leq t_{j} < 2^{1-k}}B(x_{j},t_{j}) \quad (k = 1, 2, \dots).
\]
由 Kakeya 集的性质,对任意 $\omega \in S^{n-1}$,存在包含于 $E$ 的单位线段 $l_{\omega}$.把一维 Hausdorff 测度记为 $|\cdot|$,则
\[
1 = |l_{\omega} \cap E| \leq \sum_{k}|l_{\omega} \cap E_{k}| = \sum_{k}\frac{|l_{\omega} \cap E_{k}|k^{2}}{k^{2}}.
\]
因此,用 $\sum_{k}\frac{1}{k^{2}} = c$,可选取正整数 $k_{\omega}$ 使得
\begin{align*}
|l_{\omega} \cap E_{k_{\omega}}| &> \frac{1}{ck_{\omega}^{2}}. \\
\intertext{同理可选取正整数 $\kappa$,使得}
\left|\left\{\omega \in S^{n-1}:\ |l_{\omega} \cap E_{\kappa}| > \frac{1}{c\kappa^{2}}\right\}\right| &> |S^{n-1}|\frac{1}{c\kappa^{2}}.
\end{align*}
其中左端外侧的 $|\cdot|$ 与右端的 $|\cdot|$ 是 $S^{n-1}$ 的 $n-1$ 维面积测度.

把 $E_{\kappa}$ 的特征函数记为 $f$,取 $\delta = 2^{-\kappa}$,由简单的几何考察,由上式得
\begin{align*}
\|K_{\delta}f\|_{q} &> C(1/\kappa^{2})^{1/q+1}. \\
\intertext{应用假设,则对某个 $\alpha$,}
(2^{-\kappa})^{n-p}\kappa^{-\alpha}(1/\kappa^{2})^{p/q+p} &\leq C|E_{\kappa}|.
\end{align*}
因此,$B(x_{j},t_{j})$($2^{-\kappa} \leq t_{j} < 2^{1-\kappa}$)的个数可由下界 $C^{-1}2^{p\kappa}\kappa^{-\alpha}(1/\kappa^{2})^{p/q+p}$ 估计,
\[
\sum_{j}t_{j}^{s} \geq C^{-1}2^{\kappa(p-s)}\kappa^{-\alpha}(1/\kappa^{2})^{p/q+p} > a > 0.
\]
其中 $a$ 是不依赖于 $\kappa$ 的数.因此 $\mathcal{H}_{s}(E) \geq a > 0$.

由该引理与定理~\ref{thm:4.7.1},如下定理成立.

\begin{theorem}\label{thm:4.8.7}
$\mathbb{R}^{n}$ 的所有 Kakeya 集至少具有 $(n+2)/2$ 的上 Hausdorff 维数以及上 Minkowski 维数.
\end{theorem}

% ============ 第4章 §4.9 Bochner-Riesz平均与Nikodym极大函数(原书 pp.371-378) ============
\section{Bochner--Riesz 平均与 Nikodym 极大函数}\label{sec:4.9}
\index{Bochner--Riesz 平均}
\index{Fefferman 定理}
\index{乘子!乘子算子(Rn 上的函数的)}


本节利用 Nikodym 极大函数的估计,讨论与 Fourier 反演公式相关的 Bochner--Riesz 平均问题.

正如在 4.3 节所确认的,对 $f \in \mathcal{S}(\mathbb{R}^{n})$,Fourier 反演公式
\[
\int_{\mathbb{R}^{n}}e^{2\pi i x\cdot\xi}\hat{f}(\xi)\dd\xi = f(x)
\]
成立.若 $f$ 是 $L^{p}(\mathbb{R}^{n})$ 中更一般的函数,则超函数意义下的 Fourier 反演公式成立,但问题在于了解更定量的收敛.于是,我们考察部分和
\[
S_{R}f(x) := \int_{|\xi| \leq R}e^{2\pi i x\cdot\xi}\hat{f}(\xi)\dd\xi \quad (R > 0)
\]
当 $R \to \infty$ 时在 $L^{p}(\mathbb{R}^{n})$ 中的收敛问题.

由一致有界性原理与扩张变换的不变性,该问题归结为取 $R = 1$ 时如下不等式是否成立:
\begin{equation}\tag{4.9.1}
\|S_{1}f\|_{L^{p}(\mathbb{R}^{n})} \leq C\|f\|_{L^{p}(\mathbb{R}^{n})}.
\label{eq:4.9.1}\end{equation}

\eqref{eq:4.9.1}在一维情形,由 Hilbert 变换的 $L^{p}(\mathbb{R})$ 有界性可知在 $1 < p < \infty$ 时成立.但是,在二维以上,只知道由 Plancherel 定理得到的 $p = 2$ 时成立.$S_{1}$ 的核的绝对值在远方以 $|x|^{-(n+1)/2}$ 衰减,因此考虑到对偶性,曾期望在
\[
p \in \left(\frac{2n}{n+1}, \frac{2n}{n-1}\right)
\]
时有界.然而,1971 年,C. Fefferman 证明了如下著名定理,表明该猜想不成立.

\begin{theorem}\label{thm:4.9.1}
在 $n > 1$,$p \neq 2$ 时\eqref{eq:4.9.1}不成立.
\end{theorem}

C. Fefferman 的定理是利用二维 Kakeya 集的 Minkowski 维数为 2 的反例证明的.这里略去证明,继续前进.

在 Fefferman 的反例出现以后,兴趣的对象转移到证明比 $S_{1}$ 奇异性更小的 Bochner--Riesz 平均
\[
S_{1}^{\lambda}f(x) := \int_{\mathbb{R}^{n}}e^{2\pi i x\cdot\xi}(1-|\xi|^{2})_{+}^{\lambda}\hat{f}(\xi)\dd\xi, \quad t_{+}^{\lambda} := \max(t,0)^{\lambda} \quad (\lambda > 0)
\]
的 $L^{p}(\mathbb{R}^{n})$ 有界性问题.由驻相法可以证明 $S_{1}^{\lambda}$ 的核的绝对值在远方以 $|x|^{-(n+1)/2-\lambda}$ 衰减,并且猜想在
\[
p \in \left[\frac{2n}{n+1}, \frac{2n}{n-1}\right], \quad \lambda > 0
\]
时 $S_{1}^{\lambda}$ 在 $L^{p}(\mathbb{R}^{n})$ 上有界.

更精确地,已经得到 Bochner--Riesz 平均在 $L^{p}(\mathbb{R}^{n})$ 上有界的必要条件.即,令
\[
\lambda(p) := \max\left\{n\left|\frac{1}{2} - \frac{1}{p}\right| - \frac{1}{2}, 0\right\},
\]
则在 $p \in [1, \infty]$ 时 $S_{1}^{\lambda}$ 在 $L^{p}(\mathbb{R}^{n})$ 上有界的必要条件为
\[
\lambda > \lambda(p).
\]
关于 Bochner--Riesz 平均的猜想是:这个必要条件也是充分条件.下面确认该猜想在二维时正确.在三维以上,只知道部分结果,仍悬而未决.

\textbf{C\'ordoba 对二维情形的解法.} 1977 年 C\'ordoba 用二维 Nikodym 极大函数的 $L^{2}(\mathbb{R}^{2})$ 估计,给出了 Bochner--Riesz 平均猜想在二维成立的另一证明.该问题最初的证明出自 Carleson 与 Sj\"olin,用的是驻相法.

这里我们概观一下 C\'ordoba 的证明.

以下,把由函数 $\phi$ 给出的 Fourier 乘子算子记为 $D(\phi)$.即
\[
D(\phi)f(x) := \int_{\mathbb{R}^{n}}e^{2\pi i x\cdot\xi}\phi(\xi)\hat{f}(\xi)\dd\xi.
\]
在 $p \in [4/3, 4]$,$\lambda > 0$ 时,目标是确认
\[
\|S_{1}^{\lambda}f\|_{L^{p}(\mathbb{R}^{2})} \leq C\|f\|_{L^{p}(\mathbb{R}^{2})}.
\]
考虑到对偶性以及与由 Plancherel 定理推出的 $L^{2}(\mathbb{R}^{2})$ 有界性之间的插值,只需证明
\begin{equation}\tag{4.9.2}
\|S_{1}^{\lambda}f\|_{L^{4}(\mathbb{R}^{2})} \leq C\|f\|_{L^{4}(\mathbb{R}^{2})}
\label{eq:4.9.2}\end{equation}
即可.

设 $\epsilon \in (0, 1/2]$.取满足
\[
\operatorname{supp}\phi_{\epsilon} \subset (1-\epsilon, 1-\epsilon/2), \quad |(d^{k}/dt^{k})\phi_{\epsilon}| \leq C_{k}\epsilon^{-k} \quad (k = 1, 2, \dots)
\]
的光滑实值函数 $\phi_{\epsilon}$,定义
\[
M_{\epsilon} := D\left(\phi_{\epsilon}(|\xi|)\right).
\]
此时,用单位分解把 $(1-|\xi|^{2})_{+}^{\lambda}$ 的支集作二进制分解,可得展开
\[
S_{1}^{\lambda}f(x) = M_{0}f(x) + \sum_{j=1}^{\infty}2^{-\lambda j}M_{2^{-j}}f(x).
\]
这里,$M_{0}$ 是以 $B(0, 1/2)$ 为支集的光滑函数给出的 Fourier 乘子算子,在 $L^{4}(\mathbb{R}^{2})$ 上有界.

现在,若对充分小的 $\epsilon > 0$ 一致地得到估计
\begin{equation}\tag{4.9.3}
\|M_{\epsilon}f\|_{L^{4}(\mathbb{R}^{2})} \leq C\epsilon^{-\lambda/2}\|f\|_{L^{4}(\mathbb{R}^{2})},
\label{eq:4.9.3}\end{equation}
则由上面的展开式显然可得\eqref{eq:4.9.2}.把 $\phi_{\epsilon}(|\xi|)$ 的支集
\[
\{\xi \in \mathbb{R}^{2}:\ 1-\epsilon < |\xi| < 1-\epsilon/2\}
\]
在角方向用长度 $\sim\sqrt{\epsilon}$ 的弯曲长方形 $\tilde{R}_{\epsilon,i}$($i = 1, 2, \ldots$)分解,把 $M_{\epsilon}$ 进一步展开为
\[
M_{\epsilon} = \sum_{i=1}^{\sim\epsilon^{-1/2}}D(\psi_{\epsilon,i}).
\]
这里,$\psi_{\epsilon,i}$ 是以 $\tilde{R}_{\epsilon,i}$ 为支集的光滑函数.

把通过各 $\tilde{R}_{\epsilon,i}$ 中心的单位向量记为 $\omega_{i} \in S^{1}$,不失一般性,以下假定 $\omega_{i}$ 的第一坐标至少为 $1/\sqrt{2}$ 来进行证明.另外,把这样的 $i$ 全体的集合记为 $I$,为简单起见虽稍显粗糙,但设
\[
M_{\epsilon} = \sum_{i \in I}D(\psi_{\epsilon,i}).
\]

由 Fourier 变换的基本公式,
\begin{equation}\tag{4.9.4}
\begin{aligned}
\|M_{\epsilon}f\|_{4}^{4} & = \int_{\mathbb{R}^{2}}\left|\sum_{i \in I}D(\psi_{\epsilon,i})f\right|^{4}\dd x \\
& = \int_{\mathbb{R}^{2}}\left|\sum_{i,j \in I}D(\psi_{\epsilon,i})f\,D(\psi_{\epsilon,j})f\right|^{2}\dd x \\
& = \int_{\mathbb{R}^{2}}\left|\sum_{i,j \in I}\left(D(\psi_{\epsilon,i})f\,D(\psi_{\epsilon,j})f\right)^{\wedge}\right|^{2}\dd\xi \\
& = \int_{\mathbb{R}^{2}}\left|\sum_{i,j \in I}\left(D(\psi_{\epsilon,i})f\right)^{\wedge}*\left(D(\psi_{\epsilon,j})f\right)^{\wedge}\right|^{2}\dd\xi.
\end{aligned}
\label{eq:4.9.4}\end{equation}
为估计该右端,准备如下几何引理.

\begin{lemma}\label{lem:4.9.2}
以 $c$ 为不依赖于 $\epsilon$ 的常数,
\[
\left\|\sum_{i,j \in I}\chi_{\tilde{R}_{\epsilon,i}} * \chi_{\tilde{R}_{\epsilon,j}}\right\|_{\infty} \leq c
\]
成立.
\end{lemma}

\begin{proof}
该引理本质上由如下事实推出:若 $\omega_{1}, \omega_{2}, \omega'_{1}, \omega'_{2} \in S^{1}$ 满足
\begin{align*}
0 &\neq \omega_{1} + \omega_{2} = \omega'_{1} + \omega'_{2}, \\
\intertext{则要么 $\omega_{1} = \omega'_{1}, \omega_{2} = \omega'_{2}$,要么 $\omega_{1} = \omega'_{2}, \omega_{2} = \omega'_{1}$. 设}
g(\xi) &:= \sum_{i,j \in I}\left(\chi_{\tilde{R}_{\epsilon,i}} * \chi_{\tilde{R}_{\epsilon,j}}\right)(\xi).
\end{align*}
若取 $g$ 的支集上的点为 $\zeta_{0}$,则存在 $\xi_{0} \in \tilde{R}_{\epsilon,i_{0}}$,$\eta_{0} \in \tilde{R}_{\epsilon,j_{0}}$ 使得
\[
\zeta_{0} = \xi_{0} + \eta_{0}.
\]
现在设存在 $\xi_{1} \in \tilde{R}_{\epsilon,i_{1}}$,$\eta_{1} \in \tilde{R}_{\epsilon,j_{1}}$ 使得
\begin{align*}
\zeta_{0} &= \xi_{1} + \eta_{1}. \\
\intertext{令}
\zeta'_{0} &:= \frac{\xi_{0}}{|\xi_{0}|} + \frac{\eta_{0}}{|\eta_{0}|}, \quad \zeta'_{1} := \frac{\xi_{1}}{|\xi_{1}|} + \frac{\eta_{1}}{|\eta_{1}|}, \\
\intertext{则}
&\begin{array}{r l}
|\zeta'_{0} - \zeta'_{1}| & \leq |\zeta'_{0} - \zeta_{0}| + |\zeta_{0} - \zeta'_{1}| \\
& = \left|\frac{\xi_{0}}{|\xi_{0}|} - \xi_{0}\right| + \left|\frac{\eta_{0}}{|\eta_{0}|} - \eta_{0}\right| + \left|\frac{\xi_{1}}{|\xi_{1}|} - \xi_{1}\right| + \left|\frac{\eta_{1}}{|\eta_{1}|} - \eta_{1}\right| \\
& \leq 4\epsilon.
\end{array}
\end{align*}
由此,满足上述条件的 $\xi_{1}, \eta_{1}$ 的个数不多是显然的.特别地,对 $\zeta'_{0} \sim (2, 0)$ 的情形,用 $1 - \cos\theta \sim \theta^{2}/2$($\theta \sim 0$)的关系考察,可知其个数不依赖于 $\epsilon$.
\end{proof}

用该引理估计\eqref{eq:4.9.4},得
\begin{equation}\tag{4.9.5}
\|M_{\epsilon}f\|_{4}^{4} \leq C\int_{\mathbb{R}^{2}}\left|\sum_{i \in I}|D(\psi_{\epsilon,i})f|^{2}\right|^{2}\dd x
\label{eq:4.9.5}\end{equation}
成立.

来估计 $D(\psi_{\epsilon,i})f(x)$.改写
\[
D(\psi_{\epsilon,i})f(x) = \iint e^{2\pi i(x-y)\cdot\xi}\psi_{\epsilon,i}(\xi)\dd\xi f(y)\dd y =: \int K_{i}(x-y)f(y)\dd y,
\]
来估计
\begin{equation}\tag{4.9.6}
K_{i}(x) = \int e^{2\pi i x\cdot\xi}\psi_{\epsilon,i}(\xi)\dd\xi.
\label{eq:4.9.6}\end{equation}

对每个 $\omega_{i}$,各选取一个满足 $U_{i}\omega_{i} = (1, 0)$ 的 $\mathbb{R}^{2}$ 的直交变换 $U_{i}$,令
\begin{align*}
R_{\epsilon} &:= \left\{x \in \mathbb{R}^{2}:\ |x_{1}| \leq \frac{1}{\epsilon},\ |x_{2}| \leq \frac{1}{\sqrt{\epsilon}}\right\}, \\
\intertext{设}
R_{\epsilon,i} &:= U_{i}^{-1}R_{\epsilon}.
\end{align*}
此时,可以证明只需证明
\begin{equation}\tag{4.9.7}
|K_{i}(x)| \leq C\frac{\chi_{R_{\epsilon,i}}(x)}{|R_{\epsilon,i}|}
\label{eq:4.9.7}\end{equation}
即可.这一点留到最后确认,先继续前进.

对每个 $\tilde{R}_{\epsilon,i}$($i \in I$),选取满足 $\tilde{R}_{\epsilon,i} \subset Q_{\epsilon,i}$ 的边长为 $\sim\sqrt{\epsilon}$ 的正方形 $Q_{\epsilon,i}$,令
\begin{align*}
f_{i} &:= D\left(\chi_{Q_{\epsilon,i}}\right)f, \\
\intertext{则显然}
D(\psi_{\epsilon,i})f &= D(\psi_{\epsilon,i})f_{i}.
\end{align*}
由该式与\eqref{eq:4.9.5},\eqref{eq:4.9.7}得
\begin{equation}\tag{4.9.8}
\|M_{\epsilon}f\|_{4}^{2} \leq C\left\|\sum_{i \in I}\left(\frac{\chi_{R_{\epsilon,i}}}{|R_{\epsilon,i}|} * |f_{i}|\right)^{2}\right\|_{2}.
\label{eq:4.9.8}\end{equation}
用对偶性来解析该右端.

任取非负函数 $w \in L^{2}(\mathbb{R}^{2})$,$\|w\|_{2} = 1$.利用 $\left\|\frac{\chi_{R_{\epsilon,i}}}{|R_{\epsilon,i}|}\right\|_{1} = 1$,
\begin{align*}
&\begin{array}{r l}
\int\left(\frac{\chi_{R_{\epsilon,i}}}{|R_{\epsilon,i}|} * |f_{i}|\right)^{2}w\dd x & \leq \int\left(\frac{\chi_{R_{\epsilon,i}}}{|R_{\epsilon,i}|} * |f_{i}|^{2}\right)w\dd x \\
& = \int|f_{i}|^{2}\left(\frac{\chi_{R_{\epsilon,i}}}{|R_{\epsilon,i}|} * w\right)\dd x.
\end{array} \\
\intertext{用此式,}
&\begin{array}{r l}
\int\sum_{i \in I}\left(\frac{\chi_{R_{\epsilon,i}}}{|R_{\epsilon,i}|} * |f_{i}|\right)^{2}w\dd x & \leq \int\left(\sum_{i \in I}|f_{i}|^{2}\right)\left(\sup_{i \in I}\frac{\chi_{R_{\epsilon,i}}}{|R_{\epsilon,i}|} * w\right)\dd x \\
& \leq \int\left(\sum_{i \in I}|f_{i}|^{2}\right)K_{N,a}w\dd x \\
& \leq \left\|\sum_{i}|f_{i}|^{2}\right\|_{2}\cdot\|K_{N,a}w\|_{2}.
\end{array}
\end{align*}
这里,取 $N = a = \frac{1}{\sqrt{\epsilon}}$,$K_{N,a}w$ 表示 4.7 节定义的 Nikodym 极大函数.

这里承认(不加证明地使用)如下事实:
\begin{align*}
\left\|\sum_{i}|f_{i}|^{2}\right\|_{2} &\leq C\|f\|_{4}^{2}, \\
\intertext{因此问题归结为 Nikodym 极大函数 $K_{N,a}w$ 的 $L^{2}(\mathbb{R}^{2})$ 估计.正如在 4.7 节所确认的,它可估计为}
\|K_{N,a}w\|_{2} &\leq C(\log N)^{1/2}\|w\|_{2} = C(\log N)^{1/2},
\end{align*}
由\eqref{eq:4.9.8}得
\[
\|M_{\epsilon}f\|_{4} \leq C(\log(1/\epsilon))^{1/4}\|f\|_{4}.
\]
由此\eqref{eq:4.9.3}成立.

\textbf{验证\eqref{eq:4.9.7}.} 设 $\psi_{0} \in C_{0}^{\infty}(\mathbb{R}^{2})$ 满足
\begin{align*}
\operatorname{supp}\psi_{0} &\subset \left\{\xi = (\xi_{1}, \xi_{2}) \in \mathbb{R}^{2}:\ 1-\epsilon < |\xi| < 1-\epsilon/2,\ -\sqrt{\epsilon} < \xi_{2} < \sqrt{\epsilon}\right\}, \\
\intertext{且}
\left|(\partial/\partial\xi_{1})^{\alpha_{1}}(\partial/\partial\xi_{2})^{\alpha_{2}}\psi_{0}\right| &\leq C_{\alpha_{1},\alpha_{2}}\epsilon^{-\alpha_{1}}\epsilon^{-\alpha_{2}/2}. \\
\intertext{不失一般性,可以设}
\psi_{\epsilon,i}(\xi) &= \psi_{0}(U_{i}\xi). \\
\intertext{由\eqref{eq:4.9.6},}
&\begin{array}{r l}
K_{i}(x) & = \int e^{2\pi i x\cdot\xi}\psi_{\epsilon,i}(\xi)\dd\xi \\
& = \int e^{2\pi i U_{i}^{-1}x\cdot U_{i}\xi}\psi_{0}(U_{i}\xi)\dd\xi \\
& = \int e^{2\pi i U_{i}^{-1}x\cdot\xi}\psi_{0}(\xi)\dd\xi.
\end{array} \\
\intertext{因此,为证明主张,只需考察}
K_{0}(x) &:= \int e^{2\pi i x\cdot\xi}\psi_{0}(\xi)\dd\xi. \\
\intertext{利用对 $\psi_{0}$ 的假设作分部积分,可知}
(\epsilon|x_{1}|)^{m}(\sqrt{\epsilon}|x_{2}|)^{m}|K_{0}(x)| &\leq C_{m}\epsilon^{3/2} \quad (m = 0, 1, \dots)
\end{align*}
成立.由此计算出
\[
|K_{0}(x)| \leq C_{m}\sum_{k=0}^{\infty}2^{-(m-2)k}\frac{\chi_{2^{k}R_{\epsilon}}}{|2^{k}R_{\epsilon}|},
\]
由 Minkowski 不等式使主张得到验证.


% ============ 第4章 参考文献(原书 pp.378-380) ============
\section*{参考文献}\addcontentsline{toc}{section}{参考文献}
\begin{thebibliography}{99}
% 条目 [13],[14] 为日文原条目,其中个别汉字(如 掛,ー 等)为 Fandol 字体所无,
% 此处临时切换至系统日文字体渲染(条目前后不改变).
\newCJKfontfamily\japonfont{MS Gothic}
\begin{sloppypar}
\item[] 关于用驻相法处理振荡积分算子以及 Fourier 限制问题(Tomas--Stein 定理)的部分,依照 \cite{ch4:1} 的最初部分整理而成.作为这一领域的基本教材,也列出 \cite{ch4:2}.
\bibitem[1]{ch4:1} C. D. Sogge, Fourier Integrals in Classical Analysis, Cambridge Tracts in Math. 105 (1993).
\bibitem[2]{ch4:2} E. M. Stein, Harmonic Analysis: Real-Variable Methods, Orthogonality, and Oscillatory Integrals, Princeton University Press, 1993.
\item[] 关于 Nikodym 极大函数(Wolff 定理)的部分,基本上依据 Wolff 的原始论文 \cite{ch4:3},但按照将其整理改编后的 \cite{ch4:4} 编写.
\bibitem[3]{ch4:3} T. Wolff, An improved bound for Kakeya type maximal functions, Rev. Mat. Iberoamericana 11 (1995), 651--674.
\bibitem[4]{ch4:4} H. Tanaka, The Fefferman-Stein type inequality for the Kakeya maximal operator in Wolff's range, Proc. Amer. Math. Soc. 133 (2005), 763--772.
\item[] 作为与 Bochner--Riesz 平均和 Nikodym 极大函数部分相关的论文,列出如下.
\bibitem[5]{ch4:5} A. C\'ordoba, The Kakeya maximal function and the spherical summation multiplier, Amer. J. Math. 99 (1977), 1-22.
\bibitem[6]{ch4:6} A. C\'ordoba, A note on Bochner-Riesz operators, Duke Math. J. 46 (1979), 505--511.
\bibitem[7]{ch4:7} J. Bourgain, Besicovitch type maximal operators and applications to Fourier analysis, Geom. Funct. Anal. 1 (1991), no. 2, 147--187.
\item[] 作为关于 Fourier 限制问题的综述,列出如下.
\bibitem[8]{ch4:8} T. Tao, Recent progress on the Restriction conjecture, Fourier analysis and convexity, 217--243, Appl. Numer. Harmon. Anal., 2004.
\item[] 作为关于 Kakeya 问题的综述,列出如下.
\bibitem[9]{ch4:9} J. Bourgain, Harmonic analysis and combinatorics: How much may they contribute to each other?, Mathematics: Frontiers and perspectives, IMU/Amer. Math. Society, 2000, 13--32.
\bibitem[10]{ch4:10} T. Wolff, Recent work connected with the Kakeya problem, Prospects in mathematics (Princeton, NJ, 1996), 129--162, Amer. Math. Soc., Providence, RI, 1999.
\bibitem[11]{ch4:11} T. Tao, From rotating needles to stability of waves: emerging connections between combinatorics, analysis, and PDE, Notices Amer. Math. Soc. 48 (2001), no. 3, 294--303.
\bibitem[12]{ch4:12} N. H. Katz and T. Tao, Recent progress on the Kakeya conjecture, Publ. Mat. 2002, Vol. Extra, 161--179.
\bibitem[13]{ch4:13} {\japonfont 田中仁, 掛谷問題について,} 数学 57 (2005), no. 2, 113--129.
\item[] 作为论述 Kakeya 集构造的日语教材,列出如下.
\bibitem[14]{ch4:14} {\japonfont 新井仁之, ルベーグ積分講義, 日本評論社,} 2003.
\item[] 作为用组合论方法研究 Nikodym 极大函数所得的结果,列出如下论文.
\bibitem[15]{ch4:15} J. Bourgain, On the dimension of Kakeya sets and related maximal inequalities, Geom. Funct. Anal. 9 (1999), no. 2, 256--282.
\bibitem[16]{ch4:16} N. H. Katz and T. Tao, Bounds on arithmetic projections, and applications to the Kakeya conjecture, Math. Res. Lett. 6 (1999), 625--630.
\bibitem[17]{ch4:17} N. H. Katz and T. Tao, New bounds for Kakeya problems, J. Anal. Math. 87 (2002), 221--263.
\item[] 作为用双线性方法研究 Fourier 限制问题以及 Bochner--Riesz 平均问题所得的结果,列出如下论文.
\bibitem[18]{ch4:18} T. Tao, A. Vargas and L. Vega, A bilinear approach to the restriction and Kakeya conjectures, J. Amer. Math. Soc. 11 (1998), 967--1000.
\bibitem[19]{ch4:19} T. Wolff, A sharp bilinear cone restriction estimate, Annals of Math. 153 (2001), 661-698.
\bibitem[20]{ch4:20} T. Tao, Endpoint bilinear restriction theorems for the cone, and some sharp null form estimates, Math. Z. 238 (2001), no. 2, 215-268.
\bibitem[21]{ch4:21} T. Tao, A sharp bilinear restriction estimate on paraboloids, Geom. Funct. Anal. 13 (2003), 1359--1384.
\bibitem[22]{ch4:22} S. Lee, Improved bounds for Bochner-Riesz and maximal Bochner-Riesz operators, Duke Math. J. 122 (2004), no. 1, 205-232.
\bibitem[23]{ch4:23} S. Lee, Bilinear restriction estimates for surfaces with curvatures of different signs. Trans. Amer. Math. Soc. 358 (2006), no. 8, 3511-3533.
\end{sloppypar}
\end{thebibliography}


% 索引
\printindex

\end{document}
